\documentclass[11pt,oneside]{article}
\usepackage[margin=25mm,headheight=15pt]{geometry}
\usepackage[T1]{fontenc}
\usepackage{lmodern}
\usepackage{amsmath,amssymb,amsthm,mathtools,mathrsfs}
\usepackage{booktabs,longtable,array,enumitem}
\usepackage{xcolor,xurl,hyperref,fancyhdr,microtype}
\hypersetup{colorlinks=true,linkcolor=blue!45!black,citecolor=blue!45!black,urlcolor=blue!45!black,
pdftitle={Points as sections: incomplete research return},
pdfauthor={GPT-6 Astra (OpenAI, ChatGPT Pro mode)}}
\setlength{\emergencystretch}{2em}
\setlength{\parskip}{3pt}
\setlength{\parindent}{1.2em}
\setlist{itemsep=2pt,topsep=4pt}
\allowdisplaybreaks
\newcommand{\Sha}{\operatorname{Sha}}
\newcommand{\Q}{\mathbb Q}
\newcommand{\Z}{\mathbb Z}
\newcommand{\R}{\mathbb R}
\newcommand{\C}{\mathbb C}
\newcommand{\F}{\mathbb F}
\DeclareMathOperator{\Aut}{Aut}
\DeclareMathOperator{\Out}{Out}
\DeclareMathOperator{\Hom}{Hom}
\DeclareMathOperator{\Gal}{Gal}
\DeclareMathOperator{\Spec}{Spec}
\DeclareMathOperator{\Pic}{Pic}
\DeclareMathOperator{\Jac}{Jac}
\DeclareMathOperator{\GL}{GL}
\DeclareMathOperator{\SL}{SL}
\DeclareMathOperator{\im}{im}
\DeclareMathOperator{\coker}{coker}
\DeclareMathOperator{\rank}{rank}
\DeclareMathOperator{\tr}{tr}
\DeclareMathOperator{\End}{End}
\DeclareMathOperator{\Inn}{Inn}
\DeclareMathOperator{\Sect}{Sect}
\newcommand{\enquote}[1]{\textquotedblleft#1\textquotedblright}
\newcounter{statementanchor}
\makeatletter
\newenvironment{statement}[3]{%
 \refstepcounter{statementanchor}\def\@currentlabel{#1}%
 \par\addvspace{9pt}\noindent
 \textbf{#1. #2. #3.}\par\nobreak\noindent\ignorespaces
}{\par\addvspace{5pt}}
\makeatother
\pagestyle{fancy}
\fancyhf{}
\fancyhead[L]{\small Points as sections}
\fancyhead[R]{\small Research return --- INCOMPLETE}
\fancyfoot[C]{\thepage}
\begin{document}
\begin{center}
{\LARGE\bfseries Points as sections}\\[4pt]
{\large The anabelian texts of 1981 and 1983, real structures,\\
Mordell, and the Fermat tower}\\[7pt]
AI-produced by GPT-6 Astra\\
OpenAI, ChatGPT Pro mode\\[4pt]
14 November 2026 \qquad \textbf{INCOMPLETE research return}
\end{center}
\noindent
We give proofs, explicit comparison maps, counterexamples, and
independently executed finite certificates for the requested manuscript
passages. The unresolved statements remain numbered, with the precise
obstruction and the attempted approaches. We claim no priority.
Results using the supplied unrefereed local section theorem are
identified as deductions from that stipulated input.

\smallskip
\noindent\textbf{Notation and manuscript concordance.}
Symbols are fixed locally when a source uses the same letter differently.
The following table governs the arithmetic curve comparisons.
\begin{center}\small
\begin{tabular}{@{}p{.17\textwidth}p{.39\textwidth}p{.36\textwidth}@{}}
\toprule
Our notation & Mathematical object & Manuscript notation and introduction\\
\midrule
$G_K,\Gamma$ & Absolute Galois group, with profinite topology &
$\Gamma$; \S5, scans 21--24; cote 149, scans 43--48\\
$U=X\setminus S$ & Smooth hyperbolic curve and its proper completion &
$U,X,S$; \S6, scans 25--31\\
$E_U,\Pi_U$ & Arithmetic fundamental group &
$\mathcal E,E$; \S7, scans 43--54; cote 149, scan 44\\
$\pi,\Delta_U$ & Geometric fundamental group &
$\pi,\Pi_{U/K}$; same passages\\
$I_c,L_c,D_c$ & Cusp inertia and arithmetic decomposition &
$L_i,\mathcal E_i,N(L_i)$; scans 28--35, 69, 75--80\\
$T,\chi,\Gamma^\circ$ & $\widehat{\mathbb Z}(1)$, cyclotomic character, its kernel &
$T,\chi,\Gamma^\circ$; cote 149, scans 45, 54--56\\
$\operatorname{Sec}(U/K)$ & Continuous sections modulo geometric conjugacy &
Classes of splittings; scans 61--65, 75--79\\
\bottomrule
\end{tabular}
\end{center}
\noindent
A field embedding and a base-point path are specified wherever they
change a genuine homomorphism. A quotient map always specifies the
induced map of section classes. All closures are taken in the named
profinite or prime-adic topology.
For a multiplicative group, the Kummer completion means
$\varprojlim_n K^*/K^{*n}$ with its inverse-limit topology.
For a number field this group need not be compact.
Every finite quotient of $K^*$ has an exponent $n$ and therefore
factors through $K^*/K^{*n}$; these maps give the comparison
from the Kummer completion to the completion over all finite
quotients. It is injective, because finite quotients separate
every nonzero class in $K^*/K^{*n}$, and it has dense image,
because its image contains $K^*$.
\newpage
\tableofcontents
\newpage
% ---- A ----
\section{The anabelian conjectures and their exact groupoid form}
% Part A.  Included below the root document's lettered section heading.
% Every numbered environment has a status in its title.
\subsection{The typed groupoid construction and a correction on scan 51}

In this part $K$ is a field of characteristic zero, $\overline K$ is a
chosen separable closure, $\Gamma=G_K$, and
$T=\varprojlim_n\mu_n(\overline K)=\widehat{\mathbf Z}(1)$.
For a smooth proper geometrically connected curve $X/K$ and a reduced
finite divisor $S$, put $U=X\setminus S$, $E=\pi_1(U)$ and
$\pi=\pi_1(U_{\overline K})$.  A choice of geometric base point is
understood; its change acts by geometric inner conjugacy.  A loop group
$L_{\widetilde s}$ is the inertia at a lift of a geometric cusp to the
completed universal covering.  Its arithmetic decomposition group is
$E_{\widetilde s}$.  We use $\operatorname{Int}(a)(b)=aba^{-1}$.
The letter $\Delta$ in the CM calculation below denotes
$G_{\mathbf Q(i)^{\rm ab}}$; it never denotes the geometric group in
that calculation.  Grothendieck's scans 43--54 use $\mathcal E,\pi,
\Gamma,\mathcal E_i,L_i,T$ for the respective objects.

\begin{statement}{A.1}{Proved}{the two functors and the exact boundary diagram}
\label{st:A.1}
\label{A:typed-functor}
The literal functor of scan 23 sends a smooth projective anabelian
$K$-scheme of dimension at most one to
$\Pi_1(X)\longrightarrow\Pi_1(K)$.  Its pointed version has morphisms
$(f,\gamma):(X',\xi')\to(X,\xi)$, where $f$ is a $K$-morphism and
$\gamma$ is an \'etale path from $f(\xi')$ to $\xi$.  The path supplies
the actual homomorphism of based groups.  Passing to paths modulo target
geometric conjugacy recovers the unpointed groupoid formulation.

For a pair $(X,S)$ the boundary enhancement of scans 28--35 is the
following diagram of profinite groupoids, with its displayed
commutation isomorphism and its Kummer isomorphism:
\[
 \Pi_{D^*}\xrightarrow{\rho}\Pi_U\xrightarrow{\varphi}\Pi_e,
 \qquad
 \Pi_{D^*}\xrightarrow{\sigma}\Pi_D\xrightarrow{\psi}\Pi_e,
 \qquad \alpha:\varphi\rho\simeq\psi\sigma,
 \qquad \ker(\sigma)\simeq\psi^*T .
\]
Here $D_s=\operatorname{Spec}\widehat{\mathcal O}_{X,s}$,
$D_s^*=D_s\setminus\{s\}$, $D=\coprod_{s\in S}D_s$ and
$D^*=\coprod_{s\in S}D_s^*$.  The map $\sigma$ induces a bijection
on components and surjections on automorphism groups, its kernel is
abelian, $\psi$ is faithful, and $\rho$ is faithful for hyperbolic
curves.  The images in $\Pi_e$ are open.

After choosing one object in each component, the same data are
$p:E\to\Gamma$, its kernel $\pi$, a finite index set $I$, injections
$\rho_i:E_i\hookrightarrow E$, open images
$\Gamma_i=p(E_i)$, and equivariant isomorphisms
$k_i:T\xrightarrow{\sim}L_i=E_i\cap\pi$.
The image of $p$ is permitted to be open rather than all of $\Gamma$.
Disjoint unions and the separate constant fields recover the
non-geometrically-connected version.
\end{statement}
\begin{proof}
The \'etale fibre functors on a scheme and their isomorphisms form its
profinite fundamental groupoid.  Pullback of finite \'etale covers
gives the stated covariant functor on fundamental groupoids.  In the
pointed construction the specified path identifies the target fibre
functor, which gives the asserted actual homomorphism.  Any two choices
of that path differ by an element of the target geometric group; this
is exactly the displayed conjugacy relation.

The complete local ring at $s$ is a complete discrete valuation ring
with residue field $K(s)$.  In characteristic zero the geometric
punctured local field has tame inertia $T$, by the covers
$t^{1/n}$.  Its conjugation action is the cyclotomic action.  This
proves the local kernel and equivariance statements, including the
nonrational case with quotient $G_{K(s)}$.  The inclusion of local
inertia and decomposition in $\pi_1(U)$ is injective for a hyperbolic
curve: over an algebraic closure this is the injection of a peripheral
profinite cyclic group in the surface group; the arithmetic assertion
then follows by the exact sequences and the injection of
$G_{K(s)}$ in $G_K$.  Choosing objects and paths in a connected
groupoid identifies it with the one-object groupoid of its
automorphism group.  Applying that equivalence simultaneously to the
diagram gives precisely the listed group data.  Reversing this
construction gives the original diagram up to equivalence.
\end{proof}

% Additional case analysis is placed here so that every case of section 8
% has an explicit statement, independently of the subsequent normalizer work.
\subsection{The individual cases of section 8}

\begin{statement}{A.2}{Proved}{proper source and affine target}
\label{st:A.2}
\label{A:proper-affine}
Let $K$ be a field finitely generated over $\mathbf Q$, or a finite
extension of $\mathbf Q_p$. Let $X'/K$ be a smooth proper
geometrically connected curve and $U/K$ an affine hyperbolic curve.
Every continuous homomorphism of extensions
\[
 f:E_{X'}\longrightarrow E_U,\qquad p_Uf=p_{X'},
\]
is trivial on $P_{X'}=\pi_1(X'_{\overline K})$. Consequently it
factors uniquely as $E_{X'}\to G_K\xrightarrow{s}E_U$, where
$s$ is a continuous section. This proves the group-theoretic assertion
of case IV(b), scan 66, manuscript folio 31. The input used to exclude
an infinite-index noncuspidal image is exactly the local section
theorem stipulated in the prompt, with its stated unrefereed status.
\end{statement}
\begin{proof}
First let $K=k$ be a finite extension of $\mathbf Q_p$. The actual
local construction in Proposition~\ref{F-local-dichotomy} gives
three possibilities for $H=f(P_{X'})$: it is open in $P_U$, it is
trivial, or it lies nontrivially in one cusp inertia $I_c$.
This construction uses the supplied point-or-cusp theorem over
finite extensions of $k$, Carlson--Stix's point-theoretic
Hom theorem, and compactifications of arbitrarily high genus.
In particular it does not require an arithmetic map to preserve
point sections before that local theorem has been applied.

The open possibility is impossible. Mochizuki's Theorem A, also
in its formulation with an arbitrary smooth source, would give a
dominant morphism $X'\to U$. Every regular function on the
geometrically connected proper curve $X'$ is constant, so a
morphism to the affine curve $U$ is constant. This is incompatible
with dominance. Equivalently, for genus zero or one sources one
can first puncture $X'$ to a hyperbolic curve, apply Theorem A,
and use triviality on the newly introduced inertia to extend the
result across those punctures.

Suppose instead that $1\ne H\subset I_c$. Normality of $H$ in
$f(E_{X'})$ and the inertia-normalizer identity put the whole
image in $D_c$. Since it projects onto $G_k$, the cusp is
$k$-rational. The resulting geometric map therefore factors as
the following $G_k$-equivariant continuous map of abelian groups:
\[
 P_{X'}\longrightarrow P_{X'}^{\rm ab}
   =\prod_\ell T_\ell J_{X'}
 \longrightarrow I_c=\widehat{\mathbf Z}(1).
\]
For each prime $\ell$, the final coordinate vanishes. Indeed the
Weil pairing identifies
\[
 \operatorname{Hom}_{G_k}(T_\ell J_{X'},\mathbf Z_\ell(1))
       =(T_\ell J_{X'}^\vee)^{G_k}=0.
\]
The last equality follows from finiteness of the rational torsion
of an abelian variety over a $p$-adic local field: the inverse
limit, under multiplication by $\ell$, of its bounded groups
of rational $\ell$-power torsion is zero. A homomorphism from a
pro-$q$ group to a pro-$\ell$ group is zero for $q\ne\ell$;
hence these coordinate calculations kill the entire displayed
map. This contradicts $H\ne1$. Thus $H=1$, and the quotient
property of $E_{X'}\to G_k$ gives the asserted factorization.

For $K$ finitely generated over $\mathbf Q$, use the embedding
$K\hookrightarrow k$ into a finite $p$-adic extension constructed
in Lemma~\ref{A:local-embedding}. Base change gives the cartesian
maps
\[
 E_{X'_k}=E_{X'}\times_{G_K}G_k,
 \qquad E_{U_k}=E_U\times_{G_K}G_k.
\]
Their geometric groups identify with the original geometric
groups, since geometric fundamental groups of smooth curves in
characteristic zero are invariant under algebraically closed
field extension. The induced $f_k$ thus has exactly the same
geometric map as $f$. The local calculation makes that map zero.
No injectivity or openness of the homomorphism $G_k\to G_K$
is used. This establishes the stated field range by an explicit
base-change morphism, without an additional specialization claim.
\end{proof}

\begin{statement}{A.3}{Proved}{cases I, II, IV(a), IV(c), IV(d), and the exact point-source remainder}
\label{st:A.3}
\label{A:section-eight-cases}
In the category constructed in Proposition~\ref{A:typed-functor},
the cases of scans 55--67 have the following precise reductions.
Cases I and II are finite Galois theory; case III is the point
section map; case IV(b) is Proposition~\ref{A:proper-affine};
case IV(c) has empty Hom set in the literal pair category;
and case IV(d) is the dominant relative Hom bijection of
Proposition~\ref{A:relative-hom}. For case IV(a), over a field
finitely generated over $\mathbf Q$ or a finite $p$-adic
extension, the groupoid Hom set is the disjoint union of the
dominant-morphism classes and the section classes of the target.
The section classes give the remaining constant maps exactly
when those sections are point-induced. Over finite $p$-adic
extensions and proper hyperbolic targets this last assertion
follows from the stipulated local theorem; over number fields it
retains precisely the global point-section surjectivity problem.
\end{statement}
\begin{proof}
In case I, write $X=\operatorname{Spec}L$ and
$X'=\operatorname{Spec}L'$ with the chosen finite subextensions
of $\overline K/K$. A pair $(f,\alpha)$ in scans 55--56 has
$f=\operatorname{Int}(\alpha)|_{G_{L'}}$ and exists exactly when
$\alpha G_{L'}\alpha^{-1}\subset G_L$. The quotient by its
2-morphisms is the actual transporter quotient
\[
 G_L\backslash
 \{\alpha\in G_K:\alpha G_{L'}\alpha^{-1}\subset G_L\}
    \longrightarrow \operatorname{Hom}_K(L,L'),
 \qquad [\alpha]\longmapsto\alpha^{-1}|_L .
\]
It is bijective by finite Galois theory. The inverse assigns an
extension of the given embedding to $\overline K$; any two such
extensions differ by the displayed left $G_L$-action.

In case II, let $L'=H^0(X',\mathcal O_{X'})$ for the complete
source. Its arithmetic fundamental group projects onto $G_{L'}$.
The target projection $G_L\hookrightarrow G_K$ is injective,
so the same equation determines $f$ from $\alpha$, kills the
geometric subgroup, and gives the same transporter quotient.
The actual morphism is the composite
$X'\to\operatorname{Spec}L'\to\operatorname{Spec}L$.
The literal pair condition requires $S'=\varnothing$; otherwise
both Hom sets are empty because neither $f^{-1}(\varnothing)=S'$
nor a map $I'\to\varnothing$ exists. The relaxation recorded
at scans 59--61 forgets the source boundary for a zero-dimensional
target, and its groupoid version does exactly the same.

In case III, conjugate $(f,\alpha)$ by an element of $E_U$
above $\alpha^{-1}$. It becomes $(s,1)$ with $p_Us=1$.
The remaining conjugations are exactly those by $P_U$, giving
the section classes, as in scans 62--64. Thus no dominant Hom
theorem proves all of this case: its exact comparison morphism
is $U(K)\to\operatorname{Sec}(E_U/G_K)/P_U$. For affine
targets this map misses the tangential packets, as the explicit
Kummer counterexample in Part B proves; the corrected version
has to exclude those packets. For proper targets its global
surjectivity is the remaining question just stated.

For IV(a), the trichotomy of Proposition~\ref{F-local-dichotomy}
has no cusp alternative because the target is proper. It gives
either an open geometric image, handled by Theorem A and the
reduction in Proposition~\ref{A:relative-hom}, or the unique
factorization through $G_K$ just written. These alternatives
are disjoint, since a constant map kills the geometric group
and a dominant map has open geometric image. The same
finite-$p$-adic base change used in the preceding proposition
proves this classification for fields finitely generated over
$\mathbf Q$.

For IV(c), $I'\ne\varnothing$ and $I=\varnothing$. The map
of finite index sets required by the original category does not
exist. Its geometric condition has the identical failure
$f^{-1}(\varnothing)=\varnothing\ne S'$. The later enlargement
in section 13 is implemented by filling the source punctures
on fundamental groups; its cohomological compatibility is
spelled out in Proposition~\ref{A:kummer-ambiguity}.

Finally IV(d) has both boundary sets nonempty. Its positive
integral local degrees force an open geometric image by
Proposition~\ref{A:positive-degree-open}, with the common
cyclotomic multiplier supplied by the boundary diagram.
Theorem A now applies over every sub-$p$-adic field.
Propositions~\ref{A:relative-hom} and
\ref{A:boundary-recovery} recover the unique map of pairs,
its local degrees, its nonrational boundary points, and every
2-morphism. Tamagawa's Theorem 0.3 gives the corresponding
affine Isom assertion over fields finitely generated over
$\mathbf Q$; taking an isomorphism and its inverse in the
same construction gives precisely his formulation.
\end{proof}

\begin{statement}{A.4}{Refuted}{the scope of Lead 1}
\label{st:A.4}
\label{A:lead-one}
The dominant part of Conjectures 5.5 and 7.1 is proved by the
explicit chain of maps in Proposition~\ref{A:relative-hom}.
The assertion that the Hom theorems settle all of Conjecture 5.5
is false for the uncorrected affine point-source instance: a
tangential section gives an actual extra groupoid morphism.
After the section-9 correction, the proper point-source
instance is still exactly the global section map of case III,
and the all-source comparison is the proved equivalence of
Proposition~\ref{F-full-curve-reduction}. Thus the relation
between the two formulations is the displayed functor and
these explicit reductions; it does not remove that surviving
point-source assertion.
\end{statement}

\begin{statement}{A.5}{Proved}{degrees, natural transformations, and the missing factor}
\label{st:A.5}
\label{A:two-morphisms}
With the convention above, a morphism between the connected boundary
systems consists of
\[
 f:E'\to E,\quad \tau:I'\to I,\quad f_j:E'_j\to E_{\tau(j)},
 \quad a_j\in E,\quad \alpha\in\Gamma
\]
subject to
\[
 pf=\operatorname{Int}(\alpha)p',\qquad
 \rho_{\tau(j)}f_j=\operatorname{Int}(a_j)f\rho'_j,
 \qquad \beta_j=p(a_j)\alpha,
\]
and
\[
 (k_{\tau(j)})^{-1}f_jk'_j
       =d_j\chi(\beta_j)\operatorname{id}_T,
       \qquad d_j\in\mathbf Z_{>0}.
\]
A natural transformation from these data to their primed versions is
given by $\mu\in E$ and $\mu_j\in E_{\tau(j)}$ with
\[
\begin{aligned}
 f'&=\operatorname{Int}(\mu)f,&
 f'_j&=\operatorname{Int}(\mu_j)f_j,\\
 \alpha'&=p(\mu)\alpha,&
 a'_j&=\rho_{\tau(j)}(\mu_j)a_j\mu^{-1}.
\end{aligned}                                                  \tag{A1}
\]
The last factor $\mu^{-1}$ is required.  It is absent from the boxed
formula on the actual scan 51, not just from its transcription.
\end{statement}
\begin{proof}
A natural transformation between homomorphisms of one-object
groupoids is an element conjugating the first homomorphism into the
second.  Naturality on the boundary square gives
\[
 \rho_{\tau(j)}(\mu_j)a_j=a'_j\mu,
\]
which proves (A1).  Naturality on the base square gives the formula
for $\alpha'$.  Applying $p$ to the boundary square gives
$p_{\tau(j)}f_j=\operatorname{Int}(\beta_j)p'_j$.  The action on
$T$ of the base-change path $\beta_j$ is $\chi(\beta_j)$; the
local map $t\mapsto u(t')^{d_j}$ contributes multiplication by the
positive ramification index $d_j$.  These two operations give the
displayed degree formula.  Conversely those equations are exactly
the four naturality equations and the Kummer compatibility in the
groupoid prism, so they construct a morphism of that diagram.

There is a finite certificate for the missing factor.  Take all
groups in the boundary square to be $S_3$, take $\rho=f=f_j$ to be
the identity and $a_j=1$, and take $\mu=(12)$ and $\mu_j=1$.
Then $f'=\operatorname{Int}(\mu)$ and $f'_j=1_{S_3}$.
The corrected $a'_j=\mu^{-1}$ satisfies the square; the printed
$a'_j=1$ would assert $\operatorname{Int}(\mu)=1$.  The independent
program \texttt{code/A/exact\_checks.py} checks both equations on
all six elements.  This finite example tests the naturality formula,
and does not claim that its groups are geometric boundary systems.
\end{proof}

\begin{statement}{A.6}{Refuted}{unrestricted rigidity of the preliminary category}
\label{st:A.6}
\label{A:abstract-rigidity}
The preliminary axioms for $\mathcal B$ on scans 30--31 do not imply
$\operatorname{Aut}(f,\alpha)=1$.  For geometric hyperbolic target
systems over a field finitely generated over $\mathbf Q$, this
automorphism group is trivial.  For any hyperbolic geometric target,
it is trivial for morphisms whose image on the geometric group is
open.  Thus the corresponding Hom groupoids are equivalent to
discrete sets, with unique isomorphisms between isomorphic objects.
\end{statement}
\begin{proof}
For the first assertion take no boundary objects, take
$E=G_K\times C_2\to G_K$, and take the homomorphism
$f:G_K\to E$, $\sigma\mapsto(\sigma,1)$, with its identity
base-commutation isomorphism.  All the preliminary conditions hold.
Every element of the second factor is a natural automorphism over
$BG_K$, so the automorphism group contains $C_2$.

For a geometric target, an automorphism in the correctly pinned
Hom groupoid has $p(\mu)=1$ by (A1).  The element $\mu\in\pi$
centralizes $f(E')$.  The elementary geometric-centralizer argument
below proves it is one when $K$ is finitely generated over
$\mathbf Q$ and $p(f(E'))$ is open.  If the geometric image is
open, geometric slimness already gives this conclusion.  The
boundary components $\mu_j$ are then determined by (A1) and the
injectivity of $\rho_i$, and equal one.  In a groupoid, trivial
automorphism groups imply that any isomorphism, when present, is
unique.  This proves the precise discreteness assertion.
\end{proof}

\begin{statement}{A.7}{Proved}{arithmetic geometric-centralizer rigidity}
\label{st:A.7}
\label{A:geometric-centralizer}
Let $U/K$ be hyperbolic with $K$ finitely generated over
$\mathbf Q$.  For every closed $H\subset E=\pi_1(U)$ whose image
in $G_K$ is open, one has $C_\pi(H)=1$.
\end{statement}
\begin{proof}
Suppose $1\ne z\in C_\pi(H)$.  A finite quotient of $\pi$ in
which $z$ has nontrivial image, followed by the inverse image of the
cyclic subgroup generated by that image, gives an open subgroup
$P\subset\pi$ in which the image of $z$ in $P^{\rm ab}$ is
nonzero.  More precisely choose a prime $\ell$ and an open normal
$N\subset\pi$ with $N\cap\overline{\langle z\rangle}$ contained
in the subgroup of $\ell$th powers of its $\ell$-component, and put
$P=N\overline{\langle z\rangle}$.  The cyclic quotient $P/N$
detects a nonzero $\ell$-component.  The group $P^{\rm ab}\otimes
\mathbf Z_\ell$ is a free module, by the surface presentation of
the corresponding finite \'etale cover.  Its class of $z$ is
therefore a nonzero vector, not merely a torsion class.

The cover descends to a finite extension $L/K$.  Intersect $H$ with
its arithmetic covering group and replace $L$ by the finite
extension corresponding to the image of that intersection.  The
resulting group still centralizes $z$ and surjects onto $G_L$.
It therefore fixes this nonzero homology vector.  Spread the cover
and its smooth compactification and boundary over a normal
finite-type arithmetic model, and take a closed point of good
reduction whose residue characteristic differs from $\ell$.
Arithmetic Frobenius at that point acts on the boundary homology
with eigenvalues of complex absolute value $q$, and on the
compactification's homology with absolute value $q^{1/2}$.
Here $q>1$.  Neither part has eigenvalue one, contradicting the
fixed vector.  This uses only the usual weight-one theorem for the
Jacobian and the explicit cyclotomic boundary module.  It applies
also to a proper cover, when the boundary part is zero.
\end{proof}

\subsection{Full and primewise commutants; a new metabelian counterexample}

Write $U_0=\mathbf P^1_{\mathbf Q}\setminus\{0,1,\infty\}$,
$F=\pi_1(U_{0,\overline{\mathbf Q}})$ and
$F_\ell=F^{(\ell)}$.  A peripheral outer automorphism preserves
the unordered conjugacy classes of the three cyclic inertia groups.
Its pure subgroup preserves each of the three labels.  On a quotient
of $F$, the peripheral groups in this definition mean their actual
images in that quotient.

\begin{statement}{A.8}{Known}{the full and every pro-prime tripod commutant}
\label{st:A.8}
\label{A:tripod-commutant}
For every number field $K$ and every prime $\ell$, including $2$,
\[
 C_{\operatorname{Out}(F)}(G_K)=S_3,
 \qquad
 C_{\operatorname{Out}(F_\ell)}(G_K)=S_3.
\]
The embeddings of $S_3$ are those induced by the six fractional
linear transformations permuting $0,1,\infty$.  Intersecting with
the peripheral outer automorphism groups leaves these equalities
unchanged.  More generally the same argument gives
$C_{\operatorname{Out}(\pi)}(G_K)=\operatorname{Aut}_K(U)$ in
the full and the geometrically pro-$\ell$ setting for every
hyperbolic $U/K$ over a sub-$\ell$-adic field.
\end{statement}
\begin{proof}
For a center-free profinite group $D$ and a continuous outer action
$r:G_K\to\operatorname{Out}(D)$, the associated extension is
canonically
\[
 \operatorname{Aut}(D)\mathop{\times}_{\operatorname{Out}(D)}G_K.
                                                               \tag{A2}
\]
Indeed conjugation and projection map the given extension into
(A2); injectivity is $Z(D)=1$, and every lift of an outer class
differs from the conjugation action of a chosen lift by a unique
inner automorphism.  If $[a]\in\operatorname{Out}(D)$ commutes
with $r(G_K)$, the rule $(b,\sigma)\mapsto(aba^{-1},\sigma)$
is an automorphism of (A2) over $G_K$.  Changing $a$ by an inner
automorphism changes this rule by geometric inner conjugacy.

The geometric full surface groups and geometric pro-$\ell$
hyperbolic groups are center-free.  Apply Mochizuki's Theorem A,
equivalently Theorem 16.5, to this arithmetic isomorphism
\cite{A_Mochizuki1999}.  It is an open continuous homomorphism
over the specified base field.  Applying the same theorem to its
inverse shows that the resulting dominant morphism is a unique
$K$-automorphism.  The theorem explicitly allows every prime in
its geometrically pro-prime version; no odd-prime assumption
enters this argument.  For $U_0$, its automorphisms are exactly
the six transformations permuting its deleted set, all defined
over $\mathbf Q$.  Each induces the stated commuting outer
automorphism.  This proves both surjectivity and injectivity.
\end{proof}

\begin{statement}{A.9}{Proved}{the abelian peripheral commutant}
\label{st:A.9}
\label{A:abelian-commutant}
On $F^{\rm ab}=\widehat{\mathbf Z}^2$ the peripheral commutant is
$\widehat{\mathbf Z}^{\times}\times S_3$.  In the whole linear
automorphism group it is $\operatorname{GL}_2(\widehat{\mathbf Z})$.
In particular the smallest derived length giving commutant $S_3$
is greater than one.
\end{statement}
\begin{proof}
Choose the first two cusp classes $e_1,e_2$; the third is
$-e_1-e_2$.  A matrix preserving each of the three corresponding
primitive rank-one direct summands has the form
$\operatorname{diag}(a,b)$, and preservation of
$\widehat{\mathbf Z}(e_1+e_2)$ forces $a=b$.
Allowing permutations adds exactly the displayed copy of $S_3$;
its intersection with the scalar subgroup is trivial.
The arithmetic action is $\chi\operatorname{id}$ and hence
commutes with every matrix.  This gives the two asserted groups.
\end{proof}

\begin{statement}{A.10}{Proved}{infinitely many arithmetic-centralizing metabelian involutions}
\label{st:A.10}
\label{A:metabelian-commutant}
Let $M=F/\overline{[F',F']}$ be the full profinite metabelian
tripod group.  Its pure peripheral outer automorphism group
contains a closed subgroup
\[
       \prod_{p\geq5\ \mathrm{prime}} C_2
       \ \subset\ C_{\operatorname{Out}_{\rm lac}(M)}(G_K)
                                                               \tag{A3}
\]
for every number field $K$.  This subgroup also centralizes the
geometric $S_3$, and all of its elements act trivially on every
maximal pro-prime quotient of $M$.  Consequently the metabelian
commutant is strictly larger than $S_3$, and the requested
smallest derived length, when defined, is at least three.
\end{statement}
\begin{proof}
Put $A=M^{\rm ab}=\widehat{\mathbf Z}^2$,
$R=\widehat{\mathbf Z}[[A]]$, and $c=[x,y]$.
The profinite Alexander calculation gives $M'=Rc$, freely of
rank one as a continuous $R$-module.  Chen--Deligne prove both
this statement and the precise outer-automorphism identity
\[
 \operatorname{IOut}(M)\simeq R^{\times,\epsilon=1}/A,
 \qquad
 \operatorname{Out}(M)\simeq
       (R^{\times,\epsilon=1}/A)\rtimes\operatorname{GL}(A).
                                                               \tag{A4}
\]
The semidirect action is the natural action on $R$, not a
new arithmetic action; these are their Corollary 4.1.5,
Proposition 4.3.7 and Theorem 5.2.4
\cite{A_ChenDeligne2017}.  We now specify the units, actual
automorphisms, peripheral conjugators and continuity needed
to use (A4).

Fix a prime $p\geq5$.  In $A^{(p')}$ let
$\theta:A^{(p')}\to C_3$ send both $X$ and $Y$ to the same
generator $t$, and put $H=\ker\theta$.  Normalized averaging
over $H$ defines an idempotent $e_H$ in
$\mathbf Z_p[[A^{(p')}]]$: in every finite quotient $B$ its
image is $|H_B|^{-1}\sum_{h\in H_B}h$.  The denominators are
$p$-adic units, and these averages are compatible under
quotient maps.  Define
\[
 e_p=e_H\left(1-\frac{1+t+t^2}{3}\right)
\]
in the coefficient-$p$ factor of $R$, and set its other
coefficient components equal to zero.  This is a nonzero
continuous idempotent.  On its factor, the prime-to-$p$
parts of $X,Y$ are primitive cube roots of unity.  Hence
$X-1$, $Y-1$ and $XY-1$ are units on $e_pR$, whether the
quadratic cyclotomic factor splits over $\mathbf Z_p$ or not.

Set
\[
 r_p=\frac{-2e_p}{(X-1)(Y-1)},\quad
 s_p=\frac{X(Y-1)r_p}{XY-1},\quad
 f_p=c^{r_p},\quad h_p=c^{s_p},\quad u_p=1-2e_p.
\]
All divisions are in $e_pR$ and are then extended by zero
on its complement.  The rule
\[
       \phi_p(x)=x,\qquad \phi_p(y)=f_p^{-1}yf_p             \tag{A5}
\]
defines a continuous endomorphism by the universal property
of $M$.  Direct multiplication in the extension
$1\to Rc\to M\to A\to1$ gives
\[
 \phi_p(c)=c^{u_p},\qquad
 \phi_p(xy)=c^{X(Y-1)r_p}xy=h_p^{-1}xyh_p.
\]
Moreover $(1+u_p)r_p=(1+u_p)s_p=0$, so
$f_p\phi_p(f_p)=h_p\phi_p(h_p)=1$ and $\phi_p^2=1$.
The action on $e_pM'$ is $-1$, so $\phi_p$ has order
exactly two and preserves the three labelled peripheral
conjugacy classes.  Its determinant in (A4) is $u_p$.

Here is the arithmetic check.  The action of every
$\sigma\in G_K$ on $A$ is the scalar $\chi(\sigma)$.
Every scalar unit preserves $H$, and it permutes the two
nontrivial characters of $C_3$.  Thus
\[
      \sigma_R(e_p)=e_p,
      \qquad \sigma_R(u_p)=u_p.                            \tag{A6}
\]
For an arbitrary element $(v,g)$ of the right-hand side of
(A4), conjugation on $(u,1)$ is
$(v,g)(u,1)(v,g)^{-1}=(g(u),1)$, because the unit group
is abelian.  Applying this exact formula with the
arithmetic $g=\chi(\sigma)\operatorname{id}$ and (A6)
proves $[\phi_p]\rho(\sigma)[\phi_p]^{-1}=\rho(\sigma)$.
No claim that an arithmetic representative commutes in
$\operatorname{Aut}(M)$ is needed or used.
The generators of $S_3$ act on the character $(1,1)$
modulo $3$ by the same character or its negative:
permuting $x,y$ fixes it, and replacing $(x,y)$ by
$((xy)^{-1},y)$ sends it to $(-2,1)=(1,1)$ modulo $3$.
They preserve $e_p$ as well, proving the $S_3$ assertion.

For any subset $J$ of the primes at least five, the
coefficientwise element $e_J=\sum_{p\in J}e_p$ is an
idempotent of $R$.  The same formulas give $\phi_J$.
Orthogonality $e_pe_q=0$ for $p\ne q$ makes the product
law exactly symmetric difference of subsets.  Every
finite quotient of $M$ uses only finitely many coefficient
primes in its finite derived subgroup; consequently this
map from $\prod C_2$ to $\operatorname{Aut}(M)$ is
continuous.  Its composite to $\operatorname{Out}(M)$ is
injective.  Indeed, an inner automorphism has determinant
in $A$ by (A4), whereas $u_J$ is one on its entire
coefficient-$3$ factor.  The embedding
$A\hookrightarrow\mathbf Z_3[[A]]^\times$ is injective,
as is seen in every finite group algebra.  Thus
$u_J\in A$ forces that monomial to be one; its nonzero
$p$-component then forces $J=\varnothing$.
Compactness of $\prod C_2$ gives the asserted closed image.

Finally every pro-$\ell$ quotient kills $e_p$ for
$p\ne\ell$.  For $p=\ell$ it kills $A^{(p')}$, and
$1-(1+t+t^2)/3$ becomes zero.  Thus (A5) becomes the
identity at every prime.  Every nonidentity member of
(A3) is IA; the faithful action of $S_3$ on $A$ shows
that it is outside $S_3$.
\end{proof}

\begin{statement}{A.11}{Proved}{an independent finite certificate for the metabelian construction}
\label{st:A.11}
\label{A:finite-metabelian}
The prime-$5$ involution of Proposition~\ref{A:metabelian-commutant}
has a noninner peripheral image on the group
\[
 G=\mathbf F_{25,\rm add}\rtimes C_3,
 \qquad (v,i)(w,j)=(v+\zeta^i w,i+j),
 \qquad \zeta^2+\zeta+1=0.
\]
Take $x=(0,1)$ and $y=(1,1)$.  They generate $G$, and
the involution is $(v,i)\mapsto(-v,i)$.
\end{statement}
\begin{proof}
In this quotient
\[
 c=(\zeta-1,0),\quad r=1+\zeta,\quad s=\zeta,
 \quad f=(3+4\zeta,0),\quad h=(4+3\zeta,0).
\]
These identities give $\phi(x)=x$,
$\phi(y)=f^{-1}yf$, $\phi(xy)=h^{-1}xyh$, and both
norm equations $f\phi(f)=h\phi(h)=1$.  Conjugates of
$\zeta-1$ span $\mathbf F_{25}$ over $\mathbf F_5$, so
$x,y$ generate all $75$ elements.  An inner automorphism
acts on the normal additive subgroup by $1,\zeta$ or
$\zeta^2$, whereas $\phi$ acts by $-1$.  Hence it is
noninner.

The independent program \texttt{code/A/exact\_checks.py}
constructs the multiplication, checks all $75^3$
associativity triples, all $75^2$ multiplication pairs
for $\phi$, all possible inner conjugators, the generator
claim, and the displayed peripheral and norm equations.
It also constructs all $1200$ automorphisms from a
semilinear part $v\mapsto a v$ or $a v^5$ and a
translation of the complement, and verifies that every
commutator with $\phi$ is inner.  All assertions pass.
The infinite arithmetic-centralizer conclusion uses
(A4)--(A6); this finite calculation is its concrete
certificate of nontriviality.
\end{proof}

\begin{statement}{A.12}{Open}{the same-level commutant at larger finite derived length}
\label{st:A.12}
\label{A:solvable-open}
The computation above settles the abelian and metabelian failures,
and the full and pro-prime equalities have been proved.  The
smallest $m\geq3$ with same-level commutant $S_3$ is not determined
here.  A theorem reconstructing the isomorphism class of a curve
from its $m$-step quotient is insufficient: it supplies a map
from isomorphism classes, whereas the required assertion concerns
every element of an automorphism group.  The precise additional
map in Mochizuki's Theorem A$'$ forgets three derived steps.
Yamaguchi's strong finite-step statements likewise specify the
image in a lower-step quotient \cite{A_Yamaguchi2024}; they do
not assert that the kernel of this forgetting map is trivial.
The involutions just constructed make this distinction visible by
an actual element and its images.  The remaining task is to
control that kernel or produce an arithmetic-centralizing lift
through the next derived quotient. Proposition~\ref{A:solvable-projection-limit}
below determines the exact lower-step images and the
prosolvable limit; it does not determine this same-level kernel.
Theorem~\ref{A:third-derived-IA} below eliminates the entire
IA image from length three to two, and
Corollary~\ref{A:third-derived-projection} identifies the
whole projected peripheral commutant with the geometric
$S_3$. Thus the remaining same-level question is the
kernel of the third-to-second-derived map, not the
lifting of the metabelian unit family.
\end{statement}

\subsection{Cyclotomic recognition and the full profinite obstruction}

\begin{statement}{A.13}{Refuted}{the literal maximal-subgroup criterion of scan 483}
\label{st:A.13}
\label{A:literal-normalizer}
For a hyperbolic curve, the condition
$p(N_E(A))\subset G_K$ open does not characterize inertia among
all closed nontrivial subgroups $A\subset\pi$.  It is satisfied
by $A=\pi$, by every characteristic open subgroup of $\pi$, and
by $\overline{[\pi,\pi]}$.  None of these assertions requires
a section or a rational cusp.
\end{statement}
\begin{proof}
Each displayed subgroup is normal in $E$, so its normalizer is
$E$, whose arithmetic image is open.  The geometric group $\pi$
of a hyperbolic curve is nonabelian, whereas cusp inertia is
procyclic.  Thus $A=\pi$ already contradicts the asserted
containment in a loop group.  An open characteristic subgroup
exists by intersecting the finitely many subgroups of any given
bounded index in the finitely generated profinite group $\pi$.
Such a subgroup is again a hyperbolic geometric fundamental
group, so is nonabelian.  The closure of the commutator is
characteristic and nontrivial.  This proves each asserted
normalizer equality and the counterexample.
\end{proof}

\begin{statement}{A.14}{Known}{the precise cyclotomic criterion and transporter}
\label{st:A.14}
\label{A:cyclotomic-normalizer}
Let $U/K$ be a hyperbolic curve over a number field.  For a
nontrivial closed subgroup $A\subset\pi$, the following three
conditions are equivalent:
\begin{enumerate}
\item $A$ is contained in a cusp inertia group $L$;
\item $A$ is procyclic, $p(N_E(A))$ is open, and the action
of $N_E(A)$ on $A$ is the cyclotomic action;
\item $A$ is procyclic and some closed subgroup of $N_E(A)$
has open arithmetic image and acts on $A$ cyclotomically.
\end{enumerate}
The containing inertia is unique.  The maximal subgroups in
this class are precisely the full cusp inertia groups, and
\[
 N_E(L)=E_{\widetilde s},\qquad
 \operatorname{Transp}_E(L^d,L)=E_{\widetilde s}\quad(d\geq1),
                                                               \tag{A7}
\]
where the transporter means $\{g:gL^dg^{-1}\subset L\}$.
This includes nonrational cusps.
\end{statement}
\begin{proof}
Nakamura's Theorem 2.1.1 is formulated for elements and a
cyclotomic normalizer, for every full class of finite groups
\cite{A_Nakamura1994}.  Stix's Theorem 20 writes exactly the
three displayed conditions for arbitrary nontrivial closed
subgroups \cite{A_StixCuspidal}.  Its field hypothesis is
$\operatorname{Hom}_{G_L}(T_\ell B,T_\ell J)=0$ for tori
$B$ and abelian varieties $J$, over every finite extension
$L/K$.  For number fields this follows by choosing a prime
of good reduction away from $\ell$: after a further finite
extension the torus is split, its eigenvalues have absolute
value $q$, and those on $J$ have absolute value $q^{1/2}$.
Thus all hypotheses of that theorem are verified here.

For clarity, its finite-cover step detects a noncuspidal
element in the proper homology of a finite cover, after
quotienting out the boundary classes.  The cyclotomic
normalizer supplies a nonzero $\mathbf Z_\ell(1)$ line in
that proper homology.  The preceding Frobenius comparison
excludes it.  This is a statement about an actual finite
cover of the full group; it does not assume that maximal
pro-prime quotients jointly separate that group.

Distinct cusp inertia groups have trivial intersection;
the same assertion holds for their nontrivial closed
subgroups, including a subgroup supported at a single
prime.  One proof passes to a finite cover separating the
two lifted cusps and then takes an abelian cover ramified
at one and at a third cusp but not the other.  This is the
peripheral centralizer/separation input in Nakamura's
Lemma 2.1.2 and Stix's Corollary 7.  It proves uniqueness,
and forces any element normalizing $A\subset L$ to fix
the same lifted cusp.  The normalizer is therefore its
decomposition group and acts cyclotomically.  Local
Kummer theory proves the converse.  Maximality now gives
the assertion for full inertia groups.

For (A7), $gL^dg^{-1}\subset L$ has nontrivial intersection
with $L$, hence the lifted cusp of $gLg^{-1}$ is the
original one.  Thus $g\in E_{\widetilde s}$.  Every element
of this decomposition group preserves $L$ and its subgroup
$L^d$, giving the reverse inclusion.  This proof follows
the lifted cusp and so does not require its residue field
to equal $K$.
\end{proof}

\begin{statement}{A.15}{Proved}{uniqueness without a previously fixed surface type}
\label{st:A.15}
\label{A:unique-loops}
For an affine hyperbolic curve over a number field, a
hyperbolic surface loop structure on its geometric full
profinite group that is invariant under an open subgroup
of $G_K$ and satisfies scan 84(c) is the geometric loop
structure.  In particular its genus and number of holes
equal those of the curve.  The same conclusion holds
over a field finitely generated over $\mathbf Q$.
\end{statement}
\begin{proof}
Pass to a finite extension that fixes each of the finitely
many proposed cusp classes.  Condition (c) supplies, for
each proposed loop $L'_j$, a cyclotomic normalizing subgroup
with full arithmetic image.  Proposition
\ref{A:cyclotomic-normalizer} puts $L'_j$ inside a genuine
inertia group.  A peripheral group of any hyperbolic
profinite surface presentation is its own normalizer in
the geometric group, by commensurable terminality
\cite[Lemma 1.3.7]{A_Mochizuki2004}.  Since a containing
procyclic group normalizes it, this inclusion is equality.
Distinct proposed loop classes remain distinct, so the
proposed cusps form a subset of the actual cusp classes.

Write $(g,r)$ for the actual type and $(h,s)$ for the
proposed one.  Thus $s\leq r$, and abelianization gives
$2h+s-1=2g+r-1$.  If $h=0$, this yields
$s=2g+r\geq r$, and hence $s=r$, $g=0$.
If $h>0$, quotienting by the closed normal subgroup
generated by the proposed loops gives the proper
surface group of genus $h$.  If any actual cusp
remained, the same quotient computed from the actual
presentation would be a free profinite group and
would have $H^2(-,\mathbf F_\ell)=0$.  A proper
surface group of positive genus has
$H^2(-,\mathbf F_\ell)=\mathbf F_\ell$, a contradiction.
Thus $s=r$ in this case too and the rank equality gives
$h=g$.  The actual quotient maps and cohomology groups
establish the comparison; no identification of differently
presented groups is presumed.

Over a finitely generated field, spread out the curves
in the finite-cover weight argument in the preceding
proposition and specialize at a good finite-residue-field
point.  This verifies the identical no-map-from-a-torus
hypothesis.  The remainder of the proof is unchanged.
\end{proof}

\begin{statement}{A.16}{Known}{the cyclotomic-kernel centralizer form}
\label{st:A.16}
\label{A:cyclotomic-centralizer}
Let $K$ be finitely generated over $\mathbf Q$, and put
$\Gamma^\circ=\ker\chi$.  For a germ of section $s$, the
group $C_\pi(s(\Gamma^\circ))$ is either trivial or one
full cusp inertia group.  Every cusp inertia group
occurs, after a finite extension, as such a centralizer.
For a section whose image is already contained in a
cyclotomic normalizer of $A\ne1$, this centralizer contains
$A$ and identifies its unique cusp.
\end{statement}
\begin{proof}
The full finite-cover proof is Proposition~\ref{F-centralizer}.
Its extra arithmetic input is the absence of nontrivial
divisible torsion in the Tate module of an abelian variety
over $K(\mu_\infty)$, obtained by specialization and Ribet's
theorem.  Thus Stix's recognition theorem applies over
$K(\mu_\infty)$ to every nonidentity fixed element.
The peripheral separation argument puts all such
elements in the same loop, and normality of
$\Gamma^\circ$ shows that the section normalizes that
loop.  Since $\chi=1$ on $\Gamma^\circ$, its whole
inertia group is fixed, giving equality.  Conversely a
nonzero rational tangent at any rationalized cusp
gives a section through its local decomposition group;
the preceding equality gives exactly that inertia.
The last assertion follows from
$s(\Gamma^\circ)$ centralizing $A$ and this equality.
\end{proof}

\begin{statement}{A.17}{Proved}{the proof-level comparison of the two cyclotomic forms}
\label{st:A.17}
\label{A:normalizer-implications}
The cyclotomic normalizer recognition proof, applied over
$K(\mu_\infty)$ with its Ribet hypothesis, gives the
section-centralizer form (1 bis).  Conversely the same
finite-cover centralizer proof applies to every closed
$B\subset E$ whose image is open in $\ker\chi$ and
which lies over that kernel.  This strengthened
centralizer assertion gives (1) for arbitrary arithmetic
normalizers, without a splitting hypothesis.
\end{statement}
\begin{proof}
The first assertion is the proof of
Proposition~\ref{A:cyclotomic-centralizer}.  For the
second, use the strengthened assertion in
Proposition~\ref{F-centralizer}.  Its proof needs only
that the image of $B$ is open in $\ker\chi$: after
descending a finite geometric cover with group $P$, the
intersection of $B$ with its arithmetic covering group
still has open image, and its action on $P^{\rm ab}$
is the ordinary arithmetic action.  Its kernel acts
inner and hence acts trivially on the abelianization.
Thus the Ribet and boundary-recognition steps are
unchanged.

Let $A\ne1$ have an open arithmetic normalizer $N$
acting cyclotomically.  Put $B=N\cap p^{-1}(\ker\chi)$.
Then $p(B)=p(N)\cap\ker\chi$ is open in $\ker\chi$,
and $A\subset C_\pi(B)$.  The strengthened assertion
puts $A$ in its unique inertia group.  The converse
follows from the local decomposition extension.
This is a reduction between the proofs with their stated
arithmetic input.  The statement quantified only over
sections does not itself supply a section of an arbitrary
normalizer extension; none is used here.
\end{proof}

\begin{statement}{A.18}{Known}{a full profinite replacement over a p-adic local field}
\label{st:A.18}
\label{A:padic-cusp-reconstruction}
Let $K_i/\mathbf Q_{p_i}$ be finite and $U_i/K_i$
hyperbolic.  Every continuous isomorphism
$\pi_1(U_1)\xrightarrow{\sim}\pi_1(U_2)$ preserves
the geometric subgroups, the types $(g_i,r_i)$, and
the full cusp inertia groups and their arithmetic
normalizers.  The inertia reconstruction uses all
finite covers and the following formula, with the
usual integral weight convention:
\[
 r_V-1=\dim(V_\ell^{\mathrm{wt}\,2})
            -\dim(V_\ell^{\mathrm{wt}\,0})
 \quad(r_V>0),
 \qquad V_\ell=H_1(V_{\overline K},\mathbf Q_\ell),
                                                               \tag{A8}
\]
where $\ell\ne p_i$.  This gives a reconstruction even
when a Tate curve contributes a cyclotomic line to
proper homology.
\end{statement}
\begin{proof}
Mochizuki's Lemma 1.3.8 first recovers the arithmetic
quotient.  The local-field group theory in Proposition
1.2.1 recovers its residue cardinality and Frobenius.
The argument of Lemma 1.3.9 then gives (A8)
\cite{A_Mochizuki2004}.  Here are its precise weight
and covering steps.  Boundary homology is of weight
two.  Proper homology is self-dual with a Tate twist,
so its weight-zero and weight-two dimensions agree.
The weight-zero part is unaffected by removing
punctures.  Subtracting the two dimensions therefore
leaves exactly the boundary dimension $r_V-1$.
Quasi-unipotence of local inertia makes these weights
independent of the chosen lift of Frobenius.  The
free/proper alternative and
$\dim H_1=2g+r-1$ then recover the type.

For a Galois $\ell$-power cover $Z\to V$, the equation
\[
          r_Z=[Z:V](r_V-1)+1                              \tag{A9}
\]
recognizes total ramification over exactly one cusp.
Indeed every ramified cusp contributes a deficit
$d(1-1/e)$ to the unramified count $dr_V$, with
$e$ a nontrivial $\ell$-power.  Two such deficits
already sum to at least $d$, whereas (A9) has
deficit $d-1$.  Thus there is one, and its index
is $d$.  The same property on pro-$\ell$ covers,
together with commensurable terminality, recovers
their maximal inertia subgroups.  This works for
every $\ell$, including $p_i$, because the cusp
counts were already obtained using an auxiliary
prime.  Repeating on every finite \'etale cover
recovers the cusps of the completed universal
cover; their stabilizers are the actual full
inertia and decomposition groups.  Every map
used in this construction is induced by the
given isomorphism, so it is functorial.
\end{proof}

\begin{statement}{A.19}{Proved}{the degree-five CM obstruction, independently checked}
\label{st:A.19}
\label{A:cm-degree-five}
In the rational free Lie algebra on $x,y$, put
$\omega=[x,y]$, let $F_0(x)=y$, $F_0(y)=0$, and set
\[
 \epsilon_4(x)=\operatorname{ad}_x^4y,\qquad
 \epsilon_4(y)=[y,\operatorname{ad}_x^3y]
                 -[[x,y],\operatorname{ad}_x^2y],
 \qquad T_4=[F_0,\epsilon_4].
\]
With standard Lyndon bracketing for $x<y$,
\[
 T_4(x)=-2b_{xxxyy}-b_{xxyxy},\qquad
 T_4(y)=2b_{xxyyy}+3b_{xyxyy},\qquad T_4(\omega)=0.
                                                               \tag{A10}
\]
In bidegree $(3,2)$, $[x,L_4]$ has dimension one,
whereas adjoining $T_4(x)$ raises the dimension to
two.  Thus a CM eigenline with first term $x$
cannot satisfy the conjugacy equation for a
nonzero degree-four arithmetic term proportional
to $T_4$.
\end{statement}
\begin{proof}
Expand the definitions in the free associative
algebra.  The two Lyndon words in (A10) are
linearly independent because their least words
have coefficient one.  The degree-four
bidegree-$(2,2)$ Lie space has its single
Lyndon basis vector $b_{xxyy}$, and its bracket
with $x$ is $b_{xxxyy}$.  The coefficient
$-1$ of $b_{xxyxy}$ proves the rank assertion.
The program \texttt{code/A/exact\_checks.py}
independently expands these integral polynomials,
checks $T_4(\omega)=0$, and computes both exact
matrix ranks.  It also checks the kernel of
$T_4$ in each degree at most seven; the only
kernel is the boundary line in degree two.

For the conjugacy assertion, an IA arithmetic
element normalizing a procyclic lift with
nonzero primitive abelianization has exponent
one on that lift.  In degree five its first
outer error must therefore be the bracket
$[x,v]$ for an inner correction $v\in L_4$.
Terms of the lift above degree one contribute
only from degree six onward.  Equation (A10)
contradicts the necessary degree-five equation.
\end{proof}

\begin{statement}{A.20}{Known}{the audited pro-prime CM normalizer theorem}
\label{st:A.20}
\label{A:cm-prime-normalizer}
Put $K=\mathbf Q(i)$, $E_0:y^2=x^3-x$,
$U=E_0\setminus\{O\}$, and
$P_\ell=\pi_1(U_{\overline K})^{(\ell)}$.
For every prime $\ell$, every nontrivial closed
procyclic subgroup of $P_\ell$ with open arithmetic
normalizer is contained in a cusp inertia group.
At $\ell=5$, a primitive CM eigenline lifts to
such a subgroup of the class-four nilpotent
quotient, but has no such lift at class five.
\end{statement}
\begin{proof}
We spell out the reduction audited in LM8 G.3--G.13.
For a generator $u$ and its first lower-central
term $u_d$, the normalizer's character has
Frobenius absolute value $q^{d/2}$, since
$H_1(E_0)$ has weight one.  Choose an open
subgroup $N\overline{\langle u\rangle}$ in
which $u$ has nonzero abelian image, descend
that finite cover, and intersect with the
normalizer.  The same character now occurs
in the cover's homology.  Its weight is one
or two, so $d=1$ or $2$.

For an inert prime, and also for $\ell=2$,
$H_1(E_0,\mathbf Q_\ell)$ has no line stable
under an open arithmetic subgroup.  A positive
power of Frobenius at $5$ has eigenvalues
$(-1+2i)^n,(-1-2i)^n$; their ratio is not
a root of unity and $i\notin\mathbf Q_\ell$.
Thus $d=1$ is impossible.  At a split prime
the two possible lines are the CM eigenlines.
The order-four automorphism $[i]$ kills the
outer degree-four characters $(3,1),(1,3)$.
The remaining character $(2,2)$ is nonzero:
this is Ishii's Corollary 4.10, whose hypotheses
are an imaginary quadratic field of class
number one, full CM over that field, and a
split prime at least five
\cite{A_Ishii2025}.  All hold here at every
split prime.  The degree-four comparison
is injective on the three special character
lines, so that nonzero character is a multiple
of $T_4$.  Proposition~\ref{A:cm-degree-five}
excludes $d=1$.

Consequently $d=2$ and $u_2=a\omega$, $a\ne0$.
The equality
$a\eta(n)\omega=a\chi_\ell(p(n))\omega$
for every normalizer element gives
$\eta=\chi_\ell\circ p$ by cancellation in
$\mathbf Z_\ell$.  Nakamura's Theorem 2.1.1
applied to the full class of finite
$\ell$-groups now identifies actual inertia
in $P_\ell$; it is not merely a claim in a
Malcev completion.

For the exact stopping level at $5$, take
CM-equivariant symplectic logarithmic coordinates
through class four.  The needed denominators
are units in $\mathbf Z_5$; straightening the
boundary through degree five also uses only
these units.  The special degree-one and
degree-three modules vanish, and the
degree-two module is
$\mathbf Z_5\operatorname{ad}\omega$.
Hence the arithmetic action differs from
its diagonal homology action, on this
quotient, by an inner action of a power
of the boundary.  Correcting the arithmetic
lift makes $\exp(\mathbf Z_5x)$ stable.
Its primitive abelianization makes it
maximal procyclic.  The nonzero degree-four
term and (A10) exclude a class-five lift.
This verifies the obstruction and every
prime-specific hypothesis used in LM8's
normalizer reduction.  The exact
$\kappa_{11}/48$ normalization is immaterial
for this nonvanishing argument.
\end{proof}

\begin{statement}{A.21}{Proved}{full profinite recognition from one detected prime}
\label{st:A.21}
\label{A:cm-detected-normalizer}
With the preceding CM curve, let
$\pi=\pi_1(U_{\overline K})$ and
$R=\bigcap_\ell\ker(\pi\to P_\ell)$.
A nontrivial closed procyclic subgroup
$C\subset\pi$ with open arithmetic
normalizer is contained in actual full
inertia as soon as its image in one
$P_\ell$ is nontrivial.
\end{statement}
\begin{proof}
Its $\ell$-component is a copy of
$\mathbf Z_\ell$ and maps injectively to
its nonzero image in $P_\ell$.
Proposition~\ref{A:cm-prime-normalizer}
therefore makes the action on that actual
component cyclotomic.  Apply Nakamura's
theorem with the full class of finite
groups, now to this component inside
$\pi$.  It lies in a full inertia group
$I$.  The centralizer in $\pi$ of any
nonidentity element of $I$ is $I$, even
when its exponent is supported at only
one prime.  Since all of $C$ commutes
with its nontrivial $\ell$-component,
it follows that $C\subset I$.

For this one-puncture group the cited
centralizer assertion also has a short
reduction to a primitive peripheral
generator.  An index-two cover has two
cusps, and its geometric group contains
the original boundary generator as a
free basis member.  If $w$ centralizes
its nontrivial power, then $w^2$ lies
in that cover and in its peripheral
centralizer.  The boundary has zero
image in $\pi^{\rm ab}$, so
$2w^{\rm ab}=0$.  Torsion-freeness of
$\widehat{\mathbf Z}^2$ gives
$w^{\rm ab}=0$, whence $w$ itself
lies in the cover.  Nakamura's
primitive-generator centralizer lemma
then gives $w\in I$.
\end{proof}

\begin{statement}{A.22}{Open}{the common-kernel normalizer problem}
\label{st:A.22}
\label{A:cm-common-kernel}
For the full CM group the remaining problem is
the existence or exclusion of a nontrivial
closed procyclic $C\subset R$ whose normalizer
in the arithmetic full profinite group has
open image in $G_{\mathbf Q(i)}$. Proposition~\ref{A:deeper-common-kernel}
below further forces any such $C$ into the strictly smaller
group $R_{\rm nil}(F')\subset F''$. No such
subgroup has been constructed or excluded
here.  For a general affine hyperbolic curve
the same bare-normalizer problem also remains
unresolved here.  It is not settled by the
primewise theorem or by the cyclotomic-kernel
centralizer theorem.
The additional restrictions in
Propositions~\ref{A:full-support-normalizer} and
\ref{A:cm-normalizer-field} below settle cofinite prime support
and all normalizer fields containing no CM subfield. Over an
imaginary quadratic normalizer field the residual support
can contain only split primes. These restrictions apply
to the actual full profinite subgroup in this remark.

The actual attacks and their exact stopping
points are as follows.  Lower central
quotients and all maximal pro-prime quotients
leave precisely $R$.  Passing to a finite
cover detects every nonidentity element in
homology, but its line can have weight one;
the proof does not make that character
cyclotomic.  For a CM detecting factor the
restriction to $G_{K^{\rm ab}}$ is trivial,
so replacing $K(\mu_\infty)$ by $K^{\rm ab}$
invalidates the no-torus-to-abelian-variety
hypothesis: both $T_\ell E_0$ and
$\mathbf Z_\ell(1)$ become trivial.

The Prym refinement has also been checked
directly.  For $Y\to U$ finite \'etale,
an element in $R$ has proper homology in
$\ker(V_\ell J_{\overline Y}\to V_\ell E_0)$.
In genus one this kernel is zero, so a
weight-one residual detector requires genus
at least two.  Sending its Jacobian factor
to an elliptic curve does not automatically
give a punctured map.  A map
$f:\overline Y\to E'$ restricts to
$Y\to E'\setminus\{q\}$ precisely when
$\operatorname{Supp}f^{-1}(q)$ lies in the
existing boundary.  Riemann--Hurwitz then
forces $\deg f\leq\deg(Y/U)$.  The concrete
map $[3]:E_0\to E_0$ on the cover
$E_0\setminus E_0[2]\to U$ still surjects
onto all of $E_0$, although it detects
proper Tate homology.  Thus it supplies
no such punctured map.  Adding punctures
requires an arithmetic-equivariant lift;
the class-five obstruction above is an
explicit case in which that lift fails.

A further approach using a normalizer
extension as a source for the Hom theorem
stops at the same precise issue: it is
an infinite-index arithmetic subgroup,
not the fundamental group of a finite-type
source for which the required open Hom
theorem applies.  Its projection onto an
open Galois group supplies no section by
itself.  The finite metabelian example
above confirms, by an actual mixed-prime
element, why primewise tests leave
additional information.  It is not a
counterexample inside $\pi$ to the
normalizer assertion.  A genuinely new
argument on the residual Prym part or
a compatible full profinite construction
is still required.
\end{statement}

\subsection{Slimness from global fields and the positive-degree Hom reduction}

\begin{statement}{A.23}{Proved}{an elementary local embedding lemma}
\label{st:A.23}
\label{A:local-embedding}
Let $k$ be a global field and let $M/k$ be a finitely
 generated separably generated extension.  There is a
nonarchimedean place $v$ of $k$ and an embedding
$M\hookrightarrow k_v$ fixing $k$.  The valuation on $M$
pulled back from $k_v$ has completion $k_v$.  In particular
no nonidentity $k$-automorphism of $M$ preserves this
valuation.
\end{statement}
\begin{proof}
Choose a smooth integral $k$-variety with function field
$M$ and an etale coordinate chart of dimension
$d=\operatorname{trdeg}(M/k)$.  A rational point of the
nonempty open image of this chart gives a closed point
with finite separable residue field $L/k$.  Choose a
nonarchimedean place of $k$ splitting completely in the
Galois closure of $L/k$.  Chebotarev supplies such places.
The chosen point is then a $k_v$-point.  The etale
implicit-function theorem gives a neighborhood of it
mapping bijectively to an open ball in $k_v^d$.
Choose coordinates in that ball algebraically independent
over $k$.  Such coordinates exist successively because
$k$ and each finitely generated extension previously
chosen are countable whereas each nonempty open ball
is uncountable.  The resulting evaluation has zero
kernel on the function field and embeds $M$ in $k_v$.
For $d=0$ the splitting place directly supplies the
embedding.  Since $k$ is dense in $k_v$, the completion
of $M$ is $k_v$.  An automorphism preserving the valuation
extends continuously to that completion and fixes its
dense subfield $k$, so is the identity.
\end{proof}

\begin{statement}{A.24}{Proved}{the exact finite-type slimness statement}
\label{st:A.24}
\label{A:slim}
For $K$ finitely generated over $\mathbf Q$, the
centralizer in $G_K$ of every open subgroup is trivial.
The same holds for a field finitely generated over
$\mathbf F_p$ of positive transcendence degree.  For a
finite field it is false: $G_{\mathbf F_q}=\widehat{
\mathbf Z}$ and every open subgroup has centralizer
$\widehat{\mathbf Z}$.
\end{statement}
\begin{proof}
Choose a global subfield $k\subset K$ over which $K$
is separably generated: take $k=\mathbf Q$ in
characteristic zero and $k=\mathbf F_q(t)$ with $t$
in a separating transcendence basis in positive
characteristic.  Suppose $\sigma$ centralizes an open
subgroup and $\sigma\ne1$.  Choose a finite Galois
extension $M/K$ such that $G_M$ is contained in that
open subgroup and $g=\sigma|_M\ne1$.  Enlarge it to
contain $\mu_\ell$ for an odd prime $\ell\ne
\operatorname{char}K$.  Lemma~\ref{A:local-embedding}
gives a discrete valuation $w$ on $M$ moved by $g$.
Its value group is discrete and nonzero, so normalize
it to $\mathbf Z$.  Weak approximation for the two
distinct valuations $w$ and $w\circ g$ gives $a\in M^*$
with
\[
             w(a)=1,\qquad w(g(a))=0.
\]
On $H^1(G_M,\mu_\ell)=M^*/M^{*\ell}$, conjugation
by $\sigma$ acts by the scalar $\chi_\ell(\sigma)$:
indeed it fixes every element of $G_M$ by assumption
and acts on the coefficients by that scalar.  Naturality
of Kummer theory therefore gives
\[
       [g(a)]=\chi_\ell(\sigma)[a]
                   \quad\text{in }M^*/M^{*\ell}.
\]
Applying $w$ modulo $\ell$ contradicts $0=
\chi_\ell(\sigma)$ in $\mathbf F_\ell$, since the
cyclotomic value is a unit.  This proves the two
positive assertions using only local embeddings,
Kummer theory, weak approximation and the global-field
splitting-prime theorem.  The finite-field assertion
follows from its Frobenius description.
\end{proof}

\begin{statement}{A.25}{Proved}{positive peripheral degrees force an open image}
\label{st:A.25}
\label{A:positive-degree-open}
Let $P'$ and $P$ be the profinite geometric groups of
connected hyperbolic punctured curves in characteristic
zero, with nonempty boundary.  Choose positively oriented
peripheral generators $c'_j$ and $c_i$ from surface
presentations.  A continuous homomorphism $h:P'\to P$
satisfying
\[
 h(c'_j)=b_j^{-1}c_{\tau(j)}^{d_j\rho}b_j,
 \qquad d_j\in\mathbf Z_{>0},\quad
 \rho\in\widehat{\mathbf Z}^{\times}
\]
has open image.  The same multiplier $\rho$ is required
at every cusp; this is exactly the multiplier supplied
by the base path and the Kummer condition in a
morphism of $\mathcal B$.
\end{statement}
\begin{proof}
First suppose the target has at least two cusps.
In $P^{\rm ab}$ the peripheral classes have the unique
relation $\sum_i[c_i]=0$ and span a free direct summand
of rank one less than the number of cusps.  Apply $h$
to $\sum_j[c'_j]=0$.  Since $\rho$ is a unit, the
positive integer totals
\[
                   D_i=\sum_{\tau(j)=i}d_j
\]
are all the same integer $D$.  Some source cusp exists,
so $D>0$; in particular every target cusp is hit.

Put $H=h(P')$, which is closed by compactness, and
let $Q\subset P$ be any open subgroup containing $H$.
Its corresponding finite covering has cusps $y$ above
the original cusps.  Write $e_y$ for the ramification
index at $y$.  A source cusp determines a unique lifted
target cusp: its nontrivial image lies in the
intersection of $Q$ with a conjugate of the original
inertia, and distinct lifted inertia groups have
trivial intersection.  For this lifted cusp,
$e_y\mid d_j$ because $\rho$ is a unit, and the local
exponent becomes $(d_j/e_y)\rho$.  The same boundary
relation in $Q^{\rm ab}$ forces the totals
\[
             D_y'=\sum_{j\mapsto y}d_j/e_y
\]
to one integer $D_Q>0$.  This equality also includes
cusps not initially known to be hit: their zero total
would force all totals to be zero, contradicting a
positive source exponent.  For any original cusp $i$,
\[
 D=\sum_{y\mid i} e_yD_Q
   =D_Q[P:Q],
\]
since the ramification indices in its finite fiber sum
to the covering degree.  Consequently $[P:Q]\le D$.
A closed subgroup whose overgroups of finite index all
have bounded index has finite index: it is the
intersection of those overgroups and the indices of
finite quotients of $P/H$ are bounded.  Hence $H$ is
open.

If the target has one cusp, its genus is positive.
Choose a degree-two abelian cover, using a nonzero
homology character $P\to\mathbf Z/2$.  Its peripheral
generator is a product of commutators, so maps to zero;
the covering therefore has two cusps.  Restrict $h$ to
the finite-index preimage of the covering group in
$P'$.  This is the group of a finite covering of the
source, with nonempty boundary.  A lifted source
exponent is $d_j e'_j/e_y$, a positive integer, where
$e'_j$ is the source ramification index.  The preceding
argument proves openness of the restricted image and
hence of the original image.
\end{proof}

\begin{statement}{A.26}{Known}{the exact dominant relative Hom-form in the boundary category}
\label{st:A.26}
\label{A:relative-hom}
Let $K$ be sub-$p$-adic and let the source and target
be hyperbolic curves.  A morphism of the groupoid
systems of $\mathcal B$ with open geometric image is
induced by a unique dominant $K$-morphism of curves
satisfying its specified boundary condition.  Its local
integers $d_j$ are the actual ramification indices.
For case IV(d), where both boundaries are nonempty,
all morphisms of $\mathcal B$ have open geometric
image, so this proves the entire boundary-decorated
case.  The result applies to full profinite groups and,
for the open-image assertion of the Hom theorem, to
geometrically pro-$p$ groups.
\end{statement}
\begin{proof}
After a finite extension splitting all cusps, use
$\alpha$ and a target arithmetic conjugation to pin the
base map.  In the boundary formula of
Proposition~\ref{A:two-morphisms}, write
$a_j=e_jb_j$ with $e_j$ in the target cusp decomposition
group having the same arithmetic projection as $a_j$
and $b_j$ geometric.  Conjugation by $e_j$ removes
$\chi(p(a_j))$ from the exponent.  Thus the geometric
map has the displayed form in
Proposition~\ref{A:positive-degree-open}, with the
single factor $\chi(\alpha)$, which is one in the
pinned presentation.  This proves openness in case
IV(d), including the one-cusp case.  Openness descends
through the finite base extension.

In all stated open cases Mochizuki's Theorem A
\cite{A_Mochizuki1999} gives the unique dominant map.
For smooth curves a dominant map extends uniquely to
the smooth projective compactifications.  Its local
map at a cusp is $t\mapsto u(t')^e$, with $e>0$;
Kummer theory therefore sends its oriented inertia
generator to the $e$th power, with the same base-path
factor as in Proposition~\ref{A:two-morphisms}.
Equality of the induced homomorphisms makes $e=d_j$
in $\widehat{\mathbf Z}$ and hence in $\mathbf Z$.
A source cusp cannot map to an interior point, since
its inertia would then be killed, contrary to $d_j>0$.
An interior source point cannot map to a deleted target
point because the map is defined on the source open
curve.  Thus the inverse image of the target boundary
is exactly the source boundary, as required by the
literal category.  The boundary decomposition map,
its conjugator and its Kummer map are recovered up to
the unique natural transformation of
Proposition~\ref{A:boundary-recovery} below.  This
proves full faithfulness in the specified category,
not merely equality of its underlying group maps.
\end{proof}

\subsection{Recovery of the Kummer decoration and the surface multipliers}

\begin{statement}{A.27}{Proved}{recovery of the boundary morphism without rational cusps}
\label{st:A.27}
\label{A:boundary-recovery}
For geometric boundary systems over a field finitely
generated over $\mathbf Q$, the geometric-conjugacy
class of the global homomorphism $f:E'\to E$ in a
morphism of $\mathcal B$ determines that morphism up
to its unique natural transformation.  No rationality
assumption on the cusps is needed.  For every cusp
and every positive integer $d$,
\[
                   \operatorname{Transp}_E(L_i^d,L_i)=E_i.
\]
\end{statement}
\begin{proof}
The transporter equality was proved in
Proposition~\ref{A:cyclotomic-normalizer}: any
transporter gives two cusp inertia groups with a
nontrivial intersection, so their lifted cusps coincide;
the resulting stabilizer is exactly the decomposition
group $E_i$.  This argument uses the decomposition
extension over $K(i)$ and does not require $K(i)=K$.

The image of each source inertia is nontrivial by its
positive degree.  It consequently determines its
unique target cusp inertia, up to the relevant
conjugacy, and hence the map of closed cusp orbits.
Since $E'_j$ normalizes that source inertia, its image
normalizes the nontrivial subgroup of the determined
target inertia.  The transporter assertion puts it
in the associated decomposition group.  Thus, once a
conjugator $a_j$ is chosen, the local homomorphism is
forced to be $\rho_{\tau(j)}^{-1}
\operatorname{Int}(a_j)f\rho'_j$.  Two choices of
$a_j$ differ on the left by an element of
$E_{\tau(j)}$; this is exactly $\mu_j$ in (A1).
Their arithmetic projections change the cyclotomic
factor by that same conjugation, so comparison of
$k_i$ gives the same positive integer $d_j$.

Two base paths for the same $f$ differ by an element
centralizing the open image of $p'$.  Proposition~
\ref{A:slim} says that they coincide.  Changing $f$
by geometric conjugation is the global $\mu$ in
(A1).  These statements construct the required natural
transformation, and Lemma~\ref{A:geometric-centralizer}
proves its uniqueness.  The argument simultaneously
recovers the residue-field homomorphisms, since they
are the projections of the forced local maps; no
extra choice appears at a nonrational cusp.
\end{proof}

\begin{statement}{A.28}{Proved}{the exact ambiguity of the Kummer data}
\label{st:A.28}
\label{A:kummer-ambiguity}
A loop structure with cyclotomic compatibility determines
its equivariant Kummer isomorphisms up to one independent
factor in $\widehat{\mathbf Z}^{\times}$ for each
arithmetic orbit of geometric cusps.  The set of such
choices is a torsor under that product.  Requiring a
common positively oriented surface presentation reduces
this to the diagonal factor
$\widehat{\mathbf Z}^{\times}$ when the boundary is
nonempty.

For a proper hyperbolic curve the replacement is its
trace orientation
\[
 \operatorname{tr}:H^2(\pi,\widehat{\mathbf Z}(1))
                       \xrightarrow{\sim}\widehat{\mathbf Z}.
\]
Equivariant orientations form a
$\widehat{\mathbf Z}^{\times}$-torsor.  A dominant
map of proper curves has pullback compatible with the
traces by multiplication by its positive degree.
For a map from an affine curve to a proper curve this
condition is imposed on the filled source group
$\bar\pi'=\pi'/\overline{\langle\!\langle L'_j
\rangle\!\rangle}$; equivalently it is the relative
orientation of the source group with its boundary.
\end{statement}
\begin{proof}
Two continuous isomorphisms from the rank-one module
$T$ to a fixed inertia group differ uniquely by a unit
of $\widehat{\mathbf Z}$.  Every such multiplication
commutes with the cyclotomic action.  Transport by
arithmetic conjugation forces exactly equality of
these scalars within a geometric cusp orbit, and
imposes no relation between different orbits.  This
proves the first torsor assertion by an explicit action
and its inverse.

In abelianization the peripheral classes have the
single primitive relation $\sum_i[c_i]=0$.
If replacement generators $c_i^{u_i}$ extend to a
surface presentation, their same relation forces
$u_1=\cdots=u_\nu$ for $\nu\ge2$.  For one cusp
there is already one scalar.  Conversely every common
unit is realized by Proposition~\ref{A:all-multipliers}
below, proving both necessity and sufficiency.

For a proper curve the comparison with a compact
oriented surface, or etale Poincare duality, identifies
$H^2(\pi,T)$ with $\widehat{\mathbf Z}$ and makes
the trace equivariant for the trivial action on the
right.  Its generators differ by precisely the stated
units.  Pullback of the fundamental cohomology class
under a finite map multiplies its trace by the degree:
this follows on every finite coefficient ring from
the pullback of a closed point divisor, whose total
multiplicity is the degree, and then by passage to the
inverse limit.

A map to a proper curve kills each source cusp
inertia, because it extends at that cusp.  It hence
factors through $\bar\pi'$, the fundamental group
of the smooth compactification.  Apply the preceding
degree identity to that factor.  In relative group
cohomology the same orientation is the trace on
$H^2(\pi',\{L'_j\};T)$; the boundary homomorphism
from $\bigoplus_jH^1(L'_j,T)$ has trace the sum of
the local residues.  The surface presentation proves
this by the single boundary relation.  Filling the
cusps identifies its degree-two orientation with the
proper orientation, giving the asserted compatibility.
This extension of the category is necessary: in the
literal category of scan 35 an index map from a
nonempty source boundary to an empty target boundary
does not exist.
\end{proof}

\begin{statement}{A.29}{Proved}{every profinite unit is a surface multiplier}
\label{st:A.29}
\label{A:all-multipliers}
Let $2g-2+\nu>0$ and $\nu\ge1$.  In the profinite
surface presentation
\[
 P_{g,\nu}=\left\langle a_1,b_1,\ldots,a_g,b_g,
 c_1,\ldots,c_\nu\ \middle|\
 \prod_{j=1}^g[a_j,b_j]\prod_{i=1}^\nu c_i=1
                          \right\rangle^{\wedge},
\]
for every $\rho\in\widehat{\mathbf Z}^{\times}$
there is an automorphism taking every $c_i$ to a
conjugate of $c_i^\rho$.  Its images of the handle
generators, together with those conjugates, form a
privileged system.  When $g\ge1$ the same original
handle generators can be retained exactly for
$\rho=1$.  When $g=0$ all units are possible and
there are no handle generators to retain.
\end{statement}
\begin{proof}
Realize the surface group by a curve and a reduced
rational divisor of degree $\nu$ over $\mathbf Q$.
For $g=0$ use the projective line.  For $g\ge1$ use
$y^2=x^{2g+1}-1$, which is smooth of genus $g$ and
has a rational point at infinity.  Distinct rational
$x$-values away from the branch locus give reduced
divisors of degree two; their union, with the point
at infinity when needed, gives every positive $\nu$.
The cusps need not individually be rational.  The
cyclotomic character of $G_{\mathbf Q}$ is surjective,
so an element with value $\rho$ acts on each inertia
as the $\rho$th power and permutes the cusp labels.
Compose its outer action with a surface homeomorphism
realizing the inverse permutation of the cusps.
That homeomorphism preserves the orientation and has
multiplier one.  Any representative of the resulting
outer automorphism has the asserted conjugacy property.
Transporting it through the complex surface comparison
proves the statement about the profinite presentation
alone; no arithmetic structure remains in its conclusion.

Here is an explicit multiplier $-1$.  Put
$Q=\prod_j[a_j,b_j]$ and $C=\prod_i c_i$.  The
substitution
\[
 R(a_j)=b_{g+1-j},\qquad R(b_j)=a_{g+1-j},\qquad
 R(c_i)=c_{\nu+1-i}^{-1}
\]
sends $Q$ to $Q^{-1}$ and $C$ to $C^{-1}$.
Since $C=Q^{-1}$, it respects the relation and
$R^2=1$.  Adjacent Hurwitz substitutions
\[
 h_i(c_i)=c_{i+1},\quad
 h_i(c_{i+1})=c_{i+1}^{-1}c_i c_{i+1}
\]
fix all other generators and preserve $C$.
Their product
$B=h_1(h_2h_1)\cdots(h_{\nu-1}\cdots h_1)$,
with the rightmost factor applied first, reverses
the cusp labels.  Thus $BR$ takes every $c_i$ to
a conjugate of $c_i^{-1}$ and supplies the requested
explicit element.

For the restriction on the handles, pass to the
class-two quotient in which all $c_i$ are central.
There $Q$ has class
$\omega=\sum_j a_j\wedge b_j$ in
$\bigwedge^2\widehat{\mathbf Z}^{2g}$, and the
relation is $C=Q^{-1}$.  The vector $\omega$ is
primitive and has no nonzero annihilator in
$\widehat{\mathbf Z}$.  Keeping all handle generators
fixed leaves $Q$ fixed, whereas replacing each cusp
by a conjugate of its $\rho$th power sends $C$ to
$C^\rho$.  Their relation therefore gives
$(1-\rho)\omega=0$, and hence $\rho=1$.
The identity realizes that value.  For $g=0$ the first
part already proves every value with no handles.
\end{proof}

\begin{statement}{A.30}{Proved}{the geometric systems satisfy all five conditions}
\label{st:A.30}
\label{A:five-conditions}
For a smooth hyperbolic curve over a field finitely
generated over $\mathbf Q$, its system satisfies
(a)--(e) of scans 82--88, in their literal sense:
a section after a finite base extension; split local
decomposition extensions; cyclotomic compatibility;
the privileged surface presentation; and, in every
finite cover, purity of weight one and independence
of $\ell$ on homology after killing boundary homology.
Condition (c) is already a consequence of the local
Kummer equivariance and the surjectivity of the local
projection after splitting the cusp fields.  When the
boundary is nonempty, (b) implies (a).
\end{statement}
\begin{proof}
A closed point of the nonempty open curve has finite
residue field over the base and gives a section after
that extension; it is an interior point and therefore
also gives the extra admissibility intended in (a).
At a cusp, a uniformizer defines the local section
and the Kummer description gives (b).  After the cusp
field is made rational, any $g\in E$ has the same
arithmetic image as some $e\in E_i$.  Writing
$g=\alpha e$ with $\alpha\in\pi$ gives
\[
 \operatorname{Int}(g)k_i(t)=
       \operatorname{Int}(\alpha)k_i(\chi(p(g))t),
\]
which is exactly (c), with the conjugator genuinely
geometric.  Finite base change gives the general
orbit form.  A section of $E_i$ composed with its
inclusion in $E$ gives the implication (b)$\Rightarrow$
(a), without its additional admissibility adjective.

An embedding of the base field in $\mathbf C$ and
the comparison theorem give (d) and its oriented
presentation.  Every finite arithmetic covering
subgroup comes from a finite etale cover; its
compactification has homology equal to the quotient
by the sum of its cusp homology.  This homology is
the Tate module of its Jacobian.  Spread out to a
good-reduction model.  At every closed point of good
reduction its characteristic polynomial is integral,
independent of $\ell$ different from the residue
characteristic, and all eigenvalues have absolute
value the square root of the residue cardinality.
These are exactly the requested clauses of (e).
\end{proof}

\begin{statement}{A.31}{Refuted}{two proposed redundancies among the five conditions}
\label{st:A.31}
\label{A:condition-countermodels}
In the abstract group-data setting, (a), (b), (c),
and (d) do not imply (e).  Nor do (a), (b), (c), and
(e) imply (d).
\end{statement}
\begin{proof}
For the first assertion take the proper genus-two
profinite surface group $P$, the empty boundary,
and $E=P\times G_K$ over a number field.  The product
section gives (a), (b) and (c) are vacuous, and the
surface presentation gives (d).  On $P^{\rm ab}$
every Frobenius is the identity, so its eigenvalues
have absolute value one and violate weight one.

For the second, take $P=T=\widehat{\mathbf Z}(1)$,
$E=T\rtimes G_K$, and one declared cusp with
$E_1=E$, $L_1=P$, and $k_1$ the identity.  The
splitting and the cyclotomic action give (a)--(c).
In every finite covering the inertia spans its entire
geometric abelianization, so the quotient required
in (e) is zero; its purity and independence statements
hold.  A surface presentation with one cusp has
abelian rank $2g$, whereas $P$ has rank one.  No
integer $g$ gives (d).
\end{proof}

\begin{statement}{A.32}{Open}{the remaining independence lattice}
\label{st:A.32}
\label{A:condition-lattice-open}
The implications and two countermodels just proved
settle those specific edges of the requested lattice.
A proof deciding every remaining implication among
(a)--(e), with the precise ambient axioms fixed, has
not been obtained here.  In particular no assertion
that a general outer arithmetic action virtually lifts
to an effective action is being inferred from the
weight condition.  The existence of a section after
finite extension is extra data in that question.
\end{statement}

\subsection{The added two-section stabilizer task}

To define the action in this task, retain the standard faithful
embedding $G_{\mathbf Q}\hookrightarrow\operatorname{Out}(F)$
and put
\[
 \mathcal A=\operatorname{Aut}(F)
       \mathop{\times}_{\operatorname{Out}(F)}
                   \widehat{\mathrm{GT}}_0.
\]
Its kernel is $F$, since $Z(F)=1$.  The pullback over
$G_L$ is the arithmetic fundamental group of the
tripod over $L$.  For an algebraic point $z$, write
$[S_z]$ for the $F$-conjugacy class of the image of
its section over its field of definition, together
with that projected arithmetic domain.  Conjugation
by a lift in $\mathcal A$ gives an action of
$\widehat{\mathrm{GT}}_0$ on these ambient subgroup
classes: two lifts differ by $F$, so the action is
independent of the lift.  This definition does not
assume that every element normalizes $G_{\mathbf Q}$.

\begin{statement}{A.33}{Proved}{the stabilizer of the two specified point sections}
\label{st:A.33}
\label{A:two-section-stabilizer}
Let $z$ be one chosen root of $z^2-z+1=0$ and
$L=\mathbf Q(z)=\mathbf Q(\sqrt{-3})$.  For the
ambient action just defined, the simultaneous stabilizer
of $[S_{1/2}]$ and $[S_z]$ is exactly
\[
                         G_L.
\]
It is contained in $\ker\Delta$ for the character
of LM7\_PUSH E.65.  It does not contain $C^\varphi$:
complex conjugation belongs to $C^\varphi$ and
exchanges the two distinct root sections.
\end{statement}
\begin{proof}
First we verify the arithmetic normalizer used in the
argument.  An outer automorphism commensurating the
faithful image of $G_{\mathbf Q}$ gives an isomorphism
between absolute Galois groups of two number fields.
The prescribed-isomorphism theorem of Uchida
\cite{A_Uchida1979} realizes it by an automorphism
$\sigma$ of $\overline{\mathbf Q}$.  Multiplication
by $\sigma^{-1}$ then centralizes $G_M$ for one
number field $M$.  Proposition~\ref{A:tripod-commutant}
identifies this centralizer with the geometric $S_3$.
Thus the full outer commensurator is
$G_{\mathbf Q}S_3$, and its subgroup preserving each
cusp label is exactly $G_{\mathbf Q}$: a nonidentity
member of that $S_3$ moves a label.  This supplies
the complete reduction behind LM3b D.8.

An element fixing $[S_{1/2}]$ normalizes its projection,
which is all of $G_{\mathbf Q}$.  The preceding result
puts it in $G_{\mathbf Q}$.  Arithmetic conjugation
transports $[S_z]$ to $[S_{\sigma z}]$.  Since $L$
is normal, both sections have domain $G_L$.
Injectivity of the point-section map makes the classes
equal exactly when $\sigma z=z$.  For these two
points this injectivity can be seen already by Kummer
theory: their ratio is a nontrivial third root of
unity, which is not a cube in $L$, since $L$ does
not contain a ninth root of unity.  Hence their
classes under the coordinate unit $t$ differ in
$H^1(L,\mu_3)$.  The stabilizer is therefore $G_L$.
The same conclusion holds for germs: after any finite
extension a sufficiently high power of three still
detects the two distinct roots.  If the phrase means
an unordered pair of these two classes with their
domains, it gives the same answer.  Interchanging
these domains would give a commensurator and hence
an arithmetic element, which cannot conjugate
$G_{\mathbf Q}$ onto its proper subgroup $G_L$.

We use only the arithmetic-vanishing part of the
character theorem E.65, not the existence assertion
E.58.  Its proof and the supplied referee's repair
were checked as follows.  On multiplier one, the two
frame signs are the restrictions to the same fourth
root $r=(3/4)^{1/4}$.  At $1/2$, the Legendre curve
is the quadratic twist by $2$ of $y^2=x^3-x$; at
$z$, it is the twist of $y^2=x^3+1$ by
$1/\sqrt{-3}$.  Over $\mathbf Q^{\rm ab}$ the first
twist is trivial, and the second twisting character
is $\sigma(r)/r$, because
$(-3)^{1/4}=\zeta_8\sqrt2\,r$.  The rational
three-torsion line of $y^2=x^3+1$ gives that same
sign; the projective three-torsion of $y^2=x^3-x$
gives the sign from its division polynomial.
The two signs therefore add to zero in $\mathbf F_2$.
The extension formula
$\Delta(a)=\delta(a\sigma^{-1})$, for an arithmetic
$\sigma$ with $\chi(\sigma)=\lambda(a)$, then gives
$\Delta|_{G_{\mathbf Q}}=0$.  In particular it
vanishes on the stabilizer just computed.

For the last assertion, write
$B=\widehat{C_2*C_3}=\langle s,r\mid s^2=r^3=1\rangle$.
The exact map of cote 140-3, scan 229, is
\[
 \varphi(c_0)=sr,\qquad \varphi(c_1)=s,
              \qquad\varphi(c_\infty)=r^{-1}.
\]
It is induced by the open immersion over $\mathbf Q$
from $\mathbf P^1_w\setminus\{0,1,\infty\}$ into
the orbifold with cusp $0$ and stacky points of orders
$2$ and $3$ at $1$ and $\infty$.  This orbifold is
$[U_0/S_3]$, with quotient coordinate
\[
        w=\frac{1728}{j}
          =\frac{27\lambda^2(1-\lambda)^2}
                  {4(\lambda^2-\lambda+1)^3}.
\]
The local ramification orders and the peripheral
product relation give precisely the three displayed
images, not merely a map with the same target.
Functoriality over $\mathbf Q$ supplies exterior
commutation for every arithmetic automorphism.
If its initial conjugator $b$ lies in $B$, the
pure-conjugator convention of equation (5) is obtained
as follows.  Changing the representative by
$\operatorname{Int}(a)$ with $a\in F$ changes $b$
to $\varphi(a)ba^{-1}$.  The map
$F\xrightarrow{\varphi}B\to S_3$ is onto, so
choose $a$ making the latter conjugator lie in $F$.
This proves membership in exactly the subgroup
$C^\varphi$ defined there.  In particular complex
conjugation belongs to it.  It exchanges $z$ and
$\bar z$, whose section classes were just shown to
be distinct, and so is outside $G_L$.
\end{proof}

\begin{statement}{A.34}{Proved}{the different stabilizer of the full root orbit}
\label{st:A.34}
\label{A:root-orbit-stabilizer}
Preserving $[S_{1/2}]$ and the unordered set of
\emph{both} root classes instead gives
$G_{\mathbf Q}$, by the same proof.  This is a
different specified set, related to the preceding
stabilizer by the exact index-two inclusion
$G_L\subset G_{\mathbf Q}$.  The result above answers
the task for one chosen root.  No conclusion that
$C^\varphi\subset G_{\mathbf Q}$ is being inferred
for this alternative task.
\end{statement}

\subsection{A further reduction of the full profinite common-kernel problem}

For a profinite group $D$, write
$R_{\rm nil}(D)=\bigcap_{j\ge1}\gamma_j(D)$,
where the lower central subgroups are closed.  This is
also the kernel of its maximal pronilpotent quotient.
For the geometric group $F$ of $E_0\setminus\{O\}$
put $P_n=\ker(F\to F^{\rm ab}/nF^{\rm ab})$.
The cover with group $P_n$ is
$E_0\setminus E_0[n]\xrightarrow{[n]}E_0\setminus\{O\}$.

\begin{statement}{A.35}{Proved}{a weight-two normalizer test in an isogeny cover}
\label{st:A.35}
\label{A:isogeny-weight-test}
Let $E/K$ be an elliptic curve over a number field,
$D$ the maximal pro-$\ell$ group of $E$ minus a
nonempty finite divisor, and $B\ne1$ a procyclic
subgroup of $D$ normalized by a closed arithmetic
subgroup with open image in $G_K$.  If the image of
$B$ in $T_\ell E$ is zero, the normalizer acts on
$B$ cyclotomically after a finite base extension.
\end{statement}
\begin{proof}
Make all punctures rational by a finite extension.
Use free pro-$\ell$ generators $x,y,c_1,\ldots,c_{r-1}$,
with the remaining cusp determined by the surface
relation.  Give $x,y$ weight one and each $c_i$
weight two in the integral Magnus expansion.  The
Magnus embedding separates the group, so a nonidentity
generator $b$ of $B\simeq\mathbf Z_\ell$ has a
first nonzero weighted degree $d$.  Its zero image in
$T_\ell E$ makes $d\ge2$.  The arithmetic action
preserves this filtration.  Its action on weight-one
generators is $T_\ell E$; its action on every cusp
generator in weight two is $\chi_\ell$.  Consequently
the eigenvalues of a good Frobenius on weight $d$
all have complex absolute value $q^{d/2}$: this
follows already in the finite free module of
homogeneous noncommutative polynomials, where the
eigenvalues are products of those on the generators.
Inner conjugation acts trivially on the first
nonzero weighted term.  If the normalizer character
is $\theta$, that term is a $\theta$-eigenvector.

There is also a finite pro-$\ell$ cover detecting
$b$ in homology.  Choose an open normal subgroup
$M$ for which $b$ survives in $D/M$, and use
$J=M\overline{\langle b\rangle}$.  The cyclic
quotient $J/M$ detects its homology class; the
abelianization of the finite covering group is free
over $\mathbf Z_\ell$, so this class remains
nonzero after tensoring with $\mathbf Q_\ell$.
Descend the cover and intersect its arithmetic group
with the original normalizer.  The projection remains
open and its homology vector again has character
$\theta$.  Good Frobenius eigenvalues on this covering
homology have absolute values $q^{1/2}$ or $q$.
Choose one good Frobenius common to both calculations.
It forces $d=1$ or $2$, and hence $d=2$.

In degree two, a group element with zero linear part
has leading term
\[
                a[X,Y]+\sum_i b_i C_i.
\]
For completeness, the group-like identity in the
Magnus Hopf algebra makes its quadratic part primitive.
The coefficients of $X\otimes X$ and $Y\otimes Y$
kill the $X^2,Y^2$ terms, and the coefficient of
$X\otimes Y$ makes the two mixed coefficients
opposite.  This argument uses torsion-freeness of
$\mathbf Z_\ell$ and works also at $\ell=2$.
On $[X,Y]$ the action is
$\det(T_\ell E)=\chi_\ell$ by the Weil pairing,
and every $C_i$ has that same character.  The
nonzero degree-two vector is torsion-free, so
$\theta=\chi_\ell$ on the indicated open Galois
subgroup, as asserted.
\end{proof}

\begin{statement}{A.36}{Proved}{the residual normalizer lies deeper than the common pro-prime kernel}
\label{st:A.36}
\label{A:deeper-common-kernel}
Let $F=\pi_1((E_0\setminus\{O\})_{\overline K})$
for $E_0:y^2=x^3-x$ and $K=\mathbf Q(i)$.
A nontrivial closed procyclic subgroup with open
arithmetic normalizer is either contained in cusp
inertia or contained in
\[
                         R_{\rm nil}(F').
\]
The latter group is strictly smaller than
$R_{\rm nil}(F)=\ker(F\to\prod_\ell F^{(\ell)})$.
More precisely,
\[
 R_{\rm nil}(F')
   =\bigcap_{n\ge1,\ \ell\ \mathrm{prime}}
                   \ker(P_n\longrightarrow P_n^{(\ell)}),
                                                               \tag{A7}
\]
where the intersection is taken inside $F'$.
Thus every subgroup still left by this problem dies
in every maximal pro-prime quotient of every
multiplication-isogeny covering of the punctured
elliptic curve, and in particular in the full
metabelian quotient of $F$.
\end{statement}
\begin{proof}
Proposition~\ref{A:cm-detected-normalizer} already
handles every subgroup detected by one $F^{(\ell)}$.
Suppose the remaining subgroup $A$ lies in
$R_{\rm nil}(F)\subset F'$.  It lies in every $P_n$.
If its image in some $P_n^{(\ell)}$ is nontrivial,
its $\ell$-component maps injectively to that
nontrivial procyclic image.  Its proper homology class
maps to zero under
\[
 T_\ell E_0\xrightarrow{[n]}T_\ell E_0,
\]
because $A$ dies in $F^{(\ell)}$.  This map is
injective, even when $\ell\mid n$, since the Tate
module is torsion-free.  Its proper homology class
is therefore zero.  Lemma~\ref{A:isogeny-weight-test}
shows that the normalizer character on this
$\ell$-component is cyclotomic after finite extension.
The full profinite recognition theorem in
Proposition~\ref{A:cyclotomic-normalizer}, including
its one-prime essential-cyclotomic version, puts that
component in one full inertia group.  Since $A$
centralizes its component, the inertia centralizer
puts all of $A$ in that inertia.  A nontrivial
inertia component survives in $F^{(\ell)}$, a
contradiction.  Thus $A$ lies in the right side of
(A7).

Here is a finite-quotient proof of (A7).  Every
map $P_n\to P_n^{(\ell)}$ kills
$R_{\rm nil}(F')$, since its restriction to $F'$
has pro-$\ell$ image.  Conversely an element of
$F'$ not in $R_{\rm nil}(F')$ survives in a finite
quotient $G$ of $F/R_{\rm nil}(F')$; the group
$G'$ is nilpotent.  Choose $n$ divisible by the
exponent of $G/G'$.  The image of $P_n$ is exactly
$G'$: it is contained there by that choice of $n$
and contains it because $F'\subset P_n$.
A Sylow projection of the finite nilpotent group
$G'$ detects the selected element and factors
through $P_n^{(\ell)}$.  This proves the reverse
inclusion.

Finally the explicit quotient
$G=\mathbf F_{25}\rtimes C_3$ of
Proposition~\ref{A:finite-metabelian} has
$G'=\mathbf F_{25}$ abelian, and
$\gamma_\infty(G)=\mathbf F_{25}\ne1$.
For a continuous finite quotient, the image of an
intersection of the nested closed lower central
groups is the intersection of their images, by
compactness.  Thus the image of
$R_{\rm nil}(F)$ in $G$ is $\mathbf F_{25}$,
whereas the image of $R_{\rm nil}(F')$ is
$\gamma_\infty(G')=1$.  This proves strictness by
an actual quotient, not by comparison of names.
\end{proof}

\subsection{Absolute Hom and the precise local comparison}

\begin{statement}{A.37}{Proved}{an open arithmetic map over number fields preserves the geometric subgroup}
\label{st:A.37}
\label{A:absolute-geometric}
Let $X/K$ and $Y/L$ be hyperbolic curves over number fields.
Every continuous open homomorphism $f:E_X\to E_Y$ satisfies
$f(P_X)\subset P_Y$. Thus it induces a continuous open
homomorphism $\alpha:G_K\to G_L$.
\end{statement}
\begin{proof}
The subgroup $H=f(E_X)$ is open in $E_Y$. Its projection
$J=p_Y(H)$ is an open subgroup of $G_L$, hence the absolute
Galois group of a number field. The group
$N=p_Yf(P_X)$ is closed, normal in $J$, and topologically
finitely generated because $P_X$ is. Mochizuki's
Theorem 2.4 in \cite{A_Mochizuki2006} states that a
finitely generated closed normal subgroup of the absolute
Galois group of a number field is trivial. Its field is
exactly the finite extension corresponding to $J$, so it
applies to $N$ and gives $N=1$. Factoring $p_Yf$ through
$E_X/P_X=G_K$ gives $\alpha$; its image is $J$, hence open.
This is the exact argument of LM3 D.11, with the requisite
normality taken in $J$, rather than assumed in all of $G_L$.
\end{proof}

\begin{statement}{A.38}{Known}{absolute open Hom over number fields}
\label{st:A.38}
\label{A:absolute-hom}
For hyperbolic curves $X/K$, $Y/L$ over number fields, the
natural map from dominant morphisms of schemes $X\to Y$ to
continuous open homomorphisms $E_X\to E_Y$, modulo inner
automorphisms of $E_Y$, is bijective. The arithmetic input
beyond the classical relative theorem is Theorem 1.1 of the
supplied October 2026 paper \cite{A_OpenAI2026Hom}; it remains
an unrefereed input, not an independently verified theorem
of this return. Its requested use of Uchida's local
correspondence for nonsurjective open homomorphisms is valid.
\end{statement}
\begin{proof}
By Proposition~\ref{A:absolute-geometric}, an open $f$ induces
an open $\alpha:G_K\to G_L$. The supplied Theorem 1.1,
applied to the algebraic closures as solvably closed Galois
extensions, produces a unique equivariant embedding
$\bar j:\overline L\to\overline K$. An embedding between
algebraic closures of number fields is onto: its image is
algebraically closed and $\overline K$ is algebraic over it.
Hence $\alpha$ is injective and identifies $G_K$ with an
open subgroup of $G_L$. Its restriction is an embedding
$j:L\to K$.

Form the actual base change $Y_j=Y\times_{L,j}K$. Its
arithmetic group is
\[
 E_{Y_j}=E_Y\times_{G_L,\alpha}G_K,
 \qquad
 \widetilde f(e)=(f(e),p_X(e)).
\]
Projection $E_{Y_j}\to E_Y$ is injective and has open image,
since $\alpha$ does. Thus $\widetilde f$ is open. More
explicitly, injectivity of $\alpha$ gives
$f(P_X)=f(E_X)\cap P_Y$, an open subgroup of $P_Y$.
Mochizuki's Theorem A \cite{A_Mochizuki1999} now produces the
unique dominant $K$-morphism $X\to Y_j$ inducing
$\widetilde f$, and composition with $Y_j\to Y$ gives the
required scheme morphism. Conversely a dominant scheme
morphism restricts on algebraic constants to an embedding
$L\to K$ and has open arithmetic fundamental-group image.
The local inertia diagram of Proposition~\ref{A:typed-functor}
recovers all its Kummer maps and positive ramification degrees.

For uniqueness, conjugating $f$ in $E_Y$ changes $\bar j$
by an element of $G_L$ and leaves $j|_L$ unchanged. Once $j$
is fixed, any remaining arithmetic conjugator centralizes the
open image of $G_K$ in $G_L$, and is trivial by slimness.
The remaining conjugation is geometric, exactly the equivalence
relation in the relative Hom theorem. This proves bijectivity
with the stated inner-automorphism convention.

Here is the requested source check. On page 595 of
\cite{A_Uchida1981}, after stating the conjecture, Uchida
explicitly fixes a continuous homomorphism with \emph{open
image}; surjectivity is not imposed. On pages 595--596,
case (ii) of the local trichotomy associates a target prime
to a source prime whose pro-$\ell$ inertia survives, for
odd $\ell$, and puts its decomposition group in the target
decomposition group. Proposition 1, page 597, proves that
all but finitely many target primes occur. Its first paragraph
also proves compatibility with replacement of the target by
any finite extension inside the solvably closed extension.
These are exactly the hypotheses and conclusions used in
Theorem 2.4 of \cite{A_OpenAI2026Hom}. In particular, passage
to an open image changes the finite target field and is
already allowed in Uchida's argument. The finite exceptional
set in Proposition 1 is for a fixed $\ell$; no uniformity in
$\ell$ is being imported from it. This checks the step named
in the task, without treating it as a check of every later
argument in the 2026 main theorem.
\end{proof}

\begin{statement}{A.39}{Refuted}{Refuted in dimension zero; Known for local curve isomorphisms}
\label{st:A.39}
\label{A:local-absolute-isom}
There are finite extensions $k_1,k_2$ of $\mathbf Q_3$ with
isomorphic absolute Galois groups and no field isomorphism.
One explicit pair is
\[
 k_1=\mathbf Q_3(\zeta_3,\sqrt[3]{3}),\qquad
 k_2=\mathbf Q_3(\zeta_3,\sqrt[3]{4}).
\]
For hyperbolic curves over finite $p$-adic extensions, however,
the supplied Mochizuki--Tsujimura RNS preprint, Theorem D
or Theorem 3.12, gives the absolute Isom bijection
\[
 \operatorname{Isom}(X,Y)\ \xrightarrow{\sim}\
 \operatorname{OutIsom}(E_X,E_Y).
\]
These two assertions are related by the actual quotient maps
$E_X\to G_{k_1}$ and $E_Y\to G_{k_2}$: the first assertion
alone does not produce an isomorphism of their source groups.
\end{statement}
\begin{proof}
The polynomials $t^3-3$ and $(t+1)^3-4$ are Eisenstein at
$3$. Their cubic fields have discriminant valuations $5$ and
$3$, respectively. Each splitting field has group $S_3$ and
unique maximal abelian subfield $\mathbf Q_3(\zeta_3)$,
because its cubic discriminant has square class $-3$.
Both splitting fields have degree six. The theorem on page 2
of Jarden--Ritter \cite{A_JardenRitter1979} therefore gives
$G_{k_1}\simeq G_{k_2}$: for local fields containing
$\zeta_p$, its exact criterion is equality of absolute degree
and of the maximal abelian subfield over $\mathbf Q_p$.

The discriminant exponents of the splitting fields are $11$
and $7$. To check this directly, decompose the regular
representation of $S_3$ as
$1\oplus\operatorname{sgn}\oplus2\operatorname{std}$.
The quadratic character has conductor exponent one; the
standard character has conductor exponent equal to the
cubic discriminant exponent, since the cubic permutation
representation is $1\oplus\operatorname{std}$. The
conductor-discriminant formula gives $1+2\cdot5=11$ and
$1+2\cdot3=7$. All the polynomial and finite character
calculations are independently checked in
\texttt{code/A/local\_field\_certificate.py} and its output.
An abstract field isomorphism between finite $p$-adic fields
preserves their canonical henselian valuation, hence is
continuous and fixes $\mathbf Q_p$; the unequal discriminants
exclude it.

For the positive curve assertion, use the full set of primes
as $\Sigma$ in the supplied RNS Theorem 3.12
\cite{A_MochizukiTsujimuraRNS}. We checked the exact
reduction used there. Theorem 3.11(i) first reconstructs the
geometric subgroup and the arithmetic quotient. Finite
residue fields satisfy its cyclotomic-kernel hypothesis;
Theorem 2.17 gives $\Sigma$-resolution of nonsingularities.
Theorem 3.11(ii) consequently makes the isomorphism preserve
closed-point decomposition groups. The last line of the
proof of Theorem 3.12 applies the point-theoretic absolute
reconstruction theorem, [AbsTopII], Corollary 2.9, to obtain
the unique scheme isomorphism. These steps apply to full
profinite groups and retain the stated preprint status.
An isomorphism of the quotients $G_{k_i}$ in the first part
is not asserted to lift through $E_X\to G_{k_1}$ and
$E_Y\to G_{k_2}$; if such a lift existed for hyperbolic
curves, this theorem would supply a field isomorphism,
contradicting the discriminant calculation.
\end{proof}

\begin{statement}{A.40}{Open}{unrestricted absolute open Hom over local fields}
\label{st:A.40}
\label{A:local-absolute-hom-open}
The preceding theorem handles isomorphisms and, by passing
to an open target subgroup, injective open homomorphisms
between local arithmetic curve groups: the target subgroup
is the group of its finite etale cover, so the given map is
an isomorphism onto that group and Theorem 3.12 applies.
For a general open homomorphism with nontrivial kernel, the
number-field proof above stops at an exact missing step.
Its projected geometric image is a finitely generated normal
subgroup of an open local Galois group; the number-field
vanishing theorem no longer applies, because that open
local Galois group is itself finitely generated. The local
Isom reconstruction uses both an isomorphism and its inverse
in Theorem 3.11(i), and supplies no morphism for this step.
No counterexample of hyperbolic curves, and no proof of this
remaining unrestricted Hom assertion, is obtained here.
The finite dimension-zero example is explicitly not substituted
for such a counterexample.
\end{statement}

\subsection{Kernels beyond number fields and beyond the full profinite group}

\begin{statement}{A.41}{Proved}{constant curves and the unavoidable geometric kernel}
\label{st:A.41}
\label{A:constant-kernel}
Let $K/k$ be a regular finitely generated extension of a number
field, of positive transcendence degree, and let $U_0/k$ be a
hyperbolic curve. For $U=U_0\times_k K$ the kernel of the full
outer action is exactly
\[
 \ker(G_K\to\operatorname{Out}(P_U))
       =\ker(G_K\to G_k)=G_{K\overline k}.
\]
It is nontrivial. More generally, for every hyperbolic curve
$U/K$, with $K$ finitely generated over $\mathbf Q$ of positive
transcendence degree, the full outer-action kernel is nontrivial
and contains the inertia of every divisorial valuation at which
a smooth proper compactification with etale boundary extends
smoothly. It is contained in $G_{K\overline{\mathbf Q}}$.
\end{statement}
\begin{proof}
For the constant curve the actual homotopy exact sequences
identify its outer representation with the composite
\[
 G_K\longrightarrow G_k\xrightarrow{\rho_{U_0/k}}
                  \operatorname{Out}(P_{U_0}).
\]
The second arrow is injective by Hoshi--Mochizuki, Theorem C,
and the first is surjective because $K/k$ is regular. This
proves the asserted equality, including its field-theoretic
identification by Galois theory.

For the general assertion take a smooth geometrically integral
model $S/k$ of $K$, shrinking it so that the compactification
$X/K$ and its reduced boundary extend to a smooth proper
curve $\mathcal X/S$ and a finite etale divisor
$\mathcal D\subset\mathcal X$. The homotopy sequence for
$\mathcal X\setminus\mathcal D$ gives a continuous outer
monodromy map $\rho_S:\pi_1(S)\to\operatorname{Out}(P_U)$
and the exact factorization
\[
          \rho_{U/K}=\rho_S\circ q_S,
       \qquad q_S:G_K\twoheadrightarrow\pi_1(S).
\]
At a codimension-one point of $S$, residue characteristic is
zero. Smooth specialization for the proper pair therefore
identifies the entire geometric fundamental group, and its
inertia is killed by $q_S$. Such inertia is nontrivial:
a root of a uniformizer gives a ramified finite extension of
the henselian discretely valued field, with nontrivial image
of inertia. A positive-dimensional smooth model has such a
codimension-one point after taking an affine open. This also
proves nontriviality in the constant case by an explicit
ramified extension. Finally Hoshi--Mochizuki, Corollary 6.2(ii)
\cite{A_HoshiMochizuki2011}, applied with the full set of primes,
puts the kernel inside the kernel of $G_K\to G_{\mathbf Q}$,
which is $G_{K\overline{\mathbf Q}}$.
\end{proof}

\begin{statement}{A.42}{Proved}{the remaining kernel maps exactly to the congruence kernel}
\label{st:A.42}
\label{A:family-kernel-quotient}
For the model $S$ and the monodromy map in the preceding
proposition, let $N_S$ be the closed normal subgroup of
$G_K$ generated by divisorial inertia at the codimension-one
points of $S$. Then
\[
 N_S=\ker q_S,\qquad
 \ker\rho_{U/K}/N_S\ \xrightarrow{\sim}\ \ker\rho_S.
\]
For a universal hyperbolic curve over a fine level cover of
$\mathcal M_{g,r}$, the group on the right is precisely the
kernel of the geometric mapping-class monodromy
\[
 \widehat\Gamma_{g,r}[N]\longrightarrow
                         \operatorname{Out}(\widehat\pi_{g,r}).
\]
Here a level $N\ge3$ and a number field of definition are
chosen so that the moduli space and universal family are
schemes; the assertion is restricted to its positive-dimensional
components. Thus this is an actual quotient of the requested
function-field kernel, with its geometric meaning specified.
\end{statement}
\begin{proof}
A finite extension of the function field is represented by
its finite normalization over $S$. Since $S$ is regular,
purity of the branch locus says that this normalization is
etale exactly when it is unramified at all codimension-one
points. The finite quotients of $G_K/N_S$ are consequently
exactly the finite etale covers of $S$, proving the first
identity. The surjectivity of $q_S$ and
$\rho_{U/K}=\rho_Sq_S$ give the second displayed isomorphism;
its map sends the class of $g$ to $q_S(g)$ and its kernel is
exactly $N_S$.

In the universal-family case the geometric fundamental group
of the fine level moduli space is the indicated profinite
mapping-class group, and its action on the fibre is its
usual geometric outer monodromy. Hoshi--Mochizuki,
Corollaries 6.4 and 6.5 \cite{A_HoshiMochizuki2011}, identify
the kernel of the full arithmetic universal monodromy with
this geometric kernel: an arithmetic monodromy element that
also lies in the geometric image has trivial image in the
number-field Galois group. Restriction to the chosen fine
level cover preserves the same assertion. Applying the
already proved quotient isomorphism now gives exactly the
claimed group, with no inference from a name or analogy.
\end{proof}

\begin{statement}{A.43}{Proved}{every single-prime outer action has nontrivial kernel}
\label{st:A.43}
\label{A:proprime-kernels}
Let $K$ be a number field or a finite extension of
$\mathbf Q_p$, and let $U/K$ be hyperbolic. For every prime
$\ell$, including $\ell=2$, the outer action
\[
 G_K\longrightarrow\operatorname{Out}(P_U^{(\ell)})
\]
has nontrivial kernel. At every place of good reduction of
the compactified pair with residue characteristic different
from $\ell$, its inertia is contained in this kernel.
For a constant curve over a regular finitely generated
extension $K/k$, its kernel is the exact inverse image of
the corresponding kernel over $k$ under $G_K\to G_k$.
\end{statement}
\begin{proof}
The group $P_U^{(\ell)}$ is finitely generated and pro-$\ell$.
The kernel of the action of its automorphism group on the
finite Frattini quotient is pro-$\ell$: on every finite
characteristic $\ell$-group quotient the analogous kernel is
an $\ell$-group, and these quotients are cofinal. Passing to
outer automorphisms preserves the conclusion that the outer
automorphism group has an open pro-$\ell$ subgroup. This is
also Proposition 4.1 of \cite{A_Mochizuki2006}.

No open subgroup of $G_K$ is pro-$\ell$. In the local case
it surjects onto $\widehat{\mathbf Z}$ through the maximal
unramified extension. In the number-field case its field
admits a cyclotomic $\mathbf Z_q$-extension for every prime
$q\ne\ell$; the intersection with that tower changes only
its finite initial portion, so the quotient is still
$\mathbf Z_q$. An injective map into the displayed outer
automorphism group would make an open subgroup of $G_K$
pro-$\ell$, a contradiction. This proves nontriviality for
all primes without an exception at two.

The inertia assertion is the prime-to-residue-characteristic
smooth specialization theorem for the proper pair; its
geometric pro-$\ell$ group is unchanged, so the inertia
outer action is trivial. For a constant curve the exact
base-change composite used in Proposition~\ref{A:constant-kernel}
also holds after maximal pro-$\ell$ completion. Taking its
kernel gives the stated inverse-image formula.
\end{proof}

\begin{statement}{A.44}{Known}{the exact tripod kernel at odd regular primes}
\label{st:A.44}
\label{A:regular-prime-kernel}
Let $U_0=\mathbf P^1_{\mathbf Q}\setminus\{0,1,\infty\}$.
For an odd regular prime $\ell$, let $\Omega_\ell$ be the
maximal pro-$\ell$ extension of $\mathbf Q(\mu_\ell)$
unramified outside $\ell$. Then
\[
 \ker\bigl(G_{\mathbf Q}\to
          \operatorname{Out}(P_{U_0}^{(\ell)})\bigr)
                   =G_{\Omega_\ell}.
\]
\end{statement}
\begin{proof}
Sharifi's Theorem 1.1 \cite{A_Sharifi2002} identifies the
fixed field of this kernel with $\Omega_\ell$ from the
freeness of the rational graded Lie algebra generated by
the odd Soule classes. Brown's Theorem 1.1 and its stated
consequence for the Deligne--Ihara action
\cite{A_Brown2012} supply this freeness. Thus both inputs
are theorems in the present application. The identification
is also stated with precisely these hypotheses in the
introduction of \cite{A_IshiiKernels2026}. Taking the subgroup
that fixes the common fixed field yields the displayed kernel
by the Galois correspondence. No step here extends the odd
regular-prime hypothesis to $\ell=2$ or to irregular primes.
\end{proof}

\begin{statement}{A.45}{Open}{the general kernel questions left after these explicit computations}
\label{st:A.45}
\label{A:kernel-open}
For a nonconstant curve over a field of positive
transcendence degree, the quotient in
Proposition~\ref{A:family-kernel-quotient} is not computed
in general here; universal families make the remaining
congruence kernel an exact instance of the problem. For a
general single-prime action, including the tripod at
$\ell=2$ and at irregular primes, nontriviality and the
specialization subgroups in Proposition~\ref{A:proprime-kernels}
do not determine the entire kernel. The Sharifi--Brown
calculation above handles its stated odd regular-prime
instances. The January 2026 once-punctured CM kernel paper
\cite{A_IshiiKernels2026} assumes a further elliptic
Deligne--Ihara assertion in its main Theorem A; its
application is therefore not an unconditional computation
of the remaining kernels, and is not used as one here.
\end{statement}

\subsection{Further consequences and limits of the finite-derived commutant}

\begin{statement}{A.46}{Refuted}{the exact same-level strong metabelian assertion}
\label{st:A.46}
\label{A:strong-metabelian-counterexample}
For $U=\mathbf P^1_{\mathbf Q}\setminus\{0,1,\infty\}$,
the natural map
\[
 \operatorname{Aut}_{\mathbf Q}(U)\longrightarrow
 \operatorname{Aut}_{G_{\mathbf Q}}(E_U/F'')/\operatorname{Inn}(F/F'')
\]
is not surjective. Thus the instance $S_{2,0,U,U}$ of
Conjecture 0.1(2) in \cite{A_Yamaguchi2024}, whose
Remark 0.2 explicitly names $(2,0)$ as the best expected
pair, has a counterexample in the full profinite setting.
The counterexamples form the closed product of involutions
in Proposition~\ref{A:metabelian-commutant}.
\end{statement}
\begin{proof}
We verify the precise passage from an outer commutant to an
arithmetic extension automorphism. The center of $M=F/F''$
is trivial. Indeed, a central element has abelian image
$a\in A$ acting trivially on the free $R$-module $M'$, so
$a=1$. It then lies in $M'$, and multiplication by $X-1$
kills it. The latter is a non-zero-divisor in $R$, as
Chen--Deligne's decomposition into coefficient power-series
rings shows; this is their Corollary 4.1.6. Thus the element
is trivial.

Consequently the map
\[
 E_U/F''\longrightarrow
 \operatorname{Aut}(M)\times_{\operatorname{Out}(M)}G_{\mathbf Q},
 \qquad e\longmapsto(\operatorname{Int}(e)|_M,p(e))
\]
is a continuous isomorphism. Its kernel is $Z(M)$; for
surjectivity, choose any lift of the specified Galois element
and correct its automorphism by the required element of $M$.
For the actual involution $\phi_J$ constructed above, the map
\[
             (a,\sigma)\longmapsto
                   (\phi_Ja\phi_J^{-1},\sigma)
\]
preserves this fibre product because $[\phi_J]$ centralizes
the outer Galois action. It is therefore an actual continuous
$G_{\mathbf Q}$-automorphism of the arithmetic metabelian
extension, restricting to $\phi_J$ geometrically. Changing
$\phi_J$ by an inner automorphism changes this map by the
allowed geometric inner automorphism.

Every geometric automorphism of the tripod belongs to $S_3$,
whose action on $M^{\rm ab}$ is faithful. The constructed
nonidentity classes act identically there and are noninner,
so none is geometric. The exact finite certificate of
Proposition~\ref{A:finite-metabelian} detects this failure.
The cited conjecture uses full geometrically metabelian
arithmetic groups in characteristic zero and its $n=0$
target is the same quotient; hence it is exactly the map
just disproved, not a statement obtained after forgetting
additional derived layers.
\end{proof}

\begin{statement}{A.47}{Known}{center-freeness of all finite-derived tripod groups}
\label{st:A.47}
\label{A:solvable-centers}
For $F$ free profinite of rank two and every $m\ge2$, put
$S_m=F/F^{(m)}$, with closed derived subgroups. Then
$Z(S_m)=1$. Its maximal prosolvable quotient
$S_\infty=\varprojlim_m S_m$ is also center-free.
The finite-step assertion is included in Theorem B of
Yamaguchi's May 2026 preprint \cite{A_Yamaguchi2026Centers};
the following transfer argument proves the case used here
independently.
\end{statement}
\begin{proof}
The case $m=2$ was proved in the preceding corollary.
Inductively, a central element $z\in S_m$ maps trivially
to $S_{m-1}$, hence lies in
$N=F^{(m-1)}/F^{(m)}$. Let $H$ range over the open normal
subgroups of $F$ containing $F^{(m-1)}$. Since
$F^{(m)}\subset H'$, the image of $z$ in $H^{\rm ab}$
is defined and invariant under the finite group $F/H$.
Its image in $F^{\rm ab}$ is zero.

Write $i:H^{\rm ab}\to F^{\rm ab}$ for inclusion and
$\operatorname{Ver}:F^{\rm ab}\to H^{\rm ab}$ for the
continuous transfer defined using finitely many coset
representatives. Its defining product gives
\[
 \operatorname{Ver}(i(v))=\sum_{q\in F/H}q(v)
 \qquad(v\in H^{\rm ab}).
\]
For the invariant image $z_H$ this says
$[F:H]z_H=\operatorname{Ver}(i(z_H))=0$.
The abelianization of $H$ is a free profinite abelian group
by profinite Nielsen--Schreier, hence has no nonzero
integer torsion. Therefore $z_H=0$ integrally at every
prime.

Finally
$\bigcap_H H'=F^{(m)}$. To verify the nontrivial inclusion,
an element of $F^{(m-1)}$ outside $F^{(m)}$ survives in
some finite quotient of $F/F^{(m)}$. In that quotient the
image of $F^{(m-1)}$ is abelian. Taking $H$ to be its
inverse image gives an open normal subgroup for which the
selected element is not in $H'$. The converse inclusion
is immediate. Hence $z=1$. A central element of
$S_\infty$ has trivial image in every $S_m$, proving the
last assertion.
\end{proof}

\begin{statement}{A.48}{Known}{lower-derived commutants and the prosolvable limit}
\label{st:A.48}
\label{A:solvable-projection-limit}
Let $C_m$ be the commutant of the tripod's $G_K$-image in
the peripheral outer automorphism group of $S_m$, for a
number field $K$, and let $C_\infty$ be the corresponding
commutant on $S_\infty$. Then, for every $m\ge5$,
\[
 \operatorname{image}\bigl(C_m\to\operatorname{Out}(S_{m-3})\bigr)
                   =S_3,
       \qquad C_\infty=S_3.
\]
In particular no nonidentity metabelian involution of
Proposition~\ref{A:metabelian-commutant} lifts to an element
of $C_5$. These exact projections do not determine whether
$C_3$, $C_4$, or any $C_m$ itself equals $S_3$.
\end{statement}
\begin{proof}
Center-freeness identifies the arithmetic extension with
$\operatorname{Aut}(S_m)\times_{\operatorname{Out}(S_m)}G_K$,
exactly as in the preceding corollary. Thus every member of
$C_m$ supplies a $G_K$-automorphism of that extension.
The truncated Hom theorem, Mochizuki's Theorem A-prime or
18.1 \cite{A_Mochizuki1999}, in its full profinite form
recorded in Theorem 0.4 of \cite{A_Yamaguchi2024}, gives a
dominant morphism $U\to U$ agreeing with that automorphism
on $S_{m-3}$. This is precisely why three derived layers
are lost; the statement of that theorem contains this
forgetting map.

A dominant self-map of the tripod has degree one. Indeed
its extension to $\mathbf P^1$ has degree $d$, the three
target punctures have inverse image contained in the three
source punctures, and the total ramification above them
is at least $3d-3$. Riemann--Hurwitz bounds total
ramification by $2d-2$, so $d\le1$. The map therefore
belongs to $S_3$. Conversely all of $S_3$ acts over $K$
and commutes with Galois, proving equality of the displayed
image. At $m=5$ the target is $S_2$, where our nontrivial
involutions are outside $S_3$; the nonlifting assertion
follows through that exact homomorphism.

For $\alpha\in C_\infty$, fix $n\ge2$ and apply the
preceding argument at $m\ge\max(5,n+3)$. It makes the
image of $\alpha$ in $\operatorname{Out}(S_n)$ a member
of $S_3$. Its faithful action on the common abelianization
makes this member independent of $n$. Compose with its
inverse, so that $\alpha$ is inner on every $S_n$.
For an actual representative of $\alpha$, the sets of
conjugators in $S_\infty$ realizing it on $S_n$ are
nonempty closed nested subsets of the compact group
$S_\infty$. Their intersection is nonempty, and an element
there realizes $\alpha$ on every quotient and hence on
$S_\infty$. Thus its outer class is trivial, proving
$C_\infty=S_3$ without asserting the same at a finite
level.
\end{proof}

\begin{statement}{A.49}{Proved}{a full profinite unit family in the metabelian commutant}
\label{st:A.49}
\label{A:metabelian-unit-family}
The pure peripheral arithmetic commutant on $M=F/F''$
contains a closed subgroup isomorphic to
\[
                    \prod_p\mathbf Z_p^\times
                         =\widehat{\mathbf Z}^{\times}.
\]
Every member acts trivially on $M^{\rm ab}$ and on every
maximal pro-prime quotient of $M$, and centralizes $S_3$.
The family includes nontrivial coefficient-$2$ and
coefficient-$3$ members. This copy acts on the commutator
module; its image on the abelianization is the identity.
Theorem~\ref{A:third-derived-IA} below proves that no
nonidentity member of this entire family lifts to the
third-derived arithmetic commutant over any number field.
The preceding coefficient-five examples are retained
as explicit numerical certificates of this obstruction.
\end{statement}
\begin{proof}
For every $p\ne3$, including $p=2$, use the idempotent
$e_p$ of Proposition~\ref{A:metabelian-commutant}.
The proof there only needs $3$ to be invertible in
$\mathbf Z_p$, so it applies at two as well. For $p=3$
take $B=(\mathbf Z/4)^2$, the quotient of $A^{(3')}$
on the two indicated generators. In $\mathbf Z_3[B]$
take the sum of the character idempotents indexed by
\[
 \mathcal S=\{(a,b)\in(\mathbf Z/4)^2:
                     a\ne0,\ b\ne0,\ a+b\ne0\}.
\]
This set has six elements. It is invariant under
$(a,b)\mapsto(b,a)$, $(a,b)\mapsto(-a-b,b)$, and
$(a,b)\mapsto(-a,-b)$. The first two generate the
peripheral $S_3$ action, and the third is the only
nonidentity scalar unit action modulo four. The Fourier
sum is defined over $\mathbf Z_3$: conjugate characters
occur in pairs and its coefficients are in
$\frac18\mathbf Z$. Lift it to
$\mathbf Z_3[[A^{(3')}]]$ by normalized averaging over
the kernel of $A^{(3')}\to B$, and then to the
coefficient-$3$ factor of $R$. Call the resulting
idempotent $e_3$. It has augmentation zero. On all
six character factors the prime-to-three parts of
$X,Y,XY$ are nonidentity fourth roots of unity;
therefore $X-1,Y-1,XY-1$ are units on $e_3R$.

For $t=(t_p)_p\in\prod_p\mathbf Z_p^\times$ put,
coefficient by coefficient,
\[
 u(t)=1+\sum_p(t_p-1)e_p,\qquad
 r(t)=\frac{u(t)-1}{(X-1)(Y-1)},\qquad
 s(t)=\frac{X(Y-1)r(t)}{XY-1}.
\]
The denominators are inverted only on the displayed
idempotent factors. Each component of $u(t)$ is $t_p$
on $e_p$ and one on its complement, hence is a unit.
Set $f(t)=c^{r(t)}$, $h(t)=c^{s(t)}$, and define
\[
 \phi_t(x)=x,\qquad \phi_t(y)=f(t)^{-1}yf(t).
\]
The same explicit multiplication as before gives
$\phi_t(c)=c^{u(t)}$ and
$\phi_t(xy)=h(t)^{-1}xyh(t)$, proving preservation
of all three labelled peripheral classes. For two
parameters the conjugator exponent in the composite is
\[
 r(t)+u(t)r(v)=\frac{u(t)u(v)-1}{(X-1)(Y-1)}=r(tv).
\]
Thus $\phi_t\phi_v=\phi_{tv}$ and
$\phi_{t^{-1}}$ is its inverse. This proves that these
are actual continuous automorphisms, without relying
on a determinant-only existence assertion.

All $e_p$ are fixed by the scalar cyclotomic action and
by $S_3$. Formula (A4) then proves arithmetic and
$S_3$ commutation for their outer classes. The
parameters $t_p\in\mathbf Z_p$ are fixed coefficients
for that natural group-ring action. Their images in
every maximal pro-prime group are the identity, since
the corresponding nontrivial prime-to-$p$ character
idempotent has zero image there.

For injectivity, suppose $u(t)$ is a monomial from
$A$, as an inner automorphism requires. For each $p$
apply the exact coefficient and group quotient map
\[
 R\longrightarrow\mathbf Z_p[[A^{(p)}]].
\]
It sends $u(t)$ to one, so it sends that monomial to
one. The embedding of $A^{(p)}$ into this completed
group ring is faithful, and hence its $p$-component
is trivial. This holds for every $p$, so the monomial
is one. The nonzero torsion-free $\mathbf Z_p$-module
$e_pR$ now gives $t_p=1$ for every $p$. Thus the outer
map is injective. Every finite quotient uses finitely
many coefficient primes and finitely many residue
digits of their parameters, proving continuity.
Compactness then gives a closed image.

The independent certificate
\texttt{code/A/metabelian\_units\_certificate.py}
checks the coefficient-three Fourier idempotent
exactly over $\mathbf Q[B]$, its three required
inverses and its symmetry. It also checks explicit
noninner peripheral images on
$((\mathbf Z/4\mathbf Z)[\zeta_3])_{\rm add}\rtimes C_3$
and $((\mathbf Z/9\mathbf Z)[i])_{\rm add}\rtimes C_4$,
of orders $48$ and $324$, respectively. In the first
the parameter is $-1$; in the second it is $2$.
All $2304$ and $104976$ multiplication pairs, respectively,
the generation claims, all candidate inner conjugators,
and the three peripheral conjugacy equations pass.
\end{proof}

\subsection{A full profinite normalizer restriction from complex-conjugation primes}

\begin{statement}{A.50}{Proved}{descent of a power of the normalizer character}
\label{st:A.50}
\label{A:normalizer-character-descent}
Let $U/K$ be an affine hyperbolic curve over a number
field, let $A\ne1$ be a closed procyclic subgroup of
$F=P_U$, and let $N\subset E_U$ be a closed subgroup
normalizing $A$, with $p(N)=G_L$ open in $G_K$.
Fix a prime $\ell$ for which $A_\ell\ne1$ and write
$\theta_\ell:N\to\mathbf Z_\ell^\times$ for the action
on $A_\ell\simeq\mathbf Z_\ell$. Put
$d=\ell-1$ for odd $\ell$, and $d=1$ for $\ell=2$.
There is a continuous character
\[
 \lambda_\ell:G_L\longrightarrow\mathbf Z_\ell^\times,
 \qquad \lambda_\ell\circ p=\theta_\ell^d.
\]
Moreover either $A$ lies in one cusp inertia group, or,
over some finite extension $L_1/L$, there is a character
$\psi:G_{L_1}\to\mathbf Z_\ell^\times$ occurring as
a nonzero line in the Tate module of the Jacobian of
a finite cover of $U$, with
\[
              \lambda_\ell|_{G_{L_1}}=\psi^d.
\]
In the latter case $\psi$ is pure of Frobenius weight one
and has Hodge--Tate weights in $\{0,1\}$, with the
convention that the cyclotomic character has
Hodge--Tate weight one.
\end{statement}
\begin{proof}
The free profinite group $F$ is torsion-free, so its
nontrivial pro-$\ell$ subgroup $A_\ell$ is
$\mathbf Z_\ell$. Put $D=N\cap F$. Its action on
$A_\ell$ has trivial pro-$\ell$ part. To see this,
take a pro-$\ell$ Sylow subgroup $S$ of $D$ containing
the normal subgroup $A_\ell$. Closed pro-$\ell$
subgroups of a free profinite group are free pro-$\ell$:
their $\ell$-cohomological dimension is at most one,
which is the free pro-$\ell$ criterion. For $g\in S$,
the group generated by $A_\ell$ and $g$ has a normal
procyclic subgroup and procyclic quotient, hence is
metabelian. A free pro-$\ell$ group on at least two
generators has finite nonmetabelian $\ell$-group
quotients. Therefore this subgroup is procyclic and
$g$ centralizes $A_\ell$. The image of a Sylow
subgroup is a Sylow subgroup of the image, proving
the assertion about the pro-$\ell$ part.

For odd $\ell$ the remaining subgroup of
$\mathbf Z_\ell^\times$ has exponent dividing
$\ell-1$; for $\ell=2$ it is trivial. Thus
$\theta_\ell^d$ kills $D$ and descends through the
actual compact quotient $N/D=G_L$, giving
$\lambda_\ell$ and its continuity. No splitting of
$N\to G_L$ is asserted.

Choose an open characteristic subgroup $V\subset F$
which does not contain a topological generator of
$A_\ell$, and put $H=VA_\ell$. The finite cyclic
quotient $H/V$ shows that $A_\ell$ has nonzero image
in $H^{\rm ab}$. This abelianization is torsion-free,
so its restriction to $A_\ell\simeq\mathbf Z_\ell$
is injective. Descend the geometric cover defined by
$H$ to a finite extension and let $E_H\subset E_U$
be its open arithmetic fundamental group. Its
geometric kernel is exactly $H$. The open subgroup
$N\cap E_H$ has open arithmetic image, say
$G_{L_1}$ after further finite base extension.
Its conjugation on $H^{\rm ab}$ is the ordinary
arithmetic action: the geometric kernel lies in $H$
and acts inner. Consequently the displayed line is
a genuine $G_{L_1}$-stable line with character
$\psi=\theta_\ell|_{N\cap E_H}$, and
$\lambda_\ell|_{G_{L_1}}=\psi^d$.

In geometric homology the boundary kernel of
$H^{\rm ab}\to T J$ is a quotient of the permutation
module on the cusps tensored with
$\widehat{\mathbf Z}(1)$. If the line has zero image
in $T_\ell J$, split the finitely many cusp fields
by a finite extension. Its character then equals
$\chi_\ell$. The exact cyclotomic-normalizer
criterion, Proposition~\ref{A:cyclotomic-normalizer},
puts $A_\ell$ inside an inertia group $I$.
Every element of $A$ normalizes $A_\ell$, and that
same proposition gives $N_F(A_\ell)=I$. Hence
$A\subset I$.

In the other case the image of the line in
$T_\ell J\otimes\mathbf Q_\ell$ is injective.
The good-reduction Frobenius eigenvalues of this
Jacobian have weight one. Tate's Hodge--Tate theorem
for abelian varieties gives weights zero and one,
and the line inherits one of these at each local
place. These are exactly the asserted properties
of $\psi$; they do not depend on good reduction
at $\ell$ itself.
The power character $\lambda_\ell$ is consequently
Hodge--Tate after finite extension and hence before
it. It is also unramified outside finitely many
places: outside the bad-reduction and ramification
places for the finite extension, its inertia group
lies in $G_{L_1}$ and $\psi^d$ is unramified there.
\end{proof}

\begin{statement}{A.51}{Proved}{one-dimensional characters at a complex-conjugation prime}
\label{st:A.51}
\label{A:conjugation-prime-character}
Let $M/\mathbf Q$ be a finite Galois number field.
Fix a complex conjugation $c\in\operatorname{Gal}(M/\mathbf Q)$,
allowing $c=1$ when $M$ is totally real. Let $\ell$
be unramified in $M$ with Frobenius conjugacy class
equal to that of $c$. Every Hodge--Tate character
$\lambda:G_M\to\mathbf Q_\ell^\times$ is a finite-order
character times an integral power $\chi_\ell^a$.
When $M$ is totally real the conclusion holds for
every prime $\ell$, without the Frobenius condition.
\end{statement}
\begin{proof}
We spell out the maps in this use of the theory of
abelian $\ell$-adic representations. Tate's local
algebraicity theorem for a one-dimensional Hodge--Tate
representation gives a $\mathbf Q_\ell$-defined
algebraic character
\[
 h:(\operatorname{Res}_{M/\mathbf Q}\mathbf G_m)_{\mathbf Q_\ell}
                              \longrightarrow\mathbf G_{m,\mathbf Q_\ell}
\]
which describes $\lambda$ on sufficiently small
local units at the places over $\ell$. See
\cite[Chapter III, Theorem on p.~III--7 and Proposition
on p.~III--9]{A_SerreAbelian}. With the convention
calibrated by $\chi_\ell$, write its character as
$h=\sum_\tau n_\tau[\tau]$, $n_\tau\in\mathbf Z$.

Choose a modulus for $\lambda$ as in Chapter III,
Proposition 2 on p.~III--11 of that source. Global
reciprocity then says that $h$ is one on a subgroup
of finite index in the global unit group: on the
other local unit factors the character is trivial,
and the product on a principal idele is one. The
equality $\prod_\tau\tau(u)^{n_\tau}=1$ is algebraic,
so it remains true after passing to the complex
embedding fixed above. Dirichlet's unit theorem
therefore gives
\[
                  n_\tau+n_{c\tau}=2a
                  \quad\text{independent of }\tau.
\]
For a totally real field it gives directly
$n_\tau=a$ for every $\tau$. This is also the
character-lattice calculation in Chapter II,
Proposition on p.~II--30 of the cited source.

In the other case choose a place over $\ell$ whose
Frobenius is the particular element $c$; the
Chebotarev conjugacy condition permits this choice.
The algebraic character $h$ is defined over
$\mathbf Q_\ell$, so its coefficient vector is fixed
by the local decomposition group, which contains
$c$. Hence $n_{c\tau}=n_\tau$. The preceding equality
forces $n_\tau=a$ for all $\tau$, and $a$ is an
integer since the $n_\tau$ are integers. Thus $h$
is exactly the $a$th power of the norm character.

The quotient $\lambda\chi_\ell^{-a}$ consequently
has trivial algebraic character on every local
unit group over $\ell$. It has a modulus at every
other place as well. The global class-field map
therefore factors this quotient through a finite
ray class group. Its image is finite, proving the
claim. This argument neither presumes that the
Frobenius values are rational numbers nor replaces
$\mathbf Q_\ell$-valuedness by a stronger coefficient
condition.
\end{proof}

\begin{statement}{A.52}{Proved}{full-support subgroups and a positive-density exclusion}
\label{st:A.52}
\label{A:full-support-normalizer}
In the situation of Lemma~\ref{A:normalizer-character-descent},
let $M/\mathbf Q$ be a finite Galois extension
containing $L$, and let $\mathcal P_M$ be the set of
primes unramified in $M$ whose Frobenius conjugacy
class is the class of a complex conjugation.
Then either $A$ is contained in one cusp inertia
group, or
\[
                       A_\ell=1
                     \qquad(\ell\in\mathcal P_M).
\]
The set $\mathcal P_M$ has positive Dirichlet density.
In particular every such $A$ with nontrivial
pro-prime component at all but finitely many primes
is contained in a unique inertia group. This includes
every $A\simeq\widehat{\mathbf Z}$. If $L$ is totally
real, every nontrivial $A$ satisfying the stated
normalizer condition is contained in a unique
inertia group, with no restriction on its prime
support.
\end{statement}
\begin{proof}
Suppose $A$ is noncuspidal and $A_\ell\ne1$ for a
prime in $\mathcal P_M$. The preceding descent
lemma supplies $\lambda_\ell$ and, after finite
extension, a weight-one character $\psi$ with
$\lambda_\ell=\psi^d$ on that open Galois subgroup.
Hodge--Tate decompositions descend through finite
extensions, so $\lambda_\ell|_{G_M}$ is Hodge--Tate.
This descent is stated in
\cite[Chapter III, Appendix, Corollary 2 on
p.~III--32]{A_SerreAbelian}.
Lemma~\ref{A:conjugation-prime-character} gives
\[
              \lambda_\ell|_{G_M}
                         =\varepsilon\chi_\ell^a
\]
with $\varepsilon$ finite and $a\in\mathbf Z$.
Pass to a common finite extension also killing
$\varepsilon$. There the exact character equality is
$\psi^d=\chi_\ell^a$. Comparison at a prime of good
reduction different from $\ell$ gives $2a=d$,
because $\psi$ has Frobenius weight one. Comparison
of Hodge--Tate weights at a place above $\ell$
gives $a\in\{0,d\}$. These equalities contradict
$d>0$. Hence $A_\ell$ must be trivial.

Chebotarev gives the Dirichlet density among
rational primes
\[
 \lim_{s\downarrow1}
 \frac{\sum_{\ell\in\mathcal P_M}\ell^{-s}}
      {\log(1/(s-1))}
          =\frac{|\operatorname{Cl}(c)|}{[M:\mathbf Q]}>0.
\]
A subgroup with cofinite prime
support therefore cannot be in the noncuspidal
case. If $L$ is totally real, choose its totally
real Galois closure for $M$. The last assertion of
Lemma~\ref{A:conjugation-prime-character} supplies
the same contradiction at any prime with
$A_\ell\ne1$. Every nontrivial torsion-free procyclic
group has such a prime. Uniqueness of the containing
inertia follows from
Proposition~\ref{A:cyclotomic-normalizer}.
\end{proof}

This restriction applies in particular to the common-kernel
subgroups of Remark~\ref{A:cm-common-kernel}: their prime
components are tested inside actual finite-cover homology,
even when their images in every maximal pro-prime quotient
of the original curve vanish. A remaining noncuspidal
subgroup must have prime support disjoint from the stated
positive-density set, in addition to lying in
$R_{\rm nil}(F')$ in the CM example. The argument does not
exclude a single-prime subgroup at a prime outside that set.

\begin{statement}{A.53}{Proved}{the necessary CM fields in the remaining normalizer case}
\label{st:A.53}
\label{A:cm-normalizer-field}
Retain the notation of
Lemma~\ref{A:normalizer-character-descent}. If $A$ is
noncuspidal, its arithmetic normalizer field $L$
contains a CM subfield. Moreover, for every prime
$\ell$ with $A_\ell\ne1$ and every finite Galois
extension $M/\mathbf Q$ containing $L$, there is a
CM subfield $E\subset M$ admitting a field embedding
$E\hookrightarrow\mathbf Q_\ell$. In particular,
when $L$ contains no CM subfield, every nontrivial
procyclic subgroup satisfying the stated open
normalizer condition is cuspidal. When $L$ itself
is imaginary quadratic, a remaining noncuspidal
subgroup can have nontrivial prime components only
at rational primes that split in $L$.
\end{statement}
\begin{proof}
Choose a nontrivial component $A_\ell$ and the
character $\lambda_\ell$ constructed above. Let
$G=\operatorname{Gal}(M/\mathbf Q)$ and
$H=\operatorname{Gal}(M/L)$. The algebraic character
of $\lambda_\ell|_{G_M}$ has a coefficient function
$n:G\to\mathbf Z$ satisfying $n(gh)=n(g)$ for
$h\in H$: restriction from $G_L$ pulls back the
character by the norm from $M$ to $L$. Write
\[
 n(g)=a+b(g),\qquad
 a=|G|^{-1}\sum_{g\in G}n(g).
\]
The unit-character calculation of
\cite[Chapter II, section 3.1, pp.~II--29--30]{A_SerreAbelian}
gives
\[
 b(cg)=-b(g)
 \quad\text{for every complex conjugation }c\in G.
\]
This is the left action, the same action as in the
preceding lemma. The function $b$ is nonzero.
Otherwise all local algebraic exponents are one
integer $a$, so $\lambda_\ell|_{G_M}$ is a finite
character times $\chi_\ell^a$. Its equality to a
power $\psi^d$ after finite extension gives both
$2a=d$ and $a\in\{0,d\}$, the contradiction already
proved above.

Let $J\lhd G$ be generated by all products $c_1c_2$
of complex conjugations. Every such product fixes
$b$, so $b$ is left $J$-invariant. Since $J$ is
normal, it is also right $J$-invariant. Together
with right $H$-invariance this makes $b$ right
$JH$-invariant. All complex conjugations have the
same image modulo $J$, and that image is central.
No complex conjugation lies in $JH$: if one did,
left multiplication by it would act on $b$ as
right multiplication modulo $J$, hence would fix
$b$, contradicting $b\ne0$ and $cb=-b$.

It follows that $F=M^{JH}\subset L$ has a
nontrivial involution induced by complex conjugation.
All complex conjugations induce this same involution;
its fixed field is totally real, whereas none of
the embeddings of $F$ is real. Thus $F$ is a
totally imaginary quadratic extension of a totally
real field, proving the asserted CM subfield of $L$.
This proves explicitly the field-theoretic content
of the norm reduction in Chapter II, section 3.3
of the cited source.

For the additional local condition put
$P=\operatorname{Stab}_G(b)$ for the left action.
It contains $J$ and contains no complex conjugation.
The same centrality argument proves that
$E=M^P$ is a CM field. Since the algebraic
character is defined over $\mathbf Q_\ell$, the
local decomposition subgroup $D_v\subset G$ at
the chosen embedding into $\overline{\mathbf Q}_\ell$
fixes $b$. Hence $D_v\subset P$. The restriction
of this embedding to $E$ takes values in
\[
                        M_v^{D_v}=\mathbf Q_\ell,
\]
giving the required field embedding. This argument
also covers primes ramified in $M$.

Finally, for imaginary quadratic $L$ take $M=L$.
Its only CM subfield is $L$ itself. A quadratic
number field embeds in $\mathbf Q_\ell$ precisely
when it has a degree-one completion there, which
in the quadratic case means that $\ell$ splits.
This proves the last assertion, including the
exclusion of ramified primes in this case.
\end{proof}

\subsection{A third-derived obstruction detected by two explicit covers}

\begin{statement}{A.54}{Proved}{the coefficient-five involution does not lift to derived length three}
\label{st:A.54}
\label{A:third-derived-obstruction}
Let $\phi_5$ be the actual metabelian involution in
Proposition~\ref{A:metabelian-commutant}, with coefficient-five
determinant $1-2e_5$, where the cubic character defining $e_5$
takes both cusp generators to the same generator of $C_3$.
There is no outer automorphism of $S_3=F/F'''$ that commutes
with $G_{\mathbf Q}$ and projects to $[\phi_5]$ on $S_2$.
In fact there is no such lift commuting with
$G_{K_0}$, where
\[
 E:y^2=x^3+16,\qquad
 K_0=\mathbf Q(\mu_{15},E[5]).
\]
Thus this nonidentity metabelian class is already excluded
at the third derived level. More generally, for the unit
family $\phi_t$ of Proposition~\ref{A:metabelian-unit-family},
the existence of a commuting third-derived lift over either
of these fields forces $t_5\equiv1\pmod5$.
These assertions do not determine the entire commutant
at that level.
\end{statement}
\begin{proof}
We construct a finite metabelian quotient and two of its
intermediate covers. The argument uses the full homology
of these covers, which is retained at the third derived
level.

The first cover is the restriction over $U_0$ of
\[
 E\xrightarrow{[5]}E\longrightarrow\mathbf P^1,
       \qquad t=(y+4)/8.
\]
The last map is a cyclic cubic cover: putting $w_1=x/4$
gives $w_1^3=t(t-1)$. Its group is
$G_1=E[5]\rtimes C_3$, where the order-three action is
$(x,y)\mapsto(\zeta_3x,y)$. It is irreducible on $E[5]$,
since $T^2+T+1$ is irreducible over $\mathbf F_5$.
Its abelian cubic character at the first two cusps is
$(1,1)$.

For the second cover take the Fermat curve
\[
 C:X^5+Y^5+Z^5=0,
 \quad r=X^5/Z^5,\quad
 w_2=\frac{r-\zeta_3}{r-\zeta_3^2},\quad
 t=\frac{w_2^3}{w_2^3-1}.
\]
The cyclic permutation $(X,Y,Z)\mapsto(Y,Z,X)$ acts
on $r$ by $r\mapsto(-1-r)/r$. Its action on $w_2$ is
multiplication by a primitive cube root of unity.
Together with the fifth-root diagonal transformations
it gives
\[
 G_2=\mathbf F_5^2\rtimes C_3,\qquad
 T(a,b)=(b-a,-a).
\]
Here the cubic intermediate extension is
$w_2^3=t/(t-1)$, so the two cusp exponents are $(1,-1)$.
The diagonal action and the cyclic permutation are defined
over $K_0$.

Let $Y\to U_{0,K_0}$ be the fibre product of these two
finite etale covers. It is connected with constant deck
group
\[
 G=G_1\times G_2
   =(V_1\times V_2)\rtimes(C_3\times C_3),
 \qquad |G|=5625,
 \quad V_1=E[5],\quad V_2=\mathbf F_5^2.
                                                               \tag{A18}
\]
Indeed the two cubic Kummer classes are independent:
their valuation vectors at $0,1$ form a matrix of
determinant $-2$, nonzero modulo $3$. Each $G_i$ has
only $C_3$ as a nontrivial proper quotient, by the
irreducibility just checked. A common nontrivial quotient
would therefore give a common cubic subextension, which
the independence excludes. This also proves connectedness
after reduction at the prime used below. In particular
$G$ is metabelian and center-free.

Let $q:F\twoheadrightarrow G$ be the geometric quotient.
On $V_1\times V_2$ the unit $1-2e_5$ acts as
\[
                 S(v_1,v_2)=(-v_1,v_2).                       \tag{A19}
\]
To verify an actual induced automorphism, rather than only
an action on a module, extend $S$ by the identity on the
displayed complement $C_3^2$. Its value on $q(x)$ can be
corrected by a unique conjugation in $V_1\times V_2$:
$q(x)-1$ is invertible on both irreducible summands.
Call the resulting automorphism $\gamma$. It fixes
$q(x)$ and has the same value on $q([x,y])$ as $q\phi_5$.
The values of $\gamma(q(y))$ and $q(\phi_5(y))$ have the
same image in $C_3^2$. Their quotient is in $V_1\times V_2$;
equality of their commutators with $q(x)$, and the same
invertibility of $q(x)-1$, makes this quotient zero.
Consequently
\[
                         q\phi_5=\gamma q.
\]
The correcting inner conjugation is trivial on $V_1\times V_2$,
so $\gamma$ still acts there by (A19).

We now specify two subgroups of order five using a place
of $K_0$ over $211$. This prime splits completely in $K_0$.
Indeed $211\equiv1\pmod{15}$, and direct enumeration of
$E(\mathbf F_{211})$ gives $225$ points and all $25$ points
of $E[5]$ rational over $\mathbf F_{211}$. The reduction
is good and the five-torsion is unramified; thus its
Frobenius acts identically on the generators of $K_0$.
Choose the place with $\zeta_3=14$ and $\zeta_5=55$ in
the residue field. A point of $E[5](K_0)$ is uniquely
specified by its reduction at this place. Put
\[
 P_0=(2,92),\quad u_0=[2]P_0=(103,117),\quad v_0=(3,1),
 \qquad
 J=\langle(u_0,v_0)\rangle,
 \quad J_-=\langle(-u_0,v_0)\rangle.                          \tag{A20}
\]
Thus $\gamma(J)=J_-$. Both quotient covers $Y/J$ and
$Y/J_-$ have degree $1125$ over $U_0$.

Here is the complete finite rule for counting their points
over $\mathbf F_{211}$. For $t\in\mathbf F_{211}\setminus\{0,1\}$
the Frobenius class can meet $J$ only if both $t(t-1)$ and
$t/(t-1)$ are cubes. There are exactly $24$ such values:
\[
\begin{aligned}
 &12,13,28,64,65,68,87,88,89,90,97,98,\\
 &114,115,122,123,124,125,144,147,148,184,199,200.
\end{aligned}                                                   \tag{A21}
\]
For each take any cube roots $w_1,w_2$, and set
\[
 \begin{aligned}
 Q&=(4w_1,8t-4)\in E(\mathbf F_{211}),& P&=[9]Q,\\
 r&=(w_2\zeta_3^2-\zeta_3)/(w_2-1),&
 v&=\bigl(\log_{55}(r^{42}),
             \log_{55}((-1-r)^{42})\bigr)\in\mathbf F_5^2.
 \end{aligned}                                                   \tag{A22}
\]
The logarithms here mean the unique exponents in
$\{0,1,2,3,4\}$, so this is a finite formula involving no
choice of analytic logarithm.

The actual elliptic Kummer class in the $[5]$-torsor is
conjugate under $C_3$ to $[2]P$. To check the factor,
the Frobenius endomorphism $\pi$ has trace $-13$ and
norm $211$. This ordinary CM reduction has endomorphism
ring $\mathbf Z[\zeta_3]$, and solving these two equations
gives $\pi=1+15\zeta_3$ or $1+15\zeta_3^2$. For a point
$R$ with $[5]R=Q$ its Kummer class is
\[
 (\pi-1)R=3\zeta_3^iQ
             =2\zeta_3^i[9]Q,
\]
since $[15]Q=0$. The Fermat Kummer class is exactly $v$:
Frobenius multiplies fifth roots of $r,-1-r$ by
$r^{42},(-1-r)^{42}$. Thus a representative, up to
conjugacy, is $\delta_t=([2]P,v)$.

For any $\delta=(u,v)\in V_1\times V_2$, its number of
fixed cosets in $G/J$ is
\[
 \#(G/J)^\delta
   =125\sum_{i,j=0}^2
       \mathbf1_J(\zeta_3^iu,T^jv).                            \tag{A23}
\]
For a fixed complement coordinate there are $125$ cosets;
translation fixes them exactly when its conjugate lies
in $J$. This proves (A23) directly. The analogous formula
holds for $J_-$. Substitution of the $24$ values (A21)
in (A22)--(A23) gives
\[
\begin{array}{c|rr}
 t&\#(G/J)^{\delta_t}&\#(G/J_-)^{\delta_t}\\ \hline
 12&125&0\\
 200&125&0\\
 \text{the other 22 values in (A21)}&0&0.
\end{array}
\]
The other $185$ rational base points have nontrivial cubic
Frobenius and contribute zero to both counts. Hence
\[
       \#(Y/J)(\mathbf F_{211})=250,
       \qquad \#(Y/J_-)(\mathbf F_{211})=0.                     \tag{A24}
\]
The integer-only program
\texttt{code/A/third\_derived\_certificate.py} verifies
every enumeration above, all nine pairs of cubic-root
choices for every row, and the coset formula by direct
subgroup membership. Its complete output is saved beside
it. The independently written
\texttt{code/F/A\_m3\_independent.py} gives the same two
counts using $\delta(Q)=3Q$ directly.

Suppose now that a lift as stated existed over $K_0$.
By Lemma~\ref{A:solvable-centers}, an outer class commuting
with $G_{K_0}$ induces an actual automorphism of the
arithmetic extension with geometric kernel $S_3$.
An inner correction makes its geometric projection to
$S_2$ equal to the specified representative $\phi_5$.
The finite quotient just constructed is preserved.
There is no undetermined arithmetic twist in this
statement. Indeed the geometrically connected constant
$G$-torsor supplies an epimorphism
$q_{\rm ar}:\pi_1(U_{0,K_0})\to G$ whose geometric
restriction is $q$. The pair $(q_{\rm ar},p)$ identifies
the quotient by $\ker(q|_F)$ with $G\times G_{K_0}$.
An automorphism of this product over $G_{K_0}$ restricting
to $\gamma$ on $G$ sends $(1,\sigma)$ to
$(z_\sigma,\sigma)$ with $z_\sigma\in Z(G)$. Since
$Z(G)=1$, it is exactly $\gamma\times\operatorname{id}$.
It therefore takes the arithmetic subgroup for $Y/J$
to that for $Y/J_-$.

Write $H=q^{-1}(J)$ and $H_-=q^{-1}(J_-)$. Because $G$
is metabelian, $F''\subset H,H_-$, and therefore
\[
                       F'''\subset H',H_-'.
\]
Abelianizing their images in $S_3$ consequently gives
the full groups $H^{\rm ab}$ and $H_-^{\rm ab}$.
The assumed arithmetic automorphism gives a
$G_{K_0}$-equivariant isomorphism between them, hence
between $H^1_{\rm et}$ of the two geometric covers with,
for example, $\mathbf Q_7$ coefficients.

Both covers have smooth marked reduction at $211$.
Their ramification is tame, its order divides $5625$,
and the three branch sections remain distinct; the
normalizations give smooth compactifications with etale
boundary after unramified extension. Smooth proper base
change together with the boundary sequence therefore
identifies the preceding representations with the
cohomology of their reductions. Poincare duality gives
isomorphic compactly supported $H^1$, while the compactly
supported $H^0$ vanishes and $H^2=\mathbf Q_7(-1)$ for
both geometrically connected affine curves. The trace
formula forces equal point counts, contradicting (A24).
This proves the assertion over $K_0$, and a fortiori
over $\mathbf Q$.

For the final assertion, the other coefficient primes of
$\phi_t$ have zero image in the derived subgroup of (A18).
On that quotient its action is $t_5$ on $V_1$ and one on
$V_2$, and the same argument defining $\gamma$ gives the
actual induced automorphism. If $t_5\not\equiv1\pmod5$,
some power of its residue in the cyclic group
$\mathbf F_5^\times$ equals $-1$. The same power of a
putative commuting lift therefore induces (A19) on
$V_1\times V_2$. Every step from (A20) to the contradiction
uses only this finite action. It applies verbatim to that
power and rules it out. Hence $t_5\equiv1\pmod5$.
\end{proof}

\begin{statement}{A.55}{Proved}{a second coefficient-five digit is forced at derived length three}
\label{st:A.55}
\label{A:third-derived-second-digit}
For the unit family $\phi_t$ of
Proposition~\ref{A:metabelian-unit-family}, a lift to the
arithmetic commutant on $F/F'''$ over $\mathbf Q$, or over
$K_0=\mathbf Q(\mu_{15},E[5])$, necessarily satisfies
\[
                            t_5\equiv1\pmod{25}.
\]
This is the explicit second-digit certificate. The
structural Theorem~\ref{A:third-derived-IA} below
excludes every nonidentity member of the family at once;
the whole third-derived commutant remains undetermined.
\end{statement}
\begin{proof}
The preceding proposition already forces
$t_5\equiv1\pmod5$. Replace its two covers by
\[
 E\xrightarrow{[25]}E\xrightarrow{(y+4)/8}\mathbf P^1,
 \qquad X^{25}+Y^{25}+Z^{25}=0,
\]
using $r=X^{25}/Z^{25}$ and the same rational functions
$w_2=(r-\zeta_3)/(r-\zeta_3^2)$ and $t=w_2^3/(w_2^3-1)$.
Over $K_{25}=\mathbf Q(\mu_{75},E[25])\supset K_0$ their fibre product
has constant metabelian deck group
\[
 G_{25}=
 \bigl(E[25]\rtimes C_3\bigr)
       \times\bigl((\mathbf Z/25)^2\rtimes C_3\bigr),
             \qquad |G_{25}|=3515625.
\]
For completeness, connectedness follows from the same
independent cubic subfields. Any nontrivial quotient
common to the two factors would be solvable and hence
have a nontrivial abelian quotient; the abelianization
of either factor is $C_3$. Thus such a common quotient
would identify their cubic subfields. The order-three
actions are fixed-point-free on both two-dimensional
$\mathbf Z/25$ modules, so $Z(G_{25})=1$.

The unit-family automorphism acts on the two derived
modules by $(t_5,1)$. The proof constructing $\gamma$
in the preceding proposition works over $\mathbf Z/25$:
$X-1$ remains invertible, and it gives an actual
automorphism of this finite quotient. All coefficient
primes other than five have zero image here.

The prime $q=5701$ is prime and splits completely in
$K_{25}$. Direct enumeration gives
\[
 \#E(\mathbf F_q)=5625=75^2,\qquad
 E[25]\subset E(\mathbf F_q),\qquad
 [75]E(\mathbf F_q)=0,
\]
and $q\equiv1\pmod{75}$. Frobenius has trace $77$ and
norm $5701$, hence is $1-75\zeta_3$ or
$1-75\zeta_3^2$. Accordingly the elliptic
$[25]$-Kummer class at $Q=(4w_1,8t-4)$ is conjugate
to $-3Q$. Choose $\zeta_3=75$ and $\zeta_{25}=303$
in $\mathbf F_q$. Reduction identifies the two
geometric torsion modules with the displayed finite
modules. Put
\[
 P_0=(25,2322)\in E[25],\qquad v_0=(0,23),\qquad
 J_a=\langle(aP_0,v_0)\rangle\quad(a\in(\mathbf Z/25)^\times).
\]
Both $P_0$ and $v_0$ have exact order $25$.
These intermediate covers have degree $140625$.

The finite point-count rule is exactly (A22)--(A23)
with $v$ taken modulo $25$, the exponent $42$
replaced by $(5701-1)/25=228$, the elliptic class
replaced by $-3Q$, and the coset multiplicity $125$
replaced by $25^4/25=15625$. Its output is
\[
\begin{array}{c|rrrrr}
 a&1&6&11&16&21\\ \hline
 \#(Y_{25}/J_a)(\mathbf F_{5701})&15625&0&0&0&15625.
\end{array}
\]
Only the fibre $t=8$ contributes to the first entry;
only $t=4104$ contributes to the last. There are $624$
base values with trivial image in both cubic quotients
and $5075$ other base values. The certificate
\texttt{code/A/third\_derived\_25\_certificate.py}
checks all $5699$ base values, all $5616$ pairs of
cubic-root choices in the relevant fibres, and all
$100$ generator choices for the five displayed
order-$25$ subgroups. Its output is saved with the code.
The independent program
\texttt{code/F/A\_m3\_25\_independent.py}, using $-3Q$
directly and separate subgroup-membership calculations,
gives the same first two counts and the unique fibre $t=8$.

If $t_5\not\equiv1\pmod{25}$, its residue, already in
$1+5\mathbf Z/25$, is a nonidentity element of a
cyclic group of order five. A suitable positive power
has residue $6$. The same power of a putative commuting
lift takes the subgroup $J_1$ to $J_6$. Restricting to
$G_{K_{25}}$ and using the center-free constant torsor
identifies the two arithmetic intermediate-cover
subgroups exactly. As before $F'''\subset H'$,
so their full first cohomology representations would
be isomorphic. Good tame reduction and the trace
formula would give equal point counts at $5701$,
contradicting the first two entries of the table.
This proves the congruence.
\end{proof}

\subsection{Fibre isolation eliminates the entire IA image at derived length three}

\begin{statement}{A.56}{Proved}{the third-derived arithmetic commutant has no nontrivial metabelian IA image}
\label{st:A.56}
\label{A:third-derived-IA}
Let $K$ be any number field and
\[
 C_3(K)=C_{\operatorname{Out}(F/F''')}(G_K).
\]
An element of $C_3(K)$ inducing the identity on
$F^{\rm ab}$ induces an inner automorphism of $M=F/F''$.
Equivalently, the actual homomorphism
\[
 \operatorname{image}\bigl(C_3(K)\longrightarrow
                      \operatorname{Out}(M)\bigr)
       \longrightarrow\operatorname{GL}_2(\widehat{\mathbf Z})
                                                               \tag{A25}
\]
is injective. No peripheral hypothesis is needed here.
In particular the entire nonidentity unit family of
Proposition~\ref{A:metabelian-unit-family}, at every
coefficient prime and every depth, admits no commuting
third-derived lift over any number field.

The theorem does not assert that
$\operatorname{Out}(F/F''')\to\operatorname{Out}(M)$
has trivial kernel on the arithmetic commutant.
\end{statement}
\begin{proof}
Put $A=M^{\rm ab}$, $R=\widehat{\mathbf Z}[[A]]$,
$I=\ker(R\xrightarrow\epsilon\widehat{\mathbf Z})$,
and $c=[x,y]$. We use $M'=Rc$ and the exact IA identity
(A4). The proof has six steps, including the arithmetic
normalizations and the point-count comparison.

\emph{1. An exactly invariant representative.}
Let the given third-derived class project to an IA class
of $M$, represented by a unit $u\in R^\times$ with
$\epsilon(u)=1$. There is a canonical continuous map
\[
 \partial:1+I\longrightarrow I/\overline{I^2}
       \simeq\widehat{\mathbf Z}^2=A,
       \qquad \partial(u)=u-1\pmod{\overline{I^2}}.
\]
The displayed identification sends $X-1,Y-1$ to the
two standard basis vectors. It follows by expanding
$X^aY^b$ modulo $I^2$: its image is
$1+a(X-1)+b(Y-1)$. This expansion extends continuously
from the integral group algebra. Conversely the ring
$\widehat{\mathbf Z}\oplus A$ with square-zero second
summand realizes these images, proving injectivity as
well as surjectivity. In particular $\partial$ is
additive on augmentation-one units, and the restriction
$A\to I/\overline{I^2}$ is the identity.

Arithmetic conjugation acts on IA outer classes by the
natural scalar action $\sigma_R$ induced by $\chi(\sigma)$
on $A$; its IA component commutes because IA outer
automorphisms form the abelian group in (A4).
Centralization therefore says
$u/\sigma_R(u)\in A$. Put $v=\partial u$ and multiply
$u$ by the monomial $a_{-v}\in A$. This is an allowed
inner change of representative. The resulting unit
$u_0$ has derivative zero and still satisfies
$u_0/\sigma_R(u_0)\in A$. That monomial has derivative
zero, so it is one. Hence
\[
       \sigma_R(u_0)=u_0\quad(\sigma\in G_K).                 \tag{A26}
\]
Choose an actual IA automorphism $\phi$ with this
determinant, and choose an inner adjustment of the
third-derived representative projecting exactly to
$\phi$. Such an adjustment lifts through the surjection
$F/F'''\to M$.

\emph{2. Exact identity on the arithmetic abelian quotient.}
Center-freeness supplies the arithmetic automorphism
$\beta$ of the extension with geometric kernel $F/F'''$.
Write $E_j=\pi_1(U_{0,K})/F^{(j)}$. We claim that its
induced automorphism of $E_1$ is the identity, not just
the identity on its geometric kernel.
For an arithmetic lift of $\sigma$, let $b_\sigma$ be
its conjugation action on $M$. On $Rc$, this action has
the form $r\mapsto\sigma_R(r)v_\sigma$ for a unit
$v_\sigma\in R^\times$. The automorphism $\phi$ acts
there by multiplication by $u_0$. Its commutator with
$b_\sigma$ thus acts by
$u_0/\sigma_R(u_0)=1$.
Since their outer classes commute, this commutator is
$\operatorname{Int}(d_\sigma)$ for some $d_\sigma\in M$.
Conjugation by $d_\sigma$ acts on $Rc$ by its monomial
$d_\sigma^{\rm ab}\in A\subset R^\times$. Faithfulness
of this monomial embedding gives $d_\sigma\in M'$.
Consequently $\beta$ changes each arithmetic lift only
by an element of $M'$, and it fixes the geometric
abelianization. This proves
\[
                         \beta|_{E_1}=\operatorname{id}.
                                                               \tag{A27}
\]

\emph{3. A finite quotient with different cyclic subgroups.}
Suppose the projected IA outer class is nontrivial.
Then $u_0\notin A$. The closed cyclic subgroups
$C=\overline{\langle c\rangle}$ and
$D=\overline{\langle\phi(c)\rangle}$ are not conjugate
in $M$. Indeed $C=\widehat{\mathbf Z}c$ and
$D=\widehat{\mathbf Z}u_0c$ inside $Rc$. A conjugator
acts by a monomial $a\in A$. Equality of these two
cyclic subgroups would give $u_0=\lambda a$ with
$\lambda\in\widehat{\mathbf Z}^{\times}$.
Applying augmentation gives $\lambda=1$, contrary to
$u_0\notin A$.

There is therefore a finite quotient of $M$ in which
their images are not conjugate. To justify this
finite detection, for every open normal subgroup $N$
consider the closed subset of $M$ of conjugators
identifying their images modulo $N$. If all these
sets were nonempty, their finite-intersection property
and compactness would give an actual conjugator in
$M$, a contradiction. Since $M$ is topologically
finitely generated, it has finitely many open subgroups
of any given index. Intersecting the finite $\phi$-orbit
of such an $N$ makes it $\phi$-stable without losing
the distinction. We obtain a finite metabelian quotient
\[
 q:M\twoheadrightarrow G,\qquad q\phi=\gamma q,
 \qquad J=\langle q(c)\rangle,
 \quad J'=\gamma(J)
\]
with $J,J'$ nonconjugate. For $g=q(c)$, the two
permutation characters satisfy
\[
                 \#(G/J)^g>0,\qquad \#(G/J')^g=0.           \tag{A28}
\]
The second equality follows because containment of the
generator $g$ in a conjugate of $J'$ would make the
cyclic groups $J,J'$ of equal order conjugate.

\emph{4. Arithmetic descent and removal of the finite twist.}
After a finite extension of $K$, the geometric quotient
is a geometrically connected constant $G$-torsor
$Y\to U_0$, with all deck transformations defined.
The map $(q_{\rm ar},p)$ identifies the corresponding
arithmetic quotient with $G\times G_K$.
The induced automorphism is of the form
\[
                 (h,\sigma)\longmapsto
                         (\gamma(h)z_\sigma,\sigma),
                 \qquad z_\sigma\in Z(G).
\]
Here $z:G_K\to Z(G)$ is a continuous homomorphism:
the $G$-factor is constant, and multiplication in the
product proves this formula directly. Replace $K$ by
the finite extension killing $z$, and then, if needed,
by its Galois closure over $\mathbf Q$. From now on
$K/\mathbf Q$ is Galois, (A27) still holds, and the
finite arithmetic action is exactly
$\gamma\times\operatorname{id}$. No center-free
assumption on $G$ is required in this step.

\emph{5. A Frobenius class and a cover isolating its fibre.}
Hilbert irreducibility gives $a\in U_0(K)$ such that
$Y_a=\operatorname{Spec}L$ is a field with
$\operatorname{Gal}(L/K)=G$; the precise specialization
statement is \cite[Proposition 3.3.1 and Theorem 3.4.1]{A_SerreTopics}.
Take the normal closure $L^*/\mathbf Q$ and lift
$g\in\operatorname{Gal}(L/K)$ to
$h\in\operatorname{Gal}(L^*/K)$.
Chebotarev applied to $L^*/\mathbf Q$
\cite[Chapter V, Theorem 3.23]{A_MilneCFT}
gives arbitrarily large rational primes $q$ and a prime
of $L^*$ with Frobenius $h$. Such a $q$ splits in $K$,
and at the resulting degree-one place $v$ the Frobenius
of $Y_a$ is $g$. Avoid the finitely many primes of bad
marked reduction, primes dividing $|G|$, and primes
at which $a$ or $1-a$ fails to be a unit. Write
$a_0\in\mathbf F_q\setminus\{0,1\}$ for its reduction.

Put $N=q-1$ and extend constants to $K'=K(\mu_N)$.
The chosen place has a degree-one extension to $K'$,
because $q\equiv1\pmod N$. Consider the actual finite
etale abelian cover
\[
 B:\qquad z_1^N=t/a,\qquad
               z_2^N=(1-t)/(1-a)                            \tag{A29}
\]
of $U_{0,K'}$. Its arithmetic monodromy factors through
$E_1$, so (A27) preserves this cover exactly, including
the two constant Kummer twists. Over $\mathbf F_q$,
the fibre $B_t$ has $N^2$ rational points at $t=a_0$
and no rational points at any other $t\in U_0(\mathbf F_q)$:
the $N$-th power of every nonzero element is one.

Form the two finite etale covers
\[
 Z=(Y/J)\times_{U_0}B,
       \qquad Z'=(Y/J')\times_{U_0}B.
\]
Their geometric monodromy is a subgroup of a product
of a metabelian group and an abelian group, hence is
metabelian. The two exact arithmetic actions already
proved make $\beta$ identify their finite-cover data.
Counting their rational points at the chosen place gives
\[
 \#Z(\mathbf F_q)=N^2\#(G/J)^g>0,
       \qquad \#Z'(\mathbf F_q)=N^2\#(G/J')^g=0.              \tag{A30}
\]

\emph{6. Why the third derived quotient forces equality.}
Neither fibre product is assumed geometrically connected.
For every geometric connected component its fundamental
group is an open subgroup $H\subset F$ containing $F''$,
because the entire geometric monodromy is metabelian.
Consequently $F'''\subset H'$ and the image of $H$ in
$F/F'''$ retains the full abelianization $H^{\rm ab}$.
The arithmetic automorphism $\beta$ gives an equivariant
bijection of the component sets and equivariant
isomorphisms of these abelianizations, with the conjugacy
identifications on stabilizers. It therefore identifies
both $H^0_{\rm et}$ and $H^1_{\rm et}$ of the geometric
curves $Z,Z'$ as $G_{K'}$-representations, including the
permutation of components.

All covers in this comparison have good marked reduction
at the chosen place: $Y$ was chosen good, and (A29) is
tame there, with constant units and branch sections only
at $0,1,\infty$. Finite products and the smooth marked
normalizations preserve this property: etale locally at
a branch section the covers have Kummer equations
$u^e=t$ with $e$ prime to the residue characteristic,
and their normalized products have the same form with
the least common multiple of the indices. Choose a
coefficient prime different from $q$. Base change and
Poincare duality identify the compactly supported first
cohomology representations as well. Compactly supported
$H^0$ vanishes, and $H^2_c$ is the component permutation
module twisted by $-1$ for both curves. Thus the trace
formula makes the two point counts equal, contradicting
(A30).

The projected IA class is therefore trivial. Applying
this to the quotient of any two elements having the
same abelian action proves the injectivity of (A25).
The unit family acts identically on $A$ and is injective
in IA outer automorphisms, so every nonidentity member
is excluded, at all primes and to all depths. The proof
does not compare different automorphisms inside the
kernel of the third-to-second-derived projection.
\end{proof}

\begin{statement}{A.57}{Proved}{the peripheral metabelian commutant has no extra scalar multiplier}
\label{st:A.57}
\label{A:metabelian-multiplier}
For every number field $K$, the image of
$C_{\operatorname{Out}_{\rm lac}(M)}(G_K)$ on
$M^{\rm ab}$ is exactly the geometric $S_3$.
In particular a pure peripheral outer automorphism
commuting with $G_K$ has common multiplier one.
This statement leaves the nontrivial IA commutant of
Proposition~\ref{A:metabelian-unit-family} present at
the metabelian level.
\end{statement}
\begin{proof}
Compose with a geometric permutation of the three cusps
to make the class pure. By
Proposition~\ref{A:abelian-commutant}, its abelian action
is $\rho\operatorname{id}$ for a unit
$\rho\in\widehat{\mathbf Z}^{\times}$. Center-freeness
gives an arithmetic automorphism $\beta$ of $E_2$.
We first normalize its action on $E_1$ exactly.

The geometric normalizer of $\langle x\rangle$ in $M$
is $\langle x\rangle$. Here is a proof with all
profinite coefficients retained. If an element commutes
with $x$, the free class-two quotient shows that its
$y$-coordinate in $M^{\rm ab}$ is zero. It is therefore
$x^a c^r$ for $a\in\widehat{\mathbf Z}$, $r\in R$,
and commutation gives $(X-1)r=0$. Multiplication by
$X-1$ is injective on $R$. Indeed in a coefficient
factor $\mathbf Z_p[[A]]$, an $X$-invariant element has
constant coefficients on the finite $X$-orbits in
every finite abelian group quotient. Refine the order
of the first cyclic factor by $p^j$ while retaining
the second factor. Projection sums $p^j$ equal
coefficients. Modulo $p^j$ the projected element is
zero. This holds for every coefficient modulus and
group quotient, so the element is zero. Thus
$C_M(x)=\langle x\rangle$. An element normalizing
$\langle x\rangle$ induces exponent one on it by
abelianization, and hence centralizes it. The same
argument applies to $y$.

Let $s_0,s_1$ be the tangential sections at $0$ pointing
towards $1$, and at $1$ pointing towards $0$,
respectively. Their images in $E_1$ are the same
section. This can be checked on the actual abelian
Kummer covers $t^{1/n},(1-t)^{1/n}$: the parameter
has leading coefficient one at its own cusp and the
other function has value one. All these leading
coefficients admit the compatible rational roots one.
Choose a path whose abelian image identifies these
root labels. Both sections then have abelian Kummer
coordinates $(0,0)$. This identifies
\[
              E_1=A\rtimes G_K.
\]
The arithmetic normalizers are exactly
$\langle x\rangle\rtimes s_0(G_K)$ and
$\langle y\rangle\rtimes s_1(G_K)$: subtract the
specified tangential lift of an arithmetic element
and use the geometric normalizer equality above.

Write the induced automorphism of $E_1$ as
\[
               (v,g)\longmapsto(\rho v+c(g),g).
\]
Peripheral preservation takes each of those two
arithmetic normalizers to a geometric conjugate of
itself. If the conjugator at cusp $i$ has abelian
coordinate $a_i$, comparison of the second coordinate
at cusp $0$ and of the first coordinate at cusp $1$
gives
\[
 c_2(g)=(1-\chi(g))(a_0)_2,
       \qquad c_1(g)=(1-\chi(g))(a_1)_1.
\]
Thus $c$ is the actual coboundary
$(1-\chi)a$, with $a=((a_1)_1,(a_0)_2)$.
An inner adjustment with abelian coordinate $-a$
now makes
\[
                   \beta|_{E_1}(v,g)=(\rho v,g).             \tag{A31}
\]

Take any odd rational prime $q$ splitting completely
in a Galois closure of $K$, put $N=q-1$, and extend
constants to $K'=K(\mu_N)$. There is a degree-one
place of $K'$ over $q$, since $q\equiv1\pmod N$.
For $a\in\{2,\ldots,q-1\}$ consider
\[
 B_a:\quad z_1^N=t/a,\qquad z_2^N=(1-t)/(1-a).
\]
Relative to the split coordinates (A31), its arithmetic
quotient is $(v,g)\mapsto v-(\kappa_N(a)(g),
\kappa_N(1-a)(g))$. If $r$ is the residue of $\rho$
modulo $N$, equation (A31) therefore sends its
arithmetic cover subgroup to that of
\[
 B'_a:\quad w_1^N=t/a^r,\qquad
                  w_2^N=(1-t)/(1-a)^r.                      \tag{A32}
\]
This verifies the constant twists as well as the
geometric monodromy; it is not merely a comparison
of untwisted abelian covers.

The monodromy of either cover is geometrically abelian.
Every geometric component subgroup $H$ thus contains
$F'$, and $F''\subset H'$. Already $E_2$ retains
its full first homology. Hence $\beta$ identifies the
$H^0$ and $H^1$ Galois representations of $B_a,B'_a$.
Both have good marked reduction at $q$: their
projective equations are diagonal Fermat equations
of degree $N$, with unit coefficients and with
$N$ invertible in the residue field. Their boundaries
are etale there. The same duality and trace-formula
argument as in the preceding theorem forces equal
affine point counts.

The first cover has exactly $N^2$ rational points,
all over $t=a$. The second has $N^2$ rational points
if
\[
                     a^r+(1-a)^r=1\quad\text{in }\mathbf F_q,
                                                               \tag{A33}
\]
and zero otherwise. Thus (A33) holds for every
$a\in\mathbf F_q\setminus\{0,1\}$.
The map $f(0)=0$, $f(x)=x^r$ on $\mathbf F_q^\times$
is multiplicative and bijective, since $r$ is a unit
modulo $q-1$. Equation (A33), applied to
$a=x/(x+y)$, proves additivity whenever
$x,y,x+y$ are nonzero. The zero cases are immediate;
when $x+y=0$ use $f(-1)=-1$, since $r$ is odd.
Therefore $f$ is an automorphism of the prime field
$\mathbf F_q$ and is the identity. It follows that
\[
                         \rho\equiv1\pmod{q-1}.
\]
For every positive integer $n$, Chebotarev supplies
such an odd prime $q$ also splitting in $\mathbf Q(\mu_n)$.
Then $n\mid q-1$, so $\rho\equiv1\pmod n$.
This proves $\rho=1$ in $\widehat{\mathbf Z}$.
Undoing the initial geometric permutation gives exactly
the displayed $S_3$ image.
\end{proof}

\begin{statement}{A.58}{Proved}{the third-to-second-derived peripheral commutant is exactly geometric}
\label{st:A.58}
\label{A:third-derived-projection}
For every number field $K$,
\[
 \operatorname{image}\left(
 C_{\operatorname{Out}_{\rm lac}(F/F''')}(G_K)
          \longrightarrow\operatorname{Out}(F/F'')\right)=S_3.
\]
The same equality holds for the projection from every
derived length $m\ge3$ to length two. In particular
all nongeometric metabelian peripheral commutant classes
are obstructed at the next derived step.
The same-level commutant at length three is reduced
to its kernel over this geometric image; that kernel
has not been proved trivial here.
\end{statement}
\begin{proof}
By Proposition~\ref{A:metabelian-multiplier}, the
abelian image of a projected class is that of one
geometric permutation. Divide by that permutation.
The quotient has trivial abelian action, and
Theorem~\ref{A:third-derived-IA} makes its metabelian
outer image trivial. Conversely all six geometric
permutations lift over $K$ at every derived level.
For $m>3$ first project to length three. This gives
the stated equality and does not assert anything
further about the kernel of that projection.
\end{proof}

\subsection{An explicit unrestricted last-layer construction and its peripheral obstruction}

\begin{statement}{A.59}{Proved}{unrestricted third-derived scalar classes fail the peripheral test}
\label{st:A.59}
\label{A:last-layer-unrestricted}
For every odd prime $\ell$, the kernel
\[
 \ker\bigl(\operatorname{Out}(F/F''')\to\operatorname{Out}(M)\bigr)
       \cap C_{\operatorname{Out}(F/F''')}(G_{\mathbf Q})
\]
contains a closed copy of $\mathbf Z_\ell^\times$.
No nonidentity member of this copy preserves the
conjugacy class of the cusp subgroup $\langle x\rangle$.
Thus this explicit construction does not give an element
of the peripheral kernel left open above.
\end{statement}
\begin{proof}
Write $L=F''/F'''$, $D=M'$, and
$\Lambda=\widehat{\mathbf Z}[[M]]$.
We use the profinite Fox relation-module sequence,
with its actual boundary map,
\[
 0\longrightarrow L\longrightarrow\Lambda^2
   \xrightarrow{\partial}\Lambda
   \xrightarrow\epsilon\widehat{\mathbf Z}\longrightarrow0,
 \qquad \partial(a,b)=a(x-1)+b(y-1).                         \tag{A34}
\]
This is the Blanchfield--Lyndon sequence, recalled
with precisely this formula in
\cite[Proposition 1.1]{A_Yamaguchi2026Centers} from
Ihara's profinite Fox calculus. More explicitly the
continuous Fox derivation $\delta:F\to\Lambda^2$
satisfies $\delta(x)=(1,0)$, $\delta(y)=(0,1)$,
$\delta(fg)=\delta(f)+\bar f\delta(g)$, and
$\partial\delta(f)=\bar f-1$.
The resulting Magnus map identifies $F/F'''$ with
\[
 \{(v,g)\in\Lambda^2\rtimes M:\partial v=g-1\}.             \tag{A35}
\]
Indeed its kernel is $F'''$ by (A34). For a prescribed
$g$, subtract the Fox vector of any lift of $g$;
(A34) says that every other vector with the same
boundary differs by an element of $L$, proving the
asserted equality of the images.

Let $D^{(\ell')}$ be the prime-to-$\ell$ factor of
the abelian group $D=Rc$. Reduction modulo two of
the augmentation $Rc\to\widehat{\mathbf Z}$ gives
an $M$-invariant surjection
$\lambda:D^{(\ell')}\to C_2$. Put $H=\ker\lambda$.
Normalized averaging over this pro-prime-to-$\ell$
group defines $e_H\in\mathbf Z_\ell[[M]]$.
For $z\in D^{(\ell')}$ with $\lambda(z)\ne1$, set
\[
                   e=e_H(1-z)/2.
\]
Regard $e$ as zero in the other coefficient factors
of $\Lambda$. It is a central idempotent of augmentation
zero: $z^2\in H$, and conjugation by $M$ fixes the
character $\lambda$. It is also arithmetic-invariant.
Indeed arithmetic acts on the augmentation of $D$
by $\chi^2$, which is one modulo two. Every finite
averaging denominator is an $\ell$-adic unit; no
averaging over a pro-$\ell$ kernel is used.

Choose explicitly $z=c^{\eta}$, where
$\eta\in\widehat{\mathbf Z}$ is zero at $\ell$
and one at every other prime, and take the same
word $\widetilde z=[x,y]^\eta$ in $F$. Since $ez=-e$,
\[
                    w=-\tfrac12e\delta(\widetilde z)
          \quad\text{satisfies}\quad \partial w=e.
\]
For $t\in\mathbf Z_\ell^\times$, put
$u_t=1+(t-1)e$ and define on (A35)
\[
 \Phi_t(v,g)=
 \bigl(u_tv+(1-t)(g-1)w,\ g\bigr).                         \tag{A36}
\]
Centrality of $e$ and the crossed-homomorphism identity
for $(g-1)w$ prove directly that this is a homomorphism
preserving the boundary condition. Also
$\Phi_t\Phi_s=\Phi_{ts}$ and
$\Phi_t^{-1}=\Phi_{t^{-1}}$. Thus these are actual
continuous automorphisms, identity on $M$, acting by
$t$ on $eL$ and by one on $(1-e)L$.

We check arithmetic commutation in the outer group.
First
\[
                         H^1(M,eL)=0.                       \tag{A37}
\]
To see this directly, a continuous cocycle restricts
on $H$ to a homomorphism from a pro-prime-to-$\ell$
group to the pro-$\ell$ group $eL$, hence is zero.
It descends to $M/H$, where $z$ is a central
involution acting by $-1$. If $b$ is such a cocycle,
the two expressions for $b(zg)=b(gz)$ give
$2b(g)=(1-g)b(z)$. This is the principal cocycle
with parameter $b(z)/2$, proving (A37).

Arithmetic preserves $eL$ and commutes with its scalar
$t$. The commutator of $\Phi_t$ with any arithmetic
automorphism therefore acts identically on $M$ and
on $L$, and is identity modulo $eL$. Such an
automorphism is a continuous cocycle $M\to eL$:
compare the image of any lift of $g\in M$ with that
lift; independence follows because $L$ is fixed.
By (A37) it is inner. Hence $[\Phi_t]$ commutes
with the arithmetic outer action.

Finally we give an exact detector of both nontriviality
and failure at the cusp. Put
\[
                    n_x=e(1,0)-(x-1)w\in eL.
\]
Map $M$ to the dihedral group of order eight using
the integral matrices
\[
 x\longmapsto\begin{pmatrix}1&0\\0&-1\end{pmatrix},
 \qquad
 y\longmapsto\begin{pmatrix}0&-1\\1&0\end{pmatrix}.
\]
The commutator maps to $-1$, so $H$ maps to the
identity and $e$ maps to the identity matrix.
The first row is fixed by $x$. Apply this row to
the first coordinate of $\Lambda^2$ in the
coefficient-$\ell$ factor. The resulting continuous
linear map kills $(x-1)eL$ and takes $n_x$ to
$(1,0)\in\mathbf Z_\ell^2$. In particular
\[
          (t-1)n_x\notin(x-1)eL\quad(t\ne1).               \tag{A38}
\]
The automorphism $\Phi_t$ changes $x$ by precisely
$(t-1)n_x$. If its image were conjugate to $x$, the
conjugator's image in $M$ would centralize $x$ and
thus be a power of $x$, by the proof of
Proposition~\ref{A:metabelian-multiplier}. Multiplying
the conjugator by the inverse power of $x$ reduces
it to $L$. The conjugation difference would then
lie in $(x-1)L$, contradicting (A38) after applying
$e$. A conjugate generator of $\langle x\rangle$
cannot use any other exponent, since $\Phi_t$ is
identity on abelianization. Thus it fails even this
single peripheral condition.

The same row detector shows that $t$ acts nontrivially
on $eL$ for $t\ne1$. An inner automorphism inducing
the identity on the center-free group $M$ has its
conjugator in $L$, and acts identically on $L$.
Therefore $t\mapsto[\Phi_t]$ is injective.
Continuity follows from (A36); compactness gives
a closed image.
\end{proof}

\begin{statement}{A.60}{Proved}{no normal prime-to-prime averaging block removes a peripheral divisor}
\label{st:A.60}
\label{A:haar-block-obstruction}
Let $\ell$ be any prime, let $N$ be a closed normal
pro-prime-to-$\ell$ subgroup of $M$, and put
$\Lambda_\ell=\mathbf Z_\ell[[M]]$.
Then $N\subset M'$. The natural map
\[
 \mathbf Z_\ell[[N]]\longrightarrow
 \Lambda_\ell/(x-1)\Lambda_\ell
       =\mathbf Z_\ell[[\langle x\rangle\backslash M]]       \tag{A39}
\]
is injective. In particular, for any nonzero
idempotent $e\in\mathbf Z_\ell[[N]]$ whose image is
central in $\Lambda_\ell$, multiplication
by $x-1$ is not surjective on $e\Lambda_\ell$.
It cannot be inverted there. The same holds for
$y$ and $xy$.

Consequently a scalar construction using Haar
idempotents of normal pro-prime-to-$\ell$ subgroups
cannot reproduce the simultaneous peripheral
inversions used in the metabelian construction.
This is an obstruction to that specified construction,
not a determination of the remaining peripheral kernel.
\end{statement}
\begin{proof}
Let $D_\ell$ be the pro-$\ell$ factor of $M'$.
Both $N$ and $D_\ell$ are normal, so
$[N,D_\ell]\subset N\cap D_\ell=1$.
The action of $M$ on
$D_\ell=\mathbf Z_\ell[[A]]c$ is its faithful
monomial action through $A$. Hence every element
of $N$ has trivial image in $A$, proving
$N\subset M'$.

Since the abelianization embeds
$\langle x\rangle$ into $A$, one has
$N\cap\langle x\rangle=1$. The map of profinite
sets $N\to\langle x\rangle\backslash M$ is therefore
a continuous injection, hence a closed embedding.
The induced map of free profinite
$\mathbf Z_\ell$-modules is injective: a finite
clopen partition of a closed subset of a profinite
space extends to a finite clopen partition of the
whole space. Thus every nonzero finite coefficient
of a measure on $N$ remains detectable in the
ambient permutation module. This proves (A39).
The equality on the right of (A39) is precisely
the coinvariant quotient for left multiplication
by the procyclic group $\langle x\rangle$;
the image of $x-1$ is closed by compactness.

In particular the image of $e$ in that quotient
is nonzero. If $x-1$ were surjective on
$e\Lambda_\ell$, one could write
$e=(x-1)v$, giving zero image, a contradiction.
Replace $x$ by $y$ or $xy$ to obtain the other
assertions: each displayed cyclic subgroup
injects into $A$ by its primitive abelian vector.
\end{proof}

\subsection{The remaining normalizer cannot admit a virtual arithmetic section}

\begin{statement}{A.61}{Proved}{virtually split procyclic normalizers are cuspidal}
\label{st:A.61}
\label{A:split-normalizer}
Let $U/K$ be an affine hyperbolic curve over a number
field, with arithmetic group $E$ and geometric group
$\pi$. Let $A\subset\pi$ be a nontrivial closed
procyclic subgroup whose arithmetic normalizer has
open image. Then the following conditions are equivalent:
\begin{enumerate}
\item $A$ is contained in a cusp inertia group;
\item for some finite extension $L/K$, the map
\[
 N_E(A)\cap p^{-1}(G_L)\longrightarrow G_L
\]
has a continuous section.
\end{enumerate}
Consequently every noncuspidal normalizer remaining
in Propositions~\ref{A:full-support-normalizer} and
\ref{A:cm-normalizer-field} fails to admit a section
after every finite arithmetic base extension.
The implication from the second condition uses the
supplied local section theorem through
Proposition~\ref{F-local-dichotomy}; its stated
unrefereed status is retained.
\end{statement}
\begin{proof}
For the first implication, pass to a finite extension
making the cusp rational. A nonzero rational tangent
gives a section of its decomposition group. The
cyclotomic action preserves every closed subgroup
of its procyclic inertia, so this section takes
values in $N_E(A)$.

Conversely, replace $K$ by the finite extension
in the second condition and write
$s:G_K\to N_E(A)$ for its actual section.
Choose a prime $\ell$ with $A_\ell\ne1$.
The torsion-freeness of the free profinite group
$\pi$ gives $A_\ell\simeq\mathbf Z_\ell$.
Choose a characteristic open subgroup $V\subset\pi$
such that a generator of $A_\ell$ has nontrivial
image in $\pi/V$, and set $H=V A_\ell$.
Then $H$ is open and is normalized by $s(G_K)$.
The actual open arithmetic subgroup
\[
                         Hs(G_K)\subset E
\]
therefore defines a geometrically connected finite
\'etale cover $Y\to U$ over $K$. Its geometric
group is $H$. The cyclic quotient $H/V$ detects
$A_\ell$ in $H^{\mathrm{ab}}$.
Since $H^{\mathrm{ab}}$ is torsion-free, the resulting
map
\[
 A_\ell\longrightarrow H^{\mathrm{ab}}_{\ell}
       =H_1(Y_{\overline K},\mathbf Z_\ell)
                                                        \tag{A40}
\]
is injective. It is equivariant for the precise
action defined by $s$: conjugation by this section
is the ordinary arithmetic action on the
abelianization of the cover.

Let $\overline Y$ be the smooth compactification
and $J=\operatorname{Jac}(\overline Y)$. The
boundary exact sequence is
\[
 0\longrightarrow W\longrightarrow H_1(Y_{\overline K},\mathbf Z_\ell)
       \longrightarrow T_\ell J\longrightarrow0.
                                                        \tag{A41}
\]
If the composite of (A40) with $T_\ell J$ is zero,
then after splitting the finitely many cusps over
a finite extension its image lies in a direct
sum of copies of $\mathbf Z_\ell(1)$. The section
$s$ thus acts cyclotomically on $A_\ell$.
Proposition~\ref{A:cyclotomic-normalizer} puts
$A_\ell$ in one cusp inertia $I$.
Every element of $A$ centralizes every element
of $A_\ell$, and the full geometric centralizer
of any nontrivial element of $I$ is $I$.
Hence $A\subset I$.

It remains to treat a nonzero composite
$i:A_\ell\hookrightarrow T_\ell J$.
It is injective because $T_\ell J$ is torsion-free.
Faltings's semisimplicity theorem over this number
field, in its exact rational Tate-module form
\cite[\S5, Satz 3]{Faltings1983}, supplies a
$G_K$-equivariant retraction of
$i\otimes\mathbf Q_\ell$. Clearing denominators
in the two finite free lattices gives a nonzero
continuous $G_K$-homomorphism
\[
                 r:T_\ell J\longrightarrow A_\ell.
                                                        \tag{A42}
\]
The action on the target of (A42) is exactly
conjugation by $s$, by (A40).

Use the zero point of $J$ to identify its actual
arithmetic extension with
\[
 E_J=\left(\prod_qT_qJ\right)\rtimes G_K.
\]
The formula
\[
 f:E_J\longrightarrow E,
       \qquad f(v,g)=r(v_\ell)s(g)                       \tag{A43}
\]
is a continuous arithmetic homomorphism: (A42)
is precisely the equivariance needed for its
multiplication law. Its geometric image is a
nonzero closed subgroup of $A_\ell$, hence is
procyclic. It is not open in the nonabelian
geometric group of a hyperbolic curve.
Proposition~\ref{F-local-dichotomy}, whose source
may be any smooth geometrically connected
variety, puts this image in one cusp inertia
$I\subset\pi$. Finally $A$ centralizes that
nonzero image. The same full inertia centralizer
identity gives $A\subset I$.
This proves the second implication and the
stated obstruction to every virtual splitting.
\end{proof}

\begin{statement}{A.62}{Proved}{the remaining normalizer section problem has prime-to-prime geometric kernel}
\label{st:A.62}
\label{A:normalizer-prime-complement}
Keep the number-field and affine hypotheses above.
Let $A\subset\pi$ be noncuspidal, procyclic and
nontrivial, with open arithmetic normalizer image,
and choose $\ell$ with $A_\ell\ne1$.
There exist an open subgroup $N_0\subset N_E(A)$,
a finite extension $L/K$ with $p(N_0)=G_L$,
and closed subgroups $P,C,B\subset N_0$ such that
\[
 \begin{gathered}
 N_0\cap\pi=P\times C,\qquad
 1\ne P\subset A_\ell,\quad P\simeq\mathbf Z_\ell,
       \quad C\text{ is pro-prime-to-}\ell,\\
 N_0=P\rtimes B,\qquad
 1\longrightarrow C\longrightarrow B
       \xrightarrow{p}G_L\longrightarrow1.
 \end{gathered}                                             \tag{A44}
\]
In particular the quotient $N_0/C$ has an actual
arithmetic section. The still required section of
$N_0\to G_L$ is the explicit lifting problem
through the last extension in (A44).
Using Proposition~\ref{A:split-normalizer}, that
extension has no section after any finite extension
of $L$; in particular $C\ne1$.
The structure and quotient splitting in (A44) use
Faltings semisimplicity and the cyclotomic recognition
theorem. Only the final non-splitting conclusion
also uses the supplied local section theorem.
\end{statement}
\begin{proof}
Initially write $N=N_E(A)$ and $D=N\cap\pi$.
The action $\theta:N\to\mathbf Z_\ell^\times$
on $A_\ell$ has finite prime-to-$\ell$ image
on $D$, by Lemma~\ref{A:normalizer-character-descent}.
Choose an open subgroup of
$\mathbf Z_\ell^\times$ meeting that finite image
only in one, and replace $N$ by its inverse
image. Then $D$ centralizes $A_\ell$.
Let $P^*$ be a pro-$\ell$ Sylow subgroup of $D$.
It contains the normal subgroup $A_\ell$.
As a closed pro-$\ell$ subgroup of the free
profinite group $\pi$, it is free pro-$\ell$.
A noncyclic free pro-$\ell$ group has trivial
center, whereas $A_\ell$ is central in $P^*$.
Thus $P^*\simeq\mathbf Z_\ell$ and $A_\ell$
is open in $P^*$.

Here is the precise shrinking needed to make the
Sylow subgroup normal; torsion-freeness alone
does not identify the different Sylow subgroups.
The compact set $P^*\setminus A_\ell$ avoids
the identity. Choose an open normal subgroup
$N_1\triangleleft N$ disjoint from this set.
Put $D_1=N_1\cap D$ and $P_1=P^*\cap D_1$.
Then $D_1$ is normal in $D$, $P_1$ is a Sylow
subgroup of $D_1$, and
\[
                 P_1\subset A_\ell\subset Z(D).
\]
Conjugacy of Sylow subgroups now makes $P_1$
the unique Sylow subgroup of $D_1$.
It is nontrivial, because it is open in $P^*$.
The normal central Sylow subgroup has a unique
prime-to-$\ell$ complement $C_1$, and
$D_1=P_1\times C_1$.
For completeness this splitting can be checked
at every finite quotient: the central
$\ell$-subgroup extension of a group of order
prime to $\ell$ splits by the averaging
calculation of $H^2$, and the complement is
unique since a homomorphism from a prime-to-$\ell$
group to an $\ell$-group is zero. These unique
complements are compatible under the finite
quotient maps. Their inverse limit is $C_1$.
Uniqueness also makes $C_1$ characteristic in
$D_1$ and normal in $N_1$.

Choose a characteristic open $V\subset\pi$
detecting a generator of $P_1$, and set
$H=V P_1$. Descend its finite geometric cover
to a finite number-field extension and choose
its actual arithmetic subgroup $E_Y\subset E$,
with $E_Y\cap\pi=H$. Intersect $N_1$ with
$E_Y$. Its geometric kernel is
\[
 (P_1\times C_1)\cap H
       =(P_1\cap H)\times(C_1\cap H).
                                                        \tag{A45}
\]
Indeed closed subgroups of a product of a
pro-$\ell$ and a pro-prime-to-$\ell$ group
split into their two factors: the two primary
components of each element are limits of its
integer powers. The same assertion holds
after every subsequent finite arithmetic
restriction. We may therefore keep the notation
$N,D=P\times C$ for this intersection, with
$P\subset P_1$ nontrivial. Its arithmetic
image is still open; replace the field by
that image's fixed field.

The detector construction gives a nonzero,
hence injective, map
$P\to H_1(Y_{\overline K},\mathbf Z_\ell)$.
Its arithmetic action is the action of $N$:
the kernel $D\subset H$ acts trivially on this
abelianization, so the action factors through
the displayed arithmetic quotient.
If its image in the compactified Tate module
were zero, after splitting the boundary it
would be cyclotomic. The third condition of
Proposition~\ref{A:cyclotomic-normalizer} would
then put $P$ in cusp inertia, and its full
centralizer would put all of $A$ there.
This contradicts the assumption. Consequently
\[
             i:P\hookrightarrow T_\ell J,
       \qquad J=\operatorname{Jac}(\overline Y),
                                                        \tag{A46}
\]
is injective and equivariant for the actual
quotient action of $G_L$ on $P$.

Pass to a further finite extension so that
$\overline Y$ has a rational point $y_0$.
The Abel--Jacobi morphism
$Y\to J$, $y\mapsto[y-y_0]$, gives an actual
arithmetic group homomorphism
\[
                  E_Y\longrightarrow
                  \left(\prod_qT_qJ\right)\rtimes G_L.
\]
Faltings semisimplicity and clearing denominators
give an equivariant $r:T_\ell J\to P$ with
$r i=\ell^k\operatorname{id}_P$ for some $k\ge0$.
After projecting to $T_\ell J$ and applying $r$,
the preceding arithmetic map restricts to
\[
                  \varphi:N\longrightarrow P\rtimes G_L.
                                                        \tag{A47}
\]
It kills $C$, since $C$ is pro-prime-to-$\ell$
and its geometric image is pro-$\ell$. It
therefore induces an injective homomorphism
\[
             \overline\varphi:N/C\hookrightarrow P\rtimes G_L:
                                                        \tag{A48}
\]
any element of its kernel lies in $D/C=P$,
where its map is multiplication by $\ell^k$.
This multiplication is injective.

The image $Q$ of (A48) projects onto $G_L$
and meets the geometric $P$ in $\ell^kP$.
Modulo $\ell^kP$ it is the graph of a
continuous finite-valued cocycle
$b:G_L\to P/\ell^kP$. On the open subgroup
where $b$ vanishes, $Q$ contains $(0,g)$
for every $g$: subtract the geometric
coordinate in $\ell^kP$ from any lift.
Taking an open normal subgroup inside this
kernel and replacing $L$ accordingly gives
an actual continuous section of $N/C$,
by the inverse of the homeomorphism (A48)
onto its compact image.

Let $B$ be the inverse image of this section
in $N$. Then $B\cap\pi=C$, $B\to G_L$
is onto, $B\cap P=1$, and $N=PB$.
As $P$ is normal, multiplication identifies
$N$ with $P\rtimes B$, proving (A44).
A section of $B\to G_L$ would be a section
of $N_E(A)\to G_L$, contradicting
Proposition~\ref{A:split-normalizer}; the same
argument applies after every finite extension.
\end{proof}

\begin{statement}{A.63}{Refuted}{a prime-to-prime kernel alone does not imply virtual splitting}
\label{st:A.63}
\label{A:prime-complement-counterexample}
For every prime $\ell$ there is a profinite extension
\[
 1\longrightarrow\mathbf Z_q(1)\longrightarrow
 \mathcal B\longrightarrow G_{\mathbf Q}
       \longrightarrow1,
       \qquad q\ne\ell,
                                                        \tag{A49}
\]
whose kernel is pro-prime-to-$\ell$ and which has
no continuous section after any finite extension
of $\mathbf Q$. Thus the last extension in (A44)
cannot be split merely from the prime support
of its kernel. No embedding of (A49) into an
arithmetic curve group is asserted.
\end{statement}
\begin{proof}
Choose an odd prime $q\ne\ell$ and two distinct
finite places $v,w$ of $\mathbf Q$.
For each $n\ge1$, the Brauer--Hasse--Noether
sequence gives a unique Brauer class with local
invariants
\[
 \operatorname{inv}_v(\beta_n)=q^{-n},\qquad
 \operatorname{inv}_w(\beta_n)=-q^{-n},\qquad
 \operatorname{inv}_u(\beta_n)=0\quad(u\ne v,w).
                                                        \tag{A50}
\]
The finite-coefficient Kummer sequence and Hilbert
90 identify this class with
$\beta_n\in H^2(G_{\mathbf Q},\mu_{q^n})$.
The coefficient transition
$\mu_{q^{n+1}}\to\mu_{q^n}$ induces multiplication
by $q$ on the Brauer group. Hence
$q\beta_{n+1}=\beta_n$ under these identifications.
The exact global sequence and the local
restriction formula used here are
\cite[IV, 2.14(e); III, (29)]{A_MilneCFT}.

We construct compatible actual continuous
cocycles, so that no finite-to-continuous
existence assertion is left implicit.
Suppose a normalized continuous two-cocycle
$z_n$ representing $\beta_n$ has been chosen.
Choose any such representative $z'_{n+1}$
of $\beta_{n+1}$. Its reduction differs from
$z_n$ by $\delta u_n$ for a continuous
one-cochain $u_n$ with values in $\mu_{q^n}$.
This finite-valued locally constant cochain
has a continuous lift $\widetilde u_n$ to
$\mu_{q^{n+1}}$, obtained by choosing a lift
of each of its finitely many values, and choosing
zero as the lift of zero. Since the discrepancy
is normalized, $u_n(1)=0$, so normalization
is preserved.
Replace $z'_{n+1}$ by
$z_{n+1}=z'_{n+1}-\delta\widetilde u_n$.
These representatives are now compatible.
Their inverse limit is a normalized continuous
two-cocycle
\[
 z:G_{\mathbf Q}^2\longrightarrow\mathbf Z_q(1).
\]
Continuity is exactly continuity of every
finite coordinate in the inverse-limit
topology. This construction lifts the specified
compatible nonzero obstruction classes; it does
not assert that vanishing finite obstruction
classes force vanishing of a continuous one.

Give $\mathbf Z_q(1)\times G_{\mathbf Q}$ the
product topology and multiplication
\[
 (a,g)(b,h)=\bigl(a+g b+z(g,h),gh\bigr).
                                                        \tag{A51}
\]
The cocycle identity proves associativity;
normalization gives the identity, and the
inverse is obtained by solving the displayed
formula. This is a compact totally disconnected
topological group, hence a profinite group,
and it gives (A49).

For a finite extension $L/\mathbf Q$ and a
place $v'\mid v$, local restriction multiplies
the invariant in (A50) by the positive integer
$[L_{v'}:\mathbf Q_v]$. For sufficiently large
$n$ the result
\[
                 [L_{v'}:\mathbf Q_v]/q^n
                         \quad\text{in }\mathbf Q/\mathbf Z
\]
is nonzero. Thus $\beta_n|_{G_L}\ne0$.
A section of (A49) over $G_L$ would make
$z|_{G_L}$ a continuous coboundary, and every
one of its finite reductions would then be
a coboundary. This contradiction proves
failure of splitting over every finite $L$.
\end{proof}

\subsection{The corresponding restrictions for proper hyperbolic curves}

\begin{statement}{A.64}{Known}{finitely generated pro-prime subgroups of a full surface group}
\label{st:A.64}
\label{A:surface-proprime-free}
Let $\Pi$ be the full profinite fundamental group
of a proper hyperbolic curve of genus $g\ge2$
over an algebraically closed field of characteristic
zero. Every finitely generated closed pro-$\ell$
subgroup $P\subset\Pi$ is free pro-$\ell$.
In particular, every finitely generated pro-$\ell$
subgroup having a nontrivial procyclic normal
subgroup is procyclic.
\end{statement}
\begin{proof}
We give a direct surface-cohomology proof; the
broader pro-prime freeness theorem for profinite
completions of torsion-free hyperbolic virtually
special groups is \cite[Theorem F]{A_WiltonZalesskii2017}.
For a closed subgroup $P$ and finite trivial
coefficients, continuous cohomology satisfies
\[
 H^i(P,\mathbf F_\ell)
       =\varinjlim_{P\subset U\subset\Pi\ \mathrm{open}}
                          H^i(U,\mathbf F_\ell).          \tag{A52}
\]
Here is the continuity argument needed for this
specific use. A finite-valued continuous cochain
on the closed profinite subset $P^i$ extends
to a continuous cochain on $\Pi^i$, by extending
its finite clopen partition. Its cocycle
equations hold on $P^{i+1}$; compactness and
local constancy make them hold on $U^{i+1}$
for an open subgroup $U=P V$, with $V$ open
normal in $\Pi$. The same extension and
shrinking argument for a cochain whose
coboundary is prescribed proves injectivity
of the direct-limit comparison. This proves
(A52), for both $i=1$ and $i=2$.

For such an open $U$, let $g_U$ be the genus
of its actual finite covering curve. The
finite-coefficient group-to-curve comparison
can be checked directly here. In the tower of
all finite \'etale covers, every $H^1$ class
is killed by its finite torsor, and every
$H^2$ class is killed by a further cover of
degree $\ell$. Such a cover exists from
the nonzero abelianization of every proper
hyperbolic covering curve. Higher curve
cohomology vanishes. The universal-cover
Cartan--Leray comparison therefore identifies
the continuous group cohomology with the
curve cohomology, compatibly with pullback.
The
orientation class and the covering trace give
\[
 \dim H^1(U,\mathbf F_\ell)=2g_U,
 \quad H^2(U,\mathbf F_\ell)=\mathbf F_\ell,
 \quad g_U-1=[\Pi:U](g-1).                                \tag{A53}
\]
For $V\subset U$, restriction on $H^2$ is
multiplication by $[U:V]$ under these orientation
identifications. These are the actual
cohomological maps of the finite \'etale
coverings, or equivalently of their complex
surface coverings.

Suppose first that the integers
$v_\ell([\Pi:U])$ for open $U\supset P$
are bounded. Choose $U_0$ attaining their
maximum. Every smaller open $V\subset U_0$
containing $P$ has index prime to $\ell$ in
$U_0$, and the same holds for every further
transition. Each restriction on $H^1$ is
injective: a nonzero homomorphism
$V\to\mathbf F_\ell$ has kernel of index
$\ell$, so cannot vanish on a subgroup of
prime-to-$\ell$ index.

The subgroup $P$ cannot be open in $\Pi$,
because every open surface group has
abelianization $\widehat{\mathbf Z}^{\,2g_U}$
and thus has nontrivial quotients at primes
different from $\ell$. Therefore the indices
$[\Pi:V]$ in this cofinal family are unbounded.
By (A53), their $H^1$ dimensions are unbounded,
and all these spaces inject into the direct
limit in (A52). This contradicts finite
generation of the pro-$\ell$ group $P$.

Thus the $\ell$-parts of the indices are
unbounded. Given any open $U\supset P$, choose
a further open $V\subset U$ containing $P$
with $\ell\mid[U:V]$; intersecting $U$ with
an overgroup of sufficiently large
$\ell$-index gives such a $V$.
Its restriction map on $H^2$ is zero by
(A53). Formula (A52) consequently gives
$H^2(P,\mathbf F_\ell)=0$, which is the
pro-$\ell$ freeness criterion.

For the last assertion, let $C\ne1$ be the
procyclic normal subgroup. For each element
$h$ the subgroup generated by $C$ and $h$
is solvable and free pro-$\ell$, hence
procyclic. Thus every $h$ centralizes $C$.
A noncyclic free pro-$\ell$ group has trivial
center, so the whole finitely generated group
is procyclic.
\end{proof}

\begin{statement}{A.65}{Proved}{proper-curve normalizers obey the same character and splitting restrictions}
\label{st:A.65}
\label{A:proper-normalizer-restrictions}
Let $X/K$ be proper hyperbolic over a number
field, and let $A\ne1$ be a closed procyclic
subgroup of its full geometric group whose
arithmetic normalizer has open image $G_L$.
Then the conclusions about arithmetic character
descent, excluded primes, and CM fields in
Lemma~\ref{A:normalizer-character-descent} and
Propositions~\ref{A:full-support-normalizer}--\ref{A:cm-normalizer-field}
hold, with the boundary alternative absent.
Explicitly, $L$ contains a CM subfield, the
prime support of $A$ omits the positive-density
complex-conjugation Frobenius set described
there, and each supported $\ell$ has the
stated CM-field embedding into $\mathbf Q_\ell$.
In particular no such $A$ has cofinite prime
support, and there is none when $L$ contains
no CM subfield.

Its normalizer has no section after any finite
arithmetic base extension, using the supplied
local theorem as in Proposition~\ref{A:split-normalizer}.
For every supported $\ell$ it also has the
exact open-subgroup structure (A44), with
the same residual pro-prime-to-$\ell$
extension and no virtual section.
These conclusions do not exclude all of
the remaining CM and restricted-support cases.
\end{statement}
\begin{proof}
We specify the replacements of the affine
inputs, so the assertions compare the same
full geometric groups and actual arithmetic
maps. For a pro-$\ell$ Sylow subgroup $S$
of the geometric normalizer and any $h\in S$,
the subgroup $\langle A_\ell,h\rangle$ is a
finitely generated pro-$\ell$ subgroup of the
full surface group. It is free by
Lemma~\ref{A:surface-proprime-free}, and
solvable because its quotient by $A_\ell$
is procyclic. It is therefore cyclic.
This proves precisely the triviality of the
pro-$\ell$ part of the geometric normalizer
character used in
Lemma~\ref{A:normalizer-character-descent}.
The remaining image is finite and prime to
$\ell$, and the same power descends to the
open arithmetic quotient.

Choose the same detector $H=V A_\ell$, with
$V$ characteristic open in the surface group.
Its abelianization is again torsion-free,
now $\widehat{\mathbf Z}^{\,2g_H}$, and the
cyclic quotient $H/V$ gives the identical
injective homology line. The cover is proper,
so this homology is already the Tate module
of its Jacobian. Thus its weights, its
Hodge--Tate weights, and the descended power
character are exactly those in the compact
case of Lemma~\ref{A:normalizer-character-descent}.
The Serre algebraic-type proof in
Propositions~\ref{A:conjugation-prime-character}--\ref{A:cm-normalizer-field}
then applies to these actual characters.
There is no toric summand and no cusp case.
This proves every stated field and prime-support
restriction.

An actual arithmetic section in the normalizer
would give, by the same detector and Faltings
retraction, the explicit arithmetic map
(A43) from a Jacobian with nonzero procyclic
geometric image. Proposition~\ref{F-local-dichotomy}
has no cusp alternative for the proper
target $X$, and such an image is not open.
It therefore forces the image to be trivial,
a contradiction. This proves failure of
every virtual normalizer section.

Finally the Sylow step of
Proposition~\ref{A:normalizer-prime-complement}
also survives. After the geometric character
has been killed, every finitely generated
subgroup of $S$ containing a generator of
the central $A_\ell$ is free pro-$\ell$ with
nontrivial center, hence cyclic. Every finite
quotient of $S$ is consequently cyclic,
so $S$ itself is procyclic. The exact open
normal shrinking, central Sylow complement,
and intersection formula (A45) now apply.
The detector is proper and its compact line
is nonzero automatically. Its Abel--Jacobi
map and the integral Faltings retraction
give (A47)--(A48), and the same finite
cocycle gives the quotient section and
(A44). The non-splitting of the remaining
extension follows from the preceding
proper-target contradiction.
\end{proof}

\subsection{Deck-twisted fibre isolation and the fourth-to-third derived projection}

\begin{statement}{A.66}{Proved}{one derived layer forces preservation of every element conjugacy class}
\label{st:A.66}\label{A:one-layer-class-preservation}
Let $K$ be a number field, let
$U=\mathbf P^1_K\setminus\{0,1,\infty\}$, and put
$S_j=F/F^{(j)}$ for its geometric derived quotients.
Write $C_j^\partial(K)$ for the centralizer of the
arithmetic outer action in the outer automorphisms of
$S_j$ preserving the three peripheral conjugacy classes
as an unordered set. For every $m\geq2$ and
$\gamma\in C_{m+1}^\partial(K)$, there is a unique
geometric cusp permutation $\sigma\in S_3$ such that
the projection of $\sigma^{-1}\gamma$ to
$\operatorname{Out}(S_m)$ has a representative
$\alpha$ satisfying
\[
 \text{for every }z\in S_m,\qquad
             \alpha(z)\text{ is conjugate to }z
             \text{ in }S_m.                         \tag{A66.1}
\]
In particular $\alpha$ preserves every closed normal
subgroup and every conjugacy class of closed procyclic
subgroups. The assertion in (A66.1) allows its conjugator
to depend on $z$.
\end{statement}
\begin{proof}
Proposition~\ref{A:third-derived-projection} determines
the metabelian image of $\gamma$ as a unique cusp
permutation: uniqueness follows from the faithful
action of $S_3$ on $F^{\mathrm{ab}}$.
Remove that permutation and conjugate a representative
$\phi$ on $S_{m+1}$ geometrically so that its induced
automorphism on $M=S_2$ is literally the identity.
The center-free arithmetic-extension construction of
Proposition~\ref{A:solvable-centers} gives an actual
automorphism $\beta$ of the arithmetic extension
$E_{m+1}$ over $G_K$. Use the tangential section at
$0$ to write
\[
 \beta(s(g))=d_gs(g),\qquad
 \phi\rho_g\phi^{-1}=\operatorname{Int}(d_g)\rho_g.
\]
The projection of $d_g$ to $M$ is central, because
$\phi_M=1$. Since $Z(M)=1$, it is trivial.
Consequently the induced arithmetic automorphism
on $E_2$, and hence on $E_1$, is exactly the
identity. This keeps the constants in every Kummer
cover fixed.

Choose a characteristic open normal subgroup of
$S_m$ and denote its quotient by $G$.
Such subgroups are cofinal because $S_m$ is finitely
generated. The automorphisms induced by $\phi$ and
the actual tangential arithmetic action both act
on $G$. After a finite extension $K_0/K$, that
arithmetic action is trivial, so the corresponding
arithmetic quotient is
$G\times G_{K_0}$. The induced $\beta$ has the form
\[
 (g,\tau)\longmapsto
             \bigl(\alpha_G(g)z(\tau),\tau\bigr),
       \qquad z:G_{K_0}\longrightarrow Z(G).
                                                        \tag{A66.2}
\]
Indeed $(1,\tau)$ centralizes $G$, so its image
does also; the homomorphism law makes $z$ a
continuous homomorphism into the finite center.
A further finite extension kills $z$.
Thus (A66.2) becomes $(\alpha_G,\mathrm{id})$.
No center-free assumption on the finite group
$G$ has been made.

Let $C\to U_{K_0}$ be the resulting geometrically
connected finite etale $G$-torsor, with its constant
deck group. Fix an element $g\in G$.
Hilbert irreducibility supplies $a\in U(K_0)$
whose fibre is a connected $G$-torsor.
Its splitting field is a $G$-extension of $K_0$.
Apply Chebotarev in a normal closure over
$\mathbb Q$, extending the chosen $g$ while
fixing $K_0$. This supplies a place of $K_0$
of residue degree one over a rational prime $q$,
at which the fibre Frobenius is conjugate to $g$.
Exclude the finitely many bad places for the
marked cover, its deck maps, and $a$, and those
dividing $|G|$. These are exactly the Hilbert and
degree-one Chebotarev steps used in
Proposition~\ref{A:third-derived-IA}.

Put $N=q-1$ and enlarge the constants to
$K_1=K_0(\mu_N)$. The chosen place has an
extension still with residue field $\mathbb F_q$,
since all $N$th roots of unity lie there.
Over this field consider
\[
 B:\qquad v_1^N=t/a,\qquad
             v_2^N=(1-t)/(1-a),\qquad
                   Z=C\times_U B.                    \tag{A66.3}
\]
These are finite etale covers of the same tripod.
They have smooth marked reduction at the chosen
place: the original marked cover has good
reduction there, and the newly added ramification
indices divide $q-1$. Normalization, with the
punctures removed, gives the indicated models.
The arithmetic automorphism $\beta$ preserves
the two equations of $B$, including their
constant Kummer twists, because it induces the
identity on $E_1$.
It acts on the finite fibre of $Z$ by applying
$\alpha_G$ to its $G$-coordinate and leaving the
$B$-coordinate unchanged.

We give the cohomological comparison for every
component, without assuming that $Z$ is
geometrically connected. Its geometric monodromy
is a subgroup of $G\times(\mathbb Z/N)^2$ and
has derived length at most $m$.
The subgroup $H\subset F$ corresponding to any
component therefore contains $F^{(m)}$.
It follows that
\[
 F^{(m+1)}=[F^{(m)},F^{(m)}]\subset H',
 \qquad
 (H/F^{(m+1)})^{\mathrm{ab}}=H^{\mathrm{ab}}.
                                                        \tag{A66.4}
\]
Thus $E_{m+1}$ retains the full first homology
of each component, together with the maps
between components induced by arithmetic and
deck transformations. More explicitly, components
are the orbits of $F$ on the finite fibre;
$H$ is the stabilizer of a point of an orbit,
and changing that point conjugates $H$.
The automorphism $\beta$ maps the orbit and
stabilizer diagrams to themselves. The canonical
identification in (A66.4) commutes with these
conjugation maps. It therefore gives an
isomorphism on the direct sum of the first
cohomology groups which commutes with arithmetic
and intertwines a deck transformation $h\in G$
with $\alpha_G(h)$.
The same finite orbit diagram identifies
$H^0$. Poincare duality on each smooth component,
with its canonical orientation, then identifies
$H_c^1$ and $H_c^2$ equivariantly as well.
For any auxiliary coefficient prime different
from $q$, good marked reduction transfers these
identifications to the special fibre.

Let $\mathcal F$ denote its Frobenius operator
in the fixed-point convention. The preceding
comparison gives equality of alternating traces
on compactly supported cohomology for
$h\mathcal F$ and $\alpha_G(h)\mathcal F$.
These alternating traces count their fixed
geometric points. To apply the usual
Lefschetz trace formula here, one may twist the
finite-field descent datum by the finite-order
deck map $h$: it is defined over $\mathbb F_q$
and commutes with Frobenius, so the twisted
descent is effective and its rational points
are precisely the fixed points of $h\mathcal F$.
This also handles the possible permutation of
components.

All such fixed points lie above $t\in U(\mathbb F_q)$.
The deck map $h$ is trivial on $B$.
Since every nonzero element of $\mathbb F_q$
has $N$th power one, (A66.3) has rational points
over exactly $t=\overline a$, with $N^2$ points
in that fibre. Label the regular $G$-fibre
of $C$ there by $G$ so that Frobenius is right
multiplication by $g$ and the deck map $h$ is
left multiplication by $h$.
Its fixed-point condition is $hkg=k$.
Consequently
\[
 \#\operatorname{Fix}(h\mathcal F\mid Z)
  =N^2
    \begin{cases}
      |C_G(g)|,&h\text{ is conjugate to }g^{-1},\\
      0,&\text{otherwise}.
    \end{cases}                                      \tag{A66.5}
\]
Take $h=g^{-1}$ and compare it with
$\alpha_G(h)$. Equality of traces implies that
$\alpha_G(g)$ is conjugate to $g$.
The argument applies to every $g\in G$.

Finally fix $z\in S_m$. For each characteristic
finite quotient $G$, the set of $u\in S_m$ such
that $u z u^{-1}$ and $\alpha(z)$ have the same
image in $G$ is nonempty and closed.
The sets have the finite intersection property
by passage to a common finer characteristic
quotient. Compactness gives one conjugator for
this $z$ in $S_m$ itself. This proves (A66.1).
Applying it to the elements of a closed normal
subgroup shows preservation of that subgroup;
applying it to a topological generator gives
preservation of each procyclic conjugacy class.
\end{proof}

\begin{statement}{A.67}{Proved}{class-preserving rigidity at derived length three and the fourth-level projection}
\label{st:A.67}\label{A:third-derived-class-rigidity}
Every continuous automorphism of $S_2$ or $S_3$
which sends each element to a conjugate of itself
is inner. Consequently, for every number field
$K$,
\[
 \operatorname{im}\bigl(
 C_4^\partial(K)\longrightarrow
                  \operatorname{Out}(S_3)\bigr)=S_3,
                                                        \tag{A67.1}
\]
where the group on the right is the geometric
group of cusp permutations. The same equality
holds with $C_j^\partial(K)$ for any $j\geq4$.
This is an equality of projected commutants;
the kernel in the same derived level is not
computed by it.
\end{statement}
\begin{proof}
We first prove the completed-module facts on
which the rigidity argument depends.
For a profinite group $P$ write
$\Lambda_P=\widehat{\mathbb Z}[[P]]$.
If $u\in P$ has infinite $p$-primary order,
then both left and right multiplication by
$u-1$ have zero kernel on $\mathbb F_p[[P]]$.
For example, an element killed by left
multiplication is invariant under $\langle u\rangle$.
At a finite quotient it has constant
coefficients on the $u$-orbits.
Refine that quotient so that each such orbit
has its length multiplied by a factor divisible
by $p$. Projection of each invariant orbit sum
is then zero in characteristic $p$.
Hence its image at every original finite
quotient is zero. The right-hand assertion
uses the right orbits in the same way.
The required refinements exist because the
closed cyclic subgroup has its full induced
profinite topology and infinite $p$-part.
Successive division by $p$ gives the same
injectivity on $\mathbb Z_p[[P]]$.
If $u$ has infinite $p$-part for every $p$,
the assertion holds on $\Lambda_P$.

Suppose two commuting elements $u,v$ generate
a closed subgroup isomorphic to
$\widehat{\mathbb Z}^2$.
The quotient
$\Lambda_P/(u-1)\Lambda_P$ is the completed
free coefficient module on the left coset
space $\langle u\rangle\backslash P$.
One sees this at every finite quotient;
compactness identifies the kernel with
the closed image of multiplication by $u-1$.
Left multiplication by $v$ is well defined
on this space and has free procyclic orbits.
In finite quotient approximations its orbit
lengths can be refined by arbitrarily large
powers of any prime: the induced topology
on $\langle u,v\rangle/\langle u\rangle
\simeq\widehat{\mathbb Z}$ supplies these
refinements. The preceding orbit-sum proof
therefore also proves
\[
 v-1:\Lambda_P/(u-1)\Lambda_P
       \longrightarrow\Lambda_P/(u-1)\Lambda_P
             \quad\hbox{is injective}.                \tag{A67.2}
\]

Here is a degree-one cohomology consequence
in the exact topology needed below.
Suppose that $P$ has a normal abelian subgroup
$D$ containing such $u,v$.
For a continuous cocycle
$f:P\to\Lambda_P$, put $a=f(u)$ and $b=f(v)$.
The commuting relation gives
$(v-1)a=(u-1)b$.
By (A67.2), $a=(u-1)t$ for some
$t\in\Lambda_P$. The injectivity of $u-1$
then gives $b=(v-1)t$.
Subtract this coboundary. The new cocycle
vanishes on $\langle u,v\rangle$.
Commuting with $u$ forces its value at every
$d\in D$ to be $u$-invariant, and hence zero.
Normality of $D$ then forces its value at
every $g\in P$ to be $D$-invariant, again zero.
This proves, for continuous cocycles with
these compact coefficients,
\[
                  H^1(P,\Lambda_P)=0.                 \tag{A67.3}
\]
In particular (A67.3) holds for
$P=\widehat{\mathbb Z}^2$, with $u,v$ its
basis.

For $M=S_2$, let
$c=[x,y]$ and $b=xcx^{-1}$.
Its derived module is the free rank-one
module
$M'=\widehat{\mathbb Z}[[M^{\mathrm{ab}}]]c$.
The vectors $c,b$ correspond to $1,X$.
They generate a copy of
$\widehat{\mathbb Z}^2$: independence of
$1,X$ already follows after setting
$X^2=1$ in the completed group algebra.
Thus $D=M'$ supplies the normal abelian
subgroup needed for (A67.3).

Let first an automorphism of $S_2$ be
class-preserving. It induces the identity on
the abelianization. Its action on the free
rank-one derived module sends $c$ to $a c$
for some $a\in M^{\mathrm{ab}}$, since the
image of $c$ is conjugate to $c$.
Remove an inner automorphism lifting $a$.
The resulting automorphism is the identity
on both the abelian quotient and the derived
module. Its difference from the identity
is a continuous cocycle from
$\widehat{\mathbb Z}^2$ to that regular
module. By (A67.3), it is a coboundary, so
the automorphism is inner.

Now put
\[
 S=S_3,\qquad M=S_2,\qquad
 L=F''/F''',\qquad
 \Lambda=\widehat{\mathbb Z}[[M]].
\]
The Fox--Magnus exact sequence used in
Proposition~\ref{A:last-layer-unrestricted}
identifies $L$ with
\[
 \ker\bigl(\partial:\Lambda^2\to\Lambda\bigr),
 \qquad
       \partial(a_1,a_2)=a_1(x-1)+a_2(y-1).
                                                        \tag{A67.4}
\]
Coefficients act on the left. The homomorphism
$\partial$ is left $M$-equivariant.
An $L$-valued cocycle, considered in
$\Lambda^2$, is a coboundary by (A67.3).
If its bounding vector is $w$, then
$(g-1)\partial w=0$ for all $g\in M$.
The invariants of $\Lambda$ are zero, by
injectivity of $c-1$. Therefore $\partial w=0$,
so $w\in L$. We have proved
\[
                         H^1(M,L)=0.                 \tag{A67.5}
\]

We next construct a specific element of $L$
with no left annihilator over $\Lambda$,
including after reduction at every coefficient
prime. Regard
$r=[c,xcx^{-1}]$ as an element of
$F''/F'''$.
The left Fox rule
$\delta(uv)=\delta(u)+\overline u\,\delta(v)$
gives
\[
 \delta_y(c)=x-c,\qquad
 \delta_y(r)=A(x-c),\qquad
 A=(c-1)x-(b-1).
                                                        \tag{A67.6}
\]
The second identity uses that $c,b$ commute
in $M$. For $d=c-1$ we have
$d^x=b-1=xdx^{-1}$, and a literal multiplication
gives
\[
                   A d=d\,d^x(x-1).                   \tag{A67.7}
\]
All of $c,b,x$ and $xc^{-1}$ have infinite
$p$-part for every prime $p$:
this follows from the free derived-module
description for $c,b$ and from the
abelianization for $x,xc^{-1}$.
Hence each factor on the right of (A67.7)
has zero left annihilator, even over
$\mathbb F_p[[M]]$.
If $zA=0$, multiply on the right by $d$
and cancel these injective right
multiplications successively; this gives
$z=0$. Since $x-c=(xc^{-1}-1)c$ also has
zero left annihilator, (A67.6) proves that
the second coordinate of $r$ has zero
left annihilator. Thus
$\Lambda\to L$, $z\mapsto zr$, is injective.
Only this specified sided regularity is used;
the completed group algebra is not declared
to be a domain.

Let $\alpha$ be a class-preserving automorphism
of $S$. Its induced automorphism on $M$ is
inner by the metabelian argument already proved.
After removing that inner automorphism, assume
$\alpha_M=1$.
Its action $T$ on $L$ is then $\Lambda$-linear.
For any $v\in L$ generating a free
$\Lambda$-submodule, class preservation says
$T(v)=g v$ for some $g\in M$.
Apply class preservation also to
$(1+2h)v$ for arbitrary $h\in M$.
Freeness of the cyclic module generated by
$v$ gives, for some $g_h\in M$,
\[
                    g+2hg=g_h+2g_hh.                  \tag{A67.8}
\]
Distinct elements of $M$ are linearly
independent in $\widehat{\mathbb Z}[[M]]$,
as finite quotients separate any finite
list of them. When $h\ne1$, the two terms
on each side of (A67.8) have distinct
coefficient vectors $1$ and $2$.
It follows that $g_h=g$ and $hg=gh$.
This also holds for $h=1$.
Thus $g\in Z(M)=1$, and $T(v)=v$.
This coefficient comparison takes place
over $\widehat{\mathbb Z}$, so it includes
the coefficient prime two.

To apply this conclusion to every element
of $L$, choose the central scalar
$\nu\in\widehat{\mathbb Z}$ whose
$p$-coordinate is $p$ for every prime $p$.
For any $l\in L$, the second coordinate
of $r+\nu l$, in the coefficient-$p$
completed group algebra, reduces modulo
$p$ to the second coordinate of $r$.
The latter has zero left annihilator by
(A67.6)--(A67.7).
Successive division by $p$ shows that
the second coordinate of $r+\nu l$ has
zero left annihilator over
$\mathbb Z_p[[M]]$.
Taking the product over all $p$ proves
that $r+\nu l$ generates a free
$\Lambda$-submodule.
The preceding paragraph fixes both $r$
and $r+\nu l$, so $\nu(T(l)-l)=0$.
Multiplication by $\nu$ is injective
on $\Lambda^2$, and hence on $L$:
each coefficient-$p$ factor is
$p$-torsion-free.
Therefore $T(l)=l$ for all $l\in L$.

An automorphism acting identically on both
$M$ and $L$ now gives the continuous cocycle
\[
 M\longrightarrow L,\qquad
   \overline s\longmapsto\alpha(s)s^{-1}.
                                                        \tag{A67.9}
\]
The expression is independent of the
lift $s$, since $L$ is abelian and
$\alpha|L=1$. It is continuous by
the quotient topology.
Equation (A67.5) makes (A67.9) a coboundary,
so $\alpha$ is inner by an element of $L$.
This proves the class-preserving rigidity
for $S_3$.

Finally apply
Proposition~\ref{A:one-layer-class-preservation}
with $m=3$. After the unique cusp permutation
is removed, the projected automorphism is
class-preserving and hence inner.
Every cusp permutation is geometric and
commutes with arithmetic, which supplies
the reverse inclusion in (A67.1).
For $j>4$, project first to level four.
\end{proof}

% ---- B ----
\section{Sections and points: local packets, abelian fibres, and descent}
\label{sec:B}

Throughout this section $U=X\setminus S$ is a smooth geometrically connected
curve, $X$ is its smooth proper completion, $\Delta_U=\pi_1(U_{\bar k})$,
and $\Pi_U=\pi_1(U)$.  A section is continuous and is taken modulo
conjugation by $\Delta_U$.  We write
\[
 \widehat{k^\times}=\varprojlim_{n\mid m}k^\times/(k^\times)^n.
\]
This is the power completion.  Its finite-level quotients need not be finite
over a number field.  Over a finite extension of $\mathbb Q_p$ they are
finite.  These two assertions use the same inverse system and do not
identify its topology with an unspecified completion.

\subsection{Injectivity, including the point germs of cote 143}

\begin{statement}{B.1}{Proved}{The corrected semi-abelian Kummer input}
\label{st:B.1}

Let $k$ be finitely generated over $\mathbb Q$ and let
$1\to T\to G\to A\to1$ be a semi-abelian variety over $k$.
Then
\[
 \bigcap_{n\geq1}nG(k)=0,
 \qquad G(k)\longrightarrow H^1_{\rm cts}(G_k,TG)
 \quad\hbox{is injective}.
\]
Consequently, for an immersion $U\hookrightarrow G$, two $k$-points with
the same abelianized section class are equal.  The corresponding map on
algebraic points and germs of sections is injective.
\end{statement}
\begin{proof}
First, $\bigcap_n(k^\times)^n=\{1\}$.  Let $F$ be the algebraic closure of
$\mathbb Q$ in $k$, a number field.  On a normal proper integral $F$-model
of $k$, the valuation of an element admitting every root is divisible by
every positive integer at every prime divisor, hence is zero.  A rational
function with zero divisor on this model is a nonzero global regular
function and therefore belongs to $F^\times$.  All its roots in $k$ then
belong to $F$, by the definition of $F$.  Valuations at finite primes of
$F$ reduce the assertion to the unit group of $F$, which is a finite group
plus a finitely generated free abelian group.  The intersection of its
$n$th powers is $1$.  A finite splitting field for $T$ therefore proves
$\bigcap_n nT(k)=0$.

The Lang--N\'eron theorem and the Mordell--Weil theorem for the constant
trace give finite generation of $A(k)$.  Let $a\in\bigcap_n nG(k)$.
Its image in $A(k)$ is zero.  Let $M$ annihilate $A(k)_{\rm tors}$.
For each $n$, choose $b_n$ with $nb_n=a$.  Then $Mb_n\in T(k)$ and
$Ma=n(Mb_n)$.  Thus $Ma=0$.  The torsion subgroup of $G(k)$ is finite:
its image is contained in the finite group $A(k)_{\rm tors}$ and its
kernel in $T(k)_{\rm tors}$, which is finite after a finite splitting
extension.  Since every $b_n$ is now torsion, taking $n$ equal to the
exponent of $G(k)_{\rm tors}$ gives $a=0$.

The multiplication-by-$n$ Kummer sequence sends $a$ to zero precisely
when $a\in nG(k)$.  Passing over all $n$ proves the asserted injectivity.
The continuous-cohomology comparison is legitimate: the groups of
$k$-rational $n$-torsion are finite, so the degree-zero inverse system
is Mittag--Leffler and contributes no $\varprojlim^1$ term.
For $x,y\in U(k)$, equality of their abelian section classes gives zero
Kummer class for their difference in $G(k)$, hence $x=y$.
Finally, equality of two section germs means equality after restriction
to $G_L$ for one finite extension $L/k$.  Apply the same argument over
$L$.  This proves exactly the germ assertion in cote 143, scans 4--5,
including its stated semi-abelian embedding hypothesis.
\end{proof}


\begin{statement}{B.2}{Refuted}{Excluded types and explicit failures over larger fields}
\label{st:B.2}

For $\mathbb P^1$ and $\mathbb A^1$ in characteristic zero the geometric
fundamental group is trivial, so every rational point gives the same
section.  On the hyperbolic curve
$U_0=\mathbb P^1\setminus\{0,1,\infty\}$ over $\mathbb R$, the distinct
points $2$ and $3$ give the same section class.  For every prime $p$, put
\[
 k=\mathbb Q_p((t^{\mathbb Q})),\qquad x=2,\qquad y=2+t.
\]
This is a $p$-adically closed field which is not sub-$p$-adic, and $x,y$
give the same full arithmetic section class on $U_0/k$.
\end{statement}
\begin{proof}
The first assertion follows from the geometric simply connectedness of
the projective and affine line in characteristic zero.  The real path
from $2$ to $3$ inside $(1,\infty)$ is fixed by conjugation.  Its induced
isomorphism of geometric fibre functors identifies the two splittings.

The Hahn field has a henselian $t$-valuation with divisible value group
$\mathbb Q$ and residue field $\mathbb Q_p$.  The characterization in
Koenigsmann's \cite[\S1, Fact on $p$-adically closed fields]{B_Koenigsmann}
therefore makes $k$ $p$-adically closed.  In residue characteristic zero
there is no defect or wild inertia; divisibility of the value group
eliminates tame ramification.  Restriction and residue specialization
give $G_k\simeq G_{\mathbb Q_p}$.
Base change of the fundamental-group sequence of the constant curve
$U_0$ is consequently an isomorphism of extensions.  Henselian
specialization of finite \emph{\'etale} covers near $2$ identifies their
fibres at $2$ and $2+t$, compatibly with morphisms of covers and with
$G_k$.  This gives an isomorphism of fibre functors, and hence equality
of the section classes.  More explicitly, a simple root in a fibre at
$2$ has a unique lift in the prescribed residue class at $2+t$;
uniqueness makes these identifications compatible over the entire
finite-\'etale category.

An embedding of $k$ into a finitely generated extension $F/\mathbb Q_p$
would send $t\ne1$ to an element with an $n$th root for every $n$.
Such an element in $F$ has zero divisor on a normal proper model over
the algebraic closure of $\mathbb Q_p$ in $F$, hence lies in that finite
extension of $\mathbb Q_p$.  In a finite $p$-adic field
$F_0^\times\simeq\mathbb Z\times\mathcal O_{F_0}^\times$ and
$\bigcap_n(F_0^\times)^n=1$.  This contradicts injectivity of the
embedding.  This also supplies a direct verification of the required
``not sub-$p$-adic'' assertion.
\end{proof}

\subsection{The local construction and its tangential interpretation}

\begin{statement}{B.3}{Proved}{Uniformizers, tangent vectors, and Deligne's fibre functor}
\label{st:B.3}

Let $y\in X(k)$ and let $R$ be its henselian local ring, with fraction
field $L$ and residue field $k$ of characteristic zero.  A uniformizer
$t$ defines a splitting of $G_L\to G_k$.  Replacing $t$ by $t(1+h)$,
$h\in\mathfrak m_R$, preserves its inertia-conjugacy class.  The
splitting classes form a torsor under
$H^1(k,\widehat{\mathbb Z}(1))=\widehat{k^\times}$.
The classes defined by nonzero tangent vectors form the
$k^\times$-Kummer orbit in this torsor and are precisely Deligne's
tangential fibre functors.
\end{statement}
\begin{proof}
Choose compatible roots $t^{1/n}$ and use the compositum of the
unramified extension with these roots.  Since the residue characteristic
is zero, every finite extension becomes a Kummer extension of this
form after an unramified extension.  Acting on the coefficients and
fixing every $t^{1/n}$ gives the splitting.  Hensel's lemma gives unique
compatible roots in $1+\mathfrak m_R$ of $1+h$, so the displayed change
of uniformizer gives an isomorphic splitting.

Fix the inertia convention
$\zeta(t^{1/n})=\zeta_n t^{1/n}$ and
$\kappa(a)(\sigma)=\sigma(a^{1/n})/a^{1/n}$.
The splitting which fixes $(at)^{1/n}$ differs from the splitting which
fixes $t^{1/n}$ by $-\kappa(a)$.  On the other hand, a tangent vector $v$
with $t(v)=a$ differs from the vector with $t(v)=1$ by $+\kappa(a)$.
This explains the inverse transformation between a basis of
$\mathfrak m/\mathfrak m^2$ and its dual tangent vector; it also fixes
the sign rather than identifying these two torsors silently.

For an abelian kernel, differences of splittings are precisely continuous
$1$-cocycles and inertia conjugation changes a cocycle by a coboundary.
The Kummer sequence and Hilbert 90 identify this $H^1$ with the stated
inverse limit.  A finite cover near $y$, after an unramified extension,
has a parameter $u$ with $t=u^e$ times a unit.  Its fibre over a tangent
vector is the set of roots of the leading coefficient.  The same roots
are obtained by the preceding compatible-root construction.  This is
an isomorphism on every object and morphism in the finite-\'etale
category and is compatible with products.  It identifies the resulting
fibre functor with the specialization to the punctured tangent line
in Deligne \cite[\S15.13--15.26]{B_Deligne}.
Stix's local packet map \cite[Proposition 8 and Theorem 14]{B_StixCuspidal}
is induced by the same inclusion $G_L\to\Pi_U$; thus it is the packet
in the 1983 letter.  Over finitely generated fields of characteristic
zero that map is injective on splitting classes.
\end{proof}

\begin{statement}{B.4}{Proved}{Ramification and the two differential maps}
\label{st:B.4}

At rational boundary points of a finite morphism of curves, choose
parameters with
$t=a u^d(1+h)$, $a\in k^\times$, $h\in(u)$.
In tangent coordinates the induced map of packets is
\[
 \widehat{k^\times}\longrightarrow\widehat{k^\times},
 \qquad c\longmapsto a c^d.
\]
It restricts to the homogeneous map of tangent vectors
$v\mapsto a v^d$.  In the pure local model $t=a u^d$, the trace and
the determinant norm of differentials satisfy respectively
\[
 \operatorname{Tr}_{k((u))/k((t))}(du/u)=dt/t,
 \qquad
 \operatorname{Nm}(du)=
   (-1)^{d-1}d^{-d}a^{-1}t^{1-d}(dt)^{\otimes d}.
\]
The latter equality is in the norm of the one-dimensional differential
space, identified with $(\Omega^1_{k((t))/k})^{\otimes d}$ by $dt$.
\end{statement}
\begin{proof}
The factor $1+h$ has unique compatible roots in its principal-unit group.
On inertia the map is multiplication by $d$; on the chosen constant
splittings its discrepancy is the Kummer cocycle of $a$.  This proves
the first formula.  Equivalently the map on conormal leading terms is
$t\mapsto a u^{\otimes d}$, whose dual description is the displayed
homogeneous tangent map.

In the pure model, $du/u=d^{-1}dt/t$, and tracing the coefficient
$1/d$ multiplies by the field degree $d$.  The minimal polynomial of
$u$ gives $\operatorname{Nm}(u)=(-1)^{d-1}t/a$.  Since
$du=(ad u^{d-1})^{-1}dt$, taking determinants yields the norm formula.
The identity
$d\log\operatorname{Nm}(u)=\operatorname{Tr}(d\log u)$ follows at once.
For a nontrivial residue-field extension these formulas are applied
after splitting the residue points and then descended by the ordinary
field norm and cohomological corestriction.  Thus neither the trace
map nor the norm line is being substituted for the tangent map.
\end{proof}

\subsection{Every fibre of abelianization on points and packets}

\begin{statement}{B.5}{Proved}{Abelianized fibres in positive genus}
\label{st:B.5}

Let $k/\mathbb Q$ be finitely generated and let $g(X)\geq1$.
Write $r=|S(\bar k)|$.  In the map from rational points and full rational
cuspidal packets to geometrically abelianized section classes:
\begin{enumerate}
\item every rational point has a singleton fibre;
\item images belonging to different rational cusps are disjoint, and
      are disjoint from all rational point images;
\item for $r=1$ the entire unique packet has one abelianized image;
\item for $r\geq2$ the map on every packet is injective.
\end{enumerate}
For $r=0$ the domain consists only of the rational points.
\end{statement}
\begin{proof}
The geometric abelianization is the full Tate module of the generalized
Jacobian $J_S$, with the compact Jacobian $J$ as quotient.  After a
finite extension choose a base point.  The image in $H^1(k,TJ)$ of an
interior point or a packet centred at $z$ is the Kummer class of
$[z-b]$.  The Kummer injection of B.1, and the injectivity of
$X\to J$, separate distinct centres.  The latter injectivity follows
directly: a principal divisor $z-z'$ for $z\ne z'$ would define a
degree-one map to $\mathbb P^1$, contradicting $g\geq1$.
An equality before finite extension remains an equality afterwards,
so the reduction has lost no possible collision.

The toric part of $J_S$ has cocharacter module
$\mathbb Z[S(\bar k)]/\mathbb Z(1,\ldots,1)$.
For $r=1$ this module is zero: inertia maps to zero in
$\Delta_U^{\rm ab}$, proving (3).
For a rational cusp $y$ and $r\geq2$ its coordinate defines the exact
sequence of torus Tate modules
\[
 0\to\widehat{\mathbb Z}(1)\to TT_S\to TT_{S\setminus\{y\}}\to0.
\]
The two relevant degree-zero groups, $H^0(k,TT_{S\setminus\{y\}})$ and
$H^0(k,TJ)$, vanish.  Indeed they are compatible systems of rational
torsion, and the rational torsion groups of a torus or abelian variety
over a finitely generated characteristic-zero field are finite.
The two cohomology exact sequences therefore make
$H^1(k,\widehat{\mathbb Z}(1))\to H^1(k,TJ_S)$ injective.
This is exactly the difference map inside one packet.
\end{proof}

\begin{statement}{B.6}{Proved}{Genus-zero abelian fibres, including the exceptional collisions}
\label{st:B.6}

Let $k/\mathbb Q$ be finitely generated, $g(X)=0$, and put
$r=|S(\bar k)|$.  An anisotropic conic has neither rational interior
points nor rational cusps, so the domain in question is empty.
Otherwise choose a rational interior point as $\infty$ and write
$S\subset\mathbb A^1(\bar k)$.  In the coordinate torsor for
$H^1(k,\Delta_U^{\rm ab})$, the point and packet maps are represented by
\[
 v(x)_s=x-s,\qquad
 v(y,a)_s=\begin{cases}a&s=y,\\y-s&s\ne y,\end{cases}
 \quad a\in\widehat{k^\times},\ y\in S(k).
\]
Two such vectors represent the same class precisely when their
quotient is the restriction of a single diagonal element of
$\widehat{k^\times}$.
For $r\geq4$ every nonempty fibre on this domain is a singleton.  For $r=3$ the
only nonsingleton fibres consist of one element in each of two distinct
rational cusp packets.  For three rational cusps $y,z,w$, their unique
pairwise intersection is
\[
 a=\frac{y-w}{z-w}(z-y),\qquad
 b=\frac{y-z}{(y-w)/(z-w)}.
\]
The three pairwise intersections are distinct.  For $r=2$ with split
boundary, every abelian class has one element in each cusp packet and
additionally one interior point exactly when the class is in
$\kappa(k^\times)$.  For $r\leq1$ the geometric group is trivial and
all elements of the domain have the same image.
\end{statement}
\begin{proof}
Let $A=\Gamma(S,\mathcal O_S)$ and
$T_S=(\operatorname{Res}_{A/k}\mathbb G_m)/\mathbb G_m$.
For a character $(n_s)$ with $\sum_s n_s=0$, the generalized Abel map
evaluates to $\prod_s(x-s)^{n_s}$.  Specialization at $y$ replaces the
vanishing factor $x-y$ by the tangent parameter and the other factors
by $y-s$.  This proves the vector formulas, including their
$G_k$-equivariance for nonsplit $S$.
The Tate-module sequence of
$1\to\mathbb G_m\to\operatorname{Res}_{A/k}\mathbb G_m\to T_S\to1$
and Shapiro's lemma show that equality of these vector classes is
exactly equality modulo a diagonal Kummer class.  There is no extra
kernel here; the preceding term $H^0(k,TT_S)$ is zero.

One may check the equalities after a finite splitting field $L/k$.
The map $\widehat{k^\times}\to\widehat{L^\times}$ is injective:
the inflation--restriction kernel is
$H^1(\operatorname{Gal}(L/k),\widehat{\mathbb Z}(1)^{G_L})$,
and the coefficient group is zero because $L$ has only finitely
many roots of unity.  One first replaces $L$ by its Galois closure.
Two distinct nonvanishing entries of a point--cusp comparison would
require
$(x-s)/(y-s)=(x-s')/(y-s')$, hence $x=y$.  Thus no such collision
exists for $r\geq3$.  Point--point injectivity also follows from B.1.
For two distinct cusps, entries outside them require
$D=(y-s)/(z-s)$ for the same $D$ and every remaining $s$.
Two remaining entries force $y=z$; one remaining entry gives precisely
the displayed $a,b$.  The two intersection parameters within any one
packet differ by $-1$, whose full Kummer class is nonzero by B.1.
There is therefore no triple intersection.  The cases $r=2$ and
$r\leq1$ follow by the same formula from $T_S=\mathbb G_m$ and
$T_S=1$, respectively.  If only some boundary points are rational,
the asserted pairwise intersections occur only for those rational
packets; with $r=3$, two rational cusps force the third to be rational.
The finite rational-coordinate verification is recorded in
\texttt{code/B/packet\_and\_extension\_checks.py}.
\end{proof}

\subsection{Local sections and their limits}

\begin{statement}{B.7}{Proved using the prompt's local theorem}{The complete neighbourhood deduction of (T1)}
\label{st:B.7}

For a finite extension $k/\mathbb Q_p$ and a hyperbolic $U/k$, the map
\[
 U(k)\ \sqcup\!\coprod_{y\in S(k)}\operatorname{Sect}(D_y/G_k)/I_y
 \longrightarrow \operatorname{Sect}(\Pi_U/G_k)/\Delta_U
\]
is bijective.  Each packet is a torsor under $\widehat{k^\times}$.
The external input used for surjectivity is precisely Theorem 1.1 of
the attached October 6, 2026 preprint \cite{B_OpenAIPadic}; its
unrefereed provenance is retained in the source audit.
\end{statement}
\begin{proof}
For a section $s$, its neighbourhoods are the finite \emph{\'etale}
covers corresponding to open subgroups containing $s(G_k)$.
They have a cofinal subsystem whose smooth completions have genus at
least two.  To see this even in initial genus zero or one, first choose
a geometric finite \emph{\'etale} cover whose completion has genus at
least two, then choose a characteristic open subgroup $N\subseteq
\Delta_U$ contained in its subgroup.  The subgroup $N s(G_k)$ gives
the required neighbourhood.  Genus zero with at least three punctures
admits a Fermat-type cover; a once-punctured genus-one curve admits
the Heisenberg cover with nontrivial boundary inertia.  Riemann--Hurwitz
then gives the required genus.  Further finite covers preserve genus
at least two.

The image of the lifted section on each such completion is the section
of a unique $k$-point by the cited theorem.  These points are compatible:
a map of neighbourhoods sends one to a point inducing the other
section, and proper-curve injectivity gives equality.  Their common
image is a point $x\in X(k)$.  When $x\in U(k)$ the compatible lifts
give a rational pro-point in the section neighbourhood tower, and
the stabilizer of that pro-point is the section image.  When
$x\in S(k)$, they give a compatible point in the boundary of the
normalized tower, whose stabilizer is a cuspidal decomposition group.
The section therefore factors through that decomposition group.
This is the exact limit argument of
\cite[\S2.1.2]{B_StixLocal}.  The disjointness and injection of packets
are \cite[Theorem 14]{B_StixCuspidal}; their local Kummer description
is B.3.  Point injection follows from the local Jacobian Kummer map.
\end{proof}

\begin{statement}{B.8}{Proved}{The $p$-adic tangential subpacket and all its point limits}
\label{st:B.8}

For $k/\mathbb Q_p$ finite, a uniformizer $\varpi$ gives
\[
 k^\times=\varpi^{\mathbb Z}\times\mathcal O_k^\times,
 \qquad \widehat{k^\times}=\varpi^{\widehat{\mathbb Z}}
                         \times\mathcal O_k^\times.
\]
The Kummer map is injective, has cokernel
$\widehat{\mathbb Z}/\mathbb Z$, and has dense image.
At a rational cusp, every packet element is a limit of point sections
of points tending to that cusp in the $p$-adic topology.  Every
convergent limit of such point sections belongs to that packet.
\end{statement}
\begin{proof}
The unit group is a finite group times a finite-rank $\mathbb Z_p$
module; its power completion is itself.  The displayed formulas and
the injection, cokernel, and density follow.
Write a desired packet parameter as $(\alpha,u)$ with
$\alpha\in\widehat{\mathbb Z}$ and $u\in\mathcal O_k^\times$.
Choose increasing positive integers $n_j\to+\infty$ with
$n_j\equiv\alpha\pmod{j!}$, and put $t_j=\varpi^{n_j}u$.
A local \emph{\'etale} coordinate at the smooth rational cusp identifies
all sufficiently small $t_j$ with points $x_j\in U(k)$.
For any fixed finite quotient of the fundamental group, normalize
its cover at the cusp.  The local maps are Kummer maps times units,
and Hensel's lemma makes the unit factors have the required roots
on a sufficiently small disc.  Thus the point section modulo this
quotient is determined by the corresponding finite Kummer parameter.
The congruences on $n_j$ make it eventually the specified packet
parameter.  This proves convergence in every finite quotient.
Conversely the same description puts every sufficiently near point
in the image of the packet at every fixed finite level.  The packet
is compact and its map to the space of sections is continuous;
the compatible limit belongs to its image.  This proves the stated
equality of accumulation sets rather than only density of tangents.
\end{proof}

\begin{statement}{B.9}{Proved}{Genus zero and genus one over a $p$-adic field}
\label{st:B.9}

For a smooth proper genus-zero curve over $k/\mathbb Q_p$ there is
exactly one arithmetic section class, including for an anisotropic
conic.  For an elliptic curve $E/k$ with its origin, the section set is
$H^1(k,TE)$ and there is an exact sequence
\[
 0\longrightarrow E(k)\longrightarrow H^1(k,TE)
 \longrightarrow T H^1(k,E)\longrightarrow0.
\]
Local Tate duality identifies the last group with
$\mathbb Z_p^{[k:\mathbb Q_p]}$ noncanonically.
The point-section image is closed; its closure is exactly $E(k)$.
For a genus-one torsor $C$ under $E$, the section set is empty exactly
when its extension class $\delta([C])\in H^2(k,TE)$ is nonzero.
This class vanishes exactly for the maximal divisible subgroup of
$H^1(k,E)$.  When it vanishes, the section set is a torsor under
$H^1(k,TE)$.
\end{statement}
\begin{proof}
A genus-zero curve has trivial geometric fundamental group.
An explicit pointless conic is the projective norm equation
$N_{k'/k}(x)=\varpi z^2$ for an unramified quadratic $k'/k$:
the valuation of a nonzero norm is even, whereas the right side has
odd valuation.  Thus the unique section need not be a point section.

For $E$, the geometric fundamental group is $TE$ and the origin
splits the extension.  Continuous cocycles classify all splittings.
The inverse Kummer sequences give the exact sequence above;
the finite cohomology groups at every torsion level eliminate the
possible inverse-limit defect.  Since $E(k)$ is compact, its own
power completion is $E(k)$.  Local Tate duality makes $H^1(k,E)$
the Pontryagin dual of the compact group $E^\vee(k)$; this has a
finite group and a free $\mathbb Z_p$-module of rank $[k:\mathbb Q_p]$.
Taking Tate modules proves the formula.  The continuous point map
from the compact group $E(k)$ to the Hausdorff inverse-limit
cohomology group has closed image.

For a torsor, the finite-level extension classes are the connecting
classes of $[C]$ for multiplication by $n$ on $E$, as is seen by
identifying the automorphism group of $[n]:E\to C$ with the pulled-back
extension; this is also the cocycle-level calculation in
\cite[Proposition 2.2]{B_HarariSzamuely}.  All vanish exactly
when $[C]$ is divisible by all $n$.  The local-duality decomposition
just given identifies these elements with the maximal divisible
subgroup and supplies compatible divisions.  Finite-level
Mittag--Leffler then gives $\delta([C])=0$ in continuous degree two.
An abelian extension splits exactly when this class vanishes;
choosing a splitting identifies its section set with $H^1(k,TE)$.
Nonzero classes in the divisible subgroup give pointed-free genus-one
torsors with sections.  None is being used as a hyperbolic example.
\end{proof}

\subsection{Fixed groups, partial splittings, and the six axioms}

For reference, the actual axioms in scans 412--421 have the following
typed content.  This transcription specifies the quantifiers before
the geometric interpretation used below.
\begin{enumerate}
\item Admissibility of an outer action survives restriction to an open
subgroup of $\Gamma$.  The converse appears as a question in the margin.
\item Filling a $\Gamma$-stable set of holes preserves admissibility,
whenever the resulting type is still hyperbolic.
\item An admissible effective action has $\pi^\Gamma=1$.
\item An admissible outer action has an admissible effective lift
after restriction to some open subgroup.
\item For a fixed hole $i$, filling it maps admissible point classes
bijectively onto the admissible point classes of the filled group
other than the canonical class $P_i$.  Clause (a) excludes $P_i$,
clause (b) is injection, and clause (c) supplies the remaining classes.
The literal last sentence of scan 418 does not repeat ``other than
$P_i$''; without that exception it contradicts the displayed bijection
and clause (a).  The display is the formulation tested here.
\item Given a partial admissible action over an open subgroup,
and an action over the full group on the filled object preserving
the finite set of removed point germs, there is a unique admissible
extension of the action on the punctured object.  Scan 421's
normalizer calculation identifies its kernel with $L^{\Gamma'}$.
\end{enumerate}
The provisional definition of scans 431--432 imposes: (a) open image
in the arithmetic quotient $\boldsymbol\Gamma_a$; (b) no normalized
remaining inertia for the effective actions obtained by filling holes;
and (c) infinitely many conjugacy classes of section germs normalizing
no inertia.  Its asserted conclusions in scans 432--433 are:
(a) trivial fixed group for an effective admissible lift, with a
separate stronger question for $\ker\chi$; (b) the distinguished
effective action after filling one hole normalizes no remaining
inertia; (c) an arbitrary effective admissible action has the same
property after filling; and (d) the deletion map is injective with
complement exactly its distinguished point.  Scans 434--435 withdraw
the last surjectivity assertion for parameter spaces because distinct
scheme sections can meet.  B.16 realizes that issue by the actual
diagonal and constant graphs and separately refutes unrestricted
group-section surjectivity.

\begin{statement}{B.10}{Proved}{Every full-section action has trivial fixed group}
\label{st:B.10}

Let $k$ be finitely generated over $\mathbb Q$, or finite over
$\mathbb Q_p$.  For every section $s$ of a hyperbolic curve and every
open subgroup $G'\subseteq G_k$,
\[
 C_{\Delta_U}(s(G'))=1.
\]
In particular this condition, appearing as axiom (3) and in scan 431,
holds for point sections and for all cuspidal sections.
\end{statement}
\begin{proof}
Replace $k$ by the finite extension belonging to $G'$.  Suppose
$z\ne1$ is fixed.  Choose a characteristic open subgroup
$N\subseteq\Delta_U$ such that $z\notin N$.  The subgroup
$H=N\overline{\langle z\rangle}$ is normalized by $s(G_k)$.
It gives a section neighbourhood with geometric group $H$.
The image of $z$ in $H^{\rm ab}$ is nonzero, since its image in the
cyclic quotient $H/N$ is nonzero, and is fixed by $G_k$.
But $H^{\rm ab}$ is the Tate module of the generalized Jacobian of
that neighbourhood, whose rational torsion is finite in both field
classes under consideration.  Its invariant Tate module is zero.
This contradiction proves the assertion.  The use of
$N\overline{\langle z\rangle}$ is essential: a fixed element need not
already survive in the abelianization of $\Delta_U$ itself.
\end{proof}

\begin{statement}{B.11}{Proved}{The exact classification of effective lifts}
\label{st:B.11}

For a centre-free profinite geometric group $\Delta$, a continuous
outer action $\rho:G\to\operatorname{Out}(\Delta)$ has the canonical
extension
\[
 E_\rho=\operatorname{Aut}(\Delta)
             \times_{\operatorname{Out}(\Delta)}G.
\]
Its effective lifts are its splittings.  After choosing one lift
$\alpha$, their classes are exactly $H^1_{\rm cts}(G,\Delta_\alpha)$.
An outer action arising from a hyperbolic affine curve with a rational
cusp has effective lifts in that cusp packet.  Every such lift
normalizes the corresponding inertia group, including the
non-tangential elements of the packet.
\end{statement}
\begin{proof}
The kernel of $\operatorname{Aut}(\Delta)\to\operatorname{Out}(\Delta)$
is identified with $\Delta$ by the injective inner-action map.
For a second lift there is a unique $c_g\in\Delta$ with
$\alpha'_g=\operatorname{Inn}(c_g)\alpha_g$.  Multiplication of
automorphisms gives
$c_{gh}=c_g\alpha_g(c_h)$; conjugating the lift changes $c$ by a
coboundary.  These operations are continuous and are inverse to one
another.  For a cusp, the image is contained in $D_y=N_{\Pi_U}(I_y)$
by construction, independent of its Kummer parameter.  Thus the
whole packet consists of the pathological effective lifts excluded
by the no-normalized-loop condition, although their outer action is
the same geometric action as that of every rational-point lift.
Moreover $G_k^\circ=\ker\chi$ centralizes $I_y$, whereas B.10 proves
that the fixed group for an open subgroup of $G_k$ is trivial.
The cyclotomic-kernel statement for point sections is the
centralizer theorem proved in Part F.
\end{proof}

\begin{statement}{B.12}{Refuted}{Vanishing fixed kernel proves uniqueness, not existence}
\label{st:B.12}

For a partial splitting over an open normal subgroup $H\triangleleft G$
of an extension $1\to L\to E\to G\to1$, the equality $L^H=1$
does not imply that the partial splitting extends.  The exact
additional condition is that $N_E(s(H))\to G$ be surjective.
When it is surjective and $L^H=1$, its inverse is the unique extension.
\end{statement}
\begin{proof}
The kernel of this normalizer map is exactly $L^H$.
Surjectivity therefore makes it a continuous isomorphism of profinite
groups.  Every extending section lies in the normalizer, proving
uniqueness as well.

For a finite counterexample take
$L=S_3$, $H=S_3\times S_3$, and
$G=H\rtimes C_2$, where $C_2$ interchanges the factors.
Set $E=L\times G$ and $s(a,b)=(a,(a,b))$ over $H$.
The fixed subgroup in $L$ is $Z(S_3)=1$.
An extension would give a homomorphism $G\to S_3$ restricting to the
first projection on $H$.  The kernel of that restriction is
$1\times S_3$, which is not stable under interchange, a contradiction.
In fact the normalizer image is exactly $H$.
The supplied program enumerates all $432$ elements of $E$, verifies
the orders and the centralizer, and computes this normalizer image.
This refutes the implication proposed in the prompt from the
one fixed-kernel condition alone. The actual text of scan 421
already says that this condition implies uniqueness and that
existence means surjectivity of the normalizer projection.
The calculation therefore agrees with that distinction in
the manuscript. It does not delete the other geometric
hypotheses of axiom (6).
\end{proof}

\begin{statement}{B.13}{Proved}{The geometric meanings of axioms (1)--(6)}
\label{st:B.13}

For actions arising from pointed curves over the fields in B.10,
with admissible points defined as the actual interior point sections,
axioms (1)--(4) of scans 412--415 hold.  The three clauses (a)--(c) of
axiom (5) are exactly exclusion of the deleted point, injectivity,
and lifting of every other point.  They hold in the geometric
subcategory.  Axiom (6) holds there by effective descent of the
curve and its finite set of points, whenever the filled curve remains
hyperbolic.  The fixed subgroup of its actual forgetting kernel is
trivial, including when the removed points are unordered.  For
intrinsic admissibility defined only by absence of a
normalized cusp, axiom (5) is the additional point-surjectivity
assertion, after the already proved point and packet injections.
\end{statement}
\begin{proof}
Restriction to an open Galois subgroup is finite base change, giving
(1).  Filling a Galois-stable set of punctures is the open immersion
of curves into their partial completion, giving (2).  B.10 gives
(3).  A nonempty finite-type curve has a closed point with finite
residue extension, giving (4).
For a rational hole $y$, the actual point sets satisfy the literal
equality $(U\cup\{y\})(k)\setminus\{y\}=U(k)$; applying the
injective point-section map proves the three clauses of (5).
On all sections the omitted packet at $y$ maps to the point $y$,
so the set equality has exactly the meaning claimed only after the
admissibility condition is specified.

A finite set of distinct geometric points stable under descent data
descends as a finite \emph{\'etale} subscheme; deleting it commutes
with descent.  This gives existence in the geometric form of (6).
Here is a proof of the fixed-kernel assertion that also retains the
finite permutation quotient for an unordered removed set.  Write
$U\subset V$ for the punctured and filled curves, both hyperbolic,
and let $\ell$ in the kernel of the outer forgetting map be fixed
by the open subgroup $G_{k'}$.  Choose an automorphism $\beta$ of
$\Delta_U$ representing $\ell$.  The equality of outer classes
$[\beta]\rho(g)[\beta]^{-1}=\rho(g)$ gives the actual arithmetic
automorphism
\[
 (a,g)\longmapsto(\beta a\beta^{-1},g)
 \quad\hbox{of}\quad
 \operatorname{Aut}(\Delta_U)
       \times_{\operatorname{Out}(\Delta_U)}G_{k'}.
\]
B.11 identifies this pullback with $\Pi_{U_{k'}}$.
The relative Isom theorem, in the precise form of
Proposition~\ref{A:relative-hom}, realizes its outer class by a
unique $k'$-automorphism $f$ of $U$.  The preserved set of removed
cusps lets $f$ extend to $V$.  Its geometric outer class on $V$ is
the identity.  After an inner conjugation its arithmetic action
is the identity on $\Delta_V$; then its value at every element of
$\Pi_V$, divided by that element, lies in $\Delta_V$, centralizes it, and
is therefore the identity.  The injective Isom map for the
hyperbolic curve $V$ forces the extended morphism to be the
identity.  Restricting to $U$ gives $\ell=1$.
Thus the fixed kernel vanishes without replacing the unordered
forgetting kernel by its ordered subgroup.  The hyperbolicity of
$V$ is used in this Isom step.
The normalizer surjectivity furnished by descent is the extra step
which B.12 shows cannot be omitted.

For intrinsically selected sections, replacing point sections in this
argument by all sections which normalize no inertia group is exactly
the missing surjectivity assertion.  The cyclotomic centralizer
theorem identifies that set with the sections of zero
$G_k^\circ$-centralizer; it does not supply a pro-point.  B.7 supplies
the pro-point over $p$-adic local fields by its neighbourhood argument.
No corresponding unrestricted number-field pro-point construction
has been obtained here.
\end{proof}

\subsection{A global descent step beyond LM3 E.17}

\begin{statement}{B.14}{Proved}{Finite support at finite places forces a point or a cusp}
\label{st:B.14}

Let $U/k$ be hyperbolic over a number field and let $s$ be a Selmer
section.  Let $Z\subset X$ be finite over $k$.  Suppose the
compactification point of $s_v$ belongs to $Z(k_v)$ for a set of
finite places of density one.  Then $s$ is induced by a point of
$Z(k)\cap U(k)$ or is cuspidal at a point of $Z(k)\cap S(k)$.
In particular, a section whose finite localizations are all
cuspidal at one fixed rational cusp is globally cuspidal at that cusp.
\end{statement}
\begin{proof}
For a neighbourhood $V\to U$ with completion $Y$, let $Z_V$ be the
finite inverse image of $Z$.  The lifted section is again Selmer:
a lifted local point or cusp in the universal neighbourhood tower
projects to a local point or cusp in every intermediate cover.
The local compactification points lie in $Z_V$ at the prescribed
density-one set.  Its proper section gives an adelic point surviving
finite descent (pull back each finite torsor and use the section to
choose its arithmetic twist).  Passing cofinally to $Y$ of positive
genus, Stoll's Theorem 8.2 \cite{B_Stoll} identifies that adelic point
with a rational point of $Z_V$.  Thus $Z_V(k)$ is nonempty for every
neighbourhood.  These are finite sets.  Every finite collection of
compatibility conditions is realized by a point in a common
dominating neighbourhood, so compactness gives a compatible system
in $\varprojlim_V Z_V(k)$.

Its stabilizer in $\Pi_U$ contains the section image.  Over an
interior point this stabilizer is a point decomposition group; at
the boundary it is the cusp decomposition group.  The projection
to $G_k$ identifies the section as the respective point or
cuspidal splitting.  This is the argument in
\cite[\S3, proof of Corollary 5]{B_StixLocal}, with its final
noncuspidality assumption removed and the resulting cusp alternative
retained.  Only density-one finite places are used; no real
connected component has been replaced by a chosen point in the proof.
\end{proof}

\begin{statement}{B.15}{Proved using the prompt's local theorem}{Affine modular cuspidality}
\label{st:B.15}

For $N$ with $g(X_1(N))\geq2$, every section of
$\pi_1(Y_1(N))\to G_{\mathbb Q}$ is cuspidal on the affine curve.
Its centre is the unique rational cusp determined on $X_1(N)$.
\end{statement}
\begin{proof}
The compact modular conclusion is
\cite[Corollary 8.5]{B_OpenAIPadic}, combined with Mazur's torsion
theorem: $X_1(N)(\mathbb Q)$ consists of cusps at these levels.
This is the precise conclusion checked in LM3 D.10 and E.17.
Let $c$ be the resulting cusp.  At every finite place, B.7 makes
the affine localization a point section or a cusp section.
Its image on the compact curve is $s_c$; proper point injection
forces its centre to be $c$.  Thus it is cuspidal at $c$ at every
finite place.  B.14 now proves global cuspidality.
This last limit argument is the additional step absent from the
earlier return; equality on compactification alone was not used to
infer it.
\end{proof}

\subsection{Families and the obstruction to lifting a tautological section}

\begin{statement}{B.16}{Refuted}{Unrestricted section surjectivity for a constant family}
\label{st:B.16}

For $U_0=\mathbb P^1\setminus\{0,1,\infty\}$ over a number field $k$
and a smooth geometrically connected base $S/k$ with a rational
point, not every section of
$\pi_1(U_0\times S)\to\pi_1(S)$ comes from a morphism $S\to U_0$.
For geometric sections, deletion of a section retains precisely the
sections disjoint from it at every geometric point of the base.
\end{statement}
\begin{proof}
The product extension is
$\Pi_{U_0}\times_{G_k}\Pi_S\to\Pi_S$.
For a tangential cusp section $s_0:G_k\to\Pi_{U_0}$ at $0$, the map
$\gamma\mapsto(s_0(p_S(\gamma)),\gamma)$ is a splitting.
Its restriction at any $b\in S(k)$ is the cuspidal section $s_0$.
A graph $f:S\to U_0$ would restrict to the point section of $f(b)$,
contradicting the point--packet disjointness theorem.  This gives
an actual group-theoretic section and a complete counterexample.

For two morphism sections $f,g$, the graph of $f$ lies in the complement
of the graph of $g$ exactly when $f(b)\ne g(b)$ for every geometric
point $b$.  For $S=U_0$, the diagonal and the constant section $2$
meet at $2$ and are distinct, proving the correction to scans
433--435.  For the corresponding group-theoretic test, pull both
sections back along each closed-point section of $S$ and use
point-section injection over its finite residue field extension.
This tests precisely equality at that point.  The assertion is
about the indicated scheme sections; no unsupported lifting of a
group splitting through deletion is inferred.
\end{proof}

\begin{statement}{B.17}{Proved}{A tautological complete section cannot lift through its deletion}
\label{st:B.17}

Let $g\geq2$, $n\geq1$, and $k$ have characteristic zero.  For the
universal curve over $\mathcal M_{g,n/k}$, none of the $n$ tautological
complete sections admits a lift to a section of the fundamental-group
extension after all $n$ markings are removed.  This remains true
after every finite \emph{\'etale} base change of degree $d$.
The obstruction can be detected by a finite $\mu_m$-extension,
for a prime $m\nmid d(2g-2)$, with $d=1$ on the original base.
\end{statement}
\begin{proof}
Fix the proposed tautological marking, indexed by $1$.  Take a smooth
proper genus-$g$ curve $C/k$, put
$K=k(C^{n-1})$, and let $p_2,\ldots,p_n\in C(K)$ be its independent
generic points.  For $n=1$ put $K=k$ and omit this list.  After a
finite constant extension we may use a rational base point of $C$
to identify its Picard components with its Jacobian $J$.
Let $T=C_K\setminus\{p_2,\ldots,p_n\}$; varying the first marking
gives a morphism $T\to\mathcal M_{g,n/k}$.  The complete pulled-back
family is $C_K\times T$, and the first section is its diagonal.
Any supposed punctured lift, followed by filling the other holes,
would lift the diagonal homomorphism through
\[
 \pi_1((C_K\times T)\setminus\operatorname{diag}(T))
 \longrightarrow\pi_1(C_K\times T).
\]

We prove $c_1(T_C|_T)\ne0$ modulo $m$.  A contrary relation in
$\operatorname{Pic}(T)/m$ would give, on $C_K$,
\[
 T_C=mL+\sum_{i=2}^n a_i[p_i]
 \tag{B.1}\label{eq:B-Picrelation}
\]
in additive Picard notation.  Varying the generic tuple gives a
rational map $C^{n-1}\dashrightarrow\operatorname{Pic}^{\deg L}(C)$.
After translation it is a rational map to $J$, and extends to a
morphism: rational maps from a smooth variety to an abelian variety
extend across codimension one by properness and across the remaining
indeterminacy by the absence of rational curves in an abelian variety.
Its induced homomorphism on the $i$th Jacobian factor is an integral
endomorphism $\phi_i\in\operatorname{End}(J)$.
Equation \eqref{eq:B-Picrelation} gives
$m\phi_i=-a_i\operatorname{id}_J$.  Restriction to $J[m]$ shows
$m\mid a_i$.  Degrees in \eqref{eq:B-Picrelation} then imply
$m\mid(2g-2)$, a contradiction.  For $n=1$ this is simply the degree
argument.  Kummer's exact sequence injects
$\operatorname{Pic}(T)/m$ into $H^2_{\rm et}(T,\mu_m)$, proving the
claimed nonzero Chern class.

Let $\mathcal G$ be the $\mu_m$-gerbe of $m$th roots of
$\mathcal O(\operatorname{diag})$ on $C_K\times T$.
Its class is the Kummer Chern class.  On the complement of the
diagonal the canonical nonvanishing divisor section trivializes
this line, hence gives an actual morphism to $\mathcal G$.
The base $C_K\times T$ is an \emph{\'etale} $K(\pi,1)$ for finite
coefficients: its geometric factors are characteristic-zero
hyperbolic curves, their discrete surface/free groups are good, and
the geometric and arithmetic Hochschild--Serre sequences preserve
this comparison.  Thus the gerbe gives the genuine continuous
extension
\[
 1\to\mu_m\to\pi_1(\mathcal G)\to
                \pi_1(C_K\times T)\to1
\]
with the cyclotomic kernel action.  The asserted lift of the diagonal
would split its pullback.  That pullback class is
$c_1(\mathcal O(\operatorname{diag})|_{\operatorname{diag}})
=c_1(T_C|_T)$, just proved nonzero.  This contradiction proves the
assertion on the original base.

For a finite \emph{\'etale} moduli cover of degree $d$, pull it back
to the test curve $T$.  The resulting $f:T'\to T$ is finite
\emph{\'etale} of degree $d$, possibly disconnected.  On finite
coefficient cohomology its trace satisfies
$\operatorname{Tr}_f f^*=d$.  Choose $m\nmid d(2g-2)$.
The nonzero class $c_1(T_C|_T)$ therefore remains nonzero on $T'$;
otherwise its trace would give $d\,c_1(T_C|_T)=0$.
A lifted tautological section on the moduli cover would split the
gerbe pullback on every component of $T'$, contradicting this
nonvanishing.  The argument exhibits the extension with finite
kernel and its pullback; it does not require goodness of a mapping
class group.
\end{proof}

\begin{statement}{B.18}{Proved using the specified universal-curve input}{No sections of the universal punctured family in genus at least three}
\label{st:B.18}

Over $k/\mathbb Q$ finitely generated, for $g\geq3$ and $n\geq1$,
the universal punctured curve over $\mathcal M_{g,n/k}$ has no
fundamental-group section.  The external input is precisely the
complete-family section classification in
\cite[Main theorem]{B_LuoWatanabe}; its status and the inspected
proof step are recorded in the source audit.  The same conclusion
holds on its geometrically connected full symplectic level-$m$
covers over fields containing $\mu_m$.
\end{statement}
\begin{proof}
Fill the $n$ marked holes.  The cited classification makes the image
of any putative section one of the $n$ tautological complete classes.
Choose a representative and conjugate so that it is that tautological
section.  B.17 proves that this section has no lift through the
deleted family, giving a contradiction.  On a full symplectic
level cover the cited paper's final Remark gives the same complete
classification, and the finite-cover form of B.17 applies.
In particular the assertion
of scan 432 that there are no admissible sections holds here with
the stronger conclusion that there are no sections at all.
Calling a tautological complete section ``cuspidal'' after deleting
it does not by itself supply a splitting of its normal-line
extension; B.17 computes the obstruction to such a splitting.
\end{proof}

\subsection{Moduli criteria, lattices, and the additional archive tasks}

\begin{statement}{B.19}{Proved using the prompt's local theorem}{Local interpretation of admissibility and the stack $\mathcal M_{1,1}$}
\label{st:B.19}

For a hyperbolic curve over a finite extension of $\mathbb Q_p$, the
sections which normalize no cusp inertia subgroup are exactly its
point sections.  The same local point-or-cusp description, with its
full automorphism groupoids, holds for a hyperbolic DM curve, including
$\mathcal M_{1,1}$ and the $S_3$-quotient of the tripod.
\end{statement}
\begin{proof}
B.7 gives the two alternatives.  Every cusp alternative normalizes
its inertia by B.11, and point sections are disjoint from the
cuspidal packets.  Stix's normalizer characterization identifies a
section normalizing inertia with the corresponding packet.
For a DM curve choose a geometric finite \emph{\'etale} curve cover,
and a characteristic open subgroup $N$ of its geometric fundamental
group contained in the subgroup of that cover.  Finite generation
makes such $N$ available.  For the given section $s$, the subgroup
$Ns(G_k)$ defines a geometrically connected curve neighbourhood
$V\to\mathcal U$ over $k$ itself: its geometric group is $N$,
and its geometric cover is a scheme.  Descent of the finite
\emph{\'etale} covering retains that property.  B.7 applies to the
lifted section on $V$.  Pushing its point or its cusp fibre functor
forward supplies the required stack point or cusp class.
Full faithfulness is tested on these curve covers over finite
extensions: point injection and trivial point centralizers give it
on interior points; uniqueness of the boundary pro-point gives it
on cusp fibre functors and separates the two sorts of centres.
These are exactly the curve-neighbourhood lemma and four steps in
the proof of
\cite[Theorems B and C]{B_Steklov}, where Theorem C is the
full-faithfulness input over sub-$p$-adic fields.
Automorphisms, including the generic $\mu_2$ of
$\mathcal M_{1,1}$, are retained in this descent.  Passing just to
the coarse $j$-line would discard them and is not the map used.
\end{proof}

\begin{statement}{B.20}{Proved}{The lattice stabilizer for the tripod}
\label{st:B.20}

For the tripod over a number field $k$ with a rational base point,
fix an embedding into $\mathbb C$ and the resulting discrete
topological lattice in the geometric fundamental group.  Its
stabilizer in $G_k$ is trivial when the chosen embedding of $k$ is
nonreal, and is generated by its complex conjugation when that
embedding is real.  The oriented lattice has trivial stabilizer
in both cases.  Thus the canonical oriented discretization claimed
in scans 489--490 already fails for the tripod.
\end{statement}
\begin{proof}
Riemann existence maps the based discrete fundamental group injectively
and densely into its profinite completion.  A Galois automorphism
preserving this subgroup restricts to a discrete automorphism.
Its action on the discrete abelianization is the scalar $\chi$;
an integral invertible scalar has to be $+1$ or $-1$.  Each discrete
peripheral generator is therefore profinitely conjugate to the
corresponding discrete generator or its inverse.  Conjugacy
separability of a free group makes this an ordinary discrete
conjugacy.  Thus it preserves the three peripheral cyclic subgroups.
In the free group on two generators,
the pure peripheral outer automorphism group consists of the
identity and the orientation-reversing reflection: this is the
Dehn--Nielsen--Baer description of the extended mapping class group
of the three-punctured sphere.  Consequently its outer action is
either the identity or the outer action of complex conjugation.
Faithfulness of the tripod Galois action identifies the Galois
element itself in either case.  Complex conjugation belongs to
$G_k$ exactly for the chosen real embedding and preserves the based
lattice then.  It reverses its orientation, so cannot stabilize the
oriented lattice.  Changing the embedding by a Galois element
transports the comparison map and hence the lattice by that element;
the orbit is therefore nontrivial, as asserted.
\end{proof}

\begin{statement}{B.21}{Proved}{What the $\overline{\mathbb Q}$-algebraic entry supplies}
\label{st:B.21}

The reading of cote 148, scan 43 as obtaining ``sections Gal$/\mathbb Q$''
has the following verified pointwise content.  A finite-type
$\overline{\mathbb Q}$-curve with an algebraic point has, after a
finite field of definition, a point-induced section; its intrinsic
datum without that field is a germ over open subgroups of
$G_{\mathbb Q}$, modulo geometric conjugacy.  This datum determines
the point under the semi-abelian hypothesis of cote 143, scans 4--5.
The stronger reading asserting a $\mathbb Q$-model for every such
curve is false.

For a hyperbolic curve $U/\overline{\mathbb Q}$, the
constants condition in cote 143, scan 2 gives the exact sequence
\[
 1\longrightarrow\operatorname{Aut}_{\overline{\mathbb Q}}(U)
 \longrightarrow\operatorname{Aut}_{\rm abs}(U)
 \xrightarrow{\;\theta\;}\Gamma_U\longrightarrow1,
 \qquad
 \theta(f)=(f^*|_{\overline{\mathbb Q}})^{-1}.
\]
Its image $\Gamma_U$ is open in $G_{\mathbb Q}$, and need not
equal $G_{\mathbb Q}$. A finite field of definition gives a
continuous partial splitting. Conversely a continuous splitting
over $G_K$ gives an effective $K$-form of the curve, with its
boundary; conjugating by
$\operatorname{Aut}_{\overline{\mathbb Q}}(U)$ gives precisely
isomorphic forms.
\end{statement}
\begin{proof}
The equations of the curve, the chosen point, and finitely many
descent isomorphisms involve finitely many algebraic numbers, hence
are defined over one number field $L$.  The structure map of the
point gives a section over $G_L$.  Enlarging $L$ restricts the
section, which is exactly the germ convention in cote 143, scan 5,
and in the groupoid calculation of LM scans 61--65.
B.1 proves its injectivity on points.
For the additional exact sequence, the elements of
$\Gamma(U,\mathcal O_U)$ algebraic over $\mathbb Q$ are exactly
$\overline{\mathbb Q}$: a polynomial over $\mathbb Q$ splits
over $\overline{\mathbb Q}$, and the coordinate ring of the
integral curve has no zero divisors. An absolute automorphism
therefore preserves this constants subring. The inverse in
$\theta$ compensates for the contravariance of pullback, so
it is a homomorphism, with the asserted kernel and image.
A model over a number field $L$ gives the semilinear
base-change automorphisms with scalar pullback $\sigma^{-1}$,
for $\sigma\in G_L$. These form a section of $\theta$ over
$G_L$, proving that its image is open.

The geometric automorphism group of a hyperbolic curve is
finite. Enlarge $L$ so that the model and all its automorphisms
are defined over $L$. The inverse image of $G_L$ is then
$\operatorname{Aut}_{\overline{\mathbb Q}}(U)\times G_L$;
this specifies the natural profinite topology on the
absolute automorphism group. Two choices give the same topology:
their models and identifications are defined over one further
finite extension, on whose Galois group the model actions agree.
A continuous splitting over
$G_K$ differs from a model splitting, on a common open
subgroup, by a continuous cocycle with values in this finite
automorphism group. The equality locus of the two section
homomorphisms is an open subgroup. Its normal core in $G_K$
contains an open normal $G_M$, with $M/K$ finite Galois,
on which the cocycle is trivial. Normality makes every element
of the given $G_K$-action descend to an $M$-semilinear
automorphism of $U_M$; modulo $G_M$ these are the actual
$\operatorname{Gal}(M/K)$ descent datum on $U_M/M$.
It is effective: extend the automorphisms to the smooth
proper completion, take the invariant function field under
the finite semilinear group, and form its smooth proper
$K$-curve. Finite Galois descent of vector spaces identifies
the base-changed function field with the original one, hence
identifies the smooth proper models. The stable finite
boundary set descends and can be removed. Conversely every
$K$-form gives this datum, and an isomorphism of forms is
exactly the indicated conjugation. This proves the
effectivity statement in this precise hyperbolic-curve
range without treating an arbitrary abstract outer action
as an already effective descent datum.

For the stronger reading, take a pointed elliptic curve with
$j$-invariant $\sqrt2$.  Its $j$-invariant is preserved by every
isomorphism and is not rational, so it cannot descend to
$\mathbb Q$.  Its field-of-definition subgroup is nevertheless
open.  The doubtful two-word reading on scan 43 is retained as a
reading uncertainty; no conclusion depends on resolving it in
favour of the stronger assertion.
\end{proof}

\begin{statement}{B.22}{Proved}{Kummer structure, crossings, and finite covers in cote 143}
\label{st:B.22}

For a curve, a chosen lift of a boundary point gives an actual
homomorphism $\widehat{\mathbb Z}(1)\to\Delta_U$; changing that lift
conjugates the map.  Arithmetic transport acts through the
cyclotomic character and the permutation of boundary points.
At a strict normal crossing with $r$ branches, the corresponding
local group is $\widehat{\mathbb Z}(1)^r$ with its branch labels;
the actual map to the global geometric fundamental group must
also be retained.  It need not be injective.
For a finite \emph{\'etale} cover defined by $H\subseteq\Pi_U$,
the boundary lifts and inertias are
\[
 H\backslash\Pi_U/D_y,
 \qquad H\cap gI_yg^{-1},
\]
with the local ramification indices giving multiplication on each
one-dimensional Kummer lattice.  Every section germ lifts after
restriction to a sufficiently small open subgroup.
\end{statement}
\begin{proof}
The curve statement is the compatible-root construction of B.3
with a chosen point of the normalized universal cover.
At a strict normal crossing the local parameters
$t_1,\ldots,t_r$ may be rooted independently; Abhyankar's lemma
in characteristic zero gives the product local inertia.  Its
global map records which local Kummer coverings extend globally.
For example, remove two coordinate lines from $\mathbb P^2$.
The complement is $\mathbb A^1\times\mathbb G_m$.  At their
intersection use $x=X/Z$, $y=Y/Z$; the global multiplicative
coordinate is $u=y/x$.  Thus the local-to-global inertia map is
\[
 \widehat{\mathbb Z}(1)^2\longrightarrow\widehat{\mathbb Z}(1),
 \qquad (a,b)\longmapsto b-a,
\]
with nonzero diagonal kernel.  This provides the explicit
replacement for the dimension-one reservation in scans 6--8.

The double-coset and intersection formulas are the stabilizer
formulas for the action on the normalized universal cover.
Intersecting a procyclic inertia with an open subgroup gives
$e\widehat{\mathbb Z}(1)$, where $e$ is exactly the ramification
index.  For a section $s:G_L\to\Pi_U$, the inverse image
$s^{-1}(H)$ is open, and restriction lands in $H$; this is the
claimed lift of germs.  These are the actual maps behind
``automatiquement'' on cote 143, scan 12.  At crossings one uses
the corresponding finite-index sublattice and its integer
monomial map, rather than a single ramification integer.
\end{proof}

\begin{statement}{B.23}{Proved using the specified universal-curve input}{The unpointed universal curve in genus at least three}
\label{st:B.23}

For $k/\mathbb Q$ finitely generated and $g\geq3$, the fundamental-group
extension of the universal complete curve over $\mathcal M_{g/k}$ has
no section.  The same argument applies to its full symplectic level
covers included in \cite[Remark on topologically ample families]{B_LuoWatanabe}.
\end{statement}
\begin{proof}
Write $\mathcal C_g\to\mathcal M_g$ for this universal curve;
$\mathcal C_g=\mathcal M_{g,1}$.  Pull a proposed section back along
$\mathcal C_g\to\mathcal M_g$.  The resulting complete curve over
$\mathcal M_{g,1}$ has its diagonal marking.  The $n=1$ case of the
specified classification makes the pulled-back section conjugate
to that diagonal section.  Restrict to the geometric fibre $C$
of $\mathcal M_{g,1}\to\mathcal M_g$ over a fixed geometric curve.
In the product extension
$\Delta_C\times\Delta_C\to\Delta_C$, the section pulled back from
$\mathcal M_g$ is the graph of the trivial homomorphism, because
the fibre group maps to $1$ in $\pi_1(\mathcal M_g)$.
The diagonal is the graph of $\operatorname{id}_{\Delta_C}$.
Conjugation cannot identify a trivial homomorphism with the
identity on the nontrivial group $\Delta_C$.  This is the required
contradiction.  Marking a point does not alter a full symplectic
level structure, so the identical pullback diagram proves the
level assertion.
\end{proof}

\begin{statement}{B.24}{Proved}{Explicit low-genus finite \emph{\'etale} multisections}
\label{st:B.24}

The proposed absence of all admissible section germs for universal
families has the following geometric counterexamples.
\begin{enumerate}
\item Over $\mathcal M_{0,4}=\mathbb P^1_\lambda\setminus\{0,1,\infty\}$,
the universal four-punctured line acquires the interior section
$x=t$ after the degree-two \emph{\'etale} cover $\lambda=t^2$.
\item Over $\mathcal M_{1,1}$, the nonzero $m$-torsion of the universal
elliptic curve, for every $m\geq2$, is a finite \emph{\'etale}
multisection of the curve with its origin deleted.
\item Over $\mathcal M_2$, the six Weierstrass points are a degree-six
finite \emph{\'etale} multisection of the universal complete curve.
\end{enumerate}
\end{statement}
\begin{proof}
In (1) the covering curve is
$\mathbb P^1_t\setminus\{0,1,-1,\infty\}$.
Ramification of $t\mapsto t^2$ occurs only at $0,\infty$, already
removed.  The section value $t$ is different from $0,1,\infty,t^2$
on this curve: the only additional equality $t=t^2$ occurs at
$t=0,1$.  This is an explicit section of the pulled-back
\emph{punctured} family.
In (2), multiplication by $m$ on a characteristic-zero elliptic
scheme is finite \emph{\'etale}.  Removing its zero section from the
kernel leaves a finite \emph{\'etale} cover of the moduli stack of
degree $m^2-1$; the universal nonzero torsion point is the desired
section after pullback.
Every smooth genus-two curve is hyperelliptic.  Its canonical
degree-two map has six distinct simple ramification points in
characteristic zero, by Riemann--Hurwitz.  The canonical involution
and this reduced ramification divisor commute with base change;
the relative divisor is therefore finite \emph{\'etale} of degree
six over the smooth genus-two moduli stack.  Its universal point
gives (3).  Each construction supplies an open subgroup and an
actual algebraic section, hence its fundamental-group splitting.
The finite polynomial equalities used in (1) are checked in the
program supplied with this section.
\end{proof}

\begin{statement}{B.25}{Proved}{The necessary moduli tests and the exact tripod test}
\label{st:B.25}

An actual marked curve over a number field satisfies the distinctness
tests of scans 487--488: in positive genus the sections of its marked
points remain pairwise distinct after all holes are filled; in genus
zero they remain distinct in every three-puncture comparison in which
they are interior points.  Its rational interior point sections
normalize no boundary inertia.  For ordered type $(0,4)$ the moduli
point map and the tripod point map are literally the same map under
\[
 \mathcal M_{0,4}(k)=k\setminus\{0,1\},\qquad
 \mathcal M_{0,4}\simeq\mathbb P^1\setminus\{0,1,\infty\},
\]
where the four markings are $(0,1,\infty,\lambda)$.
The Isom-form reconstructs at most one marked curve from an outer
action which is already realized by a marked curve; it supplies no
point of this moduli curve for an arbitrary section.
\end{statement}
\begin{proof}
The positive-genus assertion is B.1 applied to the compact Jacobian.
For genus zero, retain three distinct marked points, send them by
the unique projective transformation to $0,1,\infty$, and apply B.1
to the resulting tripod.  The non-normalization assertion follows
from the inertia normalizer description and the disjointness of
point sections and cusp packets.  These are the actual quotient
maps involved in the manuscript tests.
The unique projective transformation also proves the displayed
moduli identification, over $k$ and on finite \emph{\'etale} covers.
Taking arithmetic fundamental groups and section classes therefore
identifies the two point maps, rather than merely comparing their
cardinalities.
Finally, two geometric realizations of the same peripheral outer
action give an isomorphism of their arithmetic curve groups over
$G_k$.  The Isom-form, with the marked inertia classes retained,
gives a unique isomorphism of marked curves.  Its input is an
isomorphism between the two realized curve groups; an arbitrary
moduli section contains no second realized curve to which that
theorem can be applied.  Thus the unresolved surjectivity here is
precisely the one recorded in B.27.
\end{proof}

\begin{statement}{B.26}{Proved}{The range of the Hom-form for constant families}
\label{st:B.26}

Let $S$ and $U$ be hyperbolic curves over a number field $k$.
A section of the constant-family extension is a homomorphism
$h:\Pi_S\to\Pi_U$ over $G_k$.  Every such $h$ with open image is
induced by a unique dominant morphism $S\to U$.
The homomorphisms with $h(\Delta_S)=1$ are exactly
$s\circ p_S$ for field sections $s:G_k\to\Pi_U$; among them, the
morphism-induced classes are exactly the constant morphisms at
$k$-points of $U$.  The corresponding graphs are the actual family
sections in these two ranges.
\end{statement}
\begin{proof}
The product fundamental group is the fibre product over $G_k$, so
projecting a splitting to the first factor gives $h$, and its graph
recovers the splitting.  Mochizuki's open Hom theorem applies to
$h$ with exactly its open-image hypothesis and produces the unique
dominant morphism.  Its graph is a section of the product family;
no claim that the graph itself is dominant in the surface is used.
For $h(\Delta_S)=1$, the quotient property of
$\Pi_S\to G_k$ gives a unique factor $s$, and being over $G_k$
makes $s$ a splitting.  A nonconstant morphism of curves is
dominant and induces an open image on geometric fundamental groups.
It therefore cannot induce this trivial geometric homomorphism.
The remaining morphisms are precisely the constant rational-point
maps.  B.16 supplies cuspidal examples among the other field-section
classes.  This proves the stated classification without applying
the dominant Hom theorem to a nondominant graph.
\end{proof}

\begin{statement}{B.27}{Open in this incomplete return}{The unresolved part of Part B}
\label{st:B.27}

The following requested conclusions have not been proved or
refuted here: point-surjectivity for all zero-cyclotomic-centralizer
sections of the tripod over a number field; the same point-or-cusp
dichotomy for every nonisotrivial hyperbolic curve over a finitely
generated positive-transcendence-degree extension of $\mathbb Q_p$;
all intrinsic instances of axioms
(5)--(6) outside the geometric subcategory; section classification
over arbitrary open subgroups of $\pi_1(\mathcal M_{g,n})$ and the
remaining cases in genus at most two; the general
intrinsic realization
criterion for arbitrary admissible outer actions.
\end{statement}
\noindent
The exact unresolved step for the first problem is construction
of an interior rational pro-point in the full section neighbourhood
tower.  B.10 and the cyclotomic-kernel theorem remove automorphisms
and identify the cusp alternative, but neither constructs that
pro-point.  B.14 constructs it for finite support; the unrestricted
tripod section does not acquire finite support from the argument.
The local theorem and B.7 yield local pro-points, whose compatibility
with one global rational pro-point is the missing descent step.
Passing to Jacobian Kummer classes determines a point already
present, but does not put every nonabelian section in the image.
For positive-transcendence-degree $p$-adic fields, B.38--B.40
now prove unramification and identify every proper generic
section with a relative fundamental-group section on a smooth
model. B.41 completes the isotrivial point-or-cusp classification
in every transcendence degree. For a nonisotrivial family the
remaining step is realization of its point-theoretic relative
section by a rational generic point. The factorization and
its fibrewise points alone do not prove that realization.

For arbitrary moduli covers, the inspected Luo--Watanabe statement
and its Remark on topologically ample families cover full symplectic
levels and bases inducing the stated geometric fundamental-group
isomorphism.  They do not quantify over every open subgroup.
Replacing an arbitrary open subgroup by a full level subgroup
would require a containment assertion not supplied by that result.
B.17 continues to exclude lifts of any tautological class, but
does not classify additional complete classes on such a cover.
The lattice question initially stopped at the discrete mapping
class group, but B.31 below completes it using the precise
archimedean intersection theorem.  The other gaps are recorded
as unfinished tasks, not as completed alternatives to the proof
requirements.

\subsection{Two exact repairs to the deletion statements}

\begin{statement}{B.28}{Refuted in the unrestricted formulation}{The collision test for equation (25)}
\label{st:B.28}

The map of scan 409, equation (25), is injective for a smooth
marked curve over a field finitely generated over $\mathbb Q$.
Its assertion for an arbitrary section of the arithmetic moduli
extension, without the distinctness admissibility condition of
scan 487, is false.  Already for $\mathcal M_{0,5}$ two different
marked holes can induce the same section after they are both filled.
\end{statement}
\begin{proof}
For a smooth marked curve the filled holes are distinct rational
points of the partially completed hyperbolic curve.  The point
injection of B.1 gives the assertion in this form, with no section
surjectivity hypothesis.

For the unrestricted assertion put $k=\mathbb Q$,
$R=k[[t]]$ and $K=k((t))$.  The ordered tuple
\[
 (0,1,\infty,2,2+t)
\]
defines a morphism $a:\operatorname{Spec}K\to\mathcal M_{0,5}$.
Let $r:G_k\to G_K$ be the uniformizer splitting of B.3.
Then $a_*\circ r$ is a section of
$\pi_1(\mathcal M_{0,5/k})\to G_k$.
The universal marked curve, pulled back by $a$, is the indicated
five-punctured projective line.  Fill its last two holes.  The
resulting curve is the constant tripod $U_0/k$, and the two
canonical sections of its arithmetic group are the composites
obtained from the $K$-points $2$ and $2+t$ with $r$.

Both points extend to $R$-points of $U_0$ with closed point $2$.
The specialization equivalence for finite \emph{\'etale} covers of
the henselian local scheme $\operatorname{Spec}R$ identifies
$\pi_1(R)$ with $G_k$.  In the commutative diagrams induced by the
two $R$-points, the map $G_k\xrightarrow{r}G_K\to\pi_1(R)$ is the
identity, and the map from the closed point to $U_0$ is $2$ in both
cases.  Thus both composites are $s_2$ up to geometric conjugacy.
They are the images of two different elements of $I'$, proving
failure of unrestricted injectivity.  In moduli this section is
tangential at the divisor where the fourth and fifth points
collide, so it fails exactly the distinctness test of B.25.
Scans 408--410 explicitly retain an admissibility reservation;
the example does not remove that reservation and then attribute
the enlarged assertion to a smooth marked curve.
\end{proof}

\begin{statement}{B.29}{Proved using the prompt's local theorem}{All fibres over the geometric locus of a deletion map}
\label{st:B.29}

Let $V/k$ be a hyperbolic curve over a number field, let $y\in V(k)$,
and put $U=V\setminus\{y\}$.  For the natural map
\[
 q_*:\operatorname{Sect}(\Pi_U/G_k)/\Delta_U
       \longrightarrow\operatorname{Sect}(\Pi_V/G_k)/\Delta_V,
\]
the fibre over $s_x$, $x\in V(k)\setminus\{y\}$, is the singleton
$\{s_x\}$; the fibre over $s_y$ is the entire cusp packet at $y$.
For every individual cuspidal class at a boundary point $z$ of $V$,
its fibre is one class in the corresponding packet of $U$.
Thus $q_*$ induces a bijection from the actual interior point
classes of $U$ onto those of $V$ other than $s_y$ and induces an
isomorphism on each old cusp packet.
For every proper hyperbolic $X/k$ with
$X(\mathbb A_k)_\bullet^{\mathrm{f\text{-}cov}}=X(k)$, all its
hyperbolic open subcurves have only point and cuspidal sections.
\end{statement}
\begin{proof}
Take a class in one of the specified fibres.  B.7 classifies its
localizations at every finite place; the real classification of
Part C supplies the remaining Selmer condition.  The corresponding local
classification for $V$ and the injection and disjointness of its
points and packets force every local compactification centre to
be $x$, $y$, or $z$, respectively.  B.14 with that single rational
point as finite support proves the same global alternative.
At $x\in U(k)$ there is just the point section.  At $y$ every
class of its packet maps to $s_y$, and no other class occurs.
At an old cusp $z$, the open immersion is an isomorphism on the
henselian punctured local neighbourhood.  Its homomorphism of
decomposition extensions identifies the two inertia groups and
their $\widehat{k^\times}$-torsors, proving the singleton assertion
for each specified cusp class.

For the final assertion, the stipulated proper theorem and its
finite-descent deduction make the compactification image of every
section the section of a rational point of $X$.  Proper point
injection at each finite place and B.7 force all local centres to
be that one point.  B.14 applies again.  This proves the advertised
affine conclusion on the indicated finite-descent class of proper
curves.

These fibre statements also give the exact logical use of axiom
(5): its clauses (a), (b), (c) are exclusion of $s_y$, injectivity,
and surjectivity onto the complement, on the admissible locus.
On a filled curve whose section set already consists of points
and packets, the computed fibres give the affine point-or-packet
classification and all three clauses.  For an arbitrary filled
curve, the remaining fibres are precisely those over section
classes outside its geometric locus; none is silently included
in this computation.
\end{proof}

\begin{statement}{B.30}{Proved using the prompt's local theorem}{A configuration section above distinct point classes}
\label{st:B.30}

Let $U/k$ be hyperbolic over a number field and let
$\operatorname{Conf}_n(U)\subset U^n$ be its ordered configuration
space.  Fix pairwise distinct $x_1,\ldots,x_n\in U(k)$.
The fibre of
\[
 \operatorname{Sect}(\pi_1(\operatorname{Conf}_n(U))/G_k)
       /\pi_1(\operatorname{Conf}_n(U)_{\bar k})
 \longrightarrow
 \prod_{i=1}^n\operatorname{Sect}(\Pi_U/G_k)/\Delta_U
\]
over $([s_{x_1}],\ldots,[s_{x_n}])$ is the singleton given by the
actual configuration $(x_1,\ldots,x_n)$.
\end{statement}
\begin{proof}
Proceed by induction on $n$.  The assertion for $n=1$ is literal.
Forgetting the last point is a hyperbolic-curve fibration
$p:\operatorname{Conf}_n(U)\to\operatorname{Conf}_{n-1}(U)$.
Its full geometric fundamental-group sequence is short exact,
including injection of the fibre group; this is
\cite[Proposition 2.2(i), with the total formation]{B_MochizukiTamagawa}.
All its hypotheses hold in characteristic zero.
By induction the image of the given section under $p_*$ is the
point section of $b=(x_1,\ldots,x_{n-1})$.
Conjugate the representative to have exactly that image.

The actual fibre is $U_b=U\setminus\{x_1,\ldots,x_{n-1}\}$.
The morphism $U_b\to\operatorname{Conf}_n(U)$ induces a continuous
isomorphism, with the fibre product taken along $s_b$,
\[
 \pi_1(U_b)\ \xrightarrow{\sim}\
 \pi_1(\operatorname{Conf}_n(U))
       \times_{\pi_1(\operatorname{Conf}_{n-1}(U))}G_k.
\]
Indeed it is the identity on the quotient $G_k$, and is an
isomorphism on geometric kernels by the homotopy exact sequence;
these facts give both injectivity and surjectivity.
Thus the given section is a section of $U_b/k$ whose image in
$\Pi_U$ is $s_{x_n}$.  Applying B.29 successively to the deleted
points makes its class uniquely the point $x_n$, since this point
is different from every deleted point.  In the configuration space
it is exactly $(x_1,\ldots,x_n)$.
Different conjugations of the base representative do not create
additional classes: the above argument supplies a geometric
conjugacy to the same actual configuration in every case.
\end{proof}

\begin{statement}{B.31}{Proved}{The lattice stabilizer in every hyperbolic type}
\label{st:B.31}

Let $U/k$ be hyperbolic over a number field, $x\in U(k)$, and
$\iota:k\hookrightarrow\mathbb C$.  Every extension
$j:\bar k\hookrightarrow\mathbb C$ gives a dense discrete subgroup
\[
 \Lambda_j=\operatorname{im}\bigl(
 \pi_1^{\rm top}(U_\iota(\mathbb C),x_\iota)
       \longrightarrow\Delta_U\bigr)
\]
whose profinite completion is $\Delta_U$.
Use $x$ for the effective action $\rho_x:G_k\to\operatorname{Aut}(\Delta_U)$,
and let $\sigma\cdot j=j\circ\sigma^{-1}$.  Then
\[
 \Lambda_{\sigma\cdot j}=\rho_x(\sigma)(\Lambda_j),\qquad
 \operatorname{Stab}_{G_k}(\Lambda_j)=G_{v,j}.
\]
Here $G_{v,j}=1$ at a complex place and is the order-two group of
complex conjugation at a real place.  The stabilizer of the oriented
lattice is $1$.  The unoriented orbit is $G_k/G_{v,j}$ and the
oriented orbit is a free transitive $G_k$-set of these comparisons.
Thus the effective action does not make the discretization
independent of the embedding choice.
\end{statement}
\begin{proof}
Riemann existence and residual finiteness of surface and free
groups give the injection and completion assertion.  Applying
$j\circ\sigma^{-1}$ to a finite cover and its fibre at $x$ gives
the displayed equivariance, including its inverse convention.

The exact arithmetic input is
\cite[Corollary 2.25]{B_HoshiMochizukiIV}:
inside $\operatorname{Out}(\Delta_U)$, the intersection of the
Galois image with
$\operatorname{Out}(\pi_1^{\rm top}(U_\iota(\mathbb C)))$ is $G_{v,j}$.
An automorphism preserving the actual based subgroup $\Lambda_j$
restricts to its discrete automorphism, so its outer class belongs
to that intersection.  The stabilizer is therefore contained in
$G_{v,j}$.  At a real place, complex conjugation is an actual
homeomorphism fixing $x_\iota$, and hence preserves the based lattice.
This proves the reverse inclusion.  It reverses the surface
orientation, giving the oriented assertion.  Orbit--stabilizer
and equivariance give the two orbit descriptions.

The cited theorem's proof was checked at the needed bridge:
its cartesian discrete/profinite configuration square
(Corollary 2.20(v)) lifts the discrete outer action to the third
configuration space; compatible tripod homomorphisms then reduce
the intersection to the faithful arithmetic tripod action.
This supplies the general-type step absent from the elementary
tripod proof B.20.
\end{proof}

\begin{statement}{B.32}{Proved using the prompt's local theorem}{All sections of proper constant families over a finite $p$-adic field}
\label{st:B.32}

Let $k$ be finite over $\mathbb Q_p$, let $S/k$ be a
hyperbolic curve, and let $X/k$ be a smooth proper
geometrically connected curve of genus at least two.
Every continuous homomorphism
$h:\Pi_S\to\Pi_X$ over $G_k$, modulo conjugation by
$\Delta_X$, is induced by a unique $k$-morphism $S\to X$.
Consequently the graph construction is a bijection
\[
 \operatorname{Hom}_k(S,X)
 \ \xrightarrow{\ \sim\ }\
 \operatorname{Sec}\bigl(\pi_1(X\times_k S)\to\Pi_S\bigr),
\tag{B.32.1}
\]
where section classes on the right are taken modulo
the geometric fibre group $\Delta_X$.
This closes the intermediate-image case of B.26 for
this field and proper-fibre range.
\end{statement}
\begin{proof}
For each finite extension $L/k$ and $a\in S(L)$,
the composite of the point section of $a$ with $h$
is a section of $\Pi_{X_L}\to G_L$.
The prompt's local theorem, used in B.7 with its
stated provenance, gives a unique point
$b\in X(L)$ inducing this composite. Thus $h$ is
point-theoretic in the exact sense of
\cite[Definition 3.2]{B_CarlsonStix}: it preserves
geometric section germs over all finite extensions.
Here properness excludes any cusp alternative.

The field $k$ is Kummer-faithful. For completeness,
an abelian variety's group $A(k)$ has an open
subgroup isomorphic to $\mathbb Z_p^d$ and finite
quotient, hence has no nonzero element divisible
by every positive integer. A split torus has the
same property from
$k^\times\simeq\mathbb Z\times\mathcal O_k^\times$;
a finite splitting extension treats all tori.
For an extension $1\to T\to G\to A\to1$, the
argument of B.1 applies with the finite exponent
of $A(k)_{\mathrm{tors}}$, and then that of
$G(k)_{\mathrm{tors}}$, to give
$\bigcap_n nG(k)=0$. This is the required
Kummer-faithfulness, and the same proof applies
over every finite extension.

Carlson--Stix, Proposition 3.17
\cite{B_CarlsonStix}, now gives the precise
dichotomy: $h$ is open, or its image is contained
in one point decomposition group. We use the
hyperbolic-target range of that proposition.
Its proof was checked at the use of arbitrarily
large finite-\'etale target covers, the
Abel--Jacobi dimension estimate of Proposition 3.16,
and the normalizer argument.
In the open case, the open Hom theorem of
Mochizuki, as in Proposition~\ref{A:relative-hom}, gives a unique dominant
$k$-morphism $S\to X$ inducing $h$.

In the other case the induced map on geometric
points is constant. Its equivariance under
$G_k$ makes the image a point $x\in X(k)$.
The image of $h$ is contained in a decomposition
group $D_x$; since both $h$ and $D_x\to G_k$
are over $G_k$, and the latter map is an
isomorphism, necessarily
$h=s_x\circ p_S$. This is induced by the
constant morphism at $x$.
Uniqueness of constant morphisms follows from
point-section injectivity. A dominant morphism
has open image, whereas $D_x$ has infinite
index, so the two cases do not overlap.

Finally, the product homomorphism is the actual
isomorphism of arithmetic extensions
\[
 \pi_1(X\times_k S)
       \simeq \Pi_X\times_{G_k}\Pi_S.
\]
On geometric kernels it is the product
isomorphism from the finite-\'etale K\"unneth
theorem, and on arithmetic quotients it is the
identity. A splitting of the projection to
$\Pi_S$ is therefore exactly
$\gamma\mapsto(h(\gamma),\gamma)$.
For a morphism $f:S\to X$ this is the induced
homomorphism of its graph
$a\mapsto(f(a),a)$, and geometric-fibre
conjugation is exactly the indicated conjugation
of $h$. These inverse constructions prove
(B.32.1).
\end{proof}

\begin{statement}{B.33}{Proved using the prompt's local theorem}{The complete constant-family decomposition and its cusp parameters}
\label{st:B.33}

Let $k$ be a number field or a finite extension of
$\mathbb Q_p$, and let $S/k$ and $U=X\setminus D$ be
hyperbolic curves. Write
$\operatorname{Sec}^{\mathrm{nc}}(U/k)$ for the field
section classes outside all cusp packets.
For each $y\in D(k)$ choose its inertia and decomposition
groups $I_y\subset D_y$ and put
\[
 \mathcal C_y(S)=
 \operatorname{Hom}_{G_k}(\Pi_S,D_y)/I_y.
\]
There is a disjoint decomposition, through the actual
product and graph maps of (B.32.1),
\[
 \operatorname{Sec}\bigl(\pi_1(U\times_k S)\to\Pi_S\bigr)
 \simeq
 \operatorname{Hom}^{\mathrm{dom}}_k(S,U)
 \ \amalg\
 \operatorname{Sec}^{\mathrm{nc}}(U/k)
 \ \amalg\
 \coprod_{y\in D(k)}\mathcal C_y(S).
\tag{B.33.1}
\]
Each $\mathcal C_y(S)$ is a torsor under
$H^1_{\mathrm{cts}}(\Pi_S,\widehat{\mathbb Z}(1))$.
The choice of a nonzero tangent at $y$ supplies its
origin. If $k$ is a number field, there is the canonical
Kummer identification
\[
 H^1_{\mathrm{cts}}(\Pi_S,\widehat{\mathbb Z}(1))
 \simeq
 \varprojlim_n
  \mathcal O(S)^\times/
                      \mathcal O(S)^{\times n}.
\tag{B.33.2}
\]
The right side is a power completion and is not
asserted compact. Inflation from $G_k$ is exactly
the subgroup arising from the constant units
$k^\times$. Thus all intermediate geometric images
in a constant family are accounted for by these
cusp torsors. The remaining noncuspidal
surjectivity question in (B.33.1) is precisely the
field-section question already isolated in B.27.
\end{statement}
\begin{proof}
The graph isomorphism identifies the left side with
$\operatorname{Hom}_{G_k}(\Pi_S,\Pi_U)/\Delta_U$.
Apply the infinite-index trichotomy of
Proposition~\ref{F-local-dichotomy}, including its
finite-$p$-adic proof. Its geometric image is either
open, zero, or a nonzero subgroup of one $I_y$.
In the first case Proposition~\ref{A:relative-hom}
gives exactly the dominant morphisms in (B.33.1).
In the second case the quotient property of
$\Pi_S\to G_k$ gives a unique field-section
factorization. Its noncuspidal classes give the
middle summand; its cuspidal classes belong to the
last summands. In the third case the inertia
normalizer identity puts the entire image in
$D_y$, and projection onto $G_k$ makes $y$
$k$-rational. This proves exhaustion.

Choose a tangential splitting $s_y:G_k\to D_y$.
Every homomorphism into $D_y$ has a unique expression
\[
 h(\gamma)=c(\gamma)s_y(p_S(\gamma)).
\]
The homomorphism equation is the continuous
$1$-cocycle equation for $c:\Pi_S\to I_y$, where
the action is cyclotomic through $G_k$.
Conjugation by $I_y$ changes $c$ by a coboundary,
giving the stated torsor.
We check that passage from $I_y$-conjugacy to
$\Delta_U$-conjugacy makes no additional
identifications. For a nonzero geometric image,
the peripheral intersection theorem forces a
geometric conjugator between two maps with image
in $D_y$ to belong to $I_y$. This is the same
nonzero-subgroup inertia calculation used in
Proposition~\ref{F-local-dichotomy}.
For zero geometric image the assertion is exactly
injectivity of the field cusp packet. Different
cusps cannot meet in either case, and zero and
nonzero geometric image are preserved under
conjugation. Finally the open case has infinite
nonprocyclic geometric image, so it meets none
of these summands. The field-section summand
was defined by removing its cusp packets.
This proves disjointness and injectivity.

For (B.33.2), the Kummer sequence on $S$ gives
\[
 0\to\mathcal O(S)^\times/\mathcal O(S)^{\times n}
   \to H^1(S,\mu_n)\to\operatorname{Pic}(S)[n]\to0.
\]
The finite-coefficient group in the middle is
$H^1(\Pi_S,\mu_n)$ by the classification of
finite-\'etale torsors. The degree map on the
smooth proper completion of $S$ has image in
$\mathbb Z$ and kernel a subgroup of its
Jacobian's $k$-points. Mordell--Weil therefore
makes that completion's Picard group finitely
generated. Quotienting by the finitely many
boundary-divisor classes gives
$\operatorname{Pic}(S)$, also finitely generated.
Consequently $T\operatorname{Pic}(S)=0$.
The inverse system of unit quotients is
surjective, while the finite groups
$H^0(S,\mu_n)$ are Mittag--Leffler along a
factorial tower. Taking inverse limits in the
displayed sequence and in continuous degree-one
cohomology therefore proves (B.33.2).

Naturality of Kummer identifies inflation with
constant units. More explicitly,
$\mathcal O(S)^\times/k^\times$ embeds by its
boundary divisor into a finitely generated
free abelian group, so it is free of finite
rank. The resulting unit sequence splits as
an abstract abelian-group sequence; its power
completions remain exact and the constant
subgroup injects. The kernel of restriction
from $H^1(\Pi_S,I_y)$ to $H^1(\Delta_S,I_y)$
is exactly the inflated $H^1(G_k,I_y)$ by
the continuous Hochschild--Serre sequence.
Hence the maps in this subgroup are precisely
the field cusp sections composed with $p_S$.
All other classes in the relative cusp torsor
have the nonzero intermediate image already
described. A morphism-induced constant class
in the middle summand is exactly an actual
point of $U(k)$, by B.26. This proves the final
assertion without filling the remaining field
section term by an unproved point set.
\end{proof}

\begin{statement}{B.34}{Refuted}{The unrestricted interior-point reflection in the printed RNS Theorem F(ii)}
\label{st:B.34}

Let $k=\mathbb Q_5$, and let
\[
 Y=\mathbb P^1_k\setminus\{0,1,\infty,2\},\qquad
 Z=\mathbb P^1_k\setminus\{0,1,\infty\},\qquad
 f:Y\hookrightarrow Z.
\]
A tangential section $s$ at the cusp $2$ of $Y$ is not a
point section of $Y$, whereas $f_*s$ is the point section of
$2\in Z(k)$. This satisfies every field, prime, subgroup and
map hypothesis in the literal interior-point statement of
Mochizuki--Tsujimura, Theorem F(ii), and Proposition 3.5(iii)
\cite{A_MochizukiTsujimuraRNS}. Thus the stated equivalence for
arbitrary dominant morphisms, with point sets $Y(\Omega^H)$
and $Z(\Omega^H)$, is false. The same curves and sections
over $\mathbb Q$ refute its unrestricted number-field analogue.
This statement concerns that specific reflection assertion;
it makes no assertion against the source's RNS Theorem A or
its Isom theorems.
\end{statement}
\begin{proof}
Choose $t=x-2$ and a nonzero tangent at $2$. The compatible-root
construction of B.3 gives
\[
 s:G_k\longrightarrow D_2(Y)\longrightarrow\Pi_Y,
 \qquad
 1\longrightarrow\widehat{\mathbb Z}(1)
 \longrightarrow D_2(Y)\longrightarrow G_k\longrightarrow1.
\]
The inclusion $f$ extends across $t=0$. It kills this inertia,
and its restriction to $D_2(Y)$ factors through the residue
map $D_2(Y)\to G_k$ followed by the point decomposition
section at $2$ in $Z$. Hence $f_*s=s_2$ up to geometric
conjugacy, for every tangent and indeed every member of
the packet.

The invertible function $x$ on $Y$ gives the actual map
$x_*:\Pi_Y\to\Pi_{\mathbb G_m}$.
At $2$, write $x=2(1+t/2)$; the factor $1+t/2$ has unique
compatible roots with residue $1$. Therefore $x_*s$ has
Kummer class $\kappa_k(2)$. If $s$ were induced by
$y\in Y(k)$, functoriality would give
$\kappa_k(x(y))=\kappa_k(2)$.
The local Kummer injection proved in B.8 forces $x(y)=2$,
which is excluded from $Y$. Over $\mathbb Q$, the same
argument uses B.1, or directly the valuations and the
finite group $\{\pm1\}$ to show
$\bigcap_n(\mathbb Q^\times)^n=\{1\}$.

Here are the exact local hypotheses. Take $K=\mathbb Q_5$,
$p=5$, $\Sigma$ equal to all primes, $\ell=3$, and
$H=G_K$. The residue field is $\mathbb F_5$; both curves
and the dominant open immersion are defined over the
$p$-adic local field $K$ itself. Tame inertia gives a
continuous surjection $H\cap I_K=I_K\twoheadrightarrow
\mathbb Z_3(1)$, hence to $\mathbb Z_3$ as a profinite
group. The $3$-adic cyclotomic image is the closure of
$5^{\mathbb Z}$ in $\mathbb Z_3^\times$, which is open:
$v_3(5^2-1)=1$, so the powers of $25$ generate
$1+3\mathbb Z_3$. Finally
$\Omega=\mathbb C_5$ and $\Omega^H=\mathbb Q_5$.
For the last equality one may apply the Ax--Sen--Tate
estimate recalled and sharpened in \cite{B_LeBorgne}:
an invariant element can be approximated arbitrarily
closely by elements of the complete field $K$ and
therefore belongs to $K$.

The source's printed page 8 uses the actual point sets
$Y(\Omega^H)$ and $Z(\Omega^H)$; its scheme and curve
conventions on pages 13--14 do not replace them with
compactifications. No finiteness or monocuspidal hypothesis
is imposed on $f$ there. In the proof of Proposition
3.5(iii), printed page 79, the reduction replaces finite
section neighbourhoods by their smooth compactifications.
In this example, every such compactification contains
the centre of the lifted cusp. The resulting point
does not lie in the original open neighbourhood.
That reduction therefore does not preserve the
interior-point property needed for the claimed equivalence.
\end{proof}

\begin{statement}{B.35}{Proved, using the stipulated local theorem}{The exact fibre over a point under a dominant morphism}
\label{st:B.35}

Let $k$ be a number field or a finite extension of
$\mathbb Q_p$, and let $f:Y\to Z$ be a dominant morphism
of hyperbolic curves over $k$. Write
$\bar f:\bar Y\to\bar Z$ for its extension to smooth
proper completions. For $z\in Z(k)$, the full-section
fibre is the disjoint union
\[
 f_*^{-1}([s_z])
 =\{[s_y]:y\in Y(k),\ f(y)=z\}
 \ \sqcup\!
 \coprod_{\substack{c\in(\bar Y\setminus Y)(k)\\
                         \bar f(c)=z}}
                  \operatorname{Packet}_c(Y).
 \tag{B.35.1}
\]
In particular, dominant morphisms preserve point
geometricity, and they reflect it on noncuspidal source
sections. Finite morphisms preserve and reflect point
geometricity on all sections. For finite \emph{\'etale}
morphisms the latter assertion follows directly from
the finite-cover action on the fibre over $z$, without
the stipulated local theorem and over every
characteristic-zero field.
\end{statement}
\begin{proof}
The map $\bar f$ is finite. A point section at $y$ maps
to that at $f(y)$. For a boundary point $c$ with
$\bar f(c)=z\in Z(k)$, the henselian map extends into
$Z$ at $c$, so its cusp inertia is killed and the
entire packet maps to $s_z$, exactly as in B.34.
These are all the displayed inclusions.

Over a finite extension of $\mathbb Q_p$, B.7 classifies
every source section as a point or a cusp. If a cusp
maps to a target cusp, B.4 maps its packet into that
target packet, which is disjoint from the point
sections. The remaining centres map into $Z$; local
point-section injection then makes their image equal
to $z$. This proves (B.35.1) locally.

Now let $k$ be a number field and let $s$ belong to
the left side of (B.35.1). At every finite place,
the local result just proved puts the compactification
centre of $s_v$ in the fixed finite $k$-subscheme
\[
 W=(\bar f^{-1}(z))_{\mathrm{red}}\subset\bar Y.
\]
The stipulated local theorem, B.7, supplies all finite
Selmer localizations; the real section comparison
of C.4 supplies the real ones. Apply B.14 to this
single finite support $W$. It proves that $s$ is a
point section at a member of $W(k)\cap Y(k)$ or a
cusp section at a member of
$W(k)\cap(\bar Y\setminus Y)(k)$.
Full point/packet injection and disjointness give
exactly the disjoint union in (B.35.1).

If $f$ is finite, normalization in $k(Y)$ identifies
$Y$ with $\bar f^{-1}(Z)$. Thus no boundary point
$c$ lies over $z\in Z$, and the packet terms vanish.
For finite \emph{\'etale} $f$, there is also the direct
proof valid over an arbitrary characteristic-zero
field. Identify $\Pi_Y$ with the corresponding open
subgroup of $\Pi_Z$. After geometric conjugation,
a section mapping to $s_z$ makes the appropriate
coset in the geometric finite fibre over $z$ fixed
by $s_z(G_k)$. That fixed coset is an actual
$k$-point $y$ of the finite \emph{\'etale} scheme
$f^{-1}(z)$. The decomposition group at that point
identifies the source section with $s_y$ inside
$\Pi_Y$. Conversely $s_y$ maps to $s_z$ by
functoriality. This proves the final assertion
through the actual finite-cover map.
\end{proof}

\subsection{The closed obstruction supplied by stable boundary cycles}

\begin{statement}{B.36}{Proved}{A closed obstruction to smooth arithmetic realization}
\label{st:B.36}

Fix a hyperbolic type $(g,n)$, its geometric full profinite
curve group $\Delta$, and its prescribed peripheral conjugacy
classes. Let $Q\subset\operatorname{Out}(\Delta)$ be the closure
of the geometric mapping class group, and let
$A\subset\operatorname{Aut}(\Delta)$ be its inverse image.
For each nonempty stable multicurve type, choose oriented
generators $c_1,\ldots,c_r\in\Delta$ of its node groups. Call
their $A$-translates, with arbitrary individual conjugations,
\emph{pro-simple multicurve tuples}.

Let $k$ be a number field, let $E\twoheadrightarrow G_k$ be a
profinite extension, and let
$\rho:E\to\operatorname{Out}(\Delta)$ be continuous, with the
image of the geometric kernel contained in $Q$. The following
condition on a continuous section $s$ defines a closed,
geometric-conjugacy-invariant subset of its section space:
there is a nonempty pro-simple multicurve tuple such that,
for every $\gamma\in G_k$, there are a permutation $\sigma$
and signs $\epsilon_i\in\{1,-1\}$ satisfying
\[
 \rho(s(\gamma))(c_i)
 \ \sim_\Delta\quad
 c_{\sigma(i)}^{\epsilon_i\chi(\gamma)}
 \qquad(1\leq i\leq r).
 \tag{B.36.1}
\]
Here the left side is a conjugacy class, so the condition is
independent of the representative of the outer automorphism.
No outer action satisfying this condition is the peripheral
outer action of a smooth marked curve over $k$, or over any
finite extension of $k$ after restriction.
\end{statement}
\begin{proof}
We first make the two compactness and nonperipherality
assertions precise. A finitely generated profinite group has
a basis of open characteristic subgroups; its automorphism
group is the closed inverse limit of the compatible finite
quotient automorphism groups. Hence $\operatorname{Aut}(\Delta)$
and its closed subgroup $A$ are compact. For a fixed tuple the
image of
\[
 A\times\Delta^r\longrightarrow\Delta^r,\qquad
 (a,h_1,\ldots,h_r)\longmapsto
 (h_i a(c_i)h_i^{-1})_{i=1}^r
 \tag{B.36.2}
\]
is compact. There are only finitely many stable multicurve
types: the number of curves is at most $3g-3+n$, and cutting
along them records a finite graph, the genera of its vertices,
and the distribution of the $n$ labelled punctures. Stability
and the genus formula bound all of these data. Classification
of the cut oriented surfaces identifies the resulting finite
list of mapping class orbits. Thus the disjoint union
$\mathcal K$ of the compact tuple spaces (B.36.2), over that
finite list, is compact. For $(g,n)=(0,3)$ the list is empty
and the assertion is vacuous.

Each initial $\langle c_i\rangle$ is an injected copy of
$\widehat{\mathbb Z}$, and is not contained in a peripheral
group. One elementary finite-cover verification is useful
here. A nonseparating curve has a primitive nonzero class
in the homology of the filled surface. For a separating
essential curve, each cut side admits a connected double
cover trivial on that boundary circle: on a positive-genus
side use a nonzero handle character to $\mathbb Z/2$; on
a genus-zero side use the character nonzero on two of its
punctures. The latter side has at least two punctures because
the curve is essential and nonperipheral. Gluing the two
connected covers along their two boundary lifts gives a
double cover in which each lift of the original curve is
nonseparating, including after the punctures are filled.
Its class is therefore primitive in compact homology.
Riemann existence gives the same finite cover of the
profinite surface group. A cusp subgroup, and its intersection
with this finite-cover subgroup, maps to zero in compact
homology. This proves the claimed injection and exclusion.
The groups belonging to distinct entries of a multicurve
are pairwise nonconjugate, by the distinct-edge assertion
in \cite[Proposition 1.2]{F-MochizukiCombinatorial}.
Every member of $A$ preserves the peripheral conjugacy
classes, as does its inverse. Applying the actual
automorphism and the individual conjugations in (B.36.2)
therefore preserves this exclusion. The limiting tuples
are pro-simple; they need not be discrete simple curves.

Conjugacy in a profinite group is a closed relation:
it is the image of the compact map
$(h,x)\mapsto(x,hxh^{-1})$. Exponentiation
$(x,u)\mapsto x^u$ for $u\in\widehat{\mathbb Z}$ is
continuous. For a fixed $\gamma$, condition (B.36.1) is
therefore a finite union of closed conditions on a section
and a member of $\mathcal K$. Intersect over all $\gamma$.
The resulting incidence subset is closed, and projection
along the compact space $\mathcal K$ is closed.

This argument uses the usual section topology, with no
compactness assumption on the set of sections itself.
For completeness, view a section as its closed image
subgroup $H\subset E$, with the subgroup topology used in
\cite[Section 3]{B_CarlsonCorwin}. Evaluation at a fixed
$\gamma$ is continuous on these subgroups: for a convergent
net $H_j\to H$, every convergent subnet of the elements
$s_j(\gamma)$ has limit in $H$ lying above $\gamma$;
there is exactly one such element. Compactness of $E$
then gives $s_j(\gamma)\to s(\gamma)$. Quotienting by
the compact geometric conjugation group preserves the
closed conclusion. Conjugation changes the witness tuple
by an element of $A$, so the subset is invariant.

Suppose now that the restricted action is realized by a
smooth marked curve $U/L$, with $L/k$ finite. Centre-freeness
of $\Delta$ identifies its arithmetic extension with the
pullback
\[
 E_{s,L}=\operatorname{Aut}(\Delta)
       \mathop{\times}_{\operatorname{Out}(\Delta)}G_L
       \ \simeq\ \pi_1(U),
 \tag{B.36.3}
\]
where the second isomorphism is the specified peripheral
realization. Choose one group $I=\langle c_i\rangle$ from
the witness. Its conjugacy class has a finite orbit under
$G_L$. Its stabilizer $G_0$ is closed, since it is the image
of the compact normalizer $N_{E_{s,L}}(I)$, and is of finite
index, hence open. Within that normalizer put
\[
 N^{\pm\chi}=
 \{e:\operatorname{Ad}(e)(c_i)
            =c_i^{\pm\chi(p(e))}\},\qquad
 N^\chi=
 \{e:\operatorname{Ad}(e)(c_i)
            =c_i^{\chi(p(e))}\}.
 \tag{B.36.4}
\]
Condition (B.36.1), with its existential conjugator,
makes $N^{\pm\chi}\to G_0$ surjective. Both are closed
subgroups. The ratio of the conjugation character and
$\chi\circ p$ gives a continuous homomorphism
$N^{\pm\chi}\to\{1,-1\}$ with kernel $N^\chi$.
Consequently $p(N^\chi)$ has index at most two in $G_0$;
it is compact and thus open in $G_L$. This is an actual
closed subgroup of the arithmetic normalizer with open
arithmetic image acting exactly cyclotomically on $I$.
No splitting of this normalizer is assumed.

The cyclotomic-normalizer theorem, in the exact form
\ref{A:cyclotomic-normalizer}, or
\cite[Theorem 20]{B_StixCuspidal}, puts $I$ in a cusp
inertia group of $U$. Its number-field hypothesis is
verified by the incompatible Frobenius weights of a torus
and a Jacobian at a good finite prime. Under (B.36.3)
the cusp groups are the prescribed peripheral groups.
This contradicts the exclusion proved above, and proves
nonrealization over every finite extension.
\end{proof}

\begin{statement}{B.37}{Proved}{Stable-boundary cuspidal sections are non-admissible}
\label{st:B.37}

Let $k$ be a number field, $2g-2+n>0$, and set
$\mathcal M=\mathcal M_{g,n,k}$ and
$\overline{\mathcal M}=\overline{\mathcal M}_{g,n,k}$.
For a $k$-rational stable boundary object $x$, let
$\mathcal V_x$ be the intersection with $\mathcal M$ of
the henselian neighbourhood of its residual gerbe in
$\overline{\mathcal M}$. Define
\[
 C_{\mathcal M,\overline{\mathcal M}}
 =\bigcup_{x\in(\overline{\mathcal M}\setminus\mathcal M)(k)}
 \operatorname{Im}\bigl(
       \operatorname{Sec}(\mathcal V_x/k)
        \longrightarrow\operatorname{Sec}(\mathcal M/k)\bigr),
 \qquad
 C^{\mathrm{cl}}_{\mathcal M,\overline{\mathcal M}}
 =\overline{C_{\mathcal M,\overline{\mathcal M}}}.
 \tag{B.37.1}
\]
Isomorphic boundary objects give the same term. For schemes
this is the nearby-point definition of Carlson--Corwin;
for stacks the residual gerbe and its finite automorphisms
are retained as described below.

Every section in $C^{\mathrm{cl}}_{\mathcal M,\overline{\mathcal M}}$
has an outer action on the universal geometric punctured
curve group which is not realized by a smooth marked curve,
even after restriction to any finite extension of $k$.
In particular, these sections are non-admissible in the
constructive sense of scans 480 and 487--488, and none is
the section of an interior moduli point. The same conclusion
holds for unordered markings, with the associated finite
permutation quotient retained.
\end{statement}
\begin{proof}
We first specify the scheme and stack maps in (B.37.1).
For a smooth scheme pair $V\subset X$ with normal crossings
boundary and $x\in(X\setminus V)(k)$, Carlson--Corwin define
the packet using the actual morphism
\[
 X_{(x)}\times_X V\longrightarrow V.
 \tag{B.37.2}
\]
Their excision lemma makes this construction invariant under
\emph{\'etale} replacement with the same residue point.
Their local fundamental-group proposition identifies its
geometric kernel with the dual of the boundary-character
group, twisted by $\widehat{\mathbb Z}(1)$; their Theorem A
asserts independence of the \emph{closure} from the normal
crossings compactification \cite{B_CarlsonCorwin}.
We have read those precise statements and their local
comparison. The theorem is stated for schemes. For the
stack in this statement we keep its specified stable
compactification and now give the local stack construction.

The residual gerbe at the object $x$ is
$B_k\operatorname{Aut}(x)$. Its henselian neighbourhood
is formed by the inverse system of representable
\emph{\'etale} neighbourhoods inducing an isomorphism on
that gerbe. In characteristic zero its finite automorphism
group is \emph{\'etale} and linearly reductive. The local
quotient construction of
\cite[Theorem 3.2 and Proposition 3.6, Step 1]{B_AOV}
gives a finite \emph{\'etale} torsor
\[
 \operatorname{Spec}R\longrightarrow
 [\operatorname{Spec}R/A]
       =\overline{\mathcal M}_{(x)},
 \tag{B.37.3}
\]
in that henselian system. Here $R$ is regular henselian,
$A$ is the finite group scheme extending
$\operatorname{Aut}(x)$, and its closed fibre torsor
over $B_k\operatorname{Aut}(x)$ has total space
$\operatorname{Spec}k$. The corrected obstruction argument
in \cite{B_AOVCorrigendum} also gives (B.37.3); in the
present \emph{\'etale} characteristic-zero case no
nonreduced stabilizer occurs. Removing the pulled-back
boundary gives
\[
 V_R=\operatorname{Spec}R\setminus D_R
       \longrightarrow [V_R/A]=\mathcal V_x
       \longrightarrow\mathcal M.
 \tag{B.37.4}
\]
These are the actual maps used to define its packet.
Changing the presentation gives the same henselian
neighbourhood and the same section image. A finite
extension splitting $A$ and the boundary branches makes
the first map an ordinary finite-group torsor. A section
of its quotient lifts to that torsor after restriction to
an open subgroup of $G_k$, by taking the kernel of its
finite torsor monodromy. The finite quotient is not
discarded in (B.37.1).

Let $r>0$ be the number of nodes of the geometric stable
curve $x$. In an \emph{\'etale} deformation chart the
boundary smoothing parameters are $t_1,\ldots,t_r$;
near node $i$ the universal curve has equation
$uv=t_i$. After the finite extension which fixes the
nodes and their branches, the local base group has
geometric inertia
$T^r$, where $T=\widehat{\mathbb Z}(1)$.
The local node group has geometric inertia
$T^{r+1}$, and its map to the base is
\[
 (a,b,d_2,\ldots,d_r)
        \longmapsto (a+b,d_2,\ldots,d_r),
 \qquad
 I_i=\{(a,-a,0,\ldots,0):a\in T\}.
 \tag{B.37.5}
\]
The coordinate numbered $1$ in this display is $t_i$.
In particular the map is surjective and the kernel
is an injected copy of $T$.

Here injection into the full geometric curve group is
essential. It is exactly the node injection and addition
map of \cite[Lemma 2.11(ii)--(iii)]{F-HoshiMochizukiNodes},
with all primes allowed. The source's Definition 2.10
identifies the full geometric log fibre group with the
punctured smooth surface group and indexes the local
character monoid by all the nodes. The comparison with
the generic smooth curve and with its universal-family
outer action is the comparison of finite covering
categories: use \cite[Theorem 7.6]{F-Illusie} for the
interior and \cite[Corollary 3.9]{F-HoshiLog} for the
completed strictly local fibre. In characteristic zero
every ramification index here is tame. Passing from the
henselization to its completion preserves the residue
group and the normal-crossings character monoid, hence
induces an isomorphism on their local groups by the same
Kummer calculation. All these maps preserve the base
quotient and the node kernel (B.37.5).

For any section of the local base group, pull back the
surjective node-group map along that section. It gives
a closed subgroup of the corresponding arithmetic
curve group which projects onto $G_k$ and acts on
$I_i$ by $\chi$. Thus this assertion holds for every
parameter in the packet; it does not restrict that
parameter to the image of an actual tangent vector
or of a one-dimensional arc. Without the preliminary
field extension, arithmetic elements permute the nodes;
exchanging the two branches changes $a$ to $-a$ in
(B.37.5). An automorphism in the finite quotient $A$
has the same effect: it permutes nodes and branches,
and a branch-preserving local coordinate change
$u\mapsto u\cdot\text{unit}$ acts identically on its
inertia. Consequently the full local group, including
$A$, satisfies exactly the permutation and $\pm\chi$
formula (B.36.1).

Choose a complex embedding of $\bar k$ and a geometric
base point of $\mathcal M$. Analytically smoothing the
stable fibre identifies its node groups with a nonempty
multicurve in the punctured topological fibre. Stability
ensures that these curves are essential and
nonperipheral. Changing the path from this local fibre
to the chosen geometric base point acts through the
geometric monodromy group $Q$; changing the fibre base
paths gives the individual conjugations of (B.36.2).
The finite-cover comparisons just described identify
these as the same full profinite node groups. Therefore
each packet in (B.37.1) supplies a tuple belonging to
the compact space $\mathcal K$ of B.36, and its moduli
section lies in the closed obstruction subset proved
there. Its entire closure lies in that same subset.
This last compactness step prevents a node witness
from being lost when boundary packets vary.

Apply B.36 to $E=\pi_1(\mathcal M)$ and the outer
monodromy of the universal punctured curve. It excludes
a smooth peripheral realization after every finite
extension. A moduli point gives precisely such a
realization, by pulling back that universal curve.
The constructive admissible germs of the manuscript
are obtained from an actual tripod curve by finite
\emph{\'etale} covering and filling specified holes;
those operations produce smooth marked curves and
their actual outer actions. Thus a boundary section
cannot be one of those germs. Passing to unordered
markings only adds the finite permutations already
allowed in the incidence relation and in (B.37.4), so
the same proof applies.

This is a necessary closed obstruction. It does not
assert that every section outside the closed boundary
set is a moduli point, or that the absence of these
pro-simple cyclotomic witnesses is a sufficient
realization criterion.
\end{proof}

\begin{statement}{B.38}{Proved using the stipulated local theorem}{Unramification over a p-adic base curve}
\label{st:B.38}

Let $k/\mathbb Q_p$ be finite, let $B/k$ be a smooth
geometrically connected hyperbolic curve, put $K=k(B)$,
and let $f:X\to B$ be a smooth proper family of
geometrically connected curves of genus at least two.
Write $E_X=\pi_1(X)$, $E_B=\pi_1(B)$, and let $\Delta$
be the geometric generic fibre group. The actual
generic-fibre comparison is
\[
 \pi_1(X_K)\ \xrightarrow{\ \sim\ }\
 E_X\times_{E_B}G_K.
 \tag{B.38.1}
\]
Pullback along $G_K\twoheadrightarrow E_B$ induces a
bijection
\[
 \operatorname{Sec}(E_X\to E_B)
 \ \xrightarrow{\ \sim\ }\
 \operatorname{Sec}(\pi_1(X_K)\to G_K),
 \tag{B.38.2}
\]
where both sides are taken modulo conjugation by the
same group $\Delta$. Every section on the left is
point-theoretic over $k$: for each finite extension
$l/k$ and each $b\in B(l)$, its restriction to the
point section at $b$ is induced by a unique point of
$X_b(l)$.

The maps on algebraic sections and on fundamental-group
sections fit in the commutative diagram
\[
\begin{array}{ccc}
 \operatorname{Sect}_B(X)&\xrightarrow{\sim}&X_K(K)\\
 \downarrow&&\downarrow\\
 \operatorname{Sec}(E_X\to E_B)&\xrightarrow{\sim}&
       \operatorname{Sec}(\pi_1(X_K)\to G_K).
\end{array}
\tag{B.38.3}
\]
The upper bijection is extension by properness. This
statement does not assert surjectivity of either
vertical map for a nonconstant family.
\end{statement}
\begin{proof}
The homotopy sequence for a relative hyperbolic curve
over this base is short exact:
\[
 1\longrightarrow\Delta\longrightarrow E_X
       \longrightarrow E_B\longrightarrow1.
 \tag{B.38.4}
\]
One can obtain its left exactness from the homotopy
fibre comparison and the fact that $B$ is an \'etale
$K(\pi,1)$. More generally, the kernel of the fibre
map lies in the centre of the hyperbolic fibre group
and is therefore trivial. These are precisely the
quasifibration comparison and the centre-free argument
in the proof comparing total and generic-fibre
$\pi_1$-openness in the supplied primary source
\cite{F-HolzschuhSchmidtStix}. The proper smooth fibre
comparison identifies $\Delta$ with the full geometric
group of every fibre, not merely its abelianization.
Apply the same exact sequence over $K$. The natural
map from $\pi_1(X_K)$ to the pullback in (B.38.1)
is the identity on $\Delta$ and on the quotient
$G_K$. The two exact sequences prove that it is an
isomorphism. This also identifies the two conjugating
groups in (B.38.2).

Let $s:G_K\to\pi_1(X_K)$ be a section and let
$h:G_K\to E_X$ be its first component under (B.38.1).
We prove that $h$ kills every inertia group belonging
to a closed point of $B$. Fix such a point $b$, with
residue field $k_b$. The completed local field is
$K_b\simeq k_b((t))$; choose a coefficient field and
a uniformizer. Its decomposition map to $E_B$
factors through the actual point section
$G_{k_b}\to E_B$, because a finite \'etale cover of
$B$ is unramified at $b$. In the proper smooth family,
the actual preimage of this point-section subgroup is
\[
 E_X\times_{E_B}G_{k_b}\ \simeq\ \pi_1(X_b).
 \tag{B.38.5}
\]
Indeed the map from the fibre group gives the identity
on the proper smooth geometric fibre group and on
$G_{k_b}$ in the two exact sequences. In particular,
the restriction of $h$ gives a homomorphism
\[
 G_{k_b((t))}\longrightarrow\pi_1(X_b)
                 \quad\text{over }G_{k_b}.
 \tag{B.38.6}
\]

The valuation in (B.38.6) has residue characteristic
zero. Its entire inertia is tame, and the chosen
coefficient field and compatible roots of $t$ give
\[
 G_{k_b((t))}\simeq
 \widehat{\mathbb Z}(1)\rtimes G_{k_b}
 \simeq\pi_1(\mathbb G_{m,k_b}),
 \tag{B.38.7}
\]
with the standard section of $\mathbb G_m$ at $1$.
Now precompose (B.38.6) with the map induced by the
actual open immersion
\[
 \mathbb P^1_{k_b}\setminus\{0,1,\infty\}
           \longrightarrow\mathbb G_{m,k_b}.
 \tag{B.38.8}
\]
Its map on geometric groups surjects onto
$\widehat{\mathbb Z}(1)$. The image of that group
in the proper hyperbolic fibre group is therefore
procyclic. The finite-$p$-adic-field version of
Proposition~\ref{F-local-dichotomy}, whose proof
uses the stipulated local section theorem, says
that this image is open, trivial, or contained in a
nontrivial cusp inertia. Properness eliminates the
last alternative. The open alternative is impossible:
an open subgroup of the genus-at-least-two proper
geometric group is the fundamental group of a finite
\'etale cover of genus at least two, and its
abelianization has $\mathbb Z_\ell$-rank at least
four for every prime $\ell$. A procyclic group has
rank at most one. Thus (B.38.6) kills the entire
$t$-adic inertia.

The kernel of $G_K\twoheadrightarrow E_B$ is the
closed normal subgroup generated by these inertia
groups. This follows directly from the finite-cover
description: a finite separable extension unramified
at every closed point of the regular curve $B$ has
finite \'etale normalization over $B$, and every
finite \'etale cover has that unramified property.
Consequently $h$ factors uniquely through a continuous
map $\bar h:E_B\to E_X$. The equation $f_*h=p_B$
and surjectivity of $p_B:G_K\to E_B$ make $\bar h$
a section. Conversely, pull back any such $\bar h$
using (B.38.1). These constructions are inverse on
actual homomorphisms and commute with conjugation by
$\Delta$, proving (B.38.2). Keeping (B.38.1) is
essential here: it proves that the pulled-back section
is the original $s$, rather than only a section with
the same image in $E_X$.

After any finite extension $l/k$, restrict a relative
section to the actual point section of $b\in B(l)$.
The same fibre comparison gives a section of
$\pi_1(X_b/l)$. The stipulated proper local theorem
gives its unique realizing point in $X_b(l)$. Thus
the map $E_B\to E_X$ has the asserted point-theoretic
property over every finite extension, with fibre
membership guaranteed by (B.38.5).

Finally, a $K$-point of the proper generic fibre
extends over every local discrete valuation ring of
$B$ by properness; uniqueness by separatedness glues
these extensions to a section of $X\to B$.
Restriction is its inverse, and functoriality gives
(B.38.3). Every smooth proper hyperbolic curve over a
one-variable function field over a finite $p$-adic
field admits such a model after replacing the base
by a nonempty hyperbolic open. The missing general
step is therefore the realization of these relative,
point-theoretic sections for nonconstant families.
The dominant relative Hom theorem in
\cite{F-HolzschuhSchmidtStix} does not supply that
step: the generic map from a point to a curve is not
dominant, and the corresponding section subgroup is
not open. A map on fundamental groups also does not
by itself provide a map of the full \'etale topoi.
No such extension of its hypotheses is used here.
\end{proof}

\begin{statement}{B.39}{Proved using the stipulated local theorem}{Geometric-image trichotomy over every sub-p-adic field}
\label{st:B.39}

Let $K$ be a sub-$p$-adic field, let $V/K$ be a smooth
geometrically connected variety, let $U/K$ be a
hyperbolic curve, and let
$h:\pi_1(V)\to\pi_1(U)$ be a continuous
$G_K$-homomorphism. Write $P_V,P_U$ for the geometric
groups. Exactly one of the following holds:
\[
 h(P_V)\text{ is open in }P_U;\qquad
 h(P_V)=1;\qquad
 1\ne h(P_V)\subset I_c
 \text{ for a unique lifted cusp }c.
 \tag{B.39.1}
\]
In the third case $h(\pi_1(V))\subset D_c$, and the
underlying cusp is $K$-rational. In the second case
$h$ factors through an actual section
$G_K\to\pi_1(U)$; no point realization of that section
is asserted.

The transfer used here preserves the full geometric
homomorphism. More precisely, there exist an extension
$K\hookrightarrow F$, a countable field $F_0\subset F$,
an embedding $F_0\hookrightarrow l$ into a finite
extension of $\mathbb Q_p$, and a descended datum
$(V_0,U_0,h_0)/F_0$ whose base change to $F$ is
$(V_F,U_F,h_F)$. The base change to $l$ has the same
geometric homomorphism through the characteristic-zero
base-change isomorphisms, including the peripheral
subgroups of the target. These are the actual field
maps used; an embedding of $F$ into $l$ is unnecessary.
\end{statement}
\begin{proof}
By definition of sub-$p$-adic, choose $F$ finitely
generated over $\mathbb Q_p$ containing $K$. Base
change gives
\[
 \pi_1(V_F)=\pi_1(V)\times_{G_K}G_F,
 \qquad
 \pi_1(U_F)=\pi_1(U)\times_{G_K}G_F.
 \tag{B.39.2}
\]
Here and below the compatible geometric base points
are understood. The geometric comparison maps are
isomorphisms in characteristic zero; one precise
primary reference is \cite[Theorem 1.1 and its
characteristic-zero specialization]{B_LandesmanBaseChange}.
Thus (B.39.2) preserves the geometric homomorphism,
without requiring $G_F\to G_K$ to be injective or open.

We give the countable descent of $h_F$ explicitly.
Put $E=\pi_1(U_F)$ and $\Delta=P_{U_F}$, and choose
a decreasing cofinal sequence of characteristic open
subgroups $N_i\subset\Delta$. Such a sequence exists
because $\Delta$ is finitely generated profinite:
intersect its finitely many subgroups of index at
most $i$. Each $N_i$ is normal in $E$.
In $E/N_i$ choose an open normal subgroup $J_i$
intersecting its finite subgroup $\Delta/N_i$
trivially. At each next stage choose $J_{i+1}$ inside
the inverse image of $J_i$. Set
\[
 M_i=(E/N_i)/J_i,\qquad
 Q_i=G_F/p(J_i),
 \qquad F_i=\overline F^{\,p(J_i)}.
 \tag{B.39.3}
\]
The finite Galois group of $F_i/F$ is $Q_i$.
Restriction of $p$ identifies $J_i$ with $p(J_i)$,
and the quotient maps give an isomorphism of
extensions
\[
 E/N_i\ \simeq\ G_F\times_{Q_i}M_i.
 \tag{B.39.4}
\]
For surjectivity, lift a prescribed element of $M_i$
to $E/N_i$ and correct its arithmetic projection by
the unique element of $J_i$ over the required
element of $p(J_i)$. Injectivity follows from
$J_i\cap(\Delta/N_i)=1$.

Equivalently, (B.39.4) is the arithmetic group of
the finite \'etale gerbe
\[
 \mathcal G_i=[\operatorname{Spec}F_i/M_i],
 \tag{B.39.5}
\]
where $M_i$ acts through $Q_i$. The quotient of $E$
and its composite with $h_F$ give actual maps
$U_F\to\mathcal G_i$ and $V_F\to\mathcal G_i$.
Concretely, these are finite $M_i$-torsors on the
two varieties, with equivariant maps to
$\operatorname{Spec}F_i$. The compatibility of
$J_i$ gives transition maps of these torsors and
gerbes. The inverse limit of the gerbes is
$\Pi^{\mathrm{\acute et}}_{U_F/F}$: it has precisely
the cofinal geometric kernels $N_i$, and its base
quotient remains $G_F$ at every stage. The exact
finite-gerbe and fundamental-group comparison is
also described in \cite[\S8 and Lemma 9.4]{B_BorneVistoli}.

Every item just listed has finite presentation:
the two varieties, the smooth compactification and
boundary of $U_F$, the finite \'etale algebras $F_i$,
the finite torsors, their group actions, transition
maps, and compatibility isomorphisms. Use finite
affine covers to describe them by equations and
maps. There are countably many items and equations.
They therefore descend together to a countable
subfield $F_0\subset F$. The gerbe formulation
does not omit the coherence data: their finitely
presented groupoid torsors and the corresponding
isomorphisms are included. The finite-presentation
limit principle used in this description is
\cite[Lemma 102.3.7, Tag 0CMX]{B_StacksLimits};
its groupoid-descent proof applies to the objects,
arrows, and their equations. All torsor properties
and equations hold after the faithfully flat
extension $F/F_0$, and hence hold over $F_0$.

We can impose a useful form on this countable
field. Let $k$ be the relative algebraic closure
of $\mathbb Q_p$ in $F$; it is a finite extension
of $\mathbb Q_p$. Choose a separating transcendence
basis $t_1,\ldots,t_d$ and one primitive element
$z$ for the finite extension $F/k(t_1,\ldots,t_d)$.
Write every coefficient of the countable data as
a rational expression in $t_1,\ldots,t_d,z$ with
coefficients in $k$. Include also the coefficients
of the minimal polynomial of $z$. The subfield
$k_0\subset k$ generated by these countably many
constants is countable, and we may take
\[
 F_0=k_0(t_1,\ldots,t_d,z)\subset F.
 \tag{B.39.6}
\]
The polynomial defining $z$ remains irreducible
over $k_0(t_1,\ldots,t_d)$ because it is irreducible
over $k(t_1,\ldots,t_d)$. Thus (B.39.6) preserves
all the algebraic relations of the descended data.

Let $U_0,V_0,\mathcal G_{i,0}$ denote these descents.
Their compatible maps give
\[
 V_0\longrightarrow\varprojlim_i\mathcal G_{i,0},
 \qquad
 \Pi^{\mathrm{\acute et}}_{U_0/F_0}
       \longrightarrow\varprojlim_i\mathcal G_{i,0}.
 \tag{B.39.7}
\]
The second map is an isomorphism. To check this,
pass to an algebraic closure of $F_0$ embedded in
an algebraic closure of $F$. Characteristic-zero
invariance identifies the geometric group of
$U_0$ with $\Delta$, and the kernels of its maps
to the geometric groups of $\mathcal G_{i,0}$
are precisely the cofinal $N_i$. Consequently
(B.39.7) is an isomorphism on geometric groups.
The arithmetic quotient is the identity of
$G_{F_0}$, so the exact fundamental sequences
give the asserted isomorphism. The first map
therefore defines an actual $h_0$ over $G_{F_0}$,
whose base change to $F$ is $h_F$, up to the usual
geometric base-path conjugacy. This construction
produces one compatible homomorphism; it does not
replace it by unrelated homomorphisms at separate
finite levels.

We next embed the particular field (B.39.6) into
one finite $p$-adic field. After inverting a
nonzero polynomial and the discriminant, it is
the fraction field of an integral algebra $A$
finite \'etale over
$k_0[t_1,\ldots,t_d,1/g]$. The same presentation
over $k$ has function field $F$. Choose a closed
point of $\operatorname{Spec}(A\otimes_{k_0}k)$;
its residue field $l$ is finite over $k$. The
\'etale map is a $p$-adic analytic isomorphism on
sufficiently small neighborhoods of that point
and its image. In the resulting closed ball
in $l^d$, choose coordinates algebraically
independent over $k_0$. This is possible by the
Baire theorem: there are countably many nonzero
polynomials over the countable field $k_0$, and
each has a closed zero set with empty interior.
The lifted point defines $A\to l$ whose kernel
contracts to zero in $k_0[t_1,\ldots,t_d,1/g]$.
Since $A$ is integral and finite over that ring,
the kernel is zero: after passage to its fraction
field the finite integral domain is a field.
Thus evaluation extends to the field embedding
$F_0\hookrightarrow l$. For $d=0$ the same
argument is simply a closed point of the finite
\'etale algebra. In particular, the embedding
preserves all the countably many retained
relations simultaneously.

Base change of $h_0$ to $l$ has exactly the original
geometric homomorphism, through (B.39.2) and
(B.39.7). The descended compactification and
boundary identify the peripheral subgroups as
well. Apply the finite-local-field trichotomy
proved in Proposition~\ref{F-local-dichotomy}.
Its proof uses the stipulated local section
theorem, with the provenance recorded in B.7.
This proves the alternatives in (B.39.1).
In the nontrivial peripheral case, return to
the original group over $K$: its geometric
image $H$ is normal in $h(\pi_1(V))$, and the
full peripheral normalizer identity gives
\[
 h(\pi_1(V))\subset
 N_{\pi_1(U)}(H)=D_c.
 \tag{B.39.8}
\]
The arithmetic projection of this image is all
of $G_K$, so the underlying cusp is rational.
If $H=1$, the quotient $\pi_1(V)\twoheadrightarrow
G_K$ gives the asserted section factorization.
Finally, a nontrivial cusp inertia has infinite
index in the hyperbolic geometric group, so
the three alternatives are disjoint.
\end{proof}

\begin{statement}{B.40}{Proved using the stipulated local theorem}{Unramification over every smooth p-adic model}
\label{st:B.40}

Let $k/\mathbb Q_p$ be finite, let $S/k$ be a smooth
geometrically connected variety, put $K=k(S)$, and
let $X\to S$ be a smooth proper family of
geometrically connected curves of genus at least two.
With $\Delta=\pi_1(X_{\overline K})$, the actual
generic-fibre comparison and restriction map are
\[
 \pi_1(X_K)\simeq
       \pi_1(X)\times_{\pi_1(S)}G_K,
 \qquad
 \operatorname{Sec}(\pi_1(X)\to\pi_1(S))
       \xrightarrow{\sim}
 \operatorname{Sec}(\pi_1(X_K)\to G_K).
 \tag{B.40.1}
\]
Both section sets use conjugation by $\Delta$.
Every relative section on the left restricts at
each $s\in S(l)$, for every finite $l/k$, to the
section of a unique point in $X_s(l)$.
The statement makes no assertion that every such
relative section is induced by a morphism $S\to X$.
\end{statement}
\begin{proof}
The centre-free fibre argument used in B.38 gives
\[
 1\longrightarrow\Delta\longrightarrow\pi_1(X)
       \longrightarrow\pi_1(S)\longrightarrow1
 \tag{B.40.2}
\]
over this arbitrary-dimensional normal base; no
$K(\pi,1)$ assumption on $S$ is required for the
centre-free argument. Comparing with the exact
sequence over $K$ proves the first isomorphism in
(B.40.1), with the identity on its kernel and quotient.
Let $h:G_K\to\pi_1(X)$ be the first component of
a generic section.

For a codimension-one point $\eta_D$ of $S$, let
$F_D=k(\eta_D)$. This is a finitely generated
extension of $\mathbb Q_p$. Completion of its
regular local discrete valuation ring gives an
equal-characteristic-zero field
$K_D\simeq F_D((t))$. Its map to $\pi_1(S)$
factors through the residue map $G_{F_D}\to\pi_1(S)$.
The exact comparison needed here is the fibre
product
\[
 \pi_1(X_D)\xrightarrow{\sim}
       \pi_1(X)\times_{\pi_1(S)}G_{F_D},
 \qquad X_D=X\times_S\operatorname{Spec}F_D.
 \tag{B.40.3}
\]
It follows by proper smooth identification of the
geometric fibre groups and the two exact sequences
over $G_{F_D}$. In particular it does not require
$G_{F_D}\to\pi_1(S)$ to be injective.
Pairing the restriction of $h$ with the residue
projection gives the actual homomorphism
\[
 G_{F_D((t))}\longrightarrow\pi_1(X_D)
                                      \quad/G_{F_D}.
 \tag{B.40.4}
\]
Identify its source with $\pi_1(\mathbb G_{m,F_D})$
as in (B.38.7), and precompose the tripod-to-$\mathbb G_m$
map. B.39 applies over $F_D$. The geometric image
is procyclic, whereas every open subgroup of the
proper hyperbolic fibre group has abelianization
of rank at least four. Properness supplies no cusp
alternative. Thus (B.40.4) kills all its inertia.

By purity of branch locus, the kernel of
$G_K\twoheadrightarrow\pi_1(S)$ is the closed
normal subgroup generated by these codimension-one
inertias. Indeed, the normalization of the regular
finite-type scheme $S$ in a finite separable
extension is finite, and being unramified in
codimension one makes it finite \'etale by
\cite[Lemma 58.21.4 and Example 58.13.3]{B_StacksPurity}.
Every finite \'etale cover has the converse
unramified property. Therefore $h$ factors uniquely
through $\pi_1(S)$. The equation over $\pi_1(S)$
makes this factor a section, and the first
isomorphism in (B.40.1) recovers the original
generic section. This proves the second isomorphism,
including its conjugacy quotient.

Finally, the same exact fibre product for a point
$s:\operatorname{Spec}l\to S$ identifies the
restricted section with a section of the proper
hyperbolic curve $X_s/l$. The stipulated local
theorem gives its unique point in $X_s(l)$.
Every smooth proper hyperbolic curve over a
finitely generated extension of $\mathbb Q_p$
spreads over a nonempty smooth model of the above
kind. The remaining generic point-realization
problem is retained; the factorization and its
fibrewise point data alone do not construct that
point or a full \'etale-topos morphism.
\end{proof}

\begin{statement}{B.41}{Proved using the stipulated local theorem}{The complete isotrivial p-adic function-field case}
\label{st:B.41}

Let $K$ be finitely generated over $\mathbb Q_p$,
and let $U/K$ be a hyperbolic curve which becomes
constant after a finite extension: there exist
$K'/K$ finite, a finite $p$-adic field $k'\subset K'$,
a hyperbolic curve $U_0/k'$, and an isomorphism
\[
 U_{K'}\simeq U_0\times_{k'}K'.
 \tag{B.41.1}
\]
For affine $U$ this hypothesis concerns the whole
open curve, including its boundary. Let
$\overline U$ be its smooth compactification and
$D=\overline U\setminus U$. Then the natural map
\[
 U(K)\ \sqcup\!\coprod_{c\in D(K)}
       \operatorname{Sec}(D_c/G_K)/I_c
       \ \xrightarrow{\ \sim\ }\
       \operatorname{Sec}(\pi_1(U)\to G_K)
 \tag{B.41.2}
\]
is bijective. Each packet is the usual torsor under
$H^1(K,\widehat{\mathbb Z}(1))=
\varprojlim_n K^\times/K^{\times n}$.
In particular every proper isotrivial hyperbolic
curve over such a field satisfies point-section
surjectivity in every transcendence degree.
\end{statement}
\begin{proof}
First let $C=C_0\times_k K$ be proper, with
$C_0/k$ of genus at least two and $k$ finite over
$\mathbb Q_p$. Enlarge $k$ to its relative
algebraic closure in $K$ and base change $C_0$;
then choose a smooth geometrically connected
$k$-model $S$ of $K$. Apply B.40 to the constant
family $S\times_k C_0$. Its arithmetic product
formula identifies a relative section with an
actual homomorphism
\[
 a:\pi_1(S)\longrightarrow\pi_1(C_0)
                                          \quad/G_k.
 \tag{B.41.3}
\]
The finite-local-field trichotomy of
Proposition~\ref{F-local-dichotomy} makes its
geometric image open or trivial, since $C_0$ is
proper. In the open case Mochizuki's Hom theorem
gives a unique dominant $k$-morphism
$S\to C_0$. Its source may have any dimension:
this is Theorem 16.5 for a smooth pro-variety,
specialized to a smooth variety, followed by the
full-profinite Remark \cite{F-Mochizuki}.
In the trivial case (B.41.3) factors through
$G_k$, and the stipulated local theorem gives
the unique constant point of $C_0(k)$ inducing
that factor. In both cases the generic value of
the resulting morphism is a point of $C(K)$.
The cartesian comparison in B.40 proves that
its section is the original generic section.

Uniqueness can be checked independently on the
same model. Two $K$-points spread to morphisms
from a common nonempty open of $S$ to $C_0$.
Equality of their generic section classes gives
equality of the induced homomorphism classes on
that model, because $G_K$ surjects onto its
fundamental group. Restriction at each closed
point gives equal point sections over a finite
extension of $k$, hence equal points by local
injectivity. Closed points are dense in this
finite-type $k$-scheme, and the equality locus
of the two morphisms is closed because $C_0$ is
separated. The morphisms, and their generic
points, are therefore equal.

We also record the geometric centralizer needed
for finite descent. For a section of this constant
proper curve, a geometric element centralizing
its image centralizes the image of (B.41.3):
the map from $G_K$ to $\pi_1(S)$ is surjective,
and the map to $\pi_1(C_0)$ is the identity on
the geometric fibre group. When (B.41.3) has
open geometric image, its geometric centralizer
is trivial by slimness of the hyperbolic curve
group. When it has trivial geometric image,
it is the section of a $k$-point and has trivial
geometric centralizer by B.10 over the finite
$p$-adic field $k$. Thus every section of the
constant proper curve over $K$ has trivial
geometric centralizer.

Now let $C/K$ be proper and satisfy (B.41.1).
Replace $K'$ by a finite Galois extension over
$K$. The result just proved gives a unique point
$x'\in C(K')$ inducing the restriction of any
given section $s$ over $K$. Its uniqueness implies
that every element of $\operatorname{Gal}(K'/K)$
fixes $x'$, since the restricted section class is
invariant under that action. Thus $x'$ descends
to $x\in C(K)$. There is still a section descent
step to check. Conjugate $s$ so that it agrees
with $s_x$ on $G_{K'}$. For $g\in G_K$, the
element $s(g)s_x(g)^{-1}$ lies in the geometric
group and centralizes $s_x(G_{K'})$, by normality
of $G_{K'}$ and equality of the restrictions.
The centralizer vanishes by the constant case
over $K'$, so $s=s_x$. This proves the proper
classification and its uniqueness for every
isotrivial curve and every finite extension of
the original field.

It remains to pass to affine $U$, allowing the
compactification to have genus zero or one.
Every geometrically connected finite \'etale
cover of an isotrivial $U$ has an isotrivial
smooth compactification. To see the precise
descent, make (B.41.1), embed an algebraic closure
of $k'$ in an algebraic closure of $K'$, and
use the equivalence of finite \'etale covering
categories given by characteristic-zero base
change \cite{B_LandesmanBaseChange}. The
geometric covering descends to $\overline{k'}$,
and then to a finite extension of $k'$ because
the cover has finite presentation. Its isomorphism
with the original cover descends to a finite
extension of $K'$. Normalization of the smooth
compactification commutes with this isomorphism.
This proves exactly the needed isotriviality of
each compactified neighborhood.

For a section of $U$, take the cofinal
neighborhoods with compact genus at least two
constructed in B.7. That construction is
geometric: a high-genus geometric cover,
a characteristic subgroup contained in its
group, and the arithmetic neighborhood
$N s(G_K)$. In each such compactification the
proper isotrivial result just proved supplies
a unique $K$-point. Under a map of neighborhoods
the points agree, because they induce the same
proper section and the proper point map is
injective. Thus they form one compatible point
in the compactified neighborhood tower. Its
image in $\overline U$ belongs to $\overline U(K)$.
An interior image gives the point section of
$U$; a boundary image gives a section in its
actual cusp decomposition group, by the same
stabilizer calculation as B.7. This proves
surjectivity in (B.41.2).

For injection and disjointness, every proper
curve used in the neighborhood proof of
\cite[Theorem 14]{B_StixCuspidal} is itself the
compactification of a finite \'etale cover of
$U$. It is isotrivial by the comparison above.
The proper injectivity just proved therefore
applies at each such level and after every
finite extension of $K$. Explicitly,
equality of sections identifies the compact
centres in every neighborhood, and a conjugator
between two sections of one packet fixes the
lifted cusp at every finite level, hence lies
in its inertia. This is the actual map from
inertia-conjugacy to geometric-conjugacy classes
in (B.41.2). Its torsor is the Kummer torsor
of B.3 over $K$, with the displayed power
completion. No compactness of $K$ or of the
finite-level multiplicative quotients is used.
\end{proof}

\begin{statement}{B.42}{Proved}{A compatible tower invisible on all closed fibres}
\label{st:B.42}

Let $k$ be a finite extension of $\mathbb Q_p$,
choose $q\in k$ with $0<|q|<1$, and let $E_q/k$
be the Tate elliptic curve. Put $B=E_q\setminus\{O\}$.
For every integer $n>1$ there is a geometrically
connected finite \'etale cyclic cover
\[
 f_n:C_n\longrightarrow B,\qquad \deg(f_n)=n,
 \tag{B.42.1}
\]
whose fibre over every closed point $b\in B$ is
the disjoint union of $n$ copies of
$\operatorname{Spec}k(b)$. Nevertheless, its
generic fibre has no $k(B)$-point, and $f_n$
has no section over any nonempty open of $B$.
The covers are compatible under divisibility
and determine a continuous surjection
\[
 \omega:\pi_1(B)\longrightarrow\widehat{\mathbb Z}
 \tag{B.42.2}
\]
which is surjective on $\pi_1(B_{\bar k})$
and vanishes on every closed-point decomposition
group. Thus the map
\[
 \operatorname{Hom}_{\mathrm{cont}}
       (\pi_1(B),\widehat{\mathbb Z})
 \longrightarrow
 \prod_{b\in B^{\mathrm{cl}}}
 \operatorname{Hom}_{\mathrm{cont}}
       (G_{k(b)},\widehat{\mathbb Z})
 \tag{B.42.3}
\]
has a nonzero kernel.
\end{statement}
\begin{proof}
Tate's uniformization gives, functorially for
every finite extension $L/k$,
\[
 E_{q^n}(L)=L^\times/q^{n\mathbb Z},
 \qquad E_q(L)=L^\times/q^{\mathbb Z}.
\]
The identity on $L^\times$ induces the
$k$-isogeny
\[
 \bar f_n:E_{q^n}\longrightarrow E_q.
\]
This is the actual algebraic morphism
$\alpha_{n,1}$ in Tate's isogeny theorem:
take $t_1=q^n$, $t_2=q$, so that
$t_2^n=t_1$. The same theorem proves its
degree is $n$; its kernel is the constant
cyclic subgroup generated by the $k$-point
$[q]\in k^\times/q^{n\mathbb Z}$
\cite[Theorem 1 and the isogeny theorem]{B_TateUniformization}.
It is \'etale because $k$ has characteristic
zero. Define
\[
 C_n=E_{q^n}\setminus\bar f_n^{-1}(O),
 \qquad f_n=\bar f_n|_{C_n}.
\]
This is a geometrically connected open in
an elliptic curve, and (B.42.1) follows.

For a closed point $b$ put $L=k(b)$.
Represent its image in $E_q(L)$ by
$a\in L^\times$. Its complete inverse image
under $\bar f_n$ consists of
\[
 [a], [aq],\ldots,[aq^{n-1}]
             \quad\text{in }L^\times/q^{n\mathbb Z}.
\]
These are $n$ distinct $L$-points, all in
$C_n$ because $b\ne O$. This proves the
asserted equality of finite \'etale fibre
schemes, not merely nonemptiness.
On the other hand $C_n$ is integral and its
function field is a degree-$n$ extension
of $k(B)$. A point of its generic fibre would
give a $k(B)$-homomorphism from that extension
to $k(B)$. Such a homomorphism is injective
and is the identity on $k(B)$, which is
impossible for degree $n>1$. A section on
any nonempty open would give that generic
point, so no such section exists.

For $n\mid m$ the identity on the
uniformizing multiplicative group gives
$C_m\to C_n$, and translation by $q$ maps
to translation by $q$. Choose a point
$b_0\in B(k)$ represented by
$a_0\in k^\times\setminus q^{\mathbb Z}$,
and the compatible lifts represented by
this same $a_0$ on every $C_n$.
The deck quotients are then the compatible
maps $\pi_1(B)\to\mathbb Z/n\mathbb Z$.
Their restrictions to the geometric group
are onto because every $C_n$ is
geometrically connected. Their inverse
limit is (B.42.2); its image is compact
and maps onto every $\mathbb Z/n\mathbb Z$,
so both it and the geometric image are
all of $\widehat{\mathbb Z}$.
Every closed-point decomposition group
acts trivially on each fibre just computed,
so its image in every deck quotient,
and hence under $\omega$, is zero.
This proves (B.42.3). Tate also records
this same completely split infinite tower
explicitly \cite[p.~18 of the author manuscript]{B_TateUniformization}.
\end{proof}

\paragraph{Exact remaining realization step.}
The finite-cover local-point hypothesis in
\cite[Definition 0.2(vi)]{B_SaidiFunctionFields}
requires a generic point whenever all the
closed fibres of a finite \'etale curve cover
have residue-field points. B.42 disproves
that hypothesis for a finite $p$-adic base
field, even for a hyperbolic base curve.
Consequently it cannot be inserted into
the realization step after B.40.
The maps in B.42 are quotient maps of
fundamental groups; they are not sections
of a proper hyperbolic family. No
nongeometric section of such a family is
constructed by this example.
The unresolved assertion is still whether
the particular relative sections produced
in B.40 all arise from algebraic sections
for nonisotrivial families. B.41 resolves
that assertion in the stated isotrivial
range only.

% ---- C ----
\section{Real structures, fixed groups, and the moduli comparison}
\label{sec:C}
Here $\Pi=\pi_1^{\mathrm{top}}(U(\mathbb C))$ is discrete,
$\widehat\Pi=\pi_1(U_{\mathbb C})$ is profinite, and
$c$ denotes the nonidentity element of $G_{\mathbb R}=C_2$.
The topological extension $E^{\mathrm{top}}$ consists of
lifts of the identity and complex conjugation to the universal
cover. Finite equivariant covering comparison identifies its
profinite completion with $E^{\mathrm{et}}=\pi_1(U_{\mathbb R})$.
Section classes are taken modulo $\Pi$ or $\widehat\Pi$,
respectively. The maps to automorphism-group quotients are
specified separately.

\subsection{The two classification clauses of Lead 3}

\begin{statement}{C.1}{Known}{the corrected surface count}
\label{st:C.1}
\label{C-classification}
For a compact connected oriented surface $X$ of genus $g$,
orientation-reversing involutions are classified, up to
orientation-preserving homeomorphism, by the quotient
surface $Y=X/\langle c\rangle$. The nonorientable quotients
have nonorientable genus $q\geq1$ and $r\geq0$ boundary
components, with $g=q+r-1$; the orientable quotients have
genus $h\geq0$ and $r\geq1$, with $g=2h+r-1$.
There are $g+2+\lfloor g/2\rfloor$ types in all.
The $g+1$ types counted in the original text of scan 439
are exactly the nonorientable-quotient types.
Its margin supplies the remaining types.
\end{statement}
\begin{proof}
A collar of a fixed circle has reflection coordinates, so the
fixed circles give precisely the boundary components of the
quotient. Away from them the quotient map has degree two.
The oriented double of a nonorientable bordered surface,
and the double of an orientable bordered surface along its
nonempty boundary, recover $X$ with its involution.
These constructions are inverse on equivariant homeomorphism
classes; they type the morphism implicit on scan 437.
Euler characteristics give $2-2g=2(2-q-r)$ and
$2-2g=2(2-2h-r)$ respectively. The classification of compact
surfaces gives the stated ranges and counts.
These are also the restrictions of Gross--Harris,
Proposition 3.1, with the realization statement immediately
following it~\cite{GrossHarris1981}.
Grothendieck's $\gamma$ counts additional crosscaps after
the initial projective plane, so $q=\gamma+1$; his relation
$\gamma+r=g$ is the first formula. The margin on
folio 222, scan 439, is mathematically necessary.
\end{proof}

\begin{statement}{C.2}{Refuted}{abelian subgroups and strips}
\label{st:C.2}
\label{C-abelian}
For the orientation-preserving mapping class group of a
surface of hyperbolic type $(g,n)$, every free abelian
subgroup has rank at most $3g-3+n$. A pants decomposition
gives that rank through its Dehn twists.
The assertion that all maximal free abelian subgroups
are generated by strips is false.
\end{statement}
\begin{proof}
The bound is Birman--Lubotzky--McCarthy, Theorem A~\cite{BLM1983}.
The twists around the $3g-3+n$ pants curves commute.
Intersection arcs crossing one cuff and missing the others
show independence of their powers.
For the proposed identification on scan 470, use the
once-punctured torus and its homological identification
with $\mathrm{SL}_2(\mathbb Z)$. The matrix
\[
 A=\begin{pmatrix}2&1\\1&1\end{pmatrix}
\]
has characteristic polynomial $t^2-3t+1$. Its rational
commutant is $\mathbb Q[A]$, a real quadratic field.
The determinant-one integral units in this order have
rank one, so a maximal free abelian subgroup containing
$A$ has rank one. Every nontrivial Dehn twist has
unipotent homological action. No nonidentity unipotent
matrix commutes with $A$: its unique rational eigenline
would have to be one of the two irrational eigenlines
of $A$. Thus this subgroup contains no strip twist.
The characteristic polynomial and commutant equations
are checked in the exact program. This gives an
explicit homological comparison for the claim.
\end{proof}

\subsection{C(a): section classes and profinite centralizers}

\begin{statement}{C.3}{Known}{discrete real sections}
\label{st:C.3}
\label{C-discrete}
For every hyperbolic real curve $U$, the map
\[
 \pi_0(U(\mathbb R))\longrightarrow
 \operatorname{Sect}(E^{\mathrm{top}}\to C_2)/\Pi
\]
is bijective. For the lift $\sigma_I$ attached to a component $I$,
\[
 \Pi^{\sigma_I}=
 \begin{cases}
  1,& I\text{ is an interval},\\
  \langle a_I\rangle\simeq\mathbb Z,&I\text{ is an oval}.
 \end{cases}
\]
In the second case $a_I$ is the oval loop; choosing its
orientation accounts for $a_I\leftrightarrow a_I^{-1}$.
\end{statement}
\begin{proof}
The hyperbolic universal cover is the disc.
An order-two lift of complex conjugation is an
orientation-reversing hyperbolic isometry, hence reflection
in one geodesic. Its fixed geodesic is connected and projects
onto one real component. The lift determined by a real
point fixes a point over it. All points in one component
therefore give the same conjugacy class, and one geodesic
cannot project into two components. This proves the bijection.
Deck transformations commuting with the reflection preserve
its fixed geodesic. Their restriction is the deck group of
the covering of that real component: trivial for an interval
and infinite cyclic for an oval. Restriction is injective
since a deck transformation fixing a point is the identity.
This proves the surjectivity in (14) as well as (6)--(10),
\S33, folios 223--225, scans 441--445. It is the
specialization of Mochizuki's Theorem 3.6 and
Corollary 3.7~\cite{MochizukiTopics}.
\end{proof}

\begin{statement}{C.4}{Proved}{the profinite fixed group in every genus}
\label{st:C.4}
\label{C-profinite-centralizer}
For every hyperbolic real curve $U$, proper or affine,
completion gives a bijection of the classes in
Theorem~\ref{C-discrete} with
$\operatorname{Sect}(\pi_1(U)\to C_2)/\widehat\Pi$.
For every component $I$,
\[
 \widehat\Pi^{\sigma_I}=
 \begin{cases}
  1,&I\text{ is an interval},\\
  \overline{\langle a_I\rangle}\simeq\widehat{\mathbb Z},
       &I\text{ is an oval}.
 \end{cases}
\]
No bound on the genus of $U$ is required.
\end{statement}
\begin{proof}
We verify the hypotheses of Boggi--Zalesskii, Theorem A,
for $E^{\mathrm{top}}$ and the prime two~\cite{BoggiZalesskii2024}.
Its index-two subgroup $\Pi$ is a free group of finite rank
or a closed orientable surface group. These groups are good.
For the required cohomological verification, first trivialize
a finite coefficient module on a finite cover. A degree-one
class is killed by the kernel of its finite-valued cocycle.
For a closed surface a degree-two class is killed on a cover
of degree divisible by the exponent of the coefficient
module: restriction on the fundamental class multiplies
by the degree. Higher cohomology vanishes. For a free
group only degree one occurs. The finite-cover criterion
for goodness applies. Comparing the two
Lyndon--Hochschild--Serre spectral sequences for the
index-two extension proves goodness of $E^{\mathrm{top}}$.
Finite coefficients have finite cohomology in every degree,
and the finite-index surface subgroup has cohomological
dimension at most two. Thus finite virtual
two-cohomological type holds.
Residual finiteness of the surface or free subgroup,
followed by the intersection of the finitely many
translates of a separating normal subgroup, proves
residual finiteness of the extension. For an arbitrary
discrete continuous $\mathbb F_2[\widehat E]$-module,
take the union of its finite submodules. The finite-type
resolutions make cohomology commute with this filtered
union, as required in the source's statement.

The cited theorem gives bijectivity for conjugacy classes
of order-two subgroups and identifies their completed
centralizers with closures of discrete centralizers.
The geometric group is torsion-free by finite
cohomological dimension. An $E$-conjugator between
two sections can be made geometric by multiplication
by the source involution. We therefore obtain exactly
the required conjugacy convention and
$C_{\widehat E}(\sigma_I)=\overline{C_E(\sigma_I)}$.
Intersect with the open subgroup $\widehat\Pi$ and use
Theorem~\ref{C-discrete}.

We check the full topology on an oval subgroup.
For each $m$, the essential simple loop $a_I$
has order divisible by $m$ in a finite quotient of $\Pi$.
For a nonseparating loop use homology modulo $m$.
For a separating loop cut along it. On a positive-genus
side assign one handle pair the permutations
$A=(1\,2\,\cdots\,m)$ and the interchange $B$ of
two blocks of size $m$. Then $[A,B]$ is an
$m$-cycle on the first block times the inverse cycle
on the second; it has order $m$. Use its inverse
on the other side. On a genus-zero side there is
another puncture; assign its loop that permutation
and the other loops the identity. These representations
glue by the boundary relation. An essential oval
has no unpunctured disc side. Thus every modulus
is detected, and its closure is $\widehat{\mathbb Z}$.
\end{proof}

\begin{statement}{C.5}{Refuted}{the two signs at a real cusp}
\label{st:C.5}
\label{C-cusp-packet}
For a real cusp $d$ with decomposition extension
$1\to\widehat{\mathbb Z}(1)\to D_d\to C_2\to1$,
the local section-class set is a torsor under
$\mathbb Z/2$. The map into the global section set
sends its two classes to the components adjacent
to the two germs at $d$. These components can coincide.
The assertion of two distinct global classes for
every real cusp is false.
\end{statement}
\begin{proof}
Conjugation acts by $-1$ on $\widehat{\mathbb Z}(1)$, so
\[
 H^1(C_2,\widehat{\mathbb Z}(1))
 =\widehat{\mathbb Z}/2\widehat{\mathbb Z}=\mathbb Z/2.
\]
The real Kummer map is the sign map. The two parameter
signs approach the two real germs; transport to a nearby
point identifies the tangential section with its component
class by Theorem~\ref{C-profinite-centralizer}.

Take $E:y^2=x^3+x$ and $U=E\setminus\{O\}$.
The cubic has one real root, so $E(\mathbb R)$ is one
oval. Deleting $O$ leaves one interval. The two
elements of $\operatorname{Sect}(D_O\to C_2)/
\widehat{\mathbb Z}(1)$ therefore have the same
image in $\operatorname{Sect}(\pi_1(U)\to C_2)/
\pi_1(U_{\mathbb C})$. This is the exact map and
fibre giving the counterexample.
\end{proof}

\begin{statement}{C.6}{Proved}{the incidence graph recovers the real boundary order}
\label{st:C.6}
\label{C-boundary-order}
The arithmetic extension of a hyperbolic real curve,
together with its peripheral conjugacy classes,
recovers the real cusps, the component of $X(\mathbb R)$
containing each, and their cyclic order up to reversal.
This includes an oval carrying just one real cusp.
\end{statement}
\begin{proof}
A cusp is real exactly when its inertia conjugacy class
is fixed by $C_2$; other cusps form two-element orbits.
For each real cusp take its two-element local section
torsor and its map to the global section set.
Make an unoriented edge for the cusp, whose ends are
the two image classes. Interval vertices are also
recognized by their trivial fixed group in
Theorem~\ref{C-profinite-centralizer}.
Every such interval has two ends, so the resulting
graph is a disjoint union of cycles. A one-edge loop
occurs precisely when the two local signs have the
same global image. Traversing a cycle lists the
cusps of one compact oval in their cyclic order,
with reversal as its ambiguity. Vertices with
fixed group $\widehat{\mathbb Z}$ are the ovals
with no real cusp. Nonreal pairs give no graph ends.
This answers the margin of \S33, folio 221, scan 437.
\end{proof}

\subsection{C(b)--(c): arithmetic conjugations and the tripod}

\begin{statement}{C.7}{Proved}{prediscretizations and archimedean places}
\label{st:C.7}
\label{C-prediscretizations}
For the arithmetic prediscretization torsor $P_a$
of scans 377--378, (29)--(30) are the maps
\[
 e\longmapsto e^{-1}ce,\qquad
 P_a/\{1,c\}\longrightarrow
 \{\text{complex conjugations in }\Gamma_a\}.
\]
The second is a bijection. Under the faithful
identification $\Gamma_a\simeq G_{\mathbb Q}$, its
image is every element of order two and it
identifies $P_a/\{\pm1\}$ with the archimedean
places of $\overline{\mathbb Q}$.
The cyclotomic-kernel orbits on this set are
$\widehat{\mathbb Z}^{\,\times}/\{\pm1\}$.
\end{statement}
\begin{proof}
The quotient construction of LM3 A.8 and B.6 makes
the oriented arithmetic prediscretizations a torsor
under the arithmetic image. Choosing a marking
identifies it with embeddings of the fixed algebraic
closure into $\mathbb C$. Pulling back conjugation
gives the displayed involution. Reversing orientation
composes the embedding with $c$ and has the same
image. Before faithfulness this gives exactly the
conjugates of $c$ in that image.

We compute the stabilizer. An automorphism of
$\overline{\mathbb Q}$ commuting with $c$ preserves
its fixed field $R_{\rm alg}$ of real algebraic
numbers. Positivity is the set of nonzero squares,
so it preserves the order. Since it fixes
$\mathbb Q$ and each real algebraic number is
determined by its rational cut, it is the
identity on $R_{\rm alg}$. As
$\overline{\mathbb Q}=R_{\rm alg}(i)$, it is
$1$ or $c$. Thus $C_{G_{\mathbb Q}}(c)=\{1,c\}$.
An arbitrary involution has a real closed fixed
field by Artin--Schreier. Its ordering gives an
embedding into $\mathbb R$; extension to the
algebraic closure conjugates it to $c$.
Embeddings modulo conjugation are the archimedean
places. Hoshi--Mochizuki, Corollary 6.4, supplies
the faithful arithmetic image used here; this
quotient argument needs no assertion about the
geometric congruence kernel.
Finally $\chi$ is onto and $\chi(c)=-1$.
The double-coset calculation gives the last formula.
\end{proof}

\begin{statement}{C.8}{Proved}{restriction and the six tripod germs}
\label{st:C.8}
\label{C-tripod}
Restriction of a number-field section to a real
decomposition group is its real component section.
For $B=\mathbb P^1_{\mathbb Q}\setminus\{0,1,\infty\}$,
choose the unit tangential base point at $0$ toward $1$.
There are free profinite generators $x,y$ with
$c(x)=x^{-1}$ and $c(y)=y^{-1}$, and the three
real section involutions are
\[
 \tau_{(0,1)}=c,\qquad
 \tau_{(-\infty,0)}=xc,\qquad
 \tau_{(1,\infty)}=y^{-1}c .
\]
The cusp germs have the following global images:
\[
\begin{array}{c|c|c}
\text{cusp}&\text{local parameter sign}&\text{involution}\\ \hline
0&t>0&c\\
0&t<0&xc\\
1&t-1<0&c\\
1&t-1>0&y^{-1}c\\
\infty&1/t>0&y^{-1}c\\
\infty&1/t<0&xc
\end{array}
\]
These are actual equalities after the stated
path choices; forgetting them gives geometric
conjugacy classes.
\end{statement}
\begin{proof}
The base-change square of arithmetic fundamental
groups over $K$ and $K_v=\mathbb R$ is a pullback
over $G_K$ and $G_{\mathbb R}$. Pulling back the
section and applying
Theorem~\ref{C-profinite-centralizer} proves compatibility.
For the formulas use the real half-plane presentation
\[
 E^{\mathrm{top}}
 =\langle a,b,d\mid a^2=b^2=d^2=1\rangle,\quad
 c=a,\quad x=ba,\quad y=ad .
\]
The geometric subgroup is the kernel sending all
three reflections to $c$. The reflected sides
project to the three intervals. The third
peripheral generator is $z=(xy)^{-1}=db$.
The reflections at infinity are $d$ and
$z^{-1}d=b$; at the other cusps they are
$a,b$ and $a,d$. This proves the six entries.
The program checks the free-product identities.
The completion bijection shows the three
classes remain distinct profinitely.
\end{proof}

\begin{statement}{C.9}{Proved}{an explicit class-two real obstruction}
\label{st:C.9}
\label{C-tripod-class-two}
The abelianized real tripod has four section classes
in $(\mathbb Z/2)^2$. Exactly
$(0,0),(1,0),(0,1)$ lift fully. The fourth already
fails in a finite nilpotent quotient of class two.
\end{statement}
\begin{proof}
Inversion gives $H^1(C_2,\widehat{\mathbb Z}(1)^2)
=(\mathbb Z/2)^2$, and Proposition~\ref{C-tripod}
gives the first three classes. For the fourth use
\[
 F_2\longrightarrow
 D_8=\langle r,t\mid r^4=t^2=1,\ trt=r^{-1}\rangle,
 \qquad x\mapsto t,\quad y\mapsto rt .
\]
The real action becomes trivial since the images
of $x,y$ are involutions. Thus the extension is
$D_8\times C_2$. The geometric coordinates above
the fourth abelian class are $r$ and $r^3$;
both square to $r^2\ne1$, preventing a section.
The program exhaustively checks that coset.
The quotient maps to this class-two extension
and the abelian extension specify the comparison.
\end{proof}

\subsection{C(d): the meaning of the component counts}

\begin{statement}{C.10}{Proved}{effective lifts and the centralizer index}
\label{st:C.10}
\label{C-moduli-count}
Fix an orientation-reversing mapping class $\sigma$
of a compact genus-$g$ surface, $g\geq2$, with
$r$ real ovals. The fibre over $\sigma$ of
order-two lifts in $\operatorname{Aut}(\Pi_g)$,
modulo inner automorphisms, has $r$ elements.
The centralizer of $\sigma$ in the
orientation-preserving outer mapping class
group acts transitively on this fibre.
For one lift $u$ its stabilizer fits into
\[
\begin{split}
 1\longrightarrow\Pi_g^u
 &\longrightarrow C_{\operatorname{Aut}^{+}(\Pi_g)}(u)\\
 &\longrightarrow
 \operatorname{Stab}_{C_{\operatorname{Out}^{+}(\Pi_g)}(\sigma)}
       ([u])\longrightarrow1.
\end{split}
\]
The last group has index $r$ in the outer centralizer
and $\Pi_g^u\simeq\mathbb Z$.
Thus the count $r$ on scans 445--448 is the size of
this fibre before the centralizer acts.
The unmarked real moduli stack has
$g+2+\lfloor g/2\rfloor$ components; with one
labelled real marked point it has
$g+1+\lfloor g/2\rfloor$.
\end{statement}
\begin{proof}
Apply Theorem~\ref{C-discrete} to a representative
of $\sigma$. Its ovals give exactly the effective
lifts and their fixed groups. Homeomorphisms of the
quotient surface permute its boundary circles
transitively. Their orientation-preserving lifts
to the double centralize $\sigma$, proving
transitivity on the lift classes.
An outer centralizer element stabilizes $[u]$
precisely when a representative can be modified
by an inner automorphism to commute with $u$.
The kernel consists of inner automorphisms by
$\Pi_g^u$. This proves the sequence and the
index. The index assertion on scan 447 must
therefore be applied to the displayed image,
with its infinite cyclic kernel retained.

Mochizuki's Corollary 3.8 identifies real moduli
components with these orientation-reversing
mapping-class conjugacy classes.
Proposition~\ref{C-classification} gives the
unmarked count. One real marked point requires
$r>0$, and transitivity gives one component
for each such type. Exactly the fixed-point-free
type is removed. This quotient of the $r$ lift
classes to one centralizer orbit is the
precise comparison needed for scans 445--448.
\end{proof}

\begin{statement}{C.11}{Proved}{a finite formula for labelled real markings}
\label{st:C.11}
\label{C-marked-count}
For $g\geq2$, $n>0$, the components of the real
stack with $n$ individually labelled real marked
points are classified by an allowed type $(r,a)$
of Proposition~\ref{C-classification}, $r>0$,
and the following data.
Partition $\{1,\dots,n\}$ into at most $r$
nonempty blocks and put a cyclic order on each.
For a nonorientable quotient ($a=1$) identify
each cyclic order with its reverse independently;
for an orientable quotient ($a=0$) identify
the simultaneous reversal of all orders.
For a partition of block sizes $m_1,\dots,m_k$
the contributions are
\[
 \prod_{j=1}^k b(m_j)\quad(a=1),\qquad
 \frac{\prod_{j=1}^k(m_j-1)!+
              {\bf1}_{\max m_j\leq2}}2\quad(a=0),
\]
where $b(1)=b(2)=1$ and $b(m)=(m-1)!/2$ for $m\geq3$.
Sum over the allowed types and partitions.
\end{statement}
\begin{proof}
Marked real points lie on boundary circles of
the quotient. Boundary permutations account
for the blocks, with cyclic orders within
them. All boundary permutations are realized.
On an orientable quotient a homeomorphism
preserves all induced boundary orientations,
or reverses all of them. On a nonorientable
quotient, moving a small boundary disc along
a one-sided loop and back reverses that
disc's orientation and fixes the other
boundary circles. These moves give precisely
the independent reversals.

Circle configuration spaces with prescribed
cyclic order are connected, as are the real
Teichmüller loci for a fixed involution.
Thus these homeomorphism orbits give exactly
the connected components.
A labelled cyclic order on $m$ points has
$(m-1)!$ possibilities; reversal fixes one
precisely for $m\leq2$. Burnside's formula
gives both counts.
For the manuscript's unordered punctures,
the complete topological data instead
allow a finite subset in the boundary
and interior of the quotient; interior
points lift to conjugate pairs. This is
the double construction of scan 437.
The finite permutation quotient of the
labelled moduli stack is the comparison
map. The numerical formula here specifies
its individually real labelled objects.
\end{proof}

\begin{statement}{C.12}{Open}{full profinite moduli in genus at least three}
\label{st:C.12}
\label{C-high-genus-open}
The classification of all involutions over complex
conjugation in the full profinite extended mapping
class group, with complete centralizers, is not
established here for $g\geq3$. The smallest
uncompleted type in this calculation is $(3,0)$.
The exact fibre calculation below records the
obstruction reached by the attempted routes.
\end{statement}
\begin{proof}[Established part and attempted routes]
Mochizuki's Corollary 3.15 proves the full
statement for $g\leq2$. The goodness proof
in Theorem~\ref{C-profinite-centralizer}
uses a free or surface subgroup of finite
index. The group here is connected to
a surface group by the point-pushing
map in the Birman exact sequence.
Induction along this sequence changes
the number of markings and retains
the genus, so it does not give the
missing genus-three cohomology
comparison.

We checked Holzschuh's Theorems A and
B~\cite{Holzschuh2025}. Theorem A
concerns the geometrically pro-two
étale homotopy type of an equivariantly
triangulable scheme; Theorem B assumes
geometric étale nilpotence for the
full homotopy type. The comparison
relevant to a level cover $X$ is
\[
 \operatorname{Et}(X)\longrightarrow
 B\pi_1^{\mathrm{et}}(X).
\]
Lifting a section on the right
requires the successive Postnikov
obstructions to vanish. Their
targets are
$H^{q+1}(C_2,\pi_q(\operatorname{Et}(X_{\mathbb C})))$,
with the action furnished by the
partial lift, for $q\geq2$.
The source does not compute these
obstructions for a level cover
of $\mathcal M_3$. The need to
retain higher homotopy is visible
already from the unmarked group:
the genus-at-least-three mapping
class group is perfect, so its
maximal pro-two quotient is trivial.

For the congruence route write
\[
 1\to K_{g,n}\to\widehat\Gamma^{\,\pm}_{g,n}
       \to\check\Gamma^{\,\pm}_{g,n}\to1 .
\]
Fix a topological involution $\tau$.
The profinite lifts of its congruence
image, modulo $K_{g,n}$-conjugacy,
form exactly $H^1(C_2,K_{g,n})$
for conjugation by $\tau$.
Indeed $k\tau$ squares to one
exactly when $k\,\tau(k)=1$,
and conjugation gives the
nonabelian coboundary relation.
The kernel on centralizers is
$K_{g,n}^{\tau}$.
We have not established the required
orbit calculation or the fixed
kernel for $(g,n)=(3,0)$.

The symmetric congruence theorem
of Boggi~\cite{Boggi2026} applies
to centralizers of finite
orientation-preserving actions
with quotient genus at most two.
We checked its actual map
$\widehat{\Gamma(S)^G}\to
\operatorname{Out}(\widehat\Pi)$.
Here the reflection lies in
the orientation extension, and
the required centralizer is
inside the full completed group.
The displayed $H^1$ and fixed
kernel remain to be determined.
A next concrete step is the
two-cohomology comparison for
a genus-three level subgroup,
or a direct calculation of
these two kernel invariants.
This record asserts no
counterexample to the
genus-three statement.
\end{proof}



% ---- C_extra ----
\subsection{C(d), continued: injectivity and the exact congruence fibres}

\begin{statement}{C.13}{Proved}{real topological types remain distinct in the congruence image}
\label{st:C.13}
\label{C-real-type-injectivity}
Let $g\geq2$, and let $\tau$ and $\tau'$ be orientation-reversing
involutions of a closed oriented genus-$g$ surface. Conjugacy of their
outer actions on the profinite surface group implies conjugacy of their
mapping classes by an orientation-preserving mapping class.
Consequently the map from the real components of $\mathcal M_g$ to
conjugacy classes of involutions in the procongruence extended mapping
class group is injective. The same map to the full profinite extended
mapping class group is injective.

In genus three this distinguishes all six topological types: the
nondividing types with $r=0,1,2,3$ ovals and the dividing types with
$r=2,4$ ovals. The assertion concerns the images of these six types;
it does not classify the remaining profinite involutions.
\end{statement}
\begin{proof}
Put $\Delta=\widehat{\pi_1(S_g)}$. This profinite surface group is
centre-free. We specify the extension recovered from an outer action
$\alpha:C_2\to\operatorname{Out}(\Delta)$:
\[
 E_\alpha=\operatorname{Aut}(\Delta)
       \mathbin{\times}_{\operatorname{Out}(\Delta)} C_2.
\]
For the arithmetic fundamental group $E$ of a real curve realizing
$\tau$, the map
\[
 E\longrightarrow E_\alpha,\qquad
 e\longmapsto(\operatorname{Ad}(e)|_\Delta,\overline e)
\]
is an isomorphism. Injectivity follows from the trivial centre of
$\Delta$. For surjectivity, choose any element of $E$ over the
specified element of $C_2$; its automorphism differs from the given
one by a unique inner automorphism. Multiplication by that element
of $\Delta$ supplies the preimage. This argument also applies to a
nonsplit real extension.

Choose an automorphism $b$ representing an outer conjugator between
$\alpha$ and $\alpha'$. It induces the extension isomorphism
\[
 E_\alpha\longrightarrow E_{\alpha'},\qquad
 (\varphi,\epsilon)\longmapsto
       (b\varphi b^{-1},\epsilon).
\]
Theorem~\ref{C-profinite-centralizer} therefore recovers the number
$r$ of real ovals as the number of geometric conjugacy classes of
sections. It recovers more than this number. For each section class
choose a representative $s_i$, and form
\[
 H_i=C_\Delta(s_i)\simeq\widehat{\mathbb Z},\qquad
 w_i=\operatorname{im}\bigl(
     1\in H_i/2H_i\longrightarrow H_1(\Delta,\mathbb F_2)\bigr).
\]
Here $1$ means the unique nonzero element of the group of order two,
so no choice of generator occurs. Changing the representative
conjugates the subgroup and leaves its image in homology unchanged.
Theorem~\ref{C-profinite-centralizer} identifies $w_i$ with the
modulo-two fundamental class of the corresponding oval $A_i$.
The extension isomorphism carries the unordered collection of these
classes to the corresponding collection for $\tau'$.

For $r>0$, the dividing invariant has the intrinsic expression
\[
 X(\mathbb C)\setminus X(\mathbb R)\text{ is disconnected}
 \quad\Longleftrightarrow\quad
 \sum_{i=1}^r w_i=0
       \quad\text{in }H_1(\Delta,\mathbb F_2).
\tag{C.13.1}
\]
Here is a direct topological proof. The complement of the oval union
$A$ is the orientation double cover of the interior of the connected
quotient surface, and thus has one or two components. With two
components, the closure of either component is a modulo-two
two-chain whose boundary is $A$. With one component, triangulate the
surface with $A$ as a subcomplex. A two-chain whose boundary is
supported on $A$ must have constant coefficient on all triangles
joined across edges outside $A$. Connectedness of the complement
forces one constant coefficient. Across every edge of $A$ the two
adjacent coefficients are then equal, so its boundary is zero.
It cannot have boundary $A$. This proves (C.13.1). When $r=0$ the
complement is connected and the type is already determined.

The extension isomorphism preserves both $r$ and (C.13.1).
The quotient-surface classification of
Proposition~\ref{C-classification}, or equivalently Gross--Harris,
Proposition 3.1 and its following realization statement
\cite{GrossHarris1981}, now gives an equivariant homeomorphism.
If it reverses orientation, compose it with the target involution
to obtain an orientation-preserving conjugator. Conjugacy in the
procongruence group is in particular conjugacy of these outer
actions. Conjugacy in the full completion maps to conjugacy in the
procongruence group, proving both injectivity assertions. Substituting
$g=3$ in the two quotient-surface formulas gives the six listed types.
\end{proof}


\begin{statement}{C.14}{Proved}{the marked injectivity comparison, including relative boundary orientation}
\label{st:C.14}
\label{C-marked-injectivity}
For $g\geq2$, distinct real components of the moduli stack with
individually labelled real markings remain distinct in the
procongruence extended mapping class group and in its full profinite
completion. The analogous injectivity holds for unordered markings,
including conjugate pairs of nonreal markings.

More precisely, an isomorphism of the real arithmetic fundamental
group extensions which preserves the peripheral conjugacy classes
recovers the compact type $(g,r,a)$ and the real boundary incidence
data of Proposition~\ref{C-boundary-order}. In the dividing case,
the recovered cyclic orders on labelled marked ovals are identified
only by simultaneous reversal, as in
Proposition~\ref{C-marked-count}.
\end{statement}
\begin{proof}
Write $U=X\setminus D$. Filling the punctures is the actual quotient
\[
 \pi_1(U_{\mathbb C})\longrightarrow\pi_1(X_{\mathbb C})
 =\pi_1(U_{\mathbb C})/
       \overline{\langle\!\langle I_d:d\in D\rangle\!\rangle}.
\tag{C.14.1}
\]
The corresponding quotient of the arithmetic extension recovers
the real extension of $X$. Hence the given isomorphism induces the
comparison in Theorem~\ref{C-real-type-injectivity}. In particular it
recovers $r$ and $a$. Proposition~\ref{C-boundary-order} recovers
the real punctures, their containing ovals, and the cyclic order
on each marked oval up to its separate reversal. Its graph is
understood as a graph with flags: the two flags of a cusp are the
two distinct elements of its local section torsor, even when they
are incident with the same vertex. Thus the one-edge loop and the
two-edge cycle retain their full incidence information.

We prove the additional orientation compatibility using the maps
of these graphs for finite covers. This supplies the specific
comparison between a discrete oval generator and its profinite
centralizer. Fix a marked oval $A$, an adjacent real interval,
and its section $\widetilde s$ of the extension for $U$. Let $s$
be its image in the extension for $X$, and put
\[
 P=\pi_1(X_{\mathbb C}),\qquad H=C_P(s)
       =\overline{\langle a\rangle}\simeq\widehat{\mathbb Z},
\]
where $a$ is one oriented circuit of the actual oval.
For an open normal subgroup $N$ of $P$ invariant under $s$, the
subgroup $N s(C_2)$ defines a finite real \'etale cover $X_N\to X$
with the chosen lift of the base real point. The section $s$ in
that subgroup selects the real oval $A_N$ through this lift.
The inverse image $U_N$ of $U$ has precisely the preimages of $D$
as punctures, and the section $\widetilde s$ selects the base
interval of the incidence cycle on $A_N$.

We record the fibre and monodromy identifications explicitly.
With the chosen sheet denoted by $1$, the complex fibre of this
cover is $P/N$. Its real points are the fixed cosets for $s$.
The fibre on the selected connected real oval is the monodromy
orbit of $1$ under the discrete oval loop:
\[
 \langle a\rangle N/N=HN/N\simeq H/(H\cap N).
\tag{C.14.2}
\]
The first equality follows by lifting paths along the real circle;
the second follows by taking the closure in this finite fibre.
In particular the real oval cover has degree
$[H:H\cap N]$. For $h\in H$, normality of $N$ and the equality
$hs=sh$ show that $h$ normalizes $Ns(C_2)$, giving an actual
real deck transformation. Its action on this selected component
agrees with (C.14.2). Equivalently,
its inverse limit is the centralizer: the inverse limit of all
fixed-coset sets is $P^s=H$, and (C.14.2) identifies the chosen
monodromy orbits with its finite cyclic quotients. These quotients
are cofinal because Theorem~\ref{C-profinite-centralizer} proves
the full profinite topology on $\langle a\rangle$.

The graph construction is functorial for these covers, not merely
for their individual graphs. A cusp maps to its image cusp; each
of its two local section classes maps to the corresponding local
class; a global section class maps by composition with the covering
homomorphism. These maps commute with the local-to-global incidence
maps. Since $X_N\to X$ is \'etale also at the filled punctures,
the result on the selected cycles is the degree-$[H:H\cap N]$
cover of the base cycle, with its flags and its selected sheet.
The deck actions in (C.14.2) act on these graph covers. The quotient
map (C.14.1) ensures that a peripheral extension isomorphism sends
exactly this system of compact \'etale covers, and their inverse
images in $U$, to the corresponding system on the other curve.
It therefore gives compatible isomorphisms of the graph-cover
maps and their deck actions.

Choose orientations of the source and target base cycles. The
base graph isomorphism has degree $\epsilon\in\{1,-1\}$. Lifting
one positively oriented circuit in every graph cover now gives
\[
 b(a)\equiv (a')^{\epsilon}\pmod {H'\cap b(N)}.
\]
The same sign occurs at every level because each graph-cover
isomorphism lies over the same base graph isomorphism. Cofinality
gives the equality $b(a)=(a')^{\epsilon}$ in $H'$, after the
conjugations used to identify the chosen sections. An arbitrary
unit of $\widehat{\mathbb Z}$ is thus excluded by the actual
graph-cover maps. On homology the equality is independent of
those conjugations.

Suppose now that $X$ is dividing. Orient its $r$ ovals as the
boundary of one half of $X(\mathbb C)$. Their classes satisfy
exactly the relation
\[
 \ker\left(\bigoplus_{i=1}^r\widehat{\mathbb Z}[a_i]
      \longrightarrow H_1(X(\mathbb C),\widehat{\mathbb Z})\right)
     =\widehat{\mathbb Z}(1,\ldots,1).
\tag{C.14.3}
\]
Indeed, in the integral relative homology sequence of the pair
$(X(\mathbb C),\bigcup A_i)$, the relative two-dimensional group
is generated by the two halves. Their boundary classes are
$\sum_i[a_i]$ and its negative. Exactness gives the stated
kernel over $\mathbb Z$, and tensoring with the flat
$\mathbb Z$-module $\widehat{\mathbb Z}$ gives (C.14.3).

An isomorphism sends a generator of an unmarked oval centralizer
to a unit power of a generator of the corresponding oval. On a
marked oval the preceding graph argument identifies that unit
with its graph sign $\epsilon_i\in\{1,-1\}$. Applying the
homology isomorphism to (C.14.3) shows that all these unit
multipliers are equal: its image vector must lie in the target
diagonal kernel. In particular all marked-oval signs agree.
The statement is also valid when there is only one marked oval;
then no relative sign needs to be compared. This proves exactly
the simultaneous-reversal equivalence for labelled dividing
components. For nondividing components the independent reversals
are already the equivalence in
Proposition~\ref{C-marked-count}.

For unordered markings, the extension also gives the two-element
orbits of nonreal peripheral classes and hence the number of
interior marked points of the quotient surface. The real incidence
cycles give the numbers of boundary markings on each boundary
component. Configurations of a fixed number of unordered interior
points in a connected surface are connected, as are the boundary
configurations with fixed cyclic incidence. Homeomorphisms of the
quotient surface realize the boundary permutations, and its double
gives the equivariant homeomorphism upstairs. These data therefore
give exactly the unordered topological type. The classification
in Proposition~\ref{C-marked-count} and the full-completion to
procongruence comparison finish the two injectivity assertions.
\end{proof}

\begin{statement}{C.15}{Proved}{the congruence-fibre action and centralizer image}
\label{st:C.15}
\label{C-congruence-fibre-action}
Write
\[
 1\longrightarrow K\longrightarrow P^{\pm}
       \xrightarrow{q} Q^{\pm}\longrightarrow1
\]
for the full-completion to procongruence quotient of an extended
mapping class group, and let $P$ and $Q$ denote its
orientation-preserving subgroups. Fix a discrete real involution
$\tau\in P^{\pm}$, put $\bar\tau=q(\tau)$, and let
$\alpha=\operatorname{Ad}(\tau)$. The action
\[
 C_Q(\bar\tau)\times H^1(C_2,K)\longrightarrow H^1(C_2,K),
 \qquad q(g)\star[k]=[gk\alpha(g)^{-1}]
\tag{C.15.1}
\]
is well defined, where $g\in P$ lifts the indicated centralizer
element. The fibre over $[\bar\tau]$ of the map on conjugacy
classes of orientation-reversing involutions is exactly
\[
 C_Q(\bar\tau)\backslash H^1(C_2,K).
\tag{C.15.2}
\]
For every cocycle $k\in K$, the centralizer fits into the exact
sequence
\[
 1\longrightarrow K^{k\tau}\longrightarrow C_P(k\tau)
 \xrightarrow{q}
 \operatorname{Stab}_{C_Q(\bar\tau)}([k])
 \longrightarrow1.
\tag{C.15.3}
\]
In particular the image on centralizers at $\tau$ is the
stabilizer of the trivial cohomology class for (C.15.1).
For the general extension calculation underlying these formulas,
the action does not automatically preserve that distinguished
class, as the finite example in the proof verifies.
\end{statement}
\begin{proof}
The elements of order two above the particular element
$\bar\tau$ are $k\tau$ with $k\alpha(k)=1$.
Conjugation by $b\in K$ changes $k$ to
$bk\alpha(b)^{-1}$, giving precisely the stated
nonabelian $H^1$. If $q(g)$ centralizes $\bar\tau$, then
$gk\alpha(g)^{-1}$ lies in $K$, and
\[
 g(k\tau)g^{-1}=gk\alpha(g)^{-1}\tau.
\]
This proves that it is again a cocycle. Changing the lift to
$bg$ changes the result by the displayed $K$-conjugacy relation.
Changing a representative of $[k]$ has the same effect, and
composition of conjugations proves the action law.

An involution whose image is $Q$-conjugate to $\bar\tau$ may
first be conjugated by a lift in $P$ so that its image equals
$\bar\tau$. Two such representatives are $P$-conjugate
precisely when their $K$-conjugacy classes are in the same orbit
of (C.15.1). This proves (C.15.2).

The kernel of the centralizer map is exactly $K^{k\tau}$.
Its image fixes $[k]$. Conversely, take an element of $C_Q(\bar\tau)$
fixing $[k]$, and choose a lift $g\in P$. By the definition of
the cohomology class, there is $b\in K$ with
$g(k\tau)g^{-1}=b(k\tau)b^{-1}$. Then $b^{-1}g$ centralizes
$k\tau$ and has the prescribed image in $Q$. This proves
surjectivity onto the stated stabilizer and all of (C.15.3).

To verify the last qualification, take the finite extension with
$P^{\pm}=D_8=\langle r,t\mid r^4=t^2=1,\ trt=r^{-1}\rangle$,
$P=\langle r\rangle$, $K=\langle r^2\rangle$, and $\tau=t$.
Here $Q=P/K$ has order two, and the whole $Q$ centralizes
$\bar\tau$. The action of $\tau$ on $K$ is trivial, so
$H^1(C_2,K)=\{[1],[r^2]\}$. The element $q(r)$ sends $[1]$
to $[r\alpha(r)^{-1}]=[r^2]$. Thus (C.15.2) has one orbit,
whereas this $H^1$ has two elements, and
$C_P(t)=K$ has trivial image in $Q$. These explicit groups
verify why the orbit and stabilizer in the formulas must be
retained. They assert no nonvanishing congruence kernel for a
mapping class group.
\end{proof}

\paragraph{Scope of the C.13--C.15 advance}
Theorems~\ref{C-real-type-injectivity} and
\ref{C-marked-injectivity} establish injectivity for every real
topological type under discussion. Formula (C.15.2) records the
entire remaining fibre over each known congruence class, and
(C.15.3) records its exact centralizer map. The cardinalities of
these orbit sets, the possible additional involution classes in
$Q^{\pm}$, and the comparison of $C_P(\tau)$ with the closure
of the discrete mapping-class centralizer remain uncomputed in
genus three and higher. Thus the surjectivity and complete
centralizer request in Proposition~\ref{C-high-genus-open}
remains incomplete. No assertion of mapping-class-group goodness
or of a vanishing congruence kernel is used in this advance.


\begin{statement}{C.16}{Proved}{the hyperelliptic part of every genus-three real centralizer}
\label{st:C.16}
\label{C-hyperelliptic-reflections}
Fix a hyperelliptic involution $\iota$ of $S_3$, and put
$H^{\pm}=C_{\Gamma_3^{\pm}}(\iota)$ and
$H=H^{\pm}\cap\Gamma_3$. The map of full profinite completions
\[
 \widehat{H^{\pm}}\longrightarrow\widehat{\Gamma_3^{\pm}}
\tag{C.16.1}
\]
is injective, with image $\overline{H^{\pm}}$. Every
orientation-reversing involution of this closed subgroup is
$\overline H$-conjugate to a discrete involution in $H^{\pm}$.
For every such discrete involution $\tau$,
\[
 C_{\widehat{\Gamma_3}}(\tau)\cap\overline H
       =C_{\overline H}(\tau)
       =\overline{C_H(\tau)}.
\tag{C.16.2}
\]
Each of the six real genus-three topological types has a
representative $\tau$ in this $H^{\pm}$.
\end{statement}
\begin{proof}
The hyperelliptic congruence subgroup theorem, in the precise
form of Boggi, Theorem 7.1~\cite{CextraBoggiHyperelliptic}, says
that the composite
\[
 \widehat H\longrightarrow\widehat{\Gamma_3}
       \longrightarrow\operatorname{Out}
               (\widehat{\pi_1(S_3)})
\]
is injective. Its first map is therefore injective. The
orientation subgroups have index two, so the kernel of
(C.16.1) lies in $\widehat H$ and is trivial. Compactness and
density identify the image with the indicated closure.

We give the goodness argument for the group actually used here.
The branched double cover supplies the Birman--Hilden sequence
\[
 1\longrightarrow\langle\iota\rangle\longrightarrow H
       \longrightarrow\Gamma_{0,8}\longrightarrow1.
\tag{C.16.3}
\]
The map sends a commuting mapping class to the induced mapping
class of the quotient sphere with its eight branch points.
Every such sphere mapping class lifts: the double cover on the
punctured sphere is the one sending every peripheral meridian
to the nonidentity element of $C_2$, and that homomorphism is
preserved by every permutation of the branch points. The kernel
in (C.16.3) is the deck group; the mapping-class exactness is the
Birman--Hilden sequence (3) of the cited primary source.

Intersect the inverse image of the pure sphere group with the
kernel of $H\to\operatorname{Sp}_6(\mathbb F_3)$. This is a
finite-index subgroup $L$ of $H$ which misses $\iota$, since
$\iota$ acts by $-1$ on first homology. By (C.16.3), $L$ is
isomorphic to a finite-index subgroup of the pure group
$\Gamma_{0,8}^{\mathrm{pure}}$. The latter is good, as can be
proved directly from its configuration-space fibrations.
Fix three points at $0,1,\infty$; forgetting the last of the
remaining ordered points has a punctured-plane fibre and a
continuous section: choose the additional real point
$2+\sum_j|z_j|$. Thus the fundamental-group extension is a
split extension by a finitely generated free group. Every
finite-index subgroup of this free kernel contains a
characteristic finite-index subgroup. The action on the
corresponding finite quotient has finite image, giving a
finite quotient of the whole semidirect product which
restricts to that quotient of the kernel. Its induced
profinite topology is therefore the full one, and the
completed extension is exact.

Free groups are good; induction on the number of points,
using the two Hochschild--Serre spectral sequences, proves
goodness of the pure sphere group. The finite-type free
resolutions give finite cohomology with finite coefficients
and a finite cohomological dimension at every stage.
Goodness passes to finite-index subgroups by Shapiro's lemma,
and to finite extensions by taking a normal finite-index
subgroup and comparing the Hochschild--Serre sequences.
It follows successively for $L$, $H$, and $H^{\pm}$.
This also proves finite virtual two-cohomological type.
These groups have finite-type resolutions in each degree,
so passage from finite modules to arbitrary discrete
continuous $\mathbb F_2$-modules is obtained by filtered
unions, as in Theorem~\ref{C-profinite-centralizer}.

Residual finiteness follows from the faithful outer action
just used, with orientation detected on the fundamental
class. All hypotheses of Boggi--Zalesskii, Theorem A, now
hold for $H^{\pm}$ and the prime two~\cite{BoggiZalesskii2024}.
That theorem gives the asserted conjugacy and centralizer
identities inside $\widehat{H^{\pm}}$. An
orientation-reversing conjugator can be made
orientation-preserving by multiplying it by the source
involution. Intersecting the centralizer identity with the
open subgroup $\widehat H$ gives (C.16.2).

For completeness, the following smooth hyperelliptic models
give all six real types; $a=0$ denotes dividing and $a=1$
nondividing:
\[
\begin{array}{c|c|c}
 y^2=f(x)&r&a\\ \hline
 -(x^8+1)&0&1\\
 (x^2-1)(x^6+1)&1&1\\
 (x^2-1)(x^2-4)(x^4+1)&2&1\\
 (x^2-1)(x^2-4)(x^2-9)(x^2+1)&3&1\\
 \prod_{j=1}^{4}(x^2-j^2)&4&0\\
 x^8+1&2&0
\end{array}
\tag{C.16.4}
\]
Every polynomial is square-free of degree eight, so its smooth
completion has genus three. With $2k>0$ real branch points
and positive leading coefficient, the $k$ positive intervals
on the real projective line lift to $k$ ovals. For
$k=1,2,3$, the preimage of each open half-plane is connected
because it contains a branch point. A negative real interval
joins these two preimages through nonreal points of the curve;
hence the complement of the ovals is connected. The case
$k=4$ has four ovals and is dividing by
Proposition~\ref{C-classification}.

The negative polynomial in the first row has no real points,
including at infinity. For $x^8+1$, the real double cover of
$\mathbb P^1(\mathbb R)$ is the union of two circles: in the
infinity chart $u=1/x$, $v=y/x^4$, each sign of $y$ joins to
the same sign of $v$. Every point above the real projective
line is real. The preimages of the upper and lower half-planes
are each connected and form the two components of its
complement, proving the last row. Complex conjugation in
each model commutes with $(x,y)\mapsto(x,-y)$. All genus-three
hyperelliptic involutions are conjugate, by the branched
double-cover construction with eight branch points. Transport
to the fixed $\iota$ therefore puts each of the six real
involutions in $H^{\pm}$.

The computed group in (C.16.2) is the specified intersection
with $\overline H$. This calculation does not establish that
every ambient profinite reflection belongs to a conjugate
of $\overline{H^{\pm}}$, or that its entire centralizer lies
in this subgroup. Those are additional, unresolved steps
in the full request.
\end{proof}

\begin{statement}{C.17}{Proved}{congruence centralizers for the five genus-three types with real points}
\label{st:C.17}
\label{C-genus-three-centralizer-reduction}
Let $\tau\in\Gamma_3^{\pm}$ be a real involution with at least
one oval, and put
\[
 P=\widehat{\Gamma_3},\qquad Q=\check\Gamma_3,
 \qquad K=\ker(P\xrightarrow{q}Q),\qquad
 C=C_{\Gamma_3}(\tau).
\]
For closures in the indicated groups, restriction of $q$ gives
an isomorphism
\[
 \overline C^{\,P}\xrightarrow{\ \sim\ }
 C_Q(\bar\tau)=\overline C^{\,Q}.
\tag{C.17.1}
\]
The full profinite centralizer has the topological semidirect
product description
\[
 C_P(\tau)=K^\tau\rtimes\overline C^{\,P}.
\tag{C.17.2}
\]
Thus the centralizer action (C.15.1) fixes the trivial class
for each of these five genus-three types. The value of the
remaining group $K^\tau$ and the other cohomology orbits is
not determined here. The fixed-point-free genus-three type
is excluded from this theorem.
\end{statement}
\begin{proof}
Let $A=\{A_1,\ldots,A_r\}$ be the real oval multicurve and
$\Delta=\widehat{\pi_1(S_3)}$. The ovals are essential and
pairwise nonisotopic. Indeed a contractible oval would bound
a disc exchanged with the other side, forcing a sphere;
two parallel ovals would give an annular component of the
complement, forcing an annulus on each side in the dividing
case or the whole complement to be that annulus in the
nondividing case. Both alternatives force genus one.

First we prove the precise preservation needed to apply a
multicurve stabilizer theorem. Take $B\in C_Q(\bar\tau)$.
The pullback description in
Theorem~\ref{C-real-type-injectivity} shows that $B$ permutes
the section classes of the real surface extension and hence
their centralizers $\overline{\langle a_i\rangle}$. To identify
the actual curve classes rather than just their procyclic
subgroups, choose an auxiliary real marked point on $A_i$.
The procongruence Birman comparison is the isomorphism
\[
 \check\Gamma_{3,1}^{\pm}
 \xrightarrow{\ \sim\ }
 \operatorname{Aut}(\Delta)
       \mathbin{\times}_{\operatorname{Out}(\Delta)}
                         \check\Gamma_3^{\pm}.
\tag{C.17.3}
\]
Here the map on the left is conjugation on the point-pushing
subgroup $\Delta$. Its faithfulness and its exact kernel and
quotient are respectively Corollaries 4.5 and 4.7 of
Boggi~\cite{CextraBoggiCongruence2020}, for the class of all
finite groups. The kernel maps isomorphically to
$\operatorname{Inn}(\Delta)$ because $\Delta$ is centre-free.
Surjectivity onto the displayed pullback follows by lifting
the outer element and then multiplying by the unique needed
inner automorphism. Adjoining the actual pointed reflection
gives the same isomorphism for the orientation extensions.

Choose a lift of $B$ in (C.17.3). Its conjugate of the pointed
reflection is a section of the surface extension and therefore,
by Theorem~\ref{C-profinite-centralizer}, is $\Delta$-conjugate
to a pointed reflection on one of the actual ovals, say $A_j$.
Adjust the lift by that element of $\Delta$. Injectivity in
(C.17.3) shows that it conjugates the two pointed reflections
as elements of the punctured procongruence group, not merely
their filled automorphisms. Apply the one-marked-oval graph
argument of Theorem~\ref{C-marked-injectivity} to its action
on the punctured surface group. It gives
\[
 B([a_i])=[a_j^{\epsilon_i}],\qquad
                   \epsilon_i\in\{1,-1\},
\tag{C.17.4}
\]
where brackets mean geometric conjugacy classes. Repeating
this construction for each oval shows that $B$ preserves
the actual unoriented multicurve $A$ in the procongruence
curve complex.

Let $L=\operatorname{Stab}_{\Gamma_3}(A)$ and let $L_0$ be
the finite-index subgroup fixing each oval and its orientation.
Let $S_1,\ldots,S_d$ be the components after cutting along
$A$ and replacing their boundary circles by labelled
punctures. The ordinary cutting homomorphism has exact
sequence
\[
 1\longrightarrow T\longrightarrow L_0
       \longrightarrow\prod_{j=1}^{d}\Gamma(S_j)^{\mathrm{pure}}
       \longrightarrow1,\qquad T\simeq\mathbb Z^r,
\tag{C.17.5}
\]
where $T$ is generated by the individual oval Dehn twists.
Theorem 4.5 of Boggi~\cite{CextraBoggiProcongruence2014}
identifies $Q_A$ with $\overline L^{\,Q}$ and gives the exact
completed cutting sequence
\[
 1\longrightarrow\widehat T\longrightarrow Q_{\vec A}
       \longrightarrow\prod_{j=1}^{d}
                  \check\Gamma(S_j)^{\mathrm{pure}}
       \longrightarrow1.
\tag{C.17.6}
\]
In particular the twist subgroup here is the injected full
$\widehat{\mathbb Z}^{r}$; it is not a quotient of that group.

Every component $S_j$ in this genus-three situation has genus
at most two. The complete list is
\[
\begin{array}{c|c|c}
 (r,a)&S_3\setminus A\text{ with boundaries punctured}
       &\text{component genera}\\ \hline
 (1,1)&S_{2,2}&2\\
 (2,1)&S_{1,4}&1\\
 (3,1)&S_{0,6}&0\\
 (2,0)&S_{1,2}\amalg S_{1,2}&1,1\\
 (4,0)&S_{0,4}\amalg S_{0,4}&0,0
\end{array}
\tag{C.17.7}
\]
The genus-at-most-two congruence subgroup theorem therefore
identifies each right-hand group in (C.17.6) with the full
completion of its group in (C.17.5). For the positive-genus
component groups this is the precise range of Theorem 1.4(ii)
of \cite{CextraBoggiCongruence2020}; the genus-zero component
groups are included in Theorem 7.1 of
\cite{CextraBoggiHyperelliptic} with the underlying closed
surface a sphere. For the class of all finite groups, the
level-two isomorphism in the first source extends to the
whole group by its finite level quotient.

Completing (C.17.5) is exact at its middle and right terms.
The map of its completed kernel to (C.17.6) is the identity
on $\widehat T$, so it is also injective on the left. The
resulting diagram of exact sequences has isomorphisms at
both ends; its middle map is an isomorphism. Thus
$\widehat{L_0}\simeq Q_{\vec A}$. The finite permutation and
orientation action on the ovals has the same image for
$L$ and its closure, which gives
\[
 \widehat L\simeq Q_A.
\tag{C.17.8}
\]
It follows also that $\widehat L\to P$ is injective, since
its composite with $q$ is the injection in (C.17.8).

We check goodness for the stabilizer to which we will apply
the finite-subgroup theorem. Pure mapping class groups of
genus zero are good by the configuration argument in
Proposition~\ref{C-hyperelliptic-reflections}. The group
$\Gamma_{1,1}\simeq\operatorname{SL}_2(\mathbb Z)$ is
virtually free. The group $\Gamma_2$ is hyperelliptic, so
the same finite-index sphere argument proves its goodness.
Adding punctures preserves goodness by the exact completed
Birman sequence and the goodness of surface groups proved
in Theorem~\ref{C-profinite-centralizer}, using the
Hochschild--Serre comparison. These arguments also give
finite virtual cohomological type and finite-type resolutions.
The product of these component groups has the same
properties. The exactness just verified for the completion
of (C.17.5), and goodness of $T\simeq\mathbb Z^r$, now give
goodness of $L_0$ by the same spectral-sequence comparison.
Finite extension gives goodness and finite virtual
two-cohomological type for $L$ and for
$L^{\pm}=L\rtimes\langle\tau\rangle$. The groups are
residually finite through their faithful congruence actions.
The finite-type resolutions justify the discrete-module
form of the cohomology comparison required in Theorem A.

Because $\tau$ preserves $A$, (C.17.8) extends to
$\widehat{L^{\pm}}\simeq Q_A\rtimes\langle\bar\tau\rangle$.
Boggi--Zalesskii, Theorem A~\cite{BoggiZalesskii2024}, applied
to this particular good group and the prime two, yields
\[
 C_{Q_A}(\bar\tau)=\overline{C_L(\tau)}^{\,Q}.
\]
By (C.17.4), the entire centralizer in $Q$ lies in $Q_A$.
The same preservation statement holds for the discrete
centralizer, so $C_L(\tau)=C$. This proves the centralizer
identity in (C.17.1). The injection $\widehat L\to P$ and
its isomorphism onto $Q_A$ identify the two closures of
$C$, proving all of (C.17.1).

Finally, $\overline C^{\,P}$ centralizes $\tau$, and maps
isomorphically onto the full quotient centralizer. For an
element $u\in C_P(\tau)$, multiply $u$ by the inverse of
the unique element of $\overline C^{\,P}$ with image $q(u)$.
The result lies in $K$ and still centralizes $\tau$.
This factorization is unique because the two factors
intersect trivially. Conjugation by the second factor
preserves $K^\tau$, and compactness gives the asserted
topological semidirect product. The final assertion about
the trivial class follows from (C.15.3) and the now
surjective centralizer map.
\end{proof}

\begin{statement}{C.18}{Proved}{the fixed-point-free comparison}
\label{st:C.18}\label{C-free-reflection-comparison}
Let $\tau$ be a fixed-point-free orientation-reversing
involution of $S_3$, let $N_4=S_3/\langle\tau\rangle$, and put
\[
 E=\widehat{\pi_1(N_4)},\qquad
 \Delta=\widehat{\pi_1(S_3)},\qquad
 Q^\pm=\check\Gamma_3^\pm,\qquad Q=\check\Gamma_3.
\]
The inclusion of the orientation double cover gives the
exact sequence $1\to\Delta\to E\to C_2\to1$, whose outer
action is $\bar\tau$.
Let $\operatorname{Aut}_Q(E;\Delta)$ be the group of
automorphisms preserving $\Delta$ whose restriction to
$\Delta$ has outer class in $Q^\pm$, and define
\[
 \mathcal A=\operatorname{Aut}_Q(E;\Delta)/
                         \operatorname{Inn}(E).
\]
There is a canonical topological isomorphism
\[
 \mathcal A
 \ \xrightarrow{\ \sim\ }\
 C_{Q^\pm}(\bar\tau)/\langle\bar\tau\rangle
 \ \xrightarrow{\ \sim\ }\
 C_Q(\bar\tau).
\tag{C.18.1}
\]
The first arrow is induced by restriction. The second
sends the class of $b$ to $b$ when $b$ preserves orientation
and to $b\bar\tau$ otherwise.
Under this isomorphism the closure of the ordinary mapping
class group of $N_4$ is exactly
$\overline{C_{\Gamma_3}(\tau)}^{\,Q}$.
Consequently the remaining congruence-centralizer assertion
for this type is the explicit density problem
\[
 \overline{\operatorname{Mod}(N_4)}^{\,\mathcal A}
                       \stackrel{?}{=}\mathcal A.
\tag{C.18.2}
\]
This density, and the corresponding full-completion
centralizer, are not established here.
\end{statement}
\begin{proof}
The index-two subgroup $\pi_1(S_3)$ has its full profinite
topology in $\pi_1(N_4)$, so completion gives the asserted
extension. Its action on the completed orientation class
of $\Delta$ is negative; hence $\bar\tau$ is a nontrivial
outer involution.

Choose $e\in E\setminus\Delta$ and write
$a=\operatorname{Ad}(e)|_\Delta$ and $z=e^2\in\Delta$.
Thus $a^2=\operatorname{Inn}(z)$.
For $b\in\operatorname{Aut}(\Delta)$ whose outer class
centralizes $\bar\tau$, there is a unique $d\in\Delta$
such that
\[
 b a b^{-1}=\operatorname{Inn}(d)a.
\]
Uniqueness uses $Z(\Delta)=1$. Squaring this identity
gives
$\operatorname{Inn}(b(z))
 =\operatorname{Inn}(d\,a(d)\,z)$, whence
\[
 b(z)=d\,a(d)\,z.
\]
These two identities show that $b$ extends to an
automorphism of $E$ by the exact formula $e\mapsto de$:
they preserve, respectively, conjugation on $\Delta$
and the square of $e$. Since $E$ is the disjoint union
of $\Delta$ and $\Delta e$, the formula is continuous
and bijective on these two cosets. The extension is
unique, since an automorphism restricting to the identity
on $\Delta$ would have $e\mapsto de$ with
$d\in Z(\Delta)$.
This proves that restriction identifies
$\operatorname{Aut}_Q(E;\Delta)$ with the inverse image
of $C_{Q^\pm}(\bar\tau)$ in
$\operatorname{Aut}(\Delta)$.

The subgroup of inner automorphisms induced by $\Delta$
restricts to $\operatorname{Inn}(\Delta)$; conjugation
by $e$ contributes precisely the outer class $\bar\tau$.
Quotienting therefore gives the first isomorphism in
(C.18.1). The orientation homomorphism and the central
element $\bar\tau$ give the direct product
\[
 C_{Q^\pm}(\bar\tau)
       =C_Q(\bar\tau)\times\langle\bar\tau\rangle,
\]
which proves the displayed formula for the second arrow.
All groups here are closed subgroups or quotients of
automorphism groups of finitely generated profinite
groups. The continuous bijections are thus
homeomorphisms.

For the discrete comparison, every homeomorphism of
$N_4$ has two lifts to its orientation double cover and
exactly one preserves orientation. This choice is
compatible with composition and isotopy. The induced
map
$\operatorname{Mod}(N_4)\to C_{\Gamma_3}(\tau)$ is
an isomorphism: the quotient map from the extended
centralizer has kernel $\langle\tau\rangle$ and is onto.
This is the precise orientation-double-cover statement
recorded by Stukow~\cite[p.~1030]{CextraStukow2016}.
The induced action on $\pi_1(N_4)$, followed by completion
and (C.18.1), is exactly the action of that positive
lift on $\Delta$. Taking closures proves the final
assertion and identifies the unproved equality
(C.18.2) without asserting it.
\end{proof}

\begin{statement}{C.19}{Proved}{no finite-index hyperelliptic reduction for the free genus-three type}
\label{st:C.19}
\label{C-no-finite-index-hyperelliptic}
Let $\tau$ be the fixed-point-free orientation-reversing involution
of $S_3$, and let
\[
 j:\operatorname{Mod}(N_4)
      \xrightarrow{\;\sim\;}C_{\Gamma_3}(\tau)
\]
be the unique orientation-preserving lift of the orientation double
cover, as in Proposition~\ref{C-free-reflection-comparison}.
For every finite-index subgroup $H$ of $\operatorname{Mod}(N_4)$,
the centralizer $C_{\Gamma_3}(j(H))$ has no nonidentity element of
finite order. In particular no such $j(H)$ is contained in the
centralizer of a hyperelliptic involution of $S_3$.
The standard hyperelliptic mapping class subgroup of
$\operatorname{Mod}(N_4)$ therefore has infinite index.
\end{statement}
\begin{proof}
We give the surface and all comparison maps explicitly. Glue two
pairs of pants $P,Q$ along three pairs of boundary circles
$a,b,c$. Choose the gluing signs so that the orientation double
has the following dual graph:
\[
\begin{array}{c|c}
\text{lifted curve}&\text{incident pairs of pants}\\ \hline
a_+,b_+&P_+,Q_+\\
a_-,b_-&P_-,Q_-\\
c_+&P_+,Q_-\\
c_-&P_-,Q_+
\end{array}
\tag{C.19.1}
\]
Equivalently the voltage of the first two edges of the base
theta graph is zero and that of the third is one in $C_2$.
The cover is connected: the cycle using $a$ and $c$ has nonzero
voltage. Thus the base is nonorientable. It is closed and has
Euler characteristic $-2$, so it is $N_4$; the connected
orientation cover has Euler characteristic $-4$ and is $S_3$.
Each of $a,b,c$ is two-sided, with an annular neighbourhood
arising from the gluing collars. The six lifted curves form a
full pants decomposition of $S_3$. They are essential and
pairwise nonisotopic: a disk or an annulus bounded by some
of them would be a union of complementary pants, contradicting
its Euler characteristic.

Choose an orientation on each base annulus and let
$t_a,t_b,t_c$ be the resulting twists. The orientation-cover map
restricts to a homeomorphism on each lifted annulus. Its deck
involution reverses the surface orientation and exchanges the
two lifts. Consequently the positive-lift homomorphism gives
the exact identities
\[
 j(t_a)=T_{a_+}T_{a_-}^{-1},\qquad
 j(t_b)=T_{b_+}T_{b_-}^{-1},\qquad
 j(t_c)=T_{c_+}T_{c_-}^{-1},
\tag{C.19.2}
\]
after choosing the indicated labels using those annular
orientations. The six Dehn twists generate the injected group
$\mathbb Z^6$. The three displayed elements generate its
subgroup with exponent vectors
$(u,-u,v,-v,w,-w)$, hence an injected $\mathbb Z^3$.
Here we use the uniqueness of the support and exponents of a
multitwist, including arbitrary nonzero signed exponents;
the precise statement is Farb--Margalit, Lemma 3.17
\cite{CextraFarbMargalit}.

Finite index supplies positive integers $m_a,m_b,m_c$ such that
$t_a^{m_a},t_b^{m_b},t_c^{m_c}\in H$.
Let $\iota\in C_{\Gamma_3}(j(H))$ have finite order.
Conjugation by an orientation-preserving mapping class takes
$T_d$ to $T_{\iota(d)}$ with the same exponent.
Applying multitwist uniqueness to each identity in (C.19.2)
therefore shows that $\iota$ fixes each of
$a_+,a_-,b_+,b_-,c_+,c_-$ individually: the two exponents in
one pair are $m$ and $-m$, so they cannot be interchanged.

Every vertex in (C.19.1) has a different set of incident labelled
edges. An automorphism fixing all six edges consequently fixes
all four vertices. A representative of $\iota$ preserving the
multicurve thus preserves each complementary pair of pants
and each of its three boundary circles. Its restriction, with
boundaries replaced by punctures, lies in
$\operatorname{PMod}(S_{0,3})=1$.
The cutting homomorphism has exactly the multitwist kernel
$\mathbb Z^6$: these are Proposition 2.3 and Proposition 3.20
of the same primary text, with its explicit cutting map
\cite{CextraFarbMargalit}. Thus $\iota$ belongs to this
torsion-free kernel and is the identity.

This excludes every nonidentity finite-order positive class,
without an assumption that it already commutes with $\tau$.
For the standard hyperelliptic subgroup of $\operatorname{Mod}(N_4)$,
the positive-lift homomorphism is explicitly a map into the
genus-three hyperelliptic centralizer
\cite[p.~1030]{CextraStukow2016}. It cannot have finite-index
domain by the result just proved. The finite graph and the
three independent exponent vectors are also checked by
\texttt{code/B/n4\_pants\_check.py}; the infinite-group
conclusion follows from the written cutting argument.
\end{proof}

% ---- D ----
\section{Mordell through splittings}
\label{sec:D}
Let $C/K$ be smooth, proper, geometrically connected of
genus $g\geq2$, with $K$ a number field. Put
$\Delta=\pi_1(C_{\overline K})$, $E=\pi_1(C)$,
$J=\operatorname{Pic}^0_{C/K}$, and
\[
 \operatorname{Sec}(C/K)=
 \operatorname{Sect}_{\rm cts}(E\longrightarrow G_K)/\Delta .
\]
All statements below using the local section theorem retain
the provenance specified in Part B: Theorem 1.1 of the
October 2026 paper is a stipulated input of this problem.
They are deductions from that input and the displayed
published results; they do not claim an independent
verification of every external input of that paper.

\subsection{D(a): exact logical maps and the abelian computation}

\begin{statement}{D.1}{Proved}{the finiteness diagram}
\label{st:D.1}
\label{D-logical}
Grothendieck's assertion on \S8, folio 33, scan 65 is
the finiteness of $\operatorname{Sec}(C/K)$.
The actual maps relevant to it are
\[
 C(K)\ \xrightarrow{\ \kappa\ }\ \operatorname{Sec}(C/K)
 \ \xrightarrow{\ \operatorname{loc}\ }\
 C(\mathbb A_K)_\bullet^{\mathrm{f\text{-}cov}}
 \ \lhook\joinrel\longrightarrow\
 \prod_{v<\infty}C(K_v)\times
 \prod_{v\mid\mathbb R}\pi_0(C(K_v)).
\]
Here $\kappa$ is injective, and $\operatorname{loc}$
is bijective. Consequently the following are exact
equivalences:
\[
 \begin{split}
 \kappa\text{ is bijective}
 &\ \Longleftrightarrow\
 C(\mathbb A_K)_\bullet^{\mathrm{f\text{-}cov}}=C(K),\\
 \#\operatorname{Sec}(C/K)<\infty
 &\ \Longleftrightarrow\
 \#C(\mathbb A_K)_\bullet^{\mathrm{f\text{-}cov}}<\infty.
 \end{split}
\]
The first equality and Faltings's theorem give the second
finiteness. Finiteness of $\operatorname{Sec}(C/K)$ gives
finiteness of $C(K)$ through the displayed injection.
For a finite section set, its cardinality equals that
of $C(K)$ exactly when $\kappa$ is bijective.
\end{statement}
\begin{proof}
Part B proves Kummer injectivity. The local theorem gives
a unique point at every finite place; Part C gives a
component at every real place. Harari--Stix, Theorem 2.1,
with the normal subgroup $U=1$ and omitted set of places
empty, identifies survival of every finite torsor with
the existence of a global section inducing these local
classes~\cite{HarariStix2012}. Finite torsor evaluations
at real places factor through components. This gives
both inclusions in the asserted image of $\operatorname{loc}$,
exactly as in Proposition 8.1 of the supplied local
paper~\cite{B_OpenAIPadic}.

For injectivity of $\operatorname{loc}$, two sections
with equal local points have equal reductions at all
finite places of good reduction. They are Selmer
sections and are noncuspidal because $C$ is proper.
Porowski's 2025 Theorem 1.1 therefore identifies their
classes~\cite{Porowski2025}; the reduction used in its
application is recorded in Proposition~\ref{D-rigidity}.
The first equivalence can also avoid this last
injectivity theorem: an adelic tuple supported on a
single rational point makes the section a point
section by Stix's finite-support Corollary 5.
The remaining assertions follow from the displayed
injection and bijection. Faltings's theorem says
$\#C(K)<\infty$~\cite{Faltings1983}. We have not reversed
either implication from finiteness to surjectivity;
the actual unresolved part is the possible complement
of the image of $\kappa$, as recorded below.
Stoll's Remark 8.9 establishes the accompanying
section-property, finite-descent, and birational-property
implications in its stated ranges~\cite{Stoll2007}.
\end{proof}


\begin{statement}{D.2}{Known}{every abelian section class}
\label{st:D.2}
\label{D-abelian}
The extension with geometric kernel $\Delta^{\rm ab}=TJ$
is a class in $H^2_{\rm cts}(K,TJ)$. Its section set is
empty when that class is nonzero. A choice of one
section identifies its section set with
$H^1_{\rm cts}(K,TJ)$. In particular, a point of $C(K)$
supplies such a choice. There are exact sequences
\[
 0\longrightarrow\widehat{J(K)}
 \longrightarrow H^1_{\rm cts}(K,TJ)
 \longrightarrow T H^1(K,J)\longrightarrow0
\]
and
\[
 0\longrightarrow\widehat{J(K)}
 \longrightarrow\operatorname{Sel}(TJ)
 \longrightarrow T\Sha(K,J)\longrightarrow0.
\]
The Selmer subgroup in the second line is the inverse
image of the local Kummer images at every place.
Here the hats are inverse limits modulo multiplication
by positive integers, and $T$ is the inverse limit
of torsion with the corresponding multiplication
transition maps.
\end{statement}
\begin{proof}
After choosing a splitting, another splitting has
the form $g\mapsto(c(g),g)$, and the homomorphism
identity is exactly the continuous cocycle identity.
Conjugation by $TJ$ changes $c$ by a coboundary.
Without a chosen splitting the extension cocycle is
the stated degree-two obstruction.

The Kummer exact sequence at level $n$ is
\[
 0\longrightarrow J(K)/nJ(K)
 \longrightarrow H^1(K,J[n])
 \longrightarrow H^1(K,J)[n]\longrightarrow0.
\]
Mordell--Weil makes its left term finite. These terms
have surjective transition maps, so their first
derived inverse limit vanishes. The groups
$H^0(K,J[n])$ are finite as well; thus continuous
cohomology of $TJ$ in degree one is the inverse
limit of these degree-one groups. Taking inverse
limits proves the first line. Intersecting with the
local Kummer conditions at finite level replaces
$H^1(K,J)[n]$ by $\Sha(K,J)[n]$ and gives the second.
This is the Selmer computation used in Stoll,
\S\S3--5~\cite{Stoll2007}.
\end{proof}

\begin{statement}{D.3}{Proved}{an infinite abelian approximation and an explicit nonpoint class}
\label{st:D.3}
\label{D-abelian-examples}
The abelian section set is infinite in general,
even for curves of genus two with a rational base point.
One example is the smooth completion of
\[
 C_+:\quad y^2=x^6-x^2+1.
\]
Moreover, on the genus-two curve
\[
 C_0:\quad y^2=x^5-x+1
\]
with base point $\infty$, the class
\[
 \kappa_J\bigl(2[(0,1)-\infty]\bigr)
 \in H^1(K,TJ),\qquad K=\mathbb Q,
\]
is a global section of the abelianized extension which
is induced by no point of $C_0(K_v)$ at any finite place
$v$.
\end{statement}
\begin{proof}
The first curve maps to $A:y^2=u^3-u+1$ by
$(x,y)\mapsto(x^2,y)$. The point $P=(0,1)$ on $A$
satisfies $2P=(1/4,-7/8)$. The integral-coordinate
conclusion of the Nagell--Lutz theorem for torsion
on this integral short Weierstrass equation shows
that $P$ has infinite order. The push-pull identity
for the degree-two map is $f_*f^*=[2]$ on $A$.
Consequently $J(C_+)(\mathbb Q)$ has positive rank,
and its completion contains $\widehat{\mathbb Z}$.
Proposition~\ref{D-abelian} proves infinitude.
Both points at infinity of $C_+$ are rational.

For the second curve put $P=(0,1)$. Suppose over
an algebraically closed characteristic-zero field
that $2[P-\infty]=[Q-\infty]$. Then
$2P\sim Q+\infty$ are effective divisors of degree
two. If these divisors differ, they give a moving
degree-two system on a genus-two curve. Riemann--Roch
identifies that system with the unique hyperelliptic
system. But $2P$ is not a hyperelliptic fibre:
$P$ is not a ramification point. If the divisors
are equal, their supports give $P=\infty$, also
impossible. Thus the class is outside the
Abel--Jacobi image over every extension field.

For a finite extension of $\mathbb Q_p$, the group
$J(K_v)$ is compact and has an open subgroup
isomorphic to $\mathbb Z_p^{g[K_v:\mathbb Q_p]}$.
Its power completion is therefore itself, and
its Kummer map into $H^1(K_v,TJ)$ is injective.
A local point inducing the displayed section
would contradict the preceding divisor argument.

The exact program checks smoothness of both
polynomials, the doubling calculation, and the
reduced Mumford representation
\[
 (u,v)=(x^2,1-x/2),\qquad v^2\equiv x^5-x+1\pmod u .
\]
Its reduced degree is two, which independently
certifies the degree-one Abel--Jacobi exclusion.
These are statements about the actual quotient
$\Delta\to TJ$, with its induced section map.
\end{proof}

\subsection{D(b): the finite-descent criterion and explicit curves}

\begin{statement}{D.4}{Proved}{the Stoll range in terms of splittings}
\label{st:D.4}
\label{D-Stoll}
Every section is induced by a unique rational point
for every proper hyperbolic curve $C/K$ with a
nonconstant map to an abelian variety $A/K$ satisfying
\[
 \#A(K)<\infty,\qquad
 \Sha(K,A)_{\rm div}=0 .
\]
In particular this holds for $X_0(N)$ and $X_1(N)$
over $\mathbb Q$ of genus at least two. The section
set is finite in each case. Explicit genus-two
examples over $\mathbb Q$ are the completions of
\[
 y^2=x^6+a,\qquad a\in\{-2,-1,1,2\}.
\]
Explicit genus-three examples are
\[
 y^2=x^8+4,\qquad y^2=x^8+64.
\]
\end{statement}
\begin{proof}
Stoll's abelian descent computation identifies
$A(\mathbb A_K)_\bullet^{\rm f\text{-}ab}$ with $A(K)$
under the displayed assumptions. Push a finite-descent
adelic point of $C$ to $A$. It lands at a point
$a\in A(K)$. The fibre $Z=f^{-1}(a)$ is finite,
since $f$ is nonconstant and $C$ is a proper curve.
Stoll's finite-support theorem, Theorem 8.2, gives
\[
 C(\mathbb A_K)_\bullet^{\rm f\text{-}ab}
 \cap Z(\mathbb A_K)_\bullet=Z(K).
\]
This proves the finite-descent equality for $C$,
and Proposition~\ref{D-logical} gives precisely
the required splitting statement. We use the
corrected real-component convention in Stoll's
erratum, \S3; finite-place support is sufficient
for this argument~\cite{StollErratum}.

For the modular curves, Corollary 8.8 of Stoll
supplies this equality: the Jacobian has a
nontrivial modular quotient of analytic rank zero,
and the rank-zero finiteness results give the
required Mordell--Weil and Tate--Shafarevich
properties. This is the complete chain behind
Corollary 8.5 of the supplied 2026 local paper:
local Theorem 1.1, Harari--Stix Theorem 2.1,
Stix Corollary 5, and Stoll Corollary 8.8.

For $y^2=x^6+a$, the two maps to elliptic curves are
\[
 (x,y)\longmapsto(x^2,y)\in E_a,\qquad
 (x,y)\longmapsto(a/x^2,ay/x^3)\in E_{a^2},
 \quad E_b:\ v^2=u^3+b .
\]
They extend to the smooth completions. Stoll's
examples immediately following Theorem 8.6
verify the analytic-rank-zero criterion for
the four listed values. For genus three use
$A:y^2=x^3+4x$, the elliptic modular curve
$X_0(32)$, whose rational Mordell--Weil group has
rank zero and whose Tate--Shafarevich group is finite.
The maps are respectively
\[
 (x,y)\longmapsto(x^4,x^2y),\qquad
 (x,y)\longmapsto(x^4/4,x^2y/8).
\]
Direct substitution verifies their equations.
The squarefree polynomials of degree eight
give genus three. The identification of this
elliptic curve with $X_0(32)$, or its standard
degree-two isogeny to $y^2=x^3-x$, supplies
the same rank-zero modular input~\cite{X032Model}.
Finally Faltings's theorem makes each rational
point set finite.
\end{proof}

\subsection{D(c): finite local images and the global fibre}

\begin{statement}{D.5}{Known}{local finiteness and reduction rigidity}
\label{st:D.5}
\label{D-rigidity}
For $K$ containing no CM subfield, the image
$\operatorname{loc}_v\operatorname{Sec}(C/K)$
is finite at every finite place $v$.
For arbitrary number fields $K$, the reduction
of Selmer sections on a set of strongly positive
density determines their classes. Strongly
positive density means positive upper density
after pullback to the primes of every finite
extension of $K$.
\end{statement}
\begin{proof}
Every global section is Selmer by the stipulated
local theorem. Betts--Stix, Theorem A, applies
to exactly this Selmer set and gives the first
claim~\cite{BettsStix2025}. Its input is a finite
set of global semisimple Galois representations
from the chosen abelian-by-finite family, followed
by the finite period-map loci of their types
(a) and (b); type (c) is excluded. The proof uses
Corollaries 6.10 and 6.16 and Proposition 6.20
to obtain Theorem 6.5. The choice of auxiliary
family depends on the place. No place-independent
cardinality bound is asserted there.

For the second claim we checked the following
chain in the supplied source of Porowski 2026,
Theorem 1.5~\cite{Porowski2026}. A finite
semistable partition at a fixed place is
separated, after a finite cover and a finite
constant extension, by finitely many smooth
proper families whose Jacobians have different
toric ranks across distinct parts. The
construction is the source's Theorem 6.4,
its compactness step in \S5, and the
resolution-of-nonsingularities step in \S6.
For the actual families the semisimple
representations are integral, pure of weight
one, of bounded dimension, and unramified
outside a fixed finite set. The Faltings
finiteness theorem used in \S7 therefore
gives a finite list.

Pass to a finite extension over which the
connected algebraic monodromy groups of each
pair are connected. If two such semisimple
representations are not potentially isomorphic,
their trace-equality locus is a proper
algebraic subset of that connected monodromy
group, hence has Haar measure zero. This is
the algebraic Chebotarev argument in Rajan's
Theorem 2~\cite{Rajan1998}. Intersecting the
finitely many complementary sets produces a
set of potentially density one. Equality of
reductions at any prime in that set forces
equality of the fixed local partition, since
different toric ranks remain different under
finite local extension. This is precisely
the source's set $\Omega(v)$.

A strongly positive set meets every potentially
density-one set: pull back to an extension
where the latter has density one and use
positive upper density of the former.
Consequently equality of reductions on the
given strongly positive set gives equality
at every place. Porowski 2025, Theorem 1.1,
then identifies the sections. Its density-one
argument uses the weak-cover property of
Theorem 3.3 and the noncuspidal reduction
estimate in Theorem 5.4; the proper curve
here has no boundary and satisfies their
noncuspidal hypotheses. This records the
exact use of both recent papers. We use no
extension of the density hypothesis to
arbitrary infinite sets.
\end{proof}

\begin{statement}{D.6}{Proved}{the precise uniform reduction criterion}
\label{st:D.6}
\label{D-uniform}
Put $r(v)=\#\operatorname{red}_v\operatorname{Sel}(C/K)$.
For each integer $d\geq1$, the inequality
$\#\operatorname{Sel}(C/K)\leq d$ is equivalent
to the existence of a strongly positive set
$\Omega$ with $r(v)\leq d$ for every $v\in\Omega$.
The step in the supplied 2026 source replacing
strongly positive density by strictly positive
density has an unproved intersection assertion.
\end{statement}
\begin{proof}
The forward implication uses all finite places.
For the reverse, take $d+1$ distinct classes.
At every $v\in\Omega$ some pair has equal
reduction. Thus $\Omega$ is a union of
finitely many pair-equality sets. At least
one of those sets has strongly positive
density. Indeed, for each set failing
strong positivity there is a finite extension
where its pullback has upper density zero.
The compositum of these finitely many
extensions makes all their pullbacks have
upper density zero; finite extension preserves
density zero, since degree-one primes have
bounded multiplicity and primes of residue
degree greater than one have density zero.
Their union cannot contain the pullback
of $\Omega$. Apply Proposition~\ref{D-rigidity}
to the resulting pair to get a contradiction.

The unnumbered corollary following the source's
Corollary 7.10 states the bound with strictly
positive density, and its proof chooses
$v'\in\Omega\cap\Omega(v)$. The established
property of $\Omega(v)$ is potentially density
one. For a nontrivial quadratic extension
$L/K$, the inert primes have density $1/2$,
the completely split primes have potentially
density one, and their intersection is empty.
This explicit Chebotarev example refutes the
intersection inference in that proof.
It does not supply a counterexample curve
to the corollary itself. The strongly positive
version has the complete proof above.
\end{proof}

\begin{statement}{D.7}{Proved}{an arithmetic model of the remaining uniformity obstruction}
\label{st:D.7}
\label{D-density-model}
There is an infinite compact totally disconnected
Hausdorff space with continuous finite images at
each rational prime, with injective product
map over every strongly positive set of primes,
and with all the potentially-density-one
partition refinements used in the preceding
proof. Thus these three properties, as
properties of continuous maps of compact spaces, do not establish
finiteness of the source set.
\end{statement}
\begin{proof}
Enumerate the rational primes $p_1,p_2,\ldots$,
and put $L_m=\mathbb Q(\sqrt{p_1},\ldots,\sqrt{p_m})$,
$L_0=\mathbb Q$. Let $b(p)$ be the largest
$m$ for which $p$ is unramified and completely
split in $L_m$. It is finite, because the
step adjoining $\sqrt p$ ramifies at $p$.
Give $S=\mathbb N_0\cup\{\infty\}$ the one-point
compactification topology, and define
\[
 f_p(n)=\min\{n,b(p)\},\qquad f_p(\infty)=b(p).
\]
Every image is finite and discrete. Each map is
continuous, because it is constant on the cofinite
neighbourhood $\{n\ge b(p)\}\cup\{\infty\}$.
This proves the required compact topology explicitly.
For any distinct finite $n,n'$
and any strongly positive $\Omega$, the set
of primes completely split in
$L_{\max(n,n')}$ has potentially density one,
so it meets $\Omega$. At such a prime
$f_p(n)\ne f_p(n')$. The product is injective.
For the pair $n,\infty$, use instead the completely
split primes of $L_{n+1}$; there $b(p)>n$, so
$f_p(n)\ne f_p(\infty)$ as well. The continuous
product injection is a homeomorphism onto a closed
subspace of the product of its finite images,
since its source is compact and the target Hausdorff.

For a fixed $v$, complete splitting in
$L_{b(v)}$ implies that equality under $f_p$
implies equality under $f_v$. Hence the
partition-refining set has potentially
density one. All assertions are realized
by the explicit maps $f_p$, while
$S$ is infinite. Its sequence $n\to\infty$
also realizes a compatible limit that every fixed
finite collection of local maps eventually fails
to distinguish. The program
records the first splitting tests using
Legendre symbols. We make no identification
of this set with a section set of a curve.
The morphisms displayed here prove exactly
which purely formal inference fails.
\end{proof}

\begin{statement}{D.8}{Open}{the global finiteness fibre}
\label{st:D.8}
\label{D-open}
The finiteness of $\operatorname{Sec}(C/K)$
for every proper hyperbolic curve over a number
field remains undecided in this return.
Proposition~\ref{D-literal-minimum} gives the literal
least quotient for every locally soluble curve;
Theorem~\ref{D-no-least-quotient} disproves the existence
of a least quotient in general.
\end{statement}
\begin{proof}[State of the attack]
The maps in Proposition~\ref{D-logical} identify
the first issue with finiteness of the
finite-descent locus. Proposition~\ref{D-rigidity}
makes each localization image finite over
fields without CM subfields, while
Proposition~\ref{D-uniform} identifies a
sufficient and necessary uniform reduction
bound. Proposition~\ref{D-density-model}
proves why the placewise and density
statements alone do not supply it.

We tried four routes. The finite list of
Galois representations in a period-map
argument bounds the local image for its
chosen auxiliary family; the degree of the
separating cover and the family change
when another local partition is imposed.
No bounded common list was obtained.
The finite-cover neighbourhood argument
turns a section supported at a known
rational point into that point section,
but it supplies no such support for a
general member of a localization fibre.
The potentially-density-one refinements
control each finite partition separately;
the explicit model above shows the missing
uniform bound. Finally, compactness of an
inverse system of finite quotient data
constructs compatible limits; it gives
no finite bound on their number.

The smallest genus left by this argument
is two over $\mathbb Q$, after the
finite-descent cases of Theorem~\ref{D-Stoll}
have been removed. More precisely the
unresolved object is a fibre of
\[
 \operatorname{Sec}(C/\mathbb Q)
 \longrightarrow C(\mathbb Q_p)
\]
over one of its finitely many values, for
a genus-two curve to which no verified
finite-descent criterion has been applied.
The next mathematical step is to bound
that fibre using a fixed finite collection
of separating families, or to prove a
uniform bound as in Proposition~\ref{D-uniform};
neither bound is asserted here.

For the quotient question, quotients are
ordered by factorization of continuous
surjections over $G_K$, and local
point-inducedness means membership in
the image of the actual map
$C(K_v)\to\operatorname{Sec}(E_Q/G_{K_v})$.
The full group has the required property
under the stipulated local input.
Proposition~\ref{D-abelian-examples}
excludes the abelian quotient by an
explicit global section. The cited
local pro-prime, metabelian and
nilpotent-by-abelian counterexamples
cannot by themselves produce a global
section: the required map is the
restriction
$\operatorname{Sec}(E_Q/G_K)\to
\operatorname{Sec}(E_Q/G_{K_v})$,
whose image must be controlled.
The independent global construction
below addresses one such image.
Pop's small birational quotients start
with $G_{K(C)}$ and include valuative
inertia; the map to $\pi_1(C)$ kills
the inertia of every point of $C$.
We therefore retain the actual quotient
maps and do not infer a least quotient
of $\pi_1(C)$ from that birational theorem.
\end{proof}


\begin{statement}{D.9}{Proved}{explicit global pro-three sections failing local geometricity}
\label{st:D.9}
\label{D-global-prothree}
Put
\[
 L=\mathbb Q(\zeta_9,\sqrt[3]{3}),\qquad
 X_9:\ X^9+Y^9+Z^9=0 .
\]
The proper curve $X_9/L$ has genus $28$. Its
geometrically pro-$3$ arithmetic fundamental
group has infinitely many section classes
which are induced by no $L_w$-point at a
place $w$ above $3$. The same construction
gives such sections after passing to each
of its geometric pro-$3$ nilpotent or
solvable quotients of positive class or
derived length.

This is an explicit instance of the global
phenomenon in Hoshi's 2010 Theorems A and B,
whose proof and scope we checked; we make
no priority claim~\cite{Hoshi2010}.
\end{statement}
\begin{proof}
Let $K=\mathbb Q(\zeta_3)$ and let $G$ be
the maximal pro-$3$ Galois group unramified
outside its unique prime above $3$.
The field has class number one, contains
$\mu_3$, and has just one prime above $3$.
Shafarevich's theorem makes $G$ free pro-$3$;
the precise reduction is given in
Hoshi, Lemma 3.2, and in
\c{C}iperiani--Stix, \S4, using
NSW Corollary 10.7.14~\cite{CiperianiStix}.
Its rank is two: Kummer theory identifies
$H^1(G,\mathbb F_3)$ with the two-dimensional
space of $S$-units modulo cubes, generated
by $\zeta_3$ and $\lambda=1-\zeta_3$.

Let $E=X_3$, based at $O=[1:-1:0]$.
Its equation is $y^2+y=x^3-7$, or
$Y^2=X^3-432$. All of $E[3]$ is rational
over $K$: the nonzero division points have
$X=0$, $Y=\pm12\sqrt{-3}$, or
$X=12\zeta_3^i$, $Y=\pm36$.
The curve has good reduction outside $3$.
Thus the action on $M=T_3E$ factors through $G$.

For completeness the rank-zero arithmetic input
has an exact verification. The normalized
newform of $E$ is
\[
 f(z)=\eta(3z)^2\eta(9z)^2 ,
\]
as recorded in Alfes--Mertens, \S4.2~\cite{AlfesMertens2024}.
Every factor in its product expansion is
positive on the positive imaginary axis, so
\[
 L(E,1)=2\pi\int_0^\infty f(it)\,dt>0.
\]
The integral converges at both cusps. The
quadratic twist by $-3$ has equation
$Y^2=X^3+11664$ and is rationally
$3$-isogenous to $E$, through the standard
map from $Y^2=X^3+B$ to
$Y^2=X^3-27B$. It has the same nonzero
$L$-value. The rank-zero theorem of
Kolyvagin gives finite Mordell--Weil and
Tate--Shafarevich groups for both curves
over $\mathbb Q$~\cite{Kolyvagin1988}.
The plus and minus decomposition for
$K/\mathbb Q$, whose degree is prime to
$3$, gives
\[
 \operatorname{rank}E(K)=0,\qquad
 \#\Sha(K,E)[3^\infty]<\infty .
\]

Choose free generators $\tau_1,\tau_2$ of
$G$. Continuous cocycles into $M$ are
arbitrary pairs of elements of $M$,
and their coboundaries are
\[
 m\longmapsto
 \bigl((\tau_1-1)m,(\tau_2-1)m\bigr).
\]
The kernel is $M^G=0$, since $E(K)$ has
finite torsion. Hence $H^1(G,M)$ has
$\mathbb Z_3$-rank two. Its image in
$H^1(K,M)$ is injective by inflation.
The global $3$-adic Selmer subgroup is
finite by the rank and Tate--Shafarevich
calculation. Choose an infinite-order
class $a\in H^1(G,M)$.

At a finite place away from $3$, the
unramified cohomology of the good elliptic
curve is exactly its local Kummer image.
This follows from the good-reduction
Kummer sequence and
$M/(\operatorname{Frob}-1)M$; it is also
the calculation in
\c{C}iperiani--Stix, \S2.2.
There is just one remaining finite place
of $K$. Consequently the image
\[
 \overline a_3\in T_3H^1(K_3,E)
\]
is nonzero. This explicitly locates the
failure of local geometricity.

Now put $U_3=X_3\setminus\{XYZ=0\}$ and
$N=\pi_1(U_{3,\overline K})^{(3)}$.
The group $N$ is free pro-$3$ of rank ten.
Use the tangential section at $O$ with
integral parameter $t=Z/X$ and tangent
value one. The based action on $N$ is
unramified outside $3$: the smooth pair
and this nonzero tangent extend over
$\mathcal O_K[1/3]$, and the prime-to-residue-
characteristic specialization equivalence
preserves the chosen punctured-disc fibre
functor. This establishes the assertion
for the actual action.

On $N^{\rm ab}/3$, its submodule of
boundary classes and its elliptic
quotient both have trivial $G_K$-action:
the nine cusps and $E[3]$ are $K$-rational,
and $\mu_3\subset K$. Its residual image
is therefore a group of unipotent block
matrices of order a power of three.
The kernel of
$\operatorname{Aut}(N)\to
\operatorname{Aut}(N^{\rm ab}/3)$
is pro-$3$, by the characteristic
Frattini series. The full based action
has pro-$3$ image and factors through $G$.

The graph of a representative of $a$ is
a homomorphism $G\to M\rtimes G$.
Projectivity of the free pro-$3$ group
lifts it through the epimorphism
\[
 N\rtimes G\longrightarrow M\rtimes G.
\]
The projection of the lift to $G$ is
the identity, because the displayed
map preserves it. Pulling back along
$G_K\to G$ gives a section of the
geometrically pro-$3$ group of $U_3$,
whose image on $E$ is $a$.

The coordinate-cube map $U_9\to U_3$
is a degree-nine torsor under $\mu_3^2$.
It is defined over $K$. At the chosen
cusp its parameters satisfy $t_3=t_9^3$,
so the reference tangential section has
zero covering character. The covering
character of the newly constructed
section factors through
$G^{\rm ab}/3$. Its fixed extension is
contained in
\[
 K(\sqrt[3]{\zeta_3},\sqrt[3]{\lambda})
   =K(\zeta_9,\sqrt[3]{3})=L ,
\]
where $\lambda^2=-3\zeta_3$ verifies
the equality. Thus over this fixed
field $L$ the section lifts to $U_9$.
The geometric subgroup corresponding
to this cover is an open subgroup of
$N$, and pro-$3$ covering comparison
identifies it with the pro-$3$ group
of $U_9$. Filling its cusps gives a
section on $X_9$.

Restriction on $T_3H^1(K_3,E)$ is
injective after any finite local
extension: corestriction composed
with restriction is multiplication
by the extension degree, and this
Tate module is $\mathbb Z_3$-torsion-free.
The image of $\overline a_3$ in
$T_3H^1(L_w,E)$ remains nonzero.
A point inducing the constructed
section on $X_9/L_w$ would map under
$X_9\to X_3=E$ to a point inducing
that elliptic section, a contradiction.

The rank-two lattice modulo its
finite Selmer subgroup provides
infinitely many distinct local
non-Kummer images, and restriction
preserves them. Hence infinitely
many classes remain distinct on
$X_9$. Finally every positive
nilpotent-class or derived-length
quotient still maps to the abelian
group $T_3E$, so the same obstruction
survives under the stated quotient
maps. The genus is
$(9-1)(9-2)/2=28$.
The program checks the elliptic
equations, division polynomial,
isogeny identity, and covering ranks.
\end{proof}



% ---- D_tail ----


\begin{statement}{D.10}{Proved}{the literal least quotient on locally soluble curves}
\label{st:D.10}
\label{D-literal-minimum}
For a fixed $C/K$ with $C(K_v)\ne\varnothing$ at
every place, the least quotient of $\pi_1(C)$
over $G_K$ for which every global section is
locally point-induced is $G_K$ itself.
Its geometric kernel is trivial.
\end{statement}
\begin{proof}
The quotient is the structure map
$\pi_1(C)\twoheadrightarrow G_K$.
Its structure map to $G_K$ is the identity,
whose only section is the identity. After
pullback to $G_{K_v}$ the structure map is
again the identity. For every $x_v\in C(K_v)$
the composite
\[
 G_{K_v}\xrightarrow{s_{x_v}}\pi_1(C_{K_v})
 \longrightarrow G_{K_v}
\]
is exactly that identity. Thus the unique
global section is locally point-induced.
Every quotient over $G_K$ maps to $G_K$,
which proves minimality in the factorization
order. Conversely, at a place with no
local point the image of the point map
is empty, so this quotient does not have
the property.

The literal assertion consequently has
this answer on every locally soluble
curve, including every curve with a
rational point. The comparison sends
all its point-section classes to the
same class. Requiring recovery of
the original point-section classes
would impose an additional condition
on this very map; we have not inserted
such a condition into the question.
A uniform selection for all curves,
including the locally insoluble ones,
is the remaining interpretation
recorded in Proposition~\ref{D-open}.
\end{proof}


\begin{statement}{D.11}{Refuted}{a least quotient need not exist}
\label{st:D.11}
\label{D-no-least-quotient}
For a fixed proper curve $C/\mathbb Q$, order arithmetic
quotients by their specified continuous epimorphisms
$E_C\twoheadrightarrow Q$ over $G_{\mathbb Q}$:
a smaller quotient is obtained by a commuting further
epimorphism. Say that $Q$ has property $\mathcal P_C$
when every global section of $Q\to G_{\mathbb Q}$ is,
at every place $v$, in the image of the actual map
$C(\mathbb Q_v)\to\operatorname{Sect}(Q_v)/\ker(Q_v\to G_v)$.

There is no least quotient with $\mathcal P_C$ for
the explicit smooth projective genus-five curve
\[
 C:\quad y^2=-(u^6+1)(u^6+4).
 \tag{D11.1}
\]
There are two quotients $Q_1,Q_2$, each with geometric
kernel $C_2$, which have no global sections and hence
have $\mathcal P_C$. Their only common further quotient
over $G_{\mathbb Q}$, with the indicated maps from $E_C$,
is $G_{\mathbb Q}$; that quotient fails $\mathcal P_C$.
Together with Proposition~\ref{D-literal-minimum}, this
decides the least-quotient request in D(c) under the
literal stated property. The genus-five example also
has $C(\mathbb Q_2)=\varnothing$, so the failure persists
when only finite places are required.
\end{statement}
\begin{proof}
We first construct the quotient maps, including their
arithmetic descent. Consider the smooth projective
genus-one curve
\[
 T:\quad y^2=-(x^2+1)(x^2+4).
\]
It has no real point: the right side is negative for
finite real $x$, and the leading coefficient $-1$
excludes real points at infinity. Its Jacobian has
rational full two-torsion. Indeed the four branch
points of $T\to\mathbb P^1$ are
$i,-i,2i,-2i$. The three nonzero geometric two-torsion
classes correspond to their three unordered partitions
into two pairs. The only nontrivial Galois permutation
is $(i,-i)(2i,-2i)$, which fixes every partition.
For this identification one may choose a branch
point as origin over $\overline{\mathbb Q}$:
the differences of the other branch points give
exactly the three nonzero points of order two.
Changing the origin translates all four labels and
preserves the three pair partitions. This also proves
the asserted Galois equivariance.

Consequently $G_{\mathbb Q}$ acts trivially on
\[
 \Delta_T/2\Delta_T
   \simeq H_1(T_{\mathbb C},\mathbb F_2)
   \simeq(\mathbb F_2)^2 .
\]
The comparison here is finite-cover comparison, with
the chosen complex embedding and geometric base point.
The action on the abelian quotient is independent of
the base-point path.

Complex conjugation acts freely and reverses
orientation on the torus $T(\mathbb C)$. Its quotient
is a Klein bottle: it is a closed nonorientable surface
with Euler characteristic zero, and its orientation
double cover is the torus. Thus its topological
arithmetic extension has presentation
\[
 E_{\mathbb R}^{\rm top}
   =\langle a,b\mid aba^{-1}=b^{-1}\rangle,\qquad
 \Delta_T^{\rm top}=\langle a^2,b\rangle\simeq\mathbb Z^2,
 \qquad a\longmapsto c,\ b\longmapsto1.
 \tag{D11.2}
\]
Completion identifies this extension with the etale
extension over $\mathbb R$. Define the two continuous
characters of the geometric group by
\[
\begin{array}{c|cc}
 &a^2&b\\ \hline
 \chi_1&1&0\\
 \chi_2&1&1
\end{array}
 \qquad\text{with values in }\mathbb F_2 .
 \tag{D11.3}
\]
They are independent and are invariant under all of
$G_{\mathbb Q}$ by the preceding two-torsion calculation.
Their kernels are therefore closed normal subgroups
of $E_T$. The quotients
\[
 E_T/\ker\chi_j\longrightarrow G_{\mathbb Q}
 \tag{D11.4}
\]
have kernel $C_2$ and trivial action on that kernel.
After restriction to $G_{\mathbb R}=C_2$, each is
$C_4\to C_2$: the lift $a$ of conjugation has square
equal to the nonidentity geometric element.
Every lift of conjugation has order four, so neither
restricted extension has a section. Restriction of a
global section would give such a section; hence
neither extension in (D11.4) has a global section.

The map
\[
 q:C\longrightarrow T,\qquad (u,y)\longmapsto(u^3,y)
 \tag{D11.5}
\]
extends to the smooth projective curves and has degree
three. The degree-twelve polynomial in (D11.1) is
squarefree, so $C$ is smooth of genus five. Pullback
on geometric $H^1(-,\mathbb F_2)$ is injective: the
trace map composed with pullback is multiplication
by $\deg q=3$, hence the identity over $\mathbb F_2$.
This trace identity follows as well from topological
Poincare duality and the degree of the continuous map
of compact oriented surfaces.

Pull back $\chi_1,\chi_2$ along the arithmetic
fundamental-group map induced by (D11.5), and denote
their independent geometric characters by
$\chi_1^C,\chi_2^C$. Put
\[
 Q_j=E_C/\ker\chi_j^C .
 \tag{D11.6}
\]
The map $q_*$ induces an isomorphism from each $Q_j$
to (D11.4): it is the identity on $G_{\mathbb Q}$ and
an isomorphism on their geometric kernels of order
two. An injective and surjective map of these
extensions follows directly by lifting an arithmetic
element and correcting it in the kernel. Thus the
two $Q_j$ have no global sections, as asserted.

Let $E_C\twoheadrightarrow Q$ be a common further
quotient through both maps in (D11.6). Its geometric
kernel is a quotient of each $C_2$. A nontrivial
such kernel would make both maps $C_2\to\ker(Q\to
G_{\mathbb Q})$ isomorphisms. The two composite
geometric characters would then have the same
kernel. This contradicts the independence in
(D11.3) and the injective pullback in (D11.5).
The geometric kernel of $Q$ is therefore trivial,
and the indicated map identifies $Q$ with
$G_{\mathbb Q}$. This quotient has its identity
section. There is no point of $C(\mathbb R)$ that
can induce its real localization, because
$C(\mathbb R)$ maps to the empty set $T(\mathbb R)$.
It fails $\mathcal P_C$, proving the nonexistence
of a least quotient.

For the additional finite-place assertion we verify
$T(\mathbb Q_2)=\varnothing$. For $x$ a two-adic unit,
$-(x^2+1)(x^2+4)$ has valuation one. For
$v_2(x)=1$ it has valuation three. For
$x=4z$ its quotient by four is the odd unit
\[
 -(16z^2+1)(4z^2+1)\equiv3\text{ or }7\pmod8 .
\]
For $v_2(x)<0$, remove the square $x^4$; writing
$z=x^{-1}\in2\mathbb Z_2$, the remaining unit is
\[
 -(1+z^2)(1+4z^2)\equiv3\text{ or }7\pmod8 .
\]
These cases also include $x=0$ in the third formula.
An odd two-adic unit is a square precisely when it is
one modulo eight. The points at infinity again
require a square root of $-1$. This proves the
claim, and (D11.5) gives $C(\mathbb Q_2)=\varnothing$.

The exact program checks squarefreeness, all residue
cases, the three branch partitions, and the finite
quotient of (D11.2)
\[
 C_4\times C_2,\qquad
 (a,b)\longmapsto a,\quad (a,b)\longmapsto a+2b
 \quad\text{in }C_4.
\]
Their two kernels generate the orientation kernel.
This is the finite certificate of the common-quotient
calculation; the global maps are (D11.4)--(D11.6).
The passage being answered is \S8, folio 33, scan 65,
with the explicit quotient follow-up in D(c).
\end{proof}



% ---- E ----
\section{Fermat: the quotient in the manuscript and the Selmer calculation}
\label{sec:E}
We use the homogeneous plus convention
\[
 X_n:\ X^n+Y^n+Z^n=0,\qquad U_n=X_n\setminus\{XYZ=0\}.
\]
The morphism $[X:Y:Z]\mapsto[X^{n/m}:Y^{n/m}:Z^{n/m}]$
is a finite étale morphism $U_n\to U_m$ for $m\mid n$.
Its extension to the compactifications is ramified only above the
boundary. For odd $n$, $[X:Y:Z]\mapsto[X:Y:-Z]$ is the
$\mathbb Q$-isomorphism to $X^n+Y^n=Z^n$. For even $n$,
the same formula has target $X^n+Y^n+Z^n=0$ again; the
isomorphism to the minus convention requires a scalar
$\alpha$ with $\alpha^n=-1$ and the map $Z\mapsto\alpha Z$.
We keep these specified base fields throughout.

\subsection{E(a): extensions (371)--(379)}

\begin{statement}{E.1}{Proved}{the three comparison maps}
\label{st:E.1}
\label{E-three-maps}
Put $B=\mathbb P^1_{\mathbb Q}\setminus\{0,1,\infty\}$,
$F=\pi_1(B_{\overline{\mathbb Q}})$, and
$N_n=\pi_1(U_{n,\overline{\mathbb Q}})$, with compatible
geometric base points for the cover
$t=-X^n/Z^n$. Then
\[
 N_n=\ker\bigl(F\longrightarrow(\mathbb Z/n)^2\bigr),
 \qquad F'=\overline{[F,F]}\subset N_n.
\]
The extensions (371), (372), and the top row of (373), in
\S51.16, folios 733v--735, scans 172--175, are respectively
\begin{align*}
 1&\longrightarrow N_n\longrightarrow\pi_1(U_n)
                 \longrightarrow G_{\mathbb Q}\longrightarrow1,\\
 1&\longrightarrow N_n\longrightarrow\pi_1([U_n/S_3])
                 \longrightarrow G_{\mathbb Q}\times S_3\longrightarrow1,\\
 1&\longrightarrow N_n/F'\longrightarrow\pi_1([U_n/S_3])/F'
                 \longrightarrow G_{\mathbb Q}\times S_3\longrightarrow1.
\end{align*}
Here $S_3$ permutes homogeneous coordinates, and the second
quotient is the fundamental group of $BS_3\times\operatorname{Spec}\mathbb Q$.
The geometric kernel in the last row is
\[
 N_n/F'\simeq n\widehat{\mathbb Z}(1)^2.
\]
The intrinsic geometric abelianization and metabelian quotient
are connected to this kernel by the continuous epimorphisms
\[
 N_n/N_n''\longrightarrow N_n/N_n'\longrightarrow N_n/F'.
\]
The middle group has rank $n^2+1$ as a
$\widehat{\mathbb Z}$-module. Thus the precise quotient in
(373) is the rank-two toric quotient. This corrects the
identification of that row in Lead 4.
\end{statement}
\begin{proof}
Over $\overline{\mathbb Q}$ the cover adjoins $n$th roots of
$t$ and $1-t$, up to the fixed constants $-1$.
Its deck group is $(\mathbb Z/n)^2$ and its monodromy is
the reduction of $F^{\mathrm{ab}}=\widehat{\mathbb Z}(1)^2$.
This proves the kernel formula and the last kernel calculation.
Riemann--Hurwitz for a degree $n^2$ cover with three sets of
$n$ boundary points gives
\[
 2-2g(X_n)-3n=-n^2,\qquad
 g(X_n)=\frac{(n-1)(n-2)}2.
\]
The punctured surface is free of rank
$2g(X_n)+3n-1=n^2+1$. Finite étale covering theory gives
the first row. Applying it to the $S_3$-torsor
$U_n\to[U_n/S_3]$ and then to the structure map to
$BS_3$ gives the second row; it records the permutation
as well as the arithmetic action. The subgroup $F'$
is preserved by both actions, so quotienting gives
the third row. Finally $N_n''\subset N_n'\subset F'$
proves the two displayed epimorphisms. They state
exactly which information is removed, without identifying
their kernels. This checks LM7 E.13 and E.16 at the
steps used here.
\end{proof}

\begin{statement}{E.2}{Known}{sections of the minimum quotient}
\label{st:E.2}
\label{E-toric}
Let $T=\mathbb G_m^3/\mathbb G_m$, with coordinates
$u=X/Z$, $v=Y/Z$. The pushout of (371) by
$N_n\to N_n/F'$ is isomorphic, as an extension over
$G_{\mathbb Q}$, to $\pi_1(T)\to G_{\mathbb Q}$.
Its section classes form
\[
 H^1(\mathbb Q,\widehat{\mathbb Z}(1)^2)
       =(\widehat{\mathbb Q^\times})^2,
\]
where
\[
 \widehat{\mathbb Q^\times}
 =\{\pm1\}\times
 \left\{(a_p)_p\in\prod_p\widehat{\mathbb Z}:
 \ \text{for every }m,\ a_p\in m\widehat{\mathbb Z}
       \text{ for all but finitely many }p\right\}.
\]
The map to the base tripod's toric quotient sends a class
$(a,b)$ to $(-a^n,-b^n)$ in multiplicative Kummer notation.
For each $n$ the restriction to actual rational points of
$U_n$ is injective. The image contains formal section classes
independently of the equation defining $U_n$.
\end{statement}
\begin{proof}
The inclusion $U_n\hookrightarrow T$ induces on geometric
groups the epimorphism sending small loops to the valuations
of $u,v$. After the identification of $N_n/F'$ with
$n\widehat{\mathbb Z}^2$, division by $n$ identifies this
map with the identity of $\widehat{\mathbb Z}^2$.
Consequently its induced morphism of quotient extensions
is an isomorphism: it is the identity on the arithmetic
quotient and an isomorphism on kernels. The identity
of $T(\mathbb Q)$ splits the target extension.
Kummer theory gives
$H^1(\mathbb Q,\mu_m)=\mathbb Q^\times/\mathbb Q^{\times m}$;
the finite groups $H^0(\mathbb Q,\mu_m)$ have the
Mittag--Leffler property, so the inverse limit has no
extra $\varprojlim^1$ term. Prime factorization gives the
displayed restricted product, since every class modulo
$m$ has finite support in its valuations.
The coordinate morphism to the base torus is
$(u,v)\mapsto(-u^n,-v^n)$, which proves the arithmetic
translation as well as the multiplication on kernels.
The map $\mathbb Q^\times\to\widehat{\mathbb Q^\times}$
is injective by valuations and sign; hence the two
coordinates recover every rational point. This is the
complete calculation of (377)--(379), scans 176--178.
It independently checks the proof of LM7 E.16.
\end{proof}

\begin{statement}{E.3}{Proved}{equivariant splitting and its ambiguity}
\label{st:E.3}
\label{E-equivariant}
Write $M=\widehat{\mathbb Z}^{\,3}/\widehat{\mathbb Z}(1,1,1)$,
with the permutation action of $S_3$. The arithmetic action
on its Tate twist is scalar $\chi$.
The extension in (373) has a splitting over
$G_{\mathbb Q}\times S_3$, and any two such splittings
are conjugate by a unique element of its kernel.
Compatible choices identify its level-$n$ kernels with $nM(1)$
as in (375)--(376).
\end{statement}
\begin{proof}
We give the cohomological verification implicit in scan 175.
The torus $T$ has identity fixed by $S_3$. Its identity
section, together with the action of $S_3$, is a section
of $\pi_1([T/S_3])\to G_{\mathbb Q}\times S_3$.
The equivariant inclusion $U_n\to T$ induces an isomorphism
from the quotient of (372) to this extension, after
identifying $nM(1)$ with $M(1)$ by division by $n$.
This proves existence at every level.

For uniqueness, $M^{S_3}=0$ and $H^1(S_3,M)=0$.
Indeed in
$0\to\widehat{\mathbb Z}\to\widehat{\mathbb Z}^{\,3}
\to M\to0$, invariants give the isomorphism between
the two diagonal copies of $\widehat{\mathbb Z}$;
$H^1$ of either permutation module is zero; and the
restriction
$H^2(S_3,\widehat{\mathbb Z})\to
 H^2(S_2,\widehat{\mathbb Z})$
is an isomorphism $\mathbb Z/2\to\mathbb Z/2$.
The latter assertion follows from
$H^2(G,\widehat{\mathbb Z})
=\operatorname{Hom}(G,\mathbb Q/\mathbb Z)$
for a finite group $G$, with the finite primary
components understood. The inflation--restriction
sequence now gives
$H^1(G_{\mathbb Q}\times S_3,M(1))=0$.
The invariant module of the product group is zero,
so the conjugating element is unique.
The same calculation holds for $nM(1)$.
For a divisibility tower the unique conjugacies commute
with transition maps. Equivalently the coordinate maps
between the tori are $n/m$ power maps and preserve the
identity. This supplies compatible splittings and
proves exactly the ambiguity asserted on scan 175.
\end{proof}

\begin{statement}{E.4}{Proved}{the implication proposed on scan 173}
\label{st:E.4}
\label{E-fermat-implication}
For odd $n\geq3$ the rational boundary of $X_n$ consists
of $[0:1:-1]$, $[1:0:-1]$, and $[1:-1:0]$.
Their full cuspidal packets are disjoint from all point
sections. Therefore a proof that all sections of (371)
belong to these packets proves $U_n(\mathbb Q)=\varnothing$.
No surjectivity assertion for point sections is needed
for that implication.
\end{statement}
\begin{proof}
A rational $n$th root of $-1$ is $-1$ for odd $n$;
setting one coordinate equal to zero gives exactly
the three displayed points. A rational point of $U_n$
always induces a section by functoriality. Stix's
Theorem 14, with the cusp decomposition embeddings,
separates this section from every packet; the
hyperbolicity condition holds already for $n=1$.
Thus a point would contradict the stated description
of the full section set. This is a proof of the
logical implication in \S51.16, folio 734, scan 173,
not an invocation of the section conjecture.
\end{proof}

\subsection{E(b): intrinsic quotients and the remaining tower problem}

\begin{statement}{E.5}{Known}{the even plus curves}
\label{st:E.5}
\label{E-even}
For every positive even $n$, the intrinsic geometrically
abelianized extension of $U_n/\mathbb Q$ has no section.
The same holds for every quotient mapping onto its
geometric pro-$2$ abelianization, including the full
metabelian and full profinite extensions. The first
intrinsic nilpotent obstruction is at class one.
More precisely, for complex conjugation $c$,
\[
 H_1(U_n(\mathbb C),\mathbb Z_2)
 \simeq\mathbb Z_2\oplus
       \mathbb Z_2[C_2]^{\,n^2/2}
 \quad\text{as a }\mathbb Z_2[C_2]\text{-module},
\]
and the real extension represents the nonzero element
of $H^2(C_2,H_1)=\mathbb Z/2$.
\end{statement}
\begin{proof}
We checked LM7 E.18 and E.25 and retain their actual
Schreier calculation. Positivity gives $X_n(\mathbb R)=\varnothing$.
The free antiholomorphic action on the punctured complex
surface has a nonorientable punctured quotient of Euler
characteristic $-n^2/2$. Its fundamental group $J$ is
free of rank $1+n^2/2$. Its orientation kernel $N$
has index two and rank $n^2+1$. Choose a free basis
$t,a_1,\dots,a_{n^2/2}$ for which only $t$ maps
nontrivially to $C_2$. Then
\[
 t^2,\quad a_i,\quad t a_i t^{-1}
\]
is a Schreier basis of $N$. Conjugation by $t$ fixes
the first abelian generator and swaps each remaining
pair. The square of the chosen lift of $c$ is the
first generator. Replacing $t$ by $bt$ changes this
square in the abelian kernel by $(1+c)[b]$.
Its first coordinate is always $1+2b_0$, which cannot
vanish in $\mathbb Z_2$. On a swapped pair the norm
surjects onto the invariant diagonal, so only the
first coordinate contributes to $H^2$.
Equivariant Riemann existence identifies the completion
of this finite-index extension with $\pi_1(U_{n,\mathbb R})$.
Restriction of a rational section to $G_{\mathbb R}$
and then projection to this module gives the contradiction.
Our exact program checks the ranks for $n=2,4,6,8,10$;
the Schreier calculation proves all even $n$.

The toric pushout of Proposition~\ref{E-toric} sends
this obstruction to zero: $c$ acts by $-1$ on the
torsion-free toric module, whose invariants are zero.
This explicitly relates the split extension (373)
to the obstructed intrinsic abelian extension.
\end{proof}

\subsection{E(c): the missing exponents in the Selmer theorem}

We call a section Selmer when at \emph{every} completion
its class is point-induced or cuspidal. At a real
place the point-induced class means the component
class. This is Steklov's definition in \S3.3
of~\cite{Steklov2026}; it includes $2$ and the
coefficient prime. We first prove the Selmer result,
which uses no local section conjecture.

\begin{statement}{E.6}{Proved}{finite support forces a global point or cusp}
\label{st:E.6}
\label{E-finite-support}
Let $U/\mathbb Q$ be a hyperbolic curve with smooth
compactification $X$, and let $s$ be a Selmer section.
Let $Z\subset X$ be a finite reduced subscheme.
When all finite localizations of $s$ have underlying
points in $Z(\mathbb Q_v)$, the section belongs to
a point or cusp above an element of $Z(\mathbb Q)$.
In particular, a nonconstant morphism
$f:X\to W$ to a proper curve satisfying
$W(\mathbb A_{\mathbb Q})_\bullet^{\mathrm{f-cov}}=W(\mathbb Q)$
makes every Selmer section of $U$ point-induced or cuspidal.
\end{statement}
\begin{proof}
A neighbourhood of $s$ is a finite étale cover
$U'\to U$ together with a lift of $s$.
Every lifted section is Selmer. To see the latter
directly, a point section in the subgroup corresponding
to the cover fixes a point of the finite fibre,
and so lifts to a rational point. For a cusp use the
intersection of its decomposition group with that
subgroup; its arithmetic quotient is surjective
because it contains the local section, so the
associated lifted cusp is rational.
Write $X'$ for the smooth compactification. We may
refine to genus at least one. Such refinements are
cofinal: take a characteristic geometric cover of
sufficiently large degree, spread it over a finite
extension, and take its neighbourhood twist determined
by the section. Riemann--Hurwitz, with boundary
inertia indices made arbitrarily large, gives the
required genus.

The localizations determine an adelic point of $X'$
surviving finite descent, by Harari--Stix,
Theorem 2.1. At every finite place its component
lies in the finite inverse image $Z'\subset X'$.
Stoll's Theorem 8.2 gives a rational point of $Z'$.
The system of these finite nonempty sets over all
neighbourhoods has the finite intersection property:
a common refinement maps to any finite collection
of levels. Compactness supplies a compatible
rational pro-point. Its stabilizer contains the
section; over an interior point this is the point
decomposition group, and over a boundary point it
is the cuspidal decomposition group. This proves
the first assertion.

For the second, functoriality of finite descent
sends the adelic localization to one point
$w\in W(\mathbb Q)$. The finite fibre $f^{-1}(w)$
is a suitable $Z$.
We use only finite places in this last argument.
This is the correction specified in Stoll's erratum
to Proposition 8.5; replacing a real component by
a point of the scheme-theoretic fibre would be
unjustified.
\end{proof}

\begin{statement}{E.7}{Proved}{full level two supplies level $4\ell$}
\label{st:E.7}
\label{E-level-four}
Let $Z_\ell\to B$ be the finite étale cover parametrizing
a cyclic subgroup $C\subset E_\lambda[\ell]$ of the
Legendre family, for an odd prime $\ell$. There is
a nonconstant $\mathbb Q$-morphism
\[
 Z_\ell\longrightarrow Y_0(4\ell)
\]
whose extension to smooth compactifications is finite.
Every component on which an arithmetic section lifts
is geometrically connected over $\mathbb Q$.
For $\ell\geq5$ its compactification has a
nonconstant map to a positive-genus modular curve.
\end{statement}
\begin{proof}
The ordered points $P=(0,0)$ and $Q=(1,0)$ of order
two are independent on $E=E_\lambda$.
Let $\phi:E\to E'=E/\langle P\rangle$ be the quotient
and let $\phi':E'\to E$ be its dual, so
$\phi'\phi=[2]$ and $\phi\phi'=[2]$.
Define
\[
 H_4=(\phi')^{-1}(\langle Q\rangle)\subset E'.
\]
It has order four. For a point $R$ with
$\phi'(R)=Q$, we have $2R=\phi(Q)\ne0$,
so $R$ has order four and generates $H_4$.
Thus $H_4$ is cyclic, defined over the base, and
$\phi(C)$ has order $\ell$. Their sum is a cyclic
subgroup of order $4\ell$, giving the asserted
moduli map. All constructions commute with
base change; the finite étale subgroup scheme
$C$ is the tautological one on $Z_\ell$.
The function $j(E')$ is nonconstant: a fixed elliptic
curve has only finitely many degree-two isogenous
curves, whereas $j(E_\lambda)$ is nonconstant.
Therefore the map of smooth projective curves
is nonconstant and finite.

For completeness, a connected arithmetic cover
supporting a section onto $G_{\mathbb Q}$ has
constant field $\mathbb Q$: otherwise its
fundamental group's arithmetic image is the
proper subgroup attached to the constant field.
For $4\ell$ the modular index is $6(\ell+1)$,
there are six cusps, and there are no elliptic
points of orders two or three. Consequently
\[
 g(X_0(4\ell))=1+\frac{6(\ell+1)}{12}-\frac62
             =\frac{\ell-3}{2}.
\]
At $\ell=5$ the target is $X_0(20)$ of genus one.
The exact program checks these indices and genera
at the first seven odd primes. The displayed
formula proves the general assertion.
\end{proof}

\begin{statement}{E.8}{Known}{the local residual representations}
\label{st:E.8}
\label{E-local-representations}
Let $\ell\geq5$ be prime and let $s$ be a Selmer
section of $U_\ell/\mathbb Q$. Let
$\bar\rho_s:G_{\mathbb Q}\to\mathrm{GL}_2(\mathbb F_\ell)$
be obtained from the Legendre family.
It is reducible. There is a quadratic twist whose
semisimplification is $1\oplus\bar\chi_\ell$.
The mod-two representation is trivial before
and after that twist.
\end{statement}
\begin{proof}
We checked Steklov's Lemmas 5.4, 5.7, 5.11--5.12 and
Propositions 5.5, 5.9--5.10, 5.13, including their
uses at $2$ and $\ell$.
Here are the arithmetic steps needed for the
extension to $\ell=5$.
For a local point write
$a^\ell+b^\ell+1=0$ and $\lambda=-a^\ell$.
The formula
\[
 j=2^8\frac{(\lambda^2-\lambda+1)^3}
                {\lambda^2(\lambda-1)^2}
\]
gives $v_p(j)=8v_p(2)-2\ell|v_p(a)|$ when
$v_p(a)\ne0$. When $v_p(a)=0$ and $\lambda\equiv1$
use $1-\lambda=-b^\ell$ instead. At $2$ one of
these cases has positive absolute valuation;
$\ell\geq5$ makes $v_2(j)<0$. At odd $p$ there
is either good reduction or potential multiplicative
reduction with $\ell\mid v_p(j)$.
After a local quadratic twist, the latter is a
Tate curve. At $p\ne\ell$ its residual extension
class is the Kummer class of $q$, so divisibility
of $v_p(q)$ by $\ell$ makes it unramified.
At $p=\ell$, the same divisibility puts $q$
modulo $\ell$th powers in the unit subgroup.
The residual representation is finite flat,
hence has Serre weight two; good reduction has
the same weight. At $2$ the untwisted Tate
representation is tame, with conductor exponent
at most two (the bound four used by Steklov
also suffices).

For a local cusp, the map to the $j$-line has
ramification index divisible by $\ell$.
The Tate parameter class is therefore divisible
by $\ell$. Its residual representation is a
quadratic twist of $1\oplus\bar\chi_\ell$.
In particular the same local conditions hold,
including at $2$ and $\ell$.

Suppose $\bar\rho_s$ were irreducible.
Its determinant is $\bar\chi_\ell$, so complex
conjugation has distinct eigenvalues $1,-1$.
Any invariant line over $\overline{\mathbb F}_\ell$
would be one of those two lines, already defined
over $\mathbb F_\ell$; hence irreducibility is
absolute. The local quadratic characters on
inertia patch to a global quadratic character:
the abelian cyclotomic isomorphism
$G_{\mathbb Q}^{\mathrm{ab}}\simeq
 \prod_p\mathbb Z_p^\times$
allows us to multiply their finitely many
nontrivial inertia characters. The twisted
representation has determinant $\bar\chi_\ell$,
weight two, and level dividing $16$.
The Khare--Wintenberger theorem then produces
a weight-two cusp form of level dividing $16$.
Each such $X_0(N)$ has genus zero, a contradiction.

Thus the representation is reducible.
At $\ell$, a good supersingular representation
over $\mathbb Q_\ell$ has irreducible inertial
action, so this case is excluded. The local
semisimplifications are quadratic twists of
$1\oplus\bar\chi_\ell$. Patch again and write
the two global characters as $\omega$ and
$\bar\chi_\ell\omega^{-1}$.
They are unramified away from $\ell$, and one
of them is unramified at $\ell$. A finite
extension of $\mathbb Q$ unramified at all
finite primes is trivial by Minkowski's
discriminant theorem. Hence one character is
$1$ and the other $\bar\chi_\ell$.
Finally the three Legendre points of order
two are rational, and a quadratic twist acts
as $[-1]$, which is the identity on $E[2]$.
\end{proof}

\begin{statement}{E.9}{Proved}{the Selmer theorem includes exponent five}
\label{st:E.9}
\label{E-selmer-primes}
For every prime $\ell\geq5$, every Selmer section
of $U_\ell/\mathbb Q$ is cuspidal.
\end{statement}
\begin{proof}
The invariant line of Lemma~\ref{E-local-representations}
lifts the pushed-forward section on $B$ to the
cover $Z_\ell$ of Proposition~\ref{E-level-four}.
The lift is Selmer by the finite-fibre argument
in Lemma~\ref{E-finite-support}.
Its compactification maps nonconstantly to
$X_0(4\ell)$, whose positive genus and Stoll's
Corollary 8.8 give
\[
 X_0(4\ell)(\mathbb A_{\mathbb Q})_\bullet^{\mathrm{f-cov}}
        =X_0(4\ell)(\mathbb Q).
\]
Lemma~\ref{E-finite-support} therefore makes the
lift, and hence the section on $B$, point-induced
or cuspidal. A point-induced section on $B$
lifts to a rational point of $U_\ell$ whenever
its image is contained in the subgroup
$\pi_1(U_\ell)$: it fixes an element of the
finite étale fibre. A cuspidal section lifts
into the intersection with a cusp decomposition
group and is cuspidal upstairs.

It remains to exclude the point case.
The associated rational Legendre elliptic
curve has full rational two-torsion.
After the quadratic twist of
Lemma~\ref{E-local-representations}, its
$\ell$-torsion has a submodule $1$ or
$\bar\chi_\ell$. In the first case it has a
rational point of order $\ell$. In the second,
quotient by the subgroup $\mu_\ell$; the
dual isogeny's kernel is the constant group
$\mathbb Z/\ell$, so the quotient elliptic
curve has a rational point of order $\ell$.
The odd-degree isogeny is an isomorphism on
two-torsion. Thus one rational elliptic curve
has a subgroup
$\mathbb Z/2\oplus\mathbb Z/(2\ell)$ in its
rational torsion. Mazur's Theorem (8)
allows such a subgroup only for $\ell\leq4$.
This contradiction proves the assertion.

The point-case exclusion also repairs the
last sentence of Steklov's Theorem 5.14.
His claim that $Y_1(\ell)(\mathbb Q)$ is empty
for $\ell\geq7$ is false at seven.
The necessary extra datum is the full
two-torsion, retained throughout this proof.
\end{proof}

\begin{statement}{E.10}{Refuted}{the claimed emptiness at seven}
\label{st:E.10}
\label{E-seven-error}
The claim $Y_1(7)(\mathbb Q)=\varnothing$ used in the
last sentence of Steklov's Theorem 5.14 is false.
The elliptic curve
\[
 E:\ y^2-xy-4y=x^3-4x^2
\]
has discriminant $-1664$ and the point $(0,0)$
of exact order seven. The simultaneous requirement
of full rational two-torsion and rational
seven-torsion is impossible by Mazur's Theorem (8).
\end{statement}
\begin{proof}
The successive nonzero multiples of $(0,0)$ are
\[
 (0,0),\ (4,8),\ (2,2),\ (2,4),\ (4,0),\ (0,4);
\]
the seventh multiple is the origin.
The included exact rational-arithmetic program
checks the curve equation at every step, the
addition law, and the nonzero discriminant.
The last assertion follows from the forbidden
subgroup $\mathbb Z/2\oplus\mathbb Z/14$.
This is precisely the replacement used in
Theorem~\ref{E-selmer-primes}.
\end{proof}

\begin{statement}{E.11}{Proved}{exponent three}
\label{st:E.11}
\label{E-selmer-three}
Every Selmer section of $U_3/\mathbb Q$ is cuspidal.
\end{statement}
\begin{proof}
On $X_3$ the functions
\[
 u=-\frac{3Z}{X+Y},\qquad
 v=\frac{9(X-Y)/(X+Y)-1}{2}
\]
give a birational map to $E:v^2+v=u^3-7$.
Indeed multiplication by $(X+Y)^3$ turns
$v^2+v-u^3+7$ into
$27(X^3+Y^3+Z^3)$, as the exact symbolic
calculation verifies. The smooth projective
curves therefore are isomorphic. The map sends
the three rational boundary points to
$(3,-5)$, $(3,4)$, and the origin.

The change of variables
$(x,y)=(4u,8v+4)$ gives $y^2=x^3-432$.
The identification of this curve with $X_0(27)$
is explicit in Kriz--Li, \S10.1~\cite{KrizLi2019};
it is the curve $27a1$. Stoll's Corollary 8.8
gives its finite-descent equality and rank zero.
For the rational torsion, good reduction at
$5$ and $7$ gives respectively $6$ and $9$
points. Prime-to-residue-characteristic
torsion injects under reduction. Applying
the two reductions prime by prime bounds
the rational torsion by three; the three
displayed points give equality. Thus
$X_3(\mathbb Q)$ consists exactly of these
boundary points. A Selmer section localizes
to their diagonal image under the
finite-descent equality. Lemma~\ref{E-finite-support}
then puts it into a boundary packet.
\end{proof}

\begin{statement}{E.12}{Proved}{all exponents and the stipulated local input}
\label{st:E.12}
\label{E-all-exponents}
For every odd $n\geq3$, all Selmer sections of
$U_n/\mathbb Q$ are cuspidal. For every even
$n\geq2$ there are no sections at all.
Using the local theorem stipulated as Tool (1)
in the development prompt, every section of
$U_n/\mathbb Q$ is Selmer. Consequently,
in the stipulated framework the complete
section set for odd $n\geq3$ is the disjoint
union of the three rational cusp packets,
each a torsor under
$H^1(\mathbb Q,\widehat{\mathbb Z}(1))$.
\end{statement}
\begin{proof}
Choose an odd prime $\ell\mid n$. The finite étale
power map $U_n\to U_\ell$ carries local point
sections to point sections and local cusp
sections into cusp packets, so it preserves
the Selmer property. Its image is cuspidal
by Theorems~\ref{E-selmer-primes} and
\ref{E-selmer-three}. Intersecting its
decomposition group with $\pi_1(U_n)$ gives
a cusp decomposition group whose image in
$G_{\mathbb Q}$ is surjective because it
contains the given section. The cusp is
therefore rational and the section is
cuspidal. The even assertion is
Proposition~\ref{E-even}.

For the final assertion, Tool (1) and the
neighbourhood reduction in Part B give the
point-or-cusp description at every finite
place, with no exclusion of $2$ or $\ell$.
The real section theorem gives the point
component at the infinite place.
These are exactly all the places required
by Steklov's definition. The three packets
exist by choosing local parameters; their
disjointness and their torsor descriptions
are Stix's Theorem 14 and Proposition 8.
This proof identifies where the unrefereed
2026 local input enters. The Selmer theorem
above is independent of that input.
\end{proof}

\paragraph{Provenance of the local input.}
The final passage in Corollary~\ref{E-all-exponents}
uses the supplied October 2026 manuscript as
the tool explicitly stipulated in the question.
Our checking of its named intermediate steps
is recorded in Part B and the audit.
We do not report an independent verification
of every argument in that manuscript.
The finite-level tower problem in E(b) is
a further task: the morphism from the full
section set to a quotient's section set need
not be surjective. The next subsection keeps
its image and its obstruction groups explicit.


% ---- E_density ----
\subsection*{E(iv). The exact quartic local-order law}

\begin{statement}{E.13}{Proved}{The full CM image and the affine image of every primitive class}
\label{st:E.13}

Let $E_+ /\mathbb Q$ be $y^2=x^3+x$, let $\ell\in\{3,5\}$,
put $V=T_\ell E_+$ and $K=\mathbb Q(i)$, and write
$\rho:G_{\mathbb Q}\to\operatorname{GL}(V)$ for the actual Tate action.
Then
\[
 \rho(G_K)=C_\ell,\qquad
 \rho(G_{\mathbb Q})=N_\ell=C_\ell\rtimes\langle c\rangle,
 \tag{E13.1}
\]
where $C_5=(\mathbb Z_5^\times)^2$ acts on the two integral CM
lines, $C_3=\mathbb Z_3[i]^\times$ acts on the rank-one
$\mathbb Z_3[i]$-module $V$, and $c$ interchanges the two lines at
five and acts by conjugation at three.  This is the full normalizer
of the displayed Cartan subgroup in $\operatorname{GL}(V)$.

Every continuous $V$-valued cocycle factors through
$G_{\mathbb Q,S}$ for $S=\{2,\ell,\infty\}$.
The module $H^1(G_{\mathbb Q},V)$ is finite free over
$\mathbb Z_\ell$ of positive rank.  For every primitive class in
this module and every continuous cocycle $h$ representing it, the
actual homomorphism
\[
 G_{\mathbb Q}\longrightarrow V\rtimes N_\ell,
 \qquad g\longmapsto(h(g),\rho(g))
 \tag{E13.2}
\]
is surjective.  Here primitive means not divisible by $\ell$ in
$H^1(G_{\mathbb Q},V)$.  No class-number calculation or assertion
that this cohomology module has rank one is needed for these conclusions.
\end{statement}
\begin{proof}
The endomorphism $I:(x,y)\mapsto(-x,iy)$ is defined over $K$
and satisfies $I^2=-1$.  Its determinant on a Tate module is one,
by the Weil pairing.  At five its two eigenvalues are therefore
the distinct roots of $T^2+1$, each with multiplicity one; the
projectors $(1\pm I/i_5)/2$ are integral.  At three the same
polynomial is irreducible modulo three, and $V$ is free of rank
one over $\mathbb Z_3[i]$.  Thus the centralizer of $I$ is $C_\ell$.
For $g\in G_{\mathbb Q}$ the actual endomorphism relation is
\[
       \rho(g)I\rho(g)^{-1}=I\quad(g\in G_K),\qquad
       \rho(g)I\rho(g)^{-1}=-I\quad(g\notin G_K).
 \tag{E13.3}
\]
This places the image in the claimed normalizer and identifies
its quotient of order two with $\operatorname{Gal}(K/\mathbb Q)$.
To check that the normalizer named here is the full matrix
normalizer, the units of $\mathbb Z_\ell[i]$ generate that
quadratic etale algebra: at three use $1,i$, and at five use
$(1,1)$ and $(2,1)$ in its split realization.
A matrix normalizing the unit group therefore normalizes this
algebra.  Its induced algebra automorphism is the identity or
conjugation; the kernel is the unit centralizer $C_\ell$.
Both automorphisms are represented, giving exactly $N_\ell$.

We prove the full integral Cartan image by local inertia.
At five take the Gaussian primes $\pi=1+2i$ and
$\bar\pi=1-2i$.  At the place $\pi$ of $K$, the completion is
$\mathbb Q_5$ and $\bar\pi$ is a unit.  The curve has good
reduction, since its discriminant is $-64$.  The isogeny
$[\bar\pi]$ extends to the smooth proper elliptic model and its
differential is multiplication by the unit $\bar\pi$.
Its kernel and the kernels of all its iterates are consequently
finite etale over the valuation ring.  Inertia acts trivially on
their Tate line.  Bezout for the coprime Gaussian ideals
$(\pi)$ and $(\bar\pi)$ gives the actual decomposition
\[
 E_+[5^n]=E_+[\pi^n]\oplus E_+[\bar\pi^n].
 \tag{E13.4}
\]
The determinant of inertia on these two factors is cyclotomic.
The $5$-power cyclotomic extensions of $\mathbb Q_5$ are totally
ramified, with full unit Galois group.  Hence this inertia image
is $\mathbb Z_5^\times\times\{1\}$ in one ordering of (E13.4).
Complex conjugation supplies $\{1\}\times\mathbb Z_5^\times$.
Both inertia groups lie in $G_K$, proving $\rho(G_K)=C_5$.

At three the completion $K_3=\mathbb Q_3(i)$ is the unramified
quadratic extension, with residue field of order nine.  The
reduction of $E_+$ satisfies the literal endomorphism identity
\[
          [-3](x,y)=(x^9,y^9).
 \tag{E13.5}
\]
For an elementary verification, in characteristic three the
division polynomials give $\psi_3=-1$, $\psi_2=-y$ and
$\psi_4=y(x^6+2x^4+x^2+2)$.  Therefore
\[
 x([3]P)=x+(x^3+x)(x^6+2x^4+x^2+2)=x^9.
\]
The curve equation makes the other coordinate either $y^9$
or $-y^9$ as a rational function.  The point $(2,1)$ has double
$(0,0)$ and order four over $\mathbb F_3$, so $[3](2,1)=(2,-1)$;
the sign is negative.  This proves (E13.5), without a choice of
a Frobenius eigenvalue.

In the formal parameter $T=-x/y$, the actual series
$f(T)=[-3](T)$ has derivative $-3$ and reduces to $T^9$.
It is a Lubin--Tate series for the uniformizer $-3$ of $K_3$.
The uniqueness statement for a formal group commuting with
$f$ identifies the actual formal group of $E_+$ with the
Lubin--Tate formal group, using the identity parameter map.
The same uniqueness for an endomorphism with prescribed linear
coefficient identifies each actual $[a+bi]$ with the
Lubin--Tate endomorphism for $a+bi$; continuity gives this
identification over all of $\mathbb Z_3[i]$.
Equation (E13.5) also shows that every geometric $3$-power
torsion point reduces to the origin, so all such torsion is in
this formal group.  Lubin--Tate's theorem now gives, at each
level, the totally ramified extension with Galois action
\[
      (\mathbb Z_3[i]/3^n)^\times
             \quad\hbox{on}\quad E_+[3^n].
\]
Its inverse limit is the full $C_3$, already obtained from local
inertia.  The precise formal-group input is
\cite[Lemma 1, Theorem 1, and Theorem 2(b),(d),(e)]{E_density_LT}.
In particular (E13.1) follows at both primes.  Choosing one basis
vector on a CM line and its conjugate gives the stated $c$ at
five.  At three choose a primitive vector in the $+1$ eigenspace
of $c$; it is an integral $\mathbb Z_3[i]$-basis and gives the
stated conjugation action.  These changes of basis are actual
linear isomorphisms of the Tate lattices.

To prove the assertions about cocycles, let $p\notin S$.
The Tate action on inertia $I_p$ is trivial.  A cocycle restricted
to $I_p$ is a homomorphism to the pro-$\ell$ group $V$, hence
factors through the $\ell$-primary tame quotient.
For its value $v$ on a tame generator, Frobenius equivariance
gives $(F_{E,p}-p)v=0$.  But
\[
    \det(F_{E,p}-p)=p\,\#E_+(\mathbb F_p)\ne0.
\]
Thus $v=0$.  The cocycle vanishes on all the corresponding
conjugate inertia groups and on their closed normal closure;
the action also vanishes there.  It factors through $G_{\mathbb Q,S}$.

The finite-coefficient finiteness theorem and the global Euler
characteristic formula give the following facts, with their
hypotheses satisfied since $S$ is finite and contains $\ell$
and the archimedean place:
\cite[I.4.10(c), I.4.15, and I.5.1]{E_density_MilneADT}.
The groups $H^j(G_{\mathbb Q,S},V/\ell^n)$ are finite;
they vanish for $j>2$ because $\ell$ is odd.
Their finite inverse systems have zero derived inverse limit.
The continuous groups $H^1(G_{\mathbb Q,S},V)$ and
$H^2(G_{\mathbb Q,S},V)$ are therefore compact; their reductions
modulo $\ell$ are finite, so compact Nakayama makes them finitely
generated over $\mathbb Z_\ell$.
The residual representation has no invariant vector: at three
multiplication by $i$ has irreducible polynomial, and at five
the full diagonal subgroup and its interchange have no common
fixed vector.  The coefficient exact sequence then proves that
$H^1(G_{\mathbb Q,S},V)$ has no $\ell$-torsion.
Write $r$ and $s$ for the ranks of $H^1$ and $H^2$ respectively.
Complex conjugation has one eigenvalue of each sign, and the
finite Euler formula is
\[
 \frac{\#H^1(G_{\mathbb Q,S},V/\ell^n)}
      {\#H^2(G_{\mathbb Q,S},V/\ell^n)}=\ell^n.
\]
The exact coefficient sequences give
\[
\begin{split}
0&\longrightarrow H^1(V)/\ell^n
 \longrightarrow H^1(V/\ell^n)
 \longrightarrow H^2(V)[\ell^n]\longrightarrow0,\\
H^2(V)/\ell^n&\ \xrightarrow{\sim}\ H^2(V/\ell^n),
\end{split}
\]
where these groups have source $G_{\mathbb Q,S}$.
The torsion contributions of $H^2(V)$ cancel in the ratio.
Consequently $r-s=1$, and $r\ge1$.  Primitive classes exist.

Fix such an $h$.  Its reduction $[h_1]$ in
$H^1(G_{\mathbb Q,S},V/\ell)$ is nonzero by the same exact
coefficient sequence.  Let $\Gamma_n=\rho_n(G_{\mathbb Q,S})$
and $J_n=\ker\rho_n$.  The scalar $-I$ lies in $\Gamma_n$.
For a cocycle $a:\Gamma_n\to V/\ell$, commuting $-I$ with $g$
in the cocycle identity gives
\[
       -2a(g)=(g-1)a(-I).
\]
Since two is a unit, every such $a$ is a coboundary.
Inflation--restriction shows that $h_1|_{J_n}$ is nonzero.
The subgroup $h_n(J_n)\subset V/\ell^n$ is stable under
$\Gamma_n$, and its reduction is nonzero.  This reduction is
the whole $V/\ell$: at three the CM action is irreducible;
at five a stable line for the diagonal subgroup is a coordinate
line and complex conjugation interchanges the two lines.
Nakayama therefore gives $h_n(J_n)=V/\ell^n$.
The image of $(h_n,\rho_n)$ has the full translation kernel
and surjects onto $\Gamma_n$, so is
$(V/\ell^n)\rtimes\Gamma_n$.
The image of (E13.2) is compact and surjects onto every one
of these finite quotients.  It is the entire inverse limit.
\end{proof}

\begin{statement}{E.14}{Proved}{Exact natural densities of all local Kummer orders}
\label{st:E.14}

Retain $E_+$, $\ell$ and any primitive cocycle $h$ of E.13.
For every rational prime $p\ne2,\ell$, choose an arithmetic
Frobenius and define
\[
 u_p=(F_{E,p}-1)^{-1}h(F_p),\qquad
 \operatorname{ord}(u_p+V)=\ell^{m_p},\quad m_p\ge0.
 \tag{E14.1}
\]
This is the local cohomology class under the explicit isomorphism
\[
 V/(F_{E,p}-1)V\ \xrightarrow{\sim}\
 \bigl((V\otimes\mathbb Q_\ell)/V\bigr)^{F_{E,p}},
 \qquad [x]\longmapsto(F_{E,p}-1)^{-1}x+V.
 \tag{E14.2}
\]
The integer $m_p$ is independent of the Frobenius prolongation,
the cocycle representative, the integral Tate basis, and
multiplication of $[h]$ by a unit.

For $q>1$ and $n\ge1$ put
\[
             T_q(n)=\frac{q^{2-n}}{q^2-1}.
\]
The following are natural densities relative to the indicated
set of rational primes; all of them exist:
\[
\begin{array}{c|cc}
 &p\equiv1\pmod4&p\equiv3\pmod4\\\hline
 \ell=3:\ m_p\ge n&T_9(n)&T_3(n)\\
 \ell=5:\ m_p\ge n&2T_5(n)-T_5(n)^2&T_5(n).
\end{array}
 \tag{E14.3}
\]
In particular, among all rational primes, the exact tails are
\[
\begin{split}
 \delta_3(m_p\ge n)&=\frac{9^{2-n}}{160}
                         +\frac{3^{2-n}}{16},\\
 \delta_5(m_p\ge n)&=\frac{5^{2-n}}{16}
                         -\frac{5^{4-2n}}{1152}.
\end{split}
 \tag{E14.4}
\]
The masses at zero and at every positive integer are
\[
\begin{array}{c|cc}
 &\delta(m_p=0)&\delta(m_p=n),\quad n\ge1\\\hline
 \ell=3&121/160&9^{1-n}/20+3^{1-n}/8\\
 \ell=5&817/1152&5^{1-n}/4-5^{2-2n}/48.
\end{array}
 \tag{E14.5}
\]
Each positive exponent has positive density.  For a nonzero scalar
$r\in\mathbb Z_\ell$, with $t=\nu_\ell(r)$, the local exponents
of $r[h]$ are exactly $\max(m_p-t,0)$; their positive tails are
obtained by replacing $n$ with $n+t$ in (E14.3)--(E14.4).
For a compact vector $(r_1[h],r_2[h],r_3[h])$, the analogous
exponent is $\max(m_p-\min_i\nu_\ell(r_i),0)$, with value zero
when every $r_i$ is zero.  These assertions apply to the primitive
class used in the prior E.115 and E.117 and do not require the
prior rank-one cohomology calculation.
\end{statement}
\begin{proof}
First $\det(F_{E,p}-1)=\#E_+(\mathbb F_p)>0$, so every inverse
in (E14.1)--(E14.2) exists.  The inverse to (E14.2) sends a
fixed class $u+V$ to $(F_{E,p}-1)u$ modulo $(F_{E,p}-1)V$.
Both composites are the identity.  Since $h$ is unramified at
$p$, $h(F_p)$ is its unramified cohomology coordinate.

Write the affine group law as
$(x,F)(y,G)=(x+Fy,FG)$.
Conjugation by $(v,B)$ gives
\[
 F'=BFB^{-1},\qquad x'=Bx-(F'-1)v,\qquad
 (F'-1)^{-1}x'=B(F-1)^{-1}x-v.
 \tag{E14.6}
\]
Thus the order of the last vector modulo $V$ is unchanged.
A coboundary adds an integral vector to $(F-1)^{-1}x$, and
a unit scalar preserves its order.  This proves all stated
independences and, in particular, the conjugacy invariance
needed for Chebotarev.

By E.13 the relevant profinite quotient is the full group
$V\rtimes N_\ell$.  Its normalized Haar measure is the product
of additive Haar measure on $V$ and Haar measure on $N_\ell$;
the linear automorphisms preserve the former measure.
The two Cartan cosets have measure one half.  Their arithmetic
quotient is $\operatorname{Gal}(\mathbb Q(i)/\mathbb Q)$, so they
correspond exactly to $p\equiv1$ and $p\equiv3\pmod4$.

Here is the one-coordinate calculation, including its finite
level form.  Let $\mathcal O$ be either $\mathbb Z_\ell$ or
$\mathbb Z_3[i]$, let its residue field have size $q$, and
normalize its valuation by $v(\ell)=1$.  Let $a$ and $x$ be
independent uniform elements of $\mathcal O^\times$ and
$\mathcal O$.  For $j\ge0$ and $k\ge1$,
\[
 \Pr(v(x)=j)=(1-q^{-1})q^{-j},\qquad
 \Pr(v(a-1)\ge k)=\frac{q^{1-k}}{q-1}.
\]
The exponent of $x/(a-1)$ modulo $\mathcal O$ is
$\max(v(a-1)-v(x),0)$ outside a set of measure zero.
For $n\ge1$ its tail is therefore
\[
 \sum_{j\ge0}(1-q^{-1})q^{-j}
          \frac{q^{1-j-n}}{q-1}
   =q^{-n}\sum_{j\ge0}q^{-2j}
   =T_q(n).
 \tag{E14.7}
\]
Modulo $\ell^M$, the exponent of $[x]$ in
$\mathcal O/((a-1),\ell^M)$ has the exact tail
\[
 T_{q,M}(n)=T_q(n)\bigl(1-q^{-2(M-n+1)}\bigr)
       \quad(1\le n\le M),
 \tag{E14.8}
\]
and zero tail for $n>M$.  This is the same sum with
$0\le j\le M-n$.

In the Cartan coset at five, $F=(a,b)$ has two independent
unit coordinates and $x=(x_1,x_2)$ two independent additive
coordinates.  The exponent is the maximum of the two
one-coordinate exponents, giving $2T_5-T_5^2$.
At three the Cartan is the single unit group of the unramified
quadratic ring, so (E14.7) gives $T_9$ directly.

In the other coset write $F=a c$.  Its square is $N(a)$,
its trace is zero and its determinant is $-N(a)$.
Modulo the odd prime $\ell$, $F$ cannot equal the identity:
it anticommutes with the invertible CM operator $I$, and
$2I\ne0$ modulo $\ell$.  At least one entry of $F-1$ is
therefore a unit.  The Smith factors of this matrix over
$\mathbb Z_\ell$ are $1$ and $1-N(a)$, up to units.
For fixed $F$, the image of a uniform $x\in V$ in the
nontrivial Smith coordinate is uniform in $\mathbb Z_\ell$.
Moreover, $N(a)$ is uniform in $\mathbb Z_\ell^\times$.
At five this is the product map on two independent units.
At three, the residue-field norm is onto and the norm on
principal units is onto: the $3$-adic logarithm identifies it
with the trace $\mathbb Z_3[i]\to\mathbb Z_3$, whose value at
$1$ is the unit two.  Haar measure under this surjective
homomorphism is normalized Haar measure.
Consequently the other coset has tail $T_\ell(n)$, and the
same argument modulo $\ell^M$ gives $T_{\ell,M}(n)$.
For the actual arithmetic Frobenius at an inert rational
prime, the determinant identity reads $N(a)=-p$ and
$\det(F-1)=p+1$, retaining the sign in this calculation.

For completeness the passage to prime density uses finite
quotients with a vanishing error, not an interchange of an
uncontrolled infinite union with density.
Let $A_{n,M}$ be the inverse image in $V\rtimes N_\ell$ of
\[
 \ell^{n-1}x\notin(F-1)(V/\ell^M V).
\]
This is a conjugacy-invariant clopen set.  Put
$D_M=\{\det(F-1)\equiv0\pmod{\ell^M}\}$.
For a matrix with $v_\ell(\det(F-1))<M$, both Smith exponents
are less than $M$.  Solvability of the displayed integral
equation is consequently detected modulo $\ell^M$.
Thus, for the actual Frobenius elements,
\[
 A_{n,M}\ \subset\ \{m_p\ge n\}\ \subset\ A_{n,M}\cup D_M.
 \tag{E14.9}
\]
The measures of $D_M$ tend to zero.  In the split Cartan,
the limiting set is $a=1$ or $b=1$; in the nonsplit Cartan
it is $a=1$; and in the other coset it is $N(a)=1$.
Each has measure zero by the product unit measures or the
surjective norm map just proved.  Finite-level Chebotarev
in its natural-density form applies to both clopen sets in
(E14.9) \cite[\S2.1, Theorem 1, and equations (8)--(9)]{E_density_Serre}.
Taking lower and upper natural densities and then $M\to\infty$
proves existence and the asserted Haar values.  The same
sandwich intersected with a Cartan coset proves the relative
densities in (E14.3).

Averaging the two cosets yields (E14.4); subtracting successive
tails gives the positive masses in (E14.5), and subtracting the
first tail from one gives its masses at zero.
Multiplication by $r$ on (E14.2) gives the stated exponent
shift.  For a vector of three scalar multiples, the order in
the direct sum is the maximum of the three scalar orders,
which gives the minimum valuation formula.

The independent finite certificate
\texttt{code/B/cm\_density\_checks.py} enumerates every matrix
in both normalizer cosets and every translation vector modulo
$3$, $3^2$, $3^3$, $5$, and $5^2$.  It forms each actual image
$(F-1)(V/\ell^M V)$ and counts the order of every cokernel
class.  Every count agrees exactly with (E14.8), its
two-coordinate maximum, and the other-coset formula.  Its
output is retained in
\texttt{code/B/cm\_density\_checks.out.json}.
\end{proof}

% ---- E_extra ----
\subsection{The first obstructions and the exact additive defect}

\begin{statement}{E.15}{Proved}{the polarization law for the actual class-two obstruction}
\label{st:E.15}
\label{E-polarization}
Let $\ell$ be odd, let $N$ be a free pro-$\ell$ group of
finite rank with a specified continuous action of a
profinite group $G$, and put
$A=N^{\mathrm{ab}}$, $B=\bigwedge^2 A$.
A choice of free lifts of a basis of $A$ gives the integral
class-two coordinates
\[
 N/\gamma_3N=A\times B,\qquad
 (v,m)(w,n)=(v+w,m+n+\tfrac12v\wedge w).
 \tag{E15.1}
\]
Write the genuine arithmetic action as
\[
 g(v,m)=(A_gv,\ B_gm+D_gv),\qquad
 B_g=\bigwedge^2 A_g.
\]
For continuous one-cocycles $a,b:G\to A$, the obstruction
to lifting $a$ through the actual class-two extension is
represented by
\[
 \omega_a(g,h)=D_g a(h)+\tfrac12a(g)\wedge A_g a(h).
 \tag{E15.2}
\]
Its exact polarization is
\[
 \operatorname{ob}(a+b)-\operatorname{ob}(a)
       -\operatorname{ob}(b)=[a\cup_{\wedge}b]
       \quad\text{in }H^2_{\rm cts}(G,B).
 \tag{E15.3}
\]
Thus the sum of two liftable classes is liftable precisely
on the zero locus of their displayed cup pairing.

For the actual quartic decomposition of E.17,
$A=W\oplus j(V^3)$ at
$\ell=3,5$, where $V=T_\ell E_+$, the mixed component
for $a=b+j_{\mathbf r}h$ has at every good prime
$p\ne2,\ell$ the following actual local representative.
Let $v_p=b(t_p)$ and
$u_p=(F_{E,p}-1)^{-1}h(F_p)$, with the choices of E.14.
Its unique rational tame correction is
\[
 c_p=v_p\wedge j_{\mathbf r}u_p .
 \tag{E15.4}
\]
Writing $t=\min_i v_\ell(r_i)$ and
$\operatorname{ord}(u_p+V)=\ell^{m_p}$, this correction
is integral exactly when
\[
 v_\ell(v_p)+t\ge m_p\qquad(m_p>0).
 \tag{E15.5}
\]
After replacing $u_p$ by its fractional representative,
the sequence of corrections is restricted exactly when
$v_\ell(v_p)+t-m_p\to+\infty$ on the primes with $m_p>0$.
These are the polarization and mixed support maps in
LM7 E.115 and E.117, with the arithmetic linear term
retained. The density of the local inequality's failure
at a fixed finite threshold $d=v_\ell(v_p)+t$ is the
tail $\delta_\ell(m_p\ge d+1)$ of E.14; this statement
concerns the local mixed test with that fixed threshold.
\end{statement}
\begin{proof}
The Hall commutator isomorphism
$\gamma_2N/\gamma_3N\simeq\bigwedge^2A$ gives (E15.1).
Division by two is allowed in $\mathbb Z_\ell$.
An automorphism of this group is determined by its
linear action on $A$ and its remaining central correction.
Preservation of the multiplication shows that the
correction $D_g$ is additive. Applying the action
composition to $(v,0)$ gives
$D_{gh}=B_gD_h+D_gA_h$.
The section lift $g\mapsto(a(g),0,g)$ has multiplication
defect (E15.2); associativity makes it a two-cocycle.
Changing a continuous lift adds a continuous coboundary.
This proves the obstruction assertion for the genuine
action, without setting $D_g$ to zero.

At the cochain level the left side of (E15.3) is
\[
 \tfrac12\{a(g)\wedge A_gb(h)+b(g)\wedge A_ga(h)\}.
\]
Set $q(g)=a(g)\wedge b(g)$ and use
$dq(g,h)=B_gq(h)-q(gh)+q(g)$. The two one-cocycle
identities give
\[
 dq(g,h)=-a(g)\wedge A_gb(h)+b(g)\wedge A_ga(h).
\]
The preceding polarization is therefore
$a\cup_\wedge b+\tfrac12dq$, proving (E15.3).
Its vanishing equivalence follows by lifting the
central extension. In particular this formula
specifies the entire additive defect, including
the possible mixed terms of two general inputs.

For the local mixed formula, consider the central
extension represented by the ordered cup
$\omega(g,h)=b(g)\wedge g e(h)$, where
$e=j_{\mathbf r}h$ is unramified at $p$.
In its multiplication
$(z,g)(w,h)=(z+gw+\omega(g,h),gh)$, choose central-zero
lifts of tame inertia. The relation
$F_pt_pF_p^{-1}=t_p^p$ has central defect
\[
 -p\,v_p\wedge e(F_p).
\]
Changing the inertia lift by $z$ changes the defect by
$(F_{B,p}-p)z$. The toric action is $F_{W,p}=pP_p$,
where $P_p$ is the boundary permutation. The
one-cocycle tame relation gives $P_pv_p=v_p$.
Consequently
\[
 (p^{-1}F_{B,p}-1)
       (v_p\wedge j_{\mathbf r}u)
   =v_p\wedge j_{\mathbf r}(F_{E,p}-1)u.
\]
Every eigenvalue of this mixed module has complex
absolute value $p^{3/2}$, so the displayed operator
is invertible over $\mathbb Q_\ell$. This proves
(E15.4) as its unique rational solution.

The natural map $W\otimes j(V^3)\to W\wedge j(V^3)$
is an integral lattice isomorphism, because these
are complementary direct summands of $A$.
A nonzero pure tensor has valuation equal to the
sum of the two vector valuations: choose primitive
coordinates in each factor and inspect their
product coordinate. A fractional representative of
$u_p+V$ is zero for $m_p=0$ and has valuation $-m_p$
otherwise. This proves (E15.5) and the limit statement.
Replacing a representative by an integral vector
changes $c_p$ by an integral restricted family,
because the inertia coordinates of a continuous
global cocycle tend to zero. E.14 supplies the
claimed density at every fixed threshold.
The passage is \S51.16, folios 733v--736v,
scans 172--178, and the explicit quartic task E.23(iv).
\end{proof}

\begin{statement}{E.16}{Proved}{full pro-three and pro-five sections beyond every nilpotent stage}
\label{st:E.16}
\label{E-regular-prime-family}
Let $\ell\in\{3,5\}$, $r\ge1$, and $n=\ell^r$.
Use the rational unit tangent at
$[0:-1:1]\in X_n$ and the resulting genuine action
on
\[
 N_n=\pi_1(U_{n,\overline{\mathbb Q}})^{(\ell)},\qquad
 V_n=N_n^{\rm ab},\qquad S=\{\ell,\infty\}.
\]
Every class of $H^1(G_{\mathbb Q,S},V_n)$ lifts to a
section of the full geometrically pro-$\ell$
arithmetic extension of $U_n/\mathbb Q$.
There are infinitely many such classes whose lifts
are cuspidal at no boundary point. Each retains
this property in every positive-class intrinsic
nilpotent quotient and every positive-length
solvable quotient, through the compactified
Jacobian quotient specified below.
Here the cuspidal classes in a quotient are the
images of the actual cusp packets under the
displayed arithmetic quotient map.

Thus no intrinsic pro-$3$ nilpotent level for
$n=3^r$, and no intrinsic pro-$5$ nilpotent level
for $n=5^r$, selects only cuspidal sections. This
extends the checked mechanism of LM7 E.34 from
the cubic to these two entire towers. It gives
a specific range of the full section problem;
the passage to all coefficient primes and
mixed finite quotients is retained in E.21.
\end{statement}
\begin{proof}
Put $K=\mathbb Q(\zeta_\ell)$ and let $S_K$ contain
the unique prime above $\ell$. Both fields have
class number one. For $\ell=3$ the Eisenstein
integer norm is Euclidean. For $\ell=5$, the
degree-four field has discriminant $125$, and
Minkowski's ideal-class bound is
\[
 \frac{4!}{4^4}\left(\frac4\pi\right)^2\sqrt{125}<2.
\]
Every ideal class therefore contains an integral
ideal of norm one. The factorization
$\ell=\text{unit}\cdot(1-\zeta_\ell)^{\ell-1}$
gives the single prime in $S_K$.

Write $r_2=(\ell-1)/2$. Unique factorization,
the $S_K$-unit theorem and Kummer theory give
\[
 \dim_{\mathbb F_\ell}
 H^1(G_{K,S_K},\mathbb F_\ell)=r_2+1 .
\]
Indeed the ordinary units contribute rank $r_2-1$,
the roots of unity contribute one dimension, and
the $S_K$-valuation contributes one. The global
Euler formula, in its totally imaginary case, is
\[
 \frac{\#H^0\#H^2}{\#H^1}=\ell^{-r_2}.
\]
Since $\#H^0=\ell$, this proves
$H^2(G_{K,S_K},\mathbb F_\ell)=0$.
The finiteness and Euler inputs are the exact
I.4.10, I.4.15 and I.5.1 statements checked in
E.13~\cite{E_density_MilneADT}.

The maximal pro-$\ell$ quotient $P$ of this Galois
group is free of rank $r_2+1$. To check the
cohomological criterion here, let $J$ be the kernel
of the maximal quotient. A nonzero invariant
character $J\to\mathbb F_\ell$ would produce a
pro-$\ell$ extension of $P$ and contradict its
maximality. Inflation therefore injects
$H^2(P,\mathbb F_\ell)$ into the zero group above.
For a finitely generated pro-$\ell$ group, this
vanishing solves every finite embedding problem
with $\ell$-group kernel: filter an abelian kernel
by trivial simple modules, and then induct on
central characteristic subgroups of a general
finite kernel. A minimal free presentation splits
by compactness of its finite solution sets; its
idempotent splitting is the identity on the
Frattini quotient and hence on the free group.
This proves freeness with the stated rank.

Let $F$ be the free pro-$\ell$ tripod group, with
its standard rational tangent. The actual power
cover identifies
\[
 N_n=\ker(F\longrightarrow(\mathbb Z/n)^2).
 \tag{E16.1}
\]
At the chosen cusp, the equations are
$t=s^n$, $y/z=-(1-s^n)^{1/n}$.
Their leading constants have coherent rational
$\ell$-power roots. Coefficientwise action on
these Puiseux fibre functors therefore gives
the genuine restriction of the tripod action
in (E16.1). Smooth marked specialization away
from $\ell$, with this integral unit tangent,
makes that action unramified there.

Over $K$ the action on $F/\Phi(F)$ is trivial,
since $\mu_\ell\subset K$. The kernel of
$\operatorname{Aut}(F)\to\operatorname{GL}(F/\Phi(F))$
is pro-$\ell$: on the characteristic lower
exponent-$\ell$ quotients its successive kernels
embed in homomorphism groups of finite
$\mathbb F_\ell$-spaces, and the finite basis
theorem gives the initial step. Hence the
action on $F$ and on $N_n$ factors through $P$.
Over $\mathbb Q$ it factors through the actual
relative quotient
\[
 1\longrightarrow P\longrightarrow\Gamma
 \longrightarrow\operatorname{Gal}(K/\mathbb Q)
 \longrightarrow1 .
\]
The quotient has order $\ell-1$. Every finite
$\ell$-primary module through $\Gamma$ has zero
second cohomology over $G_{\mathbb Q,S}$:
over $G_{K,S_K}$ filter it by trivial
$\mathbb F_\ell$ factors and use the preceding
vanishing; then use the exact averaging
operator for the quotient of order $\ell-1$.

Every continuous $V_n$-cocycle factors through
this same $\Gamma$. On $G_{K,S_K}$ its graph
takes values in the pro-$\ell$ group
$V_n\rtimes\rho(G_{K,S_K})$ and therefore kills
$J$. Finite embedding problems now lift the
graph through every finite characteristic
quotient of $N_n$. There are finitely many
solutions at each finite level, since
$\Gamma$ is finitely generated; any finite
set of level constraints is solved at a
common refinement. Compactness gives
compatible solutions and hence the asserted
continuous full lift.

To produce noncuspidal lifts, use the integral
exact sequence from the generalized Jacobian:
\[
 0\longrightarrow W_n\longrightarrow V_n
 \longrightarrow M_n=T_\ell\operatorname{Jac}(X_n)
 \longrightarrow0 .
 \tag{E16.2}
\]
The toric term $W_n$ has the same pro-$\ell$
action factorization. The finite-coefficient
vanishing above and finite $H^1$ groups give
$H^2_{\rm cts}(G_{\mathbb Q,S},W_n)=0$.
Thus (E16.2) surjects on $H^1$.
The group $H^1(G_{\mathbb Q,S},M_n)$ is finitely
generated over $\mathbb Z_\ell$. Its rank is
\[
 g(X_n)=\frac{(n-1)(n-2)}2>0 .
\]
Here $H^2_{\rm cts}$ vanishes by the same argument,
$H^0(M_n)=0$ by Mordell--Weil, and complex
conjugation has $g(X_n)$ eigenvalues of each
sign. The integral Euler formula gives the
displayed rank; possible finite torsion in
$H^1$ does not affect it. Choose any infinite-order
class, lift it through (E16.2), and then apply
the full lifting construction.

For any two Fermat cusps $d,O$, the ratio of
their tangent-line equations has divisor
$n(d-O)$. Each tangent line meets the Fermat
curve only at its tangent cusp, with
multiplicity $n$. Thus every cusp class in
the proper Jacobian is $n$-torsion.
The compactified image of a cuspidal section
is the Kummer class of that cusp; changing its
tangent affects only $W_n$. Our chosen
infinite-order compact class cannot be such
a class. Distinct geometric cusp inertia
classes stay distinct in the free pro-$\ell$
group, so a section normalizing one requires
the underlying cusp to be rational. There
are exactly three rational cusps for odd $n$.
The same compactified map factors through
every quotient named in the theorem, and
retains the obstruction to cuspidality.
Varying the infinite-order class proves
infinitude. We claim no priority for the
general pro-prime phenomenon; compare the
published constructions of Hoshi~\cite{Hoshi2010}.
The manuscript comparison is \S51.16,
folios 733v--736v, scans 172--178.
\end{proof}

\begin{statement}{E.17}{Proved}{the actual quartic coefficient maps and a class-two origin}
\label{st:E.17}
\label{E-quartic-origin}
Put $C:x^4+y^4+z^4=0$, $U=C\setminus\{xyz=0\}$,
$F=\mathbb Q(\zeta_8)$, and $E_+:Y^2Z=X^3+XZ^2$.
For every odd prime $\ell$, the actual arithmetic
extension with geometric kernel
$N/\gamma_3N$, $N=\pi_1(U_{\overline{\mathbb Q}})^{(\ell)}$,
has a continuous section $S_\ell$. Its image in the
rank-two coordinate-unit extension is the zero section.
There is a specified $G_{\mathbb Q}$-equivariant integral
direct sum
\[
 N^{\rm ab}=W\oplus j((T_\ell E_+)^3),
 \qquad \rank W=11,
 \tag{E17.1}
\]
with the forward and inverse maps below. Relative to
$S_\ell$ the formulas of E.15 therefore apply to the
genuine quartic extension. These assertions verify the
precise steps of LM7 E.103 and E.105 used here; no
rank-one assertion about elliptic global cohomology is
needed for this construction or for E.13--E.15.
\end{statement}
\begin{proof}
Let $J=\Jac(C)$ and $G=J_{C,D}$, where $D=C\setminus U$.
Write $p:G\to J$. A tangent-line ratio at two cusps has
divisor four times their difference. The extension
$1\to T_D\to G\to J\to1$ is classified by these cusp
differences in the dual Jacobian. Pullback by $[4]$
therefore splits. Its splitting is unique because
$\Hom(J,T_D)=0$; hence it descends to $\mathbb Q$ and
gives $j_0:J\to G$ with $pj_0=[4]$.

The three morphisms
\[
 \begin{aligned}
 f_0[x:y:z]&=[-x^2y:z^2x:y^3],\\
 f_1[x:y:z]&=[-y^2z:x^2y:z^3],\\
 f_2[x:y:z]&=[-z^2x:y^2z:x^3]
 \end{aligned}
 \tag{E17.2}
\]
have no simultaneous zero in their coordinate triples
on $C$, and direct substitution gives the equation of
$E_+$. For $f_0$, coordinates $u=x/y$, $v=z/y$ give
$X=-u^2$, $Y=uv^2$, so a generic fibre has four points.
The cyclic formulas have the same degree. Their three
different sign characters for the coordinate group
$C_2^2$ give $f_{i*}f_j^*=0$ for $i\ne j$: such a
composite equals its negative, and a homomorphism of
abelian varieties has no nonzero torsion. The diagonal
composites are $[4]$ by push--pull. Consequently
\[
 \Phi=(f_{i*}):J\to E_+^3,\quad
 \Psi=\sum_i f_i^*:E_+^3\to J,\quad
 \Phi\Psi=\Psi\Phi=[4].
\]
The second equality follows also because the image of
$\Psi\Phi-[4]$ is connected and contained in the finite
kernel of the isogeny $\Phi$. The map denoted $j$ in
(E17.1) is $j_{0*}\Psi_*/16$. The actual inverse maps are
\[
 a\longmapsto
 \bigl(a-j_{0*}p_*a/4,\ \Phi_*p_*a\bigr),
 \qquad
 (w,e)\longmapsto w+j_{0*}\Psi_*e/16.
 \tag{E17.3}
\]
All divisions are integral for odd $\ell$.

We verify the origin at the level of actual paths.
Set $\alpha_i=\zeta_8^{1+2i}$ and use the following
points, local parameters, and tangent lines:
\[
 \begin{array}{c|c|c|c}
 d&\text{point}&s_d&L_d\\\hline
 A_i&[0:\alpha_i:1]&x/z&y-\alpha_i z\\
 B_i&[\alpha_i:0:1]&y/z&x-\alpha_i z\\
 C_i&[\alpha_i:1:0]&z/y&x-\alpha_i y
 \end{array}
 \quad O=B_0 .
 \tag{E17.4}
\]
All tangent vectors have parameter value one. The
nonzero varying coordinate is
$\alpha_i(1+s_d^4)^{1/4}$; thus the leading coefficient
of $L_d$ is $\alpha_i/4$ in order four. Frame
$z_d=\mathcal O(d-O)\in G(F)$ by $s_d^{-1}$ at $d$,
$s_O$ at $O$, and its canonical section elsewhere.
For $f_d=L_d/L_O$, the torus coordinates of $4z_d$ are
\[
 t_{q,d}=
 \frac{\operatorname{lead}_q f_d}{
       \operatorname{lead}_O f_d}
       =\zeta_8^{K_{qd}}(\sqrt2)^{N_{qd}}.
 \tag{E17.5}
\]
This is a directly evaluable formula: the denominator
is $4L_d(O)/\alpha_0$; the numerator is
$\alpha_d/(4L_O(d))$ for $q=d$ and $L_d(q)/L_O(q)$
otherwise. The independently executed program
\texttt{quartic\_frames.py} evaluates all $121$ entries
in $\mathbb Q[z]/(z^4+1)$ using rational coefficients,
and records both $N$ and the phases $K$. With the order
$A_0,\ldots,A_3,B_1,B_2,B_3,C_0,\ldots,C_3$,
the three columns used below have average
\[
 \frac1{32}\sum_{d=B_1,B_2,B_3}N_{\bullet d}
 =v=\left(-\tfrac12,-\tfrac12,-\tfrac12,-\tfrac12,
 -\tfrac58,-\tfrac34,-\tfrac58,
 -\tfrac12,-\tfrac12,-\tfrac12,-\tfrac12\right).
 \tag{E17.6}
\]
The entry formulas prove the computation without a
numerical field approximation.

Put $H=G_F$ and choose coherent odd roots of $2$.
Choose roots of $\sqrt2$ whose squares are these roots;
then $\kappa_{\sqrt2}=\kappa_2/2$. All coherent odd
roots of $\mu_8$ lie in $F$, so their $H$-cocycles
vanish. For an odd power $m$ of $\ell$, choose an
$m$-division $v_m$ of $4z_d=t_d$. Integers $a,b$ with
$4a+mb=1$ give $u_m=a v_m+bz_d$, satisfying
$mu_m=z_d$, $4u_m=v_m$. These two equations determine
$u_m$ uniquely and compatibly, because $[4]$ is an
automorphism on $G[m]$. Its proper image is the fixed
four-torsion point $b[d-O]$. Thus the actual abelian
path cocycle, in (E17.3), is
\[
  \kappa(z_d)=v_d\kappa_2,
       \qquad v_d=N_{\bullet d}/8\in W .
 \tag{E17.7}
\]
To compare these framed divisors with tangential
paths, apply each boundary-unit character to the
framed Abel--Jacobi map. Regularization at the two
endpoints multiplies its value by
$(-1)^{m_d+m_O}$ for a divisor $\sum m_q q$.
This is the character of the torus tuple having
entry $+1$ at $d$ and $-1$ at every other non-$O$
cusp. Its coherent odd divisions are $H$-fixed.
Thus it contributes zero cocycle, proving (E17.7)
for the stated tangential paths as well.

Lift these abelian paths to full paths $p_d$ from
the tangent at $d$ to the tangent at $O$, composing
paths from right to left; the
map of path torsors is onto, since the quotient
$N\to N^{\rm ab}$ is onto. Put
$d_d=p_dt_dp_d^{-1}$ for positive cusp inertia.
The genuine conjugator $u_d(g)=p_dg(p_d)^{-1}$ has
abelian coordinate $v_d\kappa_2(g)$ and satisfies
$u_dg(d_d)u_d^{-1}=d_d^{\chi(g)}$.
The eleven peripheral classes together with a basis
of the compact summand form a basis of $N^{\rm ab}$;
their lifts form a free pro-$\ell$ basis by the
Frattini criterion. In the class-two coordinates of
E.15, put $L(e_d)=e_d\wedge v_d$ and $a=\kappa_2$.
The resulting actual $H$-action on the subgroup
$W\oplus\bigwedge^2W$ is
\[
 g(x,z)=(\chi_gx,\chi_g^2z+\chi_ga(g)L(x)).
 \tag{E17.8}
\]
Indeed the central term of
$u_d^{-1}d_d^{\chi_g}u_d$ is exactly the last term.
No action on a compact free generator is used here.

We give the averaging map used twice in the descent.
For an extension of $G_{\mathbb Q}$ by an abelian
$\mathbb Z_\ell$-module $M$ with a section $s_H$
over $H$, fix four left cosets $g_iH$ and lifts $e_i$.
For a lift $e$ of $g$ put $j=g\cdot i$,
$h_i=g_j^{-1}gg_i$, and
\[
 b_i(e)=e e_i s_H(h_i)^{-1}e_j^{-1}\in M,
 \quad B(e)=\sum_i b_i(e),
 \quad s_0(g)=(-B(e)/4)e .
 \tag{E17.9}
\]
Multiplying the displayed words gives
$B(ef)=B(e)+gB(f)$ and $B(m)=4m$.
Hence $s_0$ is a homomorphism independent of $e$.
It is continuous by using the continuous lifts
$e_i s_H(g_i^{-1}g)$ on the four clopen cosets.

First apply (E17.9) to the abelianized quartic
extension using the paths from the four $B_i$.
Its restriction to $H$ differs from the tangent
section at $O$ by $a(g)v$, with the positive sign
because $b_i(s_H(h))=-u_i(h)$ in abelian coordinates.
Both coordinate-unit rows annihilate (E17.6).
The leading values at $O$ are $\zeta_8$ and $1$,
whose odd Kummer cocycles on $H$ vanish. The toric
cocycle is therefore zero on $H$, and is zero
globally: a cocycle zero on this normal subgroup
has values in $\mathbb Z_\ell(1)^{H}=0$.
This constructs the asserted toric-zero abelian
origin $s_\ell$.

Over $H$, its literal class-two lift relative to
the original tangent has cocycle
\[
 c_v(g)=\bigl(a(g)v,\tfrac12a(g)^2L(v)\bigr).
 \tag{E17.10}
\]
Substitute (E17.8) and
$a(gh)=a(g)+\chi_ga(h)$ in the group law E15.1:
the extra term is
$\chi_ga(g)a(h)L(v)$, and the wedge of $v$ with
itself vanishes. This proves the cocycle identity.
Pull back the actual class-two extension along the
fixed $s_\ell$. The kernel is the abelian module
$\bigwedge^2N^{\rm ab}$, and (E17.10) gives its
$H$-section. Applying (E17.9) to this pullback
constructs $S_\ell$ over $G_{\mathbb Q}$.
This last averaging is in an abelian-kernel
extension, so its homomorphism proof applies
literally. The anchors are \S51.16, folios
733v--736v, scans 172--178, and task LM7-E.23(iv).
\end{proof}

\subsection{The cubic class-three map and its remaining arithmetic input}

\begin{statement}{E.18}{Proved}{the exact class-three criterion and every nonempty fibre}
\label{st:E.18}
\label{E-class-three-criterion}
Let $N=\pi_1(U_{3,\overline{\mathbb Q}})^{(3)}$, with
the genuine action defined by the rational unit tangent
at $O=[1:-1:0]$. Put $N_i=N/\gamma_{i+1}N$,
$V=N^{\rm ab}$, $A=\gamma_3N/\gamma_4N=\operatorname{Lie}_3V$,
and $K=\mathbb Q(\zeta_3)$. The rank of $A$ is $330$.
Fix a continuous ordered-Hall set section $h:N_2\to N_3$
carrying one to one. For an actual continuous cocycle
$a:G_{\mathbb Q}\to N_2$ and every prime
$\mathfrak q\nmid3$ of $K$, choose a prolongation,
positive tame generator $t_{\mathfrak q}$ and arithmetic
Frobenius $F_{\mathfrak q}$. Write $q=N\mathfrak q$ and
\[
 \begin{aligned}
 n_{\mathfrak q}&=h(a(t_{\mathfrak q})),&
 x_{\mathfrak q}&=h(a(F_{\mathfrak q})),\\
 \epsilon_{\mathfrak q}
 &=x_{\mathfrak q}\rho_3(F_{\mathfrak q})(n_{\mathfrak q})
       x_{\mathfrak q}^{-1}n_{\mathfrak q}^{-q}\in A,&
 z_{\mathfrak q}
 &=-(\Phi_{\mathfrak q}-q)^{-1}\epsilon_{\mathfrak q}
       \in A\otimes\mathbb Q_3 .
 \end{aligned}
 \tag{E18.1}
\]
Here $\rho_3$ is the actual action on $N_3$ and
$\Phi_{\mathfrak q}$ its induced action on $A$.
The fibre over $[a]$ of
\[
 H^1_{\rm cts}(G_{\mathbb Q},N_3)
       \longrightarrow H^1_{\rm cts}(G_{\mathbb Q},N_2)
 \tag{E18.2}
\]
is nonempty precisely when every $z_{\mathfrak q}$ is
integral and $z_{\mathfrak q}\to0$ in the uniform
$3$-adic congruence topology. Each nonempty fibre is
a torsor under the explicitly evaluated group
\[
 \mathbb Z_3^{166}\oplus(\mathbb Z/3)^7
   \oplus(\mathbb Z/9)^{56}\oplus(\mathbb Z/27)^{25}
   \oplus(\mathbb Z/81)^3 .
 \tag{E18.3}
\]
For every class unramified outside $3$, this fibre
is nonempty. At the $K$-prime $(2)$ the integral
local obstruction module is
\[
 \coker(\Phi_{(2)}-4)=
 (\mathbb Z/3)^{70}\oplus(\mathbb Z/9)^{24}
 \oplus(\mathbb Z/27)^{76}\oplus(\mathbb Z/81)^{28}.
 \tag{E18.4}
\]
The test (E18.1) retains the higher action and the
actual Frobenius cocycle values as data. Its proof
and the two computed coefficient groups verify LM7
E.41 and E.44 in the degree used here. E.21 records
the unevaluated arithmetic data for unrestricted
inputs rather than counting this test as their
complete evaluation.
\end{statement}
\begin{proof}
The ordered commutator basis of a free pro-three
group gives $\gamma_3N/\gamma_4N=\operatorname{Lie}_3V$.
In the convention $[[x_i,x_j],x_k]$ with $i>j$ and
$k\ge j$, it has $(10^3-10)/3=330$ basis elements.
Ordered products with $3$-adic exponents supply the
continuous Hall set section; no BCH division by
three is required.

The set lift $\widehat a=h\circ a$ has central defect
\[
 \omega(g,h)=\widehat a(g)g(\widehat a(h))
                    \widehat a(gh)^{-1}\in A.
 \tag{E18.5}
\]
Associativity proves the two-cocycle identity.
A correcting continuous cochain $b$ gives a lift
exactly when $db=-\omega$. At a prime away from
three the genuine action is unramified. Wild
inertia has trivial image in the pro-three
geometric group. The lower inertia cocycle is a
homomorphism, and we lift it by
$t^u\mapsto n^u$ for $u\in\mathbb Z_3$.
The tame relation gives precisely
$\epsilon_{\mathfrak q}$ in (E18.1).
Changing $n$ to $zn$ changes this error by
$(\Phi_{\mathfrak q}-q)z$; changing $x$ centrally
has no effect. The rational weights of $A$ have
Frobenius absolute values $q^{j/2}$ for
$3\le j\le6$. None equals $q$, proving existence
and uniqueness of the inverse used in (E18.1).

Here is the continuous global comparison, including
its image rather than its closure. Define
\[
 \mathcal R(A)=\{(v_{\mathfrak q})\in
          \prod_{\mathfrak q\nmid3}A:
                   v_{\mathfrak q}\longrightarrow0\},
 \qquad d(v)_{\mathfrak q}=(q^{-1}\Phi_{\mathfrak q}-1)v_{\mathfrak q}.
\]
We have the actual isomorphism and representative
\[
 H^2_{\rm cts}(G_K,A)=\mathcal R(A)/d\mathcal R(A),
 \qquad [\omega]\longmapsto
       [(q^{-1}\epsilon_{\mathfrak q})_{\mathfrak q}].
 \tag{E18.6}
\]
To prove it, let $S=\operatorname{Spec}\mathcal O_K[1/3]$
and $A_m=A/3^mA$. The proof of E.16 gives
$H^2(G_{K,\{3\}},A_m)=0$; the action has pro-three
image and a filtration with trivial simple factors.
The number-field cohomological-dimension theorem,
the finiteness theorem, and the scheme comparison
give $H^3(S,A_m)=0$ and finite $H^1(S,A_m)$.
We use exactly Milne I.4.10(c), I.4.15 and II.2.9,
with three inverted and no real place
\cite{E_density_MilneADT}.

For a finite set $J$ of primes away from three,
localization and purity give the triangle
\[
 R\Gamma(S,A_m)\longrightarrow
 R\Gamma(S-J,A_m)\longrightarrow
 \bigoplus_{\mathfrak q\in J}
 R\Gamma(k_{\mathfrak q},A_m(-1))[-1]
 \longrightarrow R\Gamma(S,A_m)[1].
 \tag{E18.7}
\]
The local purity term is also obtained directly:
over a strict henselization wild inertia has zero
three-primary cohomology and tame inertia is
$\mathbb Z_3(1)$. Take the direct limit over finite
$J$, then the derived inverse limit in $m$.
The residue complexes are the compatible two-term
complexes $[A_m\xrightarrow{q^{-1}\Phi-1}A_m]$.
Both direct-sum term systems have surjective
transition maps. Their inverse limit is the literal
complex $[\mathcal R(A)\xrightarrow d\mathcal R(A)]$.
Finite $H^1(S,A_m)$ has vanishing first derived
limit, and the two higher unramified terms vanish.
Jannsen's Proposition 1.6 and Theorem 2.2 identify
this calculation with continuous-cochain cohomology
\cite{Jannsen1988}. This proves the first equality
in (E18.6).

The support map measures failure of the chosen
inertia splitting to be Frobenius invariant.
Evaluating its transport on $t^{q^{-1}}$ gives
$q^{-1}\epsilon_{\mathfrak q}$, which proves the
second formula and its Tate twist. Its family is
restricted: at every fixed finite characteristic
quotient the cocycle graph has finite image and
is ramified at finitely many primes. Thus all its
inertia values approach one uniformly outside
finite sets. Continuity of $h$ and of powers
$n^u$ on the compact spaces in question gives
the same property for the local null-homotopies
used in (E18.7). Hence their errors are zero
modulo $3^mA$ outside a finite set for every $m$.

Equation (E18.6), with the rational invertibility
already proved, makes its vanishing equivalent
to the integral restricted corrections of the
statement. Restriction from $G_{\mathbb Q}$ to
$G_K$ is injective on this $H^2$: restriction
followed by corestriction is multiplication by
two, an automorphism. Thus the same criterion
gives a continuous global lift. Changing Hall
lifts changes the support sequence by an actual
element of $d\mathcal R(A)$. Changing tame
generators rescales by units, changing Frobenius
by an inertia power preserves its error, and
changing prolongations transports by integral
Galois automorphisms followed by such a boundary.
The criterion therefore depends only on $[a]$.

Differences of two lifts are $A$-valued cocycles.
At a good prime their inertia values satisfy
$(\Phi-q)v=0$, so they vanish. Every such
cocycle is therefore unramified outside three.
There is also no lower geometric stabilizer
to divide out: the lowest nonzero graded
coordinate of a fixed element of the
section-twisted $N_2$ would be fixed by the
arithmetic action. At any good Frobenius all
its eigenvalues have absolute value greater
than one; inner twisting acts trivially on
this lowest coordinate. The contradiction
proves triviality of the stabilizer. The fibre
is exactly a torsor under $H^1(G_{\mathbb Q,\{3,\infty\}},A)$.
E.16's embedding-problem argument lifts every
unramified lower section through this central
step as well.

We supply the integral coefficient evaluation.
The eight nonorigin cusps are
$A_i=[0:-\zeta^i:1]$, $B_i=[-\zeta^i:0:1]$,
and $C_1,C_2=[1:-\zeta^i:0]$.
In the order $A_0,A_1,A_2,B_0,B_1,B_2,C_1,C_2$,
their elliptic three-torsion coordinates are
\[
 (u_d,v_d)=(1,0),(2,2),(0,1),(2,0),
                 (1,1),(0,2),(2,1),(1,2).
 \tag{E18.8}
\]
Tangent-line divisors make all nine cusps three-torsion;
their collinear triples give this addition table.
On $Y^2+Y=X^3-7$ the basis points are
$(3,-5),(3\zeta,-5)$, whose tangent-line ratio gives
the pairing $-9(\zeta-1)/(9(\zeta^2-1))=\zeta$.
For the framed generalized-Jacobian extension,
translation on the trivialization of the pullback
line bundle is the inverse of this pairing.
Thus, writing $U_d=(v_d,-u_d)$, the integral
Tate map of the unique $j_0$ with $pj_0=[3]$ is
$j_{0*}=\binom{3I_2}{U}$. Equivariance forces
the following actual degree-one matrices:
\[
 \begin{gathered}
 R=\begin{pmatrix}0&-1\\1&-1\end{pmatrix},\quad
 C_E=\begin{pmatrix}1&-1\\0&-1\end{pmatrix},\\
 A(\psi,q)=
 \begin{pmatrix}A_E&0\\U(A_E-qI)/3&qI_8\end{pmatrix},
       \quad A_E=aI+bR,\quad\psi=a+b\zeta,\\
 C_V=\begin{pmatrix}C_E&0\\(UC_E-C_TU)/3&C_T\end{pmatrix}.
 \end{gathered}
 \tag{E18.9}
\]
Here $C_T$ is minus the permutation interchanging
subscripts one and two. All displayed divided
blocks are integral. The elliptic Frobenius
scalars at $(2)$ and a specified prime over seven
are $-2$ and $1+3\zeta$: their norm, point count,
and primary congruence determine them. In
$1+3\mathbb Z_3[\zeta]$, the logarithms divided
by three of $1+3\zeta$ and $1+3\zeta^2$ reduce
to a basis modulo three. They therefore give
the full rank-two CM image. Since the unramified
pro-three source is free of rank two, a Frobenius
lift $\tau$ at the chosen prime and
$\tau'=c\tau c$ form a free basis by the
Frattini criterion.

An $A$-cocycle normalized to vanish at $c$ is
freely determined by its value at $\tau$.
Its value at $\tau'$ is $J$ times that value,
where $J=\operatorname{Lie}_3(C_V)$.
The coboundaries preserving this normalization
are $(T-1)A^c$, where
$T=\operatorname{Lie}_3(A(1+3\zeta,7))$.
Since two is a unit, $A^c=(1+J)A$. Hence the
actual fibre group is
\[
 \coker\bigl((T-1)(1+J):A\to A\bigr).
 \tag{E18.10}
\]
The program \texttt{class\_three\_fibre.py}
constructs both integral matrices in the Hall
basis. It verifies $J^2=1$, $\tr J=-2$, and
the Smith valuations
$0^{73},1^7,2^{56},3^{25},4^3$.
These are all nonzero elementary divisors:
$1+J$ has exact rational rank $164$, and the
calculation finds $164$ pivots at precision
$3^{10}$. Its remaining free rank is $166$.
This proves (E18.3), including its torsion.

For (E18.4), the independent program
\texttt{class\_three\_matrix.py} applies the
same Hall functor to $A(-2,4)$. It finds all
$330$ pivots of $\Phi-4$ at precision $3^9$,
with valuations
$0^{132},1^{70},2^{24},3^{76},4^{28}$.
Each pivot is strictly below the working
precision, so integral elementary operations
lift with the same valuations. As a separate
check, the rational eigenvalues have
multiplicities
$64^{168},(-32)^{128},16^{32},(-8)^2$;
their determinant difference from four has
valuation $168+256+32+2=458$, equal to the
sum of the computed valuations. This proves
the full cokernel, of order $3^{458}$.
All source maps, actions, and choices have
therefore been retained. The manuscript
anchors are \S51.16, folios 733v--736v,
scans 172--178, and LM7-E.23(i).
\end{proof}

\begin{statement}{E.19}{Proved}{a linear reduction of the entire zero-abelian slice}
\label{st:E.19}
\label{E-zero-slice}
Retain E.18 and put $M=\bigwedge^2V$ and
$L=\gamma_2N/\gamma_4N$. The fibre of
$H^1(G_{\mathbb Q},N_2)\to H^1(G_{\mathbb Q},V)$
over zero is exactly $H^1(G_{\mathbb Q},M)$.
Its class-three obstruction is the linear connecting
homomorphism for the actual $G_{\mathbb Q}$-module
sequence
\[
 0\longrightarrow A\longrightarrow L
       \longrightarrow M\longrightarrow0 .
 \tag{E19.1}
\]
This whole obstruction is determined by the action
on $N_2$. An action on $N_3$ need not be supplied
separately for this slice.

More precisely let
$w=(3f_1+U_1)\wedge(3f_2+U_2)\in M$,
in the integral coordinates (E18.9).
It is primitive and spans a $G_{\mathbb Q}$-submodule
$\mathbb Z_3(1)$. Every $[z]\in H^1(G_{\mathbb Q},M)$
has a unique decomposition
\[
 [z]=[z_0]+\beta w,\qquad
 [z_0]\in H^1(G_{\mathbb Q,\{3,\infty\}},M),
 \qquad \beta=\sum_{p\ne3}a_p\kappa_3(p),
 \quad a_p\in\mathbb Z_3,\quad a_p\to0.
 \tag{E19.2}
\]
Choose a lift $\omega\in L$ of $w$ with
$c\omega=-\omega$ for complex conjugation $c$.
The actual continuous cocycle
\[
 \eta(g)=\chi(g)^{-1}\bigl(g\omega-\chi(g)\omega\bigr)
       \in A(-1)
 \tag{E19.3}
\]
is unramified outside three. The obstruction of
(E19.2) is exactly
\[
                  [\eta\cup\beta]\in
                        H^2_{\rm cts}(G_{\mathbb Q},A).
 \tag{E19.4}
\]
Over $K$, its unique rational tame corrections are
\[
 -a_p\,(q^{-1}\Phi_{\mathfrak q}-1)^{-1}
                          \eta(F_{\mathfrak q}),
       \qquad \mathfrak q\mid p,\quad p\ne3.
 \tag{E19.5}
\]
Thus the exact image in this slice is the kernel
of the linear map (E19.4), equivalently the integral
restricted locus of (E19.5); its nonempty fibres
are (E18.3).

The remaining action class (E19.3) is a single
element of the explicitly computed group
\[
 H^1(G_{\mathbb Q,\{3,\infty\}},A(-1))
 =\mathbb Z_3^{164}\oplus(\mathbb Z/3)^{59}
       \oplus(\mathbb Z/9)^{27}\oplus(\mathbb Z/27)^7.
 \tag{E19.6}
\]
Its normalized cocycle is determined by the one
vector $\eta(\tau)\in A(-1)$, with $\tau$ as in
E.18. We have evaluated the coefficient group
and the exact reduction maps, not this vector.
The unresolved evaluation is retained in E.21.
\end{statement}
\begin{proof}
An abelianly trivial class in $H^1(N_2)$ can be
conjugated by a lift of its abelian coboundary to
a cocycle in $M$. Its further ambiguity from
$V^{G_{\mathbb Q}}$ vanishes by good Frobenius
weights, exactly as in E.18. This proves the
first identification. The group $L$ is abelian
because $[\gamma_2N,\gamma_2N]\subset\gamma_4N$.
Its two graded pieces are free $\mathbb Z_3$-modules,
and ordered Hall products make (E19.1) a sequence
of finite free modules with continuous action.
A linear Hall section $\iota:M\to L$ gives
\[
 g\iota(m)=\iota(B_gm)+J_gm,
       \qquad B_g=\bigwedge^2 A_g,
       \quad J_g:M\to A.
\]
For an $M$-cocycle $z$, its lift has defect
$J_gz(h)$, the connecting cocycle of (E19.1).
This proves linearity. The action on every
commutator $[x_i,x_j]$ modulo $\gamma_4N$
depends only on the images of $x_i,x_j$ modulo
$\gamma_3N$: changing either by a degree-three
element changes this commutator in degree four.
It follows on the Hall generators that the
entire action on $L$ is determined by the
given action on $N_2$.

The coefficient of $w$ on
$e_{A_0}\wedge e_{A_2}$ is one, so $w$ is primitive.
The map $j_{0*}$ in E.18 is equivariant, and its
exterior square has determinant character $\chi$.
At each good prime the only rational eigenline
of $M$ with Frobenius eigenvalue $p$ is this
compact determinant line: all other eigenvalues
have absolute values $p^{3/2}$ or $p^2$.
Consequently the integral inertia value of an
$M$-cocycle is $a_pw$. These coefficients tend
to zero by continuity. The uniformly convergent
sum $\sum_{p\ne3}a_p\kappa_3(p)w$ is an actual
continuous cocycle with precisely these residues.
Its subtraction kills all inertia away from
three, giving (E19.2). The residues and the
remaining unramified class also prove uniqueness.

Average an arbitrary lift of $w$ as
$\omega=(\omega_0-c\omega_0)/2$.
This still lifts $w$ and satisfies the desired
sign. Formula (E19.3) then has $\eta(c)=0$.
Applying $gh$ to $\omega$ shows
$\eta(gh)=\eta(g)+\chi(g)^{-1}g\eta(h)$,
which is exactly the cocycle law in $A(-1)$.
The action on $L$ is unramified away from
three, proving the same for $\eta$.
Changing $\omega$ adds a coboundary, so its
class is the pullback of (E19.1) along
$\mathbb Z_3(1)w\hookrightarrow M$.

The obstruction of $z_0$ vanishes by
$H^2(G_{\mathbb Q,\{3,\infty\}},A)=0$,
proved in E.16. Lift the other term by
$g\mapsto\beta(g)\omega$. Its defect is
$\chi(g)\beta(h)\eta(g)$, which is the
ordered cup $\eta\cup\beta$ under
$A(-1)\otimes\mathbb Z_3(1)\to A$.
This proves (E19.4) with its sign. Locally,
both the inertia and Frobenius lift lie in
the abelian group $L$. Their conjugation
term in (E18.1) therefore disappears, and
the error is
$a_p(g\omega-q\omega)=qa_p\eta(F)$.
Equation (E18.1) gives (E19.5).

Finally the same free source basis as in
E.18 evaluates the group containing $\eta$.
The matrices on $A(-1)$ are $T/7$ and
$-J$. The normalized cocycle is determined
freely by its value at $\tau$, and its
coboundaries form
$(T-7)(1-J)A$, since seven is a unit.
The independent computation
\texttt{class\_three\_zero\_slice.py}
constructs this matrix from E18.9. It finds
$166$ nonzero pivots with Smith valuations
$0^{73},1^{59},2^{27},3^7$ at precision
$3^{10}$. These are all its elementary
divisors because $1-J$ has exact rank
$(330+2)/2=166$. The cokernel has the
$164$ free coordinates in (E19.6).
This calculation specifies the finite
arithmetic datum still to evaluate; it
does not infer its value from the matrix
of the graded action. The manuscript
anchors are \S51.16, folios 733v--736v,
scans 172--178, task LM7-E.23(i).
\end{proof}

\begin{statement}{E.20}{Proved}{a fully checked class-two lift obstructed at class three}
\label{st:E.20}
\label{E-finite-next-obstruction}
Let $P=[0:-1:1]$, $Q=[-1:0:1]$ on the plus Fermat
cubic, and let $\alpha_P\in H^1(G_{\mathbb Q},V)$ be
the class of the rational tangent at $P$ with value
two in the parameter $-x/z$. For positive cusp
inertia $e_Q\in V$, set
\[
           \alpha=\alpha_P+27\kappa_3(19)e_Q.
 \tag{E20.1}
\]
This class lifts to the full intrinsic class-two
quotient of $N$. Every such lift fails to extend
to $N/\gamma_4N$, $N/N''$, or $N/\gamma_3I_Q$,
where $I_Q$ is the closed normal subgroup generated
by the marked cusp $Q$. The failure is detected
at $19$ by an actual arithmetic quotient of $N$
of order $3^{30}$ and nilpotency class three.
This independently verifies the exact example
of LM7 E.38 and its transfer maps.
\end{statement}
\begin{proof}
We first check existence at full intrinsic class
two. Use $K,V,M,w$ and the integral lattice of
E.18--E.19, and put $M_0=\ker(M\to\bigwedge^2T_3E)$.
Since $w$ has coefficient one at
$e_{A_0}\wedge e_{A_2}$, deleting that column gives
a basis of $M/\mathbb Z_3w$. Deleting the row
$f_1\wedge f_2$ gives a basis of $M_0$.
At a prime of norm $q$ the map
\[
 d_{\mathfrak q}:M/\mathbb Z_3w\longrightarrow M_0,
       \qquad [m]\longmapsto(q^{-1}\bigwedge^2A_{\mathfrak q}-1)m
 \tag{E20.2}
\]
is rationally invertible, by the weight computation
in E.19. The same localization proof as E.18,
now in degree two and retaining its one-dimensional
rational kernel, identifies the obstruction with
the restricted support cokernel of these maps.

The point-tangent class $\alpha_P$ lifts to the
full group. Its framed cusp-path values have
support only at three, and its tangent value
adds the prime two. This follows directly by
evaluating the ratios of the cubic tangent lines;
their nonzero values are products of roots of
unity and $1-\zeta_3$. At the prime of $K$ over
two, $19\in1+2\mathcal O_{K,2}$ has coherent
third-power roots because cubing is an automorphism
of this pro-two unit group. Thus the added
Kummer class is zero locally, and the obstruction
there is unchanged from $\alpha_P$.

At a prime over $19$, normalize
$\kappa_3(19)(t)=1$. Let $a_P=\alpha_P(F)$, and
let $k_Q$ be the abelian path conjugator from the
unit tangent at $Q$ to the fixed origin. Expanding
the class-two action on its inertia loop and the
tame relation gives the support numerator
\[
 D=(a_P+27\kappa_3(19)(F)e_Q)\wedge27e_Q
                   -27k_Q\wedge e_Q
       =27(a_P-k_Q)\wedge e_Q\in27M_0.
 \tag{E20.3}
\]
The two elliptic Frobenius scalars over $19$ are
$-2+3\zeta$ and $-5-3\zeta$. Substitution in
E18.9 and exterior squaring construct the two
actual $44$-by-$44$ matrices (E20.2).
The independently executed program
\texttt{class\_three\_counterexample.py} finds
Smith valuations $0^8,1^6,2^{18},3^{12}$ for each.
All $44$ pivots are found below its precision,
and their valuation sum is $78$. In particular
$27M_0\subset d_{\mathfrak q}(M/\mathbb Z_3w)$,
so (E20.3) has an integral correction at both
primes. Every other good-prime support term is
zero. These corrections have finite support,
and the continuous localization quotient
therefore gives a global class-two lift.

For the higher obstruction use the unit tangent
at $P$. The power-cover coordinate $t=s^3$,
$s=-x/z$, and the branch
$y/z=-(1-s^3)^{1/3}$ with rational leading
coefficient give a genuine equivariant inclusion
\[
 N=\ker(F\to(\mathbb Z/3)^2)\hookrightarrow
       F=\langle x,y\rangle^{\rm pro\text{-}3}.
 \tag{E20.4}
\]
A path to the previous tangent changes the
arithmetic origin by its actual tangential
cocycle and preserves the defined abelian
class (E20.1). Put $I_y=\ker(F\to\mathbb Z_3)$,
where $x\mapsto1$, $y\mapsto0$. With $T=1+U$,
the exact quotient map is
\[
 \begin{gathered}
 F/[I_y,I_y]=R=\mathbb Z_3\ltimes\mathbb Z_3[[U]],\\
 (a,v)(b,w)=(a+b,v+T^a w),\qquad
 x\mapsto(1,0),\quad y\mapsto(0,1),\\
 N\twoheadrightarrow H=3\mathbb Z_3\ltimes J,
            \qquad J=(3,U)\mathbb Z_3[[U]].
 \end{gathered}
 \tag{E20.5}
\]
To verify the quotient, the conjugates of $y$
generate $I_y^{\rm ab}$ over the completed
group ring $\mathbb Z_3[[\mathbb Z_3]]$.
The universal free-group map to the displayed
semidirect product is inverse on this module
generator. It proves equality of the modules,
the group law, and the topology. On
$F^{\rm ab}$ the map is the identity, with
second coordinate the augmentation $v(1)$;
this proves the stated image of $N$.

The normalized tripod action is
$x\mapsto x^\chi$, $y\mapsto f_g^{-1}y^\chi f_g$
with $f_g\in F'$. In $R$, this conjugator lies
in the abelian module and commutes with $y$.
Consequently the actual descended action is
\[
           g(a,v(T))=(\chi(g)a,\chi(g)v(T^{\chi(g)})).
 \tag{E20.6}
\]
The group $R$ is metabelian. Every loop above
the base puncture one maps into its abelian
normal module. Hence the same map from $N$
factors through $N/N''$ and $N/\gamma_3I_Q$,
with all arithmetic actions retained.

At $19$, choose $\zeta_9$ reducing to
$4\in\mathbb F_{19}$; then
$2^{(19-1)/9}=4$ gives
$\kappa_3(2)(F)\equiv1\pmod9$.
The first coordinate of any section above
(E20.1) is $3\kappa_3(2)$, and its tame inertia
has augmentation $3\cdot27=81$.
Conjugating by $x^{3c}$ normalizes its Frobenius
first coordinate to three: the change is
$3(1-19)c$, and
$(\kappa_3(2)(F)-1)/18$ is integral.
The inertia value is therefore $(0,v(T))$
with $v(1)=81$, and the tame relation becomes
\[
                         T^3v(T^{19})=v(T).
 \tag{E20.7}
\]
The Frobenius module coordinate cancels, so this
equation covers every possible choice of lift.
Over $\mathbb Q_3[[U]]$, its unique solution is
$v(T)=81T^{-1/6}$. Indeed putting $T=e^z$ and
multiplying by $e^{z/6}$ makes every positive
coefficient invariant under $z\mapsto19z$,
so it vanishes. Its fourth coefficient is
$1729/384$, of three-adic valuation $-1$.

For a finite certificate, reduce the module
modulo $(729,U^5)$ and the first coordinate
modulo $2187$. The coefficients of
$(1+U)^{2187}-1$ through degree four are
divisible by $729$, so (E20.5)--(E20.6)
descend to a finite arithmetic quotient.
Its subgroup $H_6$ has order $3^{35}$.
Writing $v=81+\sum_{i=1}^4v_iU^i$ gives
$Kv=b\pmod{729}$, where
\[
 K=\begin{pmatrix}
 18&0&0&0\\228&360&0&0\\
 1539&7581&6858&0\\7315&86640&205770&130320
 \end{pmatrix},\qquad
 b=\begin{pmatrix}-243\\-243\\-81\\0\end{pmatrix}.
 \tag{E20.8}
\]
The row $r=(1,660,558,162)$ satisfies
$rK=0$ and $rb=486$ modulo $729$.
These identities are reproduced by the program
directly from the polynomial substitution.

Let $\delta=T^3-1$. Commutators in $H_6$ give
$\gamma_{c+1}H_6=\delta^cJ$. The subgroup
$\gamma_4H_6$ is generated by
\[
       3\delta^3=81U^3+243U^4,
               \qquad U\delta^3=27U^4.
 \tag{E20.9}
\]
It has order $3^5$: the first generator has
order nine, the second order twenty-seven,
and their cyclic groups meet trivially by
their $U^3$ coordinates. The program also
enumerates all $243$ elements. The row $r$
annihilates both columns in (E20.9), so its
nonzero pairing with $b$ persists in the
quotient $H_6/\gamma_4H_6$ of order $3^{30}$.
This quotient has class exactly three,
because $\delta^2U=9U^3+18U^4$ is not in
(E20.9). The arithmetic surjection
\[
             N\twoheadrightarrow H
                 \twoheadrightarrow H_6/\gamma_4H_6
 \tag{E20.10}
\]
factors through $N/\gamma_4N$ by functoriality
of the lower central series and through the
other two specified quotients by (E20.5).
Every proposed section would yield the
excluded local solution of (E20.8).
This proves all three failures. The anchors
are \S51.16, folios 733v--736v, scans
172--178, and LM7-E.23(i)--(ii).
\end{proof}

\begin{statement}{E.21}{Open}{the unrestricted intrinsic tower image}
\label{st:E.21}
\label{E-remaining-tower}
The following requested classifications remain
unfinished in this return:
\begin{enumerate}[label=(\roman*)]
\item The evaluated image of (E18.2) for every
ramified class-two input on the cubic. Its nonempty
fibres are (E18.3), and its exact continuous test
is (E18.1). The action vector $\eta(\tau)$
of E.19 is partially evaluated in E.35--E.37:
its two rational Kummer coordinates are known,
its elliptic projection is an explicit Ross
cup edge, and one primitive integral row has
value six. The numerical Ross value and
integral torsion remain.
\item The image and all fibres of the maps from
the intrinsic cubic metabelian and marked-depth
quotients to the inherited quotient $H$ of
(E20.5), for every first class.
\item The complete metabelian section sets for
the other coefficient primes and for finite
quotients involving several primes, including
all odd Fermat levels. This includes the first
level in the full profinite nilpotent tower
capable of selecting precisely the cuspidal
classes at an odd level.
\end{enumerate}
Task LM7-E.23(iv), the quartic local Kummer-order
law and additive defect at $3$ and $5$, is
evaluated by E.13--E.15 with the actual origin
verified in E.17. That conclusion leaves
(i)--(iii) on the page.
\end{statement}
\begin{proof}[State reached and the attempted steps]
We checked the supplied LM7 E.37, E.38, E.41
and E.44 proofs where used. Their later programs
were absent from the supplied expanded state;
the matrices and finite certificate reported
here were constructed and executed independently.
The earlier toric and local class-two criteria
are retained through the exact quotient maps
in E.1, E.15 and E.20.

For (i), replacing $\rho_3(F)$ in (E18.1)
by $\operatorname{Lie}_{\le3}(A_F)$ has no
justification: the displayed support error
uses the images of actual generators, whereas
the graded representation records only their
leading coordinates. We retained that term
and computed every nonempty fibre independently.
E.20 evaluates a ramified input with a finite
negative answer; E.16 gives an entire positive
unramified family. E.19 goes further on the
whole zero-abelian slice: the obstruction
factors through the class of (E19.3), so only
the action on degree-two commutators and one
normalized global vector remain. The original next
specific computation was the image of the
primitive compact determinant $w$ under the
actual Frobenius at the chosen prime over
seven, through $\gamma_2N/\gamma_4N$.
The peripheral path values of LM7 E.37
do not supply the images of both compact
free generators needed for this vector.
Passing to their linear Tate matrices gives
the coefficient group (E19.6), but leaves
the vector undetermined.

For (ii), the kernel of the exact map
$N\to H$ is $[I_y,I_y]\subset N$.
Thus the two kernel groups to be restored are
$[I_y,I_y]/N''$ and
$[I_y,I_y]/\gamma_3I_Q$, respectively.
A section of the target defines a pulled-back
extension by that specified kernel; its
splitting is the precise next question.
The local calculations of radical-class
towers in the supplied E.111, E.114 and
E.116 compute target sections. They do
not provide such splittings, and E.20
shows that passage back can fail already
in a finite class-three quotient.
The next approach is to compute the first
central graded kernel in these pullbacks
for the target families, retaining their
Frobenius correction words in E18.1.

For (iii), the even plus levels are settled
at their integral geometric abelian
two-primary quotient by E.5. For every
$n=3^r$ and $n=5^r$, E.16 proves that no
intrinsic nilpotent quotient at that
coefficient prime selects only cusps:
the constructed full pro-prime sections
have infinite-order compactified images.
The map from the full geometric profinite
group to a product of its maximal pro-prime
quotients kills elements detected by finite
groups involving several primes. A lift
from that product to the full metabelian
quotient must therefore solve its actual
kernel extension; none of the separate
coefficient calculations supplies it.
We have neither classified this lift
image nor proved a universal first
nilpotent level for odd $n$.
E.31 below makes one further comparison explicit:
the characteristic Haar summand embeds the entire
three-primary derived module into the full
profinite derived module. It supplies an actual
$\mathbb Z_3$-family of zero-abelian noncuspidal
full metabelian sections at the cubic level.
The complete mixed-prime image and its other
fibres remain unevaluated here.
E.32 also assembles the E.16 section with
the base cusp sections at all other primes
in the full pronilpotent product. It proves
that no full nilpotent level works for
$n=3^r$ or $5^r$. E.33 constructs the larger
full metabelian zero-abelian families for
every $n=3^r$, with their exact ranks;
E.34 proves the same full-metabelian family
statement for every $n=5^r$. E.35--E.36
identify the remaining elliptic class and
evaluate both rational Kummer coordinates.
E.37 gives the zero-abelian class
$\kappa_3(19)w$ a finite coefficient
obstruction at nineteen. E.38 constructs an
intrinsic zero-abelian class with lifts at
every finite coefficient level and no
continuous class-three lift. E.39 computes
the exact finite kernel lost by the compact
and unit cohomology comparisons. The actual
integral path class still requires correlated
division-point and unit-root data, so arbitrary
cochain representatives cannot fill that step.
These results retain the remaining complete image and
fibres as the open classification task.

For (iv), the exact density calculation
uses primitive integral Tate cocycles,
proves their full affine image, counts
finite Cartan and non-Cartan strata,
and controls the determinant-zero
boundary by measure tending to zero.
E.15 proves the polarization in the
genuine arithmetic action and its
restricted mixed correction. Thus the
local-order law, its scalar shifts,
and the failure-of-additivity pairing
are theorems; they do not evaluate
the different remaining input of (i).
The manuscript anchors for all four
tasks are \S51.16, folios 733v--736v,
scans 172--178.
\end{proof}

\begin{statement}{E.22}{Proved}{the precise comparison with unipotent Chabauty}
\label{st:E.22}
\label{E-Kim-comparison}
Fix a rational point or rational tangent on a
hyperbolic curve $Y/\mathbb Q$, a coefficient
prime $\ell$, and a nilpotency class $c$.
The actual continuous coefficient maps are
\[
 \pi_1(Y_{\overline{\mathbb Q}})
 \longrightarrow N^{(\ell)}
 \longrightarrow N^{(\ell)}/\gamma_{c+1}
 \longrightarrow U_{\ell,c}(\mathbb Q_\ell),
 \tag{E22.1}
\]
where $U_{\ell,c}$ is the class-$c$ unipotent
completion. They induce maps of section
classes, compatible with every restriction
to $G_{\mathbb Q_p}$ and with the point
Kummer maps. For a regular model
$\mathcal Y/\mathbb Z_S$, $\ell\notin S$,
the image of the $S$-integral finite-descent
set at $\ell$ is consequently contained
in the refined Chabauty--Kim locus for
every unipotent quotient.

Under the coefficient map $A\to A\otimes\mathbb Q_3$
of E.18, every finite-support obstruction
vanishes. Unrestricted support retains
the boundedness and convergence test.
In particular the finite integral local
obstruction modules (E18.4) and (E20.8)
map to zero under this specified
rationalization. The local point-image
conditions used by Chabauty--Kim remain
additional conditions on the target
cohomology classes, defined by the
commuting maps just stated.
\end{statement}
\begin{proof}
For a free geometric pro-$\ell$ group, map
its free generators to $1+X_i$ in the
completed noncommutative power-series
algebra. Modulo words of degree $c+1$,
the Magnus map sends group elements to
group-like units; finite logarithm and
exponential polynomials over $\mathbb Q_\ell$
give the corresponding unipotent group.
Ordered Hall products show directly
that its map on each graded piece is
the usual lattice inclusion into the
free rational Lie algebra. This
constructs (E22.1); for a proper curve,
quotient by the closed surface relation
and its unipotent image. The universal
property extends each arithmetic
automorphism, so the maps are equivariant.
Applying them to a continuous cocycle
commutes literally with restricting
its domain. Applying them to a path
torsor gives the corresponding
unipotent path torsor. These two
constructions prove the asserted
commuting Kummer maps.

An $S$-integral adelic point surviving
finite descent is the localization of
a Selmer section by the Harari--Stix
comparison used in D.1. Push this
section through (E22.1). At every finite
place its restriction lies in the
unipotent Kummer image of the chosen
local point, and outside $S$ in that
of integral points. Thus its global
class satisfies the defining refined
Selmer local conditions. Its value
at $\ell$ is in the scheme-theoretic
localization image, proving the
inclusion. This is the exact map
in Betts--Kumpitsch--L\"udtke,
Lemma 2.19, whose statement and proof
we checked~\cite{BettsKumpitschLudtke2026}.
It applies in particular to the odd
Fermat curves with the rational
tangent used above.

For the integral-to-rational assertion,
every $\Phi_{\mathfrak q}-q$ in E.18
is rationally invertible. A support
family at finitely many primes has
only finitely many rational inverse
vectors, so one power of three makes
all of them integral. The continuous
obstruction is therefore torsion and
maps to zero after tensoring with
$\mathbb Q_3$. More explicitly, a
continuous cochain from a compact
profinite domain to $A\otimes\mathbb Q_3$
has bounded image. Hence its cochain
complex is the filtered union of
the complexes with values in $3^{-n}A$,
and its cohomology is the localization
of the integral continuous cohomology.
For an infinite support sequence,
the same map requires a common
denominator and convergence to zero;
these requirements remain visible
in $\mathcal R(A\otimes\mathbb Q_3)$.
Thus the stated conclusion follows
from the literal coefficient map,
without reversing it.

For a concrete local comparison, the
plus Fermat quartic has no
$\mathbb Q_2$-point. Clear denominators
in a putative point with at least one
odd coordinate: a fourth power is
zero or one modulo $16$, so the sum
of the three coordinates' fourth
powers cannot be zero. Its odd-primary
class-two extensions nevertheless have
the sections constructed in E.17.
Pushing such a section to unipotent
cohomology does not place its
$2$-adic restriction in the image
of an empty local point set. This
is an explicit instance of the
Selmer condition in the same
commuting diagram.

The distinction between full and
finite-support classes is essential
in applying this comparison to the
Fermat tower. E.16 constructs full
pro-prime sections; it makes no
claim that their images satisfy
every Chabauty--Kim point-image
condition. E.9--E.12 instead use
the supplied local theorem for
full profinite sections to obtain
the local geometric conditions and
then prove global cuspidality.
All three conclusions follow along
their displayed maps. The anchors
are Lead 4 and \S51.16, folios
733v--736v, scans 172--178.
\end{proof}


% ---- E_symmetry ----
\subsection{A rational symmetry reduction of the cubic action class}

\begin{statement}{E.23}{Proved}{two projected arithmetic parameters and an eight-dimensional preimage}
\label{st:E.23}
\label{E-cubic-symmetry}
Retain the actual cubic notation of E.18--E.19. Put
$H=T_3E$, where $E=X_3$, and write
$V_{\mathbb Q_3}=H_{\mathbb Q_3}\oplus W_{\mathbb Q_3}$
for its rational compact and boundary summands. Let
$w\in\bigwedge^2V$ be the primitive compact determinant
vector of E.19, and define the actual submodule and quotient
\[
 B_w=[V,w]\subset A=\operatorname{Lie}_3(V),
 \qquad \mathcal Q=A/B_w.
 \tag{E23.1}
\]
Then $B_w$ is primitive of rank ten, its bracket map identifies
it with $V(1)$ as a $G_{\mathbb Q}$-module, and
$\mathcal Q$ is free of rank $320$.

Let $K=\mathbb Q(\zeta_3)$ and let
$\mathcal A=E[3]\rtimes\mu_6$ be the order-$54$ group of
$K$-automorphisms of $U=E\setminus E[3]$. Its action on
$\mathcal Q$ is independent of base paths. The subspace
\[
 D=(\mathcal Q_{\mathbb Q_3})^{\mathcal A}
\]
is $G_{\mathbb Q}$-stable and has the following rational
$G_{\mathbb Q}$-module description:
\[
 D\simeq H_{\mathbb Q_3}(2)\oplus
        \mathbb Q_3(3)\oplus
        \mathbb Q_3(3)\varepsilon_K,
 \tag{E23.2}
\]
where $\varepsilon_K$ is the quadratic character of $K/\mathbb Q$.
In particular $\dim D=4$ and $\operatorname{tr}(c|D)=0$.

For $\mathcal G=G_{\mathbb Q,\{3,\infty\}}$, the image
$\overline\eta$ of the actual class of E.19 satisfies
\[
 \overline\eta\in
 H^1_{\rm cts}(\mathcal G,D(-1))
 \lhook\joinrel\longrightarrow
 H^1_{\rm cts}(\mathcal G,\mathcal Q_{\mathbb Q_3}(-1)),
 \qquad
 \dim_{\mathbb Q_3}H^1(\mathcal G,D(-1))=2.
 \tag{E23.3}
\]
The two summands with potentially nonzero classes in (E23.2)
after twisting by $(-1)$ are $H_{\mathbb Q_3}(1)$ and
$\mathbb Q_3(2)\varepsilon_K$; each has one-dimensional
unramified global $H^1$. The untwisted quadratic-character
summand $\mathbb Q_3(2)$ has zero $H^1$.

More explicitly, the original rational class $\eta$ lies in
an at-most-eight-dimensional subspace of
$H^1(\mathcal G,A_{\mathbb Q_3}(-1))$: it is the inverse
image of the two-dimensional subspace in (E23.3) under the
map in the exact sequence
\[
 H^1(\mathcal G,V_{\mathbb Q_3})\longrightarrow
 H^1(\mathcal G,A_{\mathbb Q_3}(-1))\longrightarrow
 H^1(\mathcal G,\mathcal Q_{\mathbb Q_3}(-1)),
 \qquad \dim H^1(\mathcal G,V_{\mathbb Q_3})=6.
 \tag{E23.4}
\]
These are constraints on the actual arithmetic class and the
actual quotient maps. They do not evaluate its two projected
coordinates, its component in the displayed kernel, or its
integral torsion components. No identification of those
coordinates with specified polylogarithmic values is asserted.
\end{statement}
\begin{proof}
Write $w=u\wedge v$, with
$u=3f_1+U_1$ and $v=3f_2+U_2$ as in E.19. A coefficient
of $w$ is a unit. Equivalently a $2\times2$ minor of the
matrix $(u,v)$ is a unit, so $u,v$ extend to a
$\mathbb Z_3$-basis of $V$. In that basis the ten elements
$[e_i,[u,v]]$ are independent in the degree-three free Lie
module. For $e_i=u,v$ they are two individual Hall basis
elements up to sign. For each remaining $e_i$, the component
of multidegree $u v e_i$ has rank two, and
$[e_i,[u,v]]$ is a vector with a unit coordinate there.
Different indices have disjoint multidegrees. These unit
coordinates provide an integral left inverse on their span;
thus the span is primitive. Galois equivariance of the bracket
and $g(w)=\chi(g)w$ identify the span with $V(1)$.
This proves (E23.1), including its integral assertion.

Let $L_w$ be the inverse image of $\mathbb Z_3w$ under
$L=\gamma_2N/\gamma_4N\to\bigwedge^2V$.
An inner automorphism by an element of $N$ with image
$a\in V$ acts trivially on $A$, and sends a lift of
$b w$ to that lift plus $b[a,w]$. Consequently it acts
trivially on the full extension
\[
 0\longrightarrow\mathcal Q\longrightarrow
 L_w/B_w\longrightarrow\mathbb Z_3(1)w
             \longrightarrow0.
 \tag{E23.5}
\]
This is the exact quotient that removes the base-path
ambiguity. In particular an automorphism of the curve
preserving $w$ acts genuinely on (E23.5), although its
action on $N$ was originally specified only up to inner
conjugation. Commutators of lifts of commuting outer
automorphisms are inner, hence disappear on (E23.5) as well.

We check preservation of the compact line used here.
The generalized-Jacobian splitting $j_0$, with proper
projection $p j_0=[3]$, is unique because a homomorphism
from an abelian variety to a torus is zero. It therefore
commutes with the automorphisms of the pair $(E,E[3])$.
Translations act identically on its compact part; the
units in $\mu_6$ have determinant one on $T_3E$.
Thus $\mathcal A$ fixes $w$ literally. Its automorphisms
are defined over $K$ and are permuted by
$\operatorname{Gal}(K/\mathbb Q)$. The preceding inner
calculation proves that its action on (E23.5) commutes
with $G_K$ and is normalized by the entire $G_{\mathbb Q}$.
No chosen path has been declared Galois invariant.

Over $\mathbb Q_3$, average any lift of $w$ in (E23.5)
over $\mathcal A$. This is legitimate because $54$ is
invertible in $\mathbb Q_3$. It gives an
$\mathcal A$-fixed lift $\omega_{\mathcal A}$.
Replace it by
$(\omega_{\mathcal A}-c\omega_{\mathcal A})/2$.
Complex conjugation normalizes $\mathcal A$, so the
new lift is still fixed by $\mathcal A$ and now satisfies
$c\omega_{\mathcal A}=-\omega_{\mathcal A}$.
For every arithmetic $g$, both $g\omega_{\mathcal A}$
and $\chi(g)\omega_{\mathcal A}$ are
$\mathcal A$-fixed. Their difference, divided by $\chi(g)$,
is consequently a cocycle valued in $D(-1)$.
Changing the lift adds a coboundary, as in E.19, so it
represents exactly the projected class $\overline\eta$.
The averaging idempotent is $G_{\mathbb Q}$-equivariant
because Galois permutes the elements of $\mathcal A$.
It splits the inclusion $D\subset\mathcal Q_{\mathbb Q_3}$,
so the cohomology arrow in (E23.3) is injective.

For completeness, we give the finite character calculation
and its arithmetic interpretation. On the nine cusp points
identify $E[3]$ with $\mathbb F_3^2$. A generator of
$\mu_6$ acts by $-R$, where $R$ is the matrix in E18.9,
and complex conjugation acts by $C_E$. An element
$(t,\alpha)$ acts on the cusp set by $x\mapsto t+\alpha x$.
Its trace on $W_{\mathbb Q_3}$ is the number of fixed
cusp points minus one. Its trace on $H_{\mathbb Q_3}$
is $\alpha+\alpha^{-1}$, interpreted in the two-dimensional
CM representation. For a complex-conjugation coset this
compact trace is zero, and the boundary trace has an
additional minus sign from the Tate twist.
For any such operator $b$, the degree-three trace is
\[
 \operatorname{tr}(b|\operatorname{Lie}_3V)
 =\frac{\operatorname{tr}(b|V)^3-
                    \operatorname{tr}(b^3|V)}{3}.
 \tag{E23.6}
\]
Subtract $\operatorname{tr}(b|V)\det(b|H)$ for $B_w$.
Averaging over the $54$ affine transformations gives
\[
 \dim V^{\mathcal A}=0,\quad
 \dim A_{\mathbb Q_3}^{\mathcal A}=4,\quad
 \dim D=4,\quad \operatorname{tr}(c|D)=0.
\]
Keeping the number of compact and boundary factors gives
\[
\begin{array}{c|rrrr}
 \text{multidegree in }\mathcal Q&H^3&H^2W&HW^2&W^3\\\hline
 \text{invariant dimension}&0&0&2&2\\
 \text{trace of }c&0&0&0&0 .
\end{array}
 \tag{E23.7}
\]
The executed program \texttt{class\_three\_symmetry.py}
records every character row using integer arithmetic.
The independent character check accompanying this proof
repeats the $54$ affine transformations and their conjugation
coset, and checks each multidegree separately. It also records
the correction to an initial guess that all four invariant
coordinates had multidegree $HW^2$; (E23.7) is the checked
replacement.

All cusps are rational over $K$, so $G_K$ acts on $W$ by
the scalar $\chi$. The two $HW^2$ invariants couple the
two distinct CM eigenlines of $H$ to their inverse
$\mu_6$-characters. Each coupling has multiplicity one:
the nonzero characters of the translation group have one
$\mu_6$-orbit of length six and one of length two; only
the length-six orbit supplies the two primitive order-six
characters of $H$. Thus their restrictions to $G_K$ are
exactly the two CM characters times $\chi^2$.
Complex conjugation exchanges the two lines, identifying
this representation over $\mathbb Q_3$ with
$H_{\mathbb Q_3}(2)$. This can also be checked after the
quadratic coefficient extension splitting the CM algebra;
changing its exchange convention gives the quadratic twist,
which is isomorphic by the CM endomorphism on which
complex conjugation acts as minus one.
The $W^3$ invariants restrict to two copies of $\chi^3$
on $G_K$; their zero conjugation trace gives one copy of
$\chi^3$ and one of $\chi^3\varepsilon_K$. This proves
(E23.2), with the actual restriction and descent maps.

Finally all these modules are unramified outside three,
and after twisting by $(-1)$ their good-Frobenius eigenvalues
have positive weights. Their invariant vectors are zero.
The free-pro-three arithmetic source and the order-two
averaging of E.16 give zero $H^2$. The global Euler
characteristic therefore makes $\dim H^1$ the dimension
of the minus eigenspace of complex conjugation.
For $H(1)$ this is one; for $\mathbb Q_3(2)$ it is zero;
for $\mathbb Q_3(2)\varepsilon_K$ it is one. This proves
(E23.3). On $V$, the compact trace is zero and the boundary
trace is $-(3-1)=-2$, since three cusps are real.
Its dimension is ten, so the same calculation gives
$\dim H^1(\mathcal G,V_{\mathbb Q_3})=6$.
The long exact sequence of (E23.1), twisted by $(-1)$,
gives (E23.4). Its kernel is the image of that
six-dimensional space, so the inverse image of the
specified two-dimensional subspace has dimension at most
eight. This proves the announced constraint without
assigning a value to the arithmetic class. The manuscript
anchor remains \S51.16, folios 733v--736v, scans 172--178,
and task LM7-E.23(i).
\end{proof}

% ---- E_integral_symmetry ----
\begin{statement}{E.24}{Proved}{the integral elliptic-involution constraint}
\label{st:E.24}
\label{E-integral-symmetry}
Retain the integral class $[\eta]\in H^1(G_{\mathbb Q,\{3,\infty\}},
A(-1))$ of E.19. Let $\iota$ be the automorphism
$[x:y:z]\mapsto[y:x:z]$ of the Fermat cubic with origin
$O=[1:-1:0]$. It induces an involution on the actual
based pro-three group, commuting with its arithmetic action.
Put $A^+=\ker(\iota-1:A\to A)$. The class $[\eta]$
is in the image of the specified inclusion
\[
 H^1(G_{\mathbb Q,\{3,\infty\}},A^+(-1))
 \longrightarrow H^1(G_{\mathbb Q,\{3,\infty\}},A(-1)),
 \tag{E24.1}
\]
and the source is
\[
 \mathbb Z_3^{81}\oplus
 (\mathbb Z/3)^{31}\oplus
 (\mathbb Z/9)^{12}\oplus
 (\mathbb Z/27)^4.\tag{E24.2}
\]
The inclusion in (E24.1) is split by the integral projector
$(1+\iota)/2$. This assertion constrains the genuine
integral class, including its torsion coordinates.
\end{statement}
\begin{proof}
On the complete elliptic curve, $\iota=[-1]$.
At $O$ the parameter $z/x$ has derivative $-1$
under $\iota$. The compatible $3^r$th roots of $-1$
can all be chosen to equal $-1$. They are rational
and multiply to $1$ when the tangent change is
performed twice. The corresponding tangential
path therefore supplies an order-two based
automorphism commuting with $G_{\mathbb Q}$.
This constructs the claimed action, including the
comparison between its two tangent base points.

Use the ordered degree-one basis and the matrix $U$
of E.18. On cusp homology $\iota$ is the permutation
matrix $P$ with column permutation
$(3,4,5,0,1,2,7,6)$, in zero-based indexing. Its
matrix on $V$ is
\[
 I_V=\begin{pmatrix}
 -I_2&0\\(-U-PU)/3&P
 \end{pmatrix}.\tag{E24.3}
\]
Every entry is integral. The map is $-1$ on the
compact lift vectors $3f_i+U_i$, so it fixes
$w=(3f_1+U_1)\wedge(3f_2+U_2)$.
Let $I_A=\operatorname{Lie}_3(I_V)$.
It commutes with the Frobenius matrix $T$ and
the conjugation matrix $J$ of E.18.

Begin with any lift $\omega_0$ of $w$ to
$L=\gamma_2N/\gamma_4N$. Apply the commuting
projectors $(1+\iota)/2$ and $(1-c)/2$ to obtain
$\omega$ with $\iota\omega=\omega$ and
$c\omega=-\omega$. Its image in $M$ is still $w$.
For every arithmetic $g$ the formula
\[
 \eta(g)=\chi(g)^{-1}(g\omega-\chi(g)\omega)
\]
now has values in $A^+(-1)$, since $g$ commutes
with $\iota$. Replacing the lift changes $\eta$
by the actual coboundary of the difference of the
two lifts in $A(-1)$. Hence (E24.1) follows for
the original cohomology class.

The degree-three free Lie trace identity gives
$\operatorname{tr}(I_A)=-2$ and
$\operatorname{rank}A^+=164$. Also
$\operatorname{tr}(J|A^+)=-2$. In the free
pro-three presentation used in E.18--E.19, the
two generators are interchanged by $c$.
Consequently the same restriction-and-coboundary
calculation as in E.19 presents the group in
(E24.1) as the quotient of $A^+$ by the image of
$(T-7I)(1-J)$ on $A^+$; the factor $7$ is a unit
and is the cyclotomic twist in that calculation.

For an exact computation in the ambient integral
lattice use the matrix
\[
             (T-7I)(1-J)(1+I_A).
\]
Multiplication by $1+I_A$ has image $A^+$, since
two is a $3$-adic unit. The executed Hall-basis
program \texttt{class\_three\_symmetry.py} gives
the following elementary-divisor valuations:
\[
\begin{array}{c|rrrr}
\text{valuation}&0&1&2&3\\\hline
\text{multiplicity}&36&31&12&4.
\end{array}
\]
All 83 nonzero pivots are certified modulo $3^{10}$.
The exact rank cannot exceed
$(164-(-2))/2=83$, the dimension of the
$(-1)$-eigenspace of $J$ in $A^+$.
Thus these are all its nonzero elementary divisors.
The complement $A^-$ is a direct summand of the
ambient lattice and the relation image lies in
$A^+$; removing that free complement leaves
free rank $164-83=81$ and exactly the torsion
factors in (E24.2).

The rational averaging over the larger group in
E.23 has a denominator divisible by three.
Its comparison map is the rationalization of
the integral extension followed by the quotient
$A\to A/[V,w]$. The integral projector used here
has denominator two. The two conclusions are
therefore recorded along their respective maps;
neither calculation evaluates the remaining
arithmetic cocycle vector. The manuscript anchor
is \S51.16, folios 733v--736v, scans 172--178.
\end{proof}

% ---- E_symmetry_refinement ----
\begin{statement}{E.25}{Proved}{the remaining rational class lies in an exact five-dimensional subspace}
\label{st:E.25}
\label{E-five-parameters}
Let $A^+$ be the integral elliptic-involution submodule of
E.24 and retain $D$ and $\mathcal Q$ from E.23. Let
$\mathcal Q^+$ and $V^+$ denote the same involution's fixed
submodules. The actual rational class $\eta$ belongs to the
inverse image
\[
 \mathcal H=\operatorname{preimage}\left(
 H^1(\mathcal G,D(-1))\subset
 H^1(\mathcal G,\mathcal Q^+_{\mathbb Q_3}(-1))\right)
\]
under the specified map from
$H^1(\mathcal G,A^+_{\mathbb Q_3}(-1))$. There is a
short exact sequence of rational vector spaces
\[
 0\longrightarrow H^1(\mathcal G,V^+_{\mathbb Q_3})
   \longrightarrow\mathcal H
   \longrightarrow H^1(\mathcal G,D(-1))
   \longrightarrow0,
 \qquad \dim\mathcal H=3+2=5.
 \tag{E25.1}
\]
Thus the rational action class remaining after E.19 is
confined to this actual five-dimensional subspace. The
integral torsion restriction remains the separate one
computed in E.24. The five coordinates are not evaluated
by this dimension calculation.
\end{statement}
\begin{proof}
The involution fixes $w$, so the equivariant isomorphism
$B_w(-1)\simeq V$ of E.23 identifies
$B_w^+(-1)$ with $V^+$. Its projector $(1+\iota)/2$ is
integral and exact, giving
\[
 0\longrightarrow V^+_{\mathbb Q_3}
   \longrightarrow A^+_{\mathbb Q_3}(-1)
   \longrightarrow\mathcal Q^+_{\mathbb Q_3}(-1)
   \longrightarrow0.
\]
The right-hand representation has zero invariants, since
its good-Frobenius eigenvalues have positive weights.
The left-hand $H^2$ is zero by the unramified arithmetic
source calculation of E.16. Its cohomology sequence is
therefore short exact in degree one.

On $V$, the elliptic involution has compact trace $-2$
and boundary trace zero; hence $\dim V^+=4$.
Complex conjugation and its product with that involution
both have trace $-2$ on $V$: on the cusp set each fixes
three points, with the boundary Tate sign, while its
compact trace is zero. Thus
\[
 \operatorname{tr}(c|V^+)=\frac{-2-2}{2}=-2.
\]
The same Euler-characteristic calculation, with zero
$H^0$ and $H^2$, gives $\dim H^1(\mathcal G,V^+)=3$.
Because $\iota$ is an element of $\mathcal A$, the
subspace $D$ lies in $\mathcal Q^+_{\mathbb Q_3}$.
The $\mathcal A$-averaging projector commutes with Galois
and with $\iota$; hence its inclusion is split and so
is injective on $H^1$. Pulling back this two-dimensional
subspace through the short exact cohomology sequence
proves (E25.1). E.24 puts the actual $\eta$ in the source,
and E.23 puts its image in the specified subspace, so it
indeed belongs to $\mathcal H$. The independent finite
character program records the two traces used here.
All operations are the displayed projector, bracket,
quotient, and inverse-image maps; no replacement of the
arithmetic class by an unspecified graded action occurs.
The anchor is again \S51.16, folios 733v--736v, scans
172--178, and LM7-E.23(i).
\end{proof}

% ---- E_cover_comparison ----
\begin{statement}{E.26}{Proved}{the exact kernel in the once-punctured comparison}
\label{st:E.26}
\label{E-elliptic-cover-comparison}
Let $P$ be the geometric pro-three group of $E\setminus\{O\}$,
with free generators $x,y$, and let
$N=\ker(P\to(\mathbb Z/3)^2)$. The covering
$[3]:E\setminus E[3]\to E\setminus\{O\}$ identifies
this $N$ with the geometric group used in E.18--E.25.
The natural comparison of its class-two quotient with
the metabelian quotient of $P$ is the diagram
\[
 N/\gamma_3N\ \longrightarrow\
 N/(P''\gamma_3N)\ \longleftarrow\ N/P''
 \ \lhook\joinrel\longrightarrow\ P/P'' .
 \tag{E26.1}
\]
The first arrow has a nonzero kernel, detected in a
Heisenberg group of order $27$. In particular, the
quotient data on the right of (E26.1) do not specify
all the elements of the quotient on the left through
these comparison maps.
\end{statement}
\begin{proof}
The degree-nine isogeny is finite \emph{\'etale} on the
displayed punctured curves. Its deck group is $E[3]$,
and the monodromy is the compact homology modulo three.
For the once-punctured elliptic group its abelianization
is precisely compact homology. This proves the kernel
description and specifies the covering identification.
Since $P'\subset N$, one has $P''\subset N$ and all
the maps in (E26.1) are the stated continuous quotient
or inclusion maps.

Put $c=xyx^{-1}y^{-1}$ and
$z=[c,xcx^{-1}]$. Both entries in this commutator
belong to $P'$, so $z\in P''$.
For an explicit free basis of $N$, use the cover of the
two-petal graph with vertices $(i,j)\in\mathbb F_3^2$.
Its positive $x$-edges increment $i$ and its positive
$y$-edges increment $j$. Choose as spanning tree the
two edges $(i,0)\to(i+1,0)$ for $i=0,1$ and the six
edges $(i,j)\to(i,j+1)$ for $i=0,1,2$, $j=0,1$.
The ten remaining oriented edges give a free basis
$e_0,\ldots,e_9$ after the tree is contracted.
Order them lexicographically by $(i,j,a)$, with
$a=1$ for $x$ and $a=2$ for $y$. Directly tracing the
two closed words gives
\[
 [c]_{N^{\rm ab}}=-e_0,\qquad
 [xcx^{-1}]_{N^{\rm ab}}=-e_3.
\]
Their exterior product is the primitive vector
$e_0\wedge e_3$. Hence $z$ is nonzero in
$\gamma_2N/\gamma_3N$.

The finite certificate sends $e_0$ and $e_3$ to
$(1,0,0)$ and $(0,1,0)$ in the group on
$\mathbb F_3^3$ with multiplication
\[
 (a,b,d)(a',b',d')=(a+a',b+b',d+d'+ab'),
\]
and sends every other basis generator to the identity.
The image of $z$ is $(0,0,1)$. The executed
\texttt{elliptic\_cover\_kernel.py} reconstructs the
Schreier words and verifies this exact value. Since
the finite target has class two, it also proves that
$z\notin\gamma_3N$ profinitely. Its image in the
middle quotient of (E26.1) is zero because $z\in P''$.
This proves the asserted nonzero kernel.

This comparison is the precise obstacle encountered
in trying to evaluate the remaining class of E.25
directly from the once-punctured elliptic power series
of Nakamura, Theorem A (4.12), and the path-torsor
formulas of Wojtkowiak~\cite{Nakamura1995Elliptic,Wojtkowiak2004Paths}.
Those primary formulas were opened and their domains
checked: the former is a Galois character construction
over the field of all prime-power division points,
and the latter begins with paths on
$E\setminus\{O\}$. Applying them to the actual
nine-puncture extension requires an additional
compatible construction retaining the kernel just
exhibited. The finite example does not assert that
no further comparison using these formulas exists;
it proves exactly what the displayed quotient maps
retain and remove. The anchor is \S51.16, folios
733v--736v, scans 172--178, and LM7-E.23(i).
The supplied LM7 tower return, E.50, gives a larger
kernel calculation for this same cover. The argument
above is an independent finite certificate for its
degree-two part; E.30 verifies its degree-three kernel
and locates the remaining arithmetic class in that kernel.
\end{proof}

% ---- E_explicit_kernel ----
\begin{statement}{E.27}{Proved}{the canonical rational splitting and the three Kummer coordinates}
\label{st:E.27}
\label{E-explicit-kernel}
Let $\mathcal G=G_{\mathbb Q,\{3,\infty\}}$, and retain
$\mathcal H$, $B_w$, $\mathcal A$, and $A^+$ from E.23--E.25.
The rational averaging operator
\[
 e_{\mathcal A}=\frac1{54}\sum_{a\in\mathcal A}a
       \quad\text{on }A_{\mathbb Q_3}
\]
gives the following specified direct-sum decomposition:
\[
 \mathcal H=
 \operatorname{im}H^1(\mathcal G,V^+_{\mathbb Q_3})
 \ \oplus\
 H^1\bigl(\mathcal G,(A_{\mathbb Q_3})^{\mathcal A}(-1)\bigr).
 \tag{E27.1}
\]
The projection to the second summand is induced by
$e_{\mathcal A}$. In the cusp basis of E.18 put
\[
 b_0=e_{A_0}+e_{B_0},\quad
 b_1=e_{A_1}+e_{B_1},\quad
 b_2=e_{A_2}+e_{B_2},\quad
 b_3=e_{C_1}+e_{C_2}.
\]
There is an integral decomposition with its indicated bases,
\[
 V^+=
 \mathbb Z_3(1)\{b_0,b_1+b_2,b_3\}
 \ \oplus\
 \mathbb Z_3(1)\varepsilon_K\{b_1-b_2\},
 \tag{E27.2}
\]
and
\[
 H^1(\mathcal G,V^+)=\mathbb Z_3^3\oplus\mathbb Z/3.
 \tag{E27.3}
\]
After rationalization the first three coordinates in
(E27.1) are therefore uniquely specified scalar multiples
of $\kappa_3(3)$ in the three displayed Tate directions,
mapped to $A(-1)$ by $v\mapsto[v,w]$.
The other two coordinates are the classes in the two
one-dimensional cohomology summands identified in E.23.
This proposition fixes their normalization without
evaluating the actual class. E.28 eliminates one
projected coordinate, E.29 relates two Tate coordinates,
and E.35--E.36 identify the elliptic class and evaluate
the remaining two rational Kummer coordinates.
\end{statement}
\begin{proof}
The action of $\mathcal A$ on the graded module $A$ is
genuine and commutes with arithmetic up to the permutation
of elements of $\mathcal A$. Thus its average commutes
with $\mathcal G$. It preserves $B_w$, on which the
fixed subspace is
$V_{\mathbb Q_3}^{\mathcal A}(1)=0$ by E.23.
Applying this exact rational projector to
$0\to B_w\to A\to\mathcal Q\to0$ gives an isomorphism
\[
       (A_{\mathbb Q_3})^{\mathcal A}
                \xrightarrow{\sim}
       (\mathcal Q_{\mathbb Q_3})^{\mathcal A}=D.
\]
The isomorphism is the actual quotient map. In E.25,
the image of a class in $\mathcal H$ lies in $H^1(D(-1))$.
Subtracting its $e_{\mathcal A}$-projection therefore
places it in the kernel of the quotient map on $H^1$.
That kernel is precisely the injected image of
$H^1(V^+_{\mathbb Q_3})$ by the exact sequence of E.25.
The two summands have zero intersection, since
$e_{\mathcal A}$ is zero on the first and the identity
on the second. This proves (E27.1) with its projection.

For (E27.2), the compact coordinates of an
$I_V$-fixed vector satisfy $-h=h$ and hence vanish.
Its toric coordinates are exactly the four paired
cusp vectors $b_i$ above. Every cusp is rational over
$K=\mathbb Q(\zeta_3)$, so $G_K$ acts on them by
$\chi$. Complex conjugation satisfies
\[
 cb_0=-b_0,\quad cb_3=-b_3,\quad
 cb_1=-b_2,\quad cb_2=-b_1.
\]
The basis change from $(b_1,b_2)$ to
$(b_1+b_2,b_1-b_2)$ has determinant $-2$, a
$3$-adic unit. The first three directions have the
ordinary Tate action, and the last has its quadratic
twist, proving the integral assertion.

Kummer theory over $\mathbb Z[1/3]$ gives
$H^1(\mathcal G,\mathbb Z_3(1))=\mathbb Z_3\kappa_3(3)$:
the units are $\{\pm3^n\}$ and the class group is zero.
For the twisted line restrict to $G_{K,\{3\}}$ and use
the exact order-two averaging projector. The ring
$\mathbb Z[\zeta_3]$ is Euclidean for its norm; its
units after inverting three have $3$-adic Kummer
completion
\[
                 \mu_3\oplus\mathbb Z_3\{1-\zeta_3\}.
\]
Set $u=\zeta_3(1-\zeta_3)$. The exact identities are
$u^2=-3$ and $c(u)=-u$. The Kummer class of $-1$
is zero at coefficient prime three, so the line of
$u$ is fixed by $c$; it is the free direct summand.
On the torsion summand $\mu_3$, conjugation acts by
inversion. The minus eigenspace is consequently
exactly $\mu_3$, and order-two descent identifies it
with $H^1(\mathcal G,\mathbb Z_3(1)\varepsilon_K)$.
This proves (E27.3).

These computations specify the coefficient group and
the bracket, quotient, and projector maps used to
separate the coordinates. The averaging denominator
$54$ is allowed only in the rational splitting
(E27.1); the integral torsion in E.24 and (E27.3)
is retained before that map. The anchor is
\S51.16, folios 733v--736v, scans 172--178,
and task LM7-E.23(i).
\end{proof}

% ---- E_no_boundary_cubic ----
\begin{statement}{E.28}{Proved}{the pure boundary cubic component vanishes}
\label{st:E.28}
\label{E-no-pure-boundary-cubic}
Use the cubic notation of Propositions~\ref{E-cubic-symmetry}
and~\ref{E-five-parameters}. Rationally write the actual
Galois-stable splitting of homology as
\[
 V_{\mathbb Q_3}=H\oplus W,
 \qquad H=T_3E\otimes\mathbb Q_3,
 \qquad w=u\wedge v\in\bigwedge^2H.
\]
The image of the arithmetic class
$\eta\in H^1(G_{\mathbb Q,S_3},A_{\mathbb Q_3}(-1))$
under the canonical projection
\[
 A_{\mathbb Q_3}=\operatorname{Lie}_3(H\oplus W)
                  \longrightarrow\operatorname{Lie}_3(W)
 \tag{E.28.1}
\]
is zero. Consequently the projected symmetry-invariant class
of E.23 lies entirely in
\[
 H^1(G_{\mathbb Q,S_3},H(1)),
 \qquad \dim_{\mathbb Q_3}H^1(G_{\mathbb Q,S_3},H(1))=1.
 \tag{E.28.2}
\]
In the canonical rational direct sum of E.27, the actual
class therefore belongs to a four-dimensional subspace:
the three Kummer coordinates in
$H^1(G_{\mathbb Q,S_3},B^+_{\mathbb Q_3}(-1))$
and the single elliptic coordinate in (E.28.2).
The $\mathbb Q_3(2)\epsilon_K$ coordinate which was allowed
by symmetry alone is zero. These assertions determine a
vanishing projection, without assigning a value to the
remaining four coordinates or to integral torsion.
\end{statement}
\begin{proof}
Let $\mathfrak n_{\leq3}$ be the rational class-three
Malcev Lie algebra of the free geometric pro-$3$ group
$N$. Its associated graded Lie algebra is the free Lie
algebra on $V_{\mathbb Q_3}$ through degree three.
The logarithm identifies
$(\gamma_2N/\gamma_4N)\otimes\mathbb Q_3$
with $\gamma_2\mathfrak n_{\leq3}$; this part is
abelian because a commutator of two degree-two terms
has degree at least four. This is an equivariant
identification, so it computes the extension defining
$\eta$ itself.

Choose lifts $a,b\in\mathfrak n_{\leq3}$ of $u,v$,
and choose lifts of a basis of $W$. They identify the
filtered free Lie algebra with its free graded model.
Set $\omega=[a,b]$. It is a lift of $w$. For
$g\in G_{\mathbb Q,S_3}$, stability of the actual
compact summand $H$ gives
\[
 ga=\alpha a+\beta b+q_a,
 \qquad gb=\gamma a+\delta b+q_b,
 \qquad q_a,q_b\in\gamma_2\mathfrak n_{\leq3},
\]
where $\alpha\delta-\beta\gamma=\chi_3(g)$.
Since $[q_a,q_b]=0$ in class three, the exact bracket
formula gives
\[
 g\omega-\chi_3(g)\omega
   =[\alpha a+\beta b,q_b]
       +[q_a,\gamma a+\delta b].
 \tag{E.28.3}
\]
Only the degree-two parts of $q_a,q_b$ contribute.
Every term of (E.28.3) therefore has at least one
factor in $H$. Under the multigrading it belongs to
\[
 \operatorname{Lie}_3(H)\oplus(H^2W)\oplus(HW^2),
\]
and its projection to $\operatorname{Lie}_3(W)$ is
zero. Dividing (E.28.3) by $\chi_3(g)$ gives the
cocycle of the extension defining $\eta$. A different
lift of $w$ changes that cocycle by a coboundary in
$A_{\mathbb Q_3}(-1)$. Thus (E.28.1) kills its
cohomology class independently of all these lifts.

The submodule $B=[V,w]$ also has at least two compact
factors in each bracket, so (E.28.1) factors through
$Q=A/B$. The geometric group
$\mathcal G=E[3]\rtimes\mu_6$ preserves the actual
homology splitting $H\oplus W$. Its rational
averaging projector therefore commutes with the
multigrading projections. E.23 identifies
\[
 (Q_{\mathbb Q_3})^{\mathcal G}
       =D_{HW^2}\oplus D_{W^3},
 \qquad
 D_{HW^2}\simeq H(2),
 \quad
 D_{W^3}\simeq\mathbb Q_3(3)
                     \oplus\mathbb Q_3(3)\epsilon_K.
\]
These are Galois-stable direct summands. The zero
projection just proved kills the whole $D_{W^3}$
component of the projected class. Its remaining
component, after twist $(-1)$, is exactly (E.28.2).
The cohomology dimension one is the same Euler
characteristic and $H^2$-vanishing computation used
for this summand in E.23.

Finally $B^{\mathcal G}_{\mathbb Q_3}=0$, so the
averaging projector on $A_{\mathbb Q_3}$ identifies
$A^{\mathcal G}_{\mathbb Q_3}$ with
$Q^{\mathcal G}_{\mathbb Q_3}$, as established in
E.27. That projector also preserves the multigrading.
Its canonical direct-sum description therefore
restricts to the three-dimensional Kummer kernel
and the one-dimensional summand (E.28.2), giving
dimension four. This uses the rational compact
splitting and rational Malcev logarithm; it imposes
no unproved integral vanishing on the torsion
recorded in E.24 and E.27.
\end{proof}

% ---- E_mu6_three_parameters ----
\begin{statement}{E.29}{Proved}{a rational based $\mu_6$ action and three remaining coordinates}
\label{st:E.29}
\label{E-mu6-three-parameters}
Let $U=E\setminus E[3]$ be the cubic curve with the
specified tangential base point at $O$, and retain its
rational class-three Malcev group. The origin-fixing
geometric group $\mu_6$ has an actual based action on
this group, normalized by the arithmetic action:
\[
       g\beta g^{-1}=\beta^{\chi(g)\bmod6}
       \qquad(g\in G_{\mathbb Q}),
 \tag{E.29.1}
\]
where $\beta$ lifts a generator of $\mu_6$ and has
order six. Its cube is the rationalization of the
based involution of E.24. Thus the actual class
$[\eta]$ lies in the image of
\[
 H^1(G_{\mathbb Q,S_3},A_{\mathbb Q_3}^{\mu_6}(-1))
       \longrightarrow
 H^1(G_{\mathbb Q,S_3},A_{\mathbb Q_3}(-1)).
 \tag{E.29.2}
\]

Together with E.23, E.27 and E.28, this places the
class in the canonical rational direct sum
\[
 \operatorname{image}
 H^1(G_{\mathbb Q,S_3},B_{\mathbb Q_3}^{\mu_6}(-1))
 \ \oplus\
 H^1(G_{\mathbb Q,S_3},H(1)),
 \tag{E.29.3}
\]
of dimensions $2+1=3$. Here
$B_{\mathbb Q_3}^{\mu_6}(-1)\simeq\mathbb Q_3(1)^2$.
The first two coordinates are multiples of the
Kummer class of $3$; the third is the remaining
elliptic cohomology coordinate. Their actual values
are not determined by this proposition. The
denominator six in the based action and averaging
is used rationally, so no integral torsion coordinate
is discarded.
\end{statement}
\begin{proof}
We construct the based action before quotienting by
inner path ambiguity. The origin-fixing automorphism
group of the $j=0$ elliptic curve is the group scheme
$\mu_6$ over $\mathbb Q$. It preserves $E[3]$.
Consequently the quotient stack
$\mathcal U=[U/\mu_6]$ is defined over $\mathbb Q$,
and $U\to\mathcal U$ is a finite etale torsor of
stacks. On geometric fundamental groups it gives
an extension by the six-element geometric group.

Near $O$, a formal parameter can be linearized for
this finite action in characteristic zero, keeping
its specified first-order tangent coordinate. To
see this explicitly, over a field containing
$\mu_6$ average any parameter $t_0$ with nonzero
derivative as
\[
       t=\frac16\sum_{j=0}^5
                 \zeta_6^{-j}\alpha^{j*}(t_0),
\]
after fixing the generator so that its derivative
is $\zeta_6$. Scaling the first coefficient to one
retains the given tangent. This is the eigenspace
projector, so the coordinate and its transformation
rule descend with the tangent character. Its
quotient coordinate is $q=t^6$. The map on nonzero
tangent directions sends the chosen direction
$t=1$ to $q=1$.

Let $d$ be the positively oriented geometric
meridian of this quotient punctured normal line,
and let $z_O$ be the meridian of the original cusp.
The actual covering $q=t^6$ gives
\[
                          d^6=z_O.
 \tag{E.29.4}
\]
The image of $d$ in the deck quotient is a generator
of $\mu_6$. Write $D$ for conjugation by $d$ on the
rational Malcev group of $N$. It preserves $z_O$
and satisfies $D^6=\operatorname{Int}(z_O)$.
The rational power $z_O^{-1/6}$ is defined by
$\exp(-\log(z_O)/6)$. Set
\[
                 \beta=\operatorname{Int}(z_O^{-1/6})D.
 \tag{E.29.5}
\]
Since $D$ fixes $z_O$, (E.29.4) gives $\beta^6=1$.
Its action on compact homology is the faithful
order-six elliptic automorphism, so its order is
exactly six.

Here is the arithmetic compatibility, including
the base-point comparison. The arithmetic action
on $N$ is the one from its specified tangent at
$O$. Sending that tangent through $q=t^6$ gives
the quotient tangent $1$, and the induced map of
root fibre functors identifies their arithmetic
splittings. Indeed $q^{1/n}$ pulls back to
$t^{6/n}$, and both fixed-root splittings fix these
compatible roots. A unit coordinate change with
constant coefficient one gives the same section.
Therefore, in this based orbifold extension, the
arithmetic section conjugates $d$ to $d^{\chi(g)}$
and $z_O$ to $z_O^{\chi_3(g)}$ in its rational
pro-three image. These equalities refer to the
actual tangential splitting, rather than an
arbitrary representative of the outer action.

Choose $r\in\{1,5\}$ congruent to $\chi(g)$
modulo six. In the local inertia extension,
\[
 \operatorname{Int}(d^{\chi(g)})
    =\operatorname{Int}
       \left(z_O^{(\chi_3(g)-r)/6}\right)D^r.
\]
This follows by subtracting the exponent $r$ in
the full local inertia group and using (E.29.4);
the remaining exponent is the stated three-adic
one after passage to the Malcev group. Combining
it with the conjugate of (E.29.5) gives
\[
 \begin{aligned}
 g\beta g^{-1}
 &=\operatorname{Int}(z_O^{-\chi_3(g)/6})
       \operatorname{Int}
          (z_O^{(\chi_3(g)-r)/6})D^r\\
 &=\operatorname{Int}(z_O^{-r/6})D^r
   =\beta^r.
 \end{aligned}
\]
This proves (E.29.1). In particular the action
commutes with $G_K$ for $K=\mathbb Q(\zeta_3)$,
and complex conjugation inverts its generator.
The cube of (E.29.5) is the order-two correction
for the same local normal-line extension. That
correction is unique because its rational inertia
has no two-torsion. It is therefore exactly the
rationalization of E.24's based involution.

The action on the associated graded module $A$
is the geometric $\mu_6$ action used in E.23:
inner conjugation acts trivially on degree three
modulo degree four. The line $\mathbb Q_3w$
is fixed by $\mu_6$, because the determinant of
the origin-fixing automorphism on compact
homology is one. Average any lift of $w$ in
$L_w\otimes\mathbb Q_3$ under the actual finite
group generated by $\beta$. Formula (E.29.1)
shows that this projector commutes with Galois.
The extension cocycle of the averaged lift
therefore has values in $A^{\mu_6}(-1)$.
Changing the lift changes only a coboundary,
proving (E.29.2).

We compute its kernel under projection to $Q=A/B$.
The geometric $\mu_6$ orbits on the nine cusps
are $\{O\}$, the two nonzero points killed by
$1-\zeta_3$, and the remaining six points.
The first subgroup is stable under conjugation
because $1-\zeta_3^2$ differs from $1-\zeta_3$
by a unit. All three orbits are consequently
Galois-stable. Their three orbit sums in cusp
homology have the single total-sum relation.
The compact summand has no $\mu_6$ invariants.
Thus
\[
 V_{\mathbb Q_3}^{\mu_6}
       =W^{\mu_6}\simeq\mathbb Q_3(1)^2,
 \qquad
 B_{\mathbb Q_3}^{\mu_6}(-1)
       \simeq V_{\mathbb Q_3}^{\mu_6}.
\]
This has an explicit basis in the notation of E.27.
On the Fermat cubic the origin is $O=C_0=[1:-1:0]$,
and the order-three automorphism
$[X:Y:Z]\mapsto[X:Y:\zeta_3 Z]$ fixes each $C_i$
and sends $A_i$ to $A_{i-1}$ and $B_i$ to $B_{i-1}$,
with subscripts modulo three. The involution interchanges
$A_i,B_i$ and interchanges $C_1,C_2$. Consequently
\[
 v_6=b_0+b_1+b_2,
       \qquad v_2=b_3
       \tag{E.29.6}
\]
are the orbit sums for the six-point and two-point orbits
and form a basis of $V_{\mathbb Q_3}^{\mu_6}$.
In the three Tate coordinates of E.27, ordered as
$b_0,b_1+b_2,b_3$, the two Kummer scalars attached
to the first two coordinates are therefore equal:
the vector of scalars has the form $(\lambda,\lambda,\mu)$.
This equality concerns the actual cohomology class under
the bracket map $v\mapsto[v,w]$ and the canonical
projector of E.27, because those maps commute with
the based $\mu_6$ action constructed above.
Rational invariants of a finite group are exact,
so $0\to B^{\mu_6}\to A^{\mu_6}\to Q^{\mu_6}
\to0$ is exact. The positive weights give
$H^0(Q^{\mu_6}(-1))=0$, and the $H^2$ vanishing
of E.16 gives $H^2(B^{\mu_6}(-1))=0$.
The corresponding $H^1$ sequence is therefore
short exact.

The full geometric averaging projector $e_{\mathcal G}$
on $A$, for $\mathcal G=E[3]\rtimes\mu_6$,
preserves $A^{\mu_6}$, kills $B$, and identifies
$A^{\mathcal G}$ with $Q^{\mathcal G}$, as in
E.27. It also preserves the compact/boundary
multigrading. E.28 kills the pure boundary
cubic component, so the allowed image in
$H^1(Q^{\mu_6}(-1))$ is precisely the embedded
$H^1(H(1))$ summand. The projector gives its
canonical lift in $H^1(A^{\mathcal G}(-1))$;
the complementary kernel is exactly the image
of $H^1(B^{\mu_6}(-1))$. This proves the direct
sum (E.29.3). Finally
$H^1(G_{\mathbb Q,S_3},\mathbb Q_3(1))$
is the one-dimensional rational $S_3$-unit
Kummer group, generated by $\kappa_3(3)$.
The two kernel coordinates and the one-dimensional
elliptic summand therefore have the stated
dimensions and meanings.
\end{proof}

% ---- E_remaining_kernel ----
\begin{statement}{E.30}{Proved}{all three remaining rational coordinates lie in the ambient-metabelian kernel}
\label{st:E.30}
\label{E-remaining-kernel}
Retain the actual covering inclusion $N\hookrightarrow P$ of E.26,
with compatible tangential base points, and put
\[
 C=N/\gamma_4N,\qquad
 \bar C=N/(P''\gamma_4N),\qquad
 A=\gamma_3N/\gamma_4N.
\]
Let $W\subset V=N^{\rm ab}$ be the primitive rank-eight
cusp lattice. The kernel on the third lower central graded piece is
the actual Galois-stable submodule
\[
 K_3=\ker\bigl(A\longrightarrow\operatorname{gr}_3\bar C\bigr)
       =[W,\operatorname{Lie}_2(V)].
 \tag{E.30.1}
\]
It is saturated of rank $304$ in $A$, which has rank $330$.
The rational coefficient module carrying the three coordinates of
E.29 lies entirely in $K_{3,\mathbb Q_3}$. Consequently
\[
 [\eta]\longmapsto0\quad\text{in}\quad
 H^1\bigl(G_{\mathbb Q,S_3},(A/K_3)_{\mathbb Q_3}(-1)\bigr),
 \tag{E.30.2}
\]
and $[\eta]$ has a unique preimage in
$H^1(G_{\mathbb Q,S_3},K_{3,\mathbb Q_3}(-1))$.
Thus the specified once-punctured metabelian comparison annihilates
the entire rational class retained in E.29. This assertion evaluates
its image under that comparison. Propositions~\ref{E-ross-edge}
and~\ref{E-kummer-values} below identify its elliptic coordinate
with an exact cup-product edge and evaluate its two Kummer coordinates.
\end{statement}
\begin{proof}
We first verify the exact kernel calculation in the supplied
LM7 tower return, E.50, rather than treating its table as an input.
Choose two handle generators $e_1,e_2$ and eight cusp generators
$t_1,\ldots,t_8$ for $N$. Such a basis is obtained by cutting the
nine-sheeted torus cover along its two handle loops and its cusp
paths; the remaining ninth cusp is eliminated by the one surface
relation. The actual filling map $N\to T_3E$ kills the $t_i$.
The composite $N\to P^{\rm ab}$ is three times that map, since the
proper covering is $[3]:E\to E$. Multiplication by three is
injective on $T_3E$, so
\[
 P'=\ker(N\longrightarrow P^{\rm ab})
    =\overline{\langle\!\langle t_1,\ldots,t_8,[e_1,e_2]
                              \rangle\!\rangle}_N.
\]
This is also a direct presentation: quotienting by those generators
leaves the free abelian pro-three group on $e_1,e_2$, mapping
isomorphically to $3P^{\rm ab}$.

In $C$, write $J$ for the image of $P'$, $T$ for the subgroup
generated by the eight cusp generators, and $C_2=\gamma_2C$.
Then $J=T C_2$, the group $C_2$ is abelian, and $[T,C_2]$
is central. Commutator collection gives
\[
 J'=[T,T]\,[T,C_2].
\]
The degree-two images of $[t_i,t_j]$ are independent basis
vectors of $\bigwedge^2W\subset\bigwedge^2V$.
It follows that $[T,T]\cap\gamma_3C=\gamma_3T$;
this latter subgroup is contained in $[T,C_2]$.
Therefore
\[
 J'\cap\gamma_3C=[T,C_2]=[W,\operatorname{Lie}_2(V)].
\]
Since $J'$ is the image of $P''$, this proves (E.30.1)
as an equality of actual kernels, including at coefficient prime three.

Here is an integral basis check. Number the two handles $0,1$
and the cusp generators $2,\ldots,9$. Use the Hall basis
$[[v_i,v_j],v_k]$ with $i>j$ and $k\ge j$.
The subgroup in (E.30.1) has as its basis exactly the $304$
Hall vectors whose last index $k$ is a cusp index.
Each is directly the negative of $[v_k,[v_i,v_j]]$.
Conversely every such bracket reduces to these vectors:
when $k<j<i$, Jacobi gives
\[
 [[v_i,v_j],v_k]
   =[[v_i,v_k],v_j]-[[v_j,v_k],v_i],
\]
and both new last indices are cusp indices whenever $k$ is.
The $26$ retained vectors are the two pure handle triples
and the $24$ vectors
$[[t,e_1],e_1]$, $[[t,e_1],e_2]$, $[[t,e_2],e_2]$.
This proves saturation and the ranks over $\mathbb Z_3$.
The independent program \texttt{remaining\_kernel.py} checks
all $360$ generating brackets and all $304$ individual basis
witnesses over the integers; it uses no program absent from the
supplied return.

Pass now to the actual rational compact splitting
$V_{\mathbb Q_3}=H\oplus W_{\mathbb Q_3}$.
The kernel consists of all degree-three terms containing at
least two boundary letters, together with
$[W_{\mathbb Q_3},\bigwedge^2H]$.
Indeed Jacobi expresses every bracket with two boundary letters
in $[W,\operatorname{Lie}_2V]$, and the remaining one-boundary
generators are exactly the latter displayed subspace.

E.29 places the class in the image of cohomology of the
Galois-stable direct sum
\[
 B^{\mu_6}_{\mathbb Q_3}
       \ \oplus\
 (A_{\mathbb Q_3}^{E[3]\rtimes\mu_6})_{HW^2}
       \ \subset\ A_{\mathbb Q_3},
\]
with twist $(-1)$ understood. The first summand is
$[W^{\mu_6}_{\mathbb Q_3},w]$, where
$w\in\bigwedge^2H$; the second has two boundary letters.
Both lie in the kernel just computed. This proves (E.30.2).
The quotient $A/K_3$ has rational summands
$\operatorname{Lie}_3(H)$ and
$\operatorname{Sym}^2H\otimes W$, of Frobenius weights
three and four before the twist. After twist $(-1)$ their
weights are one and two, so a good Frobenius has no eigenvalue
one. Thus $H^0(G_{\mathbb Q,S_3},(A/K_3)_{\mathbb Q_3}(-1))=0$.
The long exact cohomology sequence makes the map from
$H^1(K_{3,\mathbb Q_3}(-1))$ injective and identifies its image
with the kernel in (E.30.2). This proves the asserted uniqueness.

This result specifies the additional kernel data that a
theta-function or elliptic path construction must retain to
evaluate E.29. It proves no vanishing before the quotient map.
The anchors are \S51.16, folios 733v--736v, scans 172--178,
and task LM7-E.23(i).
\end{proof}

% ---- E_full_metabelian ----
\begin{statement}{E.31}{Proved}{noncuspidal sections of the full profinite metabelian cubic quotient}
\label{st:E.31}
\label{E-full-metabelian-family}
Let $U=E\setminus E[3]$ be the plus Fermat cubic with origin
$O$, over $\mathbb Q$, and choose the rational tangential section
$s_O$ used above. In this statement $N_{\mathrm{all}}$ denotes
the full geometric profinite group, and
\[
 M=N_{\mathrm{all}}/N_{\mathrm{all}}'',\qquad
 A=N_{\mathrm{all}}^{\mathrm{ab}},\qquad
 D=N_{\mathrm{all}}'/N_{\mathrm{all}}''.
\]
There is an injective continuous family of arithmetic section classes
\[
 \mathbb Z_3\longrightarrow
 \operatorname{Sec}(M\rtimes_{s_O}G_{\mathbb Q})/M,
       \qquad a\longmapsto[s_a],
 \tag{E.31.1}
\]
whose abelianized section is $[s_O]$ for every $a$.
The class $[s_0]$ is $[s_O]$. Every $a\ne0$ gives a
noncuspidal section of this full profinite metabelian quotient.
Its images in every finite quotient, including quotients with
several coefficient primes, are the images of this one actual
continuous section.

The first nonzero family member is detected in the intrinsic
class-two three-primary quotient with coefficients modulo $27$.
The construction does not give a lift to the full geometric
fundamental group. In particular it supplies actual full
metabelian classes for task LM7-E.23(iii), while the complete
section image and its other fibres remain as stated in E.21.
\end{statement}
\begin{proof}
We first construct the comparison that makes the full profinite
claim possible. Let $P=N_{\mathrm{all}}^{(3)}$ be the maximal
pro-three quotient and put
\[
 D_3=P'/P'',\qquad A_3=P^{\mathrm{ab}},\qquad
 R=\widehat{\mathbb Z}[[A]],\qquad R_3=\mathbb Z_3[[A_3]].
\]
The group $N_{\mathrm{all}}$ is free profinite of rank ten,
and $P$ is free pro-three of that rank. A chosen free basis
gives the continuous Alexander boundary sequences
\[
 \begin{aligned}
 0\longrightarrow D&\longrightarrow R^{10}
       \xrightarrow{\partial} I_A\longrightarrow0,\\
 0\longrightarrow D_3&\longrightarrow R_3^{10}
       \xrightarrow{\partial_3} I_{A_3}\longrightarrow0,
 \end{aligned}
 \qquad
 \partial(r_i)=\sum_i r_i(X_i-1).
 \tag{E.31.2}
\]
Here the ideals are the kernels of augmentation. One direct
verification uses the Cayley graphs of the finite abelian
covering quotients: their edge and vertex boundaries are the
displayed rows, and their first homology is the abelianization
of the covering group. Passing to all finite coefficients and
all such covers gives (E.31.2). To justify the limit on first
homology, a continuous finite character of $N_{\mathrm{all}}'$
has kernel containing $N_{\mathrm{all}}'\cap V$ for some open
normal subgroup $V$ of $N_{\mathrm{all}}$. It extends to
$N_{\mathrm{all}}'V$ by being trivial on $V$. This is
well-defined, and it is a homomorphism because
$[V,N_{\mathrm{all}}']$ is in that intersection. Thus these
abelian covers detect every finite character of the closed
commutator group. The same argument over pro-three quotients
gives the second sequence and their natural comparison.

In the coefficient-three factor $\mathbb Z_3[[A]]$, average
over the characteristic subgroup $A^{(3')}$, the product of
the other pro-prime factors. At each finite quotient its
order is prime to three, so its normalized Haar average
is a compatible idempotent. Extend that idempotent by zero
in the other coefficient factors of $R$ and call it $e$.
Then
\[
 eR\xrightarrow{\sim}R_3,
 \qquad eI_A\xrightarrow{\sim}I_{A_3},
 \qquad e(X_i-1)\longmapsto X_{i,3}-1.
 \tag{E.31.3}
\]
Multiplication by a central idempotent is exact.
Comparing the kernels of the two rows in (E.31.2)
therefore shows that the natural map $q:D\to D_3$
restricts to a topological isomorphism $eD\to D_3$.
Its inverse followed by inclusion is a specified continuous
section
\[
 j:D_3\lhook\joinrel\longrightarrow D,
                  \qquad qj=1.
 \tag{E.31.4}
\]
It is independent of the free basis: the map $q$ is natural,
and $e$ is preserved by every automorphism of $A$.
For the actual arithmetic action these same facts give
$g(e d)=e\,g(d)$ and $qg=gq$, hence $jg=gj$.
This is an equivariant embedding of the derived modules;
no section $P\to N_{\mathrm{all}}$ is asserted.

We now specify the cocycle. Use the actual arithmetic
quotient from E.16 and E.18,
\[
 \Gamma=P_{\mathrm{ar}}\rtimes\langle c\rangle,
 \qquad
 P_{\mathrm{ar}}=G_{\mathbb Q(\zeta_3),\{3\}}^{(3)}
                  =\langle\tau,\tau'\rangle_{\mathrm{pro}-3},
 \qquad c\tau c=\tau'.
 \tag{E.31.5}
\]
The action on the entire geometric pro-three group, hence
on $D_3$, factors through this quotient. The proof in E.16
checks the three-group residual action and the free arithmetic
source; it is not a replacement of an arbitrary arithmetic
action by its pro-three quotient.

Choose positive compatible roots for $\kappa=\kappa_3(3)$.
The Frobenius basis at seven can be chosen with
\[
                    \kappa(c)=0,\qquad
                    \kappa(\tau)=1,\qquad
                    \kappa(\tau')=-1.
 \tag{E.31.6}
\]
Here is the exact normalization. Start with the conjugate
Frobenii $\tau_0,\tau'_0$ at the primes above seven and put
$b=\kappa(\tau_0)$. The cubic residue computation
$3^2=2$ modulo seven, with $\zeta_3$ reducing to $2$,
gives $b\equiv1\pmod3$. Set
$r=\tau'_0\tau_0^{-1}$ and
$h=r^{(b^{-1}-1)/12}$. The exponent is in $\mathbb Z_3$.
One has $\chi(r)=1$, $\kappa(r)=-2b$, and
\[
 \kappa(h\tau_0h^{-1})
    =b+12b\,(b^{-1}-1)/12=1.
\]
Take $\tau=h\tau_0h^{-1}$ and $\tau'=c\tau c$.
Their Frattini classes are unchanged, so they remain a
free basis. This also checks the normalization used in
the supplied LM7 tower return, E.49, by its actual map.
The image on $V_3$ over $\mathbb Q(\zeta_3)$ is abelian:
after the explicit rational splitting of E.18 it is the
sum of the compact complex-multiplication character and
the boundary cyclotomic character. Equality on the rational
representation implies equality on its stable integral
lattice. Thus conjugating $\tau_0$ by $h$ leaves its
displayed matrix on $V_3$ unchanged; the finite certificate
below uses this same normalized generator.

Let $u,v\in V_3=P^{\mathrm{ab}}$ be the two compact
lift vectors of E.18, so $w=u\wedge v$ is primitive,
arithmetic $g$ acts on $w$ by $\chi_3(g)$, and
$c(w)=-w$. Choose free geometric generators lifting
the stated basis of $V_3$. Ordered products with the
integer exponent vectors $u,v$ give elements
$\widetilde u,\widetilde v\in P$. Put
\[
                  d=[\widetilde u,\widetilde v]\in D_3.
\]
Its image under $D_3\twoheadrightarrow
\gamma_2P/\gamma_3P=\bigwedge^2V_3$ is exactly $w$.
The free pro-three property in (E.31.5) supplies a unique
continuous crossed homomorphism on $P_{\mathrm{ar}}$
with values
\[
                       b(\tau)=d,\qquad b(\tau')=c(d).
\]
Indeed these values give a homomorphism into the pro-three
semidirect product $D_3\rtimes P_{\mathrm{ar}}$ on the two
free generators. Taking $b(c)=0$ extends it to $\Gamma$,
because the two assignments respect $c\tau c=\tau'$ and
$c^2=1$. Its degree-two projection is precisely
$\kappa w$, by (E.31.6) and equality on these generators.
Inflate it to $G_{\mathbb Q}$.

For $a\in\mathbb Z_3$ define the actual full metabelian
section by
\[
                     s_a(g)=j(a\,b(g))s_O(g).
 \tag{E.31.7}
\]
The image of $j$ is in the abelian normal subgroup $D$.
Its equivariance and the cocycle identity prove directly
that (E.31.7) is a continuous homomorphism and a splitting.
It has zero abelian cocycle. Its image in the three-primary
class-two fibre is $a[\kappa w]$.

For completeness this latter class is nonzero with an exact
finite certificate. In the ordered wedge basis of E.18,
write $T=\bigwedge^2 A_\tau$ and $J=\bigwedge^2 C_V$.
After the normalization $b(c)=0$, coboundaries are exactly
the image of $(T-I)(I+J)$. The program
\texttt{full\_metabelian\_family.py} produces an integral
row $\lambda$ with
\[
 \lambda(T-I)(I+J)=0,
                   \qquad \lambda w=1872\equiv9\pmod{27}.
 \tag{E.31.8}
\]
The row, both $45$-dimensional matrices and the wedge
vector are included in its output. Thus the first member
is already nonzero with these finite coefficients.
For any $a\ne0$, the same row modulo
$3^{v_3(a)+3}$ detects $a[\kappa w]$.
Equivalently, the rational argument uses
$(T-I)^{-1}w=w/6$ and $Jw=-w$; this vector cannot be
a $J$-fixed coboundary parameter.

There is no further identification by conjugation in $M$.
A conjugator taking two zero-abelian sections to one another
has abelian coordinate in $H^0(G_{\mathbb Q},A)$, which is zero.
For each coefficient prime this follows from the compact
weight-one and boundary weight-two eigenvalues at a good
prime distinct from it; no eigenvalue is one. Thus the
conjugator belongs to $D$, and it changes $jb$ only by a
coboundary. Projecting with $q$ and then to degree two
proves injectivity in (E.31.1). Continuity follows from
continuous scalar multiplication and evaluation of (E.31.7);
the compact parameter space and the closed profinite
conjugacy relation also give the asserted section-class topology.

Finally consider a cusp packet in this quotient. A rational
cusp $x\ne O$ has nonzero compact Kummer class of $x-O$:
it is a nonzero three-torsion point of $E(\mathbb Q)$,
and the Mordell--Weil Kummer map into $H^1(\mathbb Q,TE)$
is injective. Its packet therefore contains no zero-abelian
class. A nonrational cusp has no packet over $\mathbb Q$.
At $O$, the cusp line $I_O^{\mathrm{ab}}\simeq
\widehat{\mathbb Z}(1)$ is primitive in $A$, and
$H^0(G_{\mathbb Q},A/I_O^{\mathrm{ab}})=0$ by the same
weight calculation. The exact cohomology sequence makes
the packet map injective on abelianization. Its only
zero-abelian class is therefore the parameter-zero
section $s_O$. Equation (E.31.8) excludes that class
for every $a\ne0$. This proves noncuspidality in the
stated quotient and all assertions of the theorem.
The primitive-loop centralizer calculation in E.33
below verifies directly that the cusp normalizers
of this metabelian quotient have precisely these
packets; in particular no extra normalizer section
is omitted by the comparison with the full group.
The anchors are \S51.16, folios 733v--736v, scans172--178,
and task LM7-E.23(iii).
\end{proof}

% ---- E_pronilpotent_families ----
\begin{statement}{E.32}{Proved}{full pronilpotent sections at the three- and five-power levels}
\label{st:E.32}
\label{E-full-pronilpotent}
Let $n=\ell^r$ with $\ell\in\{3,5\}$ and $r\ge1$, and let
$N=\pi_1(U_{n,\overline{\mathbb Q}})$ be the full profinite
geometric group of the plus Fermat curve. Use the rational
unit tangent at $O=[0:-1:1]$ to specify the arithmetic action.
Its maximal pronilpotent quotient has the canonical description
\[
 N^{\mathrm{nil}}=\prod_q N^{(q)},
 \qquad
 N/\gamma_{c+1}N\simeq
       \prod_q N^{(q)}/\gamma_{c+1}N^{(q)}\quad(c\ge1).
 \tag{E.32.1}
\]
Here $q$ ranges over all primes and lower central subgroups
are closed. At each quotient, the cuspidal classes in this
statement are the images of the actual cusp packets under
$D_c\hookrightarrow\pi_1(U_n)\twoheadrightarrow
E^{(c)}$. This specifies the packet map even at class one,
where the normalizer of a boundary line can contain further
sections. The arithmetic extension with geometric kernel
$N^{\mathrm{nil}}$ has infinitely many continuous section
classes whose image in every displayed positive-class
nilpotent quotient is noncuspidal. Consequently no level of
this full profinite nilpotent tower selects precisely the
cuspidal section classes at these Fermat levels.
\end{statement}
\begin{proof}
Every finite nilpotent group is the direct product of its
Sylow subgroups. A map from $N$ to such a group factors,
in each Sylow coordinate, through $N^{(q)}$. Conversely,
the natural map from $N$ to a finite product of its finite
prime-power quotients is onto: its image is a subdirect
product, and factors of coprime orders have no common
nontrivial quotient. Compactness now proves that
$N\to\prod_qN^{(q)}$ is onto. This proves the first
identity in (E.32.1), with its actual quotient map.

The quotient $N/\gamma_{c+1}N$ is the inverse limit of
its finite nilpotent quotients of class at most $c$.
Apply the preceding factorization to each such quotient.
The reverse comparison is induced by the coordinate
maps to $N^{(q)}/\gamma_{c+1}N^{(q)}$, whose product has
class at most $c$. The universal properties give inverse
continuous maps, proving the second identity. Both
comparisons commute with every automorphism of $N$,
and in particular with the specified Galois action.

Write a section supplied by E.16 as
$s_\ell(g)=b_\ell(g)s_O(g)$, where
$b_\ell:G_{\mathbb Q}\to N^{(\ell)}$ is its actual
continuous nonabelian cocycle. Define a cocycle in the
first product of (E.32.1) by taking coordinate
$b_\ell$ at $\ell$ and the identity at every other prime.
The product topology and the coordinatewise arithmetic
action make it a continuous cocycle. It therefore gives
one actual section of the pronilpotent arithmetic
extension, whose images are compatible in every
finite nilpotent quotient involving any finite set
of coefficient primes.

Choose the E.16 sections with distinct infinite-order
images in
$H^1(\mathbb Q,T_\ell\operatorname{Jac}(X_n))$.
That compactified abelian image factors through every
quotient in (E.32.1). Every cusp image is $n$-torsion:
the ratio of two cusp tangent-line equations has
divisor $n(d-O)$, as checked in E.16. Thus none of the
chosen images can belong to a cusp packet at any
positive class. Distinct compact images also prevent
conjugacy of the corresponding sections. This proves
the assertions, including the full coefficient-prime
comparison. It does not compute section classes of
finite quotients outside the specified nilpotent
tower. The anchors are \S51.16, folios733v--736v,
scans172--178, and LM7-E.23(iii).
\end{proof}

\begin{statement}{E.33}{Proved}{full metabelian zero-abelian families at every three-power level}
\label{st:E.33}
\label{E-full-metabelian-tower}
Let $n=3^r$ with $r\ge1$ and keep the full profinite
group $N$ and tangent $s_O$ of E.32. Set
\[
 M=N/N'',\qquad A=N^{\mathrm{ab}},\qquad
 d_n=\frac{(n^2-1)(n^2+3)}4.
\]
There is a continuous injective family
\[
 \mathbb Z_3^{d_n}\longrightarrow
 \operatorname{Sec}(M\rtimes_{s_O}G_{\mathbb Q})/M
 \tag{E.33.1}
\]
all of whose abelianized classes equal $[s_O]$.
The zero parameter maps to $[s_O]$; every nonzero
parameter is noncuspidal. In particular $d_3=24$.
The family is specified by a choice of a free lattice
in the class-two cohomology quotient below and lifts
of its generators in the pro-three derived module.
E.31 supplies a separate explicit finite certificate
for one cubic direction.
\end{statement}
\begin{proof}
The geometric rank is $n^2+1$. Put
\[
 P=N^{(3)},\quad V=P^{\mathrm{ab}},\quad
 D_3=P'/P'',\quad B_2=\bigwedge^2_{\mathbb Z_3}V.
\]
The continuous surjection $D_3\twoheadrightarrow B_2$
is the degree-two commutator map, with kernel
$\gamma_3P/P''$. The Fox and Haar comparison of E.31
works in every finite free rank, by the identical
row $\sum r_i(X_i-1)$. It gives the actual
arithmetic-equivariant embedding
\[
 j:D_3\hookrightarrow N'/N'',
 \qquad qj=1,
 \tag{E.33.2}
\]
where $q$ is the natural pro-three comparison.
No change in the arithmetic action is made.

By E.16 that action on $P$ factors through
\[
 \Gamma=\langle\tau,\tau'\rangle_{\mathrm{pro}-3}
                 \rtimes\langle c\rangle,
 \qquad c^2=1,\quad c\tau c=\tau'.
 \tag{E.33.3}
\]
Use the conjugate Frobenius basis at seven; let
$T$ be the action of $\tau$ on $B_2$. Since two is
invertible, every $B_2$-cocycle class has a
representative with $b(c)=0$. Such a representative
is specified by its arbitrary value at $\tau$;
the value at $\tau'$ must be $c b(\tau)$. Freeness
extends these assignments continuously, and they
respect the two relations involving $c$. A
coboundary preserves $b(c)=0$ precisely when its
parameter belongs to $B_2^c$. Consequently
\[
 H^1_{\mathrm{cts}}(\Gamma,B_2)
       =B_2/(T-1)B_2^c.
 \tag{E.33.4}
\]
This is a quotient by the literal compact image.

The eigenvalues of good Frobenius on the compact
part of $V$ have weight one and those on the
boundary part have weight two. All eigenvalues
on $B_2$ therefore have weights two, three or
four, so none equals one. Thus $T-1$ is invertible
on $B_2\otimes\mathbb Q_3$, and (E.33.4) has
free rank $\dim(B_2\otimes\mathbb Q_3)^-$.

We compute that dimension with the actual real
action. The compact curve has genus
$g=(n-1)(n-2)/2$; complex conjugation has
$g$ eigenvalues of each sign on its first
homology. This follows from its reversal of
the nondegenerate intersection pairing: the
two eigenspaces are isotropic, and the pairing
identifies either with the dual of the other.
Of the $3n$ boundary points, exactly three
are real, and the rest form $(3n-3)/2$
conjugate pairs. The boundary module is the
quotient of their permutation module tensored
with $\mathbb Z_3(1)$ by the diagonal line
generated by the sum of all cusp generators.
Complex conjugation acts by minus one on that
line; thus taking this quotient removes one
minus eigenvalue.
Its plus and minus dimensions are respectively
$(3n-3)/2$ and $(3n+1)/2$. Hence
\[
 \dim V^+=\frac{n^2-1}2,
 \qquad
 \dim V^-=\frac{n^2+3}2,
 \qquad
 \dim B_2^-=\dim V^+\dim V^-=d_n.
 \tag{E.33.5}
\]
The last equality uses the actual direct sum
$\bigwedge^2(V^+\oplus V^-)$: its minus part is
$V^+\otimes V^-$. Thus no trace estimate replaces
the eigenspace calculation.
The program \texttt{metabelian\_family\_rank.py}
constructs the cubic boundary quotient matrix
and its full exterior-square action independently;
its exact nullspace ranks are $4$, $6$ and $24$.

Choose $d_n$ elements of $B_2$ whose classes
give a free lattice in (E.33.4), and choose
their lifts $d_i\in D_3$. For
$a=(a_1,\ldots,a_{d_n})\in\mathbb Z_3^{d_n}$
assign
\[
 b_a(\tau)=\sum_i a_i d_i,\qquad
 b_a(\tau')=c b_a(\tau),\qquad b_a(c)=0.
 \tag{E.33.6}
\]
The same free-source argument gives an actual
continuous $D_3$-cocycle. Its dependence on $a$
is continuous: equality on the free generators
gives $b_a=\sum_i a_i b_i$ for the fixed cocycles
$b_i$. This is jointly continuous in $(a,g)$,
and therefore continuous in the compact-open
topology. Inflate to $G_{\mathbb Q}$
and set $s_a(g)=j(b_a(g))s_O(g)$. This proves
existence of (E.33.1) with zero abelian cocycle.

The invariant group $H^0(G_{\mathbb Q},A)$ is
zero, coefficient prime by coefficient prime,
using a good Frobenius of different residue
characteristic and the same compact and
boundary weights. Thus a geometric conjugator
between two such sections lies in $N'/N''$.
Projecting with $q$ and then to $B_2$ proves
that their difference is a coboundary in
(E.33.4). Inflation from $\Gamma$ creates no
new coboundaries: the actual action factors
through $\Gamma$, so the same parameter
defines a coboundary before and after
inflation. The chosen free lattice therefore
proves injectivity. Compactness and the
closed conjugacy relation give the stated
topology on this family of section classes.

We check cusp normalizers and packets in the
quotient itself. A cusp generator can be
included in a free basis of $N$, by omitting
a different cusp in the surface presentation.
For such a generator $x$, class two shows
that the abelian coordinate of any element
centralizing $x$ is a multiple of $x$.
After subtracting that power, the element
lies in $D=N'/N''$ and is killed by $X-1$.
The Fox embedding $D\hookrightarrow R^{n^2+1}$
and injectivity of $X-1$ on $R$ force it
to vanish. For the last injectivity, in each
coefficient-$q$ factor refine a finite
abelian quotient by a further $q^m$ in the
primitive $X$ direction. An $X$-invariant
coefficient pushes forward as a multiple
of $q^m$, and is zero modulo every fixed
$q$-power once $m$ is large. The inverse
limit is therefore zero. Consequently
$C_M(x)=\langle x\rangle$. A normalizer
cannot change the exponent of $x$, since
conjugation fixes its primitive abelian
vector. Thus its geometric normalizer is
the same cusp group. The arithmetic
normalizer over a rational cusp has
exactly the usual Kummer packet. Over a
nonrational cusp its image is the
stabilizer of that cusp: distinct boundary
lines remain distinct in $A$ and the
Galois action permutes them. Such a
normalizer cannot contain a $G_{\mathbb Q}$
section.

For the three rational cusps, the
tangent-line ratio has divisor $n(d-O)$.
The Abel--Jacobi image of $d-O$ is nonzero
for $d\ne O$, by injectivity of the curve's
Abel map in genus at least one. It is
therefore nonzero three-primary torsion.
Its three-primary Kummer class is nonzero
by the Mordell--Weil Kummer injection.
No packet at $d\ne O$ has zero abelian
class. At $O$ the primitive cusp line $I_O$
and $H^0(G_{\mathbb Q},A/I_O)=0$ make the
packet map injective on abelianization.
Its zero class is precisely $s_O$.
Injectivity of (E.33.1) excludes this
class for every nonzero parameter.
The anchors are \S51.16, folios733v--736v,
scans172--178, and LM7-E.23(iii).
\end{proof}

% ---- E_five_metabelian ----
\begin{statement}{E.34}{Proved}{full metabelian zero-abelian families at every five-power level}
\label{st:E.34}
\label{E-five-metabelian}
Let $n=5^r$, $r\ge1$, and let
$N=\pi_1(U_{n,\overline{\mathbb Q}})$ be the full
profinite geometric group with the rational unit
tangent $s_O$ at $O=[0:-1:1]$. Put
\[
 M=N/N'',\qquad d_n=\frac{(n^2-1)(n^2+3)}4.
\]
There is a continuous injective family
\[
 \mathbb Z_5^{d_n}\longrightarrow
 \operatorname{Sec}(M\rtimes_{s_O}G_{\mathbb Q})/M
 \tag{E.34.1}
\]
with zero abelian cocycle. Its zero member is
$[s_O]$, and every nonzero member is noncuspidal,
including for the cusp normalizers in this
metabelian quotient itself. In particular
$d_5=168$.
\end{statement}
\begin{proof}
Put $P=N^{(5)}$, $V=P^{\mathrm{ab}}$,
$D_5=P'/P''$, and $B_2=\bigwedge^2V$.
The same characteristic Haar summand as in
E.31, now in the coefficient-five factor,
gives a continuous arithmetic-equivariant
embedding
$j:D_5\hookrightarrow N'/N''$ which is a
section of the natural comparison. The
surjection $D_5\twoheadrightarrow B_2$ is
the actual degree-two quotient.

E.16 identifies the actual arithmetic action
with a quotient
\[
 \Gamma=P_{\mathrm{ar}}\rtimes C_4,
 \qquad
 P_{\mathrm{ar}}=G_{\mathbb Q(\zeta_5),\{5\}}^{(5)},
 \tag{E.34.2}
\]
where $P_{\mathrm{ar}}$ is free pro-five of
rank three. The complement is chosen by
the profinite coprime splitting theorem;
changing it gives conjugate data. No lift
of that finite complement to $G_{\mathbb Q}$
is required.

We prove surjectivity on cohomology for the
precise comparison
\[
 H^1_{\mathrm{cts}}(\Gamma,D_5)
       \longrightarrow H^1_{\mathrm{cts}}(\Gamma,B_2).
 \tag{E.34.3}
\]
Let $b$ represent a class on the right.
Since four is invertible on these modules,
averaging on $C_4$ changes $b$ by a
coboundary so that $b|_{C_4}=0$.
Its restriction to $P_{\mathrm{ar}}$ is
then exactly invariant for the action
\[
 (a\cdot b)(g)=a\,b(a^{-1}ga),\qquad a\in C_4.
 \tag{E.34.4}
\]
Lift its values on the three free
generators arbitrarily to $D_5$.
Freeness gives a continuous $D_5$-cocycle
$\widetilde b$ on $P_{\mathrm{ar}}$.
The sum
\[
 \widetilde b_{\mathrm{av}}
       =\frac14\sum_{a\in C_4}a\cdot\widetilde b
 \tag{E.34.5}
\]
is meaningful in this abelian pro-five
module. It is a continuous invariant
cocycle, and its projection is exactly
$b|_{P_{\mathrm{ar}}}$ because that
cocycle was already invariant. Setting
its value on $C_4$ equal to zero extends
it to $\Gamma$. This proves (E.34.3)
on actual cocycles, without an inverse
limit of unrelated lift choices.

We compute the rank of the target using
the same arithmetic cohomology inputs
proved in E.16. Explicitly, a $B_2$-cocycle
on $G_{\mathbb Q(\zeta_5),\{5\}}$ has
graph in the pro-five group
$B_2\rtimes\rho(G_{\mathbb Q(\zeta_5),\{5\}})$.
It kills the kernel of the maximal
pro-five quotient. Consequently inflation
identifies, rather than merely embeds,
\[
 H^1(\Gamma,B_2)
       =H^1(G_{\mathbb Q,\{5,\infty\}},B_2).
 \tag{E.34.6}
\]
The finite-module filtration, the
coprime averaging and the finite-level
Mittag--Leffler argument in E.16 give
zero continuous $H^2$ here. Good-Frobenius
weights give $H^0=0$. The global Euler
formula, in its real-place form used
there, therefore gives
\[
 \operatorname{rank}_{\mathbb Z_5}H^1(\Gamma,B_2)
       =\dim_{\mathbb Q_5}(B_2\otimes\mathbb Q_5)^-.
 \tag{E.34.7}
\]
The cusp permutation and compact
intersection-pairing calculation of
(E.33.5) applies to every odd $n$.
Thus the last dimension is $d_n$.

Choose a free lattice of rank $d_n$
in this cohomology group, choose
representative cocycles, and lift
them by (E.34.3). A linear combination
with coefficients in $\mathbb Z_5$
gives a jointly continuous family of
cocycles in $D_5$. Apply $j$ and
multiply by $s_O$. The conjugacy proof
of E.33 applies with five in place
of three: $H^0(G_{\mathbb Q},N^{\rm ab})=0$
forces every conjugator into $N'/N''$,
and projection to $B_2$ then proves
injectivity of the parameter family.

The primitive-loop centralizer proof
in E.33 uses the full completed
group ring and is independent of
the coefficient prime, so it
identifies the cusp normalizers
here as well. Each nonorigin
rational cusp has a nonzero
Abel--Jacobi class killed by
$n=5^r$, by the tangent-line
divisor. Its five-primary Kummer
class is nonzero. Its packet
therefore has no zero-abelian
section. The origin packet has
only $s_O$ at zero abelianization,
and the nonrational cusps have
proper arithmetic stabilizers.
This proves every assertion.
The anchors are \S51.16,
folios733v--736v, scans172--178,
and LM7-E.23(iii).
\end{proof}

% ---- E_ross_edge ----
\subsection{The elliptic action coordinate as an explicit cup-product edge}

\begin{statement}{E.35}{Proved}{the remaining elliptic coordinate is the Ross cup-product edge}
\label{st:E.35}
\label{E-ross-edge}
Retain the actual rational cubic class $[\eta]$ of
E.19, its compact splitting $V_{\mathbb Q_3}=H\oplus W$,
and the primitive determinant vector $w=u\wedge v$.
Put $x=-X/Z$, $y=-Y/Z$, so that $x^3+y^3=1$, and
use the unit tangent for $q=Z/X$ at $O=[1:-1:0]$.
The two functions
\[
 f=1-x,\qquad h=y-1
 \tag{E.35.1}
\]
are units on $U$. Their geometric Kummer maps
$\ell_f,\ell_h:V_{\mathbb Q_3}\to\mathbb Q_3(1)$
vanish on $H$. Let
\[
 b=\ell_f\wedge\ell_h:
       \bigwedge^2W\longrightarrow\mathbb Q_3(2),
 \qquad
 J:H\otimes\bigwedge^2W\longrightarrow\operatorname{Lie}_3(V),
 \quad J(a\otimes(r\wedge s))=[a,[r,s]].
 \tag{E.35.2}
\]
The map $J$ is injective. The two-dimensional invariant
elliptic summand $D_{HW^2}$ of E.28 is contained in its
image, and contraction defines an actual
$G_{\mathbb Q}$-equivariant isomorphism
\[
 \Phi=(1_H\otimes b)J^{-1}:
       D_{HW^2}\xrightarrow{\sim}H(2).
 \tag{E.35.3}
\]
This fixes a normalization of the remaining elliptic
coordinate, without choosing a scalar generator of its
one-dimensional arithmetic cohomology group.

Here is its value as a specified cohomology class.
Write $\Gamma=P\rtimes C_2$ for the actual free
pro-three arithmetic quotient of E.16, and
$\mathcal E=N\rtimes\Gamma$ for the geometrically
pro-three extension with its chosen tangent.
For the Hochschild--Serre edge map with the convention
specified in the proof, put
\[
 \begin{aligned}
 \rho&=\kappa(f)\smile\kappa(h)
              \in H^2_{\rm cts}(\mathcal E,\mathbb Q_3(2)),\\
 \operatorname{edge}_H(\rho)&\in
 H^1_{\rm cts}\bigl(\Gamma,
                   \operatorname{Hom}(H,\mathbb Q_3(2))\bigr),\\
 \iota_w(\lambda)&=u\lambda(v)-v\lambda(u).
 \end{aligned}
 \tag{E.35.4}
\]
In the last line the determinant twist gives the
equivariant map
$\iota_w:\operatorname{Hom}(H,\mathbb Q_3(2))\to H(1)$.
In (E.35.5), $\Phi_*$ is induced by the explicitly
twisted map $\Phi(-1):D_{HW^2}(-1)\to H(1)$.
If $[\eta]_{\rm ell}$ denotes the invariant elliptic
component in E.29, then
\[
 \boxed{\quad
  \Phi_*[\eta]_{\rm ell}
      =-\iota_{w*}\operatorname{edge}_H(\rho)
       \quad}\qquad\text{in }H^1(\Gamma,H(1)).
 \tag{E.35.5}
\]
Inflation gives the same equality for the original
$G_{\mathbb Q,\{3,\infty\}}$ action. Since
$\kappa_3(-1)=0$, the cup product in (E.35.4) is
also $\kappa(1-x)\smile\kappa(1-y)$, the Kummer
cup attached to the Ross symbol.
Formula (E.35.5) computes an exact comparison class.
Its numerical three-adic value, its vanishing or
nonvanishing, and the integral torsion coordinates
remain unevaluated. Proposition~\ref{E-kummer-values}
below evaluates the two Kummer coordinates of E.29.
\end{statement}
\begin{proof}
First we fix every comparison needed by the cup
calculation. Let $\zeta=\zeta_3$ and label the cusps
\[
 A_i=[0:-\zeta^i:1],\quad
 B_i=[-\zeta^i:0:1],\quad
 C_i=[-\zeta^i:1:0]\quad(0\le i<3).
\]
Here $O=C_0$. The tangent-line equations give
\[
 \operatorname{div}(f)=3B_0-C_0-C_1-C_2,
 \qquad
 \operatorname{div}(h)=3A_0-C_0-C_1-C_2.
 \tag{E.35.6}
\]
For example, $X+Z=0$ meets the cubic with
intersection multiplicity three at $B_0$, whereas
$Z=0$ meets its three cusps simply. This proves
the first formula; the second follows from
$Y+Z=0$. The corresponding Kummer character on
a cusp meridian is its displayed order of
vanishing. The boundary module is
$\mathbb Q_3(1)^{\{A_i,B_i,C_i\}}$ modulo the
diagonal sum. Both order vectors sum to zero,
so they define precisely the maps $\ell_f,\ell_h$
on that quotient. On the compact summand the
maps are zero: they are Galois-equivariant, while
at good Frobenius the eigenvalues on $H$ have
weight one and that on $\mathbb Q_3(1)$ has
weight two. Thus the actual linear maps vanish
on $H$, not only on a chosen coordinate subspace.

At $O$ the branch with $Y/X=-1$ has
$Y/X=-(1+q^3)^{1/3}$. Hence
\[
 f=q^{-1}(1+q),\qquad
 h=q^{-1}\bigl((1+q^3)^{1/3}-q\bigr).
 \tag{E.35.7}
\]
Both leading coefficients are one. Their
Kummer cocycles therefore vanish on the specified
tangential arithmetic section. They factor through
$\mathcal E$: on the geometric group their
prime-three Kummer characters factor through $N$,
and their values on the section are zero. This
also specifies their values on the entire
semidirect product. A different rational tangent
changes them by arithmetic constant Kummer
classes; the compact edge of these added cross
terms is zero because both geometric characters
vanish on $H$. The fixed tangent in the statement
is the one used in the following literal formulas.

Choose lifts of a basis of $V$ in the rational
class-two Malcev Lie algebra. In its exponential
coordinates a geometric element is $(v,z)$,
with $z\in\bigwedge^2V$, and multiplication is
\[
 (v,z)(v',z')=
       (v+v',z+z'+\tfrac12v\wedge v').
 \tag{E.35.8}
\]
The genuine arithmetic action has the form
\[
 g(v,z)=\bigl(A_gv,
          (\bigwedge^2A_g)z+B_g(v)\bigr),
 \qquad C_g=B_gA_g^{-1}.
 \tag{E.35.9}
\]
Composition of these actual actions gives
$C_{gg'}=C_g+g\cdot C_{g'}$, with the usual
action on $\operatorname{Hom}(V,\bigwedge^2V)$.
Thus $C$ is a continuous one-cocycle. Changing
the graded coordinate lifts changes its class by
a coboundary in this same module.

The class-three bracket calculation of E.28 now
gives a more precise formula. Divide
$g[u,v]-\chi(g)[u,v]$ by $\chi(g)$. Its
$HW^2$ component is
\[
 \eta_{HW^2}(g)=
 J\bigl(u\otimes C_g^{WW}(v)
            -v\otimes C_g^{WW}(u)\bigr),
 \tag{E.35.10}
\]
where $C_g^{WW}$ is the restriction to $H$
followed by projection to $\bigwedge^2W$.
Indeed the determinant of $A_g|H$ is $\chi(g)$;
expanding the two cross brackets in (E.28.3)
cancels precisely that determinant. An inner
coordinate change contributes $a\wedge H$
to the degree-two correction, whose $WW$
projection is zero. Formula (E.35.10) therefore
also respects the specified quotient by inner
path ambiguity. The map $J$ is injective because,
in the tensor-algebra realization of the free
Lie algebra, the words whose first letter is
in $H$ recover the independent tensor
$a\otimes(r\wedge s)$ from $[a,[r,s]]$.

We compute the edge including its sign. Extend
$b$ by zero on the other multidegrees of
$\bigwedge^2V$. Write $e=(n,g)$, $e'=(m,g')$
and let $(v_n,z_n),(v_m,z_m)$ be their Malcev
coordinates. The antisymmetrized cup cocycle is
\[
 \omega(e,e')=
       \tfrac12 b(v_n\wedge A_gv_m).
\]
It represents $\rho$: graded commutativity
identifies the original cup with its
antisymmetrization, or the symmetric part is
the coboundary of the negative one-half product
of the two Kummer one-cochains. For the
continuous one-cochain $c(e)=-b(z_n)$, the
inhomogeneous differential and (E.35.8)--(E.35.9)
give exactly
\[
 \delta c(e,e')
    =\omega(e,e')+bB_g(v_m),
 \qquad
 \omega-\delta c=-bB_g(v_m)
                         =-(bC_g)(A_gv_m).
 \tag{E.35.11}
\]
The Tate action is included: $g\,b(z)=
b((\bigwedge^2A_g)z)$. The expression in the
last line is the mixed two-cocycle associated
to the left-module one-cocycle $-bC_g$.
We use this as the sign convention for the
Hochschild--Serre edge.

The edge is defined on the actual $\rho$ and
is an isomorphism here. The free pro-three
groups $N$ and $P$ have cohomology zero in
degree at least two for finite three-primary
coefficients. Their finite generating sets
make the cohomology at finite levels finite,
so the inverse-limit Mittag--Leffler argument
and rationalization give the same vanishing
for the Tate lattices and rational modules in
use. Averaging over $C_2$ preserves it.
Consequently the only nonzero term of total
degree two in the Hochschild--Serre sequence
for $N\rtimes\Gamma$ is
$H^1(\Gamma,\operatorname{Hom}(V,\mathbb Q_3(2)))$;
its possible outgoing degree-two differential
has target $H^3(\Gamma,\mathbb Q_3(2))=0$.
This proves the claimed edge comparison, with
the compact projection given by restriction
to the actual summand $H$. Applying $\iota_w$
to (E.35.11) and comparing (E.35.10) proves
(E.35.5) before the invariant normalization.

We finish by proving that the normalization
is an isomorphism on the required invariant
summand. This is a finite exact calculation
on the actual Fermat automorphisms. Let
$\mathcal A=\mu_3^2\rtimes S_3$ act by
diagonal scalings followed by coordinate
permutations. It is the same order-$54$ group
as $E[3]\rtimes\mu_6$ in E.23. In a rational
CM basis of $H$ put
\[
 R=\begin{pmatrix}0&-1\\1&-1\end{pmatrix}.
\]
The automorphism with diagonal exponents
$(a,b,0)$ and permutation $p$ acts on $H$ by
$\operatorname{sgn}(p)R^{a+b}$. This follows
by pulling back the differential
$d(X/Z)/(Y/Z)^2$: its multiplier is
$\operatorname{sgn}(p)\zeta^{a+b}$, and the
conjugate eigenvalue supplies the other CM
line. On cusp homology it is the literal
permutation of the displayed projective
points, followed by passage to the diagonal
quotient.

Use the images of
$A_0,A_1,A_2,B_0,B_1,B_2,C_0,C_1$ as a basis
of $W$, with $C_2$ their negative sum. If $P_p$
is any of the nine-point permutation matrices,
its quotient action is $\pi P_p L$, where
$L$ includes the first eight coordinate vectors
and $\pi(v)_i=v_i-v_8$. In particular
$\pi P_p=(\pi P_pL)\pi$, which verifies the
module comparison explicitly.
Let $e_{\mathcal A}$ be the average of the
$54$ operators on $H\otimes\bigwedge^2W$.
Consider the two columns
\[
 d_0=e_{\mathcal A}(a_0\otimes(A_0\wedge B_0)),
 \qquad
 d_1=e_{\mathcal A}(a_0\otimes(A_0\wedge B_1)).
\]
Here $a_0$ is the first CM basis vector used
only for this rank calculation; its scale does
not alter the determinant's nonvanishing or
the intrinsic map (E.35.3). The two order vectors
in (E.35.6) give
\[
 (1_H\otimes b)(d_0,d_1)=
     \begin{pmatrix}-1/3&1/3\\0&1/3\end{pmatrix},
 \qquad \det=-1/9.
 \tag{E.35.12}
\]
For a short check of every entry, sum first
over the nine diagonal scalings. For each of
the coordinate permutations $012$ and $102$
the resulting two-column matrix is
$\left(\begin{smallmatrix}-9&9\\0&9\end{smallmatrix}\right)$.
For each of the other four permutations it
is zero. Dividing their sum by $54$ gives
(E.35.12). The program
\texttt{ross\_contraction.py} also records the
entire integral numerator of the projector,
its rational inverse on these two columns,
and the conjugation check. All use exact
integer and rational arithmetic.

Thus the invariant part of $J$'s image has
dimension at least two. The whole $HW^2$
invariant space has dimension two by E.23,
so it equals this part. Formula (E.35.12)
proves that its contraction is an isomorphism.
The maps $J$ and $b$ are Galois-equivariant,
and Galois permutes the elements of
$\mathcal A$, so its average commutes with
Galois. The class $[\eta]_{HW^2}$ lies in the
image of this average by E.23; $B_w=[V,w]$
has at least two compact factors, hence has
zero $HW^2$ projection. Averaging therefore
does not change that cohomology class. This
completes the comparison to exactly the
elliptic component in E.29.

For the final identification, $-1$ has
compatible rational $3^r$th roots, so its
three-adic Kummer class is zero and replacing
$h$ by $1-y$ leaves $\rho$ unchanged.
The symbol $\{1-x,1-y\}$ is Ross's specified
symbol; his Theorem~1 proves that its multiple
by six is tame on the proper cubic and has
nonzero real regulator \cite{Ross1994}.
We use the explicit Kummer cup in (E.35.4),
not that real-regulator conclusion, to define
the three-adic class in (E.35.5).
The searched syntomic formula
\cite[Theorem 4.11]{Asakura2023} assumes
$p>\max(N,M)$; with $p=N=M=3$ this reads
$3>3$ and supplies no value for the present
class. Proposition~\ref{E-kummer-values} below
evaluates the two rational Kummer coordinates;
the numerical value of this exact three-adic
cup edge and the integral torsion remain.
The anchor is \S51.16,
folios 733v--736v, scans 172--178, and
LM7-E.23(i).
\end{proof}

% ---- E_kummer_values ----
\begin{statement}{E.36}{Proved}{the two rational Kummer coordinates are $1/3$ and $1/2$}
\label{st:E.36}
\label{E-kummer-values}
Use the actual unit tangent for $q=Z/X$ at
$O=[1:-1:0]$, and the bracket and projector
normalizations of E.27--E.29. Put
\[
 v_6=\sum_{i=0}^2(e_{A_i}+e_{B_i}),
 \qquad v_2=e_{C_1}+e_{C_2}.
\]
For the genuine cubic action class $[\eta]$,
the complementary projection to the elliptic
component is
\[
 (1-e_{\mathcal A})[\eta]
   =\left[\left(
          \frac13v_6+\frac12v_2\right)\kappa_3(3),w\right]
 \quad\text{in }H^1(G_{\mathbb Q,\{3,\infty\}},
                         A_{\mathbb Q_3}(-1)).
 \tag{E.36.1}
\]
The bracket in this formula means the map
$H^1(V_{\mathbb Q_3})\to H^1(A_{\mathbb Q_3}(-1))$
induced by $v\mapsto[v,w]$.
Thus $(\lambda,\lambda,\mu)$ in the three
Tate coordinates of E.27 is exactly
\[
                  (1/3,1/3,1/2).
 \tag{E.36.2}
\]
Together with E.35, this expresses the full
rational class as these evaluated Kummer
terms plus the inverse image, under the
specified isomorphism $\Phi(-1)$, of
$-\iota_{w*}\operatorname{edge}_H
(\kappa(1-x)\smile\kappa(1-y))$.
The latter three-adic class and the integral
torsion remain to be evaluated.
\end{statement}
\begin{proof}
We work first over $K=\mathbb Q(\zeta_3)$,
where every element of
$\mathcal A=E[3]\rtimes\mu_6$ is defined.
For $a\in\mathcal A$, choose a geometric
base-point path identifying its image of
the tangent at $O$ with our fixed base point.
The automorphism of the arithmetic fundamental
group supplied by the curve automorphism
then has the form
\[
 n\longmapsto\sigma_a(n),\qquad
 s_O(g)\longmapsto k_a(g)s_O(g).
 \tag{E.36.3}
\]
Here $\sigma_a$ is a continuous automorphism
of the geometric pro-three group, $k_a$ is
an actual continuous section cocycle, and
its section class is the image tangent
$a(s_O)$. The arithmetic homomorphism law
gives the exact identity
\[
 \sigma_a\rho_g\sigma_a^{-1}
        =\operatorname{Int}(k_a(g))\rho_g.
 \tag{E.36.4}
\]
Changing the chosen path conjugates these
maps geometrically. After abelianization,
$\alpha_a=[k_a^{\mathrm{ab}}]\in H^1(G_K,V_{\mathbb Q_3})$
is independent of that change.

Let $\omega$ be a lift of $w$ in the rational
degree-two-to-three extension defining $\eta$.
Every $a$ fixes $w$ on homology, so write
$\sigma_a\omega=\omega+t_a$ with $t_a\in A_{\mathbb Q_3}$.
On this extension, inner conjugation by
$k_a(g)$ adds $[k_a(g)^{\rm ab},w]$ to a lift
of $w$. Apply (E.36.4) to $\omega$ and use
$\rho_g\omega=\chi(g)(\omega+\eta_g)$.
The result, after dividing by $\chi(g)$, is
\[
 (a-1)\eta_g
    =\chi(g)^{-1}g(t_a)-t_a
                    +[k_a(g)^{\rm ab},w].
 \tag{E.36.5}
\]
Thus the first two terms form precisely the
coboundary in the twisted module $A(-1)$.
Taking classes and averaging the actual
$54$ identities gives
\[
 (1-e_{\mathcal A})[\eta]
       =-[\overline\alpha,w],
 \qquad
 \overline\alpha=\frac1{54}\sum_{a\in\mathcal A}\alpha_a.
 \tag{E.36.6}
\]
This proves the sign and identifies the
required coordinate with actual path classes.

Under the compact projection
$V_{\mathbb Q_3}\to H$, the class $\alpha_a$
is the Kummer class of $aO-O\in E[3]$.
It vanishes after tensoring with $\mathbb Q_3$.
The fixed rational splitting $V=H\oplus W$
therefore identifies $\overline\alpha$ with
a class in $H^1(G_K,W)$. This is a statement
about the cohomology class; the initially
chosen path cocycles need not have zero
compact coordinates as actual cochains.
The following unit pairings determine this
boundary class exactly.

For each cusp $c$ define its normalized
tangent-line equation by
\[
 \begin{aligned}
 L_{A_i}&=\zeta^{2i}Y+Z,&
 L_{B_i}&=\zeta^{2i}X+Z,&
 L_{C_i}&=\zeta^{2i}X+Y.
 \end{aligned}
\]
In particular $L_O=X+Y$. For $c\ne O$ the
unit $f_c=L_c/L_O$ on $U_K$ has divisor
$3(c-O)$. Its geometric Kummer linear form
on the boundary quotient is
\[
                   \ell_{f_c}=3(\delta_c-\delta_O).
 \tag{E.36.7}
\]
These eight linear forms are a basis of
$\operatorname{Hom}(W,\mathbb Q_3(1))$.
As in E.35 they vanish on the compact
summand. For a function with a zero or
pole at $O$, write $\operatorname{lt}_O$
for its first nonzero coefficient in the
specified parameter $q$.
Functoriality of the Kummer map applied to
the two sections in (E.36.3) gives
\[
 \ell_{f_c*}(\alpha_a)=
     \kappa_3\bigl(\operatorname{lt}_O(a^*f_c)\bigr)
       -\kappa_3\bigl(\operatorname{lt}_O(f_c)\bigr).
 \tag{E.36.8}
\]
Indeed the first coefficient in the
pullback is the leading value of $f_c$
on the transported tangent $a(s_O)$;
the second is its leading value on $s_O$.
Their difference is exactly the geometric
abelianized section cocycle under the
map induced by the unit $f_c$.

We evaluate the mean of (E.36.8).
The monomial automorphisms permute the nine
tangent lines up to scalars in $\mu_3$.
Every line occurs six times in each
product $\prod_{a\in\mathcal A}a^*L_c$.
Consequently
\[
       \prod_{a\in\mathcal A}a^*f_c\in\mu_3.
 \tag{E.36.9}
\]
The orders at $O$ in this product sum to
zero, so its leading coefficient is also
the product of the individual leading
coefficients. A root of unity has zero
rational three-adic Kummer class. Averaging
(E.36.8) therefore gives the exact equality
\[
 \ell_{f_c*}(\overline\alpha)
      =-\kappa_3\bigl(\operatorname{lt}_O(f_c)\bigr).
 \tag{E.36.10}
\]
The program \texttt{kummer\_coordinates.py}
checks every pullback of every tangent
line; for these normalizations all nine
ratio products in fact equal one. Only
their displayed root-of-unity property
is needed for (E.36.10).

On the branch through $O$,
$Y/X=-(1+q^3)^{1/3}$, so
\[
             L_O/X=-q^3/3+O(q^6).
\]
For $c\ne O$ this gives
\[
 \operatorname{lt}_O(f_c)=-3L_c(1,-1,0).
 \tag{E.36.11}
\]
For $A_i$ and $B_i$ these numbers are
$3\zeta^{2i}$ and $-3\zeta^{2i}$,
respectively. Their rational Kummer class
is $\kappa_3(3)$. For $C_1,C_2$ they are
$3(1-\zeta^{2i})$. The identities
\[
 (1-\zeta)^2=-3\zeta,\qquad
 1-\zeta^2=-\zeta^2(1-\zeta)
\]
show that their rational Kummer classes
are $\tfrac32\kappa_3(3)$.

To solve the resulting eight equations,
represent a boundary vector by its nine
cusp coefficients modulo the common
diagonal, and choose the representative
whose $O$ coefficient is zero. In these
coordinates (E.36.7) is three times the
$c$ coefficient. Equations
(E.36.10)--(E.36.11) thus give
\[
 \overline\alpha=
      -\left(\frac13v_6+\frac12v_2\right)\kappa_3(3)
             \quad\text{in }H^1(G_K,W).
 \tag{E.36.12}
\]
No information in another boundary direction
is discarded: the eight unit forms are a
basis and all their values have been computed.
The same program checks their eight rational
linear equations and records the resulting
vector in this exact quotient presentation.

Both $v_6$ and $v_2$ are stable under
$G_{\mathbb Q}$, with their Tate action,
and $\kappa_3(3)$ is defined over
$\mathbb Q$. The restriction map in
degree-one rational cohomology for the
quadratic extension $K/\mathbb Q$ is
injective by finite-group averaging. The inflation
map from the unramified arithmetic quotient
to $G_{\mathbb Q}$ is also injective in
degree one. Thus (E.36.6) and (E.36.12)
prove (E.36.1) over the original specified
arithmetic group. In the basis of E.27,
$v_6=b_0+b_1+b_2$ and $v_2=b_3$, giving
(E.36.2). The denominator three belongs
to the rational projector calculation;
it supplies no assertion about an
integral splitting or integral torsion.
E.28 eliminates the other projected
multidegree, and E.35 identifies the
remaining elliptic component, proving
the final assertion. The manuscript
anchor is \S51.16, folios 733v--736v,
scans 172--178, and LM7-E.23(i).
\end{proof}

% ---- E_integral_detector ----
\begin{statement}{E.37}{Proved}{an integral action detector and a zero-abelian obstruction at nineteen}
\label{st:E.37}
\label{E-integral-detector}
Retain the cubic lattice $A=\operatorname{Lie}_3(V)$,
the genuine normalized cocycle $\eta$ of E.19,
and the unit tangent at $O$ used in E.36.
There is an explicit surjective integral
$G_{\mathbb Q}$-module map
\[
 \Psi:A(-1)\longrightarrow\mathbb Z_3(1)
 \quad\text{with}\quad
 \Psi\eta(g)=6\kappa_3(3)(g)
 \quad(g\in G_{\mathbb Q}).
 \tag{E.37.1}
\]
The Kummer cocycle is normalized using the
positive real compatible roots of $3$.
In particular the class $[\eta]$ is nonzero
already after the displayed coefficient
pushout and reduction modulo nine.

The actual zero-abelian class-two section
\[
 s_{19}(g)=\kappa_3(19)(g)w\,s_O(g)
       \quad\text{in }N_2\rtimes G_{\mathbb Q}
 \tag{E.37.2}
\]
has no lift to $N_3\rtimes G_{\mathbb Q}$.
Here $\kappa_3(19)w$ is written additively
in the central subgroup
$\gamma_2N/\gamma_3N=\bigwedge^2V$.
The obstruction is detected by the actual
central coefficient pushout
$A\to\mathbb Z/9(2)$ induced by $\Psi$,
at the prime $19$, where its residue is
$6\in\mathbb Z/9$. Consequently no lift
to the geometrically pro-three metabelian
group, or to the full geometric group,
can have (E.37.2) as its class-two image.
\end{statement}
\begin{proof}
Let $p:V\to T_3E$ be the actual compact
projection of E.18. Fix its determinant
pairing so that $p(w)=9$ in
$\bigwedge^2T_3E=\mathbb Z_3(1)$.
Use the global unit
\[
                 f=(X+Z)/(X+Y),
       \qquad\operatorname{div}(f)=3(B_0-O).
\]
Its geometric Kummer map
$\ell_f:V\to\mathbb Z_3(1)$ is integral
and vanishes on the rational compact
summand. Embed the free degree-three
Lie lattice in its tensor algebra and
apply the literal map
\[
 \Phi=(\ell_f\otimes\det(p,p)):
    A\lhook\joinrel\longrightarrow V^{\otimes3}
       \longrightarrow\mathbb Z_3(2).
 \tag{E.37.3}
\]
It is integral and Galois-equivariant.
The first-letter terms of $[z,[u,v]]$
give
\[
                  \Phi([z,w])=18\ell_f(z).
 \tag{E.37.4}
\]
The other two terms have first letter in
$H$, on which $\ell_f$ is zero. The same
factor count shows that $\Phi$ kills the
$HW^2$ multidegree.

Here is its complete integral row, so
that division by three is checked on
the actual lattice. In E.18's ordered
basis write $e_0=f_1$, $e_1=f_2$,
followed by the eight nonorigin cusps;
then $e_5=e_{B_0}$. The relation
$j_{0*}=\binom{3I}{U}$ gives
\[
 \ell_f=(0,2,0,0,0,3,0,0,0,0).
\]
On the Hall basis $[[e_i,e_j],e_k]$,
with $i>j$ and $k\ge j$, the only nonzero
values of $\Phi$ are
\[
 \begin{array}{c|rrr}
 (i,j,k)&(1,0,1)&(1,0,5)&(5,0,1)\\\hline
 \Phi&6&6&3\\
 \Psi&2&2&1
 \end{array}
 \tag{E.37.5}
\]
Thus $\Phi/3$ is an integral surjection;
after twisting by $(-1)$ it is exactly
the map $\Psi$ in the statement. The
division preserves equivariance because
the target lattice is torsion-free.
The program
\texttt{zero\_slice\_inverse\_limit.py}
constructs this row from the tensor
formula and checks $\Psi T=49\Psi$ and
$\Psi J=\Psi$ for the actual integral
Frobenius and conjugation matrices.

By E.35--E.36 the rational action class
consists of its evaluated bracket term
and an $HW^2$ elliptic term. The latter
is killed by (E.37.3). Since
$\ell_f(v_6/3+v_2/2)=1$, equations
(E.36.1) and (E.37.4) give
$\Psi_*[\eta]=6\kappa_3(3)$ rationally.
Both classes are unramified outside
three. The integral group
$H^1(G_{\mathbb Q,\{3,\infty\}},\mathbb Z_3(1))$
is the torsion-free module
$\mathbb Z_3\kappa_3(3)$, as proved in
E.27. The same equality therefore holds
integrally. Its possible coboundary is
$(\chi(g)-1)b$ for a fixed $b\in\mathbb Z_3(1)$.
Both cocycles vanish at complex
conjugation $c$: this is the normalization
of E.19 and the chosen real Kummer roots.
At $c$ the coboundary is $-2b$, hence
$b=0$. This proves the actual cocycle
identity (E.37.1).
At the normalized arithmetic generator
$\tau$ of E.31 its value is six.
Modulo nine a Tate coboundary vanishing
at $c$ is again zero. Thus the coefficient
class modulo nine is nonzero as claimed.

For (E.37.2), the equivariant inclusion
$\mathbb Z_3(1)w\subset\bigwedge^2V$
makes $\kappa_3(19)w$ an actual central
class-two cocycle. E.19 identifies its
connecting obstruction with
$[\eta\smile\kappa_3(19)]$.
Apply the map (E.37.1) to the first
factor. At $19$ its scalar tame
correction would have to be
\[
       -\frac{6\kappa_3(3)(F_{19})}{19-1}
          =-\frac{\kappa_3(3)(F_{19})}{3}.
 \tag{E.37.6}
\]
The unramified Kummer value has valuation
zero: $3^6\equiv7\pmod{19}$, so three
is not a cube in the residue field.
More precisely choose the compatible
root-of-unity coordinate with
$\zeta_9\equiv4\pmod{19}$. This element
has exact order nine, and
\[
         3^2\equiv9\equiv4^4\pmod{19}.
\]
The Frobenius Kummer value is therefore
$4\pmod9$. The right side of the scalar
equation before division is
$6\cdot4=6\pmod9$, whereas multiplication
by $19-1$ is zero modulo nine.
This is a finite coefficient obstruction
to the central pushout. In particular
the correction in (E.37.6) is not
integral, and no original lift exists.
The program verifies all of these finite
field residues.

Finally $N''\subset\gamma_4N$ gives the
continuous quotient map
$N/N''\twoheadrightarrow N_3$,
compatible with the genuine arithmetic
action and the map to $N_2$.
The full geometric group maps to this
pro-three quotient as well. Each proposed
lift along either map would give the
excluded class-three lift. These are
the comparison maps used in the last
assertion. The anchor is \S51.16,
folios 733v--736v, scans 172--178,
and LM7-E.23(i)--(ii).
\end{proof}

\begin{statement}{E.38}{Proved}{every finite coefficient lifts, but the continuous class-three lift fails}
\label{st:E.38}
\label{E-zero-slice-inverse-limit}
In the same intrinsic cubic groups there
is an explicitly specified continuous
central class-two cocycle
\[
       z(g)=\beta(g)w,
 \qquad
       \beta=\sum_{n\ge1}3^{n+6}\kappa_3(p_n),
 \tag{E.38.1}
\]
with the distinct rational primes $p_n$
specified below by finite Frobenius
conditions, whose class-three connecting
class satisfies
\[
 \begin{gathered}
 0\ne[\eta\smile\beta]
             \in H^2_{\rm cts}(G_{\mathbb Q},A),\\
 [\eta\smile\beta]\longmapsto0
             \quad\text{in }H^2(G_{\mathbb Q},A/3^mA)
                    \quad\text{for every }m\ge1.
 \end{gathered}
 \tag{E.38.2}
\]
Equivalently, $z$ lifts through each
actual central extension
$N_3/(3^mA)\to N_2$, and has no
continuous lift through $N_3\to N_2$.
No compatible system of those finite
coefficient lifts is asserted. Every
finite partial sum in (E.38.1) does
have an integral continuous class-three
lift.
\end{statement}
\begin{proof}
We specify the primes and all uniform
integrality bounds. Let $T_V$ be E.18's
actual Frobenius matrix at seven, after
the normalization in E.31. Thus the
arithmetic free generators satisfy
\[
 \chi(\tau)=\chi(\tau')=7,
 \quad \kappa_3(3)(\tau)=1,
 \quad \kappa_3(3)(\tau')=-1,
 \quad \tau'=c\tau c.
 \tag{E.38.3}
\]
Their actions on $V$ commute: over
$K=\mathbb Q(\zeta_3)$ the compact CM
image is abelian and the boundary image
is cyclotomic, and the rational compact
splitting is equivariant. Equality of
the rational matrices is equality of
these integral matrices. Set
$r=[\tau,\tau']=\tau\tau'\tau^{-1}\tau'^{-1}$.
The cocycle law gives exactly
\[
           \rho_V(r)=1,\qquad
           \chi(r)=1,\qquad
           \kappa_3(3)(r)=-12.
 \tag{E.38.4}
\]
For $n\ge1$ define
\[
                  g_n=r\tau^{3^{n+1}}.
 \tag{E.38.5}
\]
These are elements of the actual
arithmetic quotient $\Gamma$.
Their finite images are thus in the
image of $G_{\mathbb Q}$.
The geometric-series cocycle formula
and the elementary lifting-the-exponent
calculation give
\[
 \begin{aligned}
 \rho_V(g_n)&=T_V^{3^{n+1}},\\
 \chi(g_n)&=7^{3^{n+1}},
       &v_3(\chi(g_n)-1)&=n+2,\\
 \kappa_3(3)(g_n)&=-12+
       \frac{7^{3^{n+1}}-1}{6},
       &v_3(\kappa_3(3)(g_n))&=1.
 \end{aligned}
 \tag{E.38.6}
\]
For the last equality the second term
has valuation $n+1\ge2$, while the
first has valuation one.

Here is the fixed matrix certificate
needed for every $n$. On $A(-1)$ put
\[
 R=7^{-1}\operatorname{Lie}_3(T_V),
       \qquad D=R^3-I.
\]
The actual residual $T_V$ is unipotent
of exponent three, so $D\in3M_{330}(\mathbb Z_3)$.
Its complete Smith valuation list is
\[
                 1^{132},2^{82},3^{52},
                       4^{49},5^{14},6^1.
 \tag{E.38.7}
\]
The executed program constructs every
entry, finds all $330$ pivots at
precision $3^{12}$, and independently
checks determinant valuation $724$.
For that second check put
$\pi=1+3\zeta$, so
$\pi+\bar\pi=-1$, $\pi\bar\pi=7$.
The contributions to the determinant
valuation from the five eigenvalue
blocks of $R^3-I$ are
\[
                  4+8\cdot4+16\cdot2
                       +64\cdot5+168\cdot2=724.
\]
These follow respectively from
$\operatorname{Norm}(\pi^3-1)$,
$\operatorname{Norm}(\pi^6-1)$,
$7^3-1$, $\operatorname{Norm}((7\pi)^3-1)$,
and $7^6-1$; the multiplicities are
the $H^3,H^2W,HW^2,W^3$ multidegrees.
All pivots lie below the working
precision, so the reported valuations
lift under integral row and column
operations. In particular $D$ is
invertible rationally and $3^6D^{-1}$
is integral.

Define
$D_n=(I+D)^{3^n}-I=R^{3^{n+1}}-I$.
For any matrix $B\in3M_{330}(\mathbb Z_3)$,
\[
 (I+B)^3-I=3B(I+B+B^2/3),
          \qquad I+B+B^2/3\in I+3M_{330}(\mathbb Z_3).
\]
Induction gives $D_n=3^nD U_n$ for an
integral invertible matrix $U_n$.
Therefore
\[
                        3^{n+6}D_n^{-1}
                             \in M_{330}(\mathbb Z_3).
 \tag{E.38.8}
\]
This proves the bound at every depth,
not merely at the computed depth.

To choose $p_n$, take the finite image
of the actual homomorphism consisting
of $\rho_V$, $\chi$, and the affine
Kummer representation $(\kappa_3(3),\chi)$,
all reduced modulo $3^{n+8}$.
The last pair is a homomorphism into
the corresponding affine semidirect
product by the cocycle law. Let $p_n$
be the least prime larger than $p_{n-1}$
and different from three whose Frobenius
conjugacy class in this finite image
is that of $g_n$; set $p_0=3$.
Chebotarev gives infinitely many such
primes at each step, so this specifies
the entire sequence. Choose a
prolongation with representative equal
to the chosen finite image of $g_n$.
The condition modulo three puts $p_n$
in the split class of $K/\mathbb Q$.
The two primes of $K$ above $p_n$ are
related by integral complex conjugation;
the estimates below hold at both.
The specification uses only the named
finite arithmetic representations and
their actual Frobenius classes. We have
not enumerated these primes numerically.

Write
$D_{p_n}=p_n^{-1}\operatorname{Lie}_3
(\rho_V(F_{p_n}))-I$.
The finite Frobenius condition gives
$D_{p_n}-D_n\in3^{n+8}M_{330}(\mathbb Z_3)$.
By (E.38.8),
$D_n^{-1}(D_{p_n}-D_n)\in9M_{330}(\mathbb Z_3)$.
The Neumann series makes the resulting
$I+9M$ matrix invertible integrally.
Thus the inverse bound (E.38.8) also
holds with $D_{p_n}$, and (E.38.6)
gives the exact valuations
\[
 v_3(p_n-1)=n+2,
 \qquad v_3(\kappa_3(3)(F_{p_n}))=1.
 \tag{E.38.9}
\]

The sum (E.38.1) converges uniformly
as a function on $G_{\mathbb Q}$:
every normalized integral Kummer
cocycle has values in $\mathbb Z_3$,
and $3^{n+6}$ tends to zero. Its limit
is therefore a continuous Tate cocycle,
and $z=\beta w$ is an actual central
class-two cocycle. The unique rational
tame correction prescribed by E.19
at $p_n$ is
\[
 x_n=-3^{n+6}D_{p_n}^{-1}\eta(F_{p_n}).
 \tag{E.38.10}
\]
It is integral by the proved matrix
bound. All other primes outside three
have zero residue and zero correction.
Nevertheless these integral corrections
do not tend to zero. Equivariance of
E.37's integral map says
$\Psi D_{p_n}=(p_n-1)\Psi$; hence
\[
 \Psi x_n=-\frac{3^{n+6}\,6\kappa_3(3)(F_{p_n})}
                            {p_n-1},
       \qquad v_3(\Psi x_n)=6
                  \quad\text{for every }n.
 \tag{E.38.11}
\]
An integral linear map takes a sequence
tending to zero to one tending to
zero. Formula (E.38.11) therefore
violates precisely the restricted
continuity condition of (E19.5).
That exact criterion proves the first
nonvanishing in (E.38.2).

For a finite partial sum of $\beta$,
(E.38.10) is integral and has finite
support, so it satisfies the same
criterion and gives an integral
continuous class-three lift. For
fixed $m$, all terms with $n+6\ge m$
vanish modulo $3^m$. The connecting
cocycle $\eta\smile\beta$ modulo
$3^mA$ therefore equals the reduction
of the obstruction of one finite
partial sum. The latter is zero.
Central-extension obstruction theory
then gives the stated lift through
$N_3/(3^mA)$ for each $m$.
The preceding nonzero continuous
obstruction proves that these
existence statements do not supply
a compatible inverse-limit lift.
This is the intrinsic cubic
zero-abelian instance of the exact
finite-coefficient comparison asked
in LM7-E.23(i), with its actual
arithmetic action retained.
\end{proof}

% ---- E_integral_path_comparison ----
\begin{statement}{E.39}{Proved}{the exact torsion lost by the compact and unit path comparisons}
\label{st:E.39}
\label{E-integral-path-comparison}
Retain the cubic lattice $V$, the cusp order and the matrix $U$
of (E18.8)--(E18.9), and put $K=\mathbb Q(\zeta_3)$.
For each nonorigin cusp $d$, let $f_d=L_d/L_O$ be the
normalized tangent-line unit of E.36, and let
$\ell_d:V\to\mathbb Z_3(1)$ be its geometric Kummer map.
Write $T=T_3E$ for the compact Tate module. Give
$B=\mathbb Z_3(1)^8$ the Galois action which permutes the
eight unit coordinates and multiplies by the cyclotomic
character. There is an actual equivariant injection
\[
 j:V\longrightarrow T\oplus B,\qquad
 v\longmapsto\bigl(p(v),(\ell_d(v))_d\bigr),\qquad
 j=\begin{pmatrix}I_2&0\\-U&3I_8\end{pmatrix}.
 \tag{E.39.1}
\]
Its cokernel $Q$ is $(\mathbb Z/3)^8$. The action on $Q$
over $K$ is trivial; over $\mathbb Q$, complex conjugation
acts by minus the permutation
\[
 (A_1\ A_2)(B_1\ B_2)(C_1\ C_2),
 \tag{E.39.2}
\]
with $A_0,B_0$ fixed by the permutation.
For continuous cohomology of the full absolute Galois
groups, and also of their quotients unramified away from
three and infinity, the exact kernels are
\[
 \begin{aligned}
 \ker H^1(G_K,j)&\simeq(\mathbb Z/3)^8,\\
 \ker H^1(G_{\mathbb Q},j)&\simeq(\mathbb Z/3)^3.
 \end{aligned}
 \tag{E.39.3}
\]
These isomorphisms are the connecting maps from $Q^{G_K}$
and $Q^{G_{\mathbb Q}}$, respectively. In particular, a
nonempty fibre of this cohomological comparison has
exactly $3^8$ elements over $K$ and $3^3$ over
$\mathbb Q$. Compact Kummer classes and all eight unit
cohomology classes therefore do not determine an integral
abelianized path cohomology class. One unresolved datum
in the integral version of E.36 is the actual class's
position in this explicitly specified fibre.
\end{statement}
\begin{proof}
The divisor of $f_d$ is $3(d-O)$, so its restriction to
the boundary lattice is the row $3e_d^*$. The compact
rational splitting of E.18 has columns $\binom{3I_2}{U}$.
The Kummer map of a geometric unit vanishes on that
compact summand: equivariance at a good Frobenius compares
its weight-one compact eigenvalues with the weight-two
cyclotomic eigenvalue. Thus
\[
 \ell_d=(-U_d,3e_d^*),
\]
which proves (E.39.1) with its actual unit and compact
maps. Galois equivariance follows either from these maps
or directly from the matrices (E18.9). Subtracting the
lower left block using integral row operations gives
$\operatorname{diag}(I_2,3I_8)$. The quotient map can be
written
\[
 T\oplus B\longrightarrow Q,\qquad
 (h,b)\longmapsto b+Uh\pmod3.
 \tag{E.39.4}
\]
Over $K$, the cusp permutation is trivial and
$\chi_3\equiv1\pmod3$. Applying (E18.9) to (E.39.4)
therefore gives the trivial action. Over $\mathbb Q$,
the same matrix identity gives (E.39.2).

The invariant subgroup $(T\oplus B)^G$ is zero for each
of the four indicated arithmetic groups. For $T$ this
follows by choosing a good Frobenius whose eigenvalues
have complex absolute value greater than one. For $B$,
first pass to $K$, where it is a direct sum of
$\mathbb Z_3(1)$; the image of the cyclotomic character
is infinite, so no nonzero vector is fixed. These proofs
apply to the stated unramified quotients because the
actual representations factor through them.
The continuous cohomology sequence of
\[
 0\longrightarrow V\xrightarrow{j}T\oplus B
       \longrightarrow Q\longrightarrow0
\]
then identifies $Q^G$ injectively with the kernel in
degree one. The finite quotient admits a continuous
set-theoretic section, so this is the ordinary long exact
sequence for these continuous coefficient modules.
The fixed space of minus the permutation (E.39.2) has
the three basis vectors
\[
 e_{A_1}-e_{A_2},\qquad e_{B_1}-e_{B_2},\qquad
 e_{C_1}-e_{C_2}.
\]
The two fixed permutation coordinates contribute no
fixed vectors in characteristic three. This proves
(E.39.3) and the fibre cardinalities.

The connecting classes can also be read as explicit
cocycles. Let $t_d$ be the $d$th boundary basis vector
of $V$, and lift the $d$th basis vector of $Q$ rationally
by $t_d/3$. Its connecting cocycle is
\[
 b_d(g)=g(t_d/3)-t_d/3.
 \tag{E.39.5}
\]
Over $K$ this is integral and, on the actual free
arithmetic generators of E.38, has
$b_d(\tau)=b_d(\tau')=2t_d$. The three displayed
differences give the invariant classes over
$\mathbb Q$, with value zero on complex conjugation.
Their nonzero order three follows from the injective
connecting map; multiplying (E.39.5) by three gives
the coboundary of $t_d$.

The program \texttt{integral\_path\_comparison.py}
checks the exact coefficient maps, both arithmetic
equivariance identities, all ten Smith factors, the
three invariant vectors, and these cocycle values.
In the attempted integral automorphism calculation,
substituting arbitrary representatives of the compact
and unit cohomology classes can change the result by
(E.39.5). No such representative is asserted to be the
actual tangential path. The remaining comparison must
retain the correlated choices of compact division
points and unit roots. The manuscript anchor is
\S51.16, folios733v--736v, scans172--178, and
LM7-E.23(i).
\end{proof}

% ---- F ----
\section{The notes of 1983: centralizers, descent, and tangential parameters}
\label{sec:F}

We write $G_K=\operatorname{Gal}(\overline K/K)$,
$H_K=\ker(\chi_K:G_K\to\widehat{\mathbb Z}^{\times})$,
$Q_K=\chi_K(G_K)$, $T=\widehat{\mathbb Z}(1)$,
$P_U=\pi_1(U_{\overline K})$, and $E_U=\pi_1(U)$.
The notation $\widehat{A}$ for an abelian group means
$\varprojlim_n A/nA$, or $\varprojlim_n A/A^n$ in multiplicative
notation. Its terms need not be finite. All section classes are
classes under geometric inner conjugation. An actual section, an
actual lifted cusp, and an actual subgroup are specified whenever a
centralizer is computed. The comparisons in this section concern cote
149, scans 44--63, 85, 100--103, and 126--127. We retain the two written
subscripts $L$ on scan 127 and separately state their mathematical
correction.

\subsection{F(a)--(c): what the centralizer calculations prove}

\begin{statement}{F.1}{Known}{the cyclotomic centralizer}
\label{st:F.1}
\label{F-centralizer}
Let $K/\mathbb Q$ be finitely generated and let $U/K$ be a smooth
geometrically connected hyperbolic curve. For every continuous section
$s:G_K\to E_U$,
\[
 C_{P_U}(s(H_K))=1
 \quad\hbox{or}\quad
 C_{P_U}(s(H_K))=I_c\simeq T,
\]
where the second case holds precisely when $s(G_K)$ is contained in
the decomposition group of a unique lifted rational cusp $c$.
The centralizer of a point section is $1$. Restricting the section to
an open subgroup of $H_K$ preserves this centralizer. This gives the
two ``pr\'ecisions'' on cote 149, scan 56, and the formula on scan 55.
More generally, for a closed subgroup $B\subset E_U^{\circ}$
whose projection to $H_K$ is open, $C_{P_U}(B)$ is trivial or
one entire cusp inertia. Consequently the same proof recognizes
every nontrivial procyclic subgroup essentially cyclotomically
normalized over an open subgroup of $G_K$; no splitting of that
normalizer is required.
\end{statement}
\begin{proof}
We checked CYCLOTOMIC\_KERNEL A1--A7 and AE1. We record the steps
which carry the arithmetic assertion to the full profinite group;
in particular, we do not discard its common kernel of pro-prime
quotients. Ribet's Appendix, Theorem 1, gives finite torsion in
$A(k(\mu_\infty))$ for an abelian variety over a number field
$k$~\cite{F-Ribet}. For a finitely generated $K/k$, spread $A$ to an
abelian scheme over a smooth integral $k$-variety $S$. Each rational
torsion point over $Kk(\mu_\infty)$ extends over $S_{k(\mu_\infty)}$:
the appropriate torsion scheme is finite \'etale, and its generic
degree-one component extends by normality. Evaluating these sections
at a closed point $x$ gives an injection into
$A_x(k(x)(\mu_\infty))_{\rm tors}$. Equality at one fibre is equality
everywhere since the equalizer of two sections of a finite \'etale
scheme is open and closed. Thus every compactified covering Jacobian
has zero $H_K$-invariants in its full Tate module.

For $z\in C_{P_U}(s(H_K))$, and every open characteristic
$N\subset P_U$, put $J=N\overline{\langle z\rangle}$. An open
subgroup of $G_K$ stabilizes $J$, and $Js(G_{K'})$ defines the
corresponding finite cover. The image of $z$ in $J^{\rm ab}$ maps
to zero in the Tate module of its compactification. The localization
sequence is
\[
 T\longrightarrow\bigoplus_{d\in D_J} T e_d
 \longrightarrow J^{\rm ab}\longrightarrow T\operatorname{Jac}(X_J)
 \longrightarrow0,
\]
with the diagonal as the first map. In the proper case the cyclic
quotient $J/N$ then forces $z=1$. In the affine case its image in
$J/N$ belongs to the subgroup generated by cusp inertia images.
As $zN$ generates the finite cyclic group $J/N$, those images generate
all of $J/N$. For a prime $\ell$, subgroups of a finite cyclic
$\ell$-group are linearly ordered. Hence one of these inertia images
contains the full $\ell$-Sylow subgroup of $J/N$, and therefore
contains the image of the $\ell$-primary component $z_\ell$ of $z$.
This inertia image is the stabilizer of a cusp in the normal cover
with subgroup $N$. Thus $z_\ell$ fixes a cusp at this level.
For any finite collection of levels a common normal refinement
gives compatible fixed cusps, so compactness gives a lifted cusp
fixed by $z_\ell$. This is the finite-cyclic argument of CK A3;
no assertion that an arbitrary sum of boundary classes belongs
to one boundary line is used. Distinct cusp inertias
intersect trivially, and the centralizer of a nontrivial element of
one inertia is that inertia (Stix, Proposition 11 and Corollary
12~\cite{F-Stix}). Since the primary components commute, they lie in
one inertia: fix any nonzero primary component, and use its
inertia as its geometric centralizer to contain every other
component. The element $z$ is the limit of its finite products
of primary components (in each finite cyclic quotient only
finitely many primes occur). The inertia is closed, so it
contains $z$. Thus compactness across covering levels and the
comparison across different primes are separate justified steps.

Here is the necessary passage from elements to the entire fixed
subgroup, checked in CK A5. A closed subgroup all of whose nonidentity
elements lie in inertia lies in a single inertia. Otherwise choose
two distinct orbits of its cusps: a finite group cannot be the union
of conjugates of one proper subgroup, by counting fixed cosets.
They remain distinct in a finite cover containing the subgroup.
Refine to at least three cusps; when there are exactly two, the
compactification has positive genus and an \'etale double cover of
it supplies four. More explicitly, if $C$ is the fixed subgroup,
choose the two cusps in distinct $C$-orbits and an open normal
$N$ at which these orbits are still distinct. Such an $N$ exists
by compactness of $C$. The cover with subgroup $CN$ contains $C$
and separates the two cusps. Every element of $C$ maps to the
identity in the fundamental group of its compactification,
since it is an inertia element in this cover. Hence the pullback
of the stated compact double cover still contains $C$; no
arbitrary refinement containing $C$ is being presumed.
In the boundary module, take nonzero images
$ae_c,be_{c'}$. Both products giving $ae_c+be_{c'}$ and
$2ae_c+be_{c'}$ would have to lie in individual boundary lines.
With at least four cusps neither can do so; with three the two
requirements imply $a=b$ and $2a=b$, hence $a=0$. This contradiction
works over $\widehat{\mathbb Z}$ without dividing by a zero divisor.
We obtain $C\subset I_c$. A nonidentity $z\in C\cap I_c$ belongs
also to every $s(h)I_cs(h)^{-1}$ for $h\in H_K$. Trivial intersection
of distinct cusp inertias therefore makes $s(H_K)$ normalize $I_c$,
and the kernel of the cyclotomic character acts trivially on it;
therefore $I_c\subset C$. Normality of $H_K$ in $G_K$ makes the whole
section normalize $I_c$, giving the asserted rational cusp.

The separation from point sections uses the Jacobian Kummer map
after a cover whose compactification has positive genus. Here is
the precise divisible-point argument. First let $L=k(\mu_\infty)$
for a number field $k$, and descend an abelian variety $A/L$ and
a divisible point $P$ to a number field $k'\subset L$ with
$L=k'(\mu_\infty)$. The finite nonempty sets of factorial-division
points of $P$ in $A(L)$ admit a compatible inverse limit. Its
Kummer cocycle $c:G_{k'}\to T(A)$ vanishes on the normal subgroup
$H=G_L$. For $h\in H$, $g\in G_{k'}$, the cocycle identity evaluated
at $hg=g(g^{-1}hg)$ gives $h c(g)=c(g)$. The torsion calculation
above gives $T(A)^H=0$, so $c=0$. All division points are already
defined over $k'$. Reducing at one place of good reduction modulo
every power of its maximal ideal sends $P$ to a divisible element
of a finite group, hence to zero. These reductions separate points,
so $P=0$.

For finitely generated $K$, every finite extension of $K(\mu_\infty)$
is $K'(\mu_\infty)$ for one finitely generated $K'/\mathbb Q$.
Descend $A$ and a proposed nonzero divisible $P$ to $K'$, choose
a smooth geometrically integral model $S$ over its number field of
constants, and spread them to an abelian scheme and a section.
Shrink $S$ once so that this section avoids the identity section.
For every $n$, the division scheme $[n]^{-1}(P)$ is then finite
\'etale over this same $S$; no further shrinking depends on $n$.
Its rational point over $K'(\mu_\infty)$ extends over
$S_{k(\mu_\infty)}$ by normality. Evaluation at any closed point
of $S$ makes a nonzero abelian-variety point over a number-field
cyclotomic extension divisible by every $n$, contradicting the
preceding paragraph. This supplies exactly the hypothesis in
Stix's Theorem 14 over every finite extension of $K(\mu_\infty)$.
For a positive-genus curve, $[x-y]=0$ implies $x=y$, since a
function with divisor $x-y$ for distinct points has degree one.
The theorem therefore separates all point and cusp
packets. Finite extensions of the cyclotomic field have the form
$K'(\mu_\infty)$; the identical argument, followed by normality,
proves the open-subgroup assertion.

For the more general $B$, let $z\in C_{P_U}(B)$ and use the same
$J=N\overline{\langle z\rangle}$. Descend its finite geometric
cover to some finite extension $K'/K$, and let $E_J\subset E_U$
be its arithmetic open subgroup, with $E_J\cap P_U=J$.
The group $B\cap E_J$ is open in $B$ and its projection is an
open subgroup of $H_K$. Its action on $J^{\rm ab}$ factors through
that projection, because its geometric kernel lies in $J$ and
acts by inner automorphisms. The image of $z$ in the compactified
Jacobian is therefore fixed by an open subgroup of $H_K$, and
vanishes by the already proved arithmetic assertion. The cyclic
and boundary-module arguments apply unchanged and give a unique
inertia containing every centralizing element. The cyclotomic
character of $B$ is trivial, so that entire inertia centralizes $B$.
Finally, for a subgroup $N$ which projects openly to $G_K$ and
normalizes a nontrivial procyclic $A$ through $\chi_K$, set
$B=N\cap E_U^{\circ}$. Its projection is open in $H_K$ and it
centralizes $A$. Hence $A$ is contained in the recognized inertia.
This also supplies the exact comparison with Stix's essential
cyclotomic-normalizer criterion without inserting a nonexistent
normalizer section.
\end{proof}




\begin{statement}{F.2}{Proved}{the necessary half of (F1)}
\label{st:F.2}
\label{F-geometric-necessary}
For an elementary hyperbolic polycurve $U/K$ over a number field and
any morphism $f:V\to U$ from a nonempty connected finite-type
$\mathbb Q$-scheme, every open subgroup $J\subset E_V^{\circ}$
satisfies
\[
 C_{P_U}(f_*(J))=C_{E_U}(f_*(J))=1.
\]
Here $E_V^{\circ}$ is the cyclotomic kernel, and $f_*$ denotes a
continuous representative formed using compatible geometric base
paths. Thus regularity of $V$ is unnecessary for this direction of
scan 48. Both readings of its word ``centralisateur'' give the same
condition.
\end{statement}
\begin{proof}
For a rational point of $U$, induct along its tower of relative
hyperbolic curves. An element centralizing its cyclotomic section
maps to $1$ in the preceding stage, and then belongs to the
geometric fibre group. The fibre group injects and its point action
is the ordinary fibre point action: these are precisely the maps
in Hoshi, Proposition 2.4(i),(ii)~\cite{F-HoshiPoly}, and in CK C5.
Proposition~\ref{F-centralizer} kills the element at the last stage.
This checks the induction in CK C6, including its left exactness.

Choose open normal $N\subset E_V$ with $N\cap E_V^{\circ}\subset J$.
The connected finite \'etale cover $V_N$ has a closed point $v$ over
a number field $k$. Its cyclotomic point homomorphism lands in $J$,
and its composite under $f_*$ is the point section of $f(v)$. The
previous paragraph proves the geometric equality. An arithmetic
centralizer element projects to an element of $G_{\mathbb Q}$
centralizing $G_{k(\mu_\infty)}$. This centralizer is trivial:
replace $k$ by its Galois closure and apply Minamide--Tsujimura,
Theorem 2.12 and Proposition 1.2~\cite{F-MinamideTsujimura} to that
nontrivial closed normal subgroup. Its nontriviality follows, for
example, by choosing a completely split prime of the finite
coefficient field and an odd prime degree ramification not occurring
in its cyclotomic closure; CK B1 gives the complete valuation
argument. Thus the element lies in $P_U$ and is already trivial.
Every change of base paths conjugates the displayed groups and
preserves the asserted equality.
\end{proof}

\begin{statement}{F.3}{Proved}{the point case and the cuspidal exclusion}
\label{st:F.3}
\label{F-point-equivalence}
For $V=\operatorname{Spec}K'$ and a hyperbolic curve $U/K$, condition
(F1)(b) holds exactly for noncuspidal section classes over $K'$, for
each finite extension $K'/K$. The point case of (F1) and (F2) are
therefore the same surjectivity assertion, with the same injective
point map. Every homomorphism through a cusp decomposition group
fails (F1)(b). Both point sections and cusp sections have infinite
index image in $E_U$.
\end{statement}
\begin{proof}
Proposition~\ref{F-centralizer}, including its open-subgroup clause,
identifies (b) with absence from all cusp packets. Removing those
packets from the set asserted to be exhausted by (F2) leaves exactly
the image sought in the point case of (F1). The equality is of sets
and maps, not an appeal to either surjectivity claim. For a map into
$D_c=I_c\rtimes G_K$, its cyclotomic-kernel image acts trivially on
$I_c$, so the latter is a nonzero subgroup of its centralizer.
After choosing a parameter at $c$, a unit $u\in\mathcal O(V)^*$
gives such a map: its compatible Kummer torsors give a continuous
cocycle $a:E_V\to T$, and the homomorphism is
$h\mapsto(a(h),p(h))\in T\rtimes G_K\hookrightarrow E_U$.
A point or cusp section has trivial intersection with $P_U$;
the quotient of $E_U$ by its image has the infinite coset space
$P_U$. Openness therefore removes both kinds of point-source map.
\end{proof}

\begin{statement}{F.4}{Known}{open maps to a curve}
\label{st:F.4}
\label{F-open-Hom}
For a normal geometrically connected variety $V/K$, a hyperbolic
curve $U/K$, and a number field $K$, continuous open homomorphisms
$E_V\to E_U$ over $G_K$, modulo conjugation by $P_U$, are precisely
the homomorphisms of unique dominant $K$-morphisms $V\to U$.
Each such open homomorphism satisfies (F1)(b).
\end{statement}
\begin{proof}
This is the ordinary finite-type specialization of Mochizuki,
Theorem 16.5 and its following remark, as checked in F149-ZCF1
and in the attached source~\cite{F-Mochizuki}. For an open map $f$,
the map $E_V^{\circ}\to f(E_V)\cap E_U^{\circ}$ is onto and open.
The image of every open subgroup therefore contains an open
subgroup of $P_U$. Its geometric centralizer is trivial by the
slimness of a hyperbolic curve group. Projection to $G_K$ and the
arithmetic centralizer calculation in
Proposition~\ref{F-geometric-necessary} give the same conclusion in
$E_U$. The groupoid comparison in Part A sends a morphism with its
base paths to precisely this outer homomorphism. Its degree on
inertia is the ramification index of the reconstructed morphism,
so the peripheral/Kummer data are the data of that same morphism.
\end{proof}

\begin{statement}{F.5}{Refuted}{open image forces dominance in higher dimension}
\label{st:F.5}
\label{F-hyperplane}
The open-image dominance assertion for higher-dimensional elementary
targets is false. For $C:y^2=x^5-x$ over $\mathbb Q$ and $Y=C^3$,
let $H$ be the inverse image in $Y$ of the trilinear divisor
\[
 1+7x_1+3x_2+2x_3+13x_1x_2+11x_1x_3
       +5x_2x_3+19x_1x_2x_3=0
\]
in $(\mathbb P^1)^3$. The closed immersion $H\hookrightarrow Y$
is nondominant and induces $E_H\simeq E_Y$ over $G_{\mathbb Q}$.
\end{statement}
\begin{proof}
We checked the entire proof of F149-ZCF4 and reran its finite
certificate after reading the program. The branch values of
$C\to\mathbb P^1$ are $0,1,-1,i,-i,\infty$. Its eight affine
singularity ideals are the unit ideal, all eighteen bilinear slice
determinants at a branch value are nonzero, and all 216 triples of
branch values give nonzero values. In local ramified coordinates
$x-b=z^2$ these checks respectively cover zero, one or two, and
three ramified coordinates, proving smoothness. The divisor is
ample, as the pullback of $\mathcal O(1,1,1)$ by a finite map.
The geometric Lefschetz theorem gives connectedness and an
isomorphism on geometric fundamental groups; the identity on
$G_{\mathbb Q}$ gives the arithmetic isomorphism. Finally
$\dim H=2<3=\dim Y$, which proves the failure of dominance through
the displayed actual morphism. The curve-target replacement is
Proposition~\ref{F-open-Hom}; no claim that it exhausts all
higher-dimensional replacements is made.
\end{proof}

\begin{statement}{F.6}{Open}{the remaining (F1) and field-reduction steps}
\label{st:F.6}
\label{F-open-main}
The unrestricted higher-dimensional sufficiency assertion in
F(b)(6) and the full reduction in F(c) from number fields remain
unproved. Proposition~\ref{F-local-dichotomy} proves the
infinite-index assertion in F(b)(5), using the local theorem
supplied in the prompt. Proposition~\ref{F-full-curve-reduction}
proves the exact universal curve-target (F1)--(F2) comparison.
The following records the obstacle in the earlier purely global
approach, now bypassed by the supplied local theorem. Its missing
group-theoretic step was this: for a continuous
$G_K$-map $f:E_V\to E_U$ satisfying (b), prove that the finitely
generated closed group $f(P_V)$ is either $1$ or open in $P_U$.
The elasticity theorem applies when this image is normal in an
\emph{open geometric subgroup}; normality in $f(E_V)$ does not
provide that hypothesis. Passing to the arithmetic normalizer leaves
precisely that unresolved input. The open Hom theorem cannot be
applied before openness has been supplied. The hyperplane example
above also prevents replacing this input by an unrestricted
higher-dimensional dominance assertion.

We checked Bresciani's published field-reduction paper
\cite{F-BrescianiFiniteType}: its proved range includes genus at most
two and curves with a suitable geometric finite cover mapping to a
constant abelian variety. Its general reduction retains an additional
Bombieri--Lang input. We have not proved that input or removed the
cover hypothesis. Spreading a section gives fibre section classes;
the unproved step in the unrestricted case is constructing a generic
rational point from that system. Accordingly the computations and
centralizer identities above supply no result labelled a proof of
this step. The first textual assertion in the 1983 notes not supplied
by our curve comparison is (F1) for arbitrary positive-dimensional
sources and elementary higher-dimensional targets, scan 48. This
identifies the scope not covered by that comparison; it is not a
claim that the entire 1981 manuscript lacks all higher-dimensional
ideas.
\end{statement}

\subsection{F(d): fixed points and descent from the cyclotomic kernel}

\begin{statement}{F.7}{Proved}{fixed pro-points and sections}
\label{st:F.7}
\label{F-fixed-dictionary}
The completed geometric universal cover of a hyperbolic $U/K$
has at most one fixed geometric pro-point under a section restricted
to $H_K$. For a section over $G_K$ its $H_K$-fixed and $G_K$-fixed
sets coincide. An interior fixed pro-point gives precisely its
rational-point section, and a boundary fixed pro-point gives precisely
a section in its cusp packet. Thus the existence assertions in (F2)
and (F3), Conjecture B with its $H_K$ strengthening, have identical
content under these comparison maps.
\end{statement}
\begin{proof}
The fibre of the universal cover over an interior geometric point
is a simply transitive $P_U$-set. Its arithmetic stabilizer maps
isomorphically to the point's Galois group. At a boundary pro-point
the stabilizer is $D_c=N_{E_U}(I_c)$. These two stabilizer
computations prove both maps of the dictionary. For uniqueness,
two pro-points differ at a finite covering level. Refine its
subgroup to a characteristic open subgroup with positive-genus
compactification. Such a cover exists: in genus zero pull back a
Fermat cubic along three distinct punctures; in positive genus use
the original curve. A further Galois and characteristic refinement
preserves positive genus by Riemann--Hurwitz. The given section
descends the cover over the cyclotomic field. The two fixed points
would give distinct rational points in its compactification with
the same section. The Jacobian argument in
Proposition~\ref{F-centralizer} excludes this. Normality of $H_K$
then makes a nonempty singleton fixed by all of $G_K$. This checks
CK D0a--D0b and uses no fixed-point existence assumption.
\end{proof}

\begin{statement}{F.8}{Refuted}{the discrete finite-group assertion}
\label{st:F.8}
\label{F-discrete}
The assertion in F(d) for \emph{every} finite group is false. The
identity splitting for $G=1$ fixes the entire universal cover.
For a nontrivial finite faithful group of orientation-preserving
hyperbolic isometries of the disk, the common fixed set is exactly
one interior point and contains no ideal boundary point.
Orientation-reversing lifts can instead fix a geodesic.
\end{statement}
\begin{proof}
For $G=1$ the extension is the surface group itself and the unique
map $1\to\pi_1(U)$ is a splitting; its action is the identity on
every point. For the repaired assertion choose any finite orbit
in the disk and minimize the sum of its squared hyperbolic
distances. Strict convexity and properness give a unique minimum,
fixed by the group. Conjugate it to the centre. Every
orientation-preserving isometry fixing the centre is a rotation;
a nonidentity rotation fixes no other disk point and no point on
the ideal circle. This proves the unique fixed-point statement
for the specified action and also proves that such a finite group
is cyclic. For an orientation-reversing example, complex
conjugation $z\mapsto\overline z$ of the disk fixes its real
diameter. Take the real tripod with real point $2$; the lift of
complex conjugation fixing a chosen lift of $2$ is exactly a
reflection of its uniformizing disk and gives the corresponding
splitting by $C_2$. These examples specify the actual actions and
comparison maps. An arbitrary outer action requires a realization
before the isometry argument applies; finite-order and
orientation hypotheses cannot be omitted from that argument.
\end{proof}

\begin{statement}{F.9}{Proved}{continuous B-prime actions acquire a quotient}
\label{st:F.9}
\label{F-continuous-action}
Let $\mathcal U$ be the geometric universal pro-cover of a hyperbolic
curve over $\overline K$. A semilinear $G_K$-action on $\mathcal U$
for which each function has open stabilizer normalizes the deck
group of a finite hyperbolic quotient defined over $K$. It is the
action of a continuous section for that quotient. The quotient can
be chosen to cover any specified initial level.
\end{statement}
\begin{proof}
We checked F149-ZCF12, including descent of the open curve. Let
$\Omega$ be the union of covering function fields and let
$\overline K(f_1,\ldots,f_m)$ be an initial level. The orbits of
the $f_i$ are finite. The field $L$ they generate together is
therefore a finite extension of the initial field inside $\Omega$,
invariant under $G_K$. Its curve is a finite \'etale cover of the
initial level, and $\operatorname{Aut}_L(\Omega)$ is its deck group.
Normalization of this deck group descends every action map to
the quotient curve: factor it through a common finite Galois
cover, then descend the resulting deck-invariant map. Thus the
action preserves the actual open curve and its smooth proper
compactification with boundary. Choose a finite coefficient field
and intersect the stabilizers of finitely many defining generators;
its normal core is an open normal subgroup acting only on scalars.
The remaining finite semilinear action is an effective descent
datum. One explicit effectiveness proof embeds the marked pair
by $\omega(D)^{\otimes3}$ and descends its section space and
homogeneous ideal. The descended quotient is defined over $K$,
and the action on $\Omega$ fixes its function field. It is a
section of its arithmetic group, continuous by the stipulated
open stabilizers. Conversely every continuous section has those
stabilizers. Without this continuity requirement the finite-orbit
argument is unavailable. Proposition~\ref{F-discontinuous-action}
exhibits an abstract action for which no continuously descended
finite \'etale quotient exists, while its own fixed point still
exists. The universal abstract fixed-point assertion remains open.
\end{proof}

\begin{statement}{F.10}{Proved}{the cyclotomic extension never splits virtually}
\label{st:F.10}
\label{F-no-cyclo-split}
For every number field $K$ and every open subgroup $Q'\subset Q_K$,
the extension
\[
 1\longrightarrow H_K\longrightarrow\chi_K^{-1}(Q')
       \longrightarrow Q'\longrightarrow1
\]
has no splitting, even as an abstract group extension.
\end{statement}
\begin{proof}
The subgroup $Q'$ is open in $\widehat{\mathbb Z}^{\times}$, so it
contains a congruence kernel modulo some positive integer $m$.
Choose a prime $q\nmid m$ with $q\equiv1\pmod3$. There are such
primes outside every finite set: for a multiple $a$ of that set
and of $3$, any prime divisor $q$ of $a^2+a+1$ is outside the set,
is not $3$, and the order of $a$ modulo $q$ is $3$.
The factor $\mathbb Z_q^{\times}$ contains an element of order
$3$, obtained by the simple-root lift of an element of order $3$
in $\mathbb F_q^{\times}$. Put this in the $q$-coordinate and put
$1$ in every other coordinate. The resulting element of $Q'$
has order $3$. A splitting would embed it into $G_K$. The
Artin--Schreier theorem on finite subgroups of an absolute Galois
group in characteristic zero gives only the trivial group and
groups of order two. In particular $G_K$ has no element of order
three, a contradiction. This is the torsion argument requested
in the addition to F(d), with the open-subgroup quantifier included.
\end{proof}

\begin{statement}{F.11}{Proved}{the exact obstruction to extending a cyclotomic section}
\label{st:F.11}
\label{F-descent-obstruction}
Let $1\to P\to E\xrightarrow{p}G_K\to1$ be a profinite
extension and $t:H_K\to E$ a continuous section over $H_K$.
Suppose its $P$-conjugacy class is $Q_K$-invariant. Put
\[
 C=C_P(t(H_K)),\qquad N=N_E(t(H_K)),\qquad W=N/t(H_K).
\]
There is an exact sequence
\[
 1\longrightarrow C\longrightarrow W\longrightarrow Q_K
       \longrightarrow1.                                      \tag{F.1}
\]
Extensions of the actual section $t$ to $G_K$ correspond bijectively
to continuous splittings of (F.1); conjugation by $C$ is respected.
In particular, $C=1$ gives a unique extension. For a cusp of a
hyperbolic curve, a choice of tangential origin identifies $C$ with
$T$, and the obstruction is the transgression
\[
 [o_t]\in H^2_{\rm cts}(Q_K,T),\qquad
 o_t(q,r)=-c\bigl(u(q)u(r)u(qr)^{-1}\bigr),                    \tag{F.2}
\]
where $t(h)=(c(h),h)$ in $T\rtimes G_K$, and $u$ is a continuous
set-theoretic transversal for $G_K\to Q_K$. Its vanishing is
necessary and sufficient; the extension classes form a torsor
under $H^1_{\rm cts}(Q_K,T)$ when nonempty.
\end{statement}
\begin{proof}
For each $g\in G_K$, invariance supplies a lift in $E$ which
normalizes $t(H_K)$, by multiplying any lift by a suitable element
of $P$. Thus $N\to G_K$ is onto. Its kernel is $C$: an element of
$P$ normalizing $t(H_K)$ must centralize each $t(h)$, since
projection to $H_K$ distinguishes its elements. The subgroup
$t(H_K)$ is normal in $N$ and meets $C$ trivially. Quotienting
gives (F.1), as an exact sequence of compact groups.

A section extending $t$ lies in $N$ and descends to a splitting
of (F.1). Conversely, the inverse image in $N$ of the graph of a
splitting maps isomorphically to $G_K$: each fibre contains one
element, since its possible discrepancy lies in both $C$ and
$t(H_K)$. The compact-to-Hausdorff bijection is a homeomorphism.
Its inverse is the required section. This proves the criterion
without assuming a splitting of $G_K\to Q_K$.

For the cusp, the normalizer formula identifies $N$ with $D_c$;
invariance of $t$ is the equivariance of the homomorphism
$c:H_K\to T$. The multiplication of representatives
$(0,u(q))$ in $D_c/t(H_K)$ gives the error
$(0,u(q)u(r)u(qr)^{-1})$, whose image in the kernel is the negative
of $c$ of that element. This proves (F.2), including its sign.
Associativity proves the cocycle identity, and changing $u$ or
the representatives adds a continuous coboundary. The groups here
are countably based; a quotient of a countably based profinite
group by a closed subgroup admits a continuous transversal. To
construct one, choose successively finite clopen partitions of
the quotient subordinate to the images of cosets at level $n$,
refine the choices compatibly, and take the unique nested coset
limit. The image of every coset is open, which supplies each
finite refinement. Hence (F.2) is a continuous cocycle. Altering
the representatives by a continuous $1$-cochain kills it exactly
when the resulting representatives multiply as a section.
Two such choices differ by a $1$-cocycle, and $T$-conjugation
changes it by a coboundary. This proves the torsor assertion.
\end{proof}

\begin{statement}{F.12}{Proved}{an invariant cuspidal class with nonzero descent obstruction}
\label{st:F.12}
\label{F-one-plus-i}
Let $F=\mathbb Q(\mu_\infty)$ and let $U=\mathbb P^1\setminus
\{0,1,\infty\}$. In the packet at $0$, the tangential class of
parameter $1+i$ over $F$ is invariant under
$Q=\operatorname{Gal}(F/\mathbb Q)$ and does not extend to a section
over $\mathbb Q$. Its class (F.2) is nonzero of order exactly two.
\end{statement}
\begin{proof}
Every root of unity has roots of every order in $F$, so its Kummer
class in $\widehat{F^*}$ is zero. The two conjugates of $1+i$
have ratio $-i$, proving invariance. Let $c=\kappa_F(1+i)$.
The exact identity $(1+i)^2=2i$ gives
$2c=\operatorname{res}\kappa_{\mathbb Q}(2)$, and therefore kills
its transgression by two.

We compute the kernel of the restriction map, rather than infer
nonextension from a finite approximation. Inflation--restriction
identifies it with $H^1(Q,T)$. In fact
\[
 H^1(Q,T)=\mu(\mathbb Q)=\{1,-1\}.                            \tag{F.3}
\]
For completeness, apply the sequence
$0\to\mathbb Z_\ell(1)\to\mathbb Q_\ell(1)\to
\mathbb Q_\ell/\mathbb Z_\ell(1)\to0$.
The invariants of $\mathbb Q_\ell(1)$ vanish. Its first cohomology
also vanishes: for a continuous cocycle $b$ and a central $q$ with
$\chi_\ell(q)\ne1$, commutativity gives
$(\chi_\ell(q)-1)b(r)=(\chi_\ell(r)-1)b(q)$, explicitly a
coboundary after division in $\mathbb Q_\ell$. Thus the first
cohomology of $\mathbb Z_\ell(1)$ is the invariant torsion in
$\mathbb Q_\ell/\mathbb Z_\ell(1)$. The full cyclotomic image
fixes no such torsion for odd $\ell$, and fixes exactly its order-two
subgroup for $\ell=2$. Taking the product proves (F.3), whose
nontrivial class is the Kummer class of $-1$.

An extension of the class $c$ would have to be cuspidal at $0$:
it normalizes the nonzero centralizer $I_0$ of its restriction.
Its Kummer class $a\in\widehat{\mathbb Q^*}$ would satisfy
$\operatorname{res}(2a-\kappa_{\mathbb Q}(2))=0$.
By (F.3) the difference is a sign. The completed valuation at $2$
would then give $2v_2(a)=1$ in $\widehat{\mathbb Z}$, impossible
already modulo two. Consequently the obstruction is nonzero;
its order is two by the first paragraph. The attached exact
program verifies the two algebraic identities and the valuation
congruence obstruction. Its $D_8$ computation records the familiar
finite non-split central-extension phenomenon, but the proof of
nonextension above is the full profinite valuation argument.
\end{proof}

\subsection{F(e): function fields and the lifting obstruction}

\begin{statement}{F.13}{Proved}{birational centralizers and the written formulas}
\label{st:F.13}
\label{F-birational}
Let $X/K$ be a smooth proper geometrically connected curve over a
finitely generated characteristic-zero field and let $L=K(X)$.
A section $s:G_K\to G_L$ has nontrivial centralizer in
$G_{\overline K(X)}$ of $s(H_K)$ precisely when it lies in the
decomposition group of a rational place. There the centralizer is
$I_v\simeq T$, both geometrically and arithmetically. Formulas
(32),(33) on scan 127 are false with their written $L$; for the
point section of a hyperbolic model $U$ containing that place,
their $U$ versions are true.
\end{statement}
\begin{proof}
We checked CK B7--B8. The maps are the continuous surjections
$q_U:G_L\to\pi_1(U)$, as $U$ runs through the nonempty opens of
$X$; their inverse limit is $G_L$, since every finite separable
extension is unramified over some such open. A nonidentity
centralizer element survives in one hyperbolic open and every
smaller open. Proposition~\ref{F-centralizer} gives its unique
fixed cusp at each of these levels. Those cusps are compatible:
a cusp mapping to an interior point would kill its nonzero
inertia image. They determine compatible points on all finite
normalizations of $X$, hence a valuation of $\overline L$, fixed
by $s(G_K)$. Its point on $X$ is $K$-rational. Conversely a
decomposition section centralizes $I_v$ on $H_K$. The local
maps identify this same $I_v$ with its inertia image on every
open deleting the place; inverse limits give equality of the
centralizers. Projection to $G_K$ kills the arithmetic centralizer
by the arithmetic calculation used above. A model containing the
place kills $I_v$ and gives a point section, whose centralizer is
zero by Proposition~\ref{F-centralizer}.

Explicitly take $L=\mathbb Q(t)$, the $t$-adic place $v$, and
$U=\mathbb P^1\setminus\{1,2,\infty\}$. Fixing compatible roots
of $t$ gives a decomposition section. Its centralizer contains the
inertia detected by $t^{1/2}$ and equals $T$; its image on $U$ is
the section of $0$. The comparison is exactly $q_U$, which kills
that ramification. Visual inspection of scan 127 supports the
subscript $L$ in both formulas. The mathematical counterexample
does not authorize silently changing that reading.
\end{proof}

\begin{statement}{F.14}{Known}{universal cuspidalization}
\label{st:F.14}
\label{F-cuspidalization}
For a hyperbolic curve $X/k$, $k/\mathbb Q$ finitely generated,
the assertion that all sections after every finitely generated
extension of $k$ are geometric or cuspidal is equivalent to the
assertion that all those sections lift through the function-field
Galois group. The references arXiv:2108.13397 and Inventiones
235 (2024), 129--150 name the same paper.
\end{statement}
\begin{proof}
We opened the publisher's version of Bresciani and checked
Theorems A and B~\cite{F-BrescianiBirational}. For a section $s$
over a finitely generated $K/k$, its base change to $K(t)$ is
the actual map
\[
 g\longmapsto(s(\operatorname{res}g),g)
 \quad\hbox{in}\quad
 \pi_1(X_K)\times_{G_K}G_{K(t)}=\pi_1(X_{K(t)}).
\]
The universal lifting assertion applies to this section. Theorem A
then gives geometricity or cuspidality of $s$. Conversely a point
section lifts through its rational-place decomposition sequence,
and a cusp section already specifies such a local splitting.
Geometric conjugators lift through the surjective generic-point
map. This proves the equivalence on classes and on chosen
representatives. The additional $K(t)$ base change is present in
Theorem A and is indispensable in this reduction; the theorem
does not establish the same conclusion from ordinary birational
liftability alone.
\end{proof}

\begin{statement}{F.15}{Open}{the remaining valuative and higher-dimensional lifting claims}
\label{st:F.15}
\label{F-open-valuations}
For the tripod over $\mathbb Q$, the exact lifting problem still
left by our work is the splitting of
\[
 1\longrightarrow\ker(G_{\mathbb Q(t)}\to E_U)
 \longrightarrow G_{\mathbb Q(t)}\times_{E_U,s}G_{\mathbb Q}
 \longrightarrow G_{\mathbb Q}\longrightarrow1
\]
for an arbitrary section $s$. All geometric and cuspidal sections
have such a lift by the local construction, but the inertia
description of its kernel gives no complement for an arbitrary
$s$. The centralizer test in Proposition~\ref{F-birational}
recognizes a decomposition group once nonzero centralizer has
been supplied; it does not supply that nonzero element.

The full-rank structure and the exact normalizer are established
in Proposition~\ref{F-valuation-structure}. The finite-residue
qualification needed to produce germs, and its failure for an
infinite algebraic residue extension, are established in
Proposition~\ref{F-valuation-germs}. Restricting a decomposition
section along $K\subset L_1\subset L$ sends its valuation to
$V\cap L_1$ and proves the forward compatibility on scans
125--126. Reversing it would require a single compatible invariant
valuation above all the one-variable restrictions, and no such
construction is supplied. In particular, selecting a place for
each subfield does not itself show that the resulting valuation
rings are compatible under inclusion.

For F(e)(4), a section of a polycurve gives a section on its base.
A birational lift of that base section has domain $G_K$, whereas
the curve-fibre lifting assertion over $K(B)$ has domain
$G_{K(B)}$. The morphisms displayed here do not extend that
domain. Applying the latter assertion without this extension
would leave precisely the missing step identified in the prompt.
Proposition~\ref{F-proper-polycurve-lifting} supplies the
implication for towers proper except possibly at their last stage,
using actual point fibres and full-flag valuations. The case with
a cuspidal projected section on an affine intermediate base remains
unproved in this return.
\end{statement}

\subsection{F(f): audit of the potential-good-reduction criterion}

\begin{statement}{F.16}{Known}{condition (e)}
\label{st:F.16}
\label{F-reduction}
Let $R$ be a henselian discrete valuation ring whose residue field
$k$ is finitely generated over $\mathbb Q$, and let $U/K$ be a
hyperbolic curve over its fraction field. The residue action on
its full geometric profinite group obtained from any splitting
$G_k\to G_K$ is realized by a smooth hyperbolic curve over $k$
precisely when $U$ has potentially good marked reduction. Allowing
a finite residue extension gives the same condition. This is the
corrected condition (e) of cote 149, scans 100--103.
\end{statement}
\begin{proof}
We checked the complete proofs of F149-F17.6 and F17.7. The
essential comparison is between the full group of the generic
curve and the logarithmic special fibre, rather than just their
homology. Illusie, Theorem 7.6~\cite{F-Illusie}, identifies covers
of the generic interior with finite Kummer \'etale covers of the
log-regular model; all boundary ramification is tame in residue
characteristic zero. Hoshi, Corollary 3.9~\cite{F-HoshiLog},
identifies the latter covers with those of the strict log special
fibre after completion. Completion induces the same residue
quotient and the same inertia $T$, hence the same absolute Galois
group here. The curve geometric groups also remain unchanged in
characteristic zero. These are compatible isomorphisms over the
base log Galois group.

At a split node of thickness $m$, the monoid chart is
$X+Y=mT_0$. Its arithmetic local group has geometric torus
\[
 \{(a,b,c)\in T^3:a+b=mc\},
\]
and projection to base inertia is $(a,b,c)\mapsto c$. The kernel
is $(a,-a,0)$, an injected node inertia. Pulling back along any
specified residue splitting preserves surjectivity to $G_k$ and
the cyclotomic action on this kernel. Mochizuki, Example 2.5 and
Proposition 1.2~\cite{F-MochizukiCombinatorial}, together with the
node injection in Hoshi--Mochizuki Lemma 2.11~\cite{F-HoshiMochizukiNodes},
identifies these as distinct maximal procyclic node groups.
Thus every node contributes a cyclotomically normalized class
in addition to the cusp classes.

We verify the numerical exclusion used to compare with an
arbitrary smooth realizing curve. After a finite residue extension
fixing the dual graph, a cycle in that graph gives a nonzero
trivial quotient of $H_1(U_{\overline K},\mathbb Q_\ell)$.
A smooth curve has none: its toric homology has character
$\chi_\ell$, while its Jacobian Tate module and its dual have
zero cyclotomic-kernel invariants by the arithmetic argument of
Proposition~\ref{F-centralizer}. For a tree with $r>0$ cusps,
eliminating end vertices from the boundary relations gives
\[
 0\to\mathbb Q_\ell(1)^{r-1}\to H_1(U_{\overline K},\mathbb Q_\ell)
 \to\bigoplus_v V_\ell\operatorname{Jac}(C_v)\to0.
\]
The dimension of maps from $\mathbb Q_\ell(1)$ is exactly $r-1$.
A smooth realizing curve would therefore have exactly $r$ cusps,
but Stix's cyclotomic normalizer criterion gives only $r$ such
classes, contradicting the $r+n$ classes from $n>0$ nodes. For
a proper curve, its nonzero $H^2(-,\mathbb Z/\ell)$ first forces
the realizing curve to be proper, and a single node class already
contradicts the same criterion. This proves nonrealization after
every finite residue extension for stable bad reduction.

For potentially good reduction choose a finite Galois extension
$K_1/K$ giving a smooth stable marked model. Uniqueness of that
model extends the Galois action to it and gives descent maps on
its special pair over $k_1$. Compose the specified splitting with
$G_K\to\operatorname{Gal}(K_1/K)$ and use these maps to twist
the pair. After a finite extension of $k$ this is an ordinary
finite Galois descent datum, effective by the marked
$\omega(D)^{\otimes3}$ embedding. Its descended curve over $k$
has the required action: the generic action factors through
\[
 \operatorname{Gal}(K_1/K)
       \times_{\operatorname{Gal}(k_1/k)}G_k,
\]
and restriction along the splitting is exactly the descent action.
Conversely, choose stable reduction after a finite field
extension, intersect the splitting's inverse image with its open
Galois subgroup, and apply the stable-bad exclusion above. Nodes
survive further stable base change. Hence realization implies a
smooth stable special fibre. This proves both directions with the
specified full-profinite comparison, and verifies the claim being
used from F149 without changing its quantifiers.
\end{proof}

\subsection{F(g): geometric tangents, local conditions, and ideles}

\begin{statement}{F.17}{Proved}{the two torsors and their coordinate ambiguity}
\label{st:F.17}
\label{F-torsors}
At a rational cusp over a finitely generated field $K/\mathbb Q$, the packet is a torsor
under $\widehat{K^*}=H^1(K,T)$, and the geometric tangent vectors
give its injective $K^*$-subtorsor. For a smooth connected separated
finite-type $\mathbb Q$-scheme $S$ and an injected normal-line
extension as in scan 85, its class is the Kummer first Chern class
of the normal line $\mathcal N$; it splits precisely when
$\mathcal N$ is trivial. In that event the geometric frame torsor
under $\mathcal O(S)^*$ embeds in the splitting torsor under
$H^1(S,T)=\widehat{\mathcal O(S)^*}$.
\end{statement}
\begin{proof}
For a local parameter $z$, compatible roots of $z$ give an origin
in the packet. Fix the convention
$\zeta:z^{1/n}\mapsto\zeta_nz^{1/n}$ and
$\kappa(a)(\sigma)=\sigma(a^{1/n})/a^{1/n}$.
The section fixing roots of $az$ differs from the section fixing
roots of $z$ by $-\kappa(a)$. A tangent vector $v$ with
$z(v)=a$ instead has coordinate $+\kappa(a)$. These explicit
inverse maps identify the conormal-frame and tangent conventions.
The Kummer kernel in
a finitely generated characteristic-zero field is trivial:
divisibility by all integers kills every divisorial valuation,
and then reduces to its number-field constant units, where it
again forces the identity. This proves injectivity.

For $S$ choose local normal frames $e_j=g_{ij}e_i$. Lifts of these
transition functions through their compatible roots have the
triple-overlap defect defining $c_1(\mathcal N)$; these are exactly
the composition defects of the local normal-tube splittings.
This is the sign convention of F149-C6. The conormal line has
the negative class. A splitting kills this class. The Kummer
sequence injects $\operatorname{Pic}(S)/n$ into
$H^2(S,\mu_n)$, so $\mathcal N$ is divisible by every $n$ in
$\operatorname{Pic}(S)$. Here $\operatorname{Pic}(S)$ is finitely
generated: choose a smooth proper compactification after
resolution, use finite generation of its rational Picard group,
and quotient by the finitely many boundary divisors. Thus
$\mathcal N$ is trivial. The converse is given by a global frame.
Taking inverse limits in the degree-one Kummer sequence gives
$H^1(S,T)=\widehat{\mathcal O(S)^*}$ because
$T\operatorname{Pic}(S)=0$; the finite root-of-unity invariants
give no derived-limit obstruction in degree one. Embedding
$\mathcal O(S)^*$ into the function field gives its injectivity
by the same Kummer-kernel argument. These are the checked steps
of F149-C6--C8, with left injection explicitly retained as part
of the normal-line extension's data.
\end{proof}

\begin{statement}{F.18}{Proved}{the abstract local extension does not select its geometric subtorsor}
\label{st:F.18}
\label{F-twist}
Automorphisms of $T\rtimes G_K$ inducing the identity on both
$T$ and $G_K$ act transitively on its section classes. Within
$E_U$, the cusp packet embeds in the full section set, so a
chosen parameter gives a uniquely determined completed Kummer
coordinate for each cuspidal class. Without that choice its
coordinate is defined up to translation by $K^*$.
\end{statement}
\begin{proof}
For a continuous cocycle $b:G_K\to T$ the formula
$(a,g)\mapsto(a+b(g),g)$ is a continuous group automorphism;
the cocycle identity is exactly multiplicativity. It carries the
zero section to the class of $b$, proving transitivity. Since
$K^*$ is a proper subset of $\widehat{K^*}$, that automorphism
group cannot preserve the proposed proper subtorsor. The packet
injection inside $E_U$ follows directly from
Proposition~\ref{F-centralizer}: a geometric conjugator between
two sections at the same lifted cusp must normalize $I_c$,
hence lies in $I_c$. Therefore precisely the local coboundary
ambiguity remains. Replacing the chosen parameter multiplies it
by a nonzero scalar in its first-order conormal line and shifts
the coordinate by that scalar's Kummer class.
\end{proof}

\begin{statement}{F.19}{Proved}{all everywhere locally geometric global Kummer classes}
\label{st:F.19}
\label{F-local-global-units}
For a number field $K$ put $E=\mathcal O_K^*$ and define
$\mathcal L_K\subset\widehat{K^*}$ by the two requirements:
its image in $\widehat{K_v^*}$ belongs to $K_v^*$ at every finite
place $v$, and it belongs to $\mathcal O_v^*$ at almost every
finite place. Then
\[
 \mathcal L_K=K^*\,\widehat E,\qquad
 \mathcal L_K/K^*\simeq\widehat E/E.                          \tag{F.4}
\]
For $K=\mathbb Q$ this is exactly $\mathcal L_{\mathbb Q}=\mathbb Q^*$.
For every infinite-order unit $\varepsilon$ and
$\alpha\in\widehat{\mathbb Z}\setminus\mathbb Z$, the class
$\varepsilon^\alpha$ is in $\mathcal L_K\setminus K^*$.
\end{statement}
\begin{proof}
At a finite place, a uniformizer gives
$\widehat{K_v^*}=\widehat{\mathbb Z}\times\mathcal O_v^*$.
The unit factor is already complete because it is profinite
abelian: its map to the inverse limit of power quotients is
injective by its finite quotients and surjective by compactness.
Thus membership in $K_v^*$ means exactly that the valuation is
an ordinary integer; the almost-all-unit condition means these
integers have finite support.

Let $D=\bigoplus_v\mathbb Z[v]$ be the finite divisor group and
$P\subset D$ its principal divisors. The valuation sequence
$1\to E\to K^*\to P\to0$ splits as an abstract abelian sequence,
since $P$ is free abelian. Consequently
$\widehat{K^*}=\widehat E\times\widehat P$ after such a choice.
The inclusion $P\to D$ induces an injection
$\widehat P\to\widehat D$. Indeed $D/P=\operatorname{Cl}(K)$
has a finite exponent $h$, and
$(hn)D\cap P\subset nP$; equivalently the induced power
topologies are cofinal. The same conclusion follows by choosing
a basis for the finite-index subgroup $P$ of the free abelian
group $D$. An integral finite divisor in the image of
$\widehat P$ is principal: its image in the finite class group
is zero, because that class-group map factors through $D/hD$.
Hence a class $a\in\mathcal L_K$ has the divisor of some
$b\in K^*$. Its quotient by $b$ lies in $\widehat E$, proving
the first equality. Conversely the image of $\widehat E$ at
every finite place lies in the compact group $\mathcal O_v^*$,
so every element of $K^*\widehat E$ has the required properties.
The intersection $K^*\cap\widehat E$ is exactly $E$, giving
the quotient statement.

The finite generation $E\simeq\mu(K)\times\mathbb Z^r$
makes the last assertion exact. An infinite-order unit has a
nonzero vector $m\in\mathbb Z^r$. The vector $\alpha m$
cannot be integral unless $\alpha\in\mathbb Z$: a nonzero
coordinate gives $m_i\alpha=n_i$, and an element of
$\widehat{\mathbb Z}$ equal to a rational number is an integer.
This proves nonglobality while the first equality proves all
local conditions. For $\mathbb Q$, $E=\{1,-1\}$ is already
complete, so no extra classes occur.
\end{proof}

\begin{statement}{F.20}{Refuted}{powers of one unit do not exhaust the defect}
\label{st:F.20}
\label{F-one-unit-false}
The added assertion in F(g)(3 bis) that all nonglobal everywhere
locally geometric classes are exactly the $\varepsilon^\alpha$
for one fixed infinite-order unit is false. The exact complete
replacement is (F.4). The associated idele classes of such locally
geometric classes, with matching real signs, belong to the
connected component of the idele class group.
\end{statement}
\begin{proof}
In $K=\mathbb Q(\sqrt2,\sqrt3)$ take
$\varepsilon=1+\sqrt2$, $\eta=2+\sqrt3$, and let
$e_2\in\widehat{\mathbb Z}$ have $2$-component $1$ and every
other prime component zero. The class $\eta^{e_2}$ is locally
geometric and nonglobal by Proposition~\ref{F-local-global-units}.
It cannot be $\varepsilon^\alpha$: the norm to
$\mathbb Q(\sqrt3)$ takes $\varepsilon$ to $-1$ and $\eta$
to $\eta^2$. Thus such an equality would make
$\eta^{2e_2}$ torsion in the completed unit group, whereas its
nonzero $2$-adic unit-lattice vector makes it infinite order.
The exact unit norms are checked by the attached program.
Even in unit rank one, the actual set in (F.4) also contains
arbitrary $K^*$ multiples; its quotient is the one described
there, with all unit directions retained.

To check the connected-component assertion without assuming a
chosen reciprocity convention, approximate a completed unit by
ordinary units simultaneously in every finite-place compact unit
group and in its finite real-sign group. The corresponding idele
classes are represented by archimedean positive components alone,
after division by the principal unit idele. Those classes lie in
the image of the connected identity component of the archimedean
idele group. Passing to the limit remains in its closed connected
component. Multiplication by a principal $K^*$ element does not
change the idele class. Different positive archimedean magnitudes
also remain in that same component. This proves exactly the
asserted membership; it does not assign an unchosen canonical
archimedean representative to a completed Kummer class.
\end{proof}

\begin{statement}{F.21}{Proved}{rational idelic gluing}
\label{st:F.21}
\label{F-ideles}
There is a canonical decomposition, using the real sign,
\[
 \varprojlim_n\mathbb A_{\mathbb Q}^{\times}/
                   (\mathbb A_{\mathbb Q}^{\times})^n
 \simeq\widehat{\mathbb Q^*}\times\widehat{\mathbb Z}^{\times}.
                                                                  \tag{F.5}
\]
A family of local completed Kummer classes belongs to its left
side precisely when, for every $n$, only finitely many finite
places have nonzero valuation modulo $n$. For such a family
$(b_v)$, let $a\in\widehat{\mathbb Q^*}$ have its real sign
and its completed valuations. Its unique second coordinate is
\[
 u_p=b_p\,\operatorname{res}_p(a)^{-1}\in\mathbb Z_p^*,
 \qquad u=(u_p)_p\in\widehat{\mathbb Z}^{\times}.             \tag{F.6}
\]
For local cuspidal sections at one fixed rational cusp with a
common rational tangent origin, $u=1$ is necessary and sufficient
for the family to be the restrictions of one global cuspidal
section.
\end{statement}
\begin{proof}
For an actual idele $b$, put
$q=\operatorname{sgn}(b_\infty)\prod_p p^{v_p(b_p)}$.
There are finitely many nonzero valuations. Then
$b=q\,u\,r$ uniquely, with $u\in\prod_p\mathbb Z_p^*$
and $r\in\mathbb R_{>0}$. This gives the topological product
$\mathbb Q^*\times\widehat{\mathbb Z}^{\times}\times
\mathbb R_{>0}$. The last factor has all roots, and the compact
middle factor is complete for power quotients by the compactness
argument in Proposition~\ref{F-local-global-units}. Taking
inverse limits proves (F.5).

At each $n$, an idele class modulo $n$ has finitely supported
valuations modulo $n$, and every such family of local classes
has an idele representative: at the other places choose unit
representatives. A family locally an $n$th power is globally an
idele $n$th power by choosing roots at all places, which are
units almost everywhere. Thus these local maps are injective
and their stated images are exact. Passing to the inverse limit
proves the criterion and (F.6). Equivalently the positive prime
exponents of $a$ form a family in $\widehat{\mathbb Z}$ tending
to zero modulo every $n$, so their formal product is an element
of $\widehat{\mathbb Q^*}$. The remaining local valuations are
zero, giving the indicated unit factors.

The cusp packet maps are the Kummer restriction maps after the
common tangent origin is chosen. A global class therefore has
the coordinate $(a,1)$, and every such coordinate defines its
global cuspidal section. The packet injection in
Proposition~\ref{F-twist} proves necessity and uniqueness also
after passage to full arithmetic section classes. Changing the
common tangent origin by $q\in\mathbb Q^*$ changes $a$ by $q$
and leaves $u$ unchanged. Identifying this unit group with
$\operatorname{Gal}(\mathbb Q(\mu_\infty)/\mathbb Q)$ uses
the cyclotomic isomorphism; an Artin-map convention can invert
that isomorphism without affecting the obstruction $u=1$.
\end{proof}

\begin{statement}{F.22}{Refuted}{arbitrary compatible local families have the idelic decomposition}
\label{st:F.22}
\label{F-ideles-unrestricted}
Compatibility of local Kummer classes modulo every $n$, by itself,
does not imply the decomposition in (F.5). Omitting the real
place also removes uniqueness of the two factors.
\end{statement}
\begin{proof}
Take at each finite prime $p$ the actual local parameter $b_p=p$,
and real sign $+1$. These give compatible local section classes
at the rational cusp $0$ for every modulus. Their valuations are
$1$ at every prime, so already modulo two they have infinite
support. Proposition~\ref{F-ideles} excludes them from the
completed idele group. This is a full counterexample to the
unrestricted family clause in the added task. For finite places
alone the pairs $(a,u)$ and $(-a,-u)$ give identical local
products: the missing real sign had distinguished these pairs.
The all-place formulation of Proposition~\ref{F-ideles} and its
restrictedness condition are the exact replacements proved here.
\end{proof}

\begin{statement}{F.23}{Proved}{geometric tangents are uniformizer specializations}
\label{st:F.23}
\label{F-specialization}
Let $K/\mathbb Q$ be finitely generated, let a smooth marked
curve over $K[[\pi]]$ have a rational cusp section, and let a
generic rational point meet this cusp with multiplicity $d>0$.
Write a parameter along the cusp as
\[
 z=a\pi^d(1+h),\qquad a\in K^*,\quad h\in\pi K[[\pi]].
                                                                  \tag{F.7}
\]
Specialization using the uniformizer splitting of the base
Galois group gives exactly the tangent packet class $\kappa_K(a)$.
Replacing $\pi$ by $b\pi(1+h')$ replaces it by
$\kappa_K(ab^{-d})$. Every geometric tangent class occurs in
this way. Thus in this precise smooth marked trait situation the
geometric subtorsor is exactly the set of such specializations.
\end{statement}
\begin{proof}
The base splitting fixes compatible roots of $\pi$ and acts on
coefficients by $G_K$. The factor $1+h$ has a unique $n$th
root congruent to $1$ for every $n$, by Hensel's lemma in residue
characteristic zero; these roots are compatible. Pulling back the
Kummer cover $z^{1/n}$ along the point of (F.7) therefore leaves
precisely the cocycle from $a^{1/n}$. The inertia homomorphism is
multiplication by $d$, while the residue section coordinate is
$\kappa_K(a)$. This proves the formula at every finite level and
hence in $T$, and substitution of the new uniformizer proves
its stated change. Conversely the formal smooth parameter chart
at the cusp admits an arc with $z=a\pi$. Its generic point is
in the open curve and its specialization is the desired class.
For the tripod this is the explicit point $x=a\pi$; formal
smoothness gives the same construction for any smooth rational
cusp. This is the exact comparison asked in scans 85--95.
\end{proof}

\begin{statement}{F.24}{Open}{a formula using only GT automorphisms}
\label{st:F.24}
\label{F-open-tangents}
The local conditions do characterize the geometric tangent
subtorsor over $\mathbb Q$ by
Proposition~\ref{F-local-global-units}. Over a general number
field their exact residual defect is $\widehat E/E$.
The cocycle-twist automorphisms in Proposition~\ref{F-twist}
prove why the isolated local extension does not remove this
defect. We have not produced or refuted a single
Grothendieck--Teichm\"uller relation selecting every prescribed
parameter $t\in\mathbb Q^*$ as demanded by F(g)(5). The
specialization characterization in
Proposition~\ref{F-specialization} is proved with the indicated
family and uniformizer data; it does not erase those data to
produce an additional intrinsic equation in the automorphism
group of the single tripod.
Theorem~\ref{F-two-Hom-tangents} below does give an intrinsic
global relation inside that arithmetic group when two open
homomorphisms from one of its open subgroups are allowed.
It recognizes the entire ordinary tangential subtorsor over
every number field, including the unit-completion defect above.
The remaining question in this remark concerns the expressly
more restrictive automorphism-only form of the relation.
\end{statement}

\subsection{Further exact comparisons and remaining reductions}

\begin{statement}{F.25}{Proved}{inertia and its normalizer in every rank}
\label{st:F.25}
\label{F-valuation-structure}
Let $L/K$ be finitely generated of characteristic zero, and let
$v$ be a valuation of $L$ trivial on $K$. Fix its prolongation to
$\overline L$, and write $\Gamma_v=v(L^*)$. There are a split
continuous exact sequence and canonical character identifications
\[
 1\longrightarrow I_v\longrightarrow D_v
   \longrightarrow G_{\kappa(v)}\longrightarrow1,
 \qquad
 I_v=\operatorname{Hom}((\Gamma_v\otimes\mathbb Q)/\Gamma_v,
                         \mu_\infty)
     =\operatorname{Hom}(\Gamma_v,T).
\]
The splitting has no preferred choice, and the action on $I_v$
is cyclotomic. Noncanonically,
\[
 I_v\simeq\prod_\ell\mathbb Z_\ell(1)^{r_\ell},\qquad
 r_\ell=\dim_{\mathbb F_\ell}\Gamma_v/\ell\Gamma_v
 \leq\operatorname{rank}_{\mathbb Q}\Gamma_v
 \leq\operatorname{trdeg}_K L-\operatorname{trdeg}_K\kappa(v).
\]
For $I_v\ne1$, let $\Delta$ be the largest divisible convex
subgroup of $\Gamma_v$ and let $w=v/\Delta$ be the coarsening.
Then, as actual subgroups for the specified prolongations,
\[
 I_w=I_v,\qquad N_{G_L}(I_v)=D_w.
\]
Thus $N_{G_L}(I_v)=D_v$ exactly when the residual valuation
$v/w$ on $\kappa(w)$ is henselian. In particular it holds for
rank one with nonzero inertia. For $I_v=1$ the normalizer is
all of $G_L$. These are the precise replacements for the
codimension and normalizer assertions on cote 149, scans
119--124 and 127.
\end{statement}
\begin{proof}
We checked F149-D4--D5 and D10 against the named sources.
Over the strict henselization the residue field is algebraically
closed and every unit has all roots, by Hensel's lemma. Residue
characteristic zero makes every finite extension tame and
defectless; this is Kuhlmann, Theorem 1.2 and Corollary
2.9~\cite{F-KuhlmannTame}. Temkin, Theorem 3.4.4, describes
the tame extensions by the finite enlargements of the value
group~\cite{F-Temkin}. Explicitly the pairing is
\[
 (\sigma,v(a)/n)\longmapsto \sigma(a^{1/n})/a^{1/n}.
\]
It is independent of the representative because the quotient
of two representatives of the same value is a unit with an
$n$th root. This gives the first character group. Its image in
$\operatorname{Hom}(\Gamma_v,T)$ sends $\gamma$ to the
compatible sequence whose $n$th term evaluates the character
on $\gamma/n\bmod\Gamma_v$. This gives the second displayed
isomorphism, rather than an unidentified abstract equivalence.
The torsion group $(\Gamma_v\otimes\mathbb Q)/\Gamma_v$ is
divisible and has $\ell$-torsion $\Gamma_v/\ell\Gamma_v$.
Choosing compatible roots of a basis gives its sum of Pr\"ufer
groups and the stated dual product. Independence modulo $\ell$
implies rational independence, by clearing denominators and
canceling a largest common power of $\ell$. Finally, elements
with rationally independent values, together with units whose
residues form a transcendence basis, are algebraically
independent: in any proposed polynomial relation the monomials
of least value have a nonzero residue polynomial. This proves
the two rank bounds.

For completeness the splitting can be constructed in the chosen
algebraic closure. Choose a maximal algebraic extension $M$ of
the henselization with unchanged residue field. Its value group
is divisible: adjoining an $\ell$th root of an element with value
outside $\ell\Gamma_M$ has ramification degree $\ell$ and
unchanged residue, contradicting maximality. The maximal
unramified extension of $M$ has divisible value group and
algebraically closed residue. Defectlessness makes it algebraically
closed. Hence $G_M\to G_{\kappa(v)}$ is an isomorphism, and
its inclusion in $D_v$ supplies the splitting. This also checks
the characteristic-zero specialization of Haran--Jarden--Pop,
Proposition 7.3(a)~\cite{F-HaranJardenPop}. Conjugation in the
Kummer pairing acts on its roots of unity by the cyclotomic
character and leaves the value group fixed.

Convex subgroups form a chain, and the union of its divisible
members is divisible. For a convex subgroup $\Delta'$ the
coarsening inertia is the character subgroup obtained by pullback
from $\Gamma_v/\Delta'$. It equals $I_v$ precisely when
$\Delta'$ is divisible. Put $F=\overline L^{I_v}$. Its induced
valuation is henselian, with value group $\Gamma_v$ and
algebraically closed residue field. The coarsenings with
algebraically closed residue are precisely those with divisible
convex kernel, by the already proved inertia computation applied
to their residual refinements. Thus $w|F$ is the coarsest of
these coarsenings. It is coarsest among all henselian valuations
of $F$ with separably closed residue: two incomparable henselian
valuations have a common coarsening with separably closed
residue, by the F.~K.~Schmidt theorem in Kuhlmann, Theorem 9.44
and Corollary 9.45~\cite{F-KuhlmannHensel}. Such a common
coarsening would contradict the choice of $w$. Comparability
then forces every such valuation to refine $w$.

Every element normalizing $I_v$ acts on $F$ and preserves this
intrinsically specified coarsest valuation. Its unique
prolongation to $\overline L$ is therefore preserved, placing
the element in $D_w$. Conversely $D_w$ normalizes $I_w=I_v$.
The inverse image of the residual decomposition group of $v/w$
in $D_w$ is $D_v$, so equality is precisely henselianity of the
residual valuation. In rank one the convex subgroups are $0$
and $\Gamma_v$; nonzero inertia excludes divisibility of the
latter. This proves the last assertions and the actual comparison
of all normalizers.
\end{proof}

\begin{statement}{F.26}{Proved}{the finite-residue requirement for germs}
\label{st:F.26}
\label{F-valuation-germs}
For a valuation $v$ as in Proposition~\ref{F-valuation-structure}
whose residue field is algebraic over $K$, the decomposition group
$D_v$ contains a section over an open subgroup of $G_K$ if and
only if $[\kappa(v):K]<\infty$. It contains a section over all
of $G_K$ precisely when $\kappa(v)=K$. Algebraicity alone
does not suffice, even for a DVR of $\mathbb Q(x,y)$.
With the finite-residue qualification, restriction to any
one-variable subfield is exactly the construction on scans
125--126: its valuation ring is $V\cap L_1$, and its section
lands in the corresponding rational-place decomposition group
after the finite residue extension.
\end{statement}
\begin{proof}
Use the embedding $\kappa(v)\subset\overline K$ fixed by the
prolongation. The image of $D_v\to G_K$ is exactly
$G_{\kappa(v)}$. It contains an open subgroup precisely when
this residue extension is finite. The splitting in
Proposition~\ref{F-valuation-structure} then supplies the section
over that subgroup. The same argument with the full image proves
the assertion over $G_K$.

Here is an explicit infinite-residue example. Choose distinct
rational primes $p_1,p_2,\ldots$, put $a_j=\sqrt{p_j}$,
$F_m=\mathbb Q(a_1,\ldots,a_m)$, and $F_\infty=\bigcup_mF_m$.
Inside $F_\infty((t))$ set
\[
 x=t,\qquad y=\sum_{j\geq1}a_jt^j,
 \qquad L=\mathbb Q(x,y),
\]
and restrict the $t$-adic valuation. Its value group is
$\mathbb Z$ because $x=t$, and its residue field is contained
in $F_\infty$. The residue of $y/x$ is $a_1$. For $m\geq1$
form the rational norm polynomial
\[
 P_m(X,Y)=\prod_{\sigma\in\operatorname{Gal}(F_m/\mathbb Q)}
 \left(Y-\sum_{j=1}^m\sigma(a_j)X^j\right)\in\mathbb Q[X,Y].
\]
The identity factor at $(t,y)$ starts with $a_{m+1}t^{m+1}$.
For a nonidentity sign vector whose first changed index is $j$,
the first term is $2a_jt^j$; there are $2^{m-j}$ such vectors.
Consequently
\[
 v(P_m(x,y))=2^{m+1}-1,
 \quad
 \overline{x^{-(2^{m+1}-1)}P_m(x,y)}
 =2^{2^m-1}\!\left(\prod_{j<m}p_j^{2^{m-j-1}}\right)
   a_m a_{m+1}.
\]
Induction gives $F_m\subset\kappa(v)$ for every $m$, hence
$\kappa(v)=F_\infty$. Were $y$ algebraic over $\mathbb Q(t)$,
the fundamental inequality for the resulting finite valued-field
extension would bound its residue degree. The displayed residue
field has degrees $2^m$ for all $m$, a contradiction. Thus this
$L$ is the rational function field in two variables, and the
valuation is a DVR with infinite algebraic residue. Its inertia
is $T\ne1$, its normalizer is $D_v$, and neither group supplies
the asserted germ. The exact program
\texttt{infinite\_residue\_checks.py} verifies the first four
norm polynomials and their leading coefficients; the induction
above proves the infinite assertion.

For a subfield $L_1$, either $b$ or $b^{-1}$ belongs to
$V\cap L_1$ for every $b\in L_1^*$, and the same alternatives
express $b$ as a quotient of two elements of that ring. It is
therefore a valuation ring with fraction field $L_1$, and its
residue injects into $\kappa(v)$. Compatible prolongations give
the actual restriction maps $D_v\to D_{v|L_1}$ and
$I_v\to I_{v|L_1}$. If $L_1/K$ has transcendence degree one,
the restricted valuation is nontrivial, since its residue embeds
in an algebraic extension of $K$. Its center on the smooth proper
curve is a closed point and its ring is the DVR there. The
section therefore restricts to exactly that decomposition group.
This proves the manuscript's forward construction. On scan 126
we visually read ``fini'' after the struck-out ``alg''; that
qualification is mathematically necessary. The supplied first-pass
transcription reports the same insertion on scan 124, whose image
is not included in the supplied selection.
\end{proof}

\begin{statement}{F.27}{Proved}{the one-variable equivalence has both lifting directions}
\label{st:F.27}
\label{F-birational-equivalence}
Fix a smooth proper geometrically connected curve $X/K$, with
$K/\mathbb Q$ finitely generated and $L=K(X)$. The conjunction
of (F4)(a) for $L K'/K'$ and (F4)(b) for every hyperbolic open
$U\subset X_{K'}$, for all finite $K'/K$, is equivalent to
(F2) for all those open curves. In this comparison (a) means
that every birational section lies in a rational-place
decomposition group; in transcendence degree one the place
and its fixed prolongation are unique. Proposition~\ref{F-birational}
identifies this (a), section by section, with nontrivial
cyclotomic centralizer and gives the margin formula on scan 127.
\end{statement}
\begin{proof}
All maps are the quotients $q_U:G_{K'(X)}\to E_U$ induced by
the generic point. A section of $E_U$ lifts by (b); applying (a)
to the lift puts it in $D_v$ for a rational place of $X_{K'}$.
The map $q_U$ kills its inertia when the place belongs to $U$,
and identifies its decomposition group with the cusp group when
the place belongs to $X_{K'}\setminus U$. Thus its image is
respectively the point section or a member of that cusp packet.
This proves one implication.

For the reverse implication, a point section at $x\in U(K')$
lifts by any splitting of the $x$-adic decomposition sequence.
For a cusp $c$, the homomorphism $D_c(G_{K'(X)})\to D_c(E_U)$
is an isomorphism: it is an isomorphism on the inertia $T$ and
on the residue quotient $G_{K'}$. A given cusp section therefore
lifts through that actual isomorphism. Conjugators lift because
$q_U$ is surjective. This constructs every lift required by (b).

Now start with a birational section $s$. For every hyperbolic
open $U$, (F2) and Proposition~\ref{F-fixed-dictionary} give
its unique fixed pro-point in the completed universal cover of
$U$. For $U'\subset U$, forgetting coverings ramified in
$U\setminus U'$ maps that fixed pro-point to the unique one for
$U$. They are therefore compatible. Every finite separable
extension of $K'(X)$ is unramified over some such $U$, so these
compatible pro-points give points on the normalizations of $X$
in every finite extension of $K'(X)$. Their local rings form the
valuation ring of a prolongation of one place; equivalently this
is the usual inverse-limit description of places on a curve.
The section fixes all these points and hence preserves that
valuation. The place on $X$ is $K'$-rational, because its
decomposition image contains all of $G_{K'}$. This proves (a).
Two invariant prolongations would give two fixed pro-points at
some finite covering level, contradicting uniqueness. Thus the
comparison includes the uniqueness assertion that is valid in
one variable.
\end{proof}

\begin{statement}{F.28}{Proved}{the proper-base part of the polycurve lifting reduction}
\label{st:F.28}
\label{F-proper-polycurve-lifting}
The universal curve-lifting assertion (F4)(b) over all finitely
generated extensions of $\mathbb Q$ implies (F4)(b) for every
hyperbolic polycurve whose tower is proper at every stage except
possibly the last. The comparison is given by taking the point
fibre of the projected section and then a valuation centered at
the resulting point or relative cusp. This proves the stated
part of the implication asked in F(e)(4); towers with a
cuspidal projected section on an affine intermediate base are
not included in the assertion.
\end{statement}
\begin{proof}
Proposition~\ref{F-cuspidalization} proves the implication from
the universal curve-lifting assertion to the universal curve
point-or-cusp assertion. For a proper hyperbolic polycurve
$B/K$, induction shows that every section is a point section.
Indeed its projection to the preceding stage is, after geometric
conjugation, the section $s_b$ of a rational point $b$. Hoshi,
Proposition 2.4(ii), identifies the pullback
\[
 E_B\times_{E_{B'},s_b}G_K=E_{B_b}.
\]
This is a pullback along an actual point, with the same domain
$G_K$. The given section is now a section of the proper
hyperbolic fibre. The curve assertion makes it a point section,
completing the induction.

Apply this to the proper base $B$ of the last map $U\to B$.
The same pullback identifies the original section with a section
of $U_b$. It is either the section of a rational point or a cusp
section of this fibre. In the former case choose a regular system
of local parameters at the point of $U$. Their successive prime
ideals define a full flag; composing the resulting discrete
valuations gives a rank-$\dim U$ valuation of $K(U)$ with
residue $K$ and center at that point. A decomposition splitting
exists by Proposition~\ref{F-valuation-structure}; its image in
$E_U$ is the prescribed point section because finite \'etale
covers are unramified at that center.

In the cuspidal case take the smooth relative compactification
of $U\to B$. The relative cusp is a $K$-point $c$ over $b$,
and its boundary is \'etale over $B$. Let $z$ be a local
equation for that boundary and complete it by parameters of
$B$ at $b$. The same full-flag valuation has residue $K$ and
inertia $T^{\dim U}$. Its map to the local punctured tube is
the coordinate projection $T^{\dim U}\to T$ belonging to
$z$: on $n$th roots it is exactly the Kummer character of $z$.
Choose any splitting supplied by
Proposition~\ref{F-valuation-structure}. The prescribed fibre
cusp section differs from its image by a cocycle in
$Z^1(G_K,T)$. Insert that cocycle in the
$z$ coordinate of $T^{\dim U}$ and put zero in the others.
Twisting the decomposition splitting by it gives a section of
$G_{K(U)}$ whose image is the prescribed cusp section. This
constructs the required lift, rather than attempting to replace
a $G_K$-section by a section with domain $G_{K(B)}$.
\end{proof}

\begin{statement}{F.29}{Proved}{point-theoretic maps and the proper-target reduction}
\label{st:F.29}
\label{F-point-theoretic}
Let $K$ be a number field or a finite extension of $\mathbb Q_p$,
$V/K$ a smooth geometrically connected variety, and $U/K$ a
hyperbolic curve. Call a continuous
$G_K$-homomorphism $f:E_V\to E_U$ \emph{point-theoretic}
when every point section over a finite extension becomes a point
section after a further finite extension. Such an $f$ is either
open or kills $P_V$ and is the point section of a unique
$K$-rational point composed with $E_V\to G_K$. In particular
every such $f$ comes from a unique morphism $V\to U$.

For proper $U$ over a number field $K$, condition (F1)(b) holds
for every continuous $G_K$-homomorphism $E_V\to E_U$. The all-source assertion
(F1) for this proper target and all finite extensions of $K$
is logically equivalent, via composition with point sections,
to (F2) for $U$ over those finite extensions. Thus the reduction
asked in F(b)(5) is proved for proper targets. For affine
targets the point-theoretic hypothesis is the additional
premise used in the proved map-level statement.
\end{statement}
\begin{proof}
We use Carlson--Stix, Proposition C and Proposition 3.17 in its
hyperbolic-target range~\cite{F-CarlsonStix}: a point-theoretic
map between smooth curves with hyperbolic target is open or
contained in one point decomposition group. We checked its
argument and its field hypothesis: it is stated over every
Kummer-faithful field of characteristic zero. Both number fields
and finite extensions of $\mathbb Q_p$ have this property.
For the latter, a semiabelian variety over every finite extension
has a subgroup of finite index described by finite-rank valuation
groups and $p$-adic Lie groups; its intersection of all $n$-multiples
is zero. Equivalently one uses the compact open formal subgroup,
its $p$-power filtration, and the finitely generated valuation
quotient. Finite covers of a hyperbolic target have Euler
characteristic tending to minus infinity as their degrees grow.
The generalized-Jacobian dimension argument makes the induced
map on algebraic points have finite image. The group it generates
is both finitely generated and a quotient of a divisible
generalized-Jacobian group, hence zero. This proves constancy;
applying it at all finite covering levels gives the decomposition
group conclusion. The broader wording of the body proposition
only assumes nontrivial geometric fundamental group, but its
written Euler-characteristic argument requires hyperbolicity.
Our application has that hypothesis.

Suppose $f$ is not open. The induced map
$V(\overline K)\to U(\overline K)$ is well defined by point
section injectivity over finite extensions. Any two geometric
points of $V$ can be joined by a chain of curves: use affine
neighborhoods with nonempty successive intersections and
Bertini complete intersections through the chosen points,
then normalize. These curves and their points are defined over
some finite extension of $K$; puncturing further makes the source
curves hyperbolic. Composition with each curve's fundamental
group cannot be open, because that would make $f$ open. The
curve result therefore makes the point map constant on each
curve and then on all of $V(\overline K)$.

We now prove triviality on $P_V$, rather than infer it merely
from this constancy at the first level. After a finite Galois
extension $K'/K$, choose $x\in V(K')$, $y\in U(K')$ and
paths for which $f\circ s_x=s_y$ on $G_{K'}$. For each open
characteristic subgroup $N\subset P_U$, the subgroup
$H_N=N s_y(G_{K'})$ defines a geometrically connected finite
\'etale cover $U_N\to U_{K'}$. Its inverse image
$J_N=f^{-1}(H_N)$ defines a geometrically connected finite
\'etale cover $V_N\to V_{K'}$: it maps onto $G_{K'}$ since
it contains $s_x(G_{K'})$. The induced map $J_N\to H_N$
is still point-theoretic. Indeed a point section in the cover
maps to a point section downstairs, and its containment in
$H_N$ is exactly the fixed-coset criterion for a lift to
$U_N$, after a finite residue extension.

The same curve-chain argument makes
$V_N(\overline K)\to U_N(\overline K)$ constant. On the
fibres above $x$ and $y$ its map is the actual coset map
\[
 P_V/(P_V\cap J_N)\longrightarrow P_U/N,
 \qquad a\longmapsto f(a)N.
\]
This follows either from the finite covering fibre functors or
from their point-section stabilizers and injectivity on $U_N$.
Constancy therefore gives $f(P_V)\subset N$. Characteristic
open subgroups are cofinal because $P_U$ is topologically
finitely generated, so $f(P_V)=1$. Thus $f$ factors through a
section of $G_K$. That section restricts to $s_y$ over $K'$.
Normality of $G_{K'}$ and the self-normalizing property of
point decomposition groups make the same pro-point fixed by the
whole section; its image $y$
is $K$-rational, and the section is its point section. In the
open case Proposition~\ref{F-open-Hom} gives the unique dominant
morphism. The point case gives the unique constant morphism.

For proper $U$ over a number field, take an open subgroup
$J\subset E_V^{\circ}$
and a finite \'etale source cover with a closed point whose
cyclotomic section lands in $J$. Its composite under $f$ is
some section of the proper curve over a number field. The
proper case of Proposition~\ref{F-centralizer} makes its
geometric cyclotomic centralizer zero, whether or not that
section is known to come from a point. This kills the
centralizer of $f(J)$; projection to $G_K$ kills the arithmetic
one as in Proposition~\ref{F-geometric-necessary}. Thus (b)
is automatic. Finally (F2) for the proper target over every
finite extension makes the composite of $f$ with every point
section a point section. The first part then proves (F1).
Conversely the point-source instance of (F1) is (F2), as in
Proposition~\ref{F-point-equivalence}. These constructions
prove the stated equivalence of assertions without asserting
either unresolved universal surjectivity statement as a theorem.
\end{proof}

\begin{statement}{F.30}{Proved}{the full curve-target reduction of (F1) to (F2)}
\label{st:F.30}
\label{F-full-curve-reduction}
Over all number fields, the following two universal assertions
are equivalent: (F2) for hyperbolic curves, and (F1) with a
hyperbolic curve as target and an arbitrary smooth connected
finite-type source. The comparison retains every non-open map
and every cusp packet. More precisely, the (F2) assertion implies
that a non-open arithmetic map to a hyperbolic curve has
geometric image either trivial or contained in one cusp inertia.
In the nontrivial case its entire image is in that cusp's
decomposition group, and condition (b) fails. This is a proved
implication between the assertions requested in F(b)(5), not
an assertion that their common universal surjectivity holds.
\end{statement}
\begin{proof}
The reverse comparison is the point-source calculation of
Proposition~\ref{F-point-equivalence}. We prove the forward
comparison as a map between the two assertions. Work with a
continuous map $f:E_V\to E_U$ over a number field $K$, where
$V$ is smooth and geometrically connected. An open map is
handled by Proposition~\ref{F-open-Hom}, so consider a
non-open map. Choose a point $x\in V(K')$ after a finite
Galois extension $K'/K$. The (F2) assertion applied to
$s=f\circ s_x$ gives a fixed lifted interior point or cusp
$c$ in the completed universal cover of $U$. Write $I_c=1$
in the interior case and use the usual cusp inertia otherwise.

Choose a characteristic open subgroup $N_0\subset P_U$ whose
cover has compactification of genus at least two. Such a
geometric cover exists for every hyperbolic curve. In genus
zero use the degree-four Fermat covering of the tripod, after
selecting three punctures. In genus one an unramified cover of
the compactification first supplies at least two punctures,
and a double cover branched at two of them has genus two;
over an algebraically closed field the required square root of
their degree-two divisor exists in the Picard group. A
characteristic refinement preserves this genus bound. Crucially,
$N_0$ is chosen as an arbitrary geometric cover; it is not
required to contain $f(P_V)$.

For every characteristic open $N\subset N_0$, put
\[
 H_N=N s(G_{K'}),\qquad J_N=f^{-1}(H_N)
       \subset E_{V_{K'}}.
\]
These define geometrically connected finite \'etale covers
$U_N$ and $V_N$ over $K'$, since both subgroups project onto
$G_{K'}$. Their chosen point or cusp $c_N$ is fixed by $s$.
Let $X_N$ be the smooth compactification of $U_N$.
We claim that the homomorphism
\[
 J_N\longrightarrow H_N\longrightarrow E_{X_N}
\]
comes from the constant morphism $V_N\to X_N$ with value
$c_N$. The (F2) assertion for proper curves makes this map
point-theoretic, so Proposition~\ref{F-point-theoretic} gives
a unique morphism. To prove it constant, hold this $N$ and
this source $V_N$ fixed. Because $f(J_N)$ is non-open in
$H_N$, there are open subgroups $A\subset H_N$ containing
$f(J_N)$ with unbounded indices. Their covers $U_A$ remain
geometrically connected over $K'$. The compactifications obey
\[
 g(X_A)\geq1+[H_N:A]\bigl(g(X_N)-1\bigr),
\]
by Riemann--Hurwitz, hence have unbounded genus. Applying the
same proper point-theoretic argument gives morphisms
$h_A:V_N\to X_A$. A nonconstant such map is dominant and
induces an open fundamental-group map. Restriction in
degree-one $\ell$-adic cohomology is consequently injective:
a homomorphism from $P_{X_A}$ to $\mathbb Q_\ell$ vanishing
on an open subgroup vanishes everywhere. Therefore
\[
 2g(X_A)\leq \dim_{\mathbb Q_\ell}
       H^1(V_{N,\overline K},\mathbb Q_\ell),
\]
an impossibility for sufficiently large index. For such $A$
the map $h_A$ is constant. Its composition to $X_N$ induces
the original homomorphism, and uniqueness gives constancy of
the original morphism. Its value is $c_N$, by evaluating at
the distinguished lift of $x$.

Now evaluate that constant morphism on all geometric points
of $V_N$ above $x$. The map on the target compactified fibre
is, with the chosen universal paths, precisely
\[
 P_V/(P_V\cap J_N)\longrightarrow N\backslash P_U/I_c,
       \qquad a\longmapsto N f(a)I_c.
\]
One can check the formula on point stabilizers: the section
of the lift labelled by $a$ maps to
$f(a)s(G_{K''})f(a)^{-1}$ after a finite extension splitting
the fibres. Its point in $X_N$ is the lift labelled by
$Nf(a)I_c$. Constancy at every $N$ now gives
\[
 f(P_V)\subset\bigcap_{N\subset N_0}N I_c=I_c,
\]
because $I_c$ is closed and these characteristic subgroups
are cofinal. This argument includes both interior and boundary
fibres, with exactly the displayed quotient between them.

If $f(P_V)\ne1$, the inertia normalizer formula gives
$N_{E_U}(f(P_V))=D_c$. Since $f(P_V)$ is normal in $f(E_V)$,
the whole image is contained in $D_c$, which fails (b) by
Proposition~\ref{F-point-equivalence}. Otherwise $f$ factors
through a section of $G_K$. Condition (b) excludes its cusp
case, and the (F2) assertion makes it the section of a unique
point. This yields the desired constant morphism. A connected
regular source over $\mathbb Q$ reduces to the geometrically
connected case over its number field of constants; compatibility
over $G_{\mathbb Q}$ gives the required embedding of the
target constants. Thus the stated all-source comparison has
exactly the domain of (F1).
\end{proof}


\begin{statement}{F.31}{Proved}{the infinite-index dichotomy}
\label{st:F.31}
\label{F-local-dichotomy}
Let $K/\mathbb Q$ be finitely generated, $V/K$ a smooth
geometrically connected variety, $U/K$ a hyperbolic curve,
and $f:E_V\to E_U$ a
continuous $G_K$-homomorphism. Then exactly one of the following
three possibilities holds for its geometric image $H=f(P_V)$:
\[
 H\text{ is open in }P_U;\qquad H=1;\qquad
 1\ne H\subset I_c\text{ for one lifted cusp }c.
\]
In the third case $f(E_V)\subset D_c$ and the cusp is
$K$-rational. Consequently condition (F1)(b) implies that $H$
is open or trivial. In the trivial case the actual factorization
is
\[
 E_V\longrightarrow G_K\xrightarrow{\,s\,}E_U,
\]
with $s$ a continuous section. This proves the first assertion
of F(b)(5) with the local section theorem explicitly supplied
as tool (1) in the prompt~\cite{B_OpenAIPadic}. Its October
2026 preprint status
and lack of a formal verification are retained as the status
of that input; no independent proof of that input is asserted
here. Together with Proposition~\ref{F-full-curve-reduction},
this completes the requested curve-target reduction, without
asserting the remaining global section-surjectivity statement.
\end{statement}
\begin{proof}
Choose an embedding $K\hookrightarrow k$, where $k$ is a
finite extension of some $\mathbb Q_p$. Such an embedding
exists: take a smooth model over the number field of constants,
a closed point after a finite extension, and an \'etale local
coordinate chart around it. Over a completion the chart contains
a $p$-adic open ball. Choose algebraically independent coordinates
in that ball; its lift is a generic point and gives an injective
evaluation of the function field into the completion. For a
number field one may simply take any finite-place completion.
Choose compatible separable closures. Base change gives the
actual cartesian
diagrams
\[
 E_{V_k}=E_V\times_{G_K}G_k,\qquad
 E_{U_k}=E_U\times_{G_K}G_k.
\]
The geometric groups are unchanged by the algebraically closed
extension in characteristic zero. No openness or injectivity
claim about $G_k\to G_K$ is needed. The induced $f_k$ therefore
has the very same geometric image $H\subset P_U$. Arithmetic
openness for a map over $G_k$ is equivalent to openness of its
geometric image.

The supplied local theorem, and its affine neighborhood
consequence (T1), say that every section of every hyperbolic
curve over every finite extension of $k$ is a point section or
a cusp section. For completeness, the non-open-image argument
in Proposition~\ref{F-full-curve-reduction} uses only this
local assertion and the point-theoretic statement of
Proposition~\ref{F-point-theoretic}, which was proved for
finite $p$-adic extensions as well. Here are its actual
constructions over $k$. After a finite extension choose
$x\in V(k')$ and let $c$ be the lifted point or cusp fixed
by $f_k s_x$. Choose an arbitrary characteristic target cover
whose compactification has genus at least two, and pull it
back to $V$; this fixes the source cover used next. Open
overgroups of the non-open image have arbitrarily large
indices. Their target compactifications have unbounded genus,
whereas a dominant map from that fixed source injects their
$H^1$ into its fixed finite-dimensional $H^1$. The local
proper section theorem makes all these maps point-theoretic,
so Proposition~\ref{F-point-theoretic} constructs them and
forces the maps of sufficiently high genus to be constant.
At every characteristic level $N$, the maps on the fibres
above $x$ are exactly
\[
 a\longmapsto N f_k(a) I_c\quad\text{in }N\backslash P_U/I_c.
\]
Their constancy proves $H\subset\bigcap_NNI_c=I_c$, where
$I_c=1$ for an interior point. This proves the local trichotomy
without needing condition (b) to remain true under localization.

Return to the global map. In the nontrivial inertia case,
$H\triangleleft f(E_V)$, and the geometric inertia-normalizer
identity is
\[
 N_{E_U}(H)=D_c.
\]
It follows that the entire global image lies in $D_c$. Its
projection is all of $G_K$, so the underlying cusp is
$K$-rational. Every element of $f(E_V^\circ)$ then centralizes
$I_c$, since the action on inertia is cyclotomic. Thus (b)
fails. In the remaining non-open case $H=1$, the quotient
property of $E_V\to G_K$ supplies the displayed continuous
map $s$, and the equation $p_Uf=p_V$ makes $s$ a section.
The cases are disjoint: a nontrivial cusp inertia has infinite
index in a hyperbolic geometric fundamental group. This last
fact also follows from its finite covers, whose Euler
characteristics cannot be those of a procyclic group.
\end{proof}

Statement B.39 extends this trichotomy to every sub-$p$-adic
field. Its explicit countable tower of finite gerbes descends
the homomorphism itself and then embeds that countable datum
in a finite $p$-adic field; the whole geometric homomorphism
and its peripheral subgroups are preserved. Statements
B.40--B.41 apply this extension to unramification over smooth
$p$-adic models and to the complete isotrivial section
classification. They retain the same supplied local input
and do not assert generic point realization for a nonconstant
relative section.

\begin{statement}{F.32}{Known}{higher-dimensional Hom and Isom ranges}
\label{st:F.32}
\label{F-higher-dimensional-ranges}
The cited higher-dimensional theorems yield the following precise
parts of F(b)(6).

For a sub-$p$-adic field $K$, a normal variety $V/K$, a hyperbolic
polycurve $U/K$ of dimension $n$, and an open arithmetic
homomorphism $f:E_V\to E_U$, Hoshi's Theorem A gives a unique
dominant morphism in each of these ranges:
\[
\begin{array}{c|l}
n&\text{additional hypotheses}\\ \hline
1&\text{none}\\
2&\ker f\text{ is topologically finitely generated}\\
3&\ker f\text{ is finite},\quad V\text{ is of LFG type},\quad\dim V\ge3\\
4&f\text{ is injective},\quad V\text{ is a hyperbolic polycurve},
       \quad\dim V\ge4 .
\end{array}
\]
Here LFG means that every nonconstant morphism from a normal
geometric variety to $V_{\overline K}$ has infinite geometric
fundamental-group image. In dimension three the resulting map
is quasi-finite and $\dim V=3$; in dimension four it is finite
\'etale and $\dim V=4$.

For $K/\mathbb Q$ finitely generated, Schmidt--Stix, Theorem 1.5
and Corollary 1.6, identify isomorphisms of strongly hyperbolic
Artin neighborhoods with outer arithmetic group isomorphisms.
The neighborhood hypothesis is a tower of elementary hyperbolic
fibrations in which every intermediate variety embeds locally
closed into a product of hyperbolic curves.

Holzschuh--Schmidt--Stix, in the sharp theorem at the end of
their paper, gives dominant relative Hom for $U\to S$ a
relative hyperbolic curve, with $S$ normal, connected,
excellent and of characteristic zero, and with sub-$p$-adic
generic field. The source $V$ is normal and connected,
$V\to S$ is flat and of finite type with connected geometric
generic fibre, and the smooth locus of $V\to S$ surjects onto
$S$. This includes the smooth surjective source in their
introductory Theorem A. Its input is a
$\pi_1$-open map of profinite \'etale homotopy types over
$\mathrm{Et}(S)$, with its commutation homotopy. When $S$ is a
hyperbolic polycurve this is exactly the corresponding relative
outer-homomorphism datum, through the canonical maps to
classifying spaces.

Every open arithmetic map between the objects of (F1) satisfies
(b). Thus the displayed Hom ranges prove both (a) and (b) for
their open maps. The Isom result is the isomorphism range of the
same comparison, through the explicit inclusion of isomorphisms
among open homomorphisms.
\end{statement}
\begin{proof}
The hypotheses and conclusions are those of Hoshi, Theorem A,
Definition 2.5 and Corollaries 3.15--3.16\cite{F-HoshiPoly},
Schmidt--Stix, Definition 6.1 and Theorem 1.5/Corollary
1.6\cite{F-SchmidtStix}, and the attached August 2026
Holzschuh--Schmidt--Stix Theorem A and its sharpened final
theorem\cite{F-HolzschuhSchmidtStix}. These primary statements
were read directly, together with the final proof section.
For the last theorem, the smooth-locus surjectivity gives
the exact sequence
\[
 \pi_1(V_{\overline\eta})\longrightarrow E_V
     \longrightarrow E_S\longrightarrow1.
\]
The corresponding target sequence has injective geometric
fibre group, since that group is centre-free. For a relative
homomorphism the two image indices consequently agree:
\[
 [E_U:f(E_V)]
 = [\pi_1(U_{\overline\eta}):
                f_\eta(\pi_1(V_{\overline\eta}))].
\]
This is the explicit comparison between total and generic-fibre
openness. Their proof then constructs the generic-fibre map,
spreads the commutation homotopy by pro-\'etale descent, and
applies the extension criterion for hyperbolic curves. The
displayed flatness and smooth-locus hypotheses are retained.

For the asserted comparison of data, an \'etale $K(\pi,1)$
object has the canonical equivalence
$\mathrm{Et}(X)\simeq B E_X$. The universal property of
$1$-truncation identifies maps from any $\mathrm{Et}(V)$ into
$B E_U$ with homomorphisms of its fundamental groupoid. Doing
this in the slice over $B E_S$ retains the specified
commutation isomorphism; it is the relative groupoid datum,
with natural transformations as its conjugacies. Hyperbolic
polycurves are $K(\pi,1)$ by the successive homotopy fibre
sequences. Thus this supplies the precise morphism from the
relative homotopy theorem to the group formulation. When $S$
is not a $K(\pi,1)$, we retain the homotopy datum itself.

Finally let $J\subset E_V^\circ$ be open and $f$ open.
Then $f(J)$ is open in $E_U^\circ$. Its intersection with
$P_U$ is open in $P_U$ and has trivial centralizer there:
geometric hyperbolic polycurve groups are slim by the fibre
sequence and the surface-group statement. The arithmetic
projection of a centralizer element centralizes an open
cyclotomic kernel in $G_K$ and is trivial by the Hilbertian
centralizer theorem used in Proposition~\ref{F-geometric-necessary}.
The element therefore belongs to $P_U$ and is trivial. This
proves both readings of (b) for open maps without presupposing
the existence of a morphism. The cited reconstruction theorems
then give (a), and Proposition~\ref{F-geometric-necessary}
gives its necessary direction.
\end{proof}

\begin{statement}{F.33}{Refuted}{the first surface case and product targets}
\label{st:F.33}
\label{F-surface-products}
The assertion that an arbitrary open arithmetic map to a
hyperbolic polycurve surface must come from a dominant morphism
is false already for a smooth proper curve as source. For
$C/\mathbb Q$ a smooth proper curve of genus at least two,
there is a smooth geometrically connected ample curve
$D\subset C\times C$ for which
\[
 E_D\longrightarrow E_{C\times C}
\]
is surjective and satisfies (b), although its unique geometric
realization is the nondominant inclusion. Its kernel is not
topologically finitely generated.

For every product $U=U_1\times\cdots\times U_r$ of hyperbolic
curves, every open arithmetic map from a normal variety is
induced by a unique morphism to $U$, with dominance determined
by that morphism's actual image. On arbitrary arithmetic maps,
condition (b) is exactly the conjunction of the conditions (b)
on all coordinate maps. Consequently universal (F1) for product
targets over number fields is equivalent to universal (F2) for
curves over number fields, through those coordinate maps.
\end{statement}
\begin{proof}
Choose a very ample line bundle on $C\times C$. Bertini gives
a nonempty open set of its sufficiently high-degree linear
system consisting of smooth geometrically connected divisors.
That open set has a rational point, since $\mathbb Q$ is
infinite; choose the corresponding $D$. For every connected
finite \'etale cover $Y\to(C\times C)_{\overline{\mathbb Q}}$,
the pullback of $D$ is ample and connected. Here connectedness
can also be checked directly by the Hodge index theorem:
a hypothetical disjoint decomposition into nonzero effective
divisors $D_1,D_2$ would give
$D_1D_2=0$ and $D_i^2=D_iD>0$, contradicting the single
positive direction of the intersection form. The finite-cover
criterion for fundamental groups therefore makes
$P_D\to P_{C\times C}$ surjective. Both arithmetic groups
surject onto $G_{\mathbb Q}$, so their arithmetic map is
surjective as well. It satisfies (b) by
Proposition~\ref{F-geometric-necessary}. The source has
dimension one and the target dimension two, so the inclusion
is not dominant: a dominant map would inject a function field
of transcendence degree two into one of degree one.
The product reconstruction below makes it the unique
realization of its group map. If the kernel were topologically
finitely generated, Hoshi's dimension-two theorem would instead
produce a dominant realization, a contradiction.

The exact arithmetic product formula is
\[
 E_U=E_{U_1}\times_{G_K}\cdots\times_{G_K}E_{U_r},
 \qquad P_U=\prod_iP_{U_i}.
\]
An open $f$ has open coordinate maps $f_i$. Mochizuki's
curve-target Hom theorem constructs their unique morphisms
$g_i$, and $g=(g_i)_i$ induces $f$ by this actual fibre-product
identity. Equality of all coordinates also proves uniqueness.
For every open $J\subset E_V^\circ$ one has
\[
 C_{P_U}(f(J))=\prod_i C_{P_{U_i}}(f_i(J)).
\]
The equivalence of geometric and arithmetic centralizer
readings proves the claimed coordinate test. Apply
Proposition~\ref{F-full-curve-reduction} to each coordinate
to obtain the implication from universal (F2) to universal
product-target (F1). Conversely take $U=X\times
(\mathbb P^1\setminus\{0,1,\infty\})$, multiply a noncuspidal
section of $X$ by the point section at $2$, and apply the
point-source instance of product-target (F1). Projection of
the resulting point gives the required point of $X$; cusp
packets are already supplied by the local construction.
This proves the reverse implication with the displayed maps.
\end{proof}

\begin{statement}{F.34}{Open}{the remaining surface Hom problem}
\label{st:F.34}
\label{F-surface-open}
After Proposition~\ref{F-surface-products}, the unrestricted
claim that openness implies dominance is refuted in the first
possible target dimension. The underlying reconstruction
question still left by F(b)(6) includes an open map to a
non-product hyperbolic polycurve surface whose kernel is not
topologically finitely generated, and for which the reconstructed
map to the base curve has no known flat presentation with
geometrically connected generic fibre whose smooth locus
surjects onto that base. This statement is
not declared proved or refuted here.

We tried the following explicit reduction. Compose with the
surface's base-curve projection and reconstruct that map by
Mochizuki. Over the largest open base where the source map is
smooth, use the fibre product
$E_{U\times_B B'}=E_U\times_{E_B}E_{B'}$ to form the
restricted relative datum. The missing input is openness on the
relative geometric fibre and extension across the removed base
points. Original openness can receive contributions from loops
around those removed points; it does not itself establish
openness of the image of the geometric generic-fibre group.
Hoshi's finite-generation condition supplies precisely a
normal-subgroup argument at this step. The HSS relative theorem
supplies an extension only under its displayed flatness,
smooth-locus surjectivity and homotopy hypotheses. No counterexample to
reconstruction in this remaining range has been constructed.
\end{statement}

\begin{statement}{F.35}{Proved}{the 1981 and 1983 formulations}
\label{st:F.35}
\label{F-manuscript-comparison}
The comparisons requested in F(a) are given by the following
actual forgetful maps, restrictions and identifications.

For hyperbolic pairs $(X',S')$, $(X,S)$ over a number field,
put $U'=X'\setminus S'$ and $U=X\setminus S$. There is a
commutative square
\[
\begin{array}{ccc}
 \operatorname{Hom}^{\rm dom,\partial}(U',U)&\lhook\joinrel\longrightarrow&
                 \operatorname{Hom}^{\rm dom}(U',U)\\
 \downarrow\scriptstyle\sim&&\downarrow\scriptstyle\sim\\
 \operatorname{Hom}^{\rm open}_{\mathcal B}(\mathcal B_{U'},\mathcal B_U)
       &\lhook\joinrel\longrightarrow&
 \operatorname{Hom}^{\rm open}_{G_K}(E_{U'},E_U)/P_U .
\end{array}
\]
The boundary superscript means the literal boundary condition
$\overline f^{-1}(S)=S'$ of the corrected 1981 category.
The bottom arrow forgets the boundary diagram, its local
conjugators and its positive integer degrees. Its image is
exactly the inverse image, under the right-hand reconstruction,
of that boundary condition. Thus the open curve part of (F1)
implies the corrected dominant 1981 assertion; the latter
recovers the same (F1) assertion on precisely this displayed
image. Every dominant geometric map is represented in the
boundary category after the explicitly described shrinkings
in the proof. No equivalence obtained by silently forgetting
that boundary restriction is asserted.

On sections of a geometric curve, the corrected admissibility
condition of Longue Marche scan 79 is
\[
 \text{normalizes no cusp inertia}
 \quad\Longleftrightarrow\quad C_{P_U}(s(H_K))=1.
\]
Consequently its assertion that all admissible effective lifts
are point sections is equivalent to (F2), and also to the
point-source case of (F1), with every cusp packet retained on
the (F2) side. The letter's conjecture (7), with its packet
correction, is the same point-or-cusp assertion. Its
cyclotomic postscript is the right-hand zero-centralizer
condition in this displayed comparison.

The stronger guess of scan 433 holds for point sections and
fails for every cusp section: their $H_K$-fixed subgroups are
respectively $1$ and $I_c$. For the full Galois group and all
open subgroups its fixed subgroup is $1$ in both cases.
Thus the comparison from the full action to its restricted
$H_K$-action is literally restriction of the same action,
and the exact change of fixed subgroup is as stated.

Axioms (1)--(6) of scans 412--421, on realized geometric
objects with actual interior points, have the meanings and
descent maps proved in B.13. On those objects, replacing
``interior point'' by the intrinsic zero-centralizer selection
adds exactly the surjectivity assertion in (F2). For arbitrary
abstract outer actions, the inclusion of the geometric objects
into that larger collection does not prove essential
surjectivity or the intrinsic forms of axioms (5)--(6);
those are the unresolved tasks identified in Part B.

The first enlargement of the explicitly stated source category
in these 1983 assertions is the arbitrary regular finite-type
source in scan 48. The relation to the 1981 dimension-at-most-one
category is its restriction along
$(X,S)\mapsto X\setminus S$. A product of two positive-dimensional
hyperbolic curves has dimension two and is permitted as a source
in (F1), whereas it is not an object of that stated 1981 category.
All-morphism questions already occur in the 1981 text; their
mere presence in 1983 is not a new scope assertion.
\end{statement}
\begin{proof}
The functor, its two-morphisms and the exact Kummer degree
condition are constructed in Propositions~\ref{A:typed-functor}
and~\ref{A:two-morphisms}. Positive peripheral degrees give
openness in the nonempty-boundary case by
Proposition~\ref{A:positive-degree-open}. The curve Hom theorem
and Proposition~\ref{A:relative-hom} give the two vertical
bijections. Proposition~\ref{A:boundary-recovery} proves that
the bottom forgetful arrow is injective and identifies its
image: a nontrivial peripheral image selects a unique target
inertia, its normalizer forces the local decomposition map,
and the Kummer formula recovers the positive integer degree.
The diagram therefore compares the actual functors of
Conjectures 5.5, its correction in section 9, and 7.1 with
(F1), without discarding their natural transformations.

Here is the asserted restriction for a dominant map whose
original punctures do not satisfy the boundary condition.
Extend it to the smooth complete curves and put
\[
 S_1=S\cup\overline f(S'),\qquad
 S'_1=\overline f^{-1}(S_1).
\]
The restrictions $X'\setminus S'_1\to X\setminus S_1$ and
the two open immersions give a commuting square of curves,
and hence a commuting square of their arithmetic groups.
The restricted map satisfies the exact boundary condition;
its exponents are its positive ramification indices. This is
the explicit relation between the unrestricted dominant map
and its boundary enhancement.

For the section comparison, every cusp-packet section
normalizes its inertia by construction, including all the
nongeometric tangential parameters. Proposition~\ref{F-centralizer}
proves the reverse recognition and the fixed-group formulas,
for sections and germs. Point-section and packet injectivity
make their images disjoint. Removing the whole union of cusp
packets from the equality asserted in (F2) gives exactly its
admissible-section assertion; restoring those packets gives
the converse. Proposition~\ref{F-point-equivalence} identifies
that same complement with (F1)(b) for a point source.
B.10 proves the full-action fixed-group statement. The
restriction to $H_K$ is the explicit homomorphism of actions
used in the display, so this also proves the requested
comparison with the stronger guess and the postscript.

For the axioms, B.13 uses finite base change for (1), deletion
and filling for (2), vanishing of fixed groups for (3), a
closed point after finite extension for (4), the equality
$(U\cup\{c\})(K)\setminus\{c\}=U(K)$ for (5), and
actual descent of the curve and its finite \'etale marked
subscheme for (6). Their uniqueness is supplied by the same
fixed-group vanishing on the configuration-space kernel.
These are maps on the realized subcategory. Extending them
to every abstract object requires precisely the realization
or surjectivity assertions still named in Part B, and is not
an additional theorem proved by this comparison. Finally the
source-category inclusion and the dimension of the displayed
product verify the claimed enlargement directly.
\end{proof}

\begin{statement}{F.36}{Proved}{Known reductions and a proved obstruction diagram for F(c)}
\label{st:F.36}
\label{F-field-reduction}
Let $\mathrm{SC}_{\rm nf}$ denote (F2) over all number fields.
There is a proved implication from this assertion to (F2)
over finitely generated characteristic-zero fields for every
hyperbolic curve $U/K$ admitting a geometric finite \'etale
cover $Y\to U_{\overline K}$ and a nonconstant map
$Y\to A_{\overline K}$, where $A$ is defined over
$\overline{\mathbb Q}$. This includes all proper hyperelliptic
curves, every hyperbolic type of genus at most two, and every
affine curve with such a cover mapping nonconstantly to the
tripod. Every curve has a nonempty open subcurve in this last
range. This is an implication between the two section
assertions, not an assertion that $\mathrm{SC}_{\rm nf}$ has
been proved.

The general implication asked in F(c), with no cover
hypothesis, remains open in this return. The precise next
map after spreading already occurs over $K=k(B)$ for a
hyperbolic affine curve $B$ over a number field $k$:
For a smooth proper hyperbolic family
$X\to B$, a generic section gives
\[
 r:E_B\longrightarrow E_X,
 \qquad p r=\mathrm{id},
 \qquad E_{X_\eta}=E_X\times_{E_B}G_{k(B)}.
\]
Every such $r$ satisfies (F1)(b) and has non-open image. A
geometric reconstruction of this particular $r$ as a map
$B\to X$ is exactly a section of the family and supplies the
generic rational point. The curve-target dichotomy of
Proposition~\ref{F-local-dichotomy} does not provide that
higher-dimensional reconstruction.
\end{statement}
\begin{proof}
The first implication and its ranges are Bresciani, Theorem B
and Corollaries A--C\cite{F-BrescianiFiniteType}. We checked
the proof: pass to a finite cover and field extension carrying
the map to a constant abelian variety; spread over a base
curve; use the coherence of the constant abelian family and
pull it back along a finite morphism by Hilbert irreducibility.
A section neighborhood whose compactification has no rational
point supplies the contradiction. The cover-to-tripod range
maps after a further cover to the constant elliptic curve of
$j=0$. The general Theorem A in that paper also uses weak
Bombieri--Lang; that additional assertion is not taken as an
established input here. Sa\"idi's earlier relative-Jacobian
approach has an additional Tate--Shafarevich-module condition;
we do not remove it by renaming the specialization problem.

For the displayed map, Bresciani's Corollary 3.4 was checked
with its proof. Inertia at a point of $B$ would give a
cyclotomic map $T\to P_{X_b}$. Its image can be detected in
the abelianization of a suitable finite cover, where weights
$-2$ and $-1$ exclude it. Thus the generic section descends
to the actual splitting $r$ of the relative fundamental-group
sequence. The fibre-product formula is the smooth hyperbolic
fibre sequence, so base change back to $G_{k(B)}$ recovers the
original section.

To prove (b), let $J\subset E_B^\circ$ be open. A geometric
centralizer of $r(J)$ projects to a centralizer of $J$ in
$P_B$, which is zero by the open case. It consequently lies
in the fibre group $P_{X_{\overline\eta}}$. Pull $J$ back
under $q:G_{k(B)}\to E_B$. Its intersection with the
cyclotomic kernel contains $H_L$ for a finite extension
$L/k(B)$. The proper-curve case of
Proposition~\ref{F-centralizer}, applied to the generic
section over $L$, kills the remaining centralizer. The
arithmetic centralizer is killed by the same projection
argument and the arithmetic centralizer theorem. On the
other hand $r(E_B)$ meets the infinite fibre group trivially,
so it has infinite index in $E_X$.

A morphism $g:B\to X$ inducing $r$ has $p g$ inducing the
identity on $E_B$. Curve Hom injectivity gives $p g=\mathrm{id}_B$,
and its generic point is the required $k(B)$-point. In the
other direction a generic rational point extends to a section
of this proper family over the regular curve $B$ by the
valuative criterion; its group map $r_g$ satisfies $r_gq=rq$ on $G_{k(B)}$,
and $q$ is surjective, so $r_g=r$. Thus this identifies
the missing reconstruction precisely.

Localizing the generic field in a finite $p$-adic extension
turns its section into a point by the supplied local theorem.
It gives a point of $X(k')$, not a proof that the point lies
in $X(K)$. The field inclusion induces the actual injection
$X(K)\to X(k')$; even on the tripod this map is not onto,
as a $p$-adic element transcendental over the countable field
$K$ gives a missing point. No extension of its section back
to $G_K$ is claimed in this example. It shows exactly why
this base-change step alone has not constructed the desired
rational point. The unresolved compatibility is the family
section reconstruction just displayed.
\end{proof}

\begin{statement}{F.37}{Proved}{the birational part of field reduction}
\label{st:F.37}
\label{F-birational-field-reduction}
The number-field assertion (F2) implies the one-variable
birational assertion (F4)(a) over every finitely generated
field of characteristic zero. Therefore, under this exact
implication of assertions, every section over such a field
which already has a birational lift is a point section or
cuspidal. The remaining lifting map is still the quotient
$G_{K(U)}\twoheadrightarrow E_U$; the comparison does not
manufacture a lift of an arbitrary section.
\end{statement}
\begin{proof}
Proposition~\ref{F-birational-equivalence} sends the universal
number-field (F2) assertion to its one-variable birational
assertion by the compatible fixed-pro-point construction over
all open subcurves. Sa\"idi--Tyler, Theorem C
\cite{F-SaidiTyler}, sends the birational assertion over number
fields to the same assertion over all finitely generated
extensions of $\mathbb Q$. We opened their primary paper and
checked the exact quantifier and the reduction in section 5.3:
reduce to the projective line; select a characteristic cover of
a constant punctured projective line whose compactification is
hyperbolic; form the subgroup generated by that characteristic
geometric subgroup and the given section. This supplies an
isotrivial hyperbolic birational neighborhood. The isotrivial
relative-Jacobian theorem supplies the needed finite
Tate--Shafarevich group in their argument. Their Theorem C
therefore has no residual Tate--Shafarevich assumption.
Finally a birational lift belongs to the rational-place
decomposition group by that assertion, and its image in
$E_U$ is the point or cusp section according to the location
of the place. These are precisely the quotient maps used in
Proposition~\ref{F-birational-equivalence}.
\end{proof}

\begin{statement}{F.38}{Known}{what inductive cuspidality gives}
\label{st:F.38}
\label{F-inductive-cuspidality}
Let $X=X_n\to\cdots\to X_0=\operatorname{Spec}K$ be a
hyperbolic polycurve over a finitely generated characteristic-zero
field. Write $\mathcal S_X$ for its section classes and
$C_{X,f}$ for the recursively defined set containing every
section whose projection is cuspidal, together with fibre
cuspidal sections above rational base points. The universal
curve section assertion over $K$ is equivalent to independence
of $C_{X,f}$ from the chosen tower, with the known
point/packet injections understood. For a normal-crossings
compactification $X\subset Z$, let $C_{X,Z}$ be the sections
coming from punctured henselian neighborhoods of rational
boundary points, and let $C_X$ be their closure. The checked
comparison is
\[
 C_X\subset\overline{C_{X,f}}.
\]
It gives no reverse inclusion or birational lift of a section
merely because its projection to an affine base is cuspidal.
Thus it does not complete the remaining case of F(e)(4).
\end{statement}
\begin{proof}
These are Carlson--Corwin's Proposition B and the proposition
labelled \texttt{prop:fibinclusion} in the attached November
2026 source\cite{F-CarlsonCorwin}; we read the entire
fibration section and its proofs. The forward implication is
induction using the actual point-fibre pullback: a section
whose base projection is not in the recursively selected cusp
set belongs to a fibre over a rational point. The curve
assertion then exhausts that fibre's sections. Consequently
$C_{X,f}=\mathcal S_X\setminus X(K)$, independently of $f$.
For the converse take a hyperbolic curve $B$ and the tripod
$T$, and use the two projections of $B\times T$. Their
section-set product and the exact two recursive subsets are
\[
\begin{split}
 \mathcal S_{B\times T}&=\mathcal S_B\times\mathcal S_T,\\
 C_{B\times T,\operatorname{pr}_B}
   &=(C_B\times\mathcal S_T)\cup(B(K)\times C_T),\\
 C_{B\times T,\operatorname{pr}_T}
   &=(\mathcal S_B\times C_T)\cup(C_B\times T(K)).
\end{split}
\]
For $s\in\mathcal S_B$ and any $c\in C_T$, equality of the
last two subsets puts $s$ in $B(K)\cup C_B$, proving the
converse with its actual projection maps.

For the inclusion, choose a compatible smooth compactification
of the tower by resolving the proper compactification of the
last fibration. A rational boundary point maps either to a
boundary point of the compactified base, or to a rational
interior point of the base. In the first case its section
projects into the recursively selected boundary set. In the
second it lies on a relative cusp, and the punctured local
fundamental group is that of the cusp in the fibre: the
relative boundary is \'etale with multiplicity one. This
gives $C_{X,Z}\subset C_{X,f}$ for that compactification;
taking closures gives the display. The proof begins with an
actual local boundary section and maps it into the recursive
set. Starting with an arbitrary section in the first term
$f_*^{-1}(C_B)$ supplies none of that local data. The
unconstructed reverse step is exactly the one needed to
extend Proposition~\ref{F-proper-polycurve-lifting} to an
affine intermediate base.
\end{proof}

\begin{statement}{F.39}{Proved}{explicit parameter transport and its GT-coordinate formula}
\label{st:F.39}
\label{F-tangent-transport}
Fix the rational tangent of parameter $1$ at $0$ on
$T=\mathbb P^1\setminus\{0,1,\infty\}$. Every parameter
$t\in\mathbb Q^*$ has the following exact characterization
using a specified four-puncture diagram. Put
\[
 W_t=\mathbb P^1\setminus\{0,1,t,\infty\},\qquad
 p(z)=z,\quad q(z)=z/t,
 \qquad p,q:W_t\longrightarrow T.
\]
Repeated punctures at $t=1$ are simply omitted. There is a
unique class in the cusp packet at $0$ of $W_t$ whose image
under $q_*$ is the parameter-$1$ class; its image under
$p_*$ is the parameter-$t$ class of $T$. In packet
coordinates the actual maps are
\[
 p_*:c\longmapsto c,\qquad
 q_*:c\longmapsto c-\kappa(t).
\]

In a normalized tripod presentation
$P_T=\langle x,y\rangle^{\wedge}$ with the parameter-$1$
section and
\[
 \rho_\sigma(x)=x^{\chi(\sigma)},\qquad
 \rho_\sigma(y)=f_\sigma^{-1}y^{\chi(\sigma)}f_\sigma,
 \qquad f_\sigma\in[P_T,P_T],
\]
changing that section by a cusp cocycle $c$ gives exactly
\[
 f^{\,c}_\sigma=f_\sigma x^{-c(\sigma)},\qquad
 (f^{\,c}_\sigma)^{\rm ab}=(-c(\sigma),0).
\]
Thus inserting the already specified Kummer cocycle
$\kappa(t)$ into these relations recognizes its parameter
class. Neither construction removes its input: the first
retains the four-puncture diagram with cross-ratio $t$, and
the second retains the Kummer character of $t$ and the
chosen origin. They do not yet answer the intrinsic
single-tripod GT-recognition question left in
Proposition~\ref{F-open-tangents}.
\end{statement}
\begin{proof}
At $0$ both maps have ramification degree one. The first local
coordinate has leading coefficient $1$, and the second has
leading coefficient $1/t$. The tangent convention of
Proposition~\ref{F-torsors} therefore gives the two displayed
Kummer maps. In particular $q_*$ identifies the actual local
decomposition extensions, on both their inertia and their
$G_{\mathbb Q}$ quotient. Its inverse carries the chosen
parameter-$1$ section to the unique local section with class
$c=\kappa(t)$. The cusp-packet injection then gives uniqueness
in the ambient section set as well.

For the second formula write the new section as
$s_c(\sigma)=x^{c(\sigma)}s_1(\sigma)$. Its action is
$\operatorname{Int}(x^{c(\sigma)})\rho_\sigma$. It has the
same image of $x$, and its image of $y$ is
\[
 x^{c(\sigma)}f_\sigma^{-1}y^{\chi(\sigma)}
                   f_\sigma x^{-c(\sigma)}
 = (f_\sigma x^{-c(\sigma)})^{-1}y^{\chi(\sigma)}
                   (f_\sigma x^{-c(\sigma)}).
\]
This proves the exact transformation. Taking abelianization
proves the last formula. Existence of an inertia conjugation
for which $f^{\,c}_\sigma x^{\kappa(t)(\sigma)}$ has zero
abelianization for every $\sigma$ is precisely
$[c]=[\kappa(t)]$; changing roots changes the cocycle by the
same coboundary. This is a useful coordinate test, but its
explicit input character is exactly the data that an intrinsic
recognition rule would have to recover, so no stronger status
is assigned to it.
\end{proof}

\begin{statement}{F.40}{Refuted}{higher-rank valuative uniqueness}
\label{st:F.40}
\label{F-valuation-nonuniqueness}
The uniqueness of a valuation normalized through its inertia,
and the inference from inertia normalization to preservation
of that valuation, fail even with nonzero inertia and residue
field exactly $\mathbb Q$. There are distinct valuations
$v_0,v_1$ of $L=\mathbb Q(z,x,y)$ and one section
$s:G_{\mathbb Q}\to G_L$ such that
\[
 \Gamma_{v_a}=\mathbb Z\oplus_{\rm lex}\mathbb Q,
 \quad\kappa(v_a)=\mathbb Q,
 \quad I_{v_0}=I_{v_1}=T\ne1,
 \quad s(G_{\mathbb Q})\subset N_{G_L}(I_{v_0})
                             =N_{G_L}(I_{v_1}),
\]
but $s(G_{\mathbb Q})\not\subset D_{v_1}$. Accordingly the
higher-dimensional version of (F4)(a) must retain actual
valuation preservation. Its written uniqueness with respect to
inertia normalization is false; the example does not assert
nonuniqueness among valuations actually preserved by this section.
\end{statement}
\begin{proof}
This is F149-ZD5, whose full construction F149-D3 we checked.
Let $w$ be the $z$-adic valuation, with residue field
$\mathbb Q(x,y)$. In the Hahn field $\mathbb Q((t^{\mathbb Q}))$
define two restrictions $r_a$, for $a=0,1$, by
\[
 x-a=t,\qquad y=2+\sum_{j\ge1}t^{j+1/j!}.
\]
Their residue field is $\mathbb Q$ and their value group is
$\mathbb Q$. To verify the latter, form the norm polynomial
of $Y-2-\sum_{j\le m}X^{j+1/j!}$ from
$\mathbb Q(X^{1/m!})$ to $\mathbb Q(X)$. At the displayed
series its identity factor has value
$m+1+1/(m+1)!$, and the other factors have total value in
$(1/m!)\mathbb Z$. Induction therefore puts $1/(m+1)!$ in
the value group. Conversely every series exponent of a
rational expression belongs to the union of those groups.
Unbounded value-group index proves that $y$ is transcendental
over $\mathbb Q(t)$, so this really defines valuations of
$\mathbb Q(x,y)$. Their centers are $(0,2)$ and $(1,2)$.

Compose each $r_a$ with $w$. The divisible final convex
subgroup is $\mathbb Q$, so
Proposition~\ref{F-valuation-structure} gives
$I_{v_a}=I_w=T$ and $N(I_{v_a})=D_w$, as actual subgroups
for prolongations refining the same $w$. Construct a section
of $G_{\mathbb Q(x,y)}\to G_{\mathbb Q}$ along the rational
flag $y=2,x=0$, then lift it through a splitting of the
$w$-decomposition sequence. Its image lies in $D_w$ and
normalizes both inertia groups.

Restricting this section to $\mathbb Q(x)$, then to the
fundamental group of
$\mathbb P^1\setminus\{2,3,\infty\}$, gives the point
section at $0$. Containment in $D_{v_1}$ would instead give
the point section at $1$, since $r_1(x-1)>0$. Their distinctness
follows from the injective point-section map; equivalently
use its two Kummer coordinates. Thus the containment fails.
This verifies both refutations by the stated restriction maps.
\end{proof}

\begin{statement}{F.41}{Refuted}{an explicit abstract action}
\label{st:F.41}
\label{F-discontinuous-action}
There is an abstract semilinear action of $G_{\mathbb Q}$ on
an algebraic universal covering $\mathcal U$ of the tripod
which is not functionwise continuous on any open subgroup.
It admits no finite-type quotient $\mathcal U\to W$ which
is a connected pro-\'etale covering and whose induced
semilinear action is functionwise continuous on an open
subgroup of $G_{\mathbb Q}$. Thus the finite-form descent
step in Proposition~\ref{F-continuous-action} is false
without continuity.

This action nevertheless has the unique fixed pro-point
asserted in (B$'$), even for $H_{\mathbb Q}$. It proves
that the forgetful map from the continuous actions to
abstract actions misses actions on this same covering, and
that the specified finite-form reduction does not extend
to them. Fixed-point existence for every abstract action
is not proved or refuted by this example.
\end{statement}
\begin{proof}
Choose distinct odd primes $p_j$. Restriction to the
multiquadratic extension gives a continuous surjection
\[
 q:G_{\mathbb Q}\longrightarrow
                  A=\prod_{j\ge1}\mathbb F_2.
\]
The independence follows from unique prime factorization:
no nonempty finite product of these primes is a rational
square. Let $D=\bigoplus_{j\ge1}\mathbb F_2\subset A$.
Choose an $\mathbb F_2$-linear functional
$\lambda_0:A/D\to\mathbb F_2$ which takes the class of
$(1,1,\ldots)$ to $1$, by extending that nonzero vector
to a basis. Put $\lambda=\lambda_0q$, with the quotient
by $D$ understood. Its kernel is dense because it contains
the inverse image of the dense subgroup $D$ and $q$ is
an open quotient. It is not open: a continuous functional
on $A$ vanishing on $D$ vanishes identically. Its
restriction to every open subgroup of $G_{\mathbb Q}$
is nonzero and has dense kernel.

Put $T=\mathbb P^1\setminus\{0,1,\infty\}$ and fix the
rational interior point $x=1/2$. Choose a geometric point
above it on the universal covering and write $s$ for the
resulting continuous point action of $G_{\mathbb Q}$.
The rational involution $a:z\mapsto1-z$ has a unique lift
$\widetilde a$ fixing the chosen pro-point: two such lifts
differ by a deck transformation fixing a geometric fibre
point, hence agree. Its square is the identity, and it
commutes with $s(G_{\mathbb Q})$, since conjugating the
lift by $s(\sigma)$ gives another lift of the same rational
involution fixing the same point. Define
\[
 r(\sigma)=\widetilde a^{\lambda(\sigma)}s(\sigma).
\]
Commutation proves directly that $r$ is an action. It is
semilinear with the original action on $\overline{\mathbb Q}$.
On the rational coordinate $z$ it acts by $z$ or $1-z$,
so its stabilizer is $\ker\lambda$; this is not open even
after restriction to any open subgroup. Every quadratic
extension of $\mathbb Q$ lies in a cyclotomic extension
(for these prime radicals one can use the quadratic Gauss
sum in $\mathbb Q(\mu_{4p_j})$). Thus $\lambda$ vanishes
on $H_{\mathbb Q}$ and $r|_{H_{\mathbb Q}}=s|_{H_{\mathbb Q}}$.
Proposition~\ref{F-fixed-dictionary} gives its unique fixed
pro-point, which is also fixed by $r(G_{\mathbb Q})$.

We prove the stronger quotient assertion, rather than
just exhibit a noncontinuous coordinate. Write $\Omega$
for the function field of $\mathcal U$. Suppose that an
open subgroup $G'\subset G_{\mathbb Q}$ preserves a
finite-type quotient field $M\subset\Omega$ and that
its induced action on $M$ is functionwise continuous.
For $f\in M$, the two maps
\[
 G'\longrightarrow\Omega,\qquad
 \sigma\longmapsto r(\sigma)f,\quad
 \sigma\longmapsto s(\sigma)f
\]
are continuous when $\Omega$ is discrete. They agree on
the dense subgroup $G'\cap\ker\lambda$, so they agree
everywhere. Choose $\sigma\in G'$ with $\lambda(\sigma)=1$.
Commutation then gives $\widetilde a f=f$. Consequently
$M\subset\Omega^{\widetilde a}$.

If $\mathcal U\to W$ were a connected pro-\'etale
covering with function field $M$, $\widetilde a$ would
be a deck automorphism fixing its chosen geometric point.
The action of the deck group of a connected pro-\'etale
cover on a geometric fibre is free: this is the usual
unique path-lifting statement on every finite \'etale
level, followed by the inverse limit. Hence
$\widetilde a=1$, contradicting its action $z\mapsto1-z$.
No such quotient exists. This gives the exact obstruction
to the specified reduction and leaves the universal
abstract fixed-point assertion with its stated open status.
\end{proof}

\begin{statement}{F.42}{Proved}{lifting above a base cusp with good marked reduction}
\label{st:F.42}
\label{F-good-boundary-lifting}
Let $U\to B$ be a hyperbolic curve fibration over a hyperbolic
curve over a finitely generated characteristic-zero field $K$.
Let $c\in(\overline B\setminus B)(K)$. Suppose that near $c$
the fibration has a smooth proper relative compactification
with finite \'etale boundary. Write $Y_c$ for its hyperbolic
special fibre with that boundary deleted.

There is the following exact extension comparison, after a
choice of specialization paths:
\[
 E_U\mathbin{\times}_{E_B}D_c
       \simeq E_{Y_c}\mathbin{\times}_{G_K}D_c.
 \tag{F.42.1}
\]
Consequently every section of $E_U\to G_K$ whose base
projection is in the packet at $c$ corresponds to a pair
consisting of that same cusp section and a section of
$E_{Y_c}\to G_K$. When the latter is a point or cusp
section, the original section has a birational lift to
$G_{K(U)}$, constructed by a rank-two valuation and its
specified inertia coordinates.

In particular the universal curve-lifting assertion (F4)(b)
over finitely generated fields implies (F4)(b) for all such
surfaces whose fibration has good marked reduction at every
rational base cusp after every finite extension. This extends
Proposition~\ref{F-proper-polycurve-lifting} to this affine-base
range. The remaining affine-base case includes genuinely bad
marked reduction at the cusp selected by the projected section.
\end{statement}
\begin{proof}
Choose a parameter $t$ at $c$ and a prolongation of its
valuation. The cusp decomposition extension $D_c$ identifies
with $G_{K((t))}$: both have kernel $T$ and quotient $G_K$,
and compatible roots of $t$ identify their kernels and the
uniformizer splittings. The comparison of the global and
local relative fundamental-group sequences gives
\[
 E_{U_{K((t))}}\xrightarrow{\sim}
            E_U\mathbin{\times}_{E_B}G_{K((t))}.
\]
On the fibre group this map is the characteristic-zero
geometric base-change isomorphism; on the quotient it is the
identity. The relative fibre sequence is injective for a
hyperbolic fibration, so these two facts prove the displayed
isomorphism of extensions.

Marked good reduction identifies that fibre group with
$P_{Y_c}$ and makes its outer $G_{K((t))}$-action the
composite of the specialization map to $G_K$ with the
outer action of $Y_c$. This is the smooth marked
specialization comparison checked in
Proposition~\ref{F-reduction}; equivalently it follows from
the Kummer-\'etale comparison for the smooth pair, with
its horizontal boundary \'etale over the trait.
Here the residue characteristic is zero, so all geometric
finite covers are in its tame range.

We spell out why outer action suffices at this point.
For a centre-free profinite group $P$, an extension $E$
of a profinite group $G$ by $P$ is identified with
\[
 E\longrightarrow
 \operatorname{Aut}(P)\mathbin{\times}_{\operatorname{Out}(P)}G,
 \qquad e\longmapsto(\operatorname{Ad}(e)|_P,p(e)).
\]
Its kernel is $Z(P)=1$. For surjectivity, choose a lift of
the specified element of $G$ and correct its automorphism
by the unique required element of $P$. Compactness gives a
homeomorphism. Applying this to the centre-free group
$P_{Y_c}$ and the just identified outer action proves
(F.42.1), as an actual map of extensions after the stated
path choices. A section whose projection is the chosen
cuspidal section lands in the left side, and its two right
components give exactly the asserted pair. Conjugacies by
the geometric fibre act simultaneously through this map.

Suppose first that the residue section is a point section
at $y\in Y_c(K)$. In the smooth surface model around $y$,
complete $t$ by a local fibre parameter $z$. The flag
valuation with successive parameters $t,z$ has value group
$\mathbb Z\oplus_{\rm lex}\mathbb Z$ and residue field
$K$. Its decomposition sequence splits by
Proposition~\ref{F-valuation-structure}. In its image in
$E_U$, the $z$-inertia is killed, since $y$ is an interior
point of the special fibre, while the $t$-inertia is exactly
the base cusp inertia. Adjust a chosen splitting by the
cocycle in its $t$-coordinate that gives the prescribed
base packet member. The image of this adjusted splitting
has the specified two components in (F.42.1), hence is the
original section.

If the residue section is cuspidal at a relative cusp
$y\in(\overline Y_c\setminus Y_c)(K)$, choose $z$ as the
local equation of that horizontal boundary. The two
boundary divisors have normal crossings, and the local
Kummer maps send the inertia of the same rank-two valuation
to the $t$ and $z$ inertia coordinates. Each is a copy of
$T$, with cyclotomic action. The prescribed two packet
members differ from the images of a chosen decomposition
splitting by two cocycles in $Z^1(G_K,T)$. Twisting that
splitting by this ordered pair gives their simultaneous
lift. These maps are the same coordinate root maps used
in Proposition~\ref{F-proper-polycurve-lifting}; no section
with domain $G_{K(B)}$ has been inferred from one with
domain $G_K$.

Finally universal curve lifting gives the universal curve
point-or-cusp assertion by
Proposition~\ref{F-cuspidalization}. Project a section to
$B$. A point projection is handled by its actual fibre
and the full-flag construction, without a properness
requirement on $B$. A cusp projection is handled by
(F.42.1) and the preceding construction. This proves the
stated implication of lifting assertions in the additional
range. At bad marked reduction the residue extension of
Proposition~\ref{F-reduction} need not be that of a smooth
hyperbolic curve; indeed its checked equivalence excludes
such a realization after every finite extension when
potentially good marked reduction fails. Thus applying
curve (F2) to an imagined smooth residue fibre would not
supply the remaining step. A fixed component or compatible
logarithmic valuation for that extension is still required.
\end{proof}

\subsection{The added LM5 comparison: tubes, inertia, and actual sections}

\begin{statement}{F.43}{Proved}{the normal-tube comparison with FLAGGED\_149}
\label{st:F.43}
\label{F-normal-tube-comparison}
Let $(X,D)$ be a smooth normal-crossings pair over a characteristic-zero
field, $U=X\setminus D$, and $S=D_J^*$ a connected open stratum with
its $r$ branches labelled. Write $N_j$ for their normal lines on $S$,
$P_S=\prod_{j\in J,S}N_j^\times$, and $\mathcal T_S$ for the
punctured henselian normal tube. Paths and geometric points are fixed
in every displayed fundamental-group map.
\begin{enumerate}
\item Normal specialization identifies the finite-covering categories
of $\mathcal T_S$ and $P_S$. Its profinite geometric homotopy sequence is
\[
 \pi_2^{\rm et}(S_{\bar K})\xrightarrow{(c_1(N_j))_j}
 T^J\longrightarrow\pi_1(\mathcal T_{S,\bar K})
 \longrightarrow\pi_1(S_{\bar K})\longrightarrow1.
 \tag{F.43.1}
\]
The first map means Chern transgression, evaluated on spherical
classes. The composite $\iota:T^J\to P_U$ kills its image. Thus
injectivity of the full global inertia implies left exactness in
(F.43.1). For a relative cusp of a hyperbolic smooth pair, the tube
extension maps isomorphically onto the cusp normalizer extension
of Proposition~\ref{F-torsors}, with extension class $c_1(N)$.

\item At a geometric boundary point, $\iota:T^J\to P_U$ is injective
precisely when every finite covering germ is an open and closed
sub-germ of the restriction of a finite global covering. If $R_{\bf n}$
is the finite root stack with orders $n_i$, and $M_{\bf n}$ is the
closed normal subgroup of $P_U$ generated by the $n_i$th powers of
boundary inertia, then
\[
 \pi_1(R_{{\bf n},\bar K})=P_U/M_{\bf n},\qquad
 \ker(T^J\to\pi_1(R_{{\bf n},\bar K}))
       =\iota^{-1}(M_{\bf n}),\qquad
 \bigcap_{\bf n}\iota^{-1}(M_{\bf n})=\ker\iota.
 \tag{F.43.2}
\]
In particular a finite root stack can enlarge this kernel.

\item Horizontal cusp inertia injects in a tower of punctured
hyperbolic smooth pairs. In the recursively compatible stable
toroidal range of LM5 F.39, the full crossing inertia injects as
well. The hypothesis here retains the stable incidence data:
the kernel of the local contact matrix is the fibre puncture
or the essential vanishing cycle. On split branches the Galois
action is cyclotomic; in general it also permutes the factors.
A monomial map $z'_j=u_j\prod_i z_i^{m_{ji}}$ induces the matrix
$(m_{ji})$ on these inertia groups.

\item For $y^2=x^3+t$, the model
\[
 [\{Y^2=X^3+1\}\times\mathbb A^1_s/\mu_6],\qquad
 \zeta(X,Y,s)=(\zeta^{-2}X,\zeta^{-3}Y,\zeta s),
 \tag{F.43.3}
\]
has open part identified with the original family by
$t=s^6$, $x=s^2X$, $y=s^3Y$. Its normal meridian injects;
the full corner meridian does too when the infinity section
is retained. On the ordinary embedded cusp resolution the
two exceptional-tail crossing maps instead have respective
kernels $\widehat{\mathbb Z}(-3,1)$ and
$\widehat{\mathbb Z}(-2,1)$. These are maps into the same
open fundamental group, with the explicit substitutions above
and in the proof.

\item The maps in (F.43.1), restricted to a relative cusp or a
trait meeting that cusp, are exactly the maps in FLAGGED\_149
C3, C4, C6, C8, and C9. C1, C2, and C5 give the corresponding
extension tests for an \emph{already given rational geometric
section}. Their geometric premise is retained; the tests do not
construct a geometric point from an arbitrary section.
\end{enumerate}
This supplies the added comparison of LM5 F.21--F.23, F.35,
F.39 with FLAGGED\_149 C1--C6, C8, C9.
\end{statement}
\begin{proof}
We read the cited return proofs directly. Here are the actual
comparison maps and the kernel checks, so that equality of names
for an inertia group is not used in place of a comparison.

In a strict normal-crossings chart the covers are dominated by
adjoining $z_j^{1/n}$. On an overlap
$z'_j=u_jz_{\sigma(j)}$, their root coordinates are identified by
roots of $u_j|_S$. These are the transition torsors of the normal
lines. Normalization after the root adjunction is unramified,
and descent for finite covers makes this local construction an
equivalence with the covering category of $P_S$. The same
transition data in the inverse system of root gerbes gives the
profinite homotopy comparison: its fibre is $BT^J$, its twisting
classes are the Kummer classes of $N_j$, and its boundary is
their Chern transgression. This is the normal-bundle instance
of the log comparison of Carchedi--Scherotzke--Sibilla--Talpo,
Theorem A and Corollary 5.5 \cite{F-CSST}; their fine saturated,
log-regular, locally noetherian hypotheses hold for this smooth
normal-crossings pair. In characteristic zero the comparison
uses all primes. This justifies the higher homotopy sequence,
including its finite-covering truncation.

Restriction of a global cover to the tube is the functor
giving $\pi_1(\mathcal T_S)\to\pi_1(U)$. Its composite with
the fibre map kills the Chern boundary by (F.43.1).
For a hyperbolic relative cusp $q$, the fibre inertia already
injects in $P_U$. The tube map is consequently injective on
that fibre, and its image is $N_{E_U}(I_q)$: it induces the
identity on the base quotient, and both geometric kernels
are $I_q$ by peripheral self-normalization. The resulting
map of short exact sequences is an isomorphism. In frames
$e_j=g_{ij}e_i$ of $N$, root coordinates satisfy
$r_i=g_{ij}^{1/n}r_j$. Their triple-overlap defect is
$g_{ij}^{1/n}g_{jk}^{1/n}g_{ki}^{1/n}$, the positive Kummer
cocycle of $N$. This is exactly FLAGGED C6 and the sign in
Proposition~\ref{F-torsors}; dual conormal frames reverse it.

For the covering-germ criterion, more generally let
$H\to G$ be a homomorphism of profinite groups. If it is
injective and $H/K$ a finite transitive $H$-set, compactness
gives an open normal $M\lhd G$ with $M\cap H\subset K$.
Then $J=KM$ is open in $G$, $J\cap H=K$, and
$hK\mapsto hJ$ embeds $H/K$ into $G/J$ as an $H$-orbit.
Disjoint unions treat every finite $H$-set. Conversely a
nonidentity kernel element is detected by a finite quotient
of $H$ and acts trivially on every restricted $G$-set, so
the regular set of that quotient cannot embed. Taking
$H=T^J$ proves the asserted equivalence, including the
qualification ``sub-germ''.

Finite covers of a finite root stack are the covers of $U$
whose inertia orders divide the specified root orders. In
the chart $[\operatorname{Spec}A[w_i]/(w_i^{n_i}-z_i)/
\prod\mu_{n_i}]$, this is exactly unramifiedness on the
covering chart. Descent proves (F.43.2). Every finite
quotient of $P_U$ has finite inertia orders, so some
$M_{\bf n}$ lies in its kernel; hence their intersection
is trivial. For $(\mathbb A^1,0)$ the map before taking
roots is $T\xrightarrow{1}T$, whereas the finite-root
map is $T\to\mu_n$ with kernel $nT$. For a finite Kummer
\emph{cover}, the local lattice maps instead by a finite-index
inclusion $T'\hookrightarrow T^J$. A nonzero kernel $A$
persists because $mA\subset A\cap T'$ is nonzero whenever
$mT^J\subset T'$. These are the precise maps relating
the finite-root and covering operations.

For a punctured hyperbolic fibration, the geometric fibre
group is a centre-free free profinite group. The homotopy
boundary from the base's $\pi_2$ is central in that fibre
group and therefore vanishes. This gives the injected
fibre sequence of FLAGGED C4; a peripheral word injects
there by the full peripheral theorem. One can also verify
the completion step directly. For a discrete free fibre
$F$, conjugation maps the total group to
$\operatorname{Aut}(\widehat F)$; completion extends that
map, and its restriction to $\widehat F$ is the injective
inner-automorphism map. Thus the completed fibre injects.
A nontrivial cyclic subgroup of $F$ has its full profinite
topology: complete its immersed labelled circle to a finite
rose cover in which that circle is embedded, then use its
primitive integral homology class and pass to cores.

At a stable node $xy=t$, the contact map is
$(a,b)\mapsto a+b$, with kernel generated by $(1,-1)$;
this generator maps to the vanishing cycle in the fibre.
That cycle is essential: an inessential side would be an
unmarked rational tree attached at one point, violating
stability by summing the special-point inequalities over
its components. At a marking the kernel is its peripheral
loop; away from either locus it is zero. Compatible
toroidal charts give the corresponding finite-index or
primitive sublattices. The diagram of local lattices over
the injected fibre sequence therefore proves full inertia
injection by induction, starting with the stipulated
injective base lattices. Smith normal form gives
$\ker\widehat M=\widehat{\ker M}$ for the integer contact
matrix, so this argument proves the profinite assertion
as well as the discrete one. The transformations of
$z_i^{1/n}$ give precisely the claimed cyclotomic and
permutation actions and the contact matrix. The units
change tangential origins by Kummer cocycles.

For (F.43.3), substitution proves the equation and the
action is free for $s\ne0$. The normal loop maps to six
times the generator under the projection to the $t$-line,
so its full profinite map is injective. At infinity the
other inertia generator is the cusp loop of the punctured
elliptic fibre; the preceding injected fibre sequence
proves rank-two injectivity. For comparison the three
ordinary blowup charts of $f=y^2-x^3$ give
\[
 f=x^2(v^2-x),\qquad f=u^2v^3(v-u),\qquad
 f=u^6w^3(w-1).
\]
The exceptional tails have squares $-3,-2$ and their
normal loops $b,a$ satisfy $h=b^3=a^2$, where $h$ is
the third exceptional loop. The crossing maps are
$(m,n)\mapsto b^{m+3n}$ and $(m,n)\mapsto a^{m+2n}$.
Both $a$ and $b$ have full procyclic order, already
detected by their respective degrees $3$ and $2$ on the
$t$-line. The displayed kernels follow. Blowing up a
smooth boundary point of the product compactification of
the tripod squared likewise gives $(a,b)\mapsto a+b$
at its new crossing, with kernel $T(1,-1)$. This explains
exactly why arbitrary compactifications are not among
the recursively compatible stable models used above.

Finally, the full FLAGGED comparison can be made directly
on sections. C1 extends a rational geometric section
across a subset of codimension at least two: resolve its
graph; projective-space fibres of the resolution map
constantly to a positive-genus curve, so the graph
descends through the normal base. In genus zero an
\'etale-local boundary section makes its complement
affine, and the normal Hartogs equality extends the
coordinate functions. The inverse image of a remaining
boundary component would be a Cartier divisor supported
in codimension at least two, hence is empty. Repeat by
actual fibre products for a tower.

For C2 and C3, properness first extends a given trait
point to the compactified curve. If its boundary equation
pulls back as $z=t^\mu u$, the map on Kummer roots is
$z^{1/n}\mapsto t^{\mu/n}u^{1/n}$; its inertia map is
$T\xrightarrow{\mu}T\hookrightarrow P_U$. It is
nontrivial exactly when $\mu>0$. C5 combines this
codimension-one test with purity: the kernel of
$G_{K(S)}\to E_S$ is the closed normal subgroup generated
by divisorial inertias. C1 then extends across the
remaining codimension-two set. C8 is the injection of
normal frames into the splitting torsor proved in
Proposition~\ref{F-torsors}. C9 reduces the same root map
modulo $t$; it gives the normal-line map $v\mapsto u v^\mu$
and the specialization $\kappa(\bar u)$ of
Proposition~\ref{F-specialization}. Replacing $t$ by $at$
changes it by $-\mu\kappa(\bar a)$. Thus all the
comparison squares commute on each finite root cover
and therefore on their inverse limits. Each extension
test starts with an actual rational section and applies
these explicit maps; it supplies no missing point in
the open assertions of Propositions~\ref{F-field-reduction}
and~\ref{F-inductive-cuspidality}.
\end{proof}

\begin{statement}{F.44}{Proved}{lifting at a potentially good base cusp}
\label{st:F.44}
\label{F-potential-boundary-lifting}
Let $K/\mathbb Q$ be finitely generated and let $U\to B$ be a
hyperbolic curve fibration over a hyperbolic curve. Let
$s:G_K\to E_U$ project to a cuspidal section $\tau$ at
$c\in(\overline B\setminus B)(K)$. Suppose the fibre over
$K((t))$, for a parameter $t$ at $c$, has potentially good
marked reduction. Let $Y_\tau/K$ be the smooth residue
curve obtained by the descent construction of
Proposition~\ref{F-reduction}; it has an identified arithmetic
extension
\[
 E_{Y_\tau}\simeq
 E_{U_{K((t))}}\mathbin{\times}_{G_{K((t))},\tau}G_K.
 \tag{F.44.1}
\]
Every $s$ whose induced section in (F.44.1) is a point or
cusp section has a birational lift $G_K\to G_{K(U)}$.
The lift can be obtained through the decomposition group
of a rank-two valuation of $K((t))(U)$ with residue $K$.

Consequently universal curve lifting implies surface
lifting for hyperbolic fibrations with potentially good
marked reduction at every rational base cusp after finite
constant extension. This extends the good-reduction range
of Proposition~\ref{F-good-boundary-lifting}; the implication
is an implication between the two universal assertions,
not an assertion of universal curve lifting.
\end{statement}
\begin{proof}
Put $F=K((t))$ and $L_F=F(U_F)$. The comparison of relative
fibre sequences used in Proposition~\ref{F-good-boundary-lifting}
identifies $E_{U_F}$ with $E_U\times_{E_B}D_c$ and
$D_c$ with $G_F$. In particular the section $s$ has the
specified lift to the pullback along $\tau$. The exact
descent in Proposition~\ref{F-reduction} gives (F.44.1).
We first explain the valuation supplied by an actual
point or cusp of its residue curve, including compatibility
on finite coverings.

Choose a finite Galois extension $F_1/F$ over which the
stable marked model is smooth, and write $k_1$ for its
residue field. Its generic Galois isomorphisms extend
uniquely to that model. With $G=\operatorname{Gal}(F_1/F)$,
the homomorphism $\alpha:G_K\xrightarrow{\tau}G_F\to G$
defines exactly the descent datum used for $Y_\tau$.
After a finite residue extension killing $\alpha$, a
point or cusp of $Y_\tau$ is therefore a closed point
$y$ of the smooth special marked model, fixed by this
specified twisted descent datum.

At this point take the flag valuation: first the vertical
divisor, then the point on its smooth special curve.
It has value group $\mathbb Z\oplus_{\rm lex}\mathbb Z$
and a finite residue extension of $K$. For a cusp the
second parameter is its horizontal boundary equation;
for an interior point it is any parameter on the special
curve. Multiplying either parameter by a unit leaves
the valuation ring unchanged. Restriction through the
finite extension of function fields gives a rank-two
valuation $v$ of $L_F$, with finite residue over $K$.
Finite extension changes its value lattice only by finite
index.

The needed prolongation of $v$ to the maximal extension
unramified over $U_F$ is fixed by $s(G_K)$. This fact
requires the full comparison on covers, not just its
outer representation on the fibre group, so we specify
the construction. A finite \'etale cover of $U_F$,
after a further finite extension of $F_1$, has a smooth
marked model induced by the good-reduction comparison.
Indeed over the strict henselian good model its geometric
cover is the lift of the corresponding special-fibre
cover; the Kummer covers at its horizontal marks give
the compactification. In residue characteristic zero
all these ramification indices are invertible. Finite
presentation descends the model to a finite extension.
Refining two such choices by their common finite
extension and a common dominating cover makes the
specialization maps commute.

The section on $Y_\tau$ provides a compatible pro-point
on the completed special-fibre covering tower, by
Proposition~\ref{F-fixed-dictionary}. For an interior
point this is its ordinary fibre functor; for a cusp it
is the point fixed by its decomposition group. The
functors in the preceding paragraph carry its coordinates
to the points of the good special models of every
generic cover. The flag valuations at these points
restrict compatibly: the first valuation is the base
trait valuation, and the second is the discrete
point valuation on its smooth normalization. Their
compatibility can also be read on the local covers
$t^{1/a},z^{1/b}$; a change of unit preserves the
same valuation. Thus they define a prolongation
$\widetilde v$ of $v$ on the maximal unramified extension.
The equivariant comparison in (F.44.1) is constructed
from these same stable-model isomorphisms. Since the
chosen pro-point is fixed by the induced section,
$s(G_K)$ fixes $\widetilde v$.

Let $D_v\subset G_{L_F}$ be the decomposition group of
a prolongation to a separable closure. Restriction maps
$D_v$ onto the stabilizer of $\widetilde v$ in
$E_{U_F}$: at every finite Galois level a decomposition
group surjects onto the decomposition group below it,
and compactness gives the inverse-limit surjection.
Write $\overline D_v$ for its image. We have proved
$s(G_K)\subset\overline D_v$. If $k(v)/K$ is the
finite residue extension, the image of $D_v$ in $G_K$
is the subgroup $G_{k(v)}$, using the residue embedding
of algebraic constants. The inclusion of $s(G_K)$
therefore forces $k(v)=K$.

It remains to lift the specified section, rather than
merely its image subgroup. The full-rank valuation
sequence of Proposition~\ref{F-valuation-structure} is
\[
 1\longrightarrow I_v\simeq T^2\longrightarrow D_v
        \longrightarrow G_K\longrightarrow1,
 \tag{F.44.2}
\]
and it splits. Its image has kernel
$\overline I_v$, the image of $T^2$, over $G_K$.
The equal-characteristic-zero henselian calculation
applies to $L_F$ although it is not finitely generated
over $K$: it uses the coefficient-field embedding of
$K$, the finite-index rank-two value lattice, and
compatible roots of a basis of that lattice, or
equivalently the iterated henselian flag model.
These are the hypotheses used for the inertia and
splitting step of Proposition~\ref{F-valuation-structure}.
The geometric group $P_U$ is torsion-free: its exact
hyperbolic fibre sequence has torsion-free curve kernel
and quotient. Consequently $\overline I_v\subset P_U$
is torsion-free. For each prime $\ell$ it is a
torsion-free finitely generated quotient of
$\mathbb Z_\ell(1)^2$, hence free and a direct summand
quotient. Choosing these module splittings and taking
their product gives a continuous
$\widehat{\mathbb Z}$-linear section
\[
                       j:\overline I_v\longrightarrow I_v.
\]
Both modules have scalar cyclotomic Galois action, so
$j$ is $G_K$-equivariant.

Fix any splitting $r$ of (F.44.2). Relative to its
image, the prescribed section $s$ has the unique form
$s(g)=c(g)\overline r(g)$ for a continuous cocycle
$c:G_K\to\overline I_v$. Then
$g\mapsto j(c(g))r(g)$ is a continuous splitting of
(F.44.2) mapping to $s$. Composing it with
$G_{L_F}\to G_{K(U)}$, the restriction homomorphism
associated to the actual field embedding
$K(U)\hookrightarrow F(U_F)$, gives the required
birational lift. Compatibility of restriction with
the finite-cover functors proves that its image in
$E_U$ is the original section.

Finally universal curve lifting is equivalent to the
universal point-or-cusp assertion by
Proposition~\ref{F-cuspidalization}. Its application to
$B$ gives a point or cusp projection. A point projection
uses its actual fibre as in
Proposition~\ref{F-proper-polycurve-lifting}. A cusp
projection uses $Y_\tau$ and the construction just
proved. The remaining case is genuinely bad marked
reduction, for which Proposition~\ref{F-reduction}
excludes any such smooth residue curve even after
finite extension. No compatibility statement for that
case has been inferred here.
\end{proof}

\subsection{A global tangential relation using two open homomorphisms}

\begin{statement}{F.45}{Proved}{intrinsic recognition by two open Hom maps}
\label{st:F.45}
\label{F-two-Hom-tangents}
Let $K$ be a number field, $U=\mathbb P^1_K\setminus\{0,1,\infty\}$,
$E=\pi_1(U)$, and $\Delta=\ker(E\to G_K)$. Fix an actual lifted
cusp inertia $I_0\subset\Delta$ above $0$ and a continuous section
$s:G_K\to N_E(I_0)$. Its class is induced by a nonzero
$K$-rational tangent vector at $0$ precisely when the following
entirely group-theoretic relation holds.

There exist an open subgroup $H\subset E$ with image $G_K$, a cusp
inertia $J$ of $H$, a continuous section
$\sigma:G_K\to N_H(J)$ of the arithmetic projection, and two
continuous open homomorphisms over $G_K$,
\[
                       p,q:H\longrightarrow E,
\]
such that
\[
 \begin{gathered}
 p|_J:J\xrightarrow{\sim} I_0,\qquad
 q(j)=p(j)^2\quad(j\in J),\\
             [p\circ\sigma]=[q\circ\sigma]=[s]
       \quad\text{modulo conjugation by }\Delta.
 \tag{F.45.1}
 \end{gathered}
\]
Here a cusp inertia of $H$ is an intersection
$H\cap aI_i a^{-1}$ for one of the three ambient peripheral
classes and an actual geometric conjugator $a$. The peripheral
classes themselves are characterized by the cyclotomic-normalizer
criterion. Thus all the quantifiers in (F.45.1) refer to the one
ambient arithmetic group, its open subgroups, their continuous
homomorphisms, and their sections; the test does not first declare
any source section geometric.

The same relation after restriction to some finite extension
$L/K$ is equivalent to (F.45.1) over $K$. For every specified
$t\in K^\times$, one can moreover choose $H,J,p,q$ so that the
common-image relation in (F.45.1), with $\sigma$ in that cusp
packet, has the unique image $[s_t]$, the tangent with coordinate
$t$ in the standard coordinate $z$ at $0$.

This supplies a uniform global relation recognizing the geometric
splittings in F(g)(5), using two open Hom maps. It does not replace
those maps by a word involving only the automorphisms of the
tripod or only the two entries of a Grothendieck--Teichm\"uller
pair. That more restrictive formula remains the separate question
in Proposition~\ref{F-open-tangents}.
\end{statement}
\begin{proof}
We first prove necessity of tangentiality from the displayed
relation, without assuming a section conjecture. The open subgroup
$H$ is the arithmetic fundamental group of an actual connected
finite \emph{\'etale} cover $V\to U$. Its image onto $G_K$ makes
$V$ geometrically connected. The subgroup $J$ is the inertia of
an actual lifted cusp $c$ of $V$. Since $\sigma(G_K)$ normalizes
$J$, the underlying cusp is $K$-rational. These identifications
come from the finite-cover correspondence and the inertia
normalizer identity, so they do not reconstruct a curve from an
unclassified abstract group.

The relative open Hom theorem
(Proposition~\ref{F-open-Hom}, or
Proposition~\ref{A:relative-hom}) realizes $p$ and $q$, up to
their specified geometric inner ambiguity, by dominant
$K$-morphisms
\[
                          P,Q:V\longrightarrow U.
\]
They extend to morphisms of smooth proper completions. Because
$p(J)=I_0$ and $q(J)=I_0^2$ are nontrivial, both send $c$ to
$0$: a morphism extending to an interior point kills the local
inertia. The index of the image of cusp inertia is the positive
ramification degree. Thus $P$ has degree one locally at $c$, and
$Q$ has degree two locally there. In a $K$-rational uniformizer
$u$ at $c$, their completed local expressions are
\[
 P=a u(1+uA(u)),\qquad
 Q=b u^2(1+uB(u)),\qquad a,b\in K^\times.
 \tag{F.45.2}
\]
This also verifies the direction of the exponent-two comparison:
under the isomorphism $p|_J$, the actual map $q|_J$ is
multiplication by two, exactly as required in (F.45.1).

Let $\beta\in H^1(G_K,T)=\widehat{K^\times}$ be the parameter
of $\sigma$ relative to the tangent $u=1$. The compatible roots
of the units $1+uA(u)$ and $1+uB(u)$ have residue one. The local
root-cover calculation of Part B and
Proposition~\ref{F-specialization} therefore gives, additively,
\[
 [p\sigma]=\kappa_K(a)+\beta,
       \qquad [q\sigma]=\kappa_K(b)+2\beta
       \quad\text{in the packet at }0.
 \tag{F.45.3}
\]
Both images lie in that packet: normalizing $J$ normalizes its
nontrivial image, whose full arithmetic normalizer is
$N_E(I_0)$. Packet injectivity in $E$ converts the equality of
section classes in (F.45.1) into equality in (F.45.3). Subtraction
gives the exact formulas
\[
 \beta=\kappa_K(a/b),\qquad
                  [s]=\kappa_K(a^2/b).
 \tag{F.45.4}
\]
Thus the source parameter and the requested target parameter are
both ordinary $K$-rational tangential parameters. No division in
$\widehat{K^\times}$ and no local-to-global saturation claim is
needed; the two ramification degrees differ by one.

Conversely take $s=s_t$ for $t\in K^\times$. On the smooth open
curve
\[
 W=\mathbb P^1_K\setminus
       \bigl(\{0,1,\infty\}\cup\{z:z^2=t\}\bigr)
\]
there are the two dominant maps
\[
               P(z)=z,\qquad Q(z)=z^2/t
                       \quad\text{from }W\text{ to }U.
 \tag{F.45.5}
\]
Both have the cusp $0$, of local degrees one and two. The tangent
$u=z$ with value $t$ maps under the first to parameter $t$, and
under the second to $(1/t)t^2=t$. Its entire root-fibre-functor
section therefore has equal images $[s_t]$ under the two maps.

To place this source inside the one ambient group, use the marked
Belyi theorem over $K$. Precisely, Scherr--Zieve, Theorem 1.1
\cite{F-ScherrZieve}, gives a finite morphism
$b:\mathbb P^1_K\to\mathbb P^1_K$ with branch locus
$\{0,1,\infty\}$ and with every point of the finite complement
of $W$ mapped into that set. The theorem allows a finite
$G_K$-stable set of geometric points, so the roots of $z^2=t$
need not be $K$-rational. We opened the primary proof: its
construction is defined over $K$, and its additional ramification
condition at the prescribed points is stronger than what is
used here. Consequently
\[
 V=b^{-1}(U)\subset W,\qquad b:V\longrightarrow U
 \tag{F.45.6}
\]
is an actual geometrically connected finite \emph{\'etale} cover.
Deleting these additional points leaves the completed local ring
and the tangent at $0$ unchanged. Restrict (F.45.5) to $V$ and
use $b$ to identify $\pi_1(V)$ with an open subgroup $H\subset E$.
The two induced homomorphisms are open: each dominant map of
this particular pair has an explicit comparison. The first is
the open immersion $V\hookrightarrow U$. The second factors as
the open immersion
$V\hookrightarrow\mathbb P^1\setminus\{0,\infty,z^2=t\}$
followed by the finite \emph{\'etale} degree-two map $z\mapsto z^2/t$
to $U$. An open immersion of these normal curves induces a
surjection on fundamental groups, and the finite cover induces
an open injection. Both maps are over $G_K$. The cusp inertia $J$ and
the specified tangential section $\sigma$ now satisfy every part
of (F.45.1). Notice that the Belyi map may be ramified at $0$;
the local degrees one and two in the test concern $P,Q$, not
the inclusion of $H$ supplied by $b$.

For these fixed $H,J,p,q$, formula (F.45.3) has $a=1$ and
$b=1/t$. Equality forces $\beta=\kappa_K(t)$, hence has the
unique common image $[s_t]$. This proves the assertion for each
specified parameter and the global existential characterization
at once. Replacing any base paths conjugates the corresponding
subgroups and maps. If desired, the equality of section classes
can be made an equality of actual homomorphisms: a geometric
conjugator between two sections in this one packet normalizes
$I_0$ by cusp uniqueness, hence belongs to $I_0$ and acts
trivially on its inertia. Such adjustments preserve the two
inertia equations in (F.45.1).

We also verify descent for the germ version. Enlarge a finite
extension in which the test holds to a finite Galois extension
$L/K$. Relative to the standard tangent, let
$\alpha\in\widehat{K^\times}$ be the parameter of $s$.
The already proved direction over $L$ gives
$\operatorname{res}_{L/K}\alpha=\kappa_L(v)$ for $v\in L^\times$.
The left side is Galois-invariant. Injectivity of the ordinary
Kummer map over $L$ gives $\gamma(v)=v$ for every
$\gamma\in\operatorname{Gal}(L/K)$, so $v\in K^\times$.
Restriction
\[
 H^1(G_K,T)\longrightarrow H^1(G_L,T)
\]
is injective because $T^{G_L}=0$. Explicitly, a cocycle whose
restriction is zero after a coboundary adjustment has values
fixed by $G_L$, by the same cocycle identity used in
Proposition~\ref{F-centralizer}, and those values vanish.
Thus $\alpha=\kappa_K(v)$ and the test holds already over $K$
by the construction above.

For a small completely explicit Belyi witness, take $K=\mathbb Q$
and $t=2$. One can use
\[
 b(z)=2z^2-z^4,
 \qquad
 b^{-1}\{0,1,\infty\}
       =\{0,1,-1,\sqrt2,-\sqrt2,\infty\}.
\]
Its derivative is $4z(1-z^2)$, so its branch values are exactly
$0,1,\infty$. With $P=z$ and $Q=z^2/2$, this degree-four
cover supplies (F.45.1) and has unique common cusp image $[s_2]$.
The exact polynomial and local-degree checks are recorded in
\texttt{code/F/two\_hom\_tangents.py} and its output.

The manuscript anchor is cote 149, scan 85 and the specialization
construction of scans 85--95, together with F(g)(5) and the last
question of B(a). Our relation uses the full Hom theorem to
interpret open group maps and Belyi to put their sources inside
one arithmetic tripod. It invokes neither global section
surjectivity nor the stipulated local section theorem.
\end{proof}

\begin{statement}{F.46}{Proved}{the finite-etale restriction selects exactly two rational tangents}
\label{st:F.46}
\label{F-finite-etale-tangents}
In Theorem~\ref{F-two-Hom-tangents} over $K=\mathbb Q$,
impose the additional requirement that both $p$ and $q$ be
injective. The resulting relation holds exactly for
\[
                         [s_{16}],\quad [s_{-16}].
 \tag{F.46.1}
\]
More generally, over a number field $K$, every tangent parameter
selected by this restricted relation belongs to $16\mu(K)$,
where $\mu(K)$ is the group of roots of unity in $K$.
This conclusion still holds when the local degrees one and two
in that relation are replaced by any two unequal positive
degrees $d,e$: require that $p$ and $q$ send the selected cusp
inertia into $I_0$ with these respective indices, and keep the
common-image equation for the section. Every common image is
then an ordinary tangent in $16\mu(K)$. In particular, over
$\mathbb Q$ even the union of all these unequal-degree
finite-\'etale tests selects exactly the two classes in
(F.46.1). The condition $d\ne e$ is essential: taking $p=q$
gives an equal-degree relation satisfied by every member of
the corresponding packet.

The additional requirement is exactly that the two reconstructed
morphisms from the common finite cover to the tripod are finite
\'etale. In that case $q\circ p^{-1}:p(H)\xrightarrow{\sim}q(H)$
is an actual isomorphism between open arithmetic subgroups.
Thus the forgetful map from these finite-cover relations to the
open-Hom relations of Theorem~\ref{F-two-Hom-tangents} has a
strictly smaller set of common section images: $[s_2]$ is an
explicit witness. This refutes replacing the two maps in that
criterion by such commensurations while retaining its inertia
degrees and equal-image equation. It makes no assertion that
every different formula using only GT automorphisms has this
same restriction.
\end{statement}
\begin{proof}
We first justify the exact comparison between the group and
curve hypotheses. By the open Hom theorem, $p$ and $q$ are
induced by dominant maps $P,Q:V\to U$, where $V$ is the
finite \'etale cover attached to $H$ and
$U=\mathbb P^1\setminus\{0,1,\infty\}$. Suppose, for example,
that $p$ is injective. Its open image corresponds to a finite
\'etale cover $W\to U$, and $P$ lifts to $R:V\to W$ by the
finite-cover lifting criterion. The map on geometric fundamental
groups is an isomorphism. Both curves are affine and hyperbolic,
so their free profinite ranks give
$\chi(V)=\chi(W)<0$.

Extend $R$ to the smooth compactifications and let $e\ge1$
be its degree. Since the preimage of the boundary of $W$ is
contained in the boundary of $V$, Riemann--Hurwitz gives
\[
 \chi(V)=e\chi(W)-R_{\rm int}-m,
 \tag{F.46.2}
\]
where $R_{\rm int}\ge0$ is ramification over the interior of
$W$, and $m\ge0$ counts the additional deleted points of $V$
mapping to that interior. Equality of the Euler characteristics
makes $R_{\rm int}+m=(e-1)\chi(W)$, so $e=1$ and both
nonnegative terms vanish. Hence $R$ is an isomorphism and $P$
is finite \'etale. Conversely a finite \'etale map induces an
injective open fundamental-group map. The same argument applies
to $q$. This proves the asserted equivalence and the displayed
open-subgroup isomorphism; no general open homomorphism has
been treated as invertible.

We now compute the restriction on the local coefficients.
Fix an embedding $K\hookrightarrow\mathbb C$. Use the standard
uniformization
\[
 U(\mathbb C)=\overline\Gamma(2)\backslash\mathbb H,
 \qquad z=\lambda(\tau),\qquad
 \lambda(\tau)=16q_0+O(q_0^2),\quad q_0=e^{\pi i\tau}.
 \tag{F.46.3}
\]
Here $\overline\Gamma(2)$ is the image of the principal
level-two subgroup in $\operatorname{PSL}_2(\mathbb R)$.
The normalization, product expansion and identification of the
three punctures in (F.46.3) are recorded together in
\cite[\S1, equation (1.0.2), p.~628]{F-CalegariDimitrovTang};
the first coefficient also follows immediately from
$\lambda=\theta_2^4/\theta_3^4$ and
$\theta_2=2q_0^{1/4}+O(q_0^{9/4})$,
$\theta_3=1+2q_0+O(q_0^4)$.

Using $P$, identify the universal cover of $V(\mathbb C)$
with $\mathbb H$ and its deck group with a finite-index
subgroup $\Lambda\subset\overline\Gamma(2)$. The second
covering map has the form
\[
                  Q(\tau)=\lambda(\alpha\tau)
\]
for an orientation-preserving automorphism
$\alpha\in\operatorname{PSL}_2(\mathbb R)$ of $\mathbb H$.
Both $\Lambda$ and $\alpha\Lambda\alpha^{-1}$ have finite
index in $\overline\Gamma(2)$. Thus $\alpha$ commensurates
$\overline\Gamma(2)$.

For clarity, the required commensurator calculation is elementary.
The fixed points of its parabolic elements are exactly
$\mathbb P^1(\mathbb Q)$: an integral parabolic matrix has a
rational eigenspace, and each rational point is the image
of $\infty$ under an integral projective matrix, and a power
of the resulting parabolic lies in the finite-index level-two
group. Passing to any finite-index subgroup preserves this set,
since a positive power of every parabolic belongs to that
subgroup. Consequently $\alpha$ permutes
$\mathbb P^1(\mathbb Q)$. Its values at $0,1,\infty$ are
three distinct rational projective points; the unique projective
transformation having those three values has rational
coefficients. Hence
$\alpha\in\operatorname{PGL}_2^+(\mathbb Q)$.
Conversely any such rational matrix commensurates the integral
group: the lattices $\mathbb Z^2$ and $\alpha\mathbb Z^2$
are commensurable, and their simultaneous stabilizer has finite
index in each stabilizer. Intersecting with level two preserves
finite index. These are the actual two inclusions used in
the commensurator identification.

Let $c$ be the selected rational cusp of $V$. Since both
$P(c)$ and $Q(c)$ equal $0$, one may change the two universal
covering lifts by elements of $\overline\Gamma(2)$ so that
the chosen lift of $c$ is $\infty$ and
$\alpha(\infty)=\infty$. Therefore
\[
                \alpha\tau=a\tau+b,
                 \qquad a\in\mathbb Q_{>0},\quad b\in\mathbb Q.
\]
For the more general asserted local degrees $d,e$, the width
of the source cusp is $2d$, its analytic parameter is
$q_c=e^{\pi i\tau/d}$, and $a=e/d$. Writing
$\zeta=e^{\pi i b}$, a root of unity, gives
\[
                 P=16q_c^d+O(q_c^{2d}),\qquad
                 Q=16\zeta q_c^e+O(q_c^{2e}).
 \tag{F.46.4}
\]
If $u$ is instead a $K$-rational local parameter at $c$ and
$q_c=\gamma u+O(u^2)$, the leading coefficients are
$A=16\gamma^d$ and $B=16\zeta\gamma^e$. Their indicated
ratio is independent of that change of parameter:
\[
               A^e/B^d=16^{e-d}\zeta^{-d}.
 \tag{F.46.5}
\]
In particular $\xi=\zeta^{-d}$ belongs to $\mu(K)$, since
$A,B\in K^\times$. For the source packet parameter $\beta$
and the common target parameter $\alpha$, the exact local
Kummer maps are
\[
 \alpha=\kappa_K(A)+d\beta=\kappa_K(B)+e\beta,
 \qquad
 (e-d)\bigl(\alpha-\kappa_K(16)\bigr)=\kappa_K(\xi).
 \tag{F.46.6}
\]
This elimination uses actual homomorphisms of the packets; it
does not divide by $e-d$ in the power completion.

We check the torsion needed to extract the conclusion.
The divisor map gives an exact sequence from
$\mathcal O_K^\times/\mu(K)$ to $K^\times/\mu(K)$ and then
onto the group of principal fractional ideals. The first
group is free abelian of finite rank; the last is a subgroup
of the free abelian group on the finite primes, and is free
abelian as well. This exact sequence splits as a sequence of
abelian groups. Thus $K^\times/\mu(K)$ is free abelian, and
so is a chosen complement $F$ in
$K^\times\simeq\mu(K)\times F$. For a basis $(e_i)_{i\in I}$
of $F$, coordinate projections give an injective map
\[
       \widehat F=\varprojlim_n F/nF
              \longrightarrow\prod_{i\in I}\widehat{\mathbb Z}.
\]
Indeed a class with every coordinate zero has zero reduction
in the direct sum $F/nF$ for every $n$. The target is
torsion-free, and hence so is $\widehat F$. Completing the
displayed product decomposition therefore proves that the
torsion subgroup of $\widehat{K^\times}$ is exactly the image
of $\mu(K)$; this description is independent of the chosen
complement. Since $e-d\ne0$ and the right side of (F.46.6)
has finite order, $\alpha-\kappa_K(16)$ is torsion. It follows
that $\alpha\in\kappa_K(16\mu(K))$, as asserted. For the
original degrees $1,2$ the stronger literal coefficient
formula is $\alpha=\kappa_K(A^2/B)$.

For $K=\mathbb Q$ the only possibilities are $16$ and $-16$.
Changes of the universal paths conjugate the group maps; they
do not change either reconstructed rational function or the
local coefficient ratio. If $e=d$, the elimination instead
has zero left side and imposes no such restriction; the
identity example in the statement verifies this exception.

Both rational possibilities occur on the same explicit curve:
\[
 V=\mathbb P^1_{\mathbb Q}\setminus\{0,1,-1,\infty\},
 \qquad P(u)=\frac{4u}{(1+u)^2},\qquad
 Q_+(u)=u^2,\quad Q_-(u)=\frac{u^2}{u^2-1}.
 \tag{F.46.7}
\]
Each is a degree-two map with complete preimage of
$\{0,1,\infty\}$ equal to the displayed four-point boundary,
and all its ramification is there. Thus each restriction to
$V$ is finite \'etale. At $u=0$ their leading terms are
$4u,u^2,-u^2$, respectively. The source tangents $4$ and $-4$
have common images $s_{16}$ for $(P,Q_+)$ and $s_{-16}$ for
$(P,Q_-)$. Taking the image of the finite \'etale map $P$
as $H\subset E$ supplies the actual ambient-subgroup witnesses.
The rational functions, boundary preimages, ramification orders
and coefficients are checked in
\texttt{code/F/finite\_etale\_tangents.py} and its output.

Finally $s_2$ satisfies the unrestricted relation by the explicit
degree-four witness of Theorem~\ref{F-two-Hom-tangents}, while
$2\notin\{16,-16\}$. This proves strictness through the
specified forgetful map. The more general GT-only question
retained in Proposition~\ref{F-open-tangents} has not been
identified with this finite-etale restriction.
\end{proof}

\begin{statement}{F.47}{Proved}{discontinuity on the cyclotomic kernel itself}
\label{st:F.47}
\label{F-cyclotomic-discontinuous-action}
Put $K=\mathbb Q(i)$, $\Gamma=G_K$,
$\Gamma^\circ=G_{K(\mu_\infty)}$, and
$U=\mathbb P^1_K\setminus\{0,1,\infty\}$.
There is an abstract semilinear $\Gamma$-action on the
geometric universal covering $\mathcal U$ of this same
tripod for which the following statements hold.
\begin{enumerate}
\item Its restriction to every open subgroup of
$\Gamma^\circ$ is not functionwise continuous. Its action
on the finite geometric homology quotient
$H_1(U_{\overline K},\mathbb F_3)$ is already discontinuous.
\item No open subgroup of $\Gamma^\circ$ admits a
finite-type quotient $\mathcal U\to W$ which is a
connected pro-\'etale covering and on which the induced
semilinear action is functionwise continuous. In particular,
no open subgroup of $\Gamma$ admits such a quotient.
\item The completed universal covering nevertheless has
exactly one fixed pro-point under this abstract
$\Gamma^\circ$-action, namely the chosen lift of $1/2$.
It is also fixed by the whole abstract $\Gamma$-action.
\end{enumerate}
Thus the inclusion of continuous actions into abstract
actions is strict already on the actual cyclotomic-kernel
domain occurring in (B$'$). The continuous finite-form
descent of Proposition~\ref{F-continuous-action} does not
extend to these actions. The universal fixed-point
existence assertion for arbitrary abstract actions is
neither proved nor refuted by this example.
\end{statement}
\begin{proof}
Choose distinct rational primes $p_j\equiv1\pmod4$ and,
for each $j$, one Gaussian prime $\pi_j$ above $p_j$.
There are infinitely many such rational primes: from any
finite list take $N$ divisible by twice their product;
every odd prime divisor of $N^2+1$ is a new prime
congruent to one modulo four. The Euclidean domain
$\mathbb Z[i]$ gives
$p_j=\pi_j\overline\pi_j$ up to a unit, with the two
prime ideals distinct. Let
\[
 \epsilon_j:\Gamma\longrightarrow\mathbb F_2
\]
be the continuous quadratic character of
$K(\sqrt{\pi_j})/K$.

These characters remain linearly independent after
restriction to $\Gamma^\circ$. To prove this, a nonempty
finite product $a=\prod_j\pi_j^{e_j}$, with
$e_j\in\{0,1\}$, is not a square in $K$ by its prime
valuations. If its character vanished on
$\Gamma^\circ$, then
\[
 K(\sqrt a)\subset K(\mu_\infty)=\mathbb Q(\mu_\infty).
 \tag{F.47.1}
\]
Every finite subextension of the abelian extension
$\mathbb Q(\mu_\infty)/\mathbb Q$ is Galois over
$\mathbb Q$. Consequently complex conjugation would
preserve $K(\sqrt a)$ and its quadratic Kummer line
over $K$. This would give
$a/\overline a\in K^{\times2}$. At any selected
$\pi_j$ its valuation is instead one, a contradiction.
Thus every finite coordinate projection of
\[
 \epsilon=(\epsilon_j)_j:\Gamma^\circ
       \longrightarrow P=\prod_{j\ge1}\mathbb F_2
 \tag{F.47.2}
\]
is onto. Its compact image is closed, so (F.47.2) is
onto. The same coordinate map is also defined on all
of $\Gamma$ and is onto there.

Let $D=\bigoplus_{j\ge1}\mathbb F_2\subset P$.
Choose an abstract $\mathbb F_2$-linear functional
$\lambda:P/D\to\mathbb F_2$ taking the class of
$(1,1,\ldots)$ to one; extension of this nonzero vector
to a basis gives such a functional. With the quotient
by $D$ understood, define
\[
 \psi=\lambda\epsilon:\Gamma\longrightarrow\mathbb F_2.
 \tag{F.47.3}
\]
The restriction of $\psi$ to $\Gamma^\circ$ is nonzero
and its kernel is dense. Indeed it contains the inverse
image of the dense subgroup $D$, and a continuous
surjection of compact groups is an open quotient. The
kernel cannot be open, since a continuous functional on
$P$ vanishing on $D$ vanishes everywhere. The restriction
to every open subgroup $H\subset\Gamma^\circ$ is still
nonzero with dense kernel: its kernel is dense in $H$;
if it contained $H$, the full kernel would be an open,
hence closed, proper dense subgroup. Continuity on any
such $H$ would likewise make the full kernel open.

Choose a geometric pro-point $\widetilde x$ over
$x=1/2$, and let $s$ be the continuous arithmetic point
action of $\Gamma$ on $\mathcal U$. The involution
$a:z\mapsto1-z$ has a unique lift $\widetilde a$ fixing
$\widetilde x$. Existence is the lift of the based
finite-\'etale fibre functor; uniqueness follows because
a deck transformation fixing a point of a geometric
fibre is the identity. It follows that
$\widetilde a^2=1$ and that $\widetilde a$ commutes with
$s(\Gamma)$: conjugation by $s(g)$ gives another lift
of the same $K$-involution fixing $\widetilde x$.
Define the actual action
\[
 r(g)=\widetilde a^{\psi(g)}s(g),\qquad g\in\Gamma.
 \tag{F.47.4}
\]
The two commuting actions and the homomorphism (F.47.3)
prove the group law. Each $r(g)$ is an algebraic
semilinear automorphism, with the original action on
$\overline K$; normalization extends it to the completed
cover. On the coordinate $z$, its value is $z$ or
$1-z$ according to $\psi(g)$. Its stabilizer is therefore
$\ker\psi$, which is not open on any $H$ above.
For a finite quotient detecting the same failure, the
natural action on
$H_1(U_{\overline K},\mathbb F_3)$ is multiplication by
the cyclotomic character: the two independent classes
are the oriented inertia classes at $0$ and $1$.
It is trivial on $\Gamma^\circ$, whereas
$\widetilde a$ interchanges these two classes.
Consequently (F.47.4) acts on this quotient by the
nontrivial interchange to the power $\psi(g)$.

Here is the precise obstruction to a continuously
descended finite quotient. Let $\Omega$ be the covering
function field, let $H\subset\Gamma^\circ$ be open,
and suppose a finite-type quotient field
$M\subset\Omega$ is $r(H)$-stable with functionwise
continuous induced action. For every $f\in M$, the maps
\[
 H\longrightarrow\Omega,\qquad
 g\longmapsto r(g)f,\quad g\longmapsto s(g)f
\]
are continuous with $\Omega$ discrete. They agree on
the dense subgroup $H\cap\ker\psi$, hence everywhere.
Choose $g\in H$ with $\psi(g)=1$. Commutation in
(F.47.4) then gives $\widetilde a f=f$, so
$M\subset\Omega^{\widetilde a}$. If
$\mathcal U\to W$ were a connected pro-\'etale covering
with field $M$, the nontrivial $\widetilde a$ would be
a deck automorphism fixing $\widetilde x$. This is
impossible: unique lifting on finite connected
\'etale levels makes the deck action free on a geometric
fibre, and the inverse limit preserves that property.
This proves assertion 2 for $H$. If an open subgroup
of $\Gamma$ supplied such a quotient, its intersection
with $\Gamma^\circ$ would supply one for an open $H$,
giving the stated further conclusion.

Finally $\widetilde x$ is fixed by (F.47.4). Any pro-point
fixed by $r(\Gamma^\circ)$ is fixed by the natural
$s$-action of the dense subgroup
$\ker(\psi|_{\Gamma^\circ})$. Its stabilizer for the
natural action is closed: this holds at every finite
covering level, and then in their inverse limit.
It is therefore fixed by all $s(\Gamma^\circ)$.
Proposition~\ref{F-fixed-dictionary} says that this fixed
set is exactly $\{\widetilde x\}$. This proves assertion
3 through the actual comparison of the two actions.
Unlike the action in Proposition~\ref{F-discontinuous-action},
the twist here remains nontrivial and discontinuous on
$\Gamma^\circ$ itself, while the fixed-point assertion
continues to hold. No statement of universal existence
for other abstract actions follows from this comparison.
\end{proof}

\clearpage
\bibliographystyle{plain}
\bibliography{LM6_POINTS_AS_SECTIONS_references}
\end{document}
